% YM4M1. Detailed proof outline; companion assembly overview YM4O1.
\documentclass[11pt,letterpaper,reqno]{amsart}
\makeatletter
\NeedsTeXFormat{LaTeX2e}[2020/10/01]

\RequirePackage[T1]{fontenc}
\RequirePackage{lmodern}
\RequirePackage{microtype}
\RequirePackage[letterpaper,margin=1in,headheight=15pt]{geometry}
\RequirePackage{amsmath,amsthm,amssymb,mathtools}
\RequirePackage{enumitem}
\RequirePackage{perpage}
\RequirePackage{xcolor}
\definecolor{manuscriptred}{rgb}{0.6,0.1,0.1}
\definecolor{manuscriptblue}{rgb}{0.1,0.3,0.5}
\definecolor{manuscriptmagenta}{HTML}{60445C}
\RequirePackage[breaklinks=true,colorlinks=true,
  linkcolor=manuscriptred,citecolor=manuscriptblue,urlcolor=manuscriptmagenta]{hyperref}
\RequirePackage{cleveref}
\hypersetup{pdfauthor={}}

\setcounter{tocdepth}{1}
\renewcommand{\thefigure}{\Roman{figure}}
\renewcommand{\thetable}{\Roman{table}}
\MakePerPage{footnote}
\renewcommand{\thefootnote}{%
  \ifnum\value{footnote}>9\arabic{footnote}\else\fnsymbol{footnote}\fi}

\theoremstyle{definition}
\newtheorem{ms@theorem}{Target theorem}
\renewcommand{\thems@theorem}{\Alph{ms@theorem}}
\newtheorem*{ms@theoremstar}{Theorem}
\newtheorem{ms@proposition}{Proposition}[section]
\newtheorem{ms@lemma}[ms@proposition]{Lemma}
\newtheorem{ms@definition}[ms@proposition]{Definition}
\newtheorem{ms@corollary}[ms@proposition]{Corollary}
\newtheorem*{ms@corollarystar}{Corollary}

\newcommand{\ms@majorstart}{%
  \pushQED{\qed}\renewcommand{\qedsymbol}{$\blacksquare\square$}}
\newcommand{\ms@minorstart}{%
  \pushQED{\qed}\renewcommand{\qedsymbol}{$\blacksquare$}}
\newenvironment{theorem}[1][]{\ms@majorstart\ms@theorem[#1]}
  {\popQED\endms@theorem}
\newenvironment{theorem*}[1][]{\ms@majorstart\ms@theoremstar[#1]}
  {\popQED\endms@theoremstar}
\newenvironment{proposition}[1][]{\ms@minorstart\ms@proposition[#1]}
  {\popQED\endms@proposition}
\newenvironment{lemma}[1][]{\ms@minorstart\ms@lemma[#1]}
  {\popQED\endms@lemma}
\newenvironment{definition}[1][]{\ms@minorstart\ms@definition[#1]}
  {\popQED\endms@definition}
\newenvironment{corollary}[1][]{\ms@majorstart\ms@corollary[#1]}
  {\popQED\endms@corollary}
\newenvironment{corollary*}[1][]{\ms@majorstart\ms@corollarystar[#1]}
  {\popQED\endms@corollarystar}

\theoremstyle{remark}
\newtheorem{remark}[ms@proposition]{Remark}
\newtheorem*{remark*}{Remark}

\crefname{ms@theorem}{target theorem}{target theorems}
\Crefname{ms@theorem}{Target theorem}{Target theorems}
\crefname{ms@proposition}{proposition}{propositions}
\Crefname{ms@proposition}{Proposition}{Propositions}
\crefname{ms@lemma}{lemma}{lemmas}
\Crefname{ms@lemma}{Lemma}{Lemmas}
\crefname{ms@definition}{definition}{definitions}
\Crefname{ms@definition}{Definition}{Definitions}
\crefname{ms@corollary}{corollary}{corollaries}
\Crefname{ms@corollary}{Corollary}{Corollaries}
\crefname{remark}{remark}{remarks}
\Crefname{remark}{Remark}{Remarks}

\makeatother

\newcommand{\R}{\mathbb R}
\newcommand{\SU}{\mathrm{SU}}
\newcommand{\cS}{\mathcal S}
\newcommand{\cO}{\mathcal O}
\newcommand{\supp}{\operatorname{supp}}
\newcommand{\Lip}{\operatorname{Lip}}
\newcommand{\tr}{\operatorname{tr}}

\title[Four-dimensional Yang--Mills: a proof outline]{A continuum action-density theory for four-dimensional
\texorpdfstring{$\mathrm{SU}(2)$}{SU(2)} Yang--Mills: a proof outline}
\date{}
\hypersetup{pdftitle={A continuum action-density theory for four-dimensional SU(2) Yang-Mills: a proof outline}}

\begin{document}
\begin{abstract}
We give a detailed proof outline for constructing the scalar action
density of four-dimensional $\SU(2)$ Yang--Mills theory from Wilson
lattice measures. The construction takes the infinite-volume limit
before removing the lattice spacing. Its target is a translation
component whose moments are Euclidean invariant, reflection positive
and clustering, with factorial bounds on Schwartz tests flat at
coincidences. The connected two-point function should have size
$d^{-8}(\log(1/d))^{-2}$ for fixed scaled bumps. We develop the
renormalisation, large-field and source arguments, the passage from
cubic symmetry to rotations, and the selection of a clustering phase.
Uniform construction estimates that still need complete verification
are identified at their point of use.
\end{abstract}
\maketitle
\tableofcontents

\section{The construction and its target}\label{sec:main}

For an oriented nearest-neighbour bond $b$ of a four-dimensional cubic
lattice, let $U_b\in\SU(2)$, with $U_{-b}=U_b^{-1}$. For an elementary
plaquette $p$, let $U_p$ be the ordered product around its boundary and set
\begin{equation}\label{eq:observable}
 a_p(U)=1-\tfrac12\operatorname{Re}\tr U_p,
 \qquad \cO_a(f)=\sum_p f(x(p))a_p(U).
\end{equation}
The sum takes each unoriented plaquette once, and $x(p)$ is its centre.
On a periodic lattice of spacing $a$ and side $\ell$, Wilson's measure is
\begin{equation}\label{eq:wilson}
 d\mu_{a,\ell,\beta}(U)
 =Z_{a,\ell,\beta}^{-1}
   \exp\!\left(-\beta\sum_p a_p(U)\right)\prod_b dU_b,
\end{equation}
where the product is over one orientation of each bond, $dU_b$ is
normalised Haar measure, and $\beta=g_0^{-2}$ defines the bare coupling
$g_0$. The continuum observable sought here is the
scalar action density, formally proportional to
$\tr(F_{\mu\nu}F_{\mu\nu})$. Its normalisation is fixed by
\eqref{eq:observable}; there is no additional power of $a$ in that sum.

We first take $\ell\to\infty$ at fixed $(a,\beta)$, and only afterwards
let $a\to0$ and $g_0\to0$. A small reference parameter $g$ determines
the terminal inverse-coupling window in \eqref{eq:first-passage}; exact
tuning to $g$ is not required. Subsequence selection is allowed. If
$\nu_a$ is the fixed-spacing infinite-volume reference state, we form
moments of $\cO_a(f)-\nu_a(\cO_a(f))$, without requiring a finite
uncentred one-point limit. A continuum limit can mix phases. An ergodic
component means one phase in its decomposition under translations.
It retains this reference-state subtraction and can therefore have a
nonzero mean. Clustering is sought in that component.

The Euclidean test space in the theorem must be specified. Let
\begin{equation}\label{eq:flat-space}
 \cS_{\ne}(\R^{4n})
 =\{\varphi\in\cS(\R^{4n}):
   \partial^\alpha\varphi(x)=0\text{ whenever }x_i=x_j
   \text{ for some }i\ne j,\text{ for every }\alpha\}.
\end{equation}
Thus the test function and every derivative vanish at coincidences.
For $n=1$ this is the usual Schwartz space, and $\cS_{\ne}(\R^0)=\mathbb C$.
Write
\[
 |\varphi|_k=
 \sup_{x,\,|\alpha|\le k}(1+|x|^2)^{k/2}|\partial^\alpha\varphi(x)|.
\]
For a separated tuple $X=(f_1,\ldots,f_n)$, the functions are smooth,
compactly supported, and have pairwise disjoint supports at positive
distance. We write $S(X)=S_n(f_1\otimes\cdots\otimes f_n)$.
Let $\theta(x_1,x_2,x_3,x_4)=(x_1,x_2,x_3,-x_4)$ and
$(\Theta\varphi)(x)=\overline{\varphi(\theta x_1,\ldots,\theta x_n)}$.

The following target organizes the proof outline. The sections below
develop the construction and the estimates leading to it; the outline
does not claim that every analytic and geometric verification is
complete.

\begin{theorem}[Euclidean action-density construction]\label{thm:main}
For every sufficiently small reference parameter $g>0$, there are bare
couplings $g_{0,j}\to0$ and spacings $a_j\to0$ for which the centred,
volume-first action-density correlations have a common subsequential
limit on all separated tuples. This limit has an ergodic component with
moments $S_n$ on $\cS_{\ne}(\R^{4n})$, $S_0=1$, with the following
properties.
\begin{enumerate}[label=(\roman*)]
\item\label{it:invariance}
The moments are invariant under translations and $\mathrm O(4)$,
symmetric under permutation of the arguments, and real on real tests.
Here Euclidean motions act on every argument by the same pullback.
The first moment is $S_1(f)=b\int f$ for a finite real constant $b$.
\item\label{it:positivity}
They are reflection positive:
\begin{equation}\label{eq:rp}
 \sum_{q,q'\ge0}S_{q+q'}(\Theta\varphi_q\otimes\varphi_{q'})\ge0
\end{equation}
for every finite sequence $\varphi_q\in\cS_{\ne}(\R^{4q})$ supported
where all time coordinates are positive, with $\varphi_0\in\mathbb C$.
\item\label{it:growth}
There are $A,B>0$, an integer $k\ge1$, and $p<\infty$, independent of
$n$, such that
\begin{equation}\label{eq:growth}
 |S_n(\varphi)|\le A B^n(n!)^p|\varphi|_{kn}
 \qquad(n\ge1).
\end{equation}
\item\label{it:cluster}
For separated compactly supported tuples $X,Y$, there are
$C_{X,Y},r_{X,Y}<\infty$ such that
\begin{equation}\label{eq:cluster}
 |S(X,\tau_vY)-S(X)S(Y)|\le C_{X,Y}|v|^{-2}
 \qquad(|v|\ge r_{X,Y}),
\end{equation}
where $(\tau_v f)(x)=f(x-v)$.
\item\label{it:nontrivial}
Fix a nonnegative smooth function $\chi$ of integral one, supported in
$B(0,1/8)$, and put
\[
 f_d(x)=d^{-4}\chi(x/d),\qquad
 h_d(x)=d^{-4}\chi((x-3de_1)/d).
\]
There are $c,C,d_0>0$, depending only on $g$ and $\chi$, such that
\begin{equation}\label{eq:nontrivial}
 \frac{c\,d^{-8}}{(\log(1/d))^2}
 \le S_2(f_d\otimes h_d)-S_1(f_d)S_1(h_d)
 \le\frac{C\,d^{-8}}{(\log(1/d))^2}
 \qquad(0<d\le d_0).
\end{equation}
The smallness threshold for $g$ may depend on the fixed choice of $\chi$;
all other constants may depend on $g$ and, in this clause, on $\chi$.
\end{enumerate}
\end{theorem}

We use the diagonal-flat test-space convention of
\cite[\S\S2--3, pp.~85--89]{OS1973} for the
Osterwalder--Schrader properties listed above.
The quantitative growth and reconstruction questions are distinguished in
\cite[\S IV.1, pp.~287--288]{OS1975}; here only the Euclidean properties
listed in the target are sought. The lower bound in
\eqref{eq:nontrivial} says that the scalar field is not a constant.
The statement concerns this observable sector; it asserts neither a
mass gap nor uniqueness of the continuum limit. It does not specify
contact terms, that is, the values of distributions supported at
coincident points. We do not include a reconstruction theorem in the
target statement.

The construction begins with one common sequence of effective action
coefficients. In four dimensions the quadratic curvature term is
marginal: its small changes accumulate over arbitrarily many scales.
It must therefore be extracted at each step and put into the running
inverse coupling. The other local terms are controlled on analytic
domains with independent field and current variables. Agreement on
nonwrapping regions makes this one coefficient sequence serve all
volumes. Stopping at the first passage into a fixed small-coupling
window then gives the lattice spacings used in the limit.

Large fields cannot be removed by imposing a global small-field
restriction. The integral is decomposed by its actual failed tests,
and every retained region carries the fields and operations stored
there. A component that has stayed small and quiet for sufficiently
many steps is prepared, integrated and replaced by local terms. This
operation prevents a rare old event from charging every subsequent
scale. Its proof has to control whole histories with an exterior
supremum inside the component integral; a bound for each individual
history would leave the history count unpaid. The effective density
and its induction are developed in \cref{sec:rg,sec:large-fields}.

Sources are inserted by freezing those real operations and changing
the weight to which they are applied. Local scalar source energies
are subtracted before taking absolute values. The same subtraction
can then be used in different volumes and Gibbs states, even when
the uncentred one-point functions diverge. This gives bounded centred
moments, estimates at each fixed order, and a graded estimate that
keeps one factor for each active insertion. Two observation scales
separate the leading free-curvature covariance from the local and
live-component errors; the latter must be small relative to the
leading $g(d)^4$ term. These are the source and two-point parts of
\cref{sec:sources,sec:two-point}.

The limiting construction initially retains only cubic symmetry.
Coordinate and diagonal reflection forms supply Hilbert spaces and
bounded transfer inequalities. Passing those bounded inequalities to
the continuum gives the sharp coordinate cones. A local insertion
bound independent of the other fields makes planar point insertions
into bounded forms. Their ordered products are holomorphic, with an
exponent in the tube margin growing only linearly in their number.
Once the ordered angular charts have been glued, the margin bound
gives only finitely many Fourier coefficients in the rotation angle. Reflection positivity gives convexity of the
highest frequency, and a polygon zero-row argument turns it into
linearity. A short-distance upper bound and a second positivity sign
at a diagonal mirror settle the two-point base case. Tracking the
coupling dependence of source stability supplies a sufficient upper
bound without the relative Gaussian comparison. Ordered tubes, an
angular differential equation and slab translations give the
continuation through different orderings in \cref{sec:rotations}.

The reference continuum correlation hierarchy may be a mixture of phases. At each
fixed lattice spacing its translation-ergodic decomposition is again
Gibbs for the same Wilson specification. A uniform estimate of the
approach to the invariant plateau is then needed to pass this
decomposition to the continuum and retain clustering. The proof uses
the transfer spectrum in four coordinate pictures and is developed
in \cref{sec:phases}. Finally, the graded correlation bounds extend
the separated distributions to tests flat at coincidences. The
extension, restoration of the finite component mean, and order of
the common construction are described in \cref{sec:extension}.

\section{Renormalisation and the choice of bare couplings}\label{sec:rg}

A renormalisation step replaces a field by its averages on blocks of
side $L$. Its purpose is to express the same integral in terms of fewer
variables, while retaining quantitative information about the
fluctuations that have been integrated out. In four dimensions this
requires a running coefficient of the action. A bound on all the other
terms, without separating that coefficient, does not remain uniform in
the number of steps.

Choose the remainder exponent $\kappa_0$ sufficiently large for the
scale sums below; in particular, $\kappa_0\ge7$.
Fix $L$ odd and sufficiently large and put $a=a_K=L^{-K}$. The auxiliary
cube sizes are chosen as powers of $L$. The torus estimates use sides
$\ell=2L^m$, with $m$ an integer large enough to satisfy the volume
floor in \cref{eq:flow-volume-floor}. Below, ``tori'' refers to this
family, except for the adapted period lattices used for diagonal
mirrors in \cref{sec:limits}. After $k$ steps
the lattice is
$\mathbb T^{(k)}=(L^ka\mathbb Z/\ell\mathbb Z)^4$, with spacing $L^ka$.
When lengths are measured in units of this lattice, the bare spacing is
$\eta_k=L^{-k}$. We write
$A(U)=\sum_p a_p(U)$ and $x_k=g_k^{-2}$.
For a scalar weight $v$, also write
$A(v,U)=\sum_p v(x(p))a_p(U)$. Supremum norms on matrix-valued
fields use a fixed matrix operator norm; on Lie-algebra fields they
use a fixed Euclidean norm. The symbol $\partial U$ denotes a
plaquette product. Covariant differences $\nabla_U$ and the
infinitesimal gauge action $D_U$ use the spacing of the lattice on
which they act; distances in a level-$k$ estimate use that level's
units unless physical units are specified. Thus $D_U$ maps site
gauge parameters to bond potentials, and $D_U^*$ is its adjoint.
Write $\mathrm d_U$ for the covariant exterior derivative from bond
potentials to plaquette two-forms, with the same spacing convention.
Write $\Delta_U=D_U^*D_U$ for the covariant Laplacian on site
functions. The symbol $\Delta$ denotes the positive flat lattice
Laplacian, acting componentwise on forms.
The construction uses two auxiliary cube sizes: $M_1$ for local
propagators and $M$ for local interactions. The constants are chosen
in the order $L$, $M_1$, the decay rate $\kappa$, then $M$ large,
the interaction bound $E_0$, and finally a small coupling ceiling
$\gamma$. In particular, the small coupling is chosen after the
constants that multiply it.
Write $W_4$ for the group of signed permutations of the four
coordinate axes, acting on the lattice and its fields.

Complex fields take values in $\mathrm{SL}(2,\mathbb C)$. For them,
$A$ and its weighted version use the holomorphic plaquette term
$a_p(\mathbf U)=1-\tfrac12\tr\mathbf U_p$.
This agrees with \eqref{eq:observable} on $\SU(2)$, whose elements
have real trace. All actions evaluated on complex fields below use
this continuation.
Adjoints are taken with respect to the fixed real-field inner products.
At complex fields the starred operator formulas denote their holomorphic
continuations, rather than Hermitian adjoints taken at the complex
argument.

The small-field construction follows Ba\l aban's
\cite{Balaban1987}. The work is to retain its locality and analytic
dependence uniformly in the cutoff and torus, and then to include the
large-field terms. We describe the small-field integral first. The
operation on the complementary regions and the exact summation over
their histories are treated in \cref{sec:large-fields}. The two parts
share the same running coupling and must close one induction.

The argument separates exact integral identities from the estimates
used to iterate them. Technical verifications still needed for the
construction are identified where they enter. Exact preservation of
an integral does not itself give a uniform bound on its local expansion.

\subsection{Averaging and gauge fixing}

Group-valued fields require an average that respects multiplication by
gauge transformations. For a family of nearby group elements $V_i$,
one first forms a mean $V$ determined locally by
$\sum_i\log(V_iV^{-1})=0$. This mean is analytic near a constant
family, commutes with left and right multiplication, and has ordinary
arithmetic averaging as its linear part in the Lie algebra. For a site
$y$ at the centre of a block $B(y)$, let $U(y,x)$ be the mean of the
transports from $y$ to $x\in B(y)$ along shortest paths that traverse
the coordinate directions in all possible orders. Write
$U(x,y)=U(y,x)^{-1}$.

Let $c$ be a coarse bond, let $[x,x']$ be its parallel copy starting at
$x\in B(c_-)$, and let $U[c]$ be straight transport along $c$.
The block average is
\begin{equation}\label{eq:rg-average}
 \mathcal M(U)(c)=
 \exp\!\left\{L^{-4}\sum_{x\in B(c_-)}
 \log\!\left[
 U(c_-,x)U([x,x'])U(x',c_+)U[c]^{-1}\right]\right\}U[c].
\end{equation}
The logarithms in this formula concern closed transports based at
$c_-$. On real fields with plaquettes within $cL^{-2}$ of the identity
in the two adjacent blocks, the contours are close enough to the
identity for this formula to be analytic. A complex field requires the
corresponding contour conditions as well: a small bound on its
plaquettes alone is insufficient, since conjugation by a complex
gauge transformation need not preserve the matrix norm. Outside the analytic domain the construction
uses covariant Borel extensions of the average. Such an extension is
needed because the exact integral is over all configurations.

The centred blocks and the mean over all coordinate orders make
$\mathcal M$ covariant under the hypercubic group. It is gauge
covariant and local in the two blocks meeting $c$.
The transformation $\mathbf T$ is the image measure under this map:
\[
 \int F\,d(\mathbf T\rho)
 =\int F(\mathcal M(U))\rho(U)\,dU,
 \qquad \int\mathbf T\rho=\int\rho.
\]
This equality is exact. It concerns the averaging transformation
$\mathbf T$; the observables accompanying the subsequent large-field
operation require the insertion identities of
\cref{sec:large-fields}.

Uniform control of repeated averages uses a gauge correction at every
step. For a regular real $U$ and a small complex potential $A$, the
expansion has the form
\begin{equation}\label{eq:rg-straightening}
 \mathcal M^k(e^{i\eta_kA}U)
 =\left(e^{i[Q_k^\sharp(U)A+C_k(U,A)]}
                   \mathcal M^kU\right)^{v_k},
 \qquad |C_k(U,A)|\le C|A|^2,
 \quad |v_k-1|\le C|A|.
\end{equation}
Here the block gauge factor produced by one average is removed before
the next is taken. The linear operator $Q_k^\sharp$ is the derivative
of this corrected average, rather than the plain derivative of
$\mathcal M^k$. Its leading term is the straight covariant average;
the difference has relative size $O(\varepsilon)$ at curvature
$\varepsilon$. The two agree at the trivial field.

This distinction accounts for uniformity in $k$. If $A^{(n)}$ is the
potential after $n$ corrected averages, set
$t_n=\sup|A^{(n)}|/(L^n\eta_k)$, and let $\varpi_n$ bound the
curvature after those averages. The corrected-average estimate is
\[
 t_{n+1}\le t_n(1+C\varpi_n+CL^n\eta_kt_n),
 \qquad \sum_n\varpi_n\le C\varepsilon,
 \qquad \sum_{n<k}L^n\eta_k\le 2L^{-1}.
\]
The error sums are geometric. By comparison, multiplying estimates for
the uncorrected linear averages produces a factor larger than $2L$
per step instead of the spacing ratio $L$. It does not give the
required uniform estimate. The corrected operator also satisfies the
gauge identity
$Q_k^\sharp(U)D_U\lambda
=D_{\mathcal M^kU}Q'\lambda$.
This is the identity used in the gauge projection below.
The site operator $Q'$ is the linearisation of the iterated
gauge-transformation average. At one step it averages a site
parameter after adjoint transport to the block centre:
$\lambda\mapsto L^{-4}\sum_{x\in B(y)}
\operatorname{Ad}(U(y,x))\lambda(x)$. Its iterate uses the
corresponding fields at successive levels.

The symmetric average here also has to be compared with the
corner-block average in Ba\l aban's background-field propagator
construction \cite{Balaban1985}.
The comparison is between their gauge-corrected linear parts.
After aligning the blocks by a half-block translation, it has the
form
$Q_k^\sharp=\operatorname{Ad}(t)(Q_k^{\mathrm B}+\mathsf D_k)$,
where $t$ transports the block base point and $\mathsf D_k$ has
relative size $C\varepsilon$ in row and column norms, uniformly
in $k$. At the trivial field these linear parts coincide.
Thus the transfer of propagator estimates uses a controlled
comparison of the actual averages.

Gauge fixing removes zero directions of the Hessian. On a block there
are $3(L^4-1)$ such directions for $\SU(2)$, corresponding to gauge
transformations equal to the identity at the coarse site.
Insert a factor $\mathfrak G_k$ whose average over these gauge
transformations is one. In terms of the mean contour variables
$V(y,x)$, its local factors have the form
\[
 \mathfrak z(x_k,\varepsilon_k)^{-1}
 \mathbf1_{\{|V(y,x)-1|<\varepsilon_k\}}
 \exp\!\left[-x_k\left(1-\tfrac12
                  \operatorname{Re}\tr V(y,x)\right)\right].
\]
The exponential term is a soft gauge fixing. With the fluctuation
scaled as $g_kB$, its quadratic part is of order one and is strictly
positive on the residual gauge directions. Its indicators equal one
on the smaller fluctuation domain used in the integral. The ratio of
these two domains is fixed after $L$. The normalising factors
$\mathfrak z$ enter the vacuum energy, and hence must be included in
the exact bookkeeping of constants.

Axial gauges choose representatives of gauge orbits. The estimates use
a Landau slice, with gauge transformations constrained by their
averages. For a regular reference field $U_0$, nearby regular fields
can be written, after such a gauge transformation, as
$e^{i\eta_kA}U_0$ with
$R_{\mathrm{comb}}D_{U_0}^*A=0$ and
$|A|+|\nabla_{U_0}A|\le C(\varepsilon+\alpha_1)$.
Here $R_{\mathrm{comb}}$ is the slice projection for the corner-comb
site averages described below, and
$\alpha_1$ bounds the change in the prescribed averages. On a cube of
side $O(M)$, curvature and current bounds of size $\varepsilon$ give
a local gauge with
$|A|+|\nabla A|\le CM\varepsilon$ when $M\varepsilon$ is small.
These local charts are the ones used to compare the background
operators with their values at the trivial field.

The plain derivative will also be needed. Put
\[
 \mathcal Q_k(U;A)=\frac1i\log\big[
       \mathcal M^k(e^{i\eta_kA}U)(\mathcal M^kU)^{-1}\big],
 \qquad Q_k(U)=\partial_A\mathcal Q_k(U;0).
\]
If $\dot v_k$ is the infinitesimal node factor in
\eqref{eq:rg-straightening}, then, with the gauge-action convention
used there,
\begin{equation}\label{eq:mrg-plain-corrected}
 Q_k=Q_k^\sharp-D\dot v_k,
 \qquad
 Q_kD_U\lambda=D_{\mathcal M^kU}
                 (\lambda|_{\mathbb T^{(k)}}).
\end{equation}
The corresponding right-hand side for $Q_k^\sharp$ contains the
average $Q'\lambda$. A gauge parameter vanishing at the coarse
sites therefore preserves the plain constraint, whereas a gauge
parameter of zero average preserves the corrected constraint. These
are different subspaces, already at the trivial field. The
propagators use the second one. The fibre coordinates below solve the
plain constraint and include its compensating node factor.

The local mean and the continuation of the average have different
purposes. Near a constant family, writing $V_i=e^{X_i}V_0$ and
$V=e^XV_0$ reduces the mean equation to
\[
 \sum_i\log(e^{X_i}e^{-X})=0.
\]
At $X_i=X=0$ its derivative in $X$ is minus the number of members
times the identity. The local implicit solution thus has linear
part $X=\frac1{\#\{i\}}\sum_iX_i$. Conjugation commutes with
the equation, and uniqueness gives local covariance. This argument
supplies no logarithm at every group element or unique mean for an
arbitrary family. In particular, the principal matrix logarithm at
$-1$ is not an $\mathfrak{su}(2)$ logarithm. Borel rules away from
the analytic neighbourhood are fixed before the expansion.

The recurrence for corrected averages gives an elementary uniform
bound once its one-step estimate is available. Fix a small upper
bound $t_\star$ for the initial potential. Up to a first index with
$t_n>2t_\star$, the product of the recurrence factors is at most
\[
 \exp\!\left(C\sum_n\varpi_n+
           2Ct_\star\sum_{n<k}L^n\eta_k\right)
 \le\exp(C\varepsilon+4Ct_\star/L).
\]
Choose $\varepsilon,t_\star$ so this is less than two. There can
then be no such first index. The uniform one-step estimate, its
quadratic error, and its row and column bounds remain the averaging
input; the displayed product explains how that input prevents an
extra factor depending on $k$.

For later use, a level-$j$ activity is holomorphic in an independent
pair $\mathbf p=(\mathbf U,\mathbf J)$ represented by a splitting
$\mathbf U=e^{i\eta_jA'}U$, with $U$ real. On each relevant cube
the analytic domain requires
\begin{equation}\label{eq:mrg-pair-domain}
 \begin{gathered}
 |\partial U-1|,|\partial\mathbf U-1|
        <\alpha_{0,j}\eta_j^2,\qquad
 \max\{|A'|,|\nabla_UA'|\}<\alpha_{1,j},\qquad
 |\mathbf J|<\alpha_{0,j},\\
 CM\alpha_{0,j}<\alpha_{1,j},\qquad
 \alpha_{1,j}^2\le C\alpha_{0,j}.
 \end{gathered}
\end{equation}
It also requires the local real gauges above and regularity of the
backgrounds formed from averages on the subdomains used by the
construction. These latter conditions are part of the domain.
Two allowed splittings of one pair must give the same value to every
activity. No saturation under arbitrary complex gauge
transformations is asserted.

Expanding the plaquette of $e^{i\eta_jA'}U$, the two plaquette
conditions give
\[
 |\mathrm d_UA'|\le C(\alpha_{0,j}+\alpha_{1,j}^2).
\]
The quadratic term has the order of the allowed curvature and scales
accordingly when the background is restricted to a later level.
Response and interpolation estimates must preserve the entire
domain, including the conditions on subdomain backgrounds.

The block factor removed at a single step is explicit.
If $W=\mathcal M^nU$ and
$W'=e^{iA^{(n)}}W$ is the already straightened field,
the factor at $y$ is
\begin{equation}\label{eq:xrg-block-factor}
 v^{(n)}(y)=\exp\!\left(
 L^{-4}\sum_{x\in B(y)}
       \log[W'(y,x)W(y,x)^{-1}]\right).
\end{equation}
It is formed from $W'$, not from the plain iterate of
the perturbed field. Removing it before the next
average is what leaves the spacing factor $L$ in
the row norm. The product of the transported block
factors is $v_k$ in \eqref{eq:rg-straightening}.
If the perturbed field has the same contour
variables as $U$ at each preceding level, it is
axial relative to $U$, and these corrections are
trivial.

For a gauge function $u$, the corresponding block
average at the reference field $W$ is
\begin{equation}\label{eq:xrg-gauge-average}
 (\overline R_Wu)(y)=u(y)
  \exp\!\left(L^{-4}\sum_{x\in B(y)}
       \log[u(y)^{-1}\operatorname{Ad}(W(y,x))u(x)]\right).
\end{equation}
At first order this gives $Q'$. The nonlinear
condition $\overline R_U^{(k)}u=1$ describes the
gauge transformations that leave the straightened
constraint unchanged at a relatively axial field.
Away from that field the condition depends on the
field itself. It should not be treated as a fixed
subgroup. The transformations equal to one at the
coarse sites do form a subgroup; the node factor
relates these two descriptions of an orbit.

This distinction also explains the choice of soft
gauge fixing. Delta functions on single tree
transports can be solved by ordering the tree.
Delta functions on the contour means in
\eqref{eq:rg-average} instead couple all the paths
and have no such triangular solution. The soft
factor retains hypercubic covariance and its
normalisation is still exact: at fixed $V$,
$V^u(y,x)=V(y,x)u(x)^{-1}$, and Haar invariance
reduces the integral of each local factor to
$\mathfrak z/\mathfrak z=1$.

Axial representatives on a sequence of layers
are pinned outside its nodes. At a seam a finer
node belongs to a block of the coarser layer, so
its transport is pinned as well. An orbit of the
group equal to one at all nodes then meets the
axial gauge once. The subsequent Landau estimate
uses transformations with trivial averages and
therefore need not keep those node values fixed;
their coarse effect is compensated in the datum.
The site averages defining $R_{\mathrm{comb}}$ are formed by
transport along combs starting at block corners. The comparison
of this slice with the Landau slice for the original averages must
be checked uniformly in the level. The two projections cannot
be identified merely because both remove infinitesimal gauge directions.

\subsection{The constrained background and its response}

The background is a regular critical configuration with prescribed
averages. Only uniqueness among regular critical configurations is
needed. No assertion about all stationary points, or all absolute
minima in the fibre, enters the construction. We first specify the
problem on a sequence of regions; the problem on one coarse lattice
is its special case in which every region is the torus.

\emph{Determining regions and the variational problem.}
Fix the last level $k$ and use its units. The bare spacing is
$\eta_k=L^{-k}$, and the spacing at level $j\le k$ is
$\xi_j=L^{j-k}$. Let
\[
 \Omega_0\supset\Omega_1\supset\cdots\supset\Omega_k,
 \qquad \Omega_{k+1}=\varnothing,
\]
where $\Omega_j$ is a union of cubes of side $M_1\xi_j$ and
\[
 \operatorname{dist}_\infty(\Omega_j^c,\Omega_{j+1})
       >\rho_{\rm geom}\xi_j,\qquad \rho_{\rm geom}=R_1M_1.
\]
The parameters $M_1,R_1$ are large powers of the fixed odd
integer $L$. The regions may be disconnected and their shapes
are unrestricted. Their separation, rather than regularity of
their boundaries, is the geometric input.

The layer of level $j$ is $\Omega_j\setminus\Omega_{j+1}$.
Its tops are the sites of $\mathbb T^{(j)}$ in that layer. A
top $y$ has scale $\xi_y=\xi_j$ and pillar $B^j(y)$.
The pillars partition the determining region. At a seam, a site
may belong to a finer lattice while lying in a pillar of a
coarser level; following its successive block centres leads to
its top. This ascent is what fixes the gauge convention at seams.
For a bare bond or plaquette $e$, write $\mathsf L(e)$ for the
largest scale attached to its corners.

Let $\Gamma_j$ consist of the bonds of $\mathbb T^{(j)}$
with both endpoints in $\Omega_0$, at least one endpoint in
$\Omega_j$, and neither endpoint in $\Omega_{j+1}$. Put
\[
 \mathbf B=\bigcup_{j=0}^k\Gamma_j,\qquad
 M_{\mathbf B}(U)=(\mathcal M^jU|_{\Gamma_j})_j.
\]
An interior row joins two tops in one layer. An outward row
has an endpoint in the adjacent finer layer; the bonds inside
that endpoint's block are retained at the finer level. There is
no row on an inward bond entering $\Omega_{j+1}$. This last
convention is important when the next step imposes its block
axial gauge: an extra inward constraint can prevent that gauge
from being reached. In the strip construction, where a different
set of nodes is retained, axial gauge is imposed outside those
nodes instead.

For datum $\mathbf V=(V_j|_{\Gamma_j})_j$, the fibre and action
are
\[
 C(\mathbf B,\mathbf V)=\{U:M_{\mathbf B}(U)=\mathbf V\},
 \qquad A_{\Omega_0}(U)=\sum_{p\subset\Omega_0}a_p(U).
\]
Criticality refers to curves in this fibre. The action contains
the inner plaquettes only. Its current includes their first
variation at the boundary bonds; it must not be replaced there
by the current of an arbitrarily continued whole-torus field.
The residual gauge group is
\[
 \mathcal G^{\rm res}(\mathbf B)
       =\{u:u=1\text{ at every node of }\mathbf B\}.
\]
The datum is regular when its test-loop products are within
$\varepsilon$ of one. Inside a layer the loops are its
plaquettes; across a seam, the finer values are averaged before
forming the loop. A regular configuration satisfies
\[
 |\partial U(p)-1|<\varepsilon_{\rm reg}\eta_k^2
                                      \mathsf L(p)^{-2},
 \qquad |J(U)(b)|<\varepsilon_{\rm reg}\mathsf L(b)^{-3}.
\]
Thus curvature has dimension two and current dimension three
in the units of the layer where they are read.

The background construction gives one orbit of regular critical
configurations when the datum is sufficiently regular. We write
it $U_{\mathbf B}(\mathbf V)$ or $U_k(V)$ on a single lattice.
The regular radius is fixed above $C\varepsilon$ and below a
small absolute ceiling. The torus must have room for the double
blocks and the enlarged local cubes used below; in particular
these cubes are shorter than half its side. Increasing the
volume introduces no new constant.

On a cube of side $O(M)$ the local regularity statement is
\begin{equation}\label{eq:rg-background-regularity}
 \begin{gathered}
  |\partial U_k-1|\le C\varepsilon\eta_k^2,
  \qquad |J(U_k)|\le C\varepsilon,
  \qquad U_k=e^{i\eta_kA},\\
  |A|+|\nabla A|+
          \|\nabla A\|_{C^{\alpha_{\mathrm H}}}
       \le CM\varepsilon,
       \qquad 0<\alpha_{\mathrm H}<1.
 \end{gathered}
\end{equation}
Here $M\varepsilon$ is small, and the gauge is chosen on a
slightly larger cube. On a layer of scale $\xi_j$ the
quantities bounded by $CM\varepsilon$ are
$\xi_j|A|$, $\xi_j^2|\nabla A|$ and
$\xi_j^{2+\alpha_{\mathrm H}}\|\nabla A\|_{C^{\alpha_{\mathrm H}}}$.
Thus their unweighted bounds contain the corresponding
inverse powers of $\xi_j$. The H\"older
exponent is fixed; constants need not remain bounded as it tends
to one. The linear and nonlinear constructions that supply these
estimates are developed below.

\emph{The distance used by the inverse estimates.}
Measure length in the units of the layer it traverses:
\[
 d_{\mathbf B}(x,x')
 =\inf_{\Gamma\ni x,x'}
 \left[\sum_{j\ge1}\xi_j^{-1}
      \operatorname{length}(\Gamma\cap
                    (\Omega_j\setminus\Omega_{j+1}))
       +\eta_k^{-1}\operatorname{length}(\Gamma\setminus\Omega_1)
 \right].
\]
The infimum is over connected paths in the torus. In particular
the distance outside $\Omega_0$ uses the bare unit; leaving the
determining region does not provide a cheaper route. The first
member $\Omega_0$ does not otherwise enter the definition.

If $y,y'$ have levels $j,j'$, each intervening layer costs
the gap:
\[
 d_{\mathbf B}(y,y')
        \ge\rho_{\rm geom}(|j-j'|-1)_+.
\]
This controls the entropy of summing over layers. Within a
layer, sums of an exponential are bounded by the corresponding
lattice integral. On moving one level outward, at most a factor
$L^4$ is lost in the count of tops. Choose the gap, for every
fixed rate $\tau>0$ and scale power $q$ that will be used, so
that $\tau\rho_{\rm geom}\ge C(1+q)\log L$. Then
\[
 \sup_y\sum_{y'}e^{-\tau d_{\mathbf B}(y,y')}\le C_{\tau,L},
 \qquad
 \left(\frac{\xi_{y'}}{\xi_y}\right)^q
                  \le L^qe^{\tau d_{\mathbf B}(y,y')}.
\]
A fixed fraction of a decay rate can therefore absorb scale
ratios without producing a factor depending on the number of
regions. This is the summation used in the Green functions,
the response, and the later replacement of global backgrounds
by local ones.

The raw distance is insufficient as a weight when second
differences are taken. Use a regularized distance $f_{y'}$ with
\[
 |f_{y'}-d_{\mathbf B}(\cdot,y')|\le C,\quad
 |\nabla f_{y'}|\le C\xi^{-1},\quad
 |\nabla^2f_{y'}|\le C\xi^{-2}
\]
in the local scale. Pairing the opposite directions in the
conjugated Laplacian gives a perturbation $O(\theta)$ for
the weight $e^{\theta f_{y'}}$. Estimating its directed bonds
separately would instead introduce the inadmissible ratio
$\xi/\eta_k$. The mass and the weight must be scaled together.

\emph{The chart and the row datum.}
Fix a real regular reference $U_0$. Put
$U=e^{i\eta_kA'}U_0$, and let $F(A')=A_{\Omega_0}(U)$.
The straightened average is the nonlinear function
$\mathcal Q_k^\sharp(A')$ whose site factors were removed in
\eqref{eq:rg-straightening}. Its derivative will be denoted
$K_{A'}$; at zero it consists of one local row for each bond
of $\mathbf B$. Each row sees only the double block beneath
that bond. The plain logarithm of the average has an additional
coarse gauge term, which is present even at the trivial field.
It cannot be discarded as a curvature error.

At an interior row the reduced datum equals $B$. At an outward
row whose outer endpoint is $z$, the node compensation in the
finer block gives
\begin{equation}\label{eq:mrg-reduced-datum}
 e^{ib(B)_c}=e^{iB_c}\operatorname{Ad}(V_{0,c})v_z(B).
\end{equation}
The expression for the other orientation is obtained by reversing
the bond. The factor $v_z$ reads only the interior rows of its
block. Consequently $B\mapsto b(B)$ is triangular, holomorphic
and locally invertible, with $\sup|b(B)|\le C\sup|B|$.
One first estimates the response to $b$ and then differentiates
through this map. The seam term in that chain rule is part of
the response.

The distinction between potential data and group-angle data also
matters at a seam. A group-angle perturbation on a row of length
$\xi_c$ corresponds to a potential of size $\xi_c^{-1}$.
If $\mathsf m$ is multiplication by $\xi_c$ and $Q$ is the
linear average per unit length, then $K_0=\mathsf mQ$.
The identities $Q\mathsf H=1$ below refer to potential data;
the response to $b$ is $\mathsf H\mathsf m^{-1}$.

For the construction one first uses the comb slice, whose
projection $R_{\rm comb}$ is defined by the corner-comb pillar
averages. The subsequent response uses the projection $R$
associated with the actual averaged contours. These slices
describe the same critical orbit, but their potentials and
projections need not coincide. The comparison uses bounds on
the rows and columns, rather than an identification of the
projections.

For one step those norms have leading sizes $L$ and $L^{-3}$.
Their relative errors at an intermediate level $m$ are bounded
by
\[
 C\big(\alpha_0L^{-2(j-m)}
             +L^{m-j}\xi_j\sup|A'|\big).
\]
The errors sum geometrically to a constant independent of $j$.
The composed rows and the accumulated site factors consequently
obey
\[
 \sum_b|K_{A'}(c,b)|\le C\xi_c,\qquad
 \sum_{c\in\Gamma_j}|K_{A'}(c,b)|\le C\eta_k^4\xi_j^{-3},\qquad
 \sum_m|Y_m|\le C\|A'\|_{(-1)},\quad v_j=\prod_me^{iY_m},
\]
with the column bound taken separately at each row level.
Here $\|A'\|_{(-1)}=\sup\mathsf L|A'|$.
These are the uniform bounds used for the transpose of
the constraint in the multiplier equation. A comparison that
introduced a fixed error at every level would not suffice.

\emph{Folding and local problems.}
A field known only on a box is placed in the local estimates
by folding onto that box. In each coordinate the folding map
is the identity inside, reflects across the adjacent face, and
is constant farther out. Collapsed bonds receive the value one.
Every plaquette of the folded field is either one or an original
plaquette, possibly with reversed orientation. Hence folding
preserves the real curvature bound exactly. It also preserves
the local dependence of the averaged quantities on double blocks.
Across a seam the change of scale permits the fixed loss $L^2$;
it does not permit a loss repeated for every preceding level.

For a history the first determining member is
$\Omega_0=\Omega_1^{\sim}$, not the entire bare torus. This
excludes uncontrolled old large fields from the regularity
requirements. At its outer boundary there are no large blocks
left to average: the rows there act at the bare scale. This
geometric fact will be used in the local inverse construction.

\subsection{Propagators on separated regions}

We distinguish estimates for a whole inverse from a decomposition
into local contributions. The first follows from a gap and an
exponential weight. The second requires a local construction of
the inverse, since later operations change the field and the
regions outside the support of one contribution. Both must be
uniform in the fine mesh and in the number of layers.
The assembly follows the local-inverse and walk method of
Ba\l aban\cite{Balaban1984,Balaban1985}; the centred averaging
rows and the scale-dependent gauge projections are retained
throughout the estimates here.

\emph{The operator class and the current variable.}
For the complex estimates use pairs $(\mathbf U,\mathbf J)$,
where $\mathbf U=e^{i\eta_kA'}U$ and $U$ is real. On a layer
of scale $\xi_j$ require
\[
 \begin{gathered}
 |\partial U-1|,|\partial\mathbf U-1|
                          <\alpha_0\eta_k^2\xi_j^{-2},\\
 \xi_j|A'|,\ \xi_j^2|\nabla_UA'|<\alpha_1,
 \qquad |\mathbf J|<\gamma_0\xi_j^{-3}.
 \end{gathered}
\]
In each local cube of side $O(M_1\xi_j)$ in two consecutive
layers there must also be a real gauge with
$\xi_j|A|,\xi_j^2|\nabla A|\le\alpha_0$. This is a bound on
the entire cube, not $M_1\alpha_0$. When a regularity theorem
produces only $CM_1\varepsilon$ in such a gauge, it is that
larger quantity which must lie below the ceiling of this class.
The field strengths of the construction are chosen after the
cube sizes for this reason.

The current is an independent variable. Small curvature and
small first differences alone do not bound the current uniformly
in the mesh: a rapidly alternating field can satisfy those
bounds while its second differences grow like a negative power
of the mesh. Taking $\mathbf J$ independently with its own
radius supplies the derivative control used by the propagators.
The pair is on shell when $\mathbf J=J(\mathbf U)$ at the
bonds in question. At the outer boundary the current is that of
the inner-plaquette action fixed above.

All operator transposes in a complex formula below are for the
bilinear pairings. We continue to write ${}^*$ for them. They
are not conjugate adjoints at complex data. Positivity statements
are made only at real pairs.

For the site average $Q'$ and the linear average $Q$ per unit
row length, put
\begin{equation}\label{eq:mrg-landau-projection}
 G'=(\Delta_{\mathbf U}+Q'^*\mathsf a Q')^{-1},\qquad
 R=I-G'Q'^*(Q'G'^2Q'^*)^{-1}Q'G'.
\end{equation}
Here $\mathsf a$ has size $a_0\xi_y^{-2}$ on tops and
$a_0\xi_c^{-2}$ on rows, with a fixed $a_0>0$. At real data
the subtracted operator is the orthogonal projection onto the
range of $G'Q'^*$. Hence $R^2=R=R^*$ and $Q'G'R=0$.
It projects onto $\Delta_U\ker Q'$, rather than onto
$\ker Q'$ itself.

The uncorrected vector inverse is
\begin{equation}\label{eq:rg-green}
 \mathsf G_0(\mathbf U)=
  (\Delta_{\mathbf U}^{\rm W}+D_{\mathbf U}RD_{\mathbf U}^*
                              +Q^*\mathsf a Q)^{-1}.
\end{equation}
On one lattice $Q=Q_k^\sharp$ and the row mass is of order
one in its units. On a layer it has size $\xi_j^{-2}$.

\emph{The flat scalar problem at all scales.}
At $\mathbf U=1$ the site form has the lower bound
\[
 \langle\lambda,(\Delta+Q'^*\mathsf a Q')\lambda\rangle
       \ge\gamma\sum_j\xi_j^{-2}
                            \|\lambda\|_{L^2(\mathcal L_j)}^2.
\]
The local Poincar\'e inequality controls the difference from the
block mean; the term $Q'^*\mathsf a Q'$ controls that mean.
Conjugate the equation for
$u=G'(f+\nabla^*F)$ by a regularized exponential weight.
The slope of the weight is $O(\xi_j^{-1})$ in layer $j$,
matching its mass. For a small fixed rate the perturbation is
at most half the gap, and the weighted energy estimate is
\[
 \|\xi^{-1}e^\varphi u\|_2
       +\|e^\varphi\nabla u\|_2
 \le C\gamma^{-1}
       \big(\|\xi e^\varphi f\|_2+\|e^\varphi F\|_2\big).
\]
No iteration in the number of layers occurs. Restricting this
estimate to local cubes and using interior estimates at scale
$\xi_y$ gives the corresponding pointwise and first-difference
estimates. In that step the averaging mass is treated as a
source at its local scale.

The one-scale interior estimates must themselves be uniform in
the number of fine sites in a block. They use the massive free
resolvent and the factorization of
$\partial(1-R)\partial^*$ through the unit lattice. The
weighted $L^2$ estimate is first improved to $L^8$ with the
resolvent kernel, and then to a bound on the field and its first
differences with the derivative kernel. The kernel exponents
are respectively $8/5<2$ and $8/7<4/3$. The coarse-image
forcing in the equation is controlled before the second step.
The detailed flat calculation for the centred straight averages
is given again at its curvature use below; it does not by itself
address the seams between different averaging orders.

There is no Dirichlet problem on a fine lattice cut through
large averaging blocks. The separated geometry places
$\partial\Omega_0$ in layer zero. There $Q=Q'=1$ and the
gauge term vanishes. A local Dirichlet inverse therefore has
\[
 a_0\le\eta_k^2G_\square^{-1}\le a_0+16.
\]
Its mass is comparable to its norm and finite range gives
decay on arbitrary sets of bare bonds. Using instead a
Dirichlet inverse of positive averaging order would require
uniform boundary regularity that need not hold, even at a flat
face or a re-entrant edge. Interior local inverses away from
layer zero are formed by periodic one-scale or two-scale
problems and cut off before their seams.

\emph{The two-scale problem at a seam.}
At a junction of orders $j$ and $j+1$, use the exact linear
change of gauge
\[
 A=A_2+H_jB-\partial H'_j\vartheta.
\]
Here $A_2\in\ker Q_j$, $B$ is a one-form on the coarse
lattice vanishing on the axial trees of the coarse blocks,
and $\vartheta$ is a coarse site function of zero mean in
each of those blocks. Its support lies in the coarse region.
The operator $H_j$ minimizes the one-form action at order
$j$; $H'_j$ is the scalar minimizer for $\|\Delta\lambda\|^2$,
not for $\|\partial\lambda\|^2$.

With these choices the representation is unique. The form
splits into a form in $A_2$, a form in $B$, and a form in
$\vartheta-\vartheta_1(A_2,B)$. The reason for retaining
the scalar square is that the Landau projection changes at
the seam. The three resulting unit-lattice forms have gaps.
In particular the axial part obeys
\[
 L^2\|Q_1^{\rm st}B\|^2+\|\partial_1B\|^2
                         \ge c_L\|B\|^2,
 \qquad c_L\ge cL^{-C}.
\]
One may retain the sharper negative power $L^{-5}$ in the
version with its specified axial coordinates. The fixed
positive constant $c_L$ suffices for the assembly here.
The averages in this inequality are the straight one-step
averages. The gauge-corrected plain average is a different
operator.

The tree restriction is essential. A gradient of a function
of zero block mean has vanishing straight average and zero
curvature; it would be a null direction without the tree
condition. For rows next to the coarse region, the norm of
$Q_1^{\rm st}B$ includes outward rows with at least one end
inside that region. Omitting them would leave the gradient
of its indicator uncontrolled.

The two-scale representation is used for one junction. It is
not iterated to express an order-$k$ gauge propagator as a
long product of lower-order ones. Different orders use
different gauge projections, so such an iteration would lose
the uniform constants. Uniform one-scale estimates, the local
two-scale calculation, and the spatial parametrix are the
three ingredients of the many-layer estimate.

\emph{Inverses on tops and rows.}
Two finite-dimensional lower bounds are needed in addition
to the scalar gap. On tops they have the scale of
$Q'G'^2Q'^*$; on rows they have the scale of
$Q\mathsf G_0Q^*$. Their normalized forms are
\[
 \Sigma_{\rm s}=\xi^{-2}Q'G'^2Q'^*\xi^{-2},\qquad
 \Sigma_{\rm r}=\mathsf m^{-1}Q\mathsf G_0Q^*\mathsf m^{-1}.
\]
The desired lower bounds are $\Sigma_{\rm s},\Sigma_{\rm r}
\ge\gamma I$, with the row and top pairings inherited from
their averaging operators. Changing from potential rows to
angle rows $\mathsf mQ$ must be accompanied by the
corresponding change of transpose and multiplier.

For the site matrix, test by block bumps $\psi_y$ with
$Q'\psi_y=\delta_y$, whose second differences, including
those across block faces, have size $C\xi_y^{-2}$.
Duality with these bumps gives
$Q'G'^2Q'^*\ge\gamma\xi^4$. This is the scaled version
of the unit-block bump calculation: the factor $\xi^4$
is the fourth power of its local length, not a constant
independent of scale.

For the row matrix, the variational formula requires a
right inverse for all rows simultaneously. On one lattice,
bumps placed near the face of a coarse bond give an average
matrix with its diagonal share greater than one half; the
remaining shares are its off-diagonal part. Diagonal
dominance inverts that matrix with bounded energy.

At a seam the correction is exact rather than perturbative.
Start with the same-level right inverses $\mathsf E_0$,
chosen biorthogonal to the two chain-count profiles that
neighbouring rows induce in a pillar. Their error has the
form
\[
 Q\mathsf E_0=I+\mathsf N.
\]
It goes only from noncrossing rows to crossing rows one
level higher. A row in that image has no further image,
because the gap separates the next seam. Thus
\[
 \mathsf N^2=0,\qquad
 \mathsf E=\mathsf E_0(I-\mathsf N),\qquad Q\mathsf E=I.
\]
Each coefficient of the cross-level error is at most
$L^{-4}$ and the coefficients in a row sum to at most one.
The weighted overlap estimate required for its energy is
$\|\mathsf m^{-1}\mathsf N\mathsf m\|\le CL^{-1}$.
With this estimate the energy of the correction is
controlled without a loss depending on the number of levels.

The remaining local calculation in this row construction is
to realize those biorthogonal profiles by the actual
centred-contour rows, with bounded overlap and the stated
weighted energy. Nilpotence gives the correction once this
calculation is verified; it does not establish the bump
estimate itself. The row normalization must be kept throughout
that calculation. This is one of the local inverse
verifications required by the layered construction.

Once these lower bounds and the decaying kernels have been
obtained, their inverses are controlled by conjugation and
the Schur test. If $\Sigma\ge\gamma I$ and
$|\Sigma(y,y')|\le Ce^{-\tau d_{\mathbf B}(y,y')}$,
conjugation by a smaller exponential changes it by an
operator of norm at most $\gamma/2$. A Neumann series gives
\[
 |\Sigma^{-1}(y,y')|\le2\gamma^{-1}
                      e^{-\tau'd_{\mathbf B}(y,y')}.
\]
Undoing the scale normalization gives, symmetrically, factors
$\xi_y^{-2}\xi_{y'}^{-2}$ for the top inverse and
$\xi_c^{-1}\xi_{c'}^{-1}$ for the row inverse. Spending a
fraction of the rate on their ratio gives the useful forms
\[
 |(Q'G'^2Q'^*)^{-1}(y,y')|
             \le C\xi_y^{-4}e^{-cd_{\mathbf B}(y,y')},\qquad
 |(Q\mathsf G_0Q^*)^{-1}(c,c')|
             \le C\xi_c^{-2}e^{-cd_{\mathbf B}(c,c')}.
\]
No coarse matrix is inverted by a bound uniform in the
unscaled entries alone.

\emph{The corrected Hessians.}
In \eqref{eq:rg-green} the curl controls nongauge high
frequencies, the Landau
term controls the remaining fine gauge modes, and the
average mass controls the coarse modes.

Let $D=D_{\mathbf U}$ and $\pi=I-DG'RD^*$. The first
correction $\Delta'_\pi[\mathbf J]$ is linear in the
independent current. On shell it has the identity
\[
 \Delta_{\mathbf U}^{\rm W}-\Delta'_\pi[J(\mathbf U)]
                     =\pi^*\Delta_{\mathbf U}^{\rm W}\pi.
\]
The gauge-corrected inverse and response are
\[
 \mathsf G=(\mathsf G_0^{-1}-\Delta'_\pi[\mathbf J])^{-1},
 \qquad \mathsf H=\mathsf GQ^*(Q\mathsf GQ^*)^{-1},
 \qquad \widetilde{\mathsf G}=\mathsf G-\mathsf HQ\mathsf G.
\]
For the second correction, let $C^{(2)}$ be the quadratic
part of the straightened average per unit row length and
define $\Delta^{(2)}[j]$ by
$\langle A,\Delta^{(2)}[j]A\rangle
=2\langle C^{(2)}(A),j\rangle$. Then
\[
 \mathsf G_1=
  (\mathsf G^{-1}-\pi^*\Delta^{(2)}[\mathsf H^*\mathbf J]\pi)^{-1},
 \qquad \mathsf H_1=\mathsf G_1Q^*(Q\mathsf G_1Q^*)^{-1}.
\]
It is this last inverse that differentiates the nonlinear
background at a critical datum: the Hessian of the action
must be corrected by the multiplier paired with the
curvature of the constraint fibre. When the current vanishes
the two corrections vanish.

The exact gauge identity starts from covariance of the
straightened average,
\[
 QD=\bar DQ'.
\]
It is not an identity for the plain logarithm of the
average. Differentiating action invariance and retaining
the independent current gives
\[
 \mathsf G^{-1}D\lambda
  =D\Delta_{\mathbf U}\lambda
          -\Delta'_\pi[\mathbf J-J(\mathbf U)]D\lambda
       \quad(\lambda\in\ker Q').
\]
Consequently
\[
 \mathsf G_{(1)}DR
 =DG'R+\mathsf G_{(1)}\Delta'_\pi[\mathbf J-J(\mathbf U)]DG'R.
\]
On shell the defect vanishes and
$RD^*\mathsf H_{(1)}=0$. Off shell it must be retained;
positivity at real pairs and holomorphy do not remove it.
In particular an exact slice identity cannot be transferred
to $\mathsf G_0$ at nonzero current.

\emph{Local inverses and the parametrix.}
For the flat many-layer operator choose
$\sum_\square h_\square^2=1$ and local periodic
one-scale or two-scale inverses $\mathsf G_\square$.
The cutoffs vanish across the seams to finer layers before
the local model is changed. In layer zero next to the
outer boundary use the massive bare Dirichlet inverse.

The delicate term is $\mathcal P=\partial(1-R)\partial^*$.
With a second cutoff $\zeta=1$ near the support of $h$,
its commutator is split as
\[
 \mathcal Ph-h\mathcal P_\square
 =\zeta[\mathcal P_\square,h]
   +\zeta(\mathcal P-\mathcal P_\square)h
   +(1-\zeta)\mathcal Ph.
\]
The first term costs $M_1^{-1}$: one derivative of the
cutoff has size $(M_1\xi_j)^{-1}$ and is paired with a
local inverse gradient of size $\xi_j$. The second term
is a change-of-domain error. Resolvent identities on the
two domains, recombined through $R=I-(1-R)$, give decay
across the collar and a factor $e^{-cM_1}$. The last term
is a tail with the same decay. Inserting the partition
of unity at its endpoint turns it into a step between
local cubes. These three estimates are what make the
spatial assembly possible despite the nonlocal gauge
projection.

At a complex regular pair the corrections introduce
additional nonlocal expressions. The local construction
unfolds them, together with $R$ and the average constraint,
into an auxiliary system of finite range. Its unknowns
are the potential, the site and row auxiliary fields,
and the multiplier. The propagators are blocks of the
inverse of that system. The local perturbation estimate
to prove in the gauges of the class has size
$C(\alpha_0+\alpha_1+\gamma_0)$ relative to the flat
local system.

The needed verification is at this last local system:
its auxiliary equations and scale norms must give a
bounded flat inverse and the stated perturbation estimate
without a loss of derivatives. In particular a source
that is a divergence must be kept in that form, and no
uncontrolled derivative may be placed on its row
auxiliary variable. The flat scalar gap alone proves
none of these facts. This is the local analytic
calculation still required for the full corrected
layered inverse.

With those local inverse bounds, the gluing argument is
explicit. Write $\mathbb K$ for the auxiliary operator,
choose source cutoffs $\theta_\square$ with
$\sum_\square\theta_\square=1$, and larger cutoffs
$\widetilde h_\square=1$ near their supports.
The local system $\mathbb K_\square$ agrees with
$\mathbb K$ on the stencil neighbourhood of
$\operatorname{supp}\widetilde h_\square$. Put
$S_k=\sum_\square\widetilde h_\square
                  \mathbb K_\square^{-1}\theta_\square$.
The derivative of $\widetilde h_\square$ is separated
from the support of $\theta_\square$ by a fixed
fraction of the local cube side. Then
\begin{equation}\label{eq:xrg-parametrix}
 \mathbb KS_k=I-\mathcal K_k,\qquad
 \mathcal K_k=-\sum_\square
 [\mathbb K_\square,\widetilde h_\square]
                      \mathbb K_\square^{-1}\theta_\square.
\end{equation}
Here the identity is first formed for the finite-range
system and then read in its propagator blocks. A local
commutator has norm, between local cubes,
\[
 CM_1^{-1}e^{-c\delta_{\rm G}M_1/L}.
\]
The power comes from differentiating the cutoff; the
exponential comes from the separation of its derivative
from the source cutoff. At fixed $L$ it outweighs the
bounded number of adjacent cubes. In the weighted
local norm $\|\mathcal K_k\|<1$, uniformly in the
regions, and
\[
 \mathbb K^{-1}=S_k\sum_{n\ge0}\mathcal K_k^n.
\]
Its propagator blocks give the required Green
functions. The Neumann series is a consequence of
the local calculation, not a replacement for it.

\emph{Walks, branches and local dependence.}
The terms of this series are indexed by overlapping
local cells. A cell consists of $O(1)$ cubes of side
$O(M_1\xi_j)$ in two consecutive layers. For the
corrected inverse $\mathsf G_1$, the coefficient
$\mathsf H^*\mathbf J$ itself has a walk expansion.
Its zero-length part stays in the local inverse; the
remainder is a small vertex bearing a branch. Thus
the resulting index is a tree walk.

List the cells depth first, recording the parent of
each new cell. For a fixed initial cell there are
at most $C^m$ lists of length $m$. Arbitrary parent
lists would count the same tree factorially many
times and would defeat the estimate. The two
depth-first conventions, inserting a branch as
soon as it is met or listing the trunk before
its branches in reverse order, give the same
tree, hull and length.

A typical propagator term satisfies
\begin{equation}\label{eq:rg-walks}
 |\mathsf G_\omega(x,y)|
 \le C(CM_1^{-1/2})^{|\omega|}M_1^{-|\omega|/2}
                   e^{-\delta_{\rm G}d_{\mathbf B}(\omega;x,y)},
\end{equation}
with the appropriate local scale factors for its
source and output. The first factor pays for
$C^{|\omega|}$ possible lists. The second remains
available when terms are transported through
coverings. Spatial decay is supplied by the local
inverses and does not deteriorate as $M_1$ grows.

The hull of a tree walk enlarges its trail by
five local cube sides, taking each cube at its
own scale. It is not cut at $\Omega_0$. Its
intersection with $\Omega_0$ is the carrier of
the field variables and can be disconnected.
Each term is holomorphic in the pair on that
carrier and reads the regions only through
their traces on the uncut hull. This is the
local dependence needed in the volume comparison.

The terms are chosen as the same functions when
a local translation carries the hull, pair and
traces of the regions between two settings.
For the fluctuation form the translation carries
also the next-level region and the conditioning
set. Equality of terms under such transport is
stronger than equality or decay of their sums.
It is imposed in their construction: real gauge
averaging and conjugation symmetrization preserve
the sum and the local bounds. Translation
averaging, push-forward from covering tori, and
a normal-family limit fix a single assignment
on all tori with room.

When a walk wraps under push-forward it can
produce at most $|\omega|+1$ additional pieces.
The reserved factor $M_1^{-|\omega|/2}$ absorbs
that loss. The entrywise bound used only to
take a normal-family subsequence may depend
on the level. It enters no decay rate, local
norm or sum of weights. The assignment is
made once, rather than choosing a new
subsequence for every later comparison.

\emph{The fluctuation form.}
The inverse used for a fluctuation is not
obtained by restricting a previously chosen
global covariance. Let $\mathsf C_V$ eliminate
the pivots in the one-step average at $V$,
and let $\mathsf G_V^{(2)}$ be the quadratic
soft-gauge form. The quadratic effective
action in the datum is
\[
 \Delta_k=(Q\mathsf G_1Q^*)^{-1}-\mathsf a
              -\mathsf E^{(2)}[\mathsf H_1^*\mathbf J],
 \qquad
 \mathsf A=\mathsf C_V^*(\mathsf G_V^{(2)}+\Delta_k)\mathsf C_V.
\]
Here $\mathsf E^{(2)}$ is the quadratic part of
the plain one-step constraint. Its correction
is not optional. Differentiating the multiplier
system identifies the first derivative of the
background action with its multiplier. A
second differentiation gives the inverse
coarse form minus the artificial mass.
Along the nonlinear one-step fibre there is
also the second derivative of that fibre;
paired with the multiplier it gives exactly
the last term in $\Delta_k$.

At the critical pair of the step this yields
$\mathsf A=\mathsf C^*\Delta^{(k)}\mathsf C$.
One can add the average equation as a multiplier
equation in the auxiliary system, so the
fluctuation covariance is another inverse
block. For the free bonds $Y$ of a union of
next-level blocks, its definition is
\[
 C^{(k)}(Y)=(\chi_Y\mathsf A\chi_Y)^{-1}.
\]
The required real lower bound is
$\chi_Y\mathsf A\chi_Y\ge\gamma^{\rm pos}I$.
At the flat field it uses the axial unit-lattice
gap and the contraction of the curvature norm
under averaging. For a real regular pair the
local inverse is continued to the flat one;
the positive gap must be retained along this
homotopy. The same local parametrix then has
error $CM_1^{-1}<1$ and yields exponential
decay, at a rate allowed to depend on $L$.
This positivity verification belongs to the
local corrected-system estimate above. Mere
continuity of a finite-volume inverse would
not give its uniform gap.

With that verification, let $x$ lie in the local cell at $y$
and let $f$ be supported in the local cell at $y'$.
The estimates used later take the forms
\[
 |\mathsf G_kf(x)|+\xi_y|\nabla\mathsf G_kf(x)|
 \le C\xi_y^2e^{-\delta_{\rm G}d_{\mathbf B}(y,y')}
                                      \|f\|_\infty,
 \qquad
 |C^{(k)}(Y)(b,b')|\le Ce^{-\delta'|b-b'|}.
\]
There are first-difference and H\"older variants,
with H\"older differences taken only over
distances at most the local scale $\xi_y$.
Complex parallel transports over larger
distances need not remain uniformly bounded.

\subsection{The nonlinear background problem on a sequence of layers}

The inverse estimates are used twice. They first construct a
regular critical orbit from a regular trial field, and then
estimate its response. A contraction on a gauge slice alone
does not prove the Euler--Lagrange equation in the full fibre.
The constraint is therefore kept nonlinear and a multiplier
is retained until the gauge condition has been recovered.

\emph{The trial field and the nonlinear equation.}
On one lattice lift a coarse configuration $V$ by putting one
on bonds within a block and $V(c)$ on the bonds crossing the
face of $c$. The lift $V^\downarrow$ satisfies
$\mathcal M(V^\downarrow)=V$; its nontrivial plaquettes are
plaquettes of $V$. The preceding background
$U_0=U_{k-1}(V^\downarrow)$ is a trial field. Write
\[
 B=\frac1i\log[V(\mathcal M^kU_0)^{-1}].
\]
This datum is zero for that particular trial field on one
lattice. It is kept because a nearby trial field, a seam
compensation, and a variation of $V$ all produce nonzero
$B$.

In the chart $U=e^{i\eta_kA'}U_0$,
\[
 \begin{split}
 F(A')&=F(0)+\langle A',J(U_0)\rangle
       +\tfrac12\langle A',\Delta^{\rm W}_{U_0}A'\rangle
       +\mathcal N(A'),\\
 \mathcal N(A')&=O(|J(U_0)|\,|A'|^2+|A'|^3)
 \end{split}
\]
after the constraint is straightened. The quadratic term
proportional to the current is the current paired with the
quadratic change of coordinates. It remains even when the
original action has already been expanded to second order.

For a potential on the slice write $A'=\mathsf H_kB+A_1$
with $QA_1=0$. Its regularity equation is
\begin{equation}\label{eq:rg-background-equation}
 A_1=-\widetilde{\mathsf G}_kJ(U_0)
       -\widetilde{\mathsf G}_k
                         \mathcal N'(\mathsf H_kB+A_1).
\end{equation}
The constrained inverse maps a force into a potential and
its first differences. The derivative of the nonlinear term
has size $C(\varepsilon+|B|)$ in those norms. This proves
the estimates once a solution on the slice is known; it does
not alone prove that the resulting potential has the
prescribed gauge.

For existence retain the multiplier equation
\begin{equation}\label{eq:xrg-multiplier-system}
 \begin{split}
 F'(A')+D_{U_0}R_{\rm comb}D_{U_0}^*A'
       &=[\mathcal Q_k^\sharp{}'(A')]^*\omega,\\
 \mathcal Q_k^\sharp(A')&=b(B).
 \end{split}
\end{equation}
In the first construction the reference rows and projection
are those of the comb slice. The actual straightened
constraint differs from these rows by a small linear term
as well as by its quadratic remainder. Write it as
$\mathsf Q A'+\mathsf C(A')$, and put
\[
 \mathsf G_0=(\Delta^{\rm W}_{U_0}+DR_{\rm comb}D^*
                                      +\mathsf Q^*\mathsf Q)^{-1},
 \quad
 H_0=\mathsf G_0\mathsf Q^*(\mathsf Q\mathsf G_0\mathsf Q^*)^{-1},
 \quad \widetilde G_0=\mathsf G_0-H_0\mathsf Q\mathsf G_0.
\]
Here the row lengths have been included in $\mathsf Q$.
Abbreviate
$F'(A')=F'(0)+\Delta^{\rm W}_{U_0}A'+\mathcal V(A')$.
The multiplier system has the fixed-point form
\[
 \begin{split}
 A'&=\widetilde G_0\mathsf N
                     +H_0(b(B)-\mathsf C(A')),\\
 \mathsf N&=-F'(0)-\mathcal V(A')+\mathsf C'(A')^*\omega,\\
 \omega&=H_0^*\big[F'(A')+DR_{\rm comb}D^*A'
                                      -\mathsf C'(A')^*\omega\big].
 \end{split}
\]
The current $F'(0)$ and the datum are the two inhomogeneous
terms. The current is not discarded when the trial field
fails to be critical.

Use
\[
 \|A'\|_X=\sup_y\max\{\xi_y|A'|,
                                \xi_y^2|\nabla_{U_0}A'|\},
 \qquad |f|_{(-3)}=\sup_y\xi_y^3|f|.
\]
Summation by parts is needed in the multiplier component:
$H_0^*\mathrm d^*\mathrm d A'=(\mathrm d H_0)^*\mathrm d A'$. It leaves first differences,
which are controlled by $\|A'\|_X$; directly differentiating
$A'$ twice would lose the required uniformity in the mesh.

The Lipschitz matrix has entries $C(\alpha_0+r)$ on a ball
of radius $r$, except for one bounded entry. Weight the
multiplier norm by a fixed $\Theta^{-1}$. First choose
$\Theta$ to reduce that bounded entry; then choose the
radii so that the small entries multiplied by $\Theta$
remain small. This gives a contraction with constant at
most one half. On a sequence one may eliminate the
multiplier: $I+H_0^*\mathsf C'(A')^*$ is invertible, so
$\omega=w(A')$. Its Lipschitz constant need not be small,
but it reenters the first equation multiplied by the
small derivative $\mathsf C'(A')^*$. The contraction can
then be run in the potential norm alone.

Two balls serve different purposes. A fixed small ball
gives uniqueness. A ball of radius
$C(|F'(0)|_{(-3)}+\sup|B|)$ gives the linear bound.
Uniform convergence of the holomorphic iterates gives
a holomorphic solution. The same estimates bound the
multiplier. The uniform comparison of the actual
straightened rows with the comb rows is needed here;
identifying their gauge projections would not give it.

\emph{Recovering the gauge condition.}
For a real solution, differentiate gauge transformations
that preserve the nonlinear constraint. In their coordinate
directions $X_\lambda$, action invariance gives
$\langle F'(A'),X_\lambda\rangle=0$, and the multiplier
term vanishes because $K_{A'}X_\lambda=0$. The remaining
equation is
\[
 \langle R_{\rm comb}D^*A',
                         R_{\rm comb}D^*X_\lambda\rangle=0.
\]
The local gauge construction makes the second factors
span the range of the projection. Therefore
$R_{\rm comb}D^*A'=0$. The equation on the slice is
then the constrained Euler--Lagrange equation. The
transversality used in this step concerns the actual
nonlinear average with its site factors; infinitesimal
gauge covariance at zero alone is not its full statement.

The later response is constructed with the true
straightened rows $K_0$ and their Landau projection $R$.
It parametrizes the same orbit. In its multiplier
equation the reference inverse is $\mathsf G_1$, since
the linearization at a critical pair contains the
multiplier paired with the second derivative of the
constraint. This explains why the two corrected
Hessians were retained separately above.

\emph{Representatives in the slice and uniqueness of the orbit.}
Uniqueness of a fixed point is useful only after every
regular configuration in the fibre can be placed in its
ball. On one lattice first align the contour variables
of the preceding averages. The rigidity estimate for
one step says that matching contours, average error
$\alpha_1$ and plaquette size $\varpi$ imply bond error
at most $\alpha_1+C\varpi$. The coefficient of
$\alpha_1$ is one. Descending through the averages
therefore adds the curvature errors rather than
multiplying the datum error at every level.

A change to the gauge slice then restores the
normalization at the nodes. It is obtained by a
finite-dimensional fixed-point argument in the local
gauge variables. Its output is an estimate, not a
holomorphic choice of gauge. The analytic potential
will instead be selected by the unique multiplier
solution.

For a sequence, a single preliminary gauge equal to
one at the tops aligns all contour variables at once.
One uses rigidity with a defect: the mismatch of
the contour variables at both endpoints is added
to the bond estimate. Fix a bond and let $j$ be
the last level containing both its endpoints.
The increments below that level have sizes
$C\alpha_0L^{2(n-j)}$ and sum geometrically. A
seam incurs the factor $2L-1$ only once, because
the gauge is frozen in the outer block. Thus
\[
 |\mathcal M^n(U^{\hat g})(b)
                       \mathcal M^nU_0(b)^{-1}-1|
                 \le(2L-1)\alpha_1+C\alpha_0
\]
for the retained bonds. This estimate passes to
the truncations of the determining sequence.

One next ascends through those truncations.
Rescaling can lose $L^2$, but the estimate in
the slice resets the potential:
\[
 \|A'\|_X+|\mathrm d^*\mathrm d A'|_{(-3)}
            \le\mathsf B_1\big(\alpha_0+
                                \sup|\mathcal Q_k^\sharp(A')|\big).
\]
The constant $\mathsf B_1$ is determined by the
inverse norm, independently of the size before
this reset. It suffices that the rescaled
potential remain within the fixed small ball.
This is the first mechanism preventing a loss
proportional to the number of levels.

Let now $U_a,U_b$ be two regular critical
configurations with the same datum, and take
$U_a$ as reference. The preceding construction
puts $U_b$ in the slice as $e^{i\eta_kA'}U_a$,
with zero reduced datum and potential in the
uniqueness ball. Criticality supplies a
multiplier of size $C\varepsilon_{\rm reg}$.
The pairs $(A',\omega_b)$ and $(0,\omega_a)$
solve the same zero-datum system; hence they
agree. Thus $A'=0$, and the compensating node
factor is one. This gives one orbit under
$\mathcal G^{\rm res}(\mathbf B)$.

\emph{The radius improvement.}
The trial lift and change of level may enlarge
the regular radius by $L^3$. One cannot let this
loss pass to the next level unchanged. The
local estimate for a critical field improves
its radius independently of the construction
that produced it.

Choose a cube of side $\sigma$, trivialize
there, cut off the reference field to one,
and align its slice with a comb gauge of the
datum. Splitting the equation into near and
far contributions gives an estimate of the
following form:
\[
 r(x_0)\le Cv_R(x_0)
    +\big(C\sigma^2e^{-c\delta_*\sigma}
                       +C\sigma^4r_R(x_0)\big)r_R(x_0).
\]
Here $r$ is the local radius of the critical
field, $v$ is the size of the datum tests,
and the suffix $R$ is a maximum over a
weighted ball of radius $R\asymp L\sigma$.
The first error is the far part of the
inverse, including the length of the comb;
the second is the local nonlinear error.
The coefficient in front of $v_R$ is
independent of $\sigma$. This is why one
can first choose $\sigma$ large and then
the regular radius small.

The remaining estimate in this local
calculation is the uniform bound for the
cutoff and nonlinear errors that gives the
two displayed coefficients. It uses the
local inverse at the trivial field and
the actual averaging rows. Once it is
verified, choose $\sigma$ so that the
exponential term is small, and then take
$\varepsilon_{\rm reg}\ll\sigma^{-4}$.
A maximum argument gives the global
improvement
\[
 \varepsilon_{\rm reg}\longmapsto
                \max\{C\varepsilon,\varepsilon_{\rm reg}/2\}.
\]
The local form uses a weight varying at
rate $\delta\asymp(L\sigma)^{-1}$.
The variation on the ball of radius $R$
is then bounded, and the same absorption
gives
\[
 r(x)\le C\max_\tau
       e^{-\delta d_{\mathbf B}(x,\tau)}
                                      |\partial\mathbf V(\tau)-1|.
\]
It concerns an arbitrary regular critical
field and does not presuppose uniqueness.

Repeated halving restores $C\varepsilon$
after every trial step, closing
$C\varepsilon\to CL^3\varepsilon
\to C'L^3\varepsilon\to C\varepsilon$.
There is also a separate reset when the
datum is descended near a seam. A section
of one step is exact away from the seam;
near it the section loses a fixed factor
$C_\downarrow$ in a collar. The gap keeps
that collar away from the errors of
preceding sections. Using a small radius
in the collars and radius
$C_\downarrow\varepsilon$ elsewhere
therefore descends with the same loss
at each stage, rather than its $k$th
power. The representative estimate,
the radius improvement, and this collar
reset address different accumulations.

\emph{Decay of the response.}
At a critical $U_0$, solve once more
with the actual rows $K_0$ and $R$.
The multiplier linearization has the
additional term
\[
 W_0\delta=
      \left.\frac d{dt}K_{t\delta}^*\omega_0\right|_{t=0},
\]
which is included in $\mathsf G_1^{-1}$.
Because the constraint was not replaced
by a global change of variables, the
remaining nonlinear terms stay local.
All nonlocality is in the already
estimated inverse operators.

The fixed-point derivative is contractive
in every norm with weight
$e^{\theta d_{\mathbf B}(\cdot,y_0)}$
for $\theta$ below a fixed fraction
of the inverse decay rate. If
$\mathbf A(b)$ is its solution with
free row datum $b$, differentiation gives
\[
 (I-\partial\Psi)(\mathsf h,\dot\omega)
                                      =(\mathsf H_1\mathsf m^{-1}e_c,0).
\]
Invert $I-\partial\Psi$ by its
convergent Neumann series in that
same weighted norm. The source
column already decays, and the norm
of the inverse is at most two.
Consequently
\[
 \xi_y|\mathsf h|+\xi_y^2|\nabla\mathsf h|
       \le Ce^{-c\delta_{\rm G}d_{\mathbf B}(y,y_c)}.
\]
Finally compose with $b(B)$.
Each input bond changes only $O(1)$
rows within $O(L)$ local spacings,
so the chain rule retains the decay.
The differentiated multiplier equation
also gives the current and appropriate
second-order members. They are obtained
from that equation, not by
differentiating a merely pointwise
kernel bound twice.

At the trivial datum on one lattice
the linear response is the periodization
of its translation-invariant kernel on
$\mathbb Z^4$. Summing its exponential
tail gives error
$Ce^{-c\delta_{\rm G}\ell_k/2}$ for
lifts within half the torus. This is
the response comparison later used
when a local computation is placed
in a different torus.

\subsection{Complex backgrounds and the domains of the activities}

The complex domain must survive taking averages, forming a local
background, and inserting a fluctuation into an older activity.
An estimate only for a real critical field is insufficient for
these operations. The two radii serve different purposes:
$\alpha_0$ bounds curvature and current, while $\alpha_1$
bounds the complex potential. In the application
$\alpha_1\gg M\alpha_0$, but $\alpha_1^2$ remains small
compared with $\alpha_0$.

\emph{Continuation with a prescribed representative.}
Fix real regular datum $\mathbf V_0$ and the representative
$U_0$ pinned axially outside its nodes. For sufficiently small
$r_B$, the multiplier contraction defines a holomorphic map
$\Phi_{\mathbf V_0}(B)$ on $\sup|B|<r_B$ agreeing at real
$B$ with the representative of the orbit over
$e^{iB}\mathbf V_0$. There is a holomorphic Landau pair
$(\mathbf H(B),u_B)$ with
\[
 \Phi_{\mathbf V_0}(B)
       =(e^{i\eta_k\mathbf H(B)}U_0)^{u_B},\qquad
 RD_{U_0}^*\mathbf H(B)=0,
\]
and with $u_Bv_j[\mathbf H(B)]=1$ at the tops of level
$j$. Its values at zero are $(0,1)$. The fixed-point
estimate gives $\|\mathbf H(B)\|_X\le C\sup|B|$.
At real $B$ the pair is real. The comb slice has its own
holomorphic pair, with the same potential estimate and
an all-site bound on its compensating gauge. At the
trivial reference the two slices agree. Away from it,
they are two gauge-related potentials, not equal maps.

The node condition is part of this statement. A Landau
gauge transformation need not be one at the nodes;
its effect on the plain averages is compensated by
the site factor. Removing that condition would change
the datum being prescribed.

\emph{Why the complex radius enters quadratically.}
Suppose $\mathbf U=e^{i\eta A'}U$, with $U$ real, and
\[
 |\partial U-1|,|\partial\mathbf U-1|<\alpha'_0\eta^2,
 \qquad |A'|<\alpha'_1.
\]
The plaquette expansion gives
\[
 \partial(e^{i\eta A'}U)
 =\big[I+i\eta^2\mathrm d_UA'
       +O(\eta^2|A'|^2+\eta|A'|\,|\partial U-1|)\big]
                                                    \partial U.
\]
Using both plaquette bounds, rather than the size of
$A'$ alone, yields
\[
 |\mathrm d_UA'|\le C(\alpha'_0+\alpha_1'^2).
\]
The real adjoint transports commute with taking the
real and imaginary parts. Thus the real fields
$e^{i\eta\operatorname{Re}A'}U$ and
$e^{i\eta\operatorname{Im}A'}U$ have curvature bounded
by $C(\alpha'_0+\alpha_1'^2)$ as well. A bound merely
linear in $\alpha'_1$ would miss this gain.

For a holomorphic component $\phi$ bounded by $M_\phi$
on a ball of radius $r$, let $w=u+iv$ with $u,v$ real
and $\|w\|\le r/2$. Its first Taylor term is complex
linear, so
\[
 \phi(w)=\phi(u)+i\phi(v)-i\phi(0)
                   +O(M_\phi\|w\|^2/r^2).
\]
Apply this to a rescaled plaquette or current of a
tower background. The three real values are controlled
by the real regularity theorem. The crude complex
supremum occurs only in the quadratic remainder.
Consequently the complex background curvature and
current have the bound
$K^\sharp(\alpha'_0+\alpha_1'^2)$ in their respective
units.

The quadratic form of this error is preserved by a
change of level. At level $n\le p$, the two input
sizes become
$L^{-2(p-n)}\alpha'_0$ and
$L^{-(p-n)}\alpha'_1$. The square of the latter has
the same factor as the former. A bound linear in
$\alpha'_1$ would introduce $L^{p-n}$ and could not
be iterated through a tower. The condition
$K^\sharp(\alpha'_0+\alpha_1'^2)\le\alpha_0$ also
leaves room for the larger potential radius.

\emph{Towers and hereditary domains.}
The activities are evaluated on a triple
$(U,A',\mathbf J)$, where $U$ is real and
$\mathbf J$ is independent. Its pair is
$(e^{i\eta_jA'}U,\mathbf J)$. On a domain $X$ at
level $j$, the simple conditions are the two
plaquette bounds, the bounds
$|A'|,|\nabla_UA'|<\alpha_1$ and
$|\mathbf J|<\alpha_0$, and real local gauges on
the cubes of side $O(LM)$ with potential and
gradient bounded by $CM\alpha_0$.

The domain must also be stable under localization
to every smaller subregion. Given $Z\subset X$
at level $p\le j$, form its tower by successive
erosion. One tower uses $M_1$-cubes and an outer
gap exceeding $\rho_{\rm geom}$; the other uses
the larger localization cubes and a gap measured
in the inner scale. Truncate either tower at
any $n\le p$. Its background is continued from
the real datum to the complex averages of
$e^{i\eta_jA'}U$.

The core condition requires all these backgrounds
to exist, and, after two boundary layers have
been removed, to satisfy
\[
 \begin{split}
 |\partial U_n(Z;\mathbf U)-1|
       &<\alpha_0L^{-(j-p)}\eta_p^2,\\
 |J(U_n(Z;\mathbf U))|
       &<\alpha_0L^{-(j-p)}L^{2(n-p)}.
 \end{split}
\]
The current in the second line is measured in
level-$n$ units. The same conditions are imposed
on the real field. The quantifiers over $Z,p,n$
make the core hereditary. They are needed when
a domain is cut into windows and then evaluated
at an earlier level; a condition on the
background of $X$ alone would not suffice.

At this stage the continued tower is defined
through the real reference $U$, hence through
the triple. The activity construction must
give the same value at compatible splittings
of the same pair. The core itself is not
saturated under arbitrary complex gauge
transformations. When gauge covariance is
used, both representations must remain in
the specified domain, with the required
condition-number bound for the gauge.

With the uniform background and response
bounds on their common complex domains,
the quadratic complex estimate gives the
core conditions from the simple conditions
at smaller radii. The application retains
fixed room in the radii and in the two
removed boundary layers. The common-domain
verification is still needed for the joined
and truncated towers used in this passage;
the quadratic Taylor estimate by itself
does not prove it. The trivial triple
belongs to every core.

To apply the inverse estimates to these
domains, the gauge clause must be converted
to the clause of a regular pair. A gauge
on a cube of side $O(LM)$ has size
$O(LM)\alpha_0$, so the regular-pair
radius is $O(LM)\alpha_0$, not merely
$\alpha_0$. On a label sequence whose
regions are unions of $LM$-cubes, the
layered conversion costs
$O(M)\max_n\alpha_{0,n}$. A cube meeting
layer $n$ uses that layer's gauge; a
smaller cube lying in layer $n+1$ uses
the containing cube and gauge at level
$n+1$, restricted and rescaled. Thus no
critical-point estimate is needed to
place a pair of the activity domain in
the local inverse class. The field
radii are chosen after $M$ so the
converted radii lie below its ceiling.

\emph{Substitution from the next fluctuation integral.}
The outer substitution is
$\mathbf U\mapsto e^{i\eta_j\mathbf A}\mathbf U$
with $|\mathbf A|,|\nabla_{\mathbf U}\mathbf A|
<c\alpha_0$. A core at half the radii is
mapped into the full core when
$\alpha_0\le\alpha_1$ and
$\alpha_1^2\le C\alpha_0$. Each localized tower
changes by its response to the perturbed
averages. The response bound is uniform
through the tower; taking $c$ smaller than
the reciprocal of its constant leaves
the unused half of the curvature and
current radii. The first-difference bounds
retain the potential radius. This is a
statement about every tower of the core,
not just its outer field.

The inner substitution is more delicate
because it evaluates an old activity
inside a new fluctuation integral. Work
on a box $\square_0$ containing $O(1)$
localization cubes. Let $\mathbf H_{k+1}$
be the Landau potential of its background,
and let $e^{iB}$ be the averages of
\[
 e^{i\eta_k\mathbf A}
                         e^{i\eta_{k+1}\mathbf H_{k+1}},
 \qquad |\mathbf A|<c\alpha_{0,k}.
\]
The fluctuation response with a cutoff
is the relevant choice of $\mathbf A$.
Take the background over the trivial
reference at datum $tB+B'$, $0\le t\le1$,
where
\[
 \sup|B'|<c\min\{\alpha_{0,j-1},\alpha_{0,j}\}.
\]
The input field has the simple conditions
and the tower conditions on a slightly
larger box, with radii larger by a fixed
factor. Once the common-domain background
and response bounds hold uniformly through
this substitution, it returns, well inside
the box, a triple in the core of every
level $1\le j\le k$. The response and
agreement margins must be verified also
for the joined towers; no unestimated
remainder may be introduced at a join.

Both the level and its radius must be
kept in this assertion. If
$\nu=k+1-j$, the flow gives
\[
 \alpha_{\iota,k+1}
       \le C(1+\nu)^{1/2}\alpha_{\iota,j},
       \qquad \iota=0,1.
\]
The change of units supplies powers of
$L^{-\nu}$, which dominate this
polynomial loss. Discarding those powers,
or replacing the two radii by one fixed
radius, would not justify the evaluation
of all the earlier terms. The resulting
Landau triple is
\[
 (1,\mathbf H,J(e^{i\eta_j\mathbf H})),
 \qquad \mathbf H=\mathcal H(tB+B'),
\]
and the actual background pair is in
its complex gauge orbit. Evaluations
are transferred only on the portions
of that orbit allowed by the domain.

In a local box the complex gauge is
controlled by its condition number,
\[
 |v|\,|v^{-1}|\le e^{CM\alpha_1}.
\]
There is no bound $|v-1|\le CM\alpha_1$
for an arbitrary real reference gauge:
a pure gauge can already have large
real values. The condition-number
bound is the one used to continue
the averages and the activities.

\emph{Restriction and comparison of backgrounds.}
The finer constraints fix the coarser
ones. If a configuration is varied
inside a smaller determining region
while its finer averages are kept,
the coarser datum outside that region
does not change. The difference of
the two inner-plaquette actions is
exactly the difference of their
actions in the smaller region.
Restriction therefore preserves
criticality. The restricted datum
is $CL^2\varepsilon$-regular, and
uniqueness of the regular orbit gives
the restriction identity.
On one lattice this is
\[
 U_j(\mathcal M^jU_k(V))=U_k(V)
                  \quad\text{modulo gauge},\qquad j\le k.
\]

For locality, a collar boundary is
required to stay within one layer.
Align gauges there and compare the
finer averages with the lift of the
datum at the collar level. The
defect-of-lifts reconstruction has
coefficient one at every descent;
the remaining plaquette errors sum
as
\[
 \sum_{q<j}2\varepsilon L^{2(q-j)}
                                     \le C\varepsilon.
\]
The aligned mismatch is a real datum
$B_1$ of size $C\varepsilon$ supported
in the collar. Its response decays,
so two backgrounds with matching
determining data on the collared
region $W$ differ there, modulo gauge,
by at most
\[
 C_\mu\varepsilon
       e^{-c\delta_{\rm G}d_{\mathbf B}(y,W^c)}.
\]
The constant depends on the fixed
collar width $\mu$; the decay starts
at the inner edge of that collar.
One must not replace it by a constant
independent of a variable collar
width. Nor can a gap of one $M$-cube
in a prepared nest be read as a gap
of $M\check R_j$: that gap supplies
fixed decay in $M$ only.

\subsection{Complex decoupling and the localized response}

Localization of the next fluctuation integral uses a response
whose propagators have been decoupled in space. This is a
separate construction from taking the actual background of a
field with altered averages. The decoupled operators retain
their bounds and locality, but generally lose the algebraic
identities of the original inverse. The fixed-point equation
has to be arranged before that loss occurs.

\emph{The operator expansion as a datum.}
For each operator fix its tree-walk terms and weights once.
The first variable of a kernel term lies near its first cell;
the second lies near some cell of its tree. The local bound
is its original scale-weighted estimate times
\[
 \operatorname{wt}(\omega)
                   e^{-\delta d_{\mathbf B}(\omega;y,y')}.
\]
The weights of all trees with a fixed first cell sum to a
constant. A term is defined on the pair restricted to its
carrier, with the same estimate on every intermediate region
between the carrier and $\Omega_0$. A separate crude
holomorphic bound is used for convergence only. That bound
may depend on the level; the local weights and rates do not.

Partition into cubes of side $\ell_{\mathsf P}$ in local
units, where $\ell_{\mathsf P}$ is a multiple of $M_1$.
Let $\mathsf P(\omega)$ be the cubes meeting the uncut
hull of a tree. For $|s(\square)|\le e^{\bar\kappa}$ set
\[
 \mathsf O(s)=\sum_\omega s^\omega\mathsf O_\omega,
 \qquad s^\omega=\prod_{\square\in\mathsf P(\omega)}s(\square).
\]
Each parameter occurs with degree one. This convention
marks a cube once even when a walk returns to it. It is
what gives the mixed Cauchy estimate used to localize an
activity in a connected set of cubes.

A tree of weighted length $t$ meets at most
$16+Ct/\ell_{\mathsf P}$ such cubes. The constant
$16=2^4$ allows its first cell to lie at a corner of the
partition. Thus
\[
 |s^\omega|e^{-\delta d_{\mathbf B}(\omega;y,y')}
 \le e^{16\bar\kappa}
       e^{-(\delta-C\bar\kappa/\ell_{\mathsf P})
                                    d_{\mathbf B}(\omega;y,y')}.
\]
Choose $\ell_{\mathsf P}$ after $\bar\kappa$ so that
the remaining rate is a fixed positive fraction of
$\delta$. The exponential in the tree length, rather
than the reserved walk factor, pays for the parameters.
The sum and all its restricted partial sums therefore
have the same local estimates, with multiplier
$Ce^{16\bar\kappa}$. The factor at the first corner
does not disappear by making $\ell_{\mathsf P}$ larger.

If parameters vary only in a subfamily $\sigma_0$,
the same calculation leaves an additional factor
\[
 e^{-\bar\kappa\#(\mathsf P(\omega)\setminus\sigma_0)}.
\]
In particular forcing a prescribed set of cubes into
the hull gives exponential smallness in their number.
The hull is connected through walls. A nonzero term
therefore belongs to a single wall-connected component
of $\{s\ne0\}$. The operator splits exactly over these
components. This assertion uses the uncut hull, not
connectedness of its carrier inside a possibly rough
$\Omega_0$.

The same construction on a region $W$ keeps only walks
with $\mathsf P(\omega)\subset W$. Their sums use the
pair only on $W\cap\Omega_0$. Agreement of the data
for traces of regions on the hull gives agreement of
these partial sums in two different settings. This is
the form of consistency needed when a shift is first
formed on a window and then transported into a larger
box.

\emph{Which operator identities are lost.}
At $s=1$ and on shell,
\[
 Q\mathsf H=I,\qquad RD^*\mathsf H=0,\qquad
 \mathsf G^{-1}\mathsf H=Q^*(Q\mathsf GQ^*)^{-1}.
\]
They are identities between full operators. They are
not identities between their walk terms. If all
parameters near one row are zero, the decoupled
response has no term there, so $Q\mathsf H(s)$ is
zero there and cannot be the identity. Even at
$s=1$ the slice identity has the off-shell defect
described above.

The force in the localized response equation is
therefore formed as follows. First rewrite the
nonquadratic gradient at the genuine critical
background, using the exact inverse identities to
remove the linear expression $\Delta_\pi\mathsf H$, where
$\Delta_\pi=\Delta^{\mathrm W}_{\mathbf U}-\Delta'_\pi[\mathbf J]$.
Only then decouple the remaining local expression.
Denote that expression by $\mathsf W_s$. For general
$s$ and an independent-current pair this is its
definition; it is an actual action gradient only
at $s=1$ on shell. The localized force calculation
must establish its quadratic estimate on the
whole complex ball. Reusing the original gradient
and the false identity $Q\mathsf H(s)=I$ would
give a different equation.

Use admissible weights $\varpi$ satisfying
\[
 \varpi(y)\le e^{\theta d_{\mathbf B}(y,y')}\varpi(y')
\]
for $\theta$ below the retained rate, and define
\[
 \|A\|_\varpi=\sup_y\varpi(y)
              \max\{\xi_y|A|,\xi_y^2|\nabla_{\mathbf U}A|\},
 \qquad |f|^\sharp_\varpi=\sup_y\varpi(y)\xi_y^3|f|.
\]
The decoupled inverse maps the force norm into
the potential norm with norm at most
$\widehat b\le Ce^{16\bar\kappa}$. The required
local force estimate is
\[
 |\mathsf W_s(A_1)-\mathsf W_s(A_2)|^\sharp_\varpi
       \le C\widehat b^2r\|A_1-A_2\|_\varpi
       \quad(\|A_i\|\le r<c/\widehat b).
\]
The term containing the independent current
receives its factor $r$ from the Schwarz estimate
on the full complex ball. This is why the
estimate is needed before restricting to
real fields.

\emph{The response with the current correction included.}
Take walk data for the two compound operators
$\mathsf G_1DRD^*\mathsf G_1$ and
$D(1-R)D^*$. Define the expansion of
\[
 \mathfrak G=\mathsf G_1-\mathsf G_1DRD^*\mathsf G_1
                                      -\mathsf H_1Q\mathsf G_1
\]
from these data. The expression $\mathfrak G(s)$
is the decoupling of this full expansion. It
must not be replaced by multiplying independently
decoupled factors in the displayed algebraic
formula.

Write the nonlinear average as
$\mathcal Q=\mathsf mQ+\mathcal C$. The localized
system has two unknown corrections $X,\mathbf A$:
\[
 \begin{split}
 A_\bullet&=\mathsf H_1(s)\mathsf m^{-1}B'+\mathbf A,\\
 \mathbf A&=-\mathfrak G(s)\mathsf W_s(A_\bullet),\\
 X&=\mathcal C(A_\bullet-\mathsf H(s)\mathsf m^{-1}X),\\
 \mathbf H&=A_\bullet-\mathsf H(s)\mathsf m^{-1}X.
 \end{split}
\]
For fixed $A_\bullet$, the equation for $X$
is solved in a small ball using the quadratic
bound on $\mathcal C$. The other equation
is contractive provided $\sup|B'|$ is small
compared with $\widehat b^{-4}$. Indeed its
Lipschitz constant has four operator factors
in size:
\[
 \widehat b\cdot\widehat b^2\cdot
                       \widehat b\sup|B'|.
\]
They come from the inverse, the reduced-force
bound, and the size of the linear response.
Choosing the datum radius after these constants
gives one small holomorphic solution, bounded
by $C\widehat b\sup|B'|$.

The same contraction works in every admissible
weight. Setting
$\varpi=e^{\theta d_{\mathbf B}(\cdot,y_c)}$
therefore gives decay of the derivative with
respect to $B'(c)$. Splitting, covariance and
transport are proved by taking the separately
constructed solutions on each component as
a candidate for the full equation, and
using uniqueness in the same ball. They
are not consequences of a total norm bound
alone.

For arbitrary complex parameters the controlled
second-order object is a reduced current,
not the actual current of a critical point.
It is built from the multiplier, using the
pairing
\[
 \langle\mathrm d^*\mathrm d A,\mathsf He_c\rangle
                         =\langle\mathrm d A,\mathrm d\mathsf He_c\rangle.
\]
Only first differences occur. This avoids
claiming that first-order bounds for arbitrary
operator data imply a uniform bound on
$\Delta_{\mathbf U}\mathbf H$.

At $s=1$ over a real critical reference,
the constructed solution satisfies the
constraint and Landau condition, and its
Euler--Lagrange defect lies initially in
the fine gauge directions. Testing it
against the constraint-preserving gauge
directions removes that defect on a real
neighbourhood of zero. The identity theorem
then extends the three equations over the
whole connected complex datum polydisc:
\[
 \mathcal Q(\mathbf H)=B',\qquad
 RD^*\mathbf H=0,\qquad
 F'(\mathbf H)=K_{\mathbf H}^*\lambda.
\]
Only this undeformed on-shell identification
uses the actual critical orbit. A small real
neighbourhood used for the gauge test need
not have a radius uniform in the volume,
because no estimate is taken from that
radius. All quantitative bounds were
proved on the fixed complex ball first.

\emph{A reduced system and its exact scope.}
There is a second useful formulation,
with unknowns the response $T$ and its
multiplier. Use angle rows
$\widehat Q=\mathsf mQ$ with their unweighted
row pairing, and set
\[
 H=\mathsf H\mathsf m^{-1},\quad
 S=(\widehat Q\mathsf G\widehat Q^*)^{-1},\quad
 H_0=\mathsf G_0\widehat Q^*
                 (\widehat Q\mathsf G_0\widehat Q^*)^{-1},\quad
 \widetilde G=\mathsf G_0-H_0\widehat Q\mathsf G_0.
\]
Let $\mathsf a_0$ denote the diagonal
row mass in this convention, and put
$\mathcal N(T)=F'(T)-F'(0)-\Delta^{\rm W}_{\mathbf U}T$.
The reduced equations are
\[
 \begin{split}
 T={}&H_0(s)(B'-\mathcal C(T))
       -\widetilde G(s)[\mathcal N(T)-\mathcal C'_T{}^*\lambda],\\
 \lambda={}&H(s)^*\mathbf J
      +(S(s)-\mathsf a_0)(B'-\mathcal C(T))
       +H(s)^*[\mathcal N(T)-\mathcal C'_T{}^*\lambda].
 \end{split}
\]
Each of the four operators has its own
fixed walk data. Their exact identities
are again used only at $s=1$.

Center the multiplier at $H(s)^*\mathbf J$.
Most terms are quadratic on a ball of
radius $2\widehat b r$ and give the
condition $\widehat b^2r\ll1$.
There are also linear terms containing
the current, notably the derivative of
$\mathcal C'_T{}^*H(s)^*\mathbf J$.
Their norm is $C\widehat b^2\gamma_0$.
Thus this formulation requires a second
current ceiling chosen after
$\bar\kappa$. With that ceiling it
is a contraction in the weighted norm
of the pair, on a datum ball of size
$c\widehat b^{-2}$.

The response splits over the parameter
components. The multiplier itself does
not: outside them it contains the
diagonal term $-\mathsf a_0B'$.
For the same reason the force is
split only after this term is retained.
Further seminorms pass from $H_0(s)$
and $\widetilde G(s)$ to the solution;
they need not be bounds for $H(s)$.

At $s=1$ over a critical real pair,
the two response systems agree on
their common small ball by uniqueness
of the constrained Euler--Lagrange
solution. Their separate force bounds
then give both $\mathrm d^*\mathrm d\mathbf H$ and
$DD^*\mathbf H$ at that value of the
parameters. At a general off-shell
pair or at $s\ne1$ they are different
functions. In particular no exact
constraint is asserted for the first
system at arbitrary parameters.
At $s=1$ the reduced system enforces
$\mathcal Q(T)=B'$ even off shell,
by $\widehat QH_0=I$ and
$\widehat Q\widetilde G=0$. The
first system can instead retain the
defect $\mathsf mQ\mathbf A$ there.

For fields given only on a region,
the first system is formed by keeping
the operator terms whose hulls lie
in that region. The local reduced
system is used on the strip sequences
for which its four partial sums
have been constructed. It is not
silently extended to all determining
structures. This distinction is
needed when the large-field operation
localizes a shift in a boundary strip.

\subsection{The fluctuation integral and the local expansion}

Write $W=V_{k+1}$ and let $V_0=\mathcal M^kU_{k+1}(W)$ be the
level-$k$ background. In this paragraph all bond potentials are
dimensionless level-$k$ variables. For every coarse bond $c$, choose
the middle bond $b_0(c)$ among the $L$ fine bonds on $c$. The oddness
of $L$ makes this choice commute with the isometries of the coarse
lattice. Let $\mathfrak X$ be the set of the remaining fine bonds.
A field on $\mathfrak X$ will always be extended by zero at the
pivots before an operator on all bonds is applied to it.

The relative constraint has the expansion
\[
 q(B'):=\frac1i\log\big[\mathcal M(e^{iB'}V_0)W^{-1}\big]
       =LQ_{V_0}B'+q_{\ge2}(B').
\]
The coefficient of the pivot variable in its linear part is an
invertible colour matrix $L\mathsf m_c$. This matrix is a small
perturbation of the corresponding coefficient at the identity.
It is uniformly invertible for the real regular backgrounds under
consideration. Denote by $\mathsf h$ the right inverse supported
on the pivots, so that $LQ_{V_0}\mathsf h=1$. If $E$ is extension
by zero from $\mathfrak X$, the map
\[
 \mathsf C=(1-\mathsf hLQ_{V_0})E
\]
identifies the free variables with the kernel of the linearised
constraint. Its restriction to $\mathfrak X$ is the identity.
In particular the use of free variables does not introduce a new
constraint on their values.

For a real scale $g=\mathsf t^{-1/2}$, solve
\begin{equation}\label{eq:deep-pivot-fixed-point}
 D=q_{\ge2}(g\mathsf CB-\mathsf hD),\qquad
 B'=g\mathsf CB-\mathsf hD.
\end{equation}
The derivative of $q_{\ge2}$ vanishes at zero. On a sufficiently
small fixed ball the first map is therefore a contraction in $D$.
Its solution is holomorphic in the background and $gB$, and
$|D|\le C_Lg^2|B|^2$. Applying $q$ to the second formula gives
$-D+q_{\ge2}(B')=0$. Thus
$V_W(B):=e^{iB'}V_0$ lies on the exact fibre, rather than on its
tangent approximation. All these assertions use the uniform
relative-average estimates established before the fluctuation
step. They do not follow from a formal expansion of the constraint.

The fluctuation test is imposed only on $\mathfrak X$:
\[
 |B(b)|<\mathfrak b(\mathsf t),\qquad
 \mathfrak b(\mathsf t)=\sqrt{\mathsf t}\,\delta(\mathsf t).
\]
The eliminated variables satisfy $|B'|\le C_L\delta(\mathsf t)$,
with $C_L$ at most a fixed multiple of $L^4$. This weaker bound
is enough. The block contour variables depend on bonds inside
their block, and the same comparison with $V_0$ bounds their
deviation by $C_L\delta(\mathsf t)$. Choosing
$\varepsilon(\mathsf t)/\delta(\mathsf t)\ge C_L$ makes every
gauge-fixing indicator equal to one on this fibre chart. The
smooth gauge-fixing weight remains in the integral. Its indicators
are not analytically continued.

There are two curved constraints in the quadratic action. The
background $U_k(V)$ is critical under the $k$-fold average, and
$V$ is restricted by the next average. For a constrained critical
point with $DF=\lambda Dq$, differentiating $q(V(z))=q(V(0))$
twice gives the restricted Hessian
\[
 D^2F[V'(0),V'(0)]-\lambda D^2q[V'(0),V'(0)].
\]
The second term cannot be omitted. For the first constraint its
multiplier is represented by the current of the fine background;
for the second it pairs with the quadratic remainder of the
one-step average. This explains both occurrences of the current
in $\Delta^{(k)}$. At $k=0$ the background map is the identity,
so the first of these corrections is absent.

Criticality removes the first derivative on the fibre. The
gauge-invariant action Hessian vanishes on its residual gauge
directions, including its mixed pairing with any tangent direction.
Split a tangent vector into a vector in the kernel of the
linearised gauge fixing and an orbit vector generated by
$\lambda$. The action acts on the first vector and the quadratic
gauge fixing acts on $\lambda$. With the normalisation of the
gauge-fixing action this gives
\[
 \langle B,\Delta^{(k)}B\rangle
     =\Delta_k(B'',B'')+\sum_x|\lambda(x)|^2.
\]
Positivity on the first summand is the constrained propagator
estimate. The second summand supplies the missing gauge mass.
Thus the free-bond form
$\mathsf K=\mathsf C^{\mathsf T}\Delta^{(k)}\mathsf C$ is
positive at real data. This positivity concerns the full
fluctuation form. Retaining selected local pieces in a large-field
operation requires the separate comparison in
\cref{sec:large-fields}.

Write $\mathrm{GF}_k$ for the soft gauge-fixing action,
whose weight at the reference inverse coupling $\mathsf t$ is
$e^{-\mathsf t\mathrm{GF}_k}$. Taylor expansion on the exact chart now gives
\[
 \mathsf t\{A(U_k(V_W(B)))+\mathrm{GF}_k(V_W(B))\}
 =\mathsf tA(U_{k+1})+\tfrac12\langle B,\mathsf KB\rangle
       +\text{terms of order at least three in }gB.
\]
The Jacobian is kept at the same precision as the action.
Disintegrating the group delta function leaves
$L^{-3}|\det\mathsf m_c|^{-1}$ at each coarse bond.
The determinants $|\det\mathsf m_c|^{-1}$ are put in the
Gaussian normaliser. The logarithm of the nonlinear Jacobian
and the Haar density in the exponential chart are put in
$\mathbb P^J$. The action and gauge-fixing terms of order above
two are denoted $\mathbb P^A$ and $\mathbb P^G$. Consequently
\[
 \mathbb P^{(k)}=\mathbb P^A+\mathbb P^G+\mathbb P^J.
\]
This division by origin will matter when the action weight is
complex while the gauge fixing remains real.

The constants can be counted without an error per bare bond.
Put $\mathsf V_j=|\mathbb T^{(j)}|$, and let $\sigma_0$ be the Haar
density at the identity in the chosen exponential coordinates.
There are $4(\mathsf V_k-\mathsf V_{k+1})$ free bonds and $\mathsf V_k-\mathsf V_{k+1}$
independent residual gauge sites. If
\[
 \mathfrak z(\mathsf t,\epsilon)
 =\int_{\SU(2)}1_{\{|u-1|<\epsilon\}}
                    e^{-\mathsf t(1-\frac12\tr u)}\,du,
\]
the field-independent factor is
\begin{equation}\label{eq:deep-fibre-number}
 \mathsf c_k(\mathsf t)
 =\mathfrak z(\mathsf t,\varepsilon(\mathsf t))^{-(\mathsf V_k-\mathsf V_{k+1})}
   (\sigma_0\mathsf t^{-3/2})^{4(\mathsf V_k-\mathsf V_{k+1})}
                  L^{-12\mathsf V_{k+1}}.
\end{equation}
The three factors come respectively from gauge averaging, scaling
the free Haar variables, and the coarse delta functions. At real
data the determinants in this calculation have their positive
values. Their holomorphic continuations are fixed from those
values; no absolute-value operation is applied at a complex pair.

For $\mathfrak b(\mathsf t)\ge1$, comparison with the quadratic
action near $u=1$ gives
$\log\mathfrak z=-\frac32\log\mathsf t+O(1)$. Hence
\[
 \log\mathsf c_k(\mathsf t)
       =-\tfrac92(\mathsf V_k-\mathsf V_{k+1})\log\mathsf t+O(\mathsf V_k).
\]
For real $W$ the full Gaussian normaliser is
\[
 Z^{(k)}(W)
  =(2\pi)^{\frac32|\mathfrak X|}
      \det(\mathsf K(W))^{-1/2}
           \prod_c|\det\mathsf m_c(W)|^{-1}.
\]
Its complex continuation is the holomorphic branch fixed by
this real formula. It has no coupling dependence at the trivial
background. The
action-dependent part of one step is therefore separated from
its numerical normalisation before the local expansion begins.

The small-field integral therefore has the form
\begin{equation}\label{eq:rg-fluctuation-integral}
 \mathsf c_k e^{-x_kA(U_{k+1})+\mathbf E_k(U_{k+1})}
 Z^{(k)}(U_{k+1})
 \int d\mu_{C^{(k)}}(B)\,\chi^{(k)}
 e^{\mathbb P^{(k)}(B)+
       \mathbf E_k(U_k(V_k))-\mathbf E_k(U_{k+1})}.
\end{equation}
The Gaussian normaliser is $Z^{(k)}$; $\mathsf c_k$ contains the
Jacobian, Haar, delta-function and gauge-fixing normalisations.
On the free-bond variables the covariance is
$C^{(k)}=(\mathsf C^{\mathsf T}\Delta^{(k)}\mathsf C)^{-1}$,
as determined by the preceding quadratic form.
Here $\chi^{(k)}$ is the product of the fluctuation indicators on
the free bonds. The factorisation defines $\mathbb P^{(k)}$ as the
remaining fluctuation exponent after the displayed quadratic form
has been put into the Gaussian measure.
The exponent $\mathbf E_k$ is the accumulated regular interaction
with its constant and quadratic action parts subtracted; its precise
form is \eqref{eq:rg-old-interaction} below.

The local terms are indexed by connected unions $X$ of $M$-cubes and
by a site $z\in X$. Let $d_j(X)$ be the length of a shortest tree
meeting these cubes, divided by $M$. Their norm is
\[
 \|E^{(j)}\|_{j,\kappa}
 =\sup_{X,z}e^{\kappa d_j(X)}
                  \sup_{\text{analytic domain}}|E^{(j)}(X,\cdot,z)|.
\]
The number of domains through a fixed cube with $d_j(X)\le n$ is
at most $e^{C(n+1)}$. Thus a sufficiently large $\kappa$ makes
$\sum_{X\supset\square}e^{-\kappa d_j(X)}$ finite uniformly in
the lattice. This elementary count is used throughout the local
expansion.

To localise the exponent, attach a parameter $s(\Delta)$ to each
cube traversed by a walk. Allowing
$|s(\Delta)|\le e^{\bar\kappa}$, with
$\bar\kappa\ge CL\kappa$, leaves the background equations
contractive when $M$ is sufficiently large. Long walks pay for these
parameters by their length, since meeting $m$ cubes costs at least
$cmM$ in the exponent. The fixed cost of short walks is absorbed by
the small fluctuation amplitude. Cauchy estimates then give decay
for the resulting local potentials.

Two small quantities occur. If the fluctuation test is
$|B|<\delta_k/g_k$, then the cubic interaction is bounded by
$C\delta_k|B|^2$. It is small under the Gaussian weight. Old
interactions are evaluated at a shifted background. Their smallness
instead contains the ratio
$\mathsf s_k=\delta_k/\alpha_{0,k}$ of that shift to its
analyticity radius. The running radius satisfies
$\alpha_{0,k}\asymp g_k(\log x_k)^q$, with $q$ chosen large
after the threshold exponents. Here $\mathfrak p_0$ is the exponent
of the threshold $p_0(g)$ in \cref{sec:large-fields}. In particular,
$q>3\mathfrak p_0+C\mathfrak r$, where $\mathfrak r$ is the
fixed exponent of the large-field cube size and $C$ covers the
geometric losses in localisation. We also take $q>21\mathfrak p_0$
for the two-point comparison. The coupling ceiling is chosen
after $q$. On running radii this ratio is only a negative
power of $\log x_k$; it should not be replaced by a power of $g_k$.
After coupling subtraction, an old term of age $n$ has the summable
weight
$\omega_n\asymp n^{3/2}L^{-\alpha_{\mathrm H}n}$.
Writing $N_\omega$ for the corresponding weighted sum of old norms,
the local old-field potential is bounded by
$C\mathsf s_kN_\omega e^{-c\kappa d_k(Y)}$.
The construction gives $N_\omega\le CE_0$; it is
$\mathsf s_kN_\omega$, rather than $N_\omega$, that is small.
Remainders left by the large-field operation have no coupling
counterterm and receive no such decreasing age weight. Once their
cubes fit in single blocks, they enter a separate part of
$N_\omega$, weighted by the square of the ratio of radii. Their
summation uses their own size $g_j^{\kappa_0}$ and the lower
flow bound. Younger remainders remain separate, as in
\cref{sec:large-fields}.

The identity for the old interaction is a difference along the
background response, integrated from the unperturbed background to
the fluctuating one. The potential and current are transported
together. At the actual endpoints they satisfy the field equation,
but the affine differentiations and decoupling interpolations between
them need not. This is why the interactions are analytic functions of
the independent pair $(\mathbf U,\mathbf J)$.
The pointwise coupling subtraction described next is performed before
the decoupling; the surviving brackets are rooted at the original
domain. Their sum recovers exactly the difference of the old
interaction in \eqref{eq:rg-fluctuation-integral}.

For completeness, the subtraction and the local interpolation
are performed in the following order. For a term of level $j$,
write the response curve as a curve of pairs
\[
 p_j(A,\mathcal J)=\big(e^{i\eta_jA}U_{k+1},
         L^{-3(k+1-j)}J(U_{k+1})+\mathcal J\big),
\]
where $\mathcal J$ is the independent current displacement.
At the actual response the
potential and current displacements are related by the critical
equation. In applying the fundamental theorem of calculus they
are differentiated as independent coordinates. A partition of
unity in the displacement localises its directions to windows.
The response of the whole lattice is next replaced by that of
a box, with their difference written as a second integral of a
derivative. Finally the walks in the resulting box response are
decoupled. The two removed differences and all differentiated
directions therefore require the off-shell pair domain.

An equivalent formula starts with the original domain as a root.
If $h_Y(s;B)$ is its localised response and $F$ is the activity
with its pointed counterterm already removed, set
\[
 \Phi_Y[F](s,\nu)=F(p_j(\nu h_Y(s;B))),
\]
where $h_Y$ contains the potential and current components.
For a finite set $K$ of decoupling cubes, the interpolation
bracket integrates $\prod_{\Delta\in K}\partial_{s(\Delta)}$
over the corresponding unit cube of parameters, with the
retained cubes fixed to one. Repeated one-variable fundamental
theorems of calculus recover the original function. A connected
piece that does not meet its root vanishes in the rooted
decomposition. Cauchy estimates in $s$ pay for the new cubes;
the Schwarz estimate in $\nu$ pays for the response amplitude.
These are algebraic sum rules before taking absolute values.
The analytic estimate to check is that the decoupled responses stay
in their independent field/current domains on one common polydisc,
uniformly in the old level. The algebraic identities alone do not
give that domain control.

Here is the bracket identity used in that construction.
For disjoint finite cube sets $\mathcal K,\mathcal K'$,
define
\begin{equation}\label{eq:xrg-bracket}
 \Psi[\mathcal K\mid\mathcal K']
  =\int_{[0,1]^{\mathcal K}}
       \prod_{\Delta\in\mathcal K}\partial_{s(\Delta)}
          \Psi(s+1_{\mathcal K'})\,ds.
\end{equation}
If $\mathcal K''$ is disjoint from the base
$\mathcal K'$, repeated fundamental theorems of
calculus give the exact resummation
\begin{equation}\label{eq:xrg-bracket-sum}
 \Psi(1_{\mathcal K'\cup\mathcal K''})
      =\sum_{\mathcal K\subset\mathcal K''}
                      \Psi[\mathcal K\mid\mathcal K'].
\end{equation}
The bracket is applied to the decoupled response
function $\Phi_Y[F]$, rather than to an arbitrary
function for which no locality has been established.
At a zero decoupling parameter no walk passes
through its cube. A connected differentiated
component separated from the base therefore
does not affect this function, and its bracket
vanishes. This is the rooting assertion.

The complex parameter disc has radius
$e^{\bar\kappa}$. A derivative in a newly
added cube pays $Ce^{-\bar\kappa}$.
The response multiplier $\nu$ is allowed up
to the ratio of analytic radius to response
size. Subtracting its value at zero and
applying Schwarz gives the independent factor
$\mathsf s_k$. Both estimates apply to a
single holomorphic function on a common
polydisc. They do not follow by separately
bounding the response at real values zero
and one.

There are two exact ways to organise the
old-field identity. In the window organisation,
the main response curve is generated by the
datum
\[
 B_j(t)=\mathcal Q_j\big(U_{k+1};
                        \eta_{k-j}t\mathbf H_k\big),
                 \qquad 0\le t\le1.
\]
It is on shell when the decoupling parameters
are one. Its endpoints are the pair of
$U_{k+1}$ and, up to gauge, that of $U_k(V_k)$.
Differentiate along it, resolve the direction
with $\sum_\square\zeta_\square=1$, and
replace the whole-lattice response by a box
response near the window. This produces two
removed integrals. The differentiating
directions and the first removed straight
segment are affine in their parameter and
can leave the shell. Decoupling is performed
after these replacements.

In the rooted organisation the bracket
\eqref{eq:xrg-bracket} uses the domain of
the term as base. Its counterterm is removed
at each base point before interpolation.
Both organisations recover exactly
\begin{equation}\label{eq:xrg-old-sum-rule}
 \sum_Y\mathsf W_k^{\mathrm{old}}(Y;B)
 = (\mathbf E_k+\mathbf R'_k)(U_k(V_k))
                  -(\mathbf E_k+\mathbf R'_k)(U_{k+1}),
\end{equation}
where $\mathbf R'_k$ denotes the aged remainder
sum. The younger remainder sum is passed
separately to the next remainder slot.

The rooted form uses a stronger piece of
information than invariance of the total
interaction: the pointed quadratic tensor
must cancel at the centre of every window.
Once that estimate is available, one base
point of level $j$ contributes at most
\[
 C\varrho_j^3\eta_{k-j}^{4+\alpha_{\mathrm H}}
             \|E^{(j)}\|_{j,\kappa},
 \qquad \varrho_j\le C(k+2-j)^{1/2}.
\]
A current-scale window has $C\eta_{k-j}^{-4}$
base points. The factor $\varrho_j^3$ is the
cost of up to three Cauchy estimates on the
older, possibly smaller radius. After
cancelling the site count, the contribution
is
$C\mathsf s_k\omega_{k+1-j}\|E^{(j)}\|_{j,\kappa}$.
An aged remainder lacks the quadratic
counterterm, and its analogous contribution
has $\varrho_j^2$ in place of $\omega_{k+1-j}$.
Thus the actual old-data load is
\begin{equation}\label{eq:xrg-old-load}
 N_\omega=
 \sum_{j\le k}\omega_{k+1-j}\|E^{(j)}\|_{j,\kappa}
       +\sum_{j\text{ aged}}\varrho_j^2
                              \|R^{(j)}\|_{j,\kappa}.
\end{equation}
The first sum uses H\"older decay and the
second uses the lower coupling drift.
This distinction survives every later
localisation of the old interaction.

\subsubsection*{The local gap and the conditional Gaussian integral.}
The Gaussian localisation uses the form $\mathsf K$ itself as
well as its inverse. The propagator construction supplies a
local analytic expansion
\[
 \mathsf K(\mathbf p)=\sum_\omega\mathsf K_\omega(\mathbf p),
 \qquad
 |\mathsf K_\omega|
 \le C(CM_1^{-1/2})^{|\omega|}e^{-c\,\mathrm{length}(\omega)}.
\]
The terms are covariant under real gauge transformations and
agree in two regions whenever their carriers and the required
background data agree there. This is a property of the actual
operator construction. A decay bound for its full kernel alone
would not define local terms with this agreement property.

Let $\gamma_0$ be a lower bound for $\mathsf K(1,0)$, and
first consider a real regular pair. On cubes of a fixed side
$\ell_0$, choose the real gauges contained in the analytic-domain
conditions. In each such gauge $U=e^{i\eta A}$ has
$|A|+|\nabla A|\le a_{\rm w}$. The local walk expansion can be
compared with its value at the identity. Schwarz's lemma on a
fixed analytic ball bounds this difference by $Ca_{\rm w}$.
For a complex pair $(e^{i\eta A'}U,\mathbf J)$, a second use
of Schwarz on
\[
 z\longmapsto (e^{iz\eta A'}U,z\mathbf J)
\]
adds $C(b_0+b_1)$ to the bound. Both comparisons use the
independent current coordinate. They do not replace it by the
current of a general complex perturbation.

Choose real cutoffs $h_i$ with $\sum_i h_i^2=1$ and
$|\partial h_i|\le C\ell_0^{-1}$. The exact localisation identity
is
\[
 \mathsf K=\sum_i h_i\mathsf Kh_i
                  +\tfrac12\sum_i[h_i,[h_i,\mathsf K]].
\]
The twice-commuted walk kernel has two factors of cutoff
variation. Its operator norm is bounded by $C\ell_0^{-2}$,
using the exponentially weighted second moment of the kernel.
Walks leaving the gauge cube contribute $Ce^{-c\ell_0}$.
Thus the real quadratic form has lower bound
\begin{equation}\label{eq:deep-full-gap}
 \mathsf K(U,0)\ge
 [\gamma_0-C\ell_0^{-2}-Ce^{-c\ell_0}-Ca_{\rm w}]\,1.
\end{equation}
The complex perturbation has operator norm at most
$C(b_0+b_1)$. Fix $\ell_0$, then $a_{\rm w}$, then $M$ large
enough to contain these cubes, and finally the radii. This leaves
a fixed real gap and a smaller fixed complex perturbation margin.
All full compressions $1_Y\mathsf K(U,0)1_Y$ inherit the real
gap. At complex pairs their inverses are obtained by a norm
Neumann series about the real compression.

The same argument gives decay for a compressed resolvent.
Conjugation by $e^{\zeta\rho}$, with $\rho$ a Lipschitz distance,
changes the form by $O(|\zeta|)$ in the weighted operator norm.
For $|\zeta|$ small relative to the gap,
\[
 (y+\mathsf K_Y)^{-1},\qquad y\ge0,
\]
therefore has an exponentially decaying kernel, with uniform
constants. Here and below $\mathsf K_Y$ is the compressed full
form, not an arbitrary truncation of its walk terms. Integrating
these resolvents in
\begin{equation}\label{eq:deep-inverse-root}
 \mathsf K_Y^{-1/2}
      =\frac1\pi\int_0^\infty
                   y^{-1/2}(y+\mathsf K_Y)^{-1}\,dy
\end{equation}
defines the inverse square root. The branch is the positive
one at real data. The gap gives convergence near zero, and the
resolvent bound gives convergence at infinity. Local resolvent
expansions give the corresponding walk expansion of this root.

The decoupling parameters are inserted in these operator
expansions. A carrier meeting $m$ cubes gains at most
$e^{m\bar\kappa}$. Apart from a bounded number of cubes at its
starting point, each additional cube requires length of order
$M$. Thus
\[
 e^{m\bar\kappa}
 \le e^{C\bar\kappa}
       e^{(C\bar\kappa/M)\mathrm{length}(\omega)}.
\]
Taking $M\ge C\bar\kappa$ leaves a fixed fraction of the walk
decay. The short-walk factor $e^{C\bar\kappa}$ is a fixed
constant; its smallness is supplied by the fluctuation amplitude
in the response equations, rather than by increasing the depth.

For conditioning on a region $Z_0$, parameters are put only on
cubes outside a collar of $Z_0$. A carrier reaching a parameter
from $Z_0$ must cross this collar. Its norm contains
$e^{C\bar\kappa}e^{-cM}$, which is exponentially small after
the preceding choice. This retains the real reference rows near
the boundary of $Z_0$. Setting parameters to zero on a boundary
without retaining this collar could instead remove whole rows
and destroy invertibility.

Let $P=1_{Z_0}$. Write
\[
 C=C(Z_0;s)=(P\mathsf K(s)P)^{-1},\qquad
 \Gamma=P\mathsf K(s)(1-P)r(s).
\]
Here $r(s)=(\mathsf K^{-1/2})(s)$ is the separately decoupled
walk series of the inverse square root. It is not the inverse
square root of $\mathsf K(s)$. At $s=1$ the two original
operators are recovered, so $r(1)=\mathsf K^{-1/2}$.
For functions depending on the bonds of $Z_0$, define
\begin{equation}\label{eq:deep-conditional-functional}
 \begin{split}
 \mathcal J(Z_0;s)[F]
 =\int d\mu_1(X)\,
      e^{-\frac12\langle\Gamma X,C\Gamma X\rangle}
       \int d\mu_C(B)\,
                    e^{-\langle B,\Gamma X\rangle}F(B).
 \end{split}
\end{equation}
The white noise $X$ has a product measure. At real data,
completing the inner square gives a Gaussian of mean
$-C\Gamma X$ and covariance $C$. The outer quadratic factor
cancels the normalising factor produced by that square. At
$s=1$ this is exactly the marginal of the original Gaussian,
as checked in \eqref{eq:mrg-covariance-recovery} below.
In particular $\mathcal J(Z_0;s)[1]=1$ at real admissible
decoupling data, and the identity continues on the common
admissible complex domain. For complex $C$, the symbol $d\mu_C$ means
the analytic Gaussian density with its continued determinant;
it is not a positive measure.

The estimate on this functional is local in the region where
parameters or test factors occur. If $Z_0^\sharp\subset Z$ and
$s=0$ outside $Z$, all white-noise variables outside $Z$ have
integral one. Compare $C,\Gamma$ with their real reference
values, whose gaps and operator norms are uniform. Suppose
their allowed perturbations have norm at most $\hat\epsilon$.
For
\[
 |F(B)|\le F_0e^{\rho|B|^2},\qquad
 u=\hat\epsilon+\rho\le c,
\]
the real reference square completion reduces the absolute
integral to a Gaussian exponential moment with quadratic
perturbation bounded by $Cu$. A reference Gaussian with
covariance norm at most $C_0$ satisfies
\[
 \int e^{v|Y|^2}\,d\mu_{C_0'}(Y)
       =\det(1-2vC_0')^{-1/2}
       \le e^{Cv n},\qquad 2v\|C_0'\|\le\tfrac12,
\]
where $n$ is its number of real components. The determinant
normalisations in the comparison obey the same estimate,
by $|\log\det(1+T)|\le Cn\|T\|$ for $\|T\|\le1/2$.
It follows that
\begin{equation}\label{eq:deep-local-gaussian-bound}
 |\mathcal J(Z_0;s)[F]|
       \le F_0\exp[CuN(Z)],\qquad
 N(Z)\le C(LM)^4\#Z.
\end{equation}
The reference square completion includes the mixed $B,X$
terms; they cannot be bounded separately at their original
coefficients without spending the Gaussian gap. The perturbation
bound above is applied after that completion.

The factor in \eqref{eq:deep-local-gaussian-bound} is bounded
per coarse cube only when $u\le c(LM)^{-4}$. This condition
applies to every contribution to $u$: the walk-decoupling error,
the analytic radii, a complex change of the action coefficient,
the quadratic growth of interaction potentials, and the Gaussian
tail parameter. Its dimensional factor is part of the argument.
Using a bound exponential in the entire torus volume at this
point would not give a local polymer estimate. Verifying the
stated operator perturbations on the actual common pair domains
remains part of the uniform one-step construction.

The covariance recovery can be checked directly at real data.
Let $P=1_{Z_0}$ and split the positive matrix $\mathsf K$ as
\[
 \mathsf K=\begin{pmatrix}A&B\\B^{\mathsf T}&D\end{pmatrix}
 \quad\text{on }P\oplus(1-P),\qquad C_Z=A^{-1}.
\]
For a fixed exterior value $y$, completing the square gives
\begin{equation}\label{eq:mrg-gaussian-square}
 \begin{split}
 \langle(x,y),\mathsf K(x,y)\rangle
  ={}&\langle x+A^{-1}By,
                       A(x+A^{-1}By)\rangle\\
   &+\langle y,(D-B^{\mathsf T}A^{-1}B)y\rangle.
 \end{split}
\end{equation}
Thus the conditional mean is $-A^{-1}By$, and the conditional
covariance is $C_Z$. Draw the exterior value as
$y=(1-P)\mathsf K^{-1/2}X$, with $X$ standard Gaussian.
Then $By=\Gamma X$. The law of total covariance gives
\begin{equation}\label{eq:mrg-covariance-recovery}
 C_Z+C_Z\Gamma\Gamma^{\mathsf T}C_Z
                      =P\mathsf K^{-1}P.
\end{equation}
This is the identity used when the decoupling parameters return
to one. It also fixes the sign of the conditional mean.
At complex parameters the identity continues algebraically on
the component where the matrices remain invertible; positivity
and probability language apply only at real data.

The Gaussian perturbation bound has a volume restriction. If
one integration cube has $O((LM)^4)$ field components, the
allowed change per component is $O((LM)^{-4})$.
The discarded tail must still obey
\begin{equation}\label{eq:mrg-tail-room}
 (\delta_k/g_k)^2
       \ge C(LM)^4(\kappa+\log M).
\end{equation}
One half of its coercive exponent is available for this tail,
and the other half remains under the Gaussian integral.
Both \eqref{eq:mrg-tail-room} and the smallness needed for
short decoupled walks are obtained by choosing the coupling
ceiling after $L,M,\kappa$.

\subsubsection*{From local potentials to a connected logarithm.}
We spell out the expansion because its local Gaussian bound, its
connected summation and its pointed sum rule are used again in the
source argument. Let $\mathfrak V$ denote the finite family of
localised potentials at one step. Besides the small-box potentials,
which are defined only where their field tests hold, it can contain
the quadratic gauge potential that occurs when the action coefficient
is varied. That potential is defined for all real $B$. Write
\[
 |V(Y;B)|\le w_Y\left(1+M^{-4}
                         \sum_{b\subset Y}|B(b)|^2\right),
 \qquad w_Y\le w\,e^{-\lambda_V d_k(Y)}.
\]
The sum over $Y$ is absolutely convergent after restriction to a
finite torus. The uniform estimates will be independent of that
restriction. In the present application
\[
 w\le C\{\delta_k+\mathsf s_kN_\omega\}
\]
for the interaction part; the extra quadratic gauge potential is
bounded separately by $C|1-\tau|M^4$.

First use the finite identity
\begin{equation}\label{eq:deep-mayer-identity}
 \exp\Big(\sum_YV(Y;B)\Big)
    =\sum_{\mathcal D}\prod_{Y\in\mathcal D}(e^{V(Y;B)}-1).
\end{equation}
Here $\mathcal D$ ranges over subsets of the potential labels;
two different origins on the same domain retain different labels.
For a chosen $\mathcal D$, let $Y_0$ be the union of the domains
of its small-box potentials. Their field tests remain in place.
The tests outside $Y_0$ may be expanded as
\begin{equation}\label{eq:deep-test-expansion}
 \prod_b\chi_b
   =\prod_{b\in Y_0}\chi_b
       \sum_{P\subset\mathfrak X\setminus Y_0}
                                (-1)^{|P|}\prod_{b\in P}\chi_b^c.
\end{equation}
This order is necessary: expanding a test required by a potential
would evaluate that potential outside its domain. The quadratic
gauge potential contributes no test to $Y_0$.

An index $\iota=(\mathcal D,P)$ has a support $Z_0(\iota)$,
formed from its potential domains and failed-test bonds. Apply the
conditional representation
\eqref{eq:deep-conditional-functional} to its product $F_\iota$.
Then interpolate all coarse-cube parameters outside the collar
$Z_0^\sharp$. For a region $Z\supset Z_0^\sharp$ the term is
\begin{equation}\label{eq:deep-polymer-term}
 T(Z,\iota)=
 \int_{[0,1]^{Z\setminus Z_0^\sharp}}
   \prod_{\Delta\in Z\setminus Z_0^\sharp}\partial_{s(\Delta)}
       \mathcal J(Z_0;s1_Z)[F_\iota]\,ds.
\end{equation}
The parameters in the collar are one. Completeness of the
conditional representation and repeated fundamental theorems of
calculus show that the sum of these terms is the original
integral. The statement is an identity before estimates.

At zero parameters, the operators, conditional integrals and
interactions split over wall-components. Therefore
$T(Z,\iota)$ splits over those same components. Regrouping its
components gives a gas with polymers equal to connected unions
of coarse cubes. Two polymers are compatible when they share
neither a cube nor a wall. Denote their activities by $H(Z)$.
The exact integral is
\[
 \Xi(H)=\sum_{\{Z_1,\ldots,Z_n\}\ \mathrm{compatible}}
                                  \prod_{i=1}^nH(Z_i),
\]
including the empty family with value one. A factorisation into
components here is justified by the decoupled operator structure,
not by decay of correlations for the original Gaussian.

For each Mayer factor, introduce a numerical parameter $t_Y\ge1$.
The elementary inequality
\[
 |e^v-1|\le t_Y^{-1}e^{2t_Y|v|}
\]
follows from $|e^v-1|\le |v|e^{|v|}$ and
$t_Y|v|\le e^{t_Y|v|}$. Taking $t_Y$ proportional to the
inverse small activity weight extracts one small factor for
each vertex. The remaining exponential has a quadratic growth
coefficient that can be made smaller than the local Gaussian
margin. A failed test is bounded by
\[
 1_{\{|B(b)|\ge\mathfrak b\}}
       \le e^{-\gamma_2\mathfrak b^2/2}
                         e^{\gamma_2|B(b)|^2/2}.
\]
Its second factor uses the same quadratic margin; its first is
the tail gain. A differentiated decoupling parameter contributes
$C e^{-\bar\kappa}$ by Cauchy on the common complex parameter
disc. These three gains are taken before applying
\eqref{eq:deep-local-gaussian-bound}.

The local Gaussian cost is at most $e^{C\#Z}$ after imposing
the per-component restriction $u\le c(LM)^{-4}$.
The tail gain is useful uniformly only if
\[
 e^{-\gamma_2\mathfrak b^2/2}(LM)^4e^{C\kappa}\ll1.
\]
Since $\gamma_2$ itself is of order $(LM)^{-4}$, this is the
floor \eqref{eq:mrg-tail-room}. Half of the coercive exponent
may be assigned to the failed test; the other half stays under
the integral. No second use of the same half is made in the
Mayer estimate.

Two geometric inequalities organise the sums. For a wall-connected
union of carriers,
\begin{equation}\label{eq:deep-carrier-length}
 d\Big(\bigcup_iX_i\Big)+5
                    \le\sum_i(d(X_i)+5).
\end{equation}
For the collar of a fine-scale carrier,
\begin{equation}\label{eq:deep-collar-scale}
 d_{k+1}(Y^\sharp)\le (C/L)d_k(Y).
\end{equation}
The latter inequality includes the case of one fine cube: the
coarse cubes in its collar share a corner, and their connecting
length is zero. It is the passage in
\eqref{eq:deep-collar-scale}, rather than an estimate on an
individual Gaussian integral, that restores the decay rate at
the next scale. The domain count through one cube is exponential
in $d_k(Y)+1$ and is paid by a fixed fraction of the incoming
rate. The remaining fraction, after rescaling, gives
\begin{equation}\label{eq:deep-gas-bound}
 |H(Z)|\le C\epsilon_k
              e^{-\lambda(d_{k+1}(Z)+5)},\qquad
 \epsilon_k=\delta_k+\mathsf s_kN_\omega
                         +e^{-c\gamma_2\mathfrak b^2},
\end{equation}
with $\lambda\ge cL\kappa$. In including the free quadratic
potential, its fixed small bound is added to $\epsilon_k$.
The constants are chosen so that $\epsilon_k\le1$. The exact
value of $c$ in the tail exponent is immaterial after the floor
has been fixed.

We use the Koteck\'y--Preiss criterion \cite{KoteckyPreiss1986},
with its pinned connected-graph bound, to take the logarithm.
For a finite hard-core polymer set, put
\[
 f(Z,Z')=-1_{\{Z\nsim Z'\}},\qquad
 \rho^T(Z_1,\ldots,Z_n)
   =\sum_{G\ \mathrm{connected}}\prod_{\{i,j\}\in G}f(Z_i,Z_j).
\]
The incompatibility relation is symmetric and reflexive, so
repeated polymer labels in the tuple are treated correctly.
Suppose nonnegative numbers $a(Z)$ obey
\begin{equation}\label{eq:deep-kp-condition}
 \sum_{Z'\nsim Z}|h(Z')|e^{a(Z')}\le a(Z).
\end{equation}
Delete vertex one from a connected incompatibility graph.
Each remaining connected component has to be attached to the
root by a nonempty set of allowed edges. Summing the signs over
all those nonempty sets gives $-1$ if at least one edge is
allowed, and zero otherwise. Taking absolute values after this
cancellation and discarding the remaining restrictions bounds
the rooted connected sum by the exponential of smaller rooted
sums. If $t_N(Z)$ is the resulting pinned majorant truncated at
$N$ other vertices, the recursive upper bound is
\[
 t_{N+1}(Z)
       \le\exp\Big(\sum_{Z'\nsim Z}|h(Z')|t_N(Z')\Big).
\]
Starting with $t_0=1$, \eqref{eq:deep-kp-condition} proves
$t_N(Z)\le e^{a(Z)}$ by induction. Monotone passage in the
nonnegative majorants gives the pinned estimate
\begin{equation}\label{eq:deep-pinned-bound}
 \sum_{n\ge1}\frac1{(n-1)!}
  \sum_{Z_2,\ldots,Z_n}|\rho^T(Z,Z_2,\ldots,Z_n)|
                        \prod_{i=2}^n|h(Z_i)|\le e^{a(Z)}.
\end{equation}
The cancellation concerns hard-core edge factors $0,-1$.
Replacing it by an arbitrary sum of absolute graph weights
would lose this recursion.

The exponential formula for finite labelled graphs identifies
the formal logarithm of $\Xi$ with the connected series.
The preceding bound makes this series absolutely convergent
on the admitted activity polydisc. Along the path $rh$,
$0\le r\le1$, the same majorant is uniform. The identity with
$\Xi$ holds near zero by the finite algebraic formula and then
throughout this path by analytic continuation and continuity.
Thus
\begin{equation}\label{eq:deep-connected-log}
 \Xi(h)=\exp\left[
  \sum_{n\ge1}\frac1{n!}
       \sum_{Z_1,\ldots,Z_n}\rho^T(Z_1,\ldots,Z_n)
                                      \prod_{i=1}^nh(Z_i)\right],
\end{equation}
and in particular $\Xi(h)\ne0$. The pinned estimate contains
more information than this nonvanishing: it controls the
response to a marked factor and will provide the local
activity bound and the source derivatives.

Apply the criterion with
\[
 h(Z)=H(Z)e^{\mu(d_{k+1}(Z)+5)},\qquad a(Z)=\#Z,
\]
where $\mu$ is the desired output rate. A polymer incompatible
with $Z$ contains a cube adjacent to one of the cubes of $Z$.
There are a bounded number of such neighbours. Also
$\#Z'\le16(d_{k+1}(Z')+1)$. Hence the left side of
\eqref{eq:deep-kp-condition}, divided by $\#Z$, is bounded
by a fixed multiple of
\[
 \epsilon_k\sum_{Z'\ni\Delta}
 e^{-(\lambda-\mu)(d_{k+1}(Z')+5)+16(d_{k+1}(Z')+1)}.
\]
The sum is at most one once the gap $\lambda-\mu$ is fixed
large enough. Its part proportional to $d_{k+1}(Z')$ pays
the domain entropy, while the spare constant
$e^{-5(\lambda-\mu)}$ pays the fixed neighbour and volume
factors. Only $\epsilon_k\le1$ is used to verify the criterion.
The smallness of $\epsilon_k$ remains available in the bound
on the output.

Assign each connected cluster in
\eqref{eq:deep-connected-log} to the union of its carriers.
Inequality \eqref{eq:deep-carrier-length} transfers the weight
$e^{\mu(d+5)}$ to its vertices. Mark one vertex whose carrier
meets a prescribed cube, and apply
\eqref{eq:deep-pinned-bound} to all the other vertices. The
marked vertex still has its actual activity $H$, so the bound
is linear in $\epsilon_k$:
\[
 |E_{\rm int}(X)|\le C\epsilon_k e^{-\mu d_{k+1}(X)}.
\]
This is why the result does not merely assert an order-one
convergent logarithm. Choosing the incoming rate and $L$ with
a reserve permits $\mu>(1+c)\kappa$, after every volume and
domain sum has been performed.

The same marked sum gives the differential of this logarithm
with respect to a potential. If a new potential
$\epsilon v(Y;B)$ is added, differentiate the expansion at
$\epsilon=0$ and keep its marked carrier. Summing its local
outputs gives the normalised expectation of $\sum_Yv(Y;B)$.
A difference of two inputs is obtained by integrating this
derivative along their segment. The list enters through the
linear old-potential rule with norm at most
$C\mathsf s_k\|\mathfrak f\|$, so
\[
 \|E_{\rm int}[\mathfrak f]-E_{\rm int}[\mathfrak f']\|
                 \le C\mathsf s_k\|\mathfrak f-\mathfrak f'\|.
\]
All these series have the same nonnegative majorant. Each
integrated term is holomorphic on the common complex domain;
uniform convergence on smaller domains preserves holomorphy.
No restriction depending on the number of preceding levels
is introduced by taking these derivatives.

\subsubsection*{Bond densities and the exact pointed sum rule.}
The logarithm just constructed must be assigned base points
before applying the marginal subtraction. Its unpointed norm
does not prescribe this assignment. A convenient construction
uses three paths from the unit integral to the desired one,
keeping a bond label in every derivative.

First turn on the free quadratic gauge potential from zero.
Differentiation produces its normalised expectation with its
original bond label. Second lower each threshold from infinity
to $\mathfrak b$. With parameters $t_b\in[0,1]$, use
$1_{\{t_b|B(b)|<\mathfrak b\}}$, and integrate the sum of its
partial derivatives along the diagonal $t_b=t$. A derivative
at a single bond can be defined after integration by rescaling
that one real variable. The test is then fixed and the
derivative acts on the Gaussian and the remaining smooth
integrand. This also proves holomorphy in the pair variables.
Its local bound contains the failed-test Gaussian tail.

Third turn on the physical fluctuation and the interaction.
The cubic action and gauge remainders are interpolated as
$t^{-2}V^{A,G}(tB)$, with their continuous values at zero;
the Jacobian terms have their own analytic expansion.
Differentiation retains the bond of the physical fluctuation
argument. Marked cluster expansions localise each resulting
normalised expectation. Denote the three bond-labelled
outputs by $G_\tau(X,b)$, $E^\chi(X,b)$ and
$E^\bullet(X,b)$. Each vanishes unless $b^\sharp\subset X$.
Telescoping along the three paths gives the original logarithm,
with the branch fixed by continuity from the unit integral.

Put the tents at bond midpoints and define
\[
 G_\tau(X,z)=\sum_bh_z(b)G_\tau(X,b),\qquad
 F(X,z)=\sum_bh_z(b)
                       [E^\chi(X,b)+E^\bullet(X,b)].
\]
Their reach is smaller than the collar, so only $z\in X$
occur. Summing over $X$ first recovers the exact bond
densities obtained before localisation, and $\sum_zh_z(b)=1$
then recovers the logarithm. These two exact sum rules preserve
the covariance of the pointed totals. They do not require an
individual polymer term to be invariant under a unit
translation that changes the auxiliary partition.

An isometry acts on the free variables by the orthogonal
change of variables induced on the fibre. On a reversed bond,
the potential is transported by
$-\mathrm{Ad}(V_0(b)^{-1})$. The action, gauge fixing, Haar
density and tests are invariant under this change. The
bond density is therefore covariant before localisation.
The midpoint identity $h_{\mathsf rz}(\mathsf rb)=h_z(b)$
transfers this covariance to its pointed total. When this
argument evaluates the older families, their backgrounds
are nested, and their full pointed sums include all domains,
including wrapping ones.

For a domain that does not wrap, all preceding operations
read only the pair and older terms in that domain and its
prescribed collar. Transport between two such domains on
different tori commutes with the construction. Increasing
the torus therefore defines one infinite-lattice companion.
Its coefficient is extracted from that companion. Agreement
of individual finite-volume integrals alone would not select
this common coefficient.

The Gaussian normaliser is localised separately. Up to local
elimination factors,
$\log Z^{(k)}=-\frac12\log\det\mathsf K$.
Its diagonal trace gives a bond label; a midpoint tent partition
$\sum_z h_z(b)=1$ converts this into the base point $z$.
The tents satisfy
$h_{\mathsf rz}(\mathsf rb)=h_z(b)$ for lattice isometries.
Localising the diagonal logarithm by inclusion--exclusion over
regions leaves only closed walk chains meeting every cube of the
domain. Their length is at least $cM d_{k+1}(X)$, so their decay
dominates both the domain count and the subsets in
inclusion--exclusion. The resulting Gaussian family has size $O(1)$
and decay rate of order $M$; the interaction family is small with
the coupling. Together their norm obeys
\[
 \|E^{(k+1)}\|_{k+1,\kappa}
 \le \tfrac12E_0+
 C\{\delta_k+\mathsf s_kN_\omega+e^{-cp_0(g_k)^2}\}
 \le E_0
\]
after the constants are chosen in the stated order.

One exact localisation formula makes the two uses of a walk
expansion particularly clear. Let $\mathsf K_Y$ be the
fluctuation form constructed in a region $Y$, compressed to
the bonds on which it is invertible. For a bond $b$ well
inside the region write
$f_b(Y)=\operatorname{tr}[(\log\mathsf K_Y)(b,b)]$, with the
interior convention giving zero otherwise. The diagonal block is taken
after the operator logarithm. On the finite
set of localisation cubes define
\begin{equation}\label{eq:mrg-logdet-mobius}
 e(X;b)=\sum_{Y\subset X}(-1)^{|X|-|Y|}f_b(Y).
\end{equation}
The same conventions are used in the regions and their
subsets. Summing this formula over $X\subset\mathcal C$
and interchanging the two finite sums gives
\[
 \sum_{X\subset\mathcal C}e(X;b)
 =\sum_{Y\subset\mathcal C}f_b(Y)
     \sum_{Y\subset X\subset\mathcal C}(-1)^{|X|-|Y|}
 =f_b(\mathcal C).
\]
This is an exact sum rule; it uses no smallness.
For its bound one expands the diagonal logarithm by resolvents
of the local form. A chain missing a cube of $X$ occurs in
two subsets with opposite signs and cancels. A surviving
closed chain meets every cube of $X$, and therefore has
length at least $cMd_{k+1}(X)$. Its walk decay pays for the
subset entropy as well as the domain entropy.

The pointed Gaussian activity is consequently formed as
\[
 N(X,z)=-\frac12\sum_bh_z(b)
           [e(X;b)-e(X;b)|_{(1,0)}]
       +\text{local elimination terms}.
\]
The minus sign is the sign of the Gaussian normaliser.
Each elimination term is placed at its pivot bond. For
bond midpoints $x_b$, the tents in the units of level $k$
can be chosen as
\[
 h_z(b)=\prod_{\mu=1}^4
       \max\{1-L^{-1}|(x_b-z)_\mu|,0\},
       \qquad z\in\mathbb T^{(k+1)}.
\]
They sum to one and commute with the lattice isometries.
The localisation domain retains a margin wider than their
reach, so only $z\in X$ occur. The bond density, summed
over its domains before applying these tents, is independent
of the auxiliary partition. This is how covariance is
returned to the pointed totals even though the individual
terms depend on the partition.

The decay of the Gaussian family uses two different expansions,
and the distinction prevents a loss at large domains. At the
real factor of an admitted pair, each regional compressed form
has spectrum in a fixed interval $[\gamma,B]$. The perturbation
at a complex pair has norm smaller than a fixed fraction of
$\gamma$. Choose a fixed $c_N>0$ such that
\[
 R_Y=1-c_N^{-1}\mathsf K_Y,
                  \qquad \|R_Y\|\le q<1
\]
throughout that complex domain. The operator logarithm is then
\[
 \log\mathsf K_Y=(\log c_N)1
                         -\sum_{m\ge1}\frac1mR_Y^m.
\]
The centre $c_N$ and $q$ depend only on the real spectral
window and the admitted perturbation margin. They are
independent of $Y$, the torus and the depth.

Apply the regional inclusion--exclusion formula
\eqref{eq:mrg-logdet-mobius} to this diagonal series. A term
of $R_Y^m$ is a closed chain of local operator carriers. When
the chain does not reach a cube of $X$, adjoining or removing
that cube leaves the term unchanged. The two subset signs
cancel. This uses the common assignment of regional operator
terms, including their fixed interior margins. A chain that
survives must reach within that margin of every cube of $X$.
Consequently its total length is at least $cMd_{k+1}(X)$.

Put $d=d_{k+1}(X)$. For the powers
$m\le\theta Md$, the walk estimate for a closed chain gives
a bound of the form
\[
 C Q_N^m e^{-c'Md},\qquad Q_N>1.
\]
The walk norm need not be smaller than one. Choose $\theta$
after $Q_N,c'$ so that
$Q_N^{\theta Md}\le e^{c'Md/2}$. Summing these powers leaves
an exponential in $-Md$. For the powers $m>\theta Md$, do
not use the walk norm. The operator estimate $\|R_Y^m\|\le q^m$
gives instead
\[
 \sum_{m>\theta Md}\frac{q^m}{m}
          \le Cq^{\theta Md}.
\]
There are at most $2^{\#X}$ subsets in the localisation, with
$\#X\le C(d+1)$. The number of relevant bonds and tent
base points is polynomial in $M$ and $d+1$. Both factors
are absorbed by a smaller fraction of the rate proportional
to $M$. The elimination determinant is local at its pivot
and obeys the same domain convention. Hence
\begin{equation}\label{eq:deep-normalizer-bound}
 |N(X,\mathbf p,z)|\le E_Ne^{-c_N'Md_{k+1}(X)}.
\end{equation}
The constants here contain no coupling. In particular the
choice $E_0\ge2E_N$ is made before the smallness conditions
on the interaction. Increasing the coupling weight cannot
supply smallness for this determinant family.

M\"obius inversion recovers the whole diagonal logarithm before
the tents are applied. Therefore the pointed sum recovers
$\log Z^{(k)}(W)-\log Z^{(k)}(1)$ exactly, rather than up to
a small error per bond. Its prelocalisation bond density is
independent of the auxiliary partition. The common regional
construction transports this family to a nonwrapping domain
of any larger torus; this also supplies its infinite-lattice
companion. The gap, regional agreement and walk estimates
used above are the uniform operator inputs. The two-range
series estimate explains how those inputs yield the required
determinant decay.

At the real action weight, an equivalent single interpolation
combines the threshold and interaction paths. Insert it by multiplying
the fluctuation field by $t$, from zero to one, including
its occurrence in the tests. Its derivative is a sum over
bonds. Localising each such summand first preserves its
bond label, and the tents then convert it to a site label.
The procedure must retain the pointed sum rule, rather
than distributing an already estimated unpointed activity
arbitrarily among its sites.

More explicitly, for the interacting
integral use
\[
 I(t)=\int d\mu_C(B)\,
          \prod_b\chi(tB(b))e^{\mathsf W_k(tg_k,B)}.
\]
On the nonvanishing component supplied
by the local expansion, its logarithmic
derivative is a sum of two bond-labelled
types:
\begin{equation}\label{eq:xrg-pointing-integral}
 \begin{split}
 \frac{d}{dt}\log I(t)
 =I(t)^{-1}\int d\mu_C(B)\,e^{\mathsf W_k(tg_k,B)}
 \bigg[&g_k\prod_{b'}\chi(tB(b'))
       \sum_b\operatorname{tr}
               \mathsf W'_k(tg_k,B,b)(\mathsf C B)(b)\\
 &+\sum_b\partial_t\chi(tB(b))
                       \prod_{b'\ne b}\chi(tB(b'))\bigg].
 \end{split}
\end{equation}
The prime is the derivative in the
physical fluctuation argument, with
its component at $b$. Products of
tests and the test-derivative sum
range over free bonds; the functional
derivative retains the bond label
of the physical field $\mathsf C B$.
Derivatives of
the characteristic functions in this
formula are understood after
integration over $B$, as derivatives
of the moving integration domains.
They are not pointwise functions of
a sharp test. The Gaussian tail
estimates control this operation.
Integrating in $t$ and localising
each bond summand before forgetting
its label gives the pointed density.
The sums over all local domains
recover these exact bond densities;
the tent conversion therefore
recovers their full lattice
covariance. This is the reason
that an arbitrary division of the
unpointed logarithm over its sites
would not be sufficient.

\subsection{Extracting the marginal action}

A bounded norm of each family is not yet a bound on their sum at a
late background. One unit volume at level $k$ contains
$\eta_{k-j}^{-4}$ sites of level $j$. The interaction at level $j$
must therefore be evaluated more accurately than its original
supremum bound. The two subtractions are its value at the trivial
field and its quadratic curvature contribution.

The coefficient of the latter is defined on a single family
$\widehat E^{(j)}$ on $\mathbb Z^4$ that agrees with all torus
terms on domains not wrapping around the torus. Let
$\mathfrak h_\infty$ be the linear response at the trivial field
in both potential and current: a unit coarse-bond perturbation
$\delta_c$ is mapped to the pair consisting of the resulting fine
potential and fine action current. For bonds $c,c'$, define the
following absolutely convergent kernel, with each activity Hessian
evaluated at the trivial pair $(\mathbf U,\mathbf J)=(1,0)$:
\[
 \Pi^{(j)}(c,c')=
 \sum_z\sum_{\widehat X\ni z}
 \operatorname{Hess}\widehat E^{(j)}(\widehat X,\cdot,z)
 [\mathfrak h_\infty\delta_c,\mathfrak h_\infty\delta_{c'}].
\]
Response decay and the tree bound give convergence. Translation and
hypercubic covariance apply to the pointed totals, and linearised
gauge invariance makes this kernel transverse. With $c=(x,\mu)$,
$c'=(y,\nu)$ and $u=y-x$, translation covariance identifies it with
the directional kernel $\Pi^{(j)}_{\mu\nu}(u)$. With
$\partial_\mu(p)=e^{ip_\mu}-1$, its small-momentum expansion is
\begin{equation}\label{eq:rg-polarisation}
 \widehat\Pi^{(j)}_{\mu\nu}(p)
 =b_j\left(\delta_{\mu\nu}\sum_\rho|\partial_\rho(p)|^2
       -\overline{\partial_\mu(p)}\partial_\nu(p)\right)
 +O(|p|^3),
 \quad
 b_j=\sum_u\Pi^{(j)}_{\mu\nu}(u)u_\mu u_\nu
 \quad(\mu\ne\nu).
\end{equation}
Simplicity of the Lie algebra makes the colour factor a scalar.
The expression in parentheses is the Hessian of the plaquette
action. The coefficient is independent of the chosen distinct
directions $\mu,\nu$.

\subsubsection*{The coefficient and the symmetry of its pointed tensor.}
The argument uses Ba\l aban's differentiated gauge identities
\cite[equations~(4.14)--(4.15), (4.33)]{Balaban1987}
and the comparison of pointed and total moments
\cite[equations~(3.58)--(3.64)]{Balaban1988}.
Here the pointing uses bond midpoints, so that it commutes
with the lattice isometries.
The Hessian defining $b_j$ is taken in independent potential
and current coordinates at the trivial pair. Let
$v=(A,J)$ and put
$\|v\|_X=\max\{\sup|A|,\sup|\nabla A|,\sup|J|\}$.
On a ball of radius $a'_j$ mapped into the activity core,
Cauchy's estimate gives
\[
 |D^2F(0)[v,v']|
        \le C(a'_j)^{-2}\sup|F|\,\|v\|_X\|v'\|_X.
\]
For each direction, the first derivative is a real-adjoint
invariant linear functional on $\mathfrak{su}(2)$ and
therefore vanishes. Differentiating real gauge invariance
once more gives
\[
 DF(0)=0,\qquad
 D^2F(0)[(\nabla\lambda,0),v]=0.
\]
The additional term from differentiating the coordinate
map vanishes because $DF(0)=0$. Thus the Hessian is
intrinsic at this point and is unchanged by adding a
linear gauge direction to either input.

The linear response $\mathfrak h_\infty$ is determined
modulo precisely these directions. A pure gauge is the
background over its own pure-gauge averages, and
uniqueness identifies the response to a gradient as
a gauge direction. It also makes colour rotations and
real conjugation commute with the response modulo
those directions. Hence
\[
 D^2F(0)[\mathfrak h_\infty(\delta_ct),
                     \mathfrak h_\infty(\delta_{c'}t')]
          =\langle t,t'\rangle\Pi_F(c,c')
\]
for one scalar kernel $\Pi_F$. Here $t,t'$ belong to
the colour algebra and $\delta_c$ is the scalar unit
bond field. Equivalently $\Pi_F$ is one third of the
sum with $t=t'$ over an orthonormal colour basis.
This specifies the scalar colour convention in the
definition of $\Pi^{(j)}$ above.

For a term on $X$ the response bound and Cauchy give
\[
 |\Pi_F(c,c')|
  \le C(a'_j)^{-2}\|E^{(j)}\|_{j,\kappa}
       e^{-\kappa d_j(X)}
       e^{-\delta_0[\operatorname{dist}(c,X)
                              +\operatorname{dist}(c',X)]}.
\]
The distance between the bonds is bounded by the sum
of these two distances and $M(d_j(X)+C)$. Half of
the exponents consequently gives decay in $|c-c'|$
at rate $\frac12\min\{\delta_0,\kappa/M\}$.
The other half pays the domain entropy and the sum
of base points in $X$. Therefore the kernel and all
its polynomial coordinate moments are absolutely
summable, uniformly in the torus size in their
periodised versions. The bound for its coefficient
has the two-radius-derivative cost
$C(a'_j)^{-2}\|E^{(j)}\|_{j,\kappa}$.

Full lattice symmetry is passed to this kernel
through the totals, rather than assumed for a term
with its auxiliary cube partition. On a finite
torus differentiate the covariance identity for
the pointed total twice at $V=1$. This point lies
in an invariant ball of block data, so all the
differentiations are legitimate on both sides.
Fix a box of radius $R$ about the base point.
Terms contained in that box agree with their
companions for every sufficiently large torus;
their response kernels converge by periodisation.
The total of all other terms has an exponentially
small bound in $R/M$, uniform in the torus.
First increase the torus at fixed $R$, then
increase $R$. Dominated convergence gives the
covariance identity for the companion Hessian
and its moments.

Translation covariance of the total makes its
kernel a function of the bond directions and
$c'_--c_-$. Signed permutations act on its
quadratic moment tensor. At a fixed base point
they act on the pointed tensor as well, because
the full pointed total was used in the preceding
limit. Differentiated gauge invariance applied
to linear and quadratic gauge parameters shows
that a linear potential contributes only
through its curl. These tests may have polynomial
growth: their boundary errors tend to zero
against the exponential response kernel.

For completeness, a pointed tensor can be written
as
\[
 T_{\mu\nu,\rho\tau}(z)
    =\sum_{c,c'}\Pi^{(j)}(c,c';z)
                 (c_--z)_\rho(c'_--z)_\tau,
\]
where the sum restricts $c$ to direction $\mu$
and $c'$ to direction $\nu$, and $\Pi^{(j)}(c,c';z)$
is the sum over all companion terms pointed at $z$. Its antisymmetries
are the curl identities just obtained. Reflection
of one coordinate eliminates a pairing of two
different coordinate plaquettes, and permutation
of coordinates identifies their diagonal
coefficients. The invariant colour pairing is
already fixed. This gives the tensor in
\eqref{eq:xrg-pointed-tensor}.

To identify its number with
\eqref{eq:rg-polarisation}, apply the Ward identity
also to cubic gauge parameters. Define
\[
 P_{\mu\nu,\rho\tau}
       =\sum_u\Pi^{(j)}_{\mu\nu}(u)u_\rho u_\tau.
\]
The moment form of the cubic identity is
\[
 P_{\rho\nu,\tau\sigma}
   +P_{\tau\nu,\rho\sigma}
    +P_{\sigma\nu,\rho\tau}=0.
\]
The relation between the two moment conventions is
$T_{\mu\nu,\rho\tau}+T_{\mu\nu,\tau\rho}
=-P_{\mu\nu,\rho\tau}$. It relates the pointed
linear-potential tensor to this second moment;
the scalar left after the signed-permutation
identities is exactly
$\sum_u\Pi^{(j)}_{\mu\nu}(u)u_\mu u_\nu$,
independent of $\mu\ne\nu$. In passing from
bond anchors to midpoint coordinates, the
extra terms multiply lower moments. The
zeroth moment vanishes by the linear Ward
identity, and inversion removes the first
moment. No additional action coefficient
comes from this change of anchors.

This argument also explains why numerical
normalisations do not affect the coefficient.
Two local logarithms of the same nonzero
integral differ by a locally constant value.
Their Hessians and coefficients agree. In
particular the total of $G_\tau$ in
\eqref{eq:rg-weight-count} is field independent,
so its coefficient is zero even when its
individual local terms vary with the field.
The coefficient is extracted after the exact
pointed sum rule, before discarding any
wrapping tail in a finite volume.

The gauge-first subtraction makes the power
counting explicit. In an axial gauge centred at the base point,
write $F_{\rho\mu}(z)$ for the background curvature component at
$z$, measured in the units of level $j$. Then
\[
 B^{\mathrm{ax}}_\mu(x)=
 \sum_{\rho<\mu}(x_\rho-z_\rho)F_{\rho\mu}(z)
 +O\!\left(|x-z|^2\eta_{k-j}^{2+\alpha_{\mathrm H}}\right).
\]
Thus this potential has size $\eta_{k-j}^2$ on the local domain.
Terms of third and higher order already have a summable gain.
The linear term vanishes by gauge invariance and the absence of
an invariant vector in the Lie algebra. The remaining quadratic
curvature form has the same coefficient as $b_jA$.
The covariance of the pointed totals is needed here: it gives this
cancellation at each base point, rather than only after summing
over a window. Individual localised terms need not be covariant
under all unit translations. The tent $h_z$ is evaluated at
plaquette centres in $A(h_z,U)$, using the same partition of unity.

For a regular background of level $k$, curvature size
$\varepsilon$, and curvature radius $\alpha_{0,j}$, the resulting
one-lattice estimate is
\begin{align}\label{eq:rg-subtraction}
 &\left|\sum_{X\ni z}
 [E^{(j)}(X,U_k,z)-E^{(j)}(X,1,z)]
                  -b_jA(h_z,U_k)\right|\notag\\
 &\hspace{8mm}\le C\|E^{(j)}\|_{j,\kappa}
 \left\{(\varepsilon/\alpha_{0,j})^2
             \eta_{k-j}^{4+\alpha_{\mathrm H}}
                  +e^{-c\kappa/\eta_{k-j}}\right\},
 \qquad \varepsilon\eta_{k-j}\le c\alpha_{0,j}.
\end{align}
Long domains supply the exponential tail. On non-wrapping local
domains the field factor is retained; domains that wrap may lose
it, and are paid for by their size. On nested regions the same
estimate has boundary tails and the counterterm carries exactly
the sum of the tents of the retained base points. These terms
are included in \cref{sec:large-fields}.

The pointed cancellation can be stated at the
quadratic-tensor level. Let $T_zF$ be the
linear axial potential at $z$ with constant
curvature tensor $F$:
\[
 (T_zF)_\mu(x)=\sum_{\rho<\mu}(x_\rho-z_\rho)F_{\rho\mu}.
\]
For each pointed companion term take its
quadratic form on the free-boundary response
to this potential. Summing the terms through
$z$ gives a quadratic form $Q_z(F)$ on
antisymmetric two-forms with colour indices.
Absolute convergence follows from response
decay and the exponential domain norm, since
the coordinate moments grow only polynomially.

At $z$ the covariance of the pointed total
acts on this form. A sign change of one
coordinate reverses the curvature components
containing that coordinate. Consequently a
mixed term between two different coordinate
plaquettes changes sign under some such
reflection and vanishes. Permuting axes
makes all the remaining diagonal coefficients
equal. The simple colour algebra supplies
one invariant colour form. With the
normalisation in \eqref{eq:rg-polarisation},
the resulting tensor is the Wilson
quadratic tensor multiplied by $b_j$.
In all ordered antisymmetric indices its
coordinate part is proportional to
\begin{equation}\label{eq:xrg-pointed-tensor}
 \delta_{\mu\rho}\delta_{\nu\tau}
                  -\delta_{\mu\tau}\delta_{\nu\rho}.
\end{equation}
Equivalently, on the components $\mu<\nu$
it is the identity tensor. Thus the leading
$F(z)$ in the pointed interaction and in
$b_jA(h_z,U)$ is exactly the same quadratic
form. The passage from a local box response
to the free-boundary response is made inside
this gauge-invariant quadratic form. Its
exponentially small error is retained
separately.

Taylor's formula in the axial block datum
then has three relevant contributions.
The zeroth order is removed. The first
order vanishes by constant gauge invariance.
The displayed quadratic tensor is removed
by the action counterterm. Because the
axial potential is built from curvature,
the third-order contribution has an extra
two powers of $\eta_{k-j}$ before the
analytic-radius factors are taken. The
remaining loss in the quadratic part is
the H\"older remainder of the potential
beyond $T_zF$, which gives precisely the
factor $\eta_{k-j}^{\alpha_{\mathrm H}}$.
This is the cancellation needed by the
rooted interpolation. It cannot be applied
to a single unpointed domain term.

\subsubsection*{The pointed cancellation at complex window fields.}
The subtraction is also needed inside the local fluctuation
expansion, where the arguments are independent complex pairs.
The real-background estimate alone does not cover these
arguments. We describe the additional window calculation,
retaining its domain requirement.

Fix an older level $j$ and a window at the current scale. Put
$\varsigma=L^{-(k-j)}$ and write $\bar a$ for a fixed
multiple, depending on $M$, of the output curvature radius.
The box response estimates give a block datum $e^{iB}$ in
level-$j$ units, with
\[
 |B|\le C\bar a\varsigma,\qquad
 |\partial B|\le C\bar a\varsigma^2,\qquad
 [\partial B]_{\alpha_{\mathrm H}}
                         \le C\bar a\varsigma^{2+\alpha_{\mathrm H}}.
\]
The last estimate is a H\"older estimate for first differences;
it supplies no bounded second derivative. The substituted
potential and reduced current must together belong, with a
fixed margin, to the core of every older term that is evaluated.
That membership uses the nonlinear background and domain
construction. The calculation below applies on that common
admitted family of pairs.

Choose an axial gauge about a base point $z$. The nonabelian
Stokes expansion on the comb gives
\begin{equation}\label{eq:deep-complex-axial-datum}
 \widehat B_\mu(x)
  =\sum_{\rho<\mu}(x_\rho-z_\rho)F_{\rho\mu}(z)
           +R_\mu(x),\qquad
 |R(x)|\le C\bar a(1+|x-z|)^2
                         \varsigma^{2+\alpha_{\mathrm H}}.
\end{equation}
Here $F(z)$ is the curvature of this block datum, with
$|F(z)|\le C\bar a\varsigma^2$. It is the same curvature
in the activity and in its pointed action counterterm. A
different strength defined from an unrelated bare gauge would
not give an exact marginal cancellation.

For small real data, uniqueness of the regular critical orbit
identifies the window field with the box Landau response to
$\widehat B$, up to a real gauge. Gauge invariance identifies
their activity values. Both expressions are holomorphic along
the specified small complex datum family. The identity theorem
therefore gives the same identity on its connected complex
continuation. It does not assert invariance under every complex
gauge transformation on an unrestricted orbit.

Let $\mathcal H$ denote that box response, including its
current entry, and let $a'_j$ be the radius used for Taylor
estimates near the trivial pair. The weighted response bound
controls its size on the domain of a term by
\begin{equation}\label{eq:deep-window-size}
 \|\mathcal H(\widehat B)\|_X\le N_X,
 \qquad N_X=CM(d_j(X)+2)\bar a\varsigma^2,
 \qquad
 \frac{N_X}{a'_j}
       \le C(M)(d_j(X)+2)\varrho_j\varsigma^2.
\end{equation}
The factor $\varrho_j$ is the ratio of output to older radius,
not a number of differentiating variables. The decay of the
response permits the linear growth in
\eqref{eq:deep-complex-axial-datum}; it is not estimated by
the supremum of that growing datum on an arbitrarily large
torus.

There are two elementary cases before Taylor expansion.
If $N_X\ge a'_j/2$, the original norm bound is bounded by
a fixed multiple of its product with $(N_X/a'_j)^3$.
This is bounded by a polynomial in $d_j(X)+2$ times
$\varrho_j^3\varsigma^6$. If the age is so small that
$C\varrho_j\varsigma>1$, use the independent norm bound.
There are only boundedly many such ages, since
$\varrho_j\le C(k+2-j)^{1/2}$ and $\varsigma$ decreases
geometrically. Enlarging the fixed constant covers them.
In the remaining case $N_X<a'_j/2$ and
$C\varrho_j\varsigma\le1$, Taylor expansion is inside the
analytic ball and its cubic remainder obeys the desired
power count.

Absorb the pointed action counterterm into the family and
write its terms as $\widetilde E^{(j)}(X,\cdot,z)$.
The absorbed family has the same decay norm up to a fixed
constant: its coefficient is bounded by
$Ca_{0,j}^{-2}\|E^{(j)}\|_{j,\kappa}$ and each pointed
plaquette action is $O(a_{0,j}^2)$ on the core. Constant
gauge invariance removes the first derivative. The Hessian
annihilates gradient directions, as follows by differentiating
real gauge invariance. Therefore replacing a box response by
the free-boundary response is allowed inside this quadratic
form. Its discrepancy is exponentially small in the distance
to the window boundary.

For the numerical tensor $\widetilde T^{X,z}$ formed from
that Hessian, the resulting expansion is
\begin{equation}\label{eq:deep-window-taylor}
 \begin{split}
 \widetilde E^{(j)}(X,\cdot,z)^\square
  ={}&\widetilde E^{(j)}(X,(1,0),z)\\
     &+\tfrac18\sum_{\mu,\nu,\rho,\tau}
          \widetilde T^{X,z}_{\mu\nu,\rho\tau}
                   \langle F_{\rho\mu}(z),F_{\tau\nu}(z)\rangle
       +\mathcal R_X,\\
 |\mathcal R_X|
  \le{}&C(M)\|E^{(j)}\|_{j,\kappa}e^{-\kappa d_j(X)}
             (d_j(X)+2)^C\\
       &\hspace{5mm}\cdot
           [\varrho_j^2\varsigma^{4+\alpha_{\mathrm H}}
                              +\varrho_j^3\varsigma^6].
 \end{split}
\end{equation}
The first error comes from pairing the linear axial potential
with the H\"older remainder. The second is the cubic Taylor
remainder. Commutator and reduced-current errors have an
additional full power and obey the same bound. Polynomial
coordinate moments are paid by a fraction of $\kappa$.
This accounts for the two distinct radius powers in the
estimate.

The cancellation is a sum over all companion domains through
$z$:
\[
             \sum_{\widehat X\ni z}\widetilde T^{X,z}=0.
\]
The tensors are taken at the trivial pair; they depend on
neither the current field nor the system of radii used to
bound it. In a window only the domains lying inside its fixed
enlargement are used. Their tensor sum is the negative of the
missing tail. Every omitted domain reaches at least
$2/\varsigma-2$ units in the domain-length convention, so its
tensor norm has an exponential tail. The quadratic contribution
of that tail is bounded by
\[
 C\|E^{(j)}\|_{j,\kappa}\varrho_j^2\varsigma^4
                                           e^{-c\kappa/\varsigma}.
\]
It is absorbed in the H\"older error. This is an equality of
the same numerical tensors before estimating a partial sum;
there is no claim that each individual local term has zero
marginal tensor.

After summing domains through $z$, and then the
$O(\varsigma^{-4})$ base points in a current-scale window,
\eqref{eq:deep-window-taylor} is bounded by
\[
 C\|E^{(j)}\|_{j,\kappa}
          [\varrho_j^2\varsigma^{\alpha_{\mathrm H}}
                                     +\varrho_j^3\varsigma^2]
 \le C\|E^{(j)}\|_{j,\kappa}
                    \varrho_j^3\varsigma^{\alpha_{\mathrm H}}.
\]
Thus the full age sum has precisely the summable weight
$(k+2-j)^{3/2}L^{-\alpha_{\mathrm H}(k-j)}$.
The gain is $\alpha_{\mathrm H}$; the supplied first-difference
regularity does not justify replacing it by a full derivative.

The cancellation is performed before decoupling the old
interaction. Each bracket then acts on a window function
whose marginal tensor has already been removed. Its response
multiplier $\nu$ has analytic radius
\[
 \rho_F\ge c\varsigma^{-1}a_{0,j}
                  \delta(\mathsf t)^{-1}e^{-C\bar\kappa}
 \ge c\varsigma^{-1}
             (\varrho_j\mathsf s_{\mathfrak a})^{-1}e^{-C\bar\kappa}.
\]
The exponential scale factor dominates the square-root radius
ratio. Choosing $\mathsf s_{\mathfrak a}$ small places
$\nu=1$ in this disc at every age. Schwarz in $\nu$, after
subtracting its value at zero, supplies the additional factor
$C e^{C\bar\kappa}\mathsf s_{\mathfrak a}$ in the old
potential estimate. Cauchy in the cube parameters supplies
the carrier decay on the same common polydisc.

An aged remainder has no pointed action counterterm and
therefore has only the first, quadratic Taylor estimate.
After the base-point count its load is $\varrho_j^2\|R^{(j)}\|$,
with no geometric gain in age. The lower coupling drift
sums that load. The regular and remainder estimates cannot
be combined into one artificially decreasing age weight.
For general fields on a torus, the wrapping-domain tail also
retains the norm-only treatment of
\eqref{eq:mrg-three-domain-bound}. The window calculation
does not remove flat holonomy from such domains.

For layered regions, split the retained
base points as $\Xi=\bigcup_{n\ge j}\Xi_n$
by the local layer. The estimate takes
the quantitative form
\begin{equation}\label{eq:xrg-layered-subtraction}
 \begin{split}
 &\left|\sum_{z\in\Xi}\sum_{X\ni z}
       [E^{(j)}(X,U_k,z)-E^{(j)}(X,1,z)]
       -b_jA\Big(\sum_{z\in\Xi}h_z,U_k\Big)\right|\\
 &\quad\le CE_0\sum_{n\ge j}
        (1+n-j)L^{-c(n-j)}L^{-4(n-j)}\#\Xi_n,
                     \qquad c\le\alpha_{\mathrm H}.
 \end{split}
\end{equation}
The retained points must lie a cell inside
the allowed summation region; this controls
the missing-domain tail at the boundary.
Here $L^{-4(n-j)}\#\Xi_n$ is current
volume, $1+n-j$ is the squared radius ratio,
and $L^{-c(n-j)}$ is the summable surplus.
If the tent sum is replaced by a smoother
profile, the difference is supported on
the layer's rim. That rim costs a constant
per site at its own scale, without the
surplus age factor.

The exponential tail can be kept sharper than the version in
\eqref{eq:rg-subtraction}. Put $\eta=\eta_{k-j}$.
There are three different classes of domains through $z$.
The local class stays in a nonwrapping window of size of order
$\eta^{-1}$ in level-$j$ units. On it the axial potential is
controlled by curvature, and Taylor's formula gives the first
term in \eqref{eq:rg-subtraction}. A second class leaves this
window but still lies in a nonwrapping box. Such a domain has
$d_j(X)\ge c/\eta$ and remains gauge equivalent to the
identity when the curvature vanishes. Its estimate retains the
quadratic field factor. Only the third class, consisting of
domains in no such box, can detect flat holonomy. For this class
$d_j(X)\ge c\ell L^{K-j}/M$, and one uses the norm alone.
Thus the required sharper estimate is
\begin{equation}\label{eq:mrg-three-domain-bound}
 C\|E^{(j)}\|_{j,\kappa}
   \left[(\varepsilon/\alpha_{0,j})^2
                    \eta^{4+\alpha_{\mathrm H}}
          +e^{-c\kappa\ell L^{K-j}/M}\right].
\end{equation}
There is no factor of $\varepsilon$ in its last summand.

For the second class the field factor is obtained without a
Taylor expansion over the whole torus. On a box containing
$X$, axial gauge bounds the potential by
$C(d_j(X)+1)\varepsilon\eta^2$.
If its ratio to the analytic radius is small, a quadratic
Cauchy estimate applies; if that ratio is large, the norm
bound is at most a constant times the square of the ratio.
In either case the polynomial in $d_j(X)+1$ is absorbed by
part of the decay. The remaining factor
$e^{-c\kappa/\eta}$ is smaller than any prescribed power
of $\eta$ and is absorbed in the first term of
\eqref{eq:mrg-three-domain-bound}. This argument uses a
nonwrapping gauge box. It does not restore the field factor
on the third class.

For a restricted set $\Xi$ of retained base points, the same
pointed subtraction gives the counterterm
$b_jA(\sum_{z\in\Xi}h_z,U)$. Replacing this tent sum by an
arbitrary smooth profile changes the action near the rim of
$\Xi$. That difference is a boundary activity and is paid
there. It has no decreasing age weight merely from changing
the name of the profile.

The factor $\eta_{k-j}^4$ is cancelled by the number of base points
per current volume. It is precisely marginal in four dimensions.
The additional H\"older factor is summable after the quadratic
action has been subtracted. Even with running radii the estimate
$(\varepsilon_k/\alpha_{0,j})^2\le C(1+k-j)$ leaves a summable
series. Hence the old interaction is kept in the form
\begin{equation}\label{eq:rg-old-interaction}
 \mathbf E_k(U)=\sum_{j\le k}
 \left\{\sum_z\sum_{X\ni z}
 [E^{(j)}(X,U,z)-E^{(j)}(X,1,z)]-b_jA(U)\right\}.
\end{equation}
The norms of previously formed terms are unchanged. Their
subtracted values at later backgrounds decrease; this is the
estimate used when the local expansion weights their ages.

With these definitions the small-field integral returns an
identity of the same form at the next level:
\[
 \int dV_k\,
 \delta(\mathcal M(V_k)V_{k+1}^{-1})\mathfrak G_k\chi^{(k)}
 e^{-x_kA(U_k)+\mathbf E_k(U_k)}
 =e^{\mathsf E^{(k)}}
   e^{-x_{k+1}A(U_{k+1})+\mathbf E_{k+1}(U_{k+1})}.
\]
The new coupling is determined from the output family, while
$\mathsf E^{(k)}$ collects the field-independent constants.
At the trivial field their logarithmic part includes
\[
 \log\mathfrak z(x_k,\varepsilon_k)
   =-\tfrac32\log x_k+O(1),
 \quad
 \mathfrak n_k=
 -\left(\tfrac92\log x_k+O(1)\right)
       (|\mathbb T^{(k)}|-|\mathbb T^{(k+1)}|)
\]
for $\SU(2)$. The remaining local Gaussian numbers and vacuum
values are counted once in $\mathsf E^{(k)}$. This exact
normalisation is used by the source bounds: an arbitrary error
per bare bond would produce a cutoff-dependent exponential
factor in the partition function.

\subsection{The one-loop sum and the closed flow}

The coefficient of the output separates into the Gaussian
normaliser, the remaining small-field interaction, and terms
descending from large-field regions:
\begin{equation}\label{eq:increment}
 x_0=\beta,\qquad x_{k+1}=x_k-b_{k+1},\qquad
 b_{k+1}=b^0_{k+1}+r_{k+1}+q_{k+1}.
\end{equation}
The required coefficient estimates are
\begin{equation}\label{eq:flow-estimates}
 \left|\sum_{l=1}^k b_l^0-c_0k\right|\le C_0,
 \qquad |r_{k+1}|\le Cg_k(\log g_k^{-2})^{\mathfrak p_0},
 \qquad |q_{k+1}|\le Cg_k^{\kappa_0-4},
\end{equation}
where, in the action normalisation fixed above,
$c_0=\frac{11}{3\pi^2}\log L$ for $\SU(2)$.

Fix positive induction constants
$B_\sharp,\lambda_\sharp,C_\sharp$. A prefix through level $k$ is
\emph{admissible} if $g_j\le\gamma$ for $j\le k$, the two-sided
bound for $x_i-x_j$ in \eqref{eq:flow} holds for
$0\le i\le j\le k$, and the torus condition
\eqref{eq:prefix-volume} holds for $j\le k$. These are requirements
on the already constructed levels. The coefficient estimates below
must preserve them at the next level, with one fixed choice of these
constants; no next-level bound is part of the premise.

\subsubsection*{Uniform estimates in the renormalisation step}
\phantomsection\label{hyp:uniform-rg}

The induction must establish
\cref{eq:rg-background-regularity,eq:rg-green,eq:rg-walks,eq:rg-subtraction}
on the specified independent field/current domains for every admissible
prefix, on the allowed tori and regular layered regions. The constants
must be independent of the cutoff, level, torus side and localisation
position. Local agreement on nonwrapping domains then gives one
coefficient sequence on all tori. The quantitative flow calculation is
\cref{eq:increment,eq:flow-estimates}; the contribution of healed regions
is estimated in \cref{sec:large-fields}. Only the already constructed
prefix enters the next increment, which is formed before the next
large-field operation.

Two analytic domains enter the estimate of the errors. For a
family of size $\epsilon$ on radius $\alpha$, taking the two
derivatives that define its coefficient costs
$C\epsilon\alpha^{-2}$. Response decay and the domain count
also give exponential decay of its Hessian kernel. Applying this
estimate only on the running radius, which is of order $g_k$
times a logarithm, would destroy the smallness of the ordinary
small-field interaction. Instead, the small-field chain is
constructed on fixed absolute radii and carries its own action
counterterm. It gives $r_{k+1}$ in
\eqref{eq:flow-estimates}. The terms left by the large-field
operation are treated on running radii. Their summed size is
$Cg_k^{\kappa_0-2}$ under the lower flow bound, and two
derivatives cost $g_k^{-2}$ up to logarithms; this gives
$q_{k+1}$. For an input list bounded on both systems of radii,
the coefficients from the two determinations agree on their
common domain, since their logarithms differ only by a
field-independent constant. Thus this split describes one
increment, rather than two different coupling prescriptions.

The two chains must also accept lists which were not produced by an
earlier small-field step. This occurs when a healed region is replaced
by local activities, and again when the action is varied while the real
tests are held fixed. The one-step construction is therefore formulated
on its analytic input families, rather than on their history of
formation. Each family must still have the covariance and the common
infinite-lattice companion used in the marginal subtraction.

To describe that input, write $b[S^{(j)}]$ for the coefficient extracted
from the companion of a pointed family $S^{(j)}$. Its contribution to
the interaction is
\begin{equation}\label{eq:rg-list-interaction}
 \sum_{j\le k}\left\{
   \sum_z\sum_{X\ni z}
     [S^{(j)}(X,\cdot,z)-S^{(j)}(X,1,z)]
                      -b[S^{(j)}]A\right\}
 +\sum_{j\in\mathcal J}\sum_X
                  [R^{(j)}(X,\cdot)-R^{(j)}(X,1)].
\end{equation}
The second sum contains the aged remainders, which have no action
counterterm. The first sum has both subtractions for each member
separately. The requirement is stronger than asking that the sum have
one quadratic coefficient: the rooted estimates use the cancellation
in each older family at its own scale. The load is
\eqref{eq:xrg-old-load}, with $S^{(j)}$ in place of $E^{(j)}$.
Thus the pointed families carry the summable weights
$\omega_{k+1-j}$, whereas the families in $\mathcal J$ carry only
the squared radius ratios $\varrho_j^2$. The lower coupling drift,
rather than geometric decay in age, is what makes that latter sum
small for the actual remainders. Younger remainders are kept outside
this input until their cubes fit in single blocks, as described in
\cref{sec:large-fields}.

Freeze the tests, field scaling and gauge fixing at a real number
$\mathsf t$, while the Wilson action has a complex weight $x$ with
$|x-\mathsf t|<\sigma_\star\mathsf t$, and put $\tau=x/\mathsf t$.
The relative radius $\sigma_\star$ is fixed after the geometric
constants so that the Gaussian coercivity margin is retained.
The fluctuation threshold is
$\delta(\mathsf t)\asymp\mathsf t^{-1/2}
(\log\mathsf t)^{\mathfrak p_0}$. Choose radii
$(a_{0,j},a_{1,j})$ so that
\[
 CMa_{0,j}<a_{1,j},\qquad a_{1,j}^2\le Ca_{0,j},\qquad
 a_{i,j}\le(1+c)a_{i,k+1}
       \le C(k+2-j)^{1/2}a_{i,j}.
\]
The norms of the list are taken on these domains. In the running
system their comparability is supplied by the flow bound. On absolute
radii it is part of the fixed analytic construction. In either case
put $\mathsf s=\delta(\mathsf t)/a_{0,k+1}$, allowing the fixed
comparison factors between consecutive radii. If
$\mathsf s(1+N_\omega)\le c$, the required one-step output besides
the Gaussian family obeys
\begin{equation}\label{eq:mrg-list-step}
 |F(X)|\le C\big[|1-\tau|+\delta(\mathsf t)
      +\mathsf sN_\omega+e^{-c\mathsf t\delta(\mathsf t)^2}\big]
                    e^{-(1+c)\kappa d_{k+1}(X)}.
\end{equation}
The output has radii $(1+c)a_{i,k+1}$ and decay rate strictly larger
than $\kappa$. The large-field operation subsequently spends this
reserve. At fixed absolute radii $\mathsf s$ is of order
$\delta(\mathsf t)$; at running radii it is only a negative
logarithmic power. The same analytic map must supply these enlarged
domains with independent potential and current variables. A bound on
the physical slice alone would not allow the later interpolations.

\emph{The complex action weight and the Gaussian count.}
The term $|1-\tau|$ in \eqref{eq:mrg-list-step} has a special
origin. Let $\Delta_A$ be the Hessian of the composite action on the
actual constrained fibre, in fluctuation coordinates. It includes the
current-dependent corrections from the constraint and the critical
field. It is not the ambient Wilson Hessian at nonzero current.
First work at a real regular constrained critical point. Let $K$ be the tangent space
of the residual gauge orbit, of dimension $n_{\mathcal O}$.
Gauge invariance makes the intrinsic action Hessian $\Delta_A$
vanish on $K$. The quadratic gauge-fixing form is a sum of squares
of the $n_{\mathcal O}$ mean-contour coordinates. Its restriction
to $K$ is positive definite by the gauge-slice property. Thus
$\Delta_{\mathrm{GF}}$ has rank $n_{\mathcal O}$: its rank is at
most the number of these coordinates and at least the dimension
of $K$. In particular,
\[
 \text{fluctuation space}=\ker\Delta_{\mathrm{GF}}\oplus K.
\]
In coordinates adapted to this decomposition, $\Delta_A$ has
only its $\ker\Delta_{\mathrm{GF}}$ block and
$\Delta_{\mathrm{GF}}$ only its $K$ block.
The change of coordinates is independent of $\tau$, so its
determinant factors cancel in the ratio. Consequently
\[
 \det(\tau\Delta_A+\Delta_{\mathrm{GF}})
   =\tau^{n_{\mathrm{phys}}}
                    \det(\Delta_A+\Delta_{\mathrm{GF}}),
 \qquad n_{\mathrm{phys}}=\dim\ker\Delta_{\mathrm{GF}}.
\]
These dimensions are fixed by the fluctuation and gauge-fixing
coordinates and are independent of the regular background.
The gauge-slice construction must retain the invertibility of these
coordinates on its chosen complex domain. There the same complementary
blocks are holomorphic, and the displayed identity continues
holomorphically; positivity was used only at the real data.

For the fluctuation expansion one instead uses a Gaussian with
covariance $\tau^{-1}C^{(k)}$. The remainder of the quadratic form
is the free potential
\[
 -\tfrac12(1-\tau)\langle B,\Delta_{\mathrm{GF}}B\rangle.
\]
It is already a sum of block terms and requires no further spatial
localisation. Its size restricts the relative complex disc for $x$;
that disc is fixed after the geometric constants so that the
Gaussian coercivity margin remains available. The construction
separates the local family produced by this potential, denoted here
by $G_\tau$, from the interaction family $F_{\mathrm{int}}$. Their
respective bounds are
\[
 \begin{split}
 |G_\tau(X)|&\le C|1-\tau|e^{-(1+c)\kappa d_{k+1}(X)},\\
 |F_{\mathrm{int}}(X)|&\le
 C\big[\delta(\mathsf t)+\mathsf sN_\omega
             +e^{-c\mathsf t\delta(\mathsf t)^2}\big]
                         e^{-(1+c)\kappa d_{k+1}(X)}.
 \end{split}
\]
Their sum is the family estimated in \eqref{eq:mrg-list-step}.

Let $n_{\mathrm f}$ be the number of free fluctuation variables and
$n_{\mathcal O}$ the number of residual gauge directions, so that
$n_{\mathrm{phys}}=n_{\mathrm f}-n_{\mathcal O}$. The determinant
identity gives the two contributions
\begin{equation}\label{eq:rg-weight-count}
 \sum_{X,z}G_\tau(X,U,z)
       =\tfrac12n_{\mathcal O}\log\tau,\qquad
 -\tfrac12n_{\mathrm f}\log\tau
       +\tfrac12n_{\mathcal O}\log\tau
       =-\tfrac12n_{\mathrm{phys}}\log\tau.
\end{equation}
The logarithm is the branch near $\tau=1$. The first equality concerns
the total family. Individual local terms can depend on the field.
The second term on the left of the second equality is combined with
the normalization of the complex Gaussian. Adding the original count
$-\tfrac12n_{\mathrm{phys}}\log\mathsf t$ leaves
$-\tfrac12n_{\mathrm{phys}}\log x$. Thus the auxiliary number
$\mathsf t$ disappears from this logarithmic normalization, and the
total $G_\tau$ has zero action coefficient. The increment still
contains the Gaussian coefficient and the coefficient of the
interaction family. Assigning a coefficient to a separate local term
of $G_\tau$ would break this cancellation.

\emph{Derivatives of a finite prefix.}
This formulation is needed for more than real coupling estimates.
Hold the real choices at numbers $\theta_0,\ldots,\theta_k$ and
allow the action weights to vary on
$|x_i-\theta_i|<\sigma_\star\theta_i$. Iterating a holomorphic
one-step map on the common domains gives a holomorphic ordinary
correction $f_{k+1}(x_0,\ldots,x_k)$ in the absolute-radius chain.
The required bound is
\[
 \sup|f_{k+1}|\le
 C\theta_k^{-1/2}(\log\theta_k)^{\mathfrak p_0}.
\]
The remainder contribution is estimated separately on running radii;
it must not be included in this bound merely because it contributes
to the same final increment.

For a point in the half-polydisc and a complex direction $v$, restrict
this function to the line $x+zv$. If
$m(v)=\max_{i\le k}|v_i|/(\sigma_\star\theta_i)>0$, the line
stays in the full polydisc for $|z|<(2m(v))^{-1}$. Cauchy's
one-variable estimate therefore gives
\begin{equation}\label{eq:rg-directional-cauchy}
 \left|\sum_{i\le k}v_i\partial_i f_{k+1}(x)\right|
 \le C\theta_k^{-1/2}(\log\theta_k)^{\mathfrak p_0}
                  \max_{i\le k}\frac{|v_i|}{\sigma_\star\theta_i}.
\end{equation}
For the zero direction the inequality is immediate. The estimate
costs no factor equal to the number of levels: it uses a single
complex line, rather than adding separate Cauchy estimates in all
coordinates. This is the form useful when a source perturbs a whole
coupling prefix at once.

There is no geometric decay in $k-i$ in
\eqref{eq:rg-directional-cauchy}. Boundedness on a polydisc alone
cannot give it; a function proportional to $x_0/\theta_0$ is a
counterexample. Any assertion that the chain forgets a particular
old perturbation has to use its recursive weighted dependence, in
addition to holomorphy. The present first-passage construction uses
the ordinary size bound and the separate remainder sum; it does not
require the stronger tuned-cutoff comparison.

\emph{One local construction on all volumes.}
The input to a step includes its local realization as well as its
norm. It consists of the decoupled covariance and background response
on each region, the localized new interactions, the rule evaluating
older families at the fluctuating background, and the localized
Gaussian logarithm. These four choices must be transported together.
For example, an old potential on $Y$ uses only family members on
subdomains of $Y$ and their coefficients; changing a representative
outside $Y$ must not change that potential.

Fix these local formulas once, with their margins and constants.
On a nonwrapping region with a common lift to $\mathbb Z^4$, the
same formulas then define the corresponding term on either torus.
The realization on smaller analytic radii is the restriction of the
one on larger radii. Evaluating those formulas on $\mathbb Z^4$
defines the companion of the new family and permits the next
coefficient to be extracted before a volume is selected. Uniqueness
here is relative to the fixed local realization. In particular the
Gaussian pointing can depend on that realization even though its
sum is the same Gaussian logarithm. No equality of every individual
local term under unrelated localization choices is needed or asserted.

The uniform analytic map and its compatible local realization are
therefore parts of the same construction. Once their estimates hold,
the preceding identities explain how it handles arbitrary input lists,
frozen real choices and common coefficients. They do not derive the
required background or walk estimates from holomorphy alone.


\subsubsection*{The two chains and the remainder load.}
The separation of absolute and running radii can be carried
through the iteration quantitatively. Denote the fixed real
inverse-coupling sequence by $\theta_m$; the real tests and
field scaling of step $m+1$ use $\theta_m$. The ordinary
chain contains only pointed families and their own
counterterms. On absolute radii its one-step small factor is
$C\delta(\theta_m)$, and the Gaussian family contributes an
order-one norm. The complex gauge correction contributes
$K_{\mathbf M}\sigma_\star$. Choose
\[
 E_0\ge\max\{2E_N,4K_{\mathbf M}\sigma_\star\},\qquad
 N_\star\ge E_0\sum_{n\ge1}\omega_n+C_{\mathrm{slot}}.
\]
Then the output budget is one half of $E_0$ for the
Gaussian family, one quarter for the fixed gauge correction,
and one quarter for the interaction, after the coupling
ceiling has been decreased. The geometric age weights make
the bound on the old pointed load independent of the number
of levels. The constants $K_{\mathbf M},\sigma_\star$ are
chosen before $E_0,N_\star$. Thus decreasing the subsequent
ceiling on $\mathsf s$ does not change the constants it has
to overcome.

The formula of a local term contains no analytic radius.
Its construction on smaller cores is the restriction of
the same construction on larger cores. Restrict the ordinary
chain to the running cores, and on those cores form a
second chain containing the aged remainders as well. Both
chains use the same coefficient $x$, the same real choices,
the same normaliser and the same free gauge correction.
Their difference therefore consists precisely of terms with
at least one marked input descended from a remainder,
including the counterterm of each marked pointed member.
The Gaussian and free gauge families cancel from the
difference. The Lipschitz bound furnished by the marked
connected sum applies to this difference.

Here the hypotheses on the real sequence are the already
constructed prefix bounds
\[
 \theta_j\ge\theta_m-c_5+\lambda_\flat(m-j),\qquad
 \theta_j\le\theta_m+C(m-j),\qquad j\le m,
\]
with a fixed lower floor. The radius ratio obeys
\[
 \varrho_{jm}^2\le
 C\frac{\theta_j}{\theta_m}
             \left(\frac{\log\theta_m}{\log\theta_j}\right)^{2q}
 \le C\frac{\theta_j}{\theta_m}.
\]
The bounded reverse drift is used in the last inequality
when $\theta_j<\theta_m$. A remainder of level $j$ has size
at most $C\theta_j^{-\kappa_0/2}$. Put
$p=(\kappa_0-2)/2>1$. Its total aged load at step $m+1$
is bounded by
\begin{equation}\label{eq:deep-aged-load}
 \begin{split}
 \sum_{j\ \mathrm{aged}}\varrho_{jm}^2
                               \theta_j^{-\kappa_0/2}
 &\le C\theta_m^{-1}\sum_{j\le m}\theta_j^{-p}\\
 &\le C\theta_m^{-1}
          \sum_{n\ge0}(\theta_m-c_5+\lambda_\flat n)^{-p}
 \le C_{\rm h}\theta_m^{-p}.
 \end{split}
\end{equation}
The last inequality follows by comparison with the integral
of $t^{-p}$; its constant is bounded by a fixed multiple of
$1+[\lambda_\flat(p-1)]^{-1}$. This is the loss of one
inverse-coupling power caused by summing remainders over
levels. No decay with age has been assigned to an individual
remainder. If the inverse coupling were held constant, that
load would grow with the number of levels.

Let $u_i$ be the running-core norm of the marked output
formed at step $i$. The one-step Lipschitz estimate gives
\begin{equation}\label{eq:deep-marked-recursion}
 u_{m+1}\le K_{\mathbf M}\mathsf s(\theta_m)
       \left[\sum_{i\le m}\omega_{m+1-i}u_i
                                    +C_{\rm h}\theta_m^{-p}\right].
\end{equation}
Suppose inductively $u_i\le A\theta_{i-1}^{-p}$.
The reverse-drift bound and the floor give
$\theta_{i-1}^{-p}\le2\theta_m^{-p}$, after a fixed
adjustment of the floor. Thus the bracket in
\eqref{eq:deep-marked-recursion} is at most
$[2A\sum_n\omega_n+C_{\rm h}]\theta_m^{-p}$.
Choose $A\ge1$ and then the ceiling so that
\[
 K_{\mathbf M}\sup\mathsf s
             (2\sum_n\omega_n+C_{\rm h})\le1.
\]
This proves $u_{m+1}\le A\theta_m^{-p}$ and closes the
induction. The factor $\mathsf s$ may be only a negative
power of a logarithm; it need only be uniformly small, since
the convolution kernel is summable.

Extracting a coefficient costs two derivatives on the radii
of creation. For a marked term that cost is at most
$C\theta_m$, up to favourable logarithms. Consequently
\[
 |b^{\rm marked}_{m+1}|
       \le C\theta_m u_{m+1}
       \le C\theta_m^{-(\kappa_0-4)/2}.
\]
On absolute radii the derivative cost is fixed, and the
ordinary correction has the bound
\[
 |r_{m+1}|\le C\theta_m^{-1/2}(\log\theta_m)^{\mathfrak p_0}.
\]
These are the two different estimates for $q_{m+1}$ and
$r_{m+1}$ in \eqref{eq:flow-estimates}. Applying the running
radius derivative cost to the ordinary chain would instead
lose its smallness.

A remainder is admitted to the second chain only after
its own partition is preserved by the next lattice. Its
cube side is $M\check R_j$, a power of $L$. Once
$M\check R_j\le L^{m+1-j}$, the next lattice is a sublattice
of its cube-centre lattice, and all required translations
preserve the remainder's partition. Before that age it is
kept outside the regular exponent and passed through the
separate remainder and boundary slots. This delay is a
symmetry requirement. It is distinct from the waiting depth
before a large-field component is healed.

The induction also retains its order of construction. The
regular output and its coefficient at step $m+1$ use only
$\theta_0,\ldots,\theta_m$. The coefficient fixes the next
inverse coupling before the next large-field cube size and
tests are chosen. The created family has a fixed enlargement
of the radius at $\theta_m$, sufficient for both adjacent
running radii. Hence neither the coefficient nor its
definition depends on a next-level construction whose
admissibility it is supposed to prove.

\subsubsection*{Varying the whole coupling prefix.}
There is a sharper way to state the dimension-independent
directional estimate. Suppose a holomorphic scalar $F$ is
bounded by $B$ on $\prod_iD(0,r_i)$. For an interior point
$\xi$, choose disc automorphisms $\phi_i$ with
$\phi_i(0)=\xi_i$. Their derivatives at zero have moduli
$r_i(1-|\xi_i/r_i|^2)$, and their phases can be chosen so
that all terms in
\[
 \frac d{dz}F(\phi_1(z),\ldots,\phi_n(z))\bigg|_{z=0}
\]
have the same phase. Cauchy's estimate on that one unit
disc gives
\begin{equation}\label{eq:deep-polydisc-slope}
 \sum_i r_i(1-|\xi_i/r_i|^2)|\partial_iF(\xi)|\le B.
\end{equation}
For a family-valued map, apply this calculation at a fixed
domain, base point and pair, and then take its weighted
supremum when estimating a direction. It introduces no
factor equal to the number of earlier couplings.

In the half-polydisc of an ordinary coefficient, take
$r_i=\sigma_\star\theta_i$ and
$B=C\theta_m^{-1/2}(\log\theta_m)^{\mathfrak p_0}$.
The slope sum is then summable along the lower-drift
sequence after one further inverse-coupling factor is
taken in a fixed-size perturbation direction. This is the
analytic estimate used when a source shifts all the
earlier action weights simultaneously. It does not prove
an individual derivative decreases geometrically with
its age; such a conclusion would require additional
recursive information.

Finally a constant source can be carried in two ways.
It may change the coefficient $x$ to $x-\zeta$, or it
may be inserted as a weight member $\zeta A$ while the
real tests are held fixed. Where both constructions are
admitted, the exact integral depends on them through the
same exponent $-xA+\zeta A$. Their Gaussian conventions
give totals differing only by the explicit number
\[
 \tfrac32|\mathfrak X|
            [\log(x/\mathsf t)-\log((x-\zeta)/\mathsf t)].
\]
Continuity along $t\zeta$, starting at zero, fixes the
logarithm branch. The coefficients of the two totals agree,
since the difference is field independent. Individual
local terms need not agree. The source induction uses this
coefficient identity together with the specified local
construction, rather than an identity between unrelated
localisations.

\subsubsection*{Constrained determinants and the one-loop sum.}
The passage from many normalisers to one determinant uses
an exact finite-dimensional identity. Let $K$ be a symmetric
form and $m$ a full-rank linear map, with $K$ positive on
$\ker m$. Choose orthonormal columns $L$ spanning $\ker m$
and $Z$ spanning its orthogonal complement. The constrained
Gaussian integral is
\[
 \int_{\mathbb R^n}\delta(mv)e^{-\langle v,Kv\rangle/2}\,dv
   =\frac{(2\pi)^{(n-r)/2}}
        {\det(L^{\mathsf T}KL)^{1/2}\det(mm^{\mathsf T})^{1/2}},
 \qquad r=\mathrm{rank}\,m.
\]
Indeed $v=La+Zb$ is an orthogonal change of variables,
$mL=0$, and integrating $\delta(mZb)$ sets $b=0$ with
Jacobian $|\det(mZ)|^{-1}$. The corresponding saddle
matrix is
\[
 \mathbb S(K,m)=
       \begin{pmatrix}K&m^{\mathsf T}\\m&0\end{pmatrix}.
\]
In the same coordinates its determinant is
\begin{equation}\label{eq:deep-saddle-determinant}
 |\det\mathbb S(K,m)|
         =\det(L^{\mathsf T}KL)\det(mm^{\mathsf T}).
\end{equation}
The mixed blocks of $K$ cancel in this determinant. Thus
the constrained integral is a constant times
$|\det\mathbb S(K,m)|^{-1/2}$. This formula keeps the
constraint local and avoids solving it by a nonlocal
projection before taking determinants.

Near the trivial field, construct the critical branch
$p_j(W)$ for one step in the axial gauge. The composite
$c_k=p_1\circ\cdots\circ p_k$ is a critical configuration
over the $k$-fold average. Its Hessian on the remaining
fibre is positive: a closed axial form with zero average
is zero. The implicit function theorem supplies this
branch on a neighbourhood that may depend on the depth
and torus. The uniqueness of the regular critical orbit
identifies it with the background branch where both are
defined. For the one-loop coefficient only derivatives
at the identity are used; the uniform estimates required
below concern their kernels, rather than a common
nonlinear neighbourhood for all depths.

Successive Gaussian integrations on this branch compose
exactly. For a positive block form,
\[
 \det\begin{pmatrix}A&B\\B^{\mathsf T}&D\end{pmatrix}
       =\det A\,\det(D-B^{\mathsf T}A^{-1}B),
\]
and the second factor is also the quadratic form obtained
by minimizing over the first variables. Repeating this
identity shows that the product of the first $k$
normalisers is one constrained Gaussian integral in the
hierarchical axial gauge. If $\Phi_j$ is the logarithm
of the $j$th normaliser relative to its value at the
identity and $\Psi_j$ is the combined logarithm, then
\[
 \Psi_j=\Phi_j+\Psi_{j-1}\circ p_j.
\]
The coefficient of $\Psi_k$ is $\sum_{j\le k}b_j^0$.
In this identity the action coefficient is exactly
marginal: each curvature has scale $L^{-2}$, so the two
curvature factors give $L^{-4}$, cancelled by the $L^4$
site count. The nonzero aliases in its symbol begin at
higher momentum order and do not change the quadratic
coefficient.

To compare with a translation-invariant massive form,
change the gauge and the constraint together. Let $Q'$
be the site average and $P_U$ the orthogonal projection
onto $\ker Q'(U)$. The gauge condition
$P_UD_U^*B=0$ removes the residual directions
$D_U\ker Q'(U)$. The exact relation
\[
 Q_k^\sharp(U)D_U=D_{\mathcal M^kU}Q'(U)
\]
puts those directions in the kernel of the straightened
constraint. The finite-dimensional gauge Jacobian and
\eqref{eq:deep-saddle-determinant} give the gauge and
constraint replacement in one formula:
\begin{equation}\label{eq:deep-one-loop-saddles}
 \begin{split}
 \Psi_k={}&\bigg[-\tfrac12\log|\det\mathbb S(
       \Delta_U^{\mathrm W}+\mathfrak J_U+D_UP_UD_U^*,
                                        Q_k^\sharp(U))|\\
       &\hspace{13mm}
       +\log|\det\mathbb S(\Delta_U,Q'(U))|\bigg]_{U=1}^{U=c_k(W)}.
 \end{split}
\end{equation}
Here $\mathfrak J_U$ includes the current correction to
the off-critical action Hessian and the multiplier paired
with the curvature of the constraint. These corrections
are required for the constrained Hessian represented by
the left side. They cannot be discarded by calling the
ambient Wilson Hessian gauge invariant. The projection
$P_U$ is also different from the propagator projection
whose range is $\Delta_U\ker Q'(U)$.

The comparison of \eqref{eq:deep-one-loop-saddles} with
the massive determinant has two sources of uniformly
bounded terms. A term containing a constraint perturbation
factors through the unit lattice. Its trace therefore
runs over a bounded number of dimensions per unit cell,
even though a Bloch fibre at bare spacing $\eta_k$ has
$O(\eta_k^{-4})$ components. The remaining resolvent
expansion is stopped only after three propagator factors
have occurred. The Schatten inequality then uses the
finite four-dimensional integral
$\int(1+|q|^2)^{-3}\,dq$. Terms with one or two
propagators are retained as the mass derivatives in
$\mathcal L^\sharp$; their traces are not uniformly
bounded by this integral.

All shifted resolvents in this calculation must lie
in one fixed complex-momentum strip. In that strip,
bounded normalised symbols and the real coercivity
give uniformly bounded inverses by the conjugation
argument below. It follows that the comparison
remainder is holomorphic in the external momentum
on one fixed polydisc. Its quadratic coefficient is
bounded by Cauchy on that polydisc. The actual
factorisation and resolvent expansion with these
uniform bounds are the substantive determinant
comparison input. The displayed saddle identities
establish the exact constrained object to which it
must apply.

The Gaussian part is handled as a sum, not by evaluating each
step independently. Along a background, the product of the first
$k$ normalisers is a single Gaussian integral over the tangent
space to $\{U:\mathcal M^kU=V\}$. Its logarithm has coefficient
$\sum_{l\le k}b_l^0$. Changing gauge in that Gaussian integral
expresses it through the gauge and site Laplacians with the
average constraints. The comparison replaces those constraints
by a mass. If $\mathcal L(m)$ is the massive determinant
functional in level-$k$ units, use
\[
 \mathcal L^\sharp=
 \mathcal L(1)-\mathcal L'(1)+\tfrac12\mathcal L''(1).
\]
This is its Taylor polynomial of order two at mass one,
evaluated at zero. Modulo a polarisation term $O(|p|^3)$, the
difference from the constrained Gaussian functional is analytic
and bounded on a momentum polydisc of radius independent of
$k$. Cauchy's estimate therefore bounds the change in its
quadratic coefficient by a constant independent of $k$.

The uniform polydisc is obtained from the flat-background
operators. Their Bloch symbols are coercive after normalisation
by $(\Delta+1)^{-1/2}$. Adding a small imaginary momentum is
equivalent to exponential conjugation in position space.
The symbols remain invertible on one fixed strip, and the
resolvent comparison is uniformly bounded there. Three
propagator powers are trace class per unit volume in dimension
four. The two lower mass-derivative terms are kept explicitly
in $\mathcal L^\sharp$, rather than estimated by a trace-class
bound that does not apply to them.

The uniformity in the momentum strip
has a concrete operator meaning. Let
$T_k(\varkappa)$ be a flat-background
Bloch symbol of an inverted operator
in the comparison and normalise it by
the two factors of
$(\Delta+1)^{-1/2}$. Write the result
as $\mathcal T_k(\varkappa)$.
For real momentum the required
coercivity estimate is
\[
 cI\le\mathcal T_k(\varkappa)\le CI,
\]
uniformly in $k$. On the initial
fixed strip the symbols are
holomorphic and bounded in operator
norm. Cauchy's estimate on a
smaller strip gives
\[
 \|\mathcal T_k(\varkappa+iq)
                    -\mathcal T_k(\varkappa)\|
                            \le C_1|q|.
\]
Take $|q|<c/(2C_1)$. The inverse
exists there by a Neumann series
and has norm at most $2/c$.
This is one fixed strip for every
level, including levels at which
the Bloch fibre has more
components. It uses the
normalised operator norm,
rather than estimates on
individual matrix entries.

After the lower mass derivatives
have been retained explicitly,
the remainder comparison has
three resolvent powers. Their
trace per unit volume is
uniformly bounded because
\[
 \int_{\mathbb R^4}(1+|q|^2)^{-3}\,dq<\infty.
\]
The same majorant bounds the
rescaled lattice integrals,
whose Brillouin regions grow
as the bare spacing decreases.
The first and second resolvent
powers would instead diverge
in that comparison. Thus the
two lower mass-derivative terms
cannot be hidden in this
uniform trace estimate.

For external momentum $p$ in
a still smaller polydisc, the
shifted factors in the
polarisation comparison remain
in the common inverse strip.
The trace majorant gives a
bound $C_{\mathrm{rem}}$ on
the analytic difference
symbol, modulo $O(|p|^3)$.
Its quadratic action
coefficient is consequently
at most
$C_{\mathrm{rem}}r_\pi^{-2}$
by Cauchy in two distinct
momentum directions.
The reduction of the
constrained determinant to
these remainders and the
massive determinant, with
this uniform majorant, is
the substantive comparison
input. The strip and trace
argument explains the
cutoff-independent constant
once that reduction is
available.

In bare-lattice units the determinant functional to be estimated
is
\[
 \log\det(\Delta_U+z)
 -\tfrac12\log\det(\Delta_U^{\mathrm W}+D_UD_U^*+z),
 \qquad z=\eta_k^2.
\]
Evaluating its quadratic action coefficient in a constant
abelian field gives a function $\widetilde\beta(z)$.
The determinant estimate used here is
\begin{equation}\label{eq:rg-mass-coefficient}
 \left|\widetilde\beta(z)
      -\frac{11}{6\pi^2}\log\frac1z\right|
 +|z\widetilde\beta'(z)|+|z^2\widetilde\beta''(z)|
 \le C_{\mathrm{uv}},\qquad 0<z\le1.
\end{equation}
The coefficient of $\mathcal L^\sharp$ is
$\widetilde\beta(\eta_k^2)
-\eta_k^2\widetilde\beta'(\eta_k^2)
+\frac12\eta_k^4\widetilde\beta''(\eta_k^2)$.
Since $\log(\eta_k^{-2})=2k\log L$, this estimate and the
bounded comparison error give the first bound in
\eqref{eq:flow-estimates}. The mass coefficient estimate and
the fixed-strip comparison are substantive analytic inputs;
the summation of the error terms is performed only after
these uniform estimates are available.

For the drift calculation, choose the coupling ceiling so that the small-field errors
are smaller than a fixed fraction of $c_0$ and the remaining
errors are uniformly small. The bounds on sums of $b_l^0$
then give constants $B_\sharp,\lambda_\sharp>0$ and
$C_\sharp<\infty$ such that
\begin{equation}\label{eq:flow}
 b_{k+1}\le B_\sharp,\qquad
 \lambda_\sharp(j-i)-C_\sharp
 \le x_i-x_j\le B_\sharp(j-i)\quad(i\le j).
\end{equation}
To see the drift estimate directly from \eqref{eq:flow-estimates}, put
$\mathsf B_n^0=\sum_{l=1}^n b_l^0$. The first estimate in
\eqref{eq:flow-estimates} implies
\[
 \big|\mathsf B_j^0-\mathsf B_i^0-c_0(j-i)\big|\le2C_0.
\]
Choose the coupling ceiling so that
$|r_l|+|q_l|\le\epsilon_\gamma<c_0/2$ at every
constructed level. Summing \eqref{eq:increment} then gives
\[
 x_i-x_j\ge(c_0-\epsilon_\gamma)(j-i)-2C_0.
\]
Also $|b_l^0|\le c_0+2C_0$, so
$b_l\le c_0+2C_0+\epsilon_\gamma$.
These two consequences provide fixed admissibility
constants. They use no sign of an individual $b_l^0$.

The lower estimate uses sums of Gaussian increments and does
not assert that every individual Gaussian increment is
positive. One may take $C_\sharp=2C_0$ and, after decreasing
the coupling ceiling, $\lambda_\sharp\ge c_0/4$.
The large-field induction uses these bounds to sum its
remainders. Conversely, the next increment is formed from
the terms at levels already constructed, before the next
large-field operation uses the updated coupling. This order
closes the induction without invoking ultraviolet stability
as a proof of its own flow bound.

\subsection{One coupling flow on all tori and first passage}

The coefficient in \eqref{eq:rg-polarisation} is defined from
local families on arbitrarily large tori. The local choices in
the background and walk constructions are made by one rule
before choosing the torus. On a domain that lifts injectively
with a margin to $\mathbb Z^4$, the formula for a term is the
same formula in either setting: it uses local backgrounds,
walks, connected polymers, earlier terms and the already
determined coefficients. Induction therefore defines one
family $\widehat E^{(j)}$ that serves all tori. In the
large-field representation the same assertion holds inside
the regular region; the history there has empty standing
trace after healed components have been summed into local
terms, as explained in \cref{sec:large-fields}.

It is important that the induction is simultaneous over all
tori. The family defining $b_j$ includes domains larger than
any one torus. There is no convergent total interaction on
the whole infinite lattice; its Hessian kernel is instead
the absolutely convergent sum defining $\Pi^{(j)}$.
The coefficient agrees with the large-volume limit of the
coefficient of torus totals. Non-wrapping terms agree
locally, wrapping terms decay exponentially in the torus
side, and the background response converges to its
infinite-lattice response. These three estimates justify
that identification. In particular, $b_j$ depends only on
the previous couplings and the fixed choices, rather than
on $\ell$.

The geometric restriction on the torus is also uniform in
the cutoff. Let $\check R_j=R(\min_{i\le j}x_i)$ be the
large-field cube factor, where $R(x)$ is the least power of
$L$ not smaller than $(\log x)^{\mathfrak r}$ for a fixed exponent
$\mathfrak r$.
The levelwise torus condition is
\begin{equation}\label{eq:prefix-volume}
 \ell L^{K-j}\ge2L^{11}M\check R_j.
\end{equation}
The flow bound gives $\check R_j\le L\check R_{j+1}$,
so the condition at the last level implies the preceding
ones. In the terminal reference window it is enough to take
\begin{equation}\label{eq:flow-volume-floor}
 \ell\ge\ell_\star(g)
 :=2L^{11}M R(g^{-2}+B_\sharp).
\end{equation}
There is no dependence on $K$ in this lower bound. This is
the volume uniformity needed before taking the
infinite-volume limits.

For $0<g\le\gamma$ and $\beta>g^{-2}+B_\sharp$, stop the flow at
\begin{equation}\label{eq:first-passage}
 K(\beta,g)=\min\{k:x_k\le g^{-2}+B_\sharp\}.
\end{equation}
Before exit, the upper increment bound implies
$x_{k+1}>g^{-2}\ge\gamma^{-2}$, so the next step stays
within the small-coupling domain. The positive lower drift
ensures that exit occurs. At exit,
$g^{-2}<x_K\le g^{-2}+B_\sharp$, and
\[
 \frac{\beta-g^{-2}}{B_\sharp}-1\le K
 \le1+\frac{\beta-g^{-2}-B_\sharp+C_\sharp}
                 {\lambda_\sharp}.
\]
Thus increasing $\beta$ gives unbounded exit depths. The
physical unit is chosen so that this terminal scale has
length one; the bare spacing is then $L^{-K(\beta,g)}$.
This selects the ultraviolet sequences used by
\cref{thm:main}. It does not fix the last inverse coupling
at a single point of the reference window.

At depth $s$ below the last level, the flow has the
cutoff-independent envelope
\[
 g^{-2}+\lambda_\sharp s-C_\sharp
 <x_{K-s}\le g^{-2}+B_\sharp(s+1).
\]
The errors in \eqref{eq:flow-estimates} can therefore be
summed in depth rather than bounded by their value at the
largest coupling. Their sum gives
\[
 |x_{K-s}-x_K-c_0s|
 \le C+C\sqrt{s}\,
                 \log^{\mathfrak p_0}(g^{-2}+s)=o(s)
 \qquad(s\longrightarrow\infty, g\text{ fixed}).
\]
In particular,
$\beta/\log(1/a_K)\longrightarrow 11/(3\pi^2)$ and
$g_0=\beta^{-1/2}\longrightarrow0$.
At a prescribed physical separation the same envelope
controls the coupling at the corresponding depth. This is
the use of the flow in the two-point comparison. The
subsequence construction and the subsequent phase
selection use uniform bounds and do not require tuning
or uniqueness at fixed renormalised coupling.

\section{Large fields, stored variables and the sum over histories}\label{sec:large-fields}

The fluctuation expansion applies only while both the background
curvature and the fluctuation are small. The Wilson measure is not
supported on that set. At each scale the construction therefore inserts
local decompositions of unity, retains the variables where a test fails,
and continues the small-field integration elsewhere. This section
describes how the retained regions are controlled. The estimates are
also needed with sources: a small probability for one bad region is
insufficient unless it remains small under later integrations and
complex changes of the exterior field.
The removal and exponentiation follow Ba\l aban's large-field method
\cite{Balaban1989}, whose Theorem~1 assumes small effective couplings.
The estimates must remain local in the recorded history, control
exterior suprema inside the component integrals, persist on the
complex-source domains, and pay a prescribed weight for the age of
every live region. We construct the operations in this order so that
these four requirements enter at the estimates where they are used.

\subsection{Tests and nested regions}

Write $x_j=g_j^{-2}$ and choose thresholds
\[
 p_0(g)=A_0(\log g^{-2})^{\mathfrak p_0},\qquad
 \varepsilon_j=g_jp_0(g_j),\qquad
 \delta_j=(A_1/A_0)\varepsilon_j.
\]
The constants $A_1/A_0$, and later the upper bound on the couplings,
are small. A plaquette test bounds the localised background curvature
by $\varepsilon_j$ in the appropriate scale. An approximate fluctuation
test compares the old block field with the average of that localised
background. Where these tests pass, the true multiscale background is
formed, and the logarithmic fluctuation can be tested as well.

An indicator includes existence of all the objects used to define it.
If a logarithm or regular background is undefined, the indicator is
zero and its complementary indicator is one. Thus each test gives an
identity on the entire configuration space. Small-field regularity is
used to estimate the terms that result, rather than assumed before
the decomposition is made. Gauge-fixing insertions are likewise
normalised identities on every configuration, with their specified
extension where local coordinates do not exist.

Let $\Omega_j$ be the region where the background and fluctuation are
defined and let $\Lambda_j\subset\Omega_j$ be the region on which
the fluctuation is integrated. The complements are enlarged around
failed tests and around previously retained regions. Their boundaries
are separated by several layers of cells. The resulting nesting has
the form
\[
 T_\ell\supset\Omega_1\supset\Lambda_1
 \supset\Omega_2\supset\Lambda_2\supset\cdots.
\]
At scale $j$ the cells have side $M\check R_j$, where
$\check R_j$ is the smallest power of $L$ at least
$(\log\min_{i\le j}x_i)^{\mathfrak r}$, with $\mathfrak r\ge2$.
Taking powers of $L$ keeps the partitions nested. The running minimum
makes $\check R_j$ nonincreasing; the bounded change of $x_j$ ensures
that consecutive drops in $\check R_j$ do not occur at sufficiently
small coupling.

The boundary layers have an analytic purpose. Let $\mathcal A_j$ be
the unscaled logarithmic fluctuation. Conditionally on its stored
part outside $\Lambda_{j+1}$, its mean inside is
\begin{equation}\label{eq:conditional-mean}
 m_j=-C^{(j)}(\Lambda_{j+1})\Delta^{(j)}
                  \mathcal A_j|_{\Omega_{j+1}\setminus\Lambda_{j+1}},
\end{equation}
where $C^{(j)}(\Lambda)$ is the inverse of the restricted quadratic
form. A large stored field close to the boundary could shift this mean
past an interior threshold. The construction therefore also tests a
collar of stored bonds and the rim touching the integration region.
The bounds there are respectively
$\delta_j/(4\Theta_1)$ and $\theta_c\delta_j$, with
$0<\theta_c<1$. Exponential decay of the conditional covariance gives
\begin{equation}\label{eq:collar-mean}
 |m_j|\le
 C\big(\Theta_1^{-1}+\theta_ce^{-c\varpi_T}
                  +e^{-cM\check R_{j+1}}\big)\delta_j
 <\tfrac12\theta_c\delta_j.
\end{equation}
Here $\varpi_T$ is the collar width. First choose
$\theta_c\Theta_1$ and $\varpi_T$ sufficiently large, and then reduce
the coupling ceiling. Away from the collar and rim, the retained
tests give $|\mathcal A_j|\le C\delta_j$; this bound is needed for
the far-field tail in \eqref{eq:collar-mean}.
Each failed collar or rim test enlarges the
retained region. Every round examines new bonds or cells, so this
enlargement terminates on a finite torus and is still a decomposition
of unity.

The geometry is formed before an estimate is applied. At step
$t=j+1$, let $P_t,Q_t,R_t$ be the failed plaquette,
approximate-fluctuation and true-fluctuation test sets.
The first two are defined using localised backgrounds.
The set $\Omega_t^c$ is obtained by enlarging their union
with the hull of the old live region, and $\Lambda_t^c$
by enlarging $\Omega_t^c\cup R_t$. A boxing rule replaces
sufficiently small components by boxes. Collar and rim
tests can require further enlargement. Each new round
examines previously unexamined cells, so the rule is total
on a finite lattice. At a point where the true background
is not yet defined, its failed existence test has already
placed that point in the retained region.

The distinction between approximate and true fluctuations
is needed for locality. The approximate test determines
whether the global regular background can be formed.
The true test, once that background exists, determines
which Gaussian variables are integrated. The two cannot
be interchanged when a history is compared on a local
window.

On a strip next to $\partial\Lambda_t$, a bondwise
change of variables replaces the true background in the
kept tests by its localised version. In local coordinates
this change has the form
\[
 W'\longmapsto
 W'e^{-i\theta_{\mathrm{co}}(W')H}e^{iH}V^{\mathrm{loc}},
       \qquad |H|\le C\varepsilon_j e^{-\check R_j}.
\]
The cutoff $\theta_{\mathrm{co}}$ equals zero in the
small-fluctuation range and one beyond the kept range.
For constant cutoff this is a translation with unit
Jacobian. Its nonconstant part is a local boundary
Jacobian. It belongs to the boundary interaction.
The construction requires that after this substitution
the retained tests depend only on the stored variables.
An old test reading a pivot bond uses the value obtained
from that bond's individual averaging constraint, rather
than the value at the averaged background.

The new tests and the carried constraints imply the old
small-field tests through the estimate
\begin{equation}\label{eq:mrg-test-implication}
 CL^{-2}+C\delta_j/\varepsilon_j
                      +Ce^{-\check R_j}<1.
\end{equation}
The three terms are, respectively, the curvature scaling,
the fluctuation displacement, and the effect of the
artificial boundary of a localised background. Choose $L$
first, the ratio $A_1/A_0$ second, and the coupling ceiling
last. This order concerns the tests that remain on the
support of a term, including its stored constraints. It
does not assert that a new plaquette test alone implies
every old test.

\subsection{Exact measures and effective actions}

A history records the outcomes of the tests, the regions at every
scale, and the operations performed on their components. Stored
variables are part of the history's state. At the measure level each
move is a test, a change of variables, an integration, or a draw from
a normalised auxiliary law. If $K_h(U,d\xi)$ is the positive kernel
associated with a complete sequence of outcomes, exactness gives
\begin{equation}\label{eq:history-kernels}
 \sum_hK_h(U,\mathcal V_k)=1,\qquad
 \tau_h[\nu_0](d\xi)=\int K_h(U,d\xi)\,\nu_0(dU).
\end{equation}
In particular,
\[
 \sum_h|\tau_h[\nu_0]|(\mathcal V_k)\le|\nu_0|(\mathcal V_0),
 \qquad
 |\tau_h[w\mu_0]|\le\|w\|_\infty\tau_h[\mu_0]
 \quad(\mu_0\ge0).
\]
Here $\mathcal V_k$ is the level-$k$ state space of the current
block field, its history and all retained stored data.
These identities do not depend on convergence of an expansion. The
expansions provide formulas and estimates for the measures on the
supports of the tests. Freezing the real construction freezes the
kernels as well, which is the reason the same procedure can later
transport a complex source.

For a test the two kernels are multiplication by $\chi$
and by $1-\chi$. A measurable map supplies a deterministic
kernel, and a reference draw supplies a probability kernel.
Integration and removal of an auxiliary coordinate are
marginalisations. Composition of these kernels shows
\eqref{eq:history-kernels} by induction, because the sum of
the outcome masses is one at each move. Undefined analytic
objects cause a complementary outcome, not a missing mass.

The total-variation inequality follows at the same level.
For every bounded measurable $F$ with $|F|\le1$,
\[
 \left|\int F\,d\tau_h[\nu_0]\right|
 \le\int K_h(U,\mathcal V_k)\,|\nu_0|(dU).
\]
Taking the supremum over such $F$, then summing over $h$,
gives the first inequality following
\eqref{eq:history-kernels}. Replacing $|\nu_0|$ by
$|w|\mu_0$ gives the domination of a bounded weight.
For a finite bare lattice the source weight $e^{A(w,U)}$
is bounded, so the same frozen kernels transport it.
No source enters their tests, normalisers or auxiliary
probability laws. Later holomorphic bounds concern the
resulting complex integrals, not complex probability
kernels.

These identities also construct the observable that must
be retained when a reference law changes a block field.
The joint measure of bare field, history and current
block state is
\[
 \mu_0(dU)K_h(U,d\xi).
\]
Conditional expectation with respect to $(h,\xi)$ gives
the transported version $F^{(h)}(\xi)$ of a bounded bare
observable. It satisfies $|F^{(h)}|\le\|F\|_\infty$
and the exact insertion formula
\begin{equation}\label{eq:mrg-history-insertion}
 \langle F\rangle_{\mu_0/\mu_0(1)}
   =\frac{\sum_h\int F^{(h)}\,d\tau_h[\mu_0]}
          {\sum_h\tau_h[\mu_0](\mathcal V_k)}.
\end{equation}
Where all the operations act away from the bonds on which
an averaged observable depends, its transported version
is that same function of the block field. Inside a
replaced box this identification is generally false.
Gauge invariance permits the usual representative changes;
it does not make a function of a replaced plaquette
independent of the replacement.

The kernel construction and the expansion describe
different aspects of the same measure. The kernel
statement is exact before estimates; the expansion
gives the following presentation of its terms.
Let $\mathcal S$ record the action regions and
their standing marks. Then
\begin{equation}\label{eq:xrg-effective-density}
 \begin{split}
 \rho_k(V_k)
 &=\sum_{\mathcal S}\mathbf T'_k(\mathcal S)
       \left[\chi_{k,\mathcal S}
           e^{A_{k,\mathcal S}(U_k(\mathbf V))}\right],\\
 A_{k,\mathcal S}(U)
 &=-A(x_k(\cdot),U)+\mathbf E_k(U)
       +\mathbf R_k(U)+\mathbf B_k(U,\underline{\mathfrak f})
        -\mathcal E_{k,\mathcal S}.
 \end{split}
\end{equation}
Here $\underline{\mathfrak f}$ denotes the stored fluctuation fields.
The operator $\mathbf T'_k(\mathcal S)$
integrates the fields stored in the live
components, with their delta constraints,
tests, gauge fixing and conditional Gaussian
forms. It retains the integral operations of
islands. Except for their core Wilson factors
it contains no Wilson weight. The weight of
the label is inside the exponential on which
the operator acts. In particular it cannot
be discarded in an estimate of the operator
on one.

The local sums in
\eqref{eq:xrg-effective-density} are
\begin{equation}\label{eq:xrg-effective-sums}
 \begin{split}
 \mathbf E_k(U)
  &=\sum_{j\le k}\sum_{(z,X)\in\mathcal Z_j^{\mathrm{pt}}(\mathcal S)}
       [E^{(j)}(X,U,z)-E^{(j)}(X,1,z)]
                           -\sum_{j\le k}b_jA(\varphi_j,U),\\
 \mathbf R_k(U)
 &=\sum_{j\le k}\sum_{X\in\mathcal R_j(\mathcal S)}
                       [R^{(j)}(X,U)-R^{(j)}(X,1)],\\
 \mathbf B_k(U,\underline{\mathfrak f})
 &=\sum_{j\le k}\sum_{X\text{ meeting the level-}j\text{ boundary}}
                       B^{(j)}(X,U,\underline{\mathfrak f}).
 \end{split}
\end{equation}
Here $\mathcal Z_j^{\mathrm{pt}}(\mathcal S)$ is the
retained pointed range, with base point in
$\Lambda_j$ and domain in its prescribed
enlargement. The remainder range
$\mathcal R_j(\mathcal S)$ stays one
level-$j$ cell inside $\Lambda_j$.
The precise enlargements and boundary
origins are those fixed by the localisation
rule; the ranges are not reselected from
the value of the density. The boundary
sum includes its distinct creation
origins, whose constants are controlled
separately below.

The same functions $E^{(j)}$ and $R^{(j)}$
serve all labels. The arguments, retained
ranges and boundary functions can depend
on the label, as can
$\mathcal E_{k,\mathcal S}$.
For the never-large label the boundary
sum is absent and the expression is
\[
 \chi_k\exp[-x_kA(U_k)
                   +\mathbf E_k(U_k)+\mathbf R_k(U_k)
                   -\mathcal E_{k,\varnothing}].
\]
This is the term retained for the floor.
It also includes the effects of histories
whose large regions were all healed and
summed into remainder terms.

The net Wilson coefficient is preserved
pointwise. Indeed,
\[
 -x_k(\cdot)-\sum_{j\le k}b_j\varphi_j=-\beta.
\]
The regular action coefficient and its
counterterm have been separated to make
the next small-field step possible.
They do not constitute two independent
changes of the bare action. The same
net-coefficient identity must hold
through partial preparations, with
their separate Wilson pieces.

The small-field integration stores $V_j$ outside $\Omega_{j+1}$.
Inside $\Omega_{j+1}$ it integrates the scaled fluctuation
$A=g_j^{-1}\mathcal A_j$ on $\Lambda_{j+1}$, conditionally on the
stored part $A'$. Completing the square leaves the form
\begin{equation}\label{eq:stored-schur}
 \Delta^{(j)}-\Delta^{(j)}
 C^{(j)}(\Lambda_{j+1})\Delta^{(j)}
\end{equation}
on $A'$. The Wilson action is not included in this conditional
operation; it remains in the function on which the operation acts.
This distinction locates the action that will pay for a failed
plaquette test.

Factorisation over disconnected retained components requires a
local choice of the form kept on their stored fields.
The tests are first recentered at localized backgrounds
by a change of variables; its Jacobian is assigned to
a boundary interaction. The pivot bonds are then solved
from their individual averaging constraints. The
conditional form is cut only after these operations.

\emph{Chains and their assignment.}
Write \(\mathsf A=\mathsf C^{\mathsf T}\Delta^{(j)}\mathsf C\)
on the free variables, and split them into integrated
and stored sets \(\mathrm I,\mathrm E\). For a stored
field \(a\), completing the square gives
\[
 \begin{split}
 \langle B,\mathsf AB\rangle
 ={}&\langle B_{\mathrm I}
       +\mathsf A_{\mathrm{II}}^{-1}\mathsf A_{\mathrm{IE}}a,\,
       \mathsf A_{\mathrm{II}}(B_{\mathrm I}
       +\mathsf A_{\mathrm{II}}^{-1}\mathsf A_{\mathrm{IE}}a)\rangle\\
 &+\langle a,\mathsf Sa\rangle,\qquad
 \mathsf S=\mathsf A_{\mathrm{EE}}
       -\mathsf A_{\mathrm{EI}}\mathsf A_{\mathrm{II}}^{-1}
                                      \mathsf A_{\mathrm{IE}} .
 \end{split}
\]
Expand \(\mathsf A\) in walks and its restricted inverse
by the small-cube parametrix. Then
\(\mathsf S=\sum_\varpi\mathsf d_\varpi\), where a chain
is either one walk term or a walk, a resolvent walk
and a walk. Its entire connected carrier belongs to
one term of the expansion.

A chain is retained when the level-\(j+1\) cells it
meets intersect the stored strip
\(\Omega_{j+1}\cap Z_{j+1}\) in a nonempty connected
set and stay within one layer of that set.
It is assigned, as a whole, to the strip component
\(D\) containing that set. The decision uses only
the carrier of the chain. Selecting all matrix
entries whose endpoints belong to one global
component would not have this property: a distant
joining could change that selection while leaving
the local field and history unchanged.

Every rejected chain has length at least that of
one large-field cell,
\(W_{j+1}=LM\check R_{j+1}\) in level-\(j\) units.
At \(|a(b)|\le Cp_0(g_j)\), the tail of the chain
expansion consequently gives
\begin{equation}\label{eq:kept-rejected-chains}
 |\mathsf B^\times(Y)|
 \le Cp_0(g_j)^2(LM)^8
       e^{-c\delta_\star W_{j+1}}e^{-\kappa'\#Y}.
\end{equation}
The rate \(\kappa'\) is prescribed before taking
\(\delta_\star LM\) large. The exponential in
\(\check R_{j+1}\) pays for the stored field squared
and any fixed logarithmic enlargement of its range.
With \(\tau_j=x_j/\theta_j\), the exact decomposition
of the remaining conditional exponent is
\[
 -\tfrac12\tau_j\langle a,\mathsf Sa\rangle
 =-\tfrac12\sum_D\langle a,\Delta_D^{\rm loc}a\rangle
                         +\sum_Y(\mathsf B^\times+\mathsf B^\tau)(Y),
 \qquad
 \sum_Y\mathsf B^\tau(Y)
       =-\tfrac12(\tau_j-1)
                          \sum_D\langle a,\Delta_D^{\rm loc}a\rangle .
\]
The companion term has only
\[
 |\mathsf B^\tau(Y)|
       \le C(M)|1-\tau_j|\,\mathfrak b^2e^{-\kappa'\#Y}
\]
when \(\mathfrak b\) bounds the stored field.
It has no long-chain factor; its smallness comes
from the frozen run's weight tube. In the unsourced
run \(\theta_j=x_j\), it vanishes.

\emph{The local coercivity comparison.}
The required bound is
\begin{equation}\label{eq:kept-coercivity}
 \langle a,\Delta_D^{\rm loc}a\rangle
                  \ge c\gamma_0\sum_{b\in D}|a(b)|^2.
\end{equation}
For the full form \(\mathsf A\ge\gamma_0\), its
genuine Schur complement satisfies this lower bound:
minimize \(\langle(B_{\mathrm I},a),
\mathsf A(B_{\mathrm I},a)\rangle\) over \(B_{\mathrm I}\).
This argument does not apply to
\(\Delta_D^{\rm loc}\), a partial chain sum.
Nor can it require a global regular extension
of the pair if the stored factor is defined only
near \(D\).

The local construction uses a quadratic partition
on small cubes, a regular gauge and cutoff extension
on each cube, and comparison with a genuine Schur
complement inside \(D^\sim\). The algebra of the
partition is explicit. For a symmetric kernel \(K\)
and real functions with \(\sum_\alpha\chi_\alpha^2=1\),
\[
 \sum_\alpha\langle\chi_\alpha a,K\chi_\alpha a\rangle
        -\langle a,Ka\rangle
 =-\tfrac12\sum_{b,b'}\langle a(b),K(b,b')a(b')\rangle
                  \sum_\alpha(\chi_\alpha(b)-\chi_\alpha(b'))^2 .
\]
Thus its error is controlled by the weighted kernel
row sums and the small-cube cutoff differences.
On each localized cube the real field is gauged
near one and cut off, so the full torus form has
its already established gap.
The Schur comparison is made with the weight
\(e^{c\,\mathrm{dist}(\cdot,D)}\), suppressing
the influence of the remote artificial boundary.
At a globally regular real pair the proposed
comparison reads
\[
 \Delta_D^{\rm loc}\ge
               \gamma-Ce^{-c\delta_\star W_{j+1}}
\]
if the full form has gap \(\gamma\).
For a pair given only on \(D^\sim\), the same argument
must be checked after the local gauge and cutoff,
with a fixed fraction of the gap left after
localization. This uniform comparison, including
the partition error, is the analytic verification
needed for \eqref{eq:kept-coercivity}. The small
scalar tail bound \eqref{eq:kept-rejected-chains}
at bounded stored fields is a different estimate;
it is not by itself an operator lower bound.

\emph{Carrying a form to later backgrounds.}
The retained form and the elimination factor are
formed once as functions of the pair near their
carrier. At a later level they are evaluated at
the background reached then. Their tests, gauge
factors and delta constraints do not move with
that background. A move therefore creates the
potential
\[
 -\tfrac12\langle a,
       [\Delta_D^{\rm loc}(\mathbf p^\sharp)
                         -\Delta_D^{\rm loc}(\mathbf p^\flat)]a\rangle ,
\]
not a newly chosen truncation rule. For a piece
with carrier \(R\), its pair oscillation is at
most
\[
 C(M)p_0(g_{l-1})^2e^{-\kappa d_l^\natural(R)}.
\]
On the enlarged far-field discs and near-field
switches this has the additional factor
\((\lambda^{\rm st})^2\), since the piece is
quadratic in those variables.

Localize the difference with the same cube parameters
as the response. Cauchy gives a factor per added cube;
Schwarz gives the relative response amplitude; the
old factor already has decay across its own carrier.
At depth \(p\), the number of old cubes per current
cube is \(L^{4p}\), and
\[
 L^{4p}e^{-cL(p-1)_+}\le CL^4.
\]
The response therefore pays this multiplicity.
After also charging the switches, the resulting
stage piece has the form
\[
 |V(Y)|\le
 C(\log\theta_j)^{-(q-3\mathfrak p_0-C\mathfrak r)}
                                      e^{-c d_k^\natural(Y)}.
\]
One power \(\mathfrak p_0\) comes from the
physical fluctuation, two from the stored form.
This is why the radius exponent must satisfy
\(q>3\mathfrak p_0+C\mathfrak r\).
The smallness of this carried difference is
compatible with a large negative stored-only
exponent; no absolute bound on that exponent
is required.

The state carries two region sequences. One determines
the averages and the current background. The other
determines the already produced action terms, their
profiles, tests and vacuum constants. They agree at
ordinary step boundaries. During a partial preparation
they need not agree, and both belong to the recorded
history. A completed integration changes the determining
sequence without retroactively changing where an old
term was created or where its counterterm was assigned.

More precisely, a history trace on a window records the
intersection with each region at that region's own
scale. It also records the unions of the marked
components, their stored regions and nests, separately
for each mark level, kind and completed stage number.
Three stored-bond sets are recorded separately: collar,
rim and cutoff-zone bonds. The regions alone do not
determine these sets. This trace, together with the
local terms and coefficients entering the window,
determines the new term formed there. Transport between
two settings requires a local translation by their cell
lattice, a common nonwrapping lift, and equality of this
entire trace.

At a real pair every stored operation is positive.
Reality of the individual analytic activities needs its
own induction. The involution is
$\mathbf U\mapsto(\mathbf U^*)^{-1}$,
$\mathbf J\mapsto\mathbf J^*$, together with complex
conjugation of scalar parameters. The real local choices,
series, integrals, unique critical solutions and the
logarithm branches based at positive real integrals
commute with it. Starting from real $\beta$ gives real
activities at real pairs. This does not identify a
complex Gaussian with a positive measure.

The averaging constraints also leave
elimination factors when their pivots
are solved. If $\mathsf m_c(W)$ is
the derivative of the average at the
pivot of a coarse bond, the real
input required by this chart is
\begin{equation}\label{eq:xrg-pivot-determinant}
 \det(L\mathsf m_c(W))
                       \ge\frac4{\pi^2}L^{-15}>0
\end{equation}
where no contour loop of that
average is $-1$. Its elimination
factor is
$|\det\mathsf m_c(W)|^{-1}$.
The bound must apply on real
large-field configurations in
this chart as well as at the
regular background. It is a
property of the actual symmetric
average, included in the pivot
construction; it is not inferred
from the sign of a Gaussian
quadratic form. Outside such
charts the total Borel rules
and complementary tests are used.

On a regular chart the unique
fixed points, means and gauges
are limits of their continuous
iterations and hence are Borel.
The local tests, deterministic
maps and normalised reference
draws therefore give Borel
kernels there. Their specified
total continuations and the
sets on which the regular
objects are declared to exist
must be Borel as well. These
must be checked before assigning
the history law or integrating
its witnesses. Local analyticity on
a regular chart does not settle
this total Borel continuation.

Positivity of the spatial Wilson weights is separate again. The
profiles $\varphi_j$ are fixed by the action sequence and defined in
\cref{sec:sources}. Away from the final prepared layers,
the telescoping array is either $\beta$ or a convex combination
of two successive $x_j$. In the final thin shells its
additional vertex is $x_k+b_{k-1}$; the one-step lower
flow bound makes this vertex positive. The full profile minimum
and its proof are given in \eqref{eq:revised-profile-minimum}. These
facts give positive reference denominators on the
supports of the real construction.

\subsection{Interrupted operations and superposition of histories}

An interrupted operation retains more than
its current regions. Suppose an operation
at level $t$ stops after $n(X)$ stages,
and write $t_\flat=t-N_t$.
Its mark stores the regions
$\Omega_l^\circ,\Lambda_l^\circ,S_l^\circ$
that were present before it began, as
well as the prepared nest $Z_l''(X)$.
For $0\le\nu(X)\le n(X)$ define the
intermediate determining regions on
$X$ by
\begin{equation}\label{eq:xrg-interrupted-sequences}
 \Omega_l^\nu\cap X=
 \begin{cases}
 X\setminus Z_l''(X),&t_\flat<l\le t_\flat+\nu(X),\\
 \Omega_l^\circ\cap X,&t_\flat+\nu(X)<l\le t_\flat+n(X).
 \end{cases}
\end{equation}
Other entries are unchanged.
The choice $\nu=n$ describes the
fields that remain and determines
the background and its propagators.
The choice $\nu=0$ describes where
the old action terms, tests,
profiles and constants were made.
The hereditary domains include
the intermediate choices as well.

A mark at level $s$ inside a class
component at level $t$ has
$s<t-N_t$. Its window is disjoint
from the present waiting window.
Marks that overlap in space are
therefore nested with disjoint
windows, and at one location and
level only one interrupted mark
is active. The intermediate choice
in \eqref{eq:xrg-interrupted-sequences}
is realised by the history in
which that operation stops at
$\nu$ stages. The later geometry
is unchanged: creations, component
classes and nests do not read how
far those interrupted stages had
progressed. Thus the intermediate
domains are tied to actual
histories, not to fictitious
interpolations of regions.

The local superposition rule has
two ingredients. First, the
margin around a non-touching
component exceeds the reach of
every bond test, pivot and
collar. Taking a hull and adding
layers commute with unions;
boxing is done one component at
a time. An ordinary step
therefore combines two histories
whose regions never touch by
acting separately on their
components.

Second, the operation class
looks back through a waiting
window. Nothing created within
that window can already have
been healed. Indeed, suppose
something created at step
$s\ge t-N_t$ were healed at
$i<t$. Healing would require
$s\le i-N_i$, or the stricter
$s<i-N_i$ for a mark. Since
$N_i\ge N_t$,
\[
 i-N_i\le i-N_t<t-N_t\le s,
\]
which is a contradiction.
The creations and marks needed
to decide membership still
lie in their respective live
regions. This window argument
is necessary because the
overall evolution is not
monotone: a healed place can
become large again later.

Consequently each history is
a union of clusters. A cluster
is rooted either at a live
component of the last level
or at a component at the level
where it is healed. It contains
all its predecessor components
and everything recorded in
them. Every component has one
last successor and hence one
cluster root. Non-touching
clusters can be interchanged
between histories. In
particular, deleting the
entire clusters of healed
components leaves a history
with the same final label.
This is its standing part
$h^\flat$. Deleting only the
last healed box and retaining
its old tests would not give
this history.

Fix standing data $\bar h$ at
level $k$. The histories that
reduce to it after the current
healed clusters are erased are
parametrised by a family
$\mathfrak X$ of pairwise
non-touching possible healed
components, disjoint from the
standing live region, and by
one internal history for each
$X\in\mathfrak X$ in which
$X$ is healed and none of its
predecessors was healed earlier.
The local choices of these
internal histories are
independent. Their stored
integrals remain coupled through
the exterior argument, which
is why the integrated supremum
majorant is needed later.
This parametrisation supplies
the hard-core component sum
used by exponentiation.

The trace on a window must
remember the action sequence,
interrupted nests, stage counts
and the separate collar, rim
and cutoff-zone bond sets.
For complete equality of the
natural analytic domains,
two further geometric facts
must match under a chart:
the sets of bare bonds
carrying the fields, and
whether two bare unit cubes
belong to one drop-out mark.
A nonconvex domain can cut
a mark into two visible
pieces, so its visible union
does not determine that
second relation. On unmarked
domains inside the first
region these extra clauses
are vacuous.

If a current window lies a
cell inside the small-field
region, its standing trace
is empty. The regular term
formed there is therefore
the term of the never-large
standing data, already including
the summed healed clusters.
This is the local identity
that makes the interior
families independent of
labels and, under a common
nonwrapping lift, of the
torus and cutoff. The
all-tori companion and the
coefficient sequence use
this identity, not only a
bound on a label-dependent
term.

\subsection{Removing a small old component}

Without a further operation, a retained region would remain in every
later label. Even very small probabilities would then be summed over
the cells of all finer scales. The large-field operation acts on a
component that has become a small box and has had a sufficiently long
interval with no new failure, no joining to another component, and
no intervening operation. Let this interval have length
$N_k=\check R_k$, ending at level $k$. These restrictions ensure that
the component has a single nested chain of boxes throughout the
interval.

The stored fields are first integrated in successive shells. This
reduces a box of side $M\check R_k$ to a prepared box of side $O(M)$,
with multiscale shell data on its boundary. The preparation uses the
sharper fluctuation threshold
\[
 p_1(g)=(\log g^{-2})^{\mathfrak p_1},\qquad
 \delta'_j=g_jA_1p_1(g_j),
 \qquad \mathfrak p_1<\mathfrak p_0.
\]
One needs both the accumulated implications between successive tests
and enough action to pay for a failed sharper test. The exponent
conditions have the form
\begin{equation}\label{eq:threshold-exponents}
 \mathfrak p_1<\mathfrak p_0-C\mathfrak r,
 \qquad 2\mathfrak p_1-C\mathfrak r>\mathfrak p_0.
\end{equation}
The exponents are chosen to leave these margins before the coupling
ceiling is fixed.

For a component that passes the preparation tests, integrate the block
field in the prepared box and replace it by a normalised local
reference law. This law contains a tree gauge, a cutoff, and the local
Wilson weight with coupling $x_k$. Its denominator is the integral
of exactly that numerator over the box. Multiplying by the previous
mass preserves the integral. A finite average over cubic frames keeps
the operation covariant; the numerator, denominator, trees and tests
must all use the same frame in each summand.

There are three outcomes. If all tests pass, the component is healed
and removed from the label. If only the final small-field test fails,
the operation is kept as an island with no inner regions. If a
preparation test fails, the component remains with the integrations
already performed; this is a drop-out. Islands and drop-outs are
the marked components. Each records the failed test and its extracted
gain. A component still retained after these operations is called
live. Write $\mathcal Z_k$ for $Z_k$ with the components healed
at step $k$ removed; it is the input region of step $k+1$.
The standing label consists of the unhealed regions and their
component marks, after histories differing only in healed regions
have been summed. Islands and drop-outs retain their local measures
and factors. The regions used to form the current background and
the regions used to index the already constructed action terms must
then be recorded separately. A partial shell integration changes the
former without authorising a new sorting of the latter.

The reference-law replacement preserves total mass and observables
independent of the replaced variables, while changing the coarse-field
distribution inside the box. Observables involving those variables are
transported by the kernels in \eqref{eq:history-kernels} and the exact
integral identities.

We describe the membership and the main integral more precisely.
There are five requirements. The depth condition is $k\ge N_k$
and availability of the compatible intermediate cube scales.
The size condition puts $X$ in a cube of side $CM\check R_k$.
Quiet means that no creation or joining occurred during
$k-N_k<t\le k$. Quarantine excludes marks placed in the
preceding waiting interval. Finally, each $X\cap Z_j$, for
$k-N_k\le j\le k$, must be one component and one box; these
boxes form its nest. Before the depth condition is met the
class is empty. The free evolution of a quiet component
contracts its linear size over two levels. It acquires the
geometric conditions after a contraction period and the
waiting interval; a new event restarts this wait.

Let $\phi_X$ be the construction's cutoff supported in $X$ and
equal to one well inside it. Let $\Lambda^{\mathrm{pr}}$ be the prepared box and
$V_\Lambda$ the critical orbit of its local action, with
the fixed exterior data, and let $U_{k,X}$ be the background localized
in $X$. Put $V'_k=V_kV_\Lambda^{-1}$.
In a fixed cubic frame the reference law is
\begin{equation}\label{eq:mrg-reference-law}
 \mathfrak s_X(V_k)=
 \frac{\delta_{\mathrm{ax}}(V'_k)\chi(V'_k)
       e^{-x_kA(\phi_X,U_{k,X}(V'_kV_\Lambda))}}
 {\displaystyle\int dV'|_{\Lambda^{\mathrm{pr}}}
       \delta_{\mathrm{ax}}(V')\chi(V')
       e^{-x_kA(\phi_X,U_{k,X}(V'V_\Lambda))}}.
\end{equation}
The cutoff restricts the free tree-gauge variables to
$|\frac1i\log V'(b)|<M_0\varepsilon_k$.
Both occurrences of $\phi_X$, both local backgrounds and
all tree choices in this quotient are identical.
If its denominator is not finite and positive, the
construction declares a complementary outcome. At real
data on the specified supports it is finite and positive.
The analytic estimates must also keep the denominator
nonzero on the required complex domain. This latter
assertion does not follow merely from real positivity.

The cubic frame average is inserted after integrating
the old box and before imposing the new tree gauge.
Subsequent choices are made in that frame. A failure
after this point retains its frame. For an earlier
drop-out, the identity
$1=|W_4|^{-1}\sum_{r\in W_4}1$ may be used in its mass
bound, but no tree draw has taken place. This placement
keeps the reference replacement linear in the input
measure and preserves every outcome mass.

The preparation integrates shells in the order
$j=k-N_k,\ldots,k-1$. Every shell receives the coupling
increments missed while it was stored. The remaining
multiscale datum consists of thin shells around the old
core, and the main integral uses one scaling $g_k$ on
all of them. The old core operation is not unfolded;
it retains the earlier failed tests and their paying
action. The least-action continuation is taken only
in the prepared box of side $O(M)$, among its regular
critical orbits. Its constant loss depends on $M$ and
is independent of $\check R_k$.

\subsection{Preparing the shells and restoring their couplings}

Put $k_\flat=k-N_k$. On a component in the
operation class the quiet geometry has an
explicit form. If $H_i$ is the hull in
level-$i$ cells of $X\cap Z_{i-1}$ and
$H_i^{[r]}$ is the box with each face
expanded by $r$ such cells, then
\begin{equation}\label{eq:xrg-quiet-boxes}
 \Omega_i^c\cap X=H_i^{[8]},\qquad
 Z_i\cap X=H_i^{[10]},
                       \qquad k_\flat<i\le k.
\end{equation}
The boxes at a level are concentric. The
single-box condition excludes a joining
within this window even if the joined
region is small and no new test fails.
Without it the preceding layer could
have two components and no nested shell
tree of this form.

The prepared box is the final nest box
with one layer of $M$-cubes added. Its
diameter is $CM$ in level-$k$ units.
The old core has diameter at most
\[
 CM\check R_{k_\flat}L^{-N_k}
\]
in those units. Shrinking the core is a
consequence of waiting; the stages
themselves thin the outer shells.
The admissible torus condition puts
every box and all its needed
neighbourhoods in a nonwrapping region.

The gauge tree on the prepared datum
uses bonds of all the shell lattices.
For a pair of boxes $P_1\supset P_2$,
take axial paths issuing from one site
and going around $P_2$. The family is
stable under enlarging $P_2$: delete
the paths of the removed sites and
obtain the chosen tree on the smaller
shell. This stability is why the
geometric class uses boxes. A shell
tree fixes every site gauge except one.
For each inner box add one bond crossing
its boundary, and add one at the
outermost box. They join the shell
trees into one tree $T_0$, rooted at
the exterior endpoint of its entering
bond. Haar invariance gives unit
Jacobian for this tree gauge.

The bare averaging chain is resolved
from the highest level down:
\[
 \int dV_i(c)\,
  \delta(\mathcal M(V_{i-1})(c)V_i(c)^{-1})=1.
\]
It leaves a free integral over the
prepared multiscale datum and the
unopened old core. Its schematic
measure, with every shell at its own
scale, is
\begin{equation}\label{eq:xrg-prepared-integral}
 \chi_k\chi_{k,\Lambda}
 \int dV_{k_\flat}|_{C_{k_\flat}}\,
      \chi\,\mathbf T_{k_\flat}(C_{k_\flat})
 \int dV''|_{\mathbf B''\setminus C_{k_\flat}}
                      \chi''_k e^{A''_k(U''_k)}.
\end{equation}
Here $C_{k_\flat}=X\cap Z_{k_\flat}$,
$\chi_{k,\Lambda}$ is the regular-extension
test, and the primes denote the data,
tests and action after the stages.
The field expanded in the main integral
is the free datum $V''$, not the stored
bare core field.
Its background datum is
\[
 V_0=(\mathcal M^jU_0
                  \text{ on the shell of level }j)_{k_\flat\le j\le k},
 \qquad V''=e^{ig_kB}V_0,
 \qquad |B|<\delta'_k/g_k.
\]
One scale $g_k$ and one test control
all these components of $B$.

The extension test also has a paying
alternative. For exterior curvature
bounded by $\varepsilon$, choose axial
gauge on the boundary of the prepared
box and continue by one inside. This
gives a continuation with curvature
$CM^2\varepsilon$. Thus absence of
an extension in the required radius
forces an exterior plaquette of size
at least $cM^{-2}\varepsilon_k$.
The failed test is controlled by the
same local action payment as a
plaquette failure. It must not be
declared impossible from exterior
smallness at an unspecified radius.

\subsubsection*{The critical extension on the prepared box.}

The stages have reduced the free box to side \(O(M)\).
This permits a fixed-radius argument, rather than an
inverse estimate growing with \(\check R_k\).
Let \(P=\Lambda^{\rm pr}\cap\mathbb T^{(k)}\), let
\(E(P)\) be the bonds having an endpoint in \(P\),
and let \(\mathsf G\) be its tree rooted at one
exterior site. Write \(\mathcal X\) for the real
Lie-algebra fields on \(E(P)\) that vanish on
\(\mathsf G\). All constants in this box argument
may depend on \(M\), fixed before the coupling
ceiling.

Start with a real regular extension \(V^\sharp\).
An exterior increment \(\mathsf b\) vanishes on
\(E(P)\); a free increment \(\delta\) is supported
there. If \(\mathsf W_{V^\sharp}\) denotes the
background continued from this datum, set
\[
 F(\mathsf b,\delta)=
 \sum_{p\subset\Omega_0}
 \left[a_p(\mathsf W_{V^\sharp}(\mathsf b+\delta))
                  -a_p(\mathsf W_{V^\sharp}(\mathsf b))\right].
\]
The continued constrained-background estimate gives
holomorphy on a fixed polydisc and
\begin{equation}\label{eq:main-box-action-difference}
 |F(\mathsf b,\delta)|
 \le K_M\bigl(\varepsilon+
             \max\{|\mathsf b|_\infty,|\delta|_\infty\}\bigr)
                                                |\delta|_\infty .
\end{equation}
For real arguments this is the difference of the
regular critical-background actions at the two data. Its derivative
in \(\delta\) is therefore the variation of that
regular critical-background action over the box datum.

\emph{Transverse Hessian.}
Let \(\mathsf H_1\delta\) be the linearized
background response to a free datum variation.
The coercivity of the fine action, the linear
averaging estimate and the tree Poincar\'e
inequality give the chain
\[
 \text{fine Hessian}\ \gtrsim\
       \|\mathrm d\mathsf H_1\delta\|^2
 \ \gtrsim\ \|dQ_k^{\rm st}(1)\mathsf H_1\delta\|^2
 \ \simeq\ \|d\delta\|^2
 \ \ge c_T|\delta|_2^2 .
\]
The averaging estimate in the middle has constant
one at the trivial field in four dimensions:
averaging the plaquette squares by Cauchy--Schwarz
and then summing fine plaquettes gives
\[
 \sum_{p'\in S}|dQ_k^{\rm st}(1)A(p')|^2
 \le\sum_{p\in S_0}\eta_k^4|dA(p)|^2.
\]
The factors from length, averaging and plaquette
count cancel. This is a free-action estimate; no
mass term is added at the bare scale.

For a regular centre the response reproduces the
coarse variation only up to the coarse node gauge.
To justify the middle comparison, differentiate
the gauge-invariant coarse Wilson action in two
ways: directly at its coarse datum, and through
the fine response and its average. The gauge
factor drops out of that action. The remaining
perturbation terms have size
\(C_M\varepsilon|\delta|_2^2\), with the box bond
count included in \(C_M\). Taking the regularity
radius small compared with \(c_T/C_M\) gives
\begin{equation}\label{eq:main-box-hessian}
 \partial_\delta^2F(0,0)[\delta,\delta]
                         \ge\tfrac12c_T|\delta|_2^2,
                    \qquad\delta\in\mathcal X .
\end{equation}
The tree condition is essential: gradients would
otherwise remain in the kernel. Its Poincar\'e
constant is fixed because the free box has
\(O(M^4)\) bonds.

\emph{Solving the box equation.}
Let \(\mathsf L=\partial_\delta^2F(0,0)|_{\mathcal X}\).
It has a bounded inverse by
\eqref{eq:main-box-hessian}. Decompose the derivative as
\[
 \partial_\delta F(\mathsf b,\delta)
       =g_{\mathsf b}+\mathsf L\delta
                                   +\mathcal N_{\mathsf b}(\delta),
 \qquad \mathcal N_{\mathsf b}(0)=0.
\]
On a smaller fixed polydisc, Cauchy estimates from
\eqref{eq:main-box-action-difference} give
\[
 |g_{\mathsf b}|\le C_M(\varepsilon+|\mathsf b|_\infty),
 \qquad
 \|\partial_\delta\mathcal N_{\mathsf b}(\delta)\|
       \le C_M(\varepsilon+|\mathsf b|_\infty+|\delta|_\infty).
\]
Choose first a transverse ball of radius \(r_P\),
then an exterior radius \(a_P\) and the regularity
ceiling so that
\(\|\mathsf L^{-1}\partial_\delta\mathcal N_{\mathsf b}\|
\le1/2\) and
\(\|\mathsf L^{-1}g_{\mathsf b}\|\le r_P/2\).
The map
\[
 \delta\longmapsto-\mathsf L^{-1}
                 [g_{\mathsf b}+\mathcal N_{\mathsf b}(\delta)]
\]
is a contraction of that ball. Its fixed point
\(B_\Lambda(\mathsf b)\) is holomorphic in the
exterior increment, real at real increments,
and obeys
\[
 |B_\Lambda(\mathsf b)|_\infty
                    \le K_P(\varepsilon+|\mathsf b|_\infty).
\]
If the centre is already critical then
\(B_\Lambda(0)=0\), and the bound improves to
\(2K_P|\mathsf b|_\infty\). Uniform convergence of
the contraction iterates proves this holomorphic
dependence on the same fixed disc.

Stationarity has so far been imposed on the
transverse tree variables. The action is invariant
under gauges at sites of \(P\). The remaining
directions are tangent to those gauge orbits,
so its derivative vanishes in all directions
supported in \(E(P)\). In particular the basic
background \(U_0\) has zero current in the free
box; this is the exact criticality used later
to cancel the linear Taylor term.

The uniqueness statement concerns regular
continuations, not arbitrary box data. In the
tree gauge a continuation whose plaquettes in
the prescribed neighbourhood of \(P\) have size
less than \(a_P\) lies in the preceding fixed
chart, after reducing \(a_P\) using the box
Poincar\'e constant. The Hessian stays positive
there by the same perturbation estimate.
Thus its real critical points form one gauge
orbit and attain the minimum on that regular
class. Its tree representative is \(V_\Lambda\),
with
\[
 |\partial V_\Lambda-1|\le C(M)\varepsilon .
\]
The argument requires the constrained-background
holomorphy and response bounds uniformly over
the stated exterior class. It is not a claim
about least action over unrestricted rough
continuations. The preparation to a fixed
box is what makes all constants above independent
of the number of shells and the waiting depth.

The main-step tree Poincar\'e bound
for a single bond is weaker:
\[
 |B(b)|^2\le C(M\check R_k)^C
                         \sum_p|\partial B(p)|^2.
\]
One power loss comes from the
diameter of a shell in its own
lattice, at most $CML\check R_k$,
and another from the number
$N_k+1$ of shells. This is the
loss used when paying for a failed
sharper test. It is distinct from
the fixed-$M$ bound used to
construct $V_\Lambda$.

To check the implications of
successive stage tests, let
$U^{(n)}_{k,X}$ be the local
background after $n$ stages and
put $j=k_\flat+n$. On a shell
of level $i\le j$, the comparison
required of the response is
\begin{equation}\label{eq:xrg-stage-curvature}
 \begin{split}
 |\partial U^{(n)}_{k,X}-1|
 &\le(1+c\beta_{\mathrm s}2^{-(j-i)})
             |\partial U^{(n+1)}_{k,X}-1|\\
 &\quad+c\beta_{\mathrm s}2^{-(j-i)}
                                     \varepsilon_i\eta_i^2.
 \end{split}
\end{equation}
Indeed, its error has the form
\[
 C\frac{\delta'_j}{\varepsilon_j}
     (1+\sqrt{j-i})e^{-c(M/M_1)(j-i)}
                                  \varepsilon_i\eta_i^2.
\]
The fixed geometric separation makes
the exponential dominate the square
root. The small threshold ratio
pays the remaining constant, leaving
$\beta_{\mathrm s}2^{-(j-i)}$.
Here $\beta_{\mathrm s}>0$ is a
fixed fraction reserved for this
operation.

The stage-$n$ shell threshold is
therefore chosen as
\[
 [1-\beta_{\mathrm s}(1-2^{-(j-i+1)})]
                                  \varepsilon_i\eta_i^2.
\]
Summing the lost margins gives
a fixed fraction rather than
$N_k$ times that fraction.
For outer shells whose collars
overlap, choose an integer
$2\le L_\flat<L/2$ and relax
the shell-$i$ bound by
$L_\flat^{2(i-k_\sharp-1)_+}$.
Here $k_\sharp+1$ is the last
level whose cell has side $M$
in level-$k$ units.
This factor depends on the shell,
not on the stage. A passage to
the next finer scale gains
$L^{-2}$ and loses only
$L_\flat^2$, so the relaxed
bound still implies the earlier
one. Prepared nests with gaps
of only one $M$-cube require
this calculation; an estimate
using $e^{-cM\check R_k}$ at
every such gap would be false.

The old fluctuation tests are
implied by
$CL\delta'_j\le\delta_j/2$.
At the main integral the sharper
threshold must also be small at
the oldest core level. The
combined requirement is
\begin{equation}\label{eq:xrg-stage-threshold-room}
 C\sqrt{N_k}\frac{p_1(g_{k_\flat})}{p_0(g_{k_\flat})}
       +C(M)L^{-N_k}\sqrt{N_k}\le c.
\end{equation}
The first term gives
$\mathfrak p_1<\mathfrak p_0-C\mathfrak r$;
the second is made small by
the wait $N_k=\check R_k$.
These implications hold on the
common support of the carried
tests and averaging constraints,
and only one current cell
inside $X$. They do not extend
to $\partial X$, where the
local background is pinned
bond by bond to the exterior.

Each completed stage restores
the missed increments on the
shell it integrates. Let
$\varphi[\Lambda]$ denote the
profile of a region and set
\begin{equation}\label{eq:xrg-stage-weights}
 \begin{gathered}
 x^{(m)}=\beta-\sum_{i\le k}b_i\varphi[\Lambda_i^{(m)}],\\
 \Lambda_i^{(m)}=
 \begin{cases}
 \Lambda_i\cup(X\setminus Z_i''),&k_\flat<i\le k_\flat+m,\\
 \Lambda_i,&\text{otherwise}.
 \end{cases}
 \end{gathered}
\end{equation}
The Wilson difference between
$x^{(m)}$ and $x^{(m-1)}$ is
kept separately, before
localisation. At each bare
plaquette the total of the
Wilson weights remains
$-\beta$. After all stages,
$x_k''=x^{(N_k)}$ equals $x_k$
outside the prepared box and
equals $x_k+O(1+k-j)$ on its
shell of level $j$.
This restoration supplies the
net weights in the following
cancellation.

\subsubsection*{The backgrounds used in the main cancellation.}

The cancellation compares several backgrounds, each solving its
own constrained problem. We describe their data before estimating
the resulting differences. Write \(Z=X\), \(\Lambda=\Lambda^{\rm pr}\),
\(N=N_k\), and \(\xi_j=L^{j-k}\) in level-\(k\) units.
The prepared determining set \(\mathbf B''\) has level \(k\)
outside its final box, descends through the thin shells, and
retains all the old levels in the innermost core \(I_2\).
The healed determining set \(\bar{\mathbf B}\) instead has
level \(k\) throughout \(Z\).

There are four local determining sets besides these two.
The towers \(\mathbf B^Z,\mathbf B^\Lambda\) are those of
\(Z,\Lambda\). Join \(\mathbf B''\) to \(\mathbf B^Z\) in the
two layers next to \(\partial Z\), obtaining \(\mathbf B^{\rm f}\).
Join it to the tower of the innermost prepared box, obtaining
\(\mathbf B^{\rm in}\). A joined determining set agrees with
the first set in the interior and with the tower in its joining
layers. In particular \(\mathbf B'',\mathbf B^{\rm f}\) agree
well inside \(Z\), and
\(\mathbf B'',\mathbf B^{\rm f},\mathbf B^{\rm in}\) agree
well inside the innermost box. These equalities of data and
determining sets are what permit local background comparisons.

Let \(V^{\rm e}\) be the exterior datum and let \(V_\Lambda\)
be its critical extension, in the tree gauge already chosen.
Its basic background \(U_0\) has zero current on the free
box. In a mixed axial and Landau gauge it has a potential
\(\mathbb A_0\) on the bulk of \(\Lambda\), with
\[
 |\mathbb A_0|+|\nabla\mathbb A_0|
              +|\mathrm d^*\mathrm d\mathbb A_0|\le C(M)\varepsilon.
\]
For an exterior increment \(\mathsf b\), the same estimate
holds with \(\varepsilon+|\mathsf b|_\infty\).
The averages of this representative on a shell of level
\(j<k\) are \(e^{i\mathsf v_0(\mathsf b)}\), where
\begin{equation}\label{eq:main-basic-shell-increment}
 |\mathsf v_0(\mathsf b)|\le
          C\xi_j(\varepsilon+|\mathsf b|_\infty),\qquad
 \left|\tfrac1i\log
       (e^{i\mathsf v_0(\mathsf b)}e^{-i\mathsf v_0(0)})\right|
          \le C\xi_j|\mathsf b|_\infty .
\end{equation}
The factor \(\xi_j\) is a potential scaling, rather than a
curvature scaling. The formula applies to every average whose
corner blocks lie in the bulk, including levels below \(k_\flat\).

Glue this basic datum to the real stored datum \(V^{\rm in}\)
at the inner rows and to the lifted exterior datum at the outer
rows. Call the result \(V^0\) on \(\mathbf B^{\rm f}\).
The mixed gauge is needed here: at the inner seam its basic
averages are \(O(L^{-N}\varepsilon)\) from one. The other
side is controlled by the last-shell test on
\((1,V^{\rm in})\), with threshold \(\varepsilon_{k_\flat}/2\).
Thus both the interior plaquettes and those crossing the seam
are regular. Gauge-invariant small-field tests on the two
unglued data would not by themselves imply this fact.

Let \(U^0\) be the background of \(V^0\). Locality and
\eqref{eq:main-basic-shell-increment} give, on an inner
test cube separated from its comparison boundary by \(D_2\),
\begin{equation}\label{eq:main-last-shell-transfer}
 |\partial U^0-1|
 \le\left[\tfrac12\varepsilon_{k_\flat}
          +C\varepsilon(L^{-N}+e^{-\delta_{\mathrm G}D_2})\right]
                                      \eta_k^2\xi_{k_\flat}^{-2}.
\end{equation}
The first margin is numerical; the other two terms tend to
zero at the chosen wait. This verifies why a half-threshold
test can supply an unchanged full threshold at the glued field.

For complex arguments use Landau forms centred at these real
data. We denote the healed background by \(\bar W\), its
tower localization in \(Z\) by \(W^Z\), the basic background
by \(W^0\), the glued local background by
\(W^{\rm f}(\mathsf b,\mathsf f)\), the prepared background
by \(W''(\bar{\mathsf b},\mathsf f)\), and the inner one
by \(W^{\rm in}(\mathsf f)\).
Here \(\bar{\mathsf b}\) is the increment of the exterior
healed datum, \(\mathsf b\) its restriction to the exterior
rows of \(Z\), and \(\mathsf f=g_kB\) the single scaled
increment on every free shell row.
At real arguments these forms represent the corresponding
critical orbits. At complex arguments they are the holomorphic
continuations with these fixed real centres.

There is one inner increment that is not quite \(\mathsf f\).
The basic shell averages enter the inner joining datum; call
the resulting increment \(\gamma^\flat\). On the inner rows
read by this comparison,
\begin{equation}\label{eq:main-inner-increment}
 |\gamma^\flat-\mathsf f|
       \le CL^{-N}(\varepsilon+|\mathsf b|_\infty).
\end{equation}
This difference need not vanish at zero exterior increment:
the basic datum itself is generally not one.

\subsubsection*{The ten Wilson differences.}

Choose \(\zeta_0=\phi_X\), equal to one away from the outer
collar, and \(\zeta_1=\phi_{\mathrm{core}}\), equal to one
on the old core, with
\(\operatorname{supp}\zeta_1\subset\{\zeta_0=1\}\). Put
\(\zeta=\zeta_0(1-\zeta_1)=\zeta_0-\zeta_1\).
Thus \(1=(1-\zeta_0)+\zeta+\zeta_1\).
The support of \(1-\zeta_0\) is at weighted distance
\(cM\check R_k\) from the free box. The support of \(\zeta\)
is that far from the unknown inner core, in the local units
of its layers.

For this calculation let \(\theta\) be the inverse coupling
used by the reference law and let \(\mathsf w\) be the
restored spatial Wilson weight. Set
\(\bar{\mathsf w}=\mathsf w\) outside \(Z\) and
\(\bar{\mathsf w}=\theta\) in \(Z\).
Write \(\mathsf H_1\) for the linear response of \(U^0\)
and define the bilinear continuation
\[
 \mathsf q_0(\mathsf f)=
 \sum_p\eta_k^4\zeta(p)
       \langle \mathrm d_{U^0}\mathsf H_1\mathsf f(p),
                         \mathrm d_{U^0}\mathsf H_1\mathsf f(p)\rangle .
\]
It is nonnegative for real \(\mathsf f\). No such sign
is asserted at arbitrary complex \(\mathsf f\).

The exact remainder consists of the following differences:
\begin{align*}
 T_1={}&A(\mathsf w(1-\zeta_0),\bar W)
                       -A(\mathsf w(1-\zeta_0),W''),\\
 T_2={}&A((\bar{\mathsf w}-\mathsf w)(1-\zeta_0),\bar W),\\
 T_3={}&\theta[A(\zeta_0,\bar W)-A(\zeta_0,W^Z)],\\
 T_4={}&\theta A(\zeta_1,W^0),\\
 T_5={}&\theta[A(\zeta,W^0)-A(\zeta,W^{\rm f}(\mathsf b,0))],\\
 T_6={}&-\theta[A(\zeta,W^{\rm f}(\mathsf b,\mathsf f))
        -A(\zeta,W^{\rm f}(\mathsf b,0))
                              -\tfrac12\mathsf q_0(\mathsf f)],\\
 T_7={}&\theta[A(\zeta,W^{\rm f}(\mathsf b,\mathsf f))
                                      -A(\zeta,W'')],\\
 T_8={}&-A((\mathsf w-\theta)\zeta,W''),\\
 T_9={}&A(\mathsf w\zeta_1,W^{\rm in}(\mathsf f))
                      -A(\mathsf w\zeta_1,W^{\rm in}(\gamma^\flat)),\\
 T_{10}={}&A(\mathsf w\zeta_1,W^{\rm in}(\gamma^\flat))
                                      -A(\mathsf w\zeta_1,W'').
\end{align*}
Adding them, and using the two partitions of the Wilson
weight, gives
\begin{equation}\label{eq:main-ten-term-identity}
 \begin{split}
 &-\theta A(\zeta_0,W^Z)-A(\mathsf w,W'')
                                 +\theta A(\zeta_0,W^0)\\
 &\quad=-A(\bar{\mathsf w},\bar W)
       -\tfrac12\theta\mathsf q_0(\mathsf f)
       -A(\mathsf w\zeta_1,W^{\rm in}(\mathsf f))
                                  +\sum_{\nu=1}^{10}T_\nu .
 \end{split}
\end{equation}
No estimate has been used in this identity. In particular the
excess \((\mathsf w-\theta)\zeta\) and the whole old core
weight remain visible.

In the unsourced run \(\theta=x_k=g_k^{-2}\) and
\(\mathsf f=g_kB\). Equation
\eqref{eq:main-ten-term-identity} is the precise form of
\begin{equation}\label{eq:mrg-core-cancellation}
 \begin{split}
 &-x_kA(\phi_X,U_{k,X}(V'_kV_\Lambda))
        -A(x'',U'')+x_kA(\phi_X,U_0)\\
 &\quad=-A(x_k(\cdot),U_k)
        -A(x''\phi_{\mathrm{core}},U^{\rm in}(g_kB))
        -\tfrac12\|\zeta^{1/2}\mathrm d_{U^0}\mathsf H_1B\|^2
        +\mathrm{Err}_X .
 \end{split}
\end{equation}
The replacement of the prepared core by the real inner
background is \(T_9+T_{10}\), not an implicit removal of
its Wilson weight. Bounding the three original actions
separately would leave the reference denominator's positive
term of order \(p_0(g_k)^2(M\check R_k)^4\).

\subsubsection*{Sizes of the errors and the surviving core action.}

Let
\[
 \mathsf s=\varepsilon+\mathsf a+\delta,\qquad
 E_1=e^{-c\delta_{\mathrm G}M\check R_k},\qquad
 \mathsf n\le C\check R_k^5,
\]
where \(\mathsf a,\delta\) bound the exterior and shell
increments and \(\mathsf n\) counts prepared unit cells
outside the core, each at its own scale. If \(y\) is such
a cell in layer \(j_y\), set
\[
 \|\mathsf w\|_\sharp=
       \sup_y\frac{\sup_{p\in\Delta(y)}|\mathsf w(p)|}
                         {1+k-j_y},\quad
 \mathsf w_2=\sup_{\zeta\ne0}|\mathsf w-\theta|,\quad
 \mathsf w'_2=\sup_{Z\cap\{\zeta_0\ne1\}}|\mathsf w-\theta|.
\]
The background response estimates give
\begin{align}
 |T_1|&\le C\|\mathsf w\|_\sharp\mathsf s^2E_1,
 &|T_2|&\le C\mathsf w'_2\mathsf s^2\check R_k^4,
 &|T_3|&\le C\theta\mathsf s^2\check R_k^4E_1,
 \label{eq:main-outer-errors}\\
 |T_4|&\le C\theta\mathsf s^2\mathsf n L^{-4N},
 &|T_5|+|T_7|&\le C\theta\mathsf s^2\mathsf n E_1,
 &|T_8|&\le C\mathsf w_2\mathsf s^2\mathsf n,
 \label{eq:main-middle-errors}\\
 |T_9|&\le C\|\mathsf w\|_\sharp N\mathsf s^2\mathsf n L^{-N},
 &|T_{10}|&\le C\|\mathsf w\|_\sharp N\mathsf s^2\mathsf n E_1.
 \label{eq:main-inner-errors}
\end{align}
Here and below a fixed \(M\)-dependent constant is harmless.
For \(T_1\), before summing cells, the bound is
\[
 C\mathsf w_y\mathsf s^2
       e^{-c\delta_{\mathrm G}d_{\bar{\mathbf B}}(y,\Lambda)}.
\]
This also pays the linear growth of the old coupling in
\(k-j_y\).

The estimates express three distinct comparisons.
For \(T_1,T_3,T_5,T_7,T_{10}\), the backgrounds have the
same data and determining set near the action being
compared, and differ only across a long collar.
For \(T_4\), the basic field in the inner region is a
level-\(k\) field viewed \(N\) levels finer; its action
there has the fourth power \(L^{-4N}\).
For \(T_9\), it is the increment difference
\eqref{eq:main-inner-increment} that supplies the small
factor. In dimension four the count of bare plaquettes
in a unit cell, \(6(\xi_j/\eta_k)^4\), cancels the
action factor \(\eta_k^4\xi_j^{-4}\). Thus the sums
are over prepared cells, not bare volume.

The Taylor term has the bound
\begin{equation}\label{eq:main-taylor-error}
 |T_6|\le C\theta\mathsf n
       [\,\mathsf s\delta(E_1+L^{-N})+\mathsf s\delta^2\,].
\end{equation}
Its linear coefficient vanishes for the full action:
minimizing first under fixed averages and then over the
free datum removes precisely those constraints. The
derivative is therefore zero in all free box directions.
After restricting the action by \(\zeta\), the omitted
outer part is separated by the long collar and the
omitted inner part has the scaling \(L^{-2N}\).
These give the first two small factors in
\eqref{eq:main-taylor-error}. The quadratic coefficient
at the perturbed exterior differs from that at the
real centre by at most
\[
 \tfrac12\theta
       |q_2(\mathsf b)[\mathsf f,\mathsf f]-\mathsf q_0(\mathsf f)|
       \le C\theta\mathsf n\mathsf s\delta^2.
\]
Higher terms have at least a third field factor.

In the application \(\mathsf a=O(\alpha_{0,k}+\alpha_{1,k})\),
\(\delta=\delta'_k\) and \(\varepsilon=O(\varepsilon_k)\).
Hence \(\theta\mathsf s^2\), \(\mathsf n\) and \(N\)
are only logarithmic powers. Both \(E_1\) and \(L^{-N}\)
beat all these powers, while
\(\theta\mathsf s\delta^2\) has a positive power of \(g_k\)
times logarithms. The weight errors \(T_2,T_8\) have no
such geometric factor. They are charged to their
displayed local volume bounds, using the restored
profiles and, in a frozen run, its weight tube.

For clarity, the complex core comparison is the cellwise
response estimate
\begin{equation}\label{eq:xrg-complex-core-error}
 \sum_j\sum_{x\in\Gamma_j}
 Ce^{-\delta_{\mathrm G}d(x)}L^{-N}\alpha_{1,k}
       [x_k+C(k-j)]\check R_j^2\varepsilon_j
                       \le(\log x_k)^C L^{-N}.
\end{equation}
Response decay is measured in the layers' own units,
and dominates both their finer-site count and polynomial
growth of the old coupling. The real term
\(A(x''\phi_{\mathrm{core}},U^{\rm in}(g_kB))\)
therefore stays in \(\Phi_\omega^\circ\), nonnegative
at real stored data, while its exterior dependence
belongs to the small error. This is the separation
needed by the later supremum under the integral.

\subsubsection*{A quantitative denominator estimate.}

The same box calculation controls the reference denominator,
which cannot be treated merely as a bounded phase.
Let \(E\setminus\mathsf G\) be the non-tree free bonds,
\(n_{\mathsf G}=|E\setminus\mathsf G|=O(M^4)\), and
let \(\sigma(Y)\,dY\) be Haar measure in the chart
\(Y\mapsto e^{iY}\), with \(\sigma_0=\sigma(0)>0\).
Centre at the continued critical point \(B_\Lambda(\mathsf b)\)
and write the free variable as \(\theta^{-1/2}B'\).
After subtracting the basic action, the denominator is
\[
 \mathcal N(\mathsf b)=
 \int_{\substack{B'|_{\mathsf G}=0\\|B'|_\infty<\theta^{1/2}\mathsf r}}
 \prod_{b\in E\setminus\mathsf G}
       \frac{\sigma(\theta^{-1/2}B'(b))}{\sigma_0}\,dB'(b)\,
       e^{-\Psi(\mathsf b,B')}.
\]
Here \(\Psi\) is \(\theta\) times the localized action
difference from its critical value. Its quadratic
reference form is strictly positive by the fixed-box
tree Poincar\'e estimate. The linear term of the full
action is zero even at complex exterior data.
Taylor's formula, the variation of this Hessian, the
Haar density and the distant cutoff error are controlled
by
\begin{equation}\label{eq:main-denominator-smallness}
 \theta^{1/2}\mathsf r\ge1,\qquad
 C\theta\mathsf r
   [\,\mathsf r(\mathsf r+\mathsf a)
                      +(\mathsf s+\mathsf r)E_1\,]\le1.
\end{equation}
Take \(C\) large enough for the comparison error to be
at most one fifth. Comparison with the positive real
Gaussian reference integral \(\mathcal N_0\) gives
\[
 |\log\mathcal N_0|\le C(M),\qquad
 \left|\frac{\mathcal N(\mathsf b)}{\mathcal N_0}-1\right|
       \le\tfrac15.
\]
Thus the integral has no zero on the prescribed exterior
disc. The branch
\(I(\mathsf b)=-\operatorname{Log}\mathcal N(\mathsf b)\)
is holomorphic, real at real data, and satisfies
\(|I|\le C(M)\), \(|\operatorname{Im}I|\le1/4\).
Restoring the removed action and chart scaling gives
\[
 \int dV'\,
       \mathbf1_{\{|\frac1i\log V'|<\mathsf r\}}
       e^{-\theta A(\zeta_0,U_{k,Z}(V'V_\Lambda))}
 =e^{-\theta A(\zeta_0,U_0)}
       (\theta^{-3/2}\sigma_0)^{n_{\mathsf G}}e^{-I(\mathsf b)}.
\]
With \(\mathsf r=M_0\varepsilon_k\), the first condition
in \eqref{eq:main-denominator-smallness} holds and the
second tends to zero: its terms are positive powers
of \(g_k\) times logarithms, or logarithms times \(E_1\).
The box has \(O(M^4)\) bonds, independent of \(N\);
this fixed size is essential to the uniform comparison.

\subsection{Localising a change on a history}

The same response device handles an
ordinary fluctuation integral, a stage,
and the main step. The data change by
$\mathbf V\mapsto e^{ig_kB}\mathbf V$.
At each old level this gives a pair
shift $(\mathbf H,\mathsf j)$ with
\[
 e^{i\eta\mathbf H}U^{\mathrm{new}}
                  =(U^{\mathrm{old}})^u,
 \qquad
 J(U^{\mathrm{new}})+\mathsf j
                  =\operatorname{Ad}(u)J(U^{\mathrm{old}}).
\]
The response is nonlocal before it is expanded: changing a datum
near a box changes the background at every old level. Localisation
must therefore be performed in the equations that determine the
response, before an old activity is evaluated. Every propagator
in those equations is replaced by
\begin{equation}\label{eq:history-parameter-propagator}
 \mathsf G(s)=\sum_\omega
       \left(\prod_{\Delta\text{ met by }\omega}s(\Delta)\right)
                         \mathsf G_\omega,
       \qquad |s(\Delta)|<e^{\bar\kappa}.
\end{equation}
The fixed choice of $\bar\kappa$ costs a factor $e^{C\bar\kappa}$
in the analytic estimates. At $s(\Delta)=0$ no walk can pass
through $\Delta$, so the fixed-point equations decouple exactly.
The cube parameters are inserted in every propagator entering the
coupled field and current equations; inserting them only in the
last displayed response would not give that decoupling. A potential and its
stored data are then evaluated through
the same decoupled response. Only the
pair is shifted; the stored data stay
fixed under the pinned compensation
of \eqref{eq:xrg-pinned-gauge}. The complex potential in a triple
is not one of these stored real variables. Its analytic variation
is performed at fixed stored data.

Let an old term $f$ on $X$ have bound
$b_f$. Suppose its relative analytic
room for this shift is $a^{-1}$, so
the multiplier $\lambda$ is allowed
in $|\lambda|<a^{-1}$. At $(s,\lambda)=(\mathbf1,1)$ the
substitution gives the old background; when all response walks
are cut it gives the new one. Subtracting the value at
$\lambda=0$ makes the difference vanish there. The bracket
resummation in all cube parameters,
followed by Cauchy in $s$ and Schwarz
in $\lambda$, gives
\begin{equation}\label{eq:xrg-history-localisation}
 f|_{U^{\mathrm{old}}}-f|_{U^{\mathrm{new}}}
       =\sum_{Y\supset X}f_Y,
 \qquad
 |f_Y|\le2ab_f
           e^{-(\bar\kappa-1)\mathfrak m(X,Y)}.
\end{equation}
The length $\mathfrak m(X,Y)$ counts
the added cubes and the least number
of cubes of $X$ needed to connect
them to a cube carrying an integrated
field. Parameters are placed on cubes
of $X$ as well as on the added cubes.
Keeping a parameter equal to one does
not give its cube free: apply Schwarz,
Cauchy and the maximum principle in
the common scale of the kept
parameters. This rooted length is
needed because ordinary joining
through an exempt hull is unavailable.

The three relative amplitudes are
\begin{equation}\label{eq:xrg-history-amplitudes}
 \begin{array}{c|c}
 \text{move}&\text{relative amplitude}\\ \hline
 \mathbf T&\mathsf s_k=\delta_k/\min_\iota\alpha_{\iota,k}\\
 \text{stage at }j&N_k\mathsf s_j\\
 \text{main step}&g_k^{-1}e^{-cM\check R_k}.
 \end{array}
\end{equation}
The factor $N_k$ is the analytic
margin spent per stage, not a count
of independent small masses. The
conditions on $q-\mathfrak p_0$
make $N_k\mathsf s_j$ small.
For an island exponent the bound
applies to its oscillation, since
constants cancel in a difference.
The product of this amplitude with
$(\log x_j)^{2\mathfrak p_0}$ must
be bounded as well. This is why
the radius exponent is chosen above
the threshold exponents before the
coupling ceiling is fixed.

The analytic room through repeated
moves is recorded layer by layer.
The coefficient of $\alpha_{\iota,n}$
contains the baseline
$1-\beta_{\mathrm s}(1-2^{-(k-n)})$,
a reserve
\[
 \mathsf g_k(n)=
 \begin{cases}
 \theta_{\mathrm s}\beta_{\mathrm s}2^{-(k-n)},&n<k,\\
 \beta_{\mathrm s},&n=k,
 \end{cases}
 \qquad 0<\theta_{\mathrm s}<1,
\]
and charges for remade levels in
interrupted operations. Remaking
level $l$ costs
$c\beta_{\mathrm s}2^{-|n-l|}$ at
layer $n$. The mark of a drop-out
therefore records the number of
completed stages; the geometry
alone does not determine the room.

One ordinary level uses only the
fraction $1-\theta_{\mathrm s}$
of its geometric loss. A stage
uses $\mathsf g_k(n)/(2N_k)$
outside the component and the
remade-level charge inside it.
The main step uses half the
reserve, except on components
whose marks are forgotten in
the final label. There the
returned domain may become
larger, and no monotonicity
argument is available.
The response estimates fit the
successive gaps because
\[
 2^\nu(1+\nu)^{1/2}
       e^{-c\delta_{\mathrm G}ML(\nu-1)}\le C,
                     \qquad \nu\ge1.
\]
Thus the geometrically small
margin at a distant layer is
still larger than the response
that reaches it.

Boundary terms have no coupling
subtraction, so their sum across
levels uses two different decays.
If $p$ layers are not reached,
response decay gives
$e^{-c\delta_{\mathrm G}MLp}$.
If $q'$ layers are crossed by
the domain, its own length gives
$e^{-c\kappa L^{q'}}$.
Outside drop-outs, $p+q'=k-j$,
which pays for the
$L^{C(k-j)}$ positions at level
$j$ per current cube. A
drop-out can have a lag $l$
between its two region sequences,
so $p+l+q'=k-j$ instead.
The separate factor
$e^{-cL\check R_k}$ pays for
this lag. Neither response
decay alone nor domain decay
alone accounts for all positions.

\subsubsection*{Substitutions next to a prepared box.}

The old analytic domain contains conditions on more than the
plaquettes of the field at which an activity is evaluated. It also
requires regularity of the backgrounds determined by averages on
subdomains. A subdomain can begin near a prepared box and extend
past the component being integrated. Evaluating the old activity
only at the background of that component would then require an
unavailable continuation of this background outside the component.
The substitution must be defined on the whole activity domain.

For an activity $f$ of level $j$ on $\mathfrak Q$, write
$P=(U,A',\mathbf J)$ for the triple representing
$(e^{i\eta_jA'}U,\mathbf J)$, and let $\xi$ denote the stored
variables of the operation. Its evaluation is made at
\[
 \mathfrak S_{\mathfrak Q}(s,\lambda;P,\xi).
\]
The domain $\mathfrak Q$ determines the rule. The real member of
the triple must remain $\SU(2)$-valued and depend only on $U$ and
the group-valued variables of $\xi$. Holomorphy is required in
$s,\lambda,A',\mathbf J$ and the complex fluctuation variables.
Keeping this distinction is necessary when a real stored variable
appears under a sharp test.

\emph{Domains away from the boxes.}
If $\mathfrak Q$ misses the kernels $\Omega_k^c\cap X$ and keeps
away from the boxes, use the restriction of the new triple,
shifted by $\lambda$ times the response to the adjacent datum.
That datum uses the collar values of $P$ and the appropriate
variables of $\xi$. The response is determined by its fixed-point
equations with \eqref{eq:history-parameter-propagator}; thus
cutting the parameters also cuts the dependence on distant data.
This rule applies to any old level $j\le k$.

\emph{The old core and its innermost shells.}
Here the real member is fixed by the stored data. Denote it by
$U_\star$. It is the regular critical configuration, in the
chosen gauge, for the action over the plaquettes of the component
inside the first region $\Omega_0$ of its determining set.
On an averaging bond that carries a group-valued variable of
$\xi$, its datum is that variable; on the other bonds of the
determining set the datum is $\mathbf1$. A bare bond stored outside
$\Omega_0$ is not added to this datum. On the test support its
plaquettes obey the prescribed smallness bound in the units of
their own levels.

Start from $(U_\star,0,J(U_\star))$ and add $\lambda$ times the
response, computed at $U_\star$, to the difference between the
old datum of the activity and this fixed datum. This is the
anchored rule. At $\lambda=0$ it has no dependence on the exterior
triple $P$. This property permits its possibly large constant part
to remain under the stored-variable integral, rather than to be
bounded by an exterior supremum.

\emph{Other domains inside a component.}
If $\mathfrak Q$ lies wholly inside a component after removing
a margin of width $CM$, use the triple generated by that
component's datum. Its base point depends on $U$. The generated
configuration is put in the analytic gauge constructed below;
its real factor is not selected by a complex change of the
stored variables.

\emph{Domains crossing the transition region.}
The remaining domain meets a kernel or comes within
$M\check R_k$ of a box, and also reaches a cube outside the
component with its margin removed. Use the generated triple
near the kernel and the shifted triple at distance at least
$M\check R_k$ from it. In the transition region their potentials
and currents are interpolated using covariant forms. A plain
bondwise interpolation of arbitrary representatives does not
preserve the relation between two splittings of one complex pair.
The covariant choice is needed for the equality at the actual
evaluation points.

Such a level-$j$ domain crosses a distance at least
$cM\check R_k$ measured at level $k$. In its own units this is
$cM\check R_kL^{k-j}$. Hence its length satisfies
\begin{equation}\label{eq:joined-long-domain}
 d_j(\mathfrak Q)\ge c\check R_kL^{k-j}.
\end{equation}
The activity's own exponential decay can therefore pay for this
case even where the response is not small. The long-domain
estimate has a different role from the small relative amplitudes
in \eqref{eq:xrg-history-amplitudes}.

\subsubsection*{The anchor, its analytic room and its inner response.}

We verify the old-core rule more quantitatively. This is also
where one must distinguish the support of a single operation
from that of a completed effective term.
Set \(h=k-N_k\). Localize the prepared sequence to the
component \(C\), using its \(M_1\)-tower:
\[
 \Omega^C_m=\Omega''_m\cap C_m,\qquad
 \mathbf B^C=\mathbf B[\Omega^C_\bullet].
\]
The first member can have holes where bare fields were
stored outside \(\Omega_0\). Those bonds are absent from
the constrained problem, rather than assigned an arbitrary
regular extension. Let \(\mathbf V_\star\) equal the real
stored datum on its inner rows and one elsewhere.
The local operation tests make this datum regular.
Denote its critical background by \(U_\star\), with the
representative adapted to the stored sites.
Let \(\varrho_0\) be the separation constant of this
localized sequence:
\[
 \operatorname{dist}_\infty((\Omega^C_m)^c,\Omega^C_{m+1})
                         >\varrho_0\xi_m.
\]
This is the separated class used by the local regularity
estimate below. It need not have the larger gaps of the
earlier inverse construction with constant \(\rho_{\rm geom}\).

Write \(l(x)\) for the prepared layer index and let
\(\mathsf e^\sharp_m=Cg_m(\log x_m)^{\mathfrak p_0+O(\mathfrak r)}\).
Write \(\mathsf e_j\) for the bound on the scaled
test-loop defects of \(\mathbf V_\star\) in layer \(j\).
The local regularity estimate for this constrained problem
is
\begin{equation}\label{eq:anchor-layer-profile}
 \eta_k^{-2}\xi_{l(x)}^2|\partial U_\star-1|,
 \quad \xi_{l(x)}^3|J(U_\star)|
 \le C\sup_j \mathsf e_j
           e^{-c\delta_{\mathrm G}d_{\mathbf B^C}(x,\text{layer }j)}
 \le\mathsf e^\sharp_{l(x)}.
\end{equation}
The last inequality uses the square-root bound on ratios
of thresholds across levels and the separation
\(d(x,\text{layer }j)\ge
\varrho_0(|j-l(x)|-2)_+\).
Since
\[
 \frac{\mathsf e^\sharp_m}{\alpha_{0,m}}
       \le C(\log x_m)^{-(q-\mathfrak p_0-O(\mathfrak r))}
                         \longrightarrow0,
\]
this anchor is much smaller than the analytic radii.

\emph{Membership in every old inner domain.}
Fix an old target domain \(\mathfrak Q\) in the inner
region, with its structure preceding the last stage and
its largest prepared layer index at most \(h+1\).
Its determining sequence is finer than the prepared
sequence. Thus its index \(m\) at a point is at most
\(l(x)\). The scaled curvature and current at this
index have the additional factors
\(L^{-2(l-m)}\), \(L^{-3(l-m)}\).
The inequalities for thresholds and radii imply
\[
 L^{-a}\mathsf e^\sharp_{m+a}\le\mathsf e^\sharp_m,
 \qquad
 \mathsf e^\sharp_{m-1}\le C_L\mathsf e^\sharp_m.
\]
Every old radius coefficient is bounded below by
\(c_*>0\), and its lift is at least one.
Consequently all the pointwise bounds at
\((U_\star,0,J(U_\star))\) hold with half their
right-hand sides once the coupling ceiling is small.

The potential bounds concern small cubes, some touching
a layer boundary. For an index \(m\ge1\), construct the
regular gauge one level below: sequence separation places
the needed neighbourhood in \(\Omega^C_{m-1}\).
The local gauge theorem gives
\[
 \xi_m|A|+\xi_m^2|\nabla A|
                   \le C_LM\mathsf e^\sharp_m.
\]
The neighbourhood may meet the boundary of a layer but
does not meet a hole. At index zero use a comb inside
the bare cube itself. In a cube of at most \(s+1\)
sites per direction, a bond in direction \(\mu\)
is obtained from the comb by at most \(3s\)
plaquette multiplications. Hence
\[
 |U_\star^v(b)-1|\le3s\bar{\mathsf e},
 \qquad
 \eta_k|A(b)|\le\tfrac{3\pi}{2}s\bar{\mathsf e},
 \qquad
 \eta_k^2|\nabla A|\le3\pi s\bar{\mathsf e},
\]
where \(\bar{\mathsf e}=\sup_m\mathsf e^\sharp_m\).
Here \(2\sin(\varphi/2)\ge2\varphi/\pi\) supplies
the logarithm bound in \(\SU(2)\). All plaquettes of
the comb lie inside the cube, so this argument still
applies to cubes next to a hole.

The tower clauses require another observation, beyond
these pointwise estimates. A regular critical configuration
is its own background over any finer subordinate system
of its averages. Its restriction is admissible for that
system, remains critical, and lies in the regular class;
uniqueness of the critical orbit identifies the new
tower background with a gauge copy of the restriction.
Its curvature and current are therefore those in
\eqref{eq:anchor-layer-profile}, rescaled to the
index read by the clause. This checks every subdomain
tower clause, rather than just towers of the whole
component. For a single-scale clause inside an old
small-field region, truncate the same sequence at that
scale and apply the same argument, with its prescribed
margin. Thus the pointwise, good-gauge, tower and
single-scale groups all hold at half radii.
This proof uses separated sequences and their local
regularity theorem; it does not apply a membership
criterion that requires the stronger wide-sequence
condition at the tower along \(\partial C\).

\emph{Regularity at holes and the pins.}
At a bond next to a hole, all adjacent present plaquettes
are in the bare layer. The current of the action on
\(\Omega^C_0\) uses only those plaquettes, at most six,
and consequently
\[
 \eta_k^3|J(U_\star)(b)|\le6\bar{\mathsf e}.
\]
At a bond with all its plaquettes present the identity
\[
 (\mathrm d^*_{U_\star}\partial U_\star)(b)
     =a(b)\mathbf1+i\eta_k^2J(U_\star)(b),\qquad
 |a(b)|\le\tfrac32\bar{\mathsf e}^{\,2}
                                 \eta_k^2\xi_{m-1}^{-3}
\]
converts the current bound to the input of the local
gauge theorem. This follows by writing
\(\partial U(p)=c_p\mathbf1+i\eta_k^2F_U(p)\),
where \(1-c_p=\tfrac12|\partial U(p)-1|^2\).
For a top of level \(m\ge1\), separation places a fixed
\(\xi_m\)-neighbourhood of its pillar inside
\(\Omega^C_0\), provided \(\varrho_0/L\) exceeds the
fixed pinning reach. All such tops are therefore
available to the pinned construction. At a bare
inner row the datum increment is zero. The constraint
then fixes the configuration on that bond exactly,
and both endpoint node factors are one. The Landau
response is zero there by continuity from zero
increment. These checks are needed before applying
the pinned response theorem to a component with holes.

\emph{A slowly varying room.}
Let \(\mathsf s(x)\) be a nonnegative profile, varying
at rate \(\theta_\flat\) in the prepared metric, with
\[
 \mathsf s(x)\le
       \kappa_{\rm R}\min_\iota\alpha_{\iota,l(x)}.
\]
A shift of this size in the prepared units has size at
most
\(\widehat{\mathsf s}_m(x)=L^{-(l(x)-m)}\mathsf s(x)\)
at an older index \(m\). The radius-ratio estimate gives
\begin{equation}\label{eq:anchor-room-rescaling}
 \widehat{\mathsf s}_m(x)
 \le C\kappa_{\rm R}L^{-(l-m)}(1+l-m)^{1/2}
                            \min_\iota\alpha_{\iota,m}
 \le C'\kappa_{\rm R}\min_\iota\alpha_{\iota,m}.
\end{equation}
Potential, curvature and current changes then spend
at most half of each pointwise radius if
\(\kappa_{\rm R}\) is chosen after these fixed constants.
The real factor and its local gauges have not changed.

For a tower quantity \(\mathsf G\), its Lipschitz bound
with decay gives
\[
 |\mathsf G(U_\star,\mathbf A;x)
                    -\mathsf G(U_\star,0;x)|
 \le K_{\rm T}\sup_b
       e^{-\delta_{\rm T}d^\flat(x,b)}
                           \widehat{\mathsf s}_{m_b}(b).
\]
The old localized metric satisfies \(d^\flat\ge d''\).
Layer indices change by at most
\(1+d^\flat/\varrho_0\), so
\[
 \widehat{\mathsf s}_{m_b}(b)
 \le L^2e^{(\theta_\flat+2\log L/\varrho_0)d^\flat(x,b)}
                                      \widehat{\mathsf s}_m(x).
\]
Choose
\(\theta_\flat+2\log L/\varrho_0\le\delta_{\rm T}\).
The last two displays and
\eqref{eq:anchor-room-rescaling} then spend less than
half the tower radius. Applying this at each index
also checks the single-scale tower clauses.
Thus \(\mathsf s\) is actual room in the whole old
domain, not merely a bound on its pointwise coordinates.

\emph{The inner datum and its depth profile.}
The increment \(B^\natural\) over the anchor is zero
on all real inner rows. On a free shell row of level
\(n<k\) it obeys
\[
 e^{iB^\natural(c)}=e^{ig_kB(c)}e^{i\nu_P(c)},\qquad
 |\nu_P(c)|\le C_1'L^{n-k}a_C,
\]
where \(a_C=\max_\iota\alpha_{\iota,k}\).
The remaining outer and tower rows have
\[
 |B^\natural(c)|\le\epsilon_{\rm out},
 \qquad \epsilon_{\rm out}=C_1'M^2\check R_k a_C\longrightarrow0.
\]
The shell estimate is
\eqref{eq:main-basic-shell-increment}; entrance rows
are bounded by closing their bond with the tree and
\(O(M)\) exterior bonds. The closed path spans
\(O(M^2)\) plaquettes of the critical extension.
For invariance under a change of splitting, the
exterior comb is formed from the complex pair.
It acts at the nodes of the datum; it is not a
complex transformation of the stored real variables.

For \(h\le n\le k\) define
\[
 \mathsf D_n(x)=
   \inf\{d''(x,z):z\in C\setminus I_2,\ l(z)\ge n-1\},
\]
and, taking a fixed \(\theta_\flat<\delta_{\mathrm G}\), set
\[
 \begin{split}
 \mathsf s_\star(x)=K_{\rm S}\bigg[
 &\delta'_ke^{-\theta_\flat\mathsf D_h(x)}
 +C_1'a_C\sum_{n=h}^{k-1}
                 L^{n-k}e^{-\theta_\flat\mathsf D_n(x)}\\
 &+\epsilon_{\rm out}e^{-\theta_\flat\mathsf D_k(x)}
 \bigg].
 \end{split}
\]
The response and current have this majorant, up to a
fixed constant. Each distance is one-Lipschitz, so the
profile varies slowly. For \(l=l(x)\le h+1\), separation
gives \(\mathsf D_n(x)\ge\varrho_0(n-2-l)_+\), and
\[
 \sum_{n=h}^{k-1}L^{n-k}
        e^{-\theta_\flat\varrho_0(n-2-l)_+}
 \le6L^{3-N_k}
              e^{-\theta_\flat\varrho_0(h-2-l)_+}.
\]
Use \(Le^{-\theta_\flat\varrho_0}\le1/2\) here.
Both
\(a_CL^{3-N_k}/\delta'_k\) and
\(\epsilon_{\rm out}e^{-\theta_\flat\varrho_0(N_k-3)}
 /\delta'_k\) tend to zero. Consequently
\begin{equation}\label{eq:anchor-response-profile}
 \mathsf s_\star(x)\le
       3K_{\rm S}e^{\theta_\flat\varrho_0}\delta'_k
             e^{-\theta_\flat\varrho_0(h-1-l(x))_+}.
\end{equation}

Let \(n_{\mathfrak Q}\) be the largest layer index of
the target and \(u_{\mathfrak Q}=(h-1-n_{\mathfrak Q})_+\).
The pinned response costs a fixed profile-transfer
constant. Combining \eqref{eq:anchor-response-profile}
with the analytic room above gives the interpolation
disc \(|\lambda|<2/a^\star_{\mathfrak Q}\), where
\begin{equation}\label{eq:anchor-explicit-amplitude}
 \begin{split}
 a^\star_{\mathfrak Q}
 &\le C(2+N_k)^{1/2}\frac{\delta'_k}{a_*}
                       e^{-\delta_\flat\varrho_0u_{\mathfrak Q}}\\
 &\le C'(\log x_k)^{-(q-\mathfrak p_1-\mathfrak r/2)}
                       e^{-\delta_\flat\varrho_0u_{\mathfrak Q}} .
 \end{split}
\end{equation}
Here \(a_*=\min_\iota\alpha_{\iota,k}\);
\(\delta_\flat<\theta_\flat\) is chosen so that
\((1+u)^{1/2}e^{-\theta_\flat\varrho_0u}
\le e^{-\delta_\flat\varrho_0u}\).
The factor \((2+N_k)^{1/2}\) comes from
\(1+k-l\le(2+N_k)(1+(h-1-l)_+)\).
Choose the separation so that
\(e^{-\delta_\flat\varrho_0}\le(4L^8)^{-1}\).
The resulting decay is faster than the growth of old
stored sites and their fixed multiplicities.

\emph{Comparison with a real stored-only point.}
One can retain more than the zero-response anchor.
Let \(b_\xi\) equal \(g_kB\) on the free inner
shell rows and zero elsewhere. It depends only on
the real variables of the operation. The interpolation
at this datum defines a real triple \(\mathfrak q_{\mathfrak Q}(\xi)\)
whose pair is the pair of the inner background
\(W^{\rm in}(g_kB)\).

To bound the exterior dependence, set
\[
 B(z)=\tfrac1i\log(e^{ig_kB}e^{iz\nu_P})
       \text{ on the free rows below }k,\qquad
 B(z)=zB^\natural\text{ on the other non-inner rows}.
\]
It is zero on the fixed inner rows,
\(B(0)=b_\xi\), and \(B(1)=B^\natural\).
It stays in the background's fixed datum polydisc on
\(|z|<r_{\rm B}/(2\epsilon_{\rm out})\).
Split \(\mathsf s_\star=\mathsf s^{(1)}+\mathsf s^{(2)}\),
where \(\mathsf s^{(1)}\) is its first summand.
The crude response bound on that disc and Schwarz'
lemma give
\[
 |\mathfrak h(B^\natural)-\mathfrak h(b_\xi)|_b
 \le\tfrac32K_1\left[
       \frac{4\epsilon_{\rm out}}{r_{\rm B}}\mathsf s^{(1)}
                                 +\mathsf s^{(2)}\right](x_b).
\]
Apply the scalar Schwarz estimate to each linear
matrix functional defining the shift norm; taking
their supremum gives this same bound for the norm.
Thus the large \(\delta'_k\) contribution acquires
the extra small factor \(\epsilon_{\rm out}\).

Interpolate between these two shifts, with the real
factor fixed at \(U_\star\). Its analytic room gives
an amplitude bounded by
\begin{equation}\label{eq:anchor-stored-comparison-amplitude}
 \frac{C(2+N_k)^{1/2}}{a_*}
 \left[
  \epsilon_{\rm out}\delta'_k+
  a_CL^{3-N_k}
  +\epsilon_{\rm out}e^{-\theta_\flat\varrho_0(N_k-3)}
 \right]e^{-\delta_\flat\varrho_0u_{\mathfrak Q}}.
\end{equation}
For a bounded old activity \(f\), Schwarz' lemma
therefore bounds its change from
\(\mathfrak q_{\mathfrak Q}(\xi)\) by this amplitude
times \(\sup|f|\). The value at that real point can
remain under its own stored-variable integral.
This refined comparison explains how the paying
real inner core action is retained while the
complex exterior dependence becomes a small local
interaction.

\subsubsection*{Checking the conditions on subdomain backgrounds.}

The domain of an old activity includes bounds on the backgrounds
of all its subdomain averages. We give the localization and gluing
argument for those conditions, isolating the uniform comparison it
uses. A bound on the plaquettes of the substituted field alone
would not check these inherited conditions.

At a layer of index $m$, write $\gamma_m$ for the additive room
reserved between the prepared and output analytic domains. If $t$
is the level of the operation and $\alpha_m^\vee$ is the smallest
radius occurring in that clause, the room obeys
\[
 c\beta_{\mathrm s}2^{-(t-m)}\alpha_m^\vee
                      \le\gamma_m\le C\alpha_m^\vee.
\]
Inside a prepared region the endpoint estimates instead leave a
fixed relative fraction of their right sides. The distinction is
essential at older layers, where the additive room can decrease
with $t-m$.

\emph{Cutting a subdomain at a distant boundary.}
Let $\widetilde Z$ be a subdomain appearing in one of the old
background conditions, and let $x$ be the point where its tower
curvature or current is tested. Here a tower is the sequence of
backgrounds obtained from its successively localized averaging
data. Choose a union $R$ of current cubes with
\[
 \operatorname{dist}_\infty(x,\widetilde Z\setminus R)
                               \ge D+3M_1,\qquad D\ge CM.
\]
Let $\widetilde Z_R$ be the connected component of
$\widetilde Z\cap R$ containing $x$. The partition at the old
level refines the current partition. Thus the cut is again an
admitted subdomain. Membership in a localized layer through the
depth being read is decided by the complement within $2M_1$
spacings. The original and cut domains therefore have identical
layers and averaging rows in the required neighborhood of $x$.
Other components of the cut are separated by a full old cube.
All rows changed by the cut lie at least distance $D$ from $x$;
the layer index used to normalize the tested quantity is unchanged.

For the two resulting towers, the analytic estimate needed here is
\begin{equation}\label{eq:joined-tower-comparison}
 \left|\,|G^{\rm full}(U,A';x)|
                 -|G^{\rm cut}(U,A';x)|\,\right|
 \le C_{\rm tw}a_*^{-1}e^{-c\delta_{\mathrm G}D}\gamma_m.
\end{equation}
The quantities $G$ are the curvature or current with the layer
normalization of its analytic-domain clause; $a_*$ is the smallest
top-layer radius. The constant $C_{\rm tw}$ depends only on the
fixed construction parameters, independently of $D$, the levels,
the subdomain and its position. This is a comparison of their norms on the
whole admitted complex domain. It must have no additional
quadratic remainder in $A'$. Such a remainder need not fit into
$\gamma_m$ at deep layers. The geometric cutting argument reduces
the inherited condition to this estimate, but does not prove its
uniform analytic bound.

\emph{Gluing the domain conditions.}
Cover the activity domain by overlapping pieces on each of which
one endpoint construction, or the controlled interpolation,
already has its complete domain bounds. Each point $x$ is assigned
a piece containing its $D+3M_1$ neighborhood. Require also that
the cubes outside the prepared region and their adjacent cubes
belong to one common exterior piece. Suppose the local endpoint
estimates leave relative room $r\beta_{\mathrm s}$ inside the
prepared region and additive room $r\gamma_m$ elsewhere.
The conditions for the real member $(U,0)$ have their own rooms,
denoted by $\kappa_3$ and $r'\gamma_m$. These do not follow
solely from the conditions for $(U,A')$.

There are four kinds of condition to check. Pointwise bounds,
potentials on small cubes and the local good-gauge conditions are
read within distance $CM$ of their base point, so they are
conditions of its assigned piece. For a tower condition, cut the
subdomain to that piece as above. Its right side and layer index
are unchanged. Apply \eqref{eq:joined-tower-comparison} to compare
the full and cut quantities and spend half the saved room.
For the old target spaces used by this substitution, a current-top
single-scale domain lies in $\Lambda_t$ of the old label, while
each treated component lies in its complement. Those domains are
therefore wholly in the common exterior piece. This assertion is
about the old targets; a newly healed output domain can include the
treated component. For an older single-scale condition,
its local parts belong to the corresponding cut components;
its subdomain-tower parts use the same cutting argument, both at
$(U,A')$ and at $(U,0)$.

More quantitatively, the right side $b_m$ of a tower clause obeys
$\gamma_m\le Cb_m$. It suffices to choose the fixed parameters so
that
\[
 C_{\rm tw}a_*^{-1}e^{-c\delta_{\mathrm G}D}
             \le c\min\{r\beta_{\mathrm s},\kappa_3,r,r'\}.
\]
Then each inherited clause holds with at least half its original
room. The proof compares one clause at its own base point; it
does not sum the truncation errors over all subdomains or layers.

The prepared-box geometry provides the required cover by an
interior piece, a generated-field piece, an interpolation strip
and an exterior piece. The strip cutoff is constant near its two
ends, with its transition confined to a shell of width comparable
to $M\check R_t$. The overlap margins can be chosen to be at
least $M\check R_t/16-M$. For $\check R_t$ above a fixed
constant this exceeds $M\check R_t/32+3M_1$. Thus one can use
$D=M\check R_t/32$, uniformly through a seam of stored layers.
The exterior piece contains the cubes outside the prepared region
together with their adjacent cubes, as required by the top-level
single-scale clause.

The factor $a_*^{-1}$ is at most an inverse coupling times a fixed
logarithmic factor. With
$\check R_t\ge(\log(x_t-C_\sharp))^{\mathfrak r}$ and
$\mathfrak r\ge2$, the exponential in $D$ beats every fixed
power $g_t^N$. Choose the required power and all room fractions
before the coupling ceiling. This makes the displayed numerical
condition possible once the uniform analytic estimate and
endpoint rooms have been established.

This also explains where a fixed enlargement of the radii can be
used. On a collar the saved room is a relative fraction, and a
small common enlargement of the two radii is harmless before
returning to the true target. Away from it the room is the
layer-dependent $\gamma_m$; a fixed relative enlargement would
exceed it at sufficiently old layers. The first remaining analytic
verification is therefore the uniform no-remainder estimate
\eqref{eq:joined-tower-comparison}, together with the endpoint
bounds furnishing the stated rooms on every permitted generated
triple. The cutting and gluing deduction applies once those
estimates hold, including at the seams.

\subsubsection*{Oscillations and compatibility of the substitutions.}

The response construction is to give an amplitude
$a_{\mathfrak Q}\le1$ for which
$\mathfrak S_{\mathfrak Q}(s,\lambda)$ remains in the domain of
$f$ on $|\lambda|<2/a_{\mathfrak Q}$. On anchored domains the
amplitude contains a negative power of $\log x_k$ and decays in
$(k-N_k-1-n)_+$, where $n$ is the highest level of a field stored
on $\mathfrak Q$. The ratio is at most $(4L^8)^{-1}$ per deeper
stored level. This decay is stronger than the $L^4$ growth of stored
sites when that level is lowered.

Apply the Schwarz estimate in $\lambda$ to the difference from
its value at zero. Once the common domain and enlarged radius
have been verified, this gives
\begin{equation}\label{eq:substituted-triple-oscillation}
 \big|f\circ\mathfrak S_{\mathfrak Q}(s,1)
             -f\circ\mathfrak S_{\mathfrak Q}(s,0)\big|
                  \le a_{\mathfrak Q}\sup|f|.
\end{equation}
At zero the rule is independent of $s$. For an anchored domain
it depends only on $\xi$; that value is absorbed into
$\Phi^\circ_\omega$. Comparing the values at two exterior
arguments therefore costs at most $2a_{\mathfrak Q}\sup|f|$.
It is this oscillation, rather than the stored-only exponent,
that is bounded uniformly in the exterior pair. The same
amplitudes enter \eqref{eq:xrg-history-localisation}.

The rules must also agree when a localisation domain is viewed
inside a larger one. This is consistency with restriction of
the local construction, not dependence on $s$ cube by cube:
a parameter carried by one cube of a component can affect a
function of the component's whole collar. The rooted length in
\eqref{eq:xrg-history-localisation} was chosen to retain that
cost. Complex conjugation and the local charts of a history
must transport the rules together with their stored data.
Real gauge transformations act on the stored variables only
through the value at the root of their gauge tree. Thus an old
term invariant under this action retains that invariance under
the substitution.

At the evaluation points, $s$ is an indicator, $\lambda=1$ and
$\xi$ is real. The substituted pair must then be the pair of
the critical field for the old datum. Equality for two splittings
of that pair is proved using the pinned gauge and the common
holomorphic family in the discussion following
\eqref{eq:xrg-pinned-gauge}. Its applicability still requires the
uniform spare radii identified above. This is a statement at the
evaluations used in the component integrals; no saturation of the
whole analytic domain by complex gauges is required.

Near a stored boundary the field is subject to
sharp tests and need not enter analytically.
Each such bond $b$ of level $l$ has a switch
$\mu_b$, with radius a large power of $\log x_l$.
The boundary function is holomorphic jointly in
these switches, the pair and far stored fields;
it is Borel in near stored fields. At $\mu_b=0$
it is independent of the field at $b$, and at
zero field it is independent of $\mu_b$.
Its actual value is at $\mu_b=1$. This structure
is used to control the successive stage integrals
and their logarithms. Uniform nonvanishing of
those complex integrals must be established with
these switches and bounds independent of depth.


\subsection{The stage expansion and the generated field}

Each stage moves a background over the whole torus, although
its change is largest at the collar being integrated.
Its exponent must therefore contain the difference of
every previously formed factor between the old and new
backgrounds. Besides the regular and boundary interactions
and the terms of earlier stages, there are seven kinds:
the forms integrated in this stage; island factors near
the treated component; the remaining island factors;
forms of levels still waiting inside the component;
forms on untreated components; forms below the whole
waiting window; and logarithms of pivot-elimination factors.
The first kind belongs to the Gaussian quadratic form.
The other kinds enter as localized differences.
Characteristic functions in stored operations are not
differenced: their genuine independence of the moving
background is necessary here.

\subsubsection*{Used variables, direct terms and Gaussian tails.}

Partition the variables into finite blocks. For a potential
\(V\) and a subset \(S\) of its variable blocks put
\[
 V_S(B)=\sum_{T\subset S}(-1)^{|S\setminus T|}V(B_T),
\]
with unlisted blocks zero. Summing over \(S\) recovers
\(V(B)\). For \(S\ne\emptyset\), the summand vanishes
when any listed block is zero. Schwarz' lemma in each
block therefore gives a factor
\[
 q_v=\frac{\text{range of block }v}
                 {\text{analytic radius of block }v}
\]
per block actually used. Its skeleton consists of those
blocks, together with the connecting walks. The full
localization domain of \(V\) is not included merely
because \(V\) was initially defined on it.

For a new field on the integrated shell the ratio is
\(\delta'_j/\delta_j\). A far stored field has a
polydisc \(\lambda^{\rm st}\) times its real range.
A near stored field is real and may stand under a sharp
test. There the analytic variable is a switch \(\mu_v\),
not that field: at \(\mu_v=0\) the function is independent
of the field, and at zero field it is independent of
\(\mu_v\). Evaluating at \(\mu_v=1\) gives the same
factor \((\lambda^{\rm st})^{-1}\) without differentiating
the characteristic function. Set
\begin{equation}\label{eq:xrg-stage-switch-size}
 \lambda^{\rm st}=(\log x_j)^{6\mathfrak r+1}.
\end{equation}
This switch property is preserved when the terms of
earlier stages re-enter the next integral.

The part \(V_\emptyset\) does not use the current
integration variables. Keep it outside that integral
as a direct term. The other parts are expanded into
a gas of configurations. A configuration records its
used blocks, its connecting walk chains, and any
blocks outside the preferred Gaussian range.
Write \(H\) for the latter blocks. Reserving a small
Gaussian weight \(\gamma_2\) assigns the tail
\[
 \mathsf p=
       \exp[-\tfrac12\gamma_2\mathfrak b_{\rm cut}^{\,2}]
\]
to each such block, where \(\mathfrak b_{\rm cut}\) is
the auxiliary field threshold. The retained part of
the quadratic estimate must still control the
remaining Gaussian integral.

The general configuration bound has the form
\[
 |\operatorname{val}(\mathbf c)|
 \le e^{c_{\rm v}\#Z(\mathbf c)}
       (C\mathsf p)^{|H|+|P_1|+|P_2|}
       e^{-(\kappa_1-1)\#\tau}
       \prod_{(\iota,S)\in\mathbf D}
                       e^{\bar u}u_\iota(S)\prod_{v\in S}q_v.
\]
Here \(P_1,P_2\) are the additional large-block lists
created by the Gaussian expansion, \(\tau\) its
nonlocal walk cubes, and \(Z(\mathbf c)\) its
integration support. The number \(u_\iota(S)\)
bounds the block component of its potential.
The full potential amplitude is bounded by \(u_\iota\),
and \(\bar u\ge\sup_\iota u_\iota\).
This majorant is independent of the interpolation
parameters and operator entries. Its extensive
factor counts variables of the configuration,
rather than the volume of every potential domain.

One residual dependence on a potential's used blocks
must be included. A large block in one configuration
can affect a potential in another. If
\(\mathfrak V_\iota^{\rm h}\) is the set of its
blocks that can be large, charge that potential
\(e^{2a_1|\mathfrak V_\iota^{\rm h}|}\).
It suffices to choose \(a_1\) as small as the
large-block activity. The configuration criterion
is that both the sum of weighted used-block
activities through one block and
\[
 C\mathsf p(e\mathsf K)^{\mathsf n_{\rm t}}
                                      e^{2\mathsf u_{\rm hot}}
\]
are small, with the latter at most \(a_1\).
The fixed \(\mathsf K,\mathsf n_{\rm t}\) account
for local variables and their multiplicity, and
\(\mathsf u_{\rm hot}\) is the load of potentials
that can meet a large block. These are the
quantities controlled by the Gaussian expansion;
a bound on the phase of an arbitrary Borel
potential would not imply this criterion.

\subsubsection*{The recursion through the stages.}

There are two useful norms. The action and primary
interactions are summed over their domains through
a fixed cell, with their exponential length weight.
On a graded cell of level \(\mu\), use the normalization
\(L^{-4(\mu-j-1)}\) for its number of level-\(j+1\)
sites. Let \(J^{\rm a}\) be this action load,
\(J^{\rm pr}\) the primary load, and
\(J_{\mathsf A}^\Sigma\) the action load summed
over the anchor cells of one component.
For direct terms retain separately their largest
single-stage norm \(\bar D_1\), their sum over
stages \(\bar D_1^\Sigma\), and their accumulated
load \(\bar D^\circ\) through a fixed cell.
A normalizer can have a controlled summed direct
norm without being assigned a small independent
constant at every stage.

At stage \(n+1\), its input includes both direct
and cluster parts from every stage \(\nu\le n\).
For \(N=N_k\) the induction uses
\[
 \bar{\mathsf J}=J^{\rm a}+CN+\bar D^\circ,\qquad
 \mathsf u_{\rm hot}\le C(\bar{\mathsf J}+J^{\rm pr}).
\]
Write \(N^*\le C\check R_k^4\) for the number of
possible anchor cubes in a component. Let
\(\hat q,\hat q^{\rm h}\) be the largest
used-variable ratios, including their fixed local
entropy factors, for all blocks and for stored
blocks respectively. The component load to control
is
\begin{equation}\label{eq:stage-component-load}
 \begin{split}
 \epsilon_{\rm comp}={}&
 CN^*\mathsf p\,e^{2\mathsf u_{\rm hot}}\\
 &+C\left[
 \hat q^{\rm h}(J_{\mathsf A}^\Sigma+N^*CN)
       +\hat qN^*(\mathsf a_{\rm st}\bar{\mathsf J}+J^{\rm pr})\right].
 \end{split}
\end{equation}
Here \(\mathsf a_{\rm st}\) is the small relative
shift amplitude entering the stage expansion. Every previous-stage
term is assigned the crude bound one when
checking this load. Thus its accumulated pinned
norm is \(O(N)\); the required smallness comes
from its used-variable ratio, rather than an
independent small mass assigned to that term.

The largest stored-field contribution, after summing
over the \(N\) stages, is
\begin{equation}\label{eq:xrg-stage-accumulation}
 CN^2\check R_k^4(\lambda^{\rm st})^{-1}
                       \le C(\log x_k)^{-1}.
\end{equation}
There is one \(N\) for the old stage load and
one for the sum of the newly made terms.
For new-shell variables the separate ratio
\[
 N^*\frac{\delta'_j}{\delta_j}
       \le C\check R_k^4
                      (\log x_j)^{-(\mathfrak p_0-\mathfrak p_1)}
\]
must be small. The source of smallness in these
two expressions is different. The first explains
the exponent \(6\mathfrak r+1\) in the switch;
the second explains the additional margin
between \(\mathfrak p_0\) and \(\mathfrak p_1\).
The logarithms throughout a waiting window
are comparable, so one coupling ceiling works
for all its stages.

The large-block term in
\eqref{eq:stage-component-load} must be controlled
at the same time, including its
\(e^{2\mathsf u_{\rm hot}}\) factor. Put
\(\mathsf L_k=\log x_k\). The crude bound one for each
earlier stage gives a pinned norm \(O(N)\); after the other
inputs are included, the bound used at a block which can be hot is
\[
 \mathsf u_{\rm hot}\le C\mathsf L_k^{4\mathfrak r}.
\]
At the characteristic functions of the integral there is no
used-variable ratio multiplying this exponent. It must be paid
by the Gaussian tail itself. The required tail bound is
\[
 \mathsf p\le
 \exp\!\left[-c\mathsf L_k^{2\mathfrak p_1-4\mathfrak r}\right].
\]
Since \(N^*\le C\mathsf L_k^{4\mathfrak r}\), it gives
\begin{equation}\label{eq:stage-hot-tail-budget}
 CN^*\mathsf p\,e^{2\mathsf u_{\rm hot}}
 \le C\mathsf L_k^{4\mathfrak r}
 \exp\!\left[-c\mathsf L_k^{2\mathfrak p_1-4\mathfrak r}
                         +C\mathsf L_k^{4\mathfrak r}\right].
\end{equation}
The right side is smaller than every negative power of
\(\mathsf L_k\) once \(2\mathfrak p_1>8\mathfrak r\) and the
coupling ceiling is sufficiently small. It therefore also pays
the fixed logarithmic entropy and variable-ratio factors in the
stage expansion. This is a separate use of the sharper exponent
\(\mathfrak p_1\), in addition to the small new-shell ratio
\(\delta'_j/\delta_j\). The tail estimate still requires the
uniform decoupled covariance and mean bounds for the actual
collar integral, on its stated characteristic functions.
Neither \eqref{eq:xrg-stage-accumulation} nor smallness of the
stored-variable ratios proves that estimate.


Closing the free source has its own analytic cost.
During the expansion an auxiliary analytic source keeps
the local parts of the background shift independent.
Closing this source substitutes their actual values.
If \(R_{b,W}>1\) is the available radius for the
closing variable attached to a bond \(b\) and
carrier \(W\), and \(\mathsf c(b)\) is the designated
cell carrying \(b\), put
\[
 \epsilon_{\rm cl}=
 \sup_\Delta\sum_{(b,W):\,\mathsf c(b)=\Delta}
       \frac{e^{C\rho_\natural\varrho_{\rm a}\#W}}{R_{b,W}-1}.
\]
The stage criterion also requires
\[
 C\hat qN^*\le1,\quad
 C\mathsf p e^{2\mathsf u_{\rm hot}}\le\hat q,\quad
 C\epsilon_{\rm comp}\le1,\quad
 CN^*\epsilon_{\rm cl}\le1,\quad
 N\ge C\varrho_{\rm a},\quad
 \kappa_1\ge C\rho_\natural\varrho_{\rm a},
\]
where \(\varrho_{\rm a}\) is the fixed geometric
anchoring ratio. These conditions control closing
the free analytic source as well as the
used-variable activities.

Spatial decay is spent in units \(1/N\).
If \(\rho_\natural\) is the input rate, set
\(\rho_n=\rho_\natural-C_\rho n/N\).
Supports of separate component operations are
separated on the scale \(NM\).
Consequently the count of components met by
a domain of length \(d\) is \(1+Cd/N\),
and the entropy costs have the form
\(e^{Cd/N}\). The budget per stage includes
one \(1/N\) for used-block support entropy,
one for passing from a summed norm to a
pinned norm, and fixed further multiples
of \(1/N\) for component sums and closing
the free source. The total loss after \(N\)
stages is the fixed \(C_\rho\).

On the configuration criterion above, the
connected expansion gives a nowhere-zero
stage integral and its logarithm. The
direct part is carried outside the integral;
the source part is closed by the same
small analytic shift. The resulting bounds
are
\[
 |\mathbb C_{\rm d}^{(n+1)}(Y)|
       \le\bar D_1e^{-\rho_\natural d_k^\natural(Y)},\qquad
 |\mathbb C_{\rm cl}^{(n+1)}(Y)|
       \le2\epsilon_{\rm comp}
              e^{-(\rho_{n+1}-C/N)d_k^\natural(Y)}.
\]
Each cluster part contains an anchor cell
of its stage. The stronger summed norm
near one component is at most
\(C\epsilon_{\rm comp}\), without the
extra pointwise component-count loss.
Induction includes the old direct and
closed parts in the next input, so these
are bounds on the actual recursion.

Choose the input rate with the budget
just described. The sum of the direct terms and
the cluster estimate is bounded by
\[
 (\bar D_1^\Sigma+2N\epsilon_{\rm comp})
       e^{-(\rho_\natural-C_\rho-C)d_k^\natural(Y)}.
\]
Thus the application must check
\(\bar D_1^\Sigma+2N\epsilon_{\rm comp}\le C_{\rm S}\);
smallness of \(\epsilon_{\rm comp}\) alone would
not suffice. Equation \eqref{eq:xrg-stage-accumulation}
is the largest stored-field contribution to this
check. With the remaining action, primary and
direct loads controlled in their stated norms,
the summed output is
\begin{equation}\label{eq:xrg-stage-logarithms}
 \begin{gathered}
 \mathbb I_k^{(n+1)}
       =\exp\left[\mathrm{const}
                       +\sum_Y\mathbb C_k^{(n+1)}(Y)\right],\\
 \sum_{n<N_k}\sup|\mathbb C_k^{(n+1)}(Y)|
       \le C_{\rm S}
               e^{-(1+3\beta_{\mathrm s})\kappa d_k^\natural(Y)},
 \qquad C_{\rm S}\text{ independent of }N_k .
 \end{gathered}
\end{equation}
If the operation stops after \(n\) stages,
the one function
\(\sum_{\nu\le n}\mathbb C_k^{(\nu)}(Y)\)
is retained as its boundary contribution.
No earlier stage integral is undone.

The analytic verification in this application is the
uniform configuration criterion for its actual
decoupled collar covariance, mean, switches and
primary terms on the successive common domains.
The displayed load, rate and recursion estimates
give its quantitative reduction and the resulting
depth-independent bound. Nonvanishing is supplied
by the connected expansion under that criterion,
rather than by boundedness of a complex phase.

Healing changes the domains
in a further way. It erases
the mark of the remade levels
and returns their analytic
charge. The new space can
therefore be larger than the
space of an old stage term.
Moreover a unit pure gauge
near a node has zero action
but unit potential, so the old
orbit cannot always be
reached by a small potential
shift. The following generated
field construction addresses
both issues.

\subsubsection*{Constructing a gauge for the generated configuration.}

A critical configuration is its own background over a finer
system of its averages. This identity, and its holomorphic
continuation, make a generated configuration suitable for
the old tower clauses. In the middle zone of \(C\), use
the field \(\mathbf G_C\) generated on \(C\) itself from
the new datum and the inner variables. At a finer index
\(n\), its curvature gains \(L^{-2(k-n)}\), whereas its
potential gains only \(L^{-(k-n)}\).
The generating domain cannot be enlarged arbitrarily
into old layers: a fixed overshoot of a potential there
would exceed the available layer-dependent room.

There is a separate difficulty at the top index. The
target has only \(c\beta_{\mathrm s}\alpha_{1,k}\)
additional room. The required estimate is
\begin{equation}\label{eq:xrg-generated-gauge-bound}
 \max\{|A^{\mathrm g}|,|\nabla_UA^{\mathrm g}|\}
          <(1+c\beta_{\mathrm s})\alpha_{1,k},\qquad
 A^{\mathrm g}=\tfrac1{i\eta_k}
                 \log(\mathbf G_C^{\mathsf k_C}U^{-1}).
\end{equation}
A bound with an unspecified multiple of \(\alpha_{1,k}\)
would not put the substituted triple in the old domain.
The construction therefore arranges its leading term
as an average of the argument's potential \(A'\).

\emph{Real local frames and the reduced background.}
Work first in level-\(k\) units. For a coarse site \(y\)
in the middle zone, let \(\mathsf Q_y,\mathsf Q_y^+\)
be the unions of the \(3^4,5^4\) surrounding blocks.
They lie in the top layer, away from both boundaries.
For a real field \(W\) define a frame with bound
\(\varrho\) to be a real gauge \(f\), equal to one
at \(y\), for which
\[
 W^f=e^{i\eta_ka_f},\qquad
 \max_{\mathsf Q_y^+}(|a_f|,|\nabla a_f|)<\varrho .
\]
The regular-gauge bound makes this open set of frames
nonempty for
\(\varrho_{\rm f}=C(M)\alpha_{0,k}\).
Average their values with normalized Haar measure
and normalize in the quaternion algebra:
\[
 \mathsf f[W](x)=\mathcal N(\langle f(x)\rangle_W),
 \qquad \mathcal N(m)=(\det m)^{-1/2}m .
\]
When the box diameter times \(\varrho_{\rm f}\) is
small, the mean is nonzero. The frame estimate gives
\[
 W^{\mathsf f[W]}=e^{i\eta_k a},\qquad
 \|a\|_{\mathsf Q_y^+}\le2\varrho_{\rm f},\qquad
 \mathsf f[W^u]=u(y)\mathsf f[W]u^{-1}.
\]
Here \(\|a\|\) includes one difference. The frame
depends on the real field alone, so its formation
does not interfere with holomorphy in \(A'\).

Use the reduced determining sequence on
\(C\setminus\Lambda_C\), descending to the bare
scale at both boundaries. Its real datum is the
average of \(U\); call its background \(U^\natural\)
and its continuation \(\Phi^\natural(B)\).
Set
\[
 \mathsf f_y=\mathsf f[U|_{\mathsf Q_y^+}],\qquad
 \mathsf f_y^\circ=\mathsf f[U^\natural|_{\mathsf Q_y^+}].
\]
The construction is made on a product domain in
\((A',B)\): these two analytic variables are
independent before the actual datum is substituted.
The increment \(B\) is \(O(\alpha_1)\) on top rows
and lies in a fixed small polydisc on the other rows.

\emph{Removing the node part locally.}
Fix a large constant radius \(\mathsf R_{\rm g}\).
For coarse sites \(y'\) in that box use combs
from \(y\) to \(y'\) and set
\[
 v_y(y')=
       V^\natural(\mathrm{cb}_{y'})^{-1}
                 (e^{iB}V^\natural)(\mathrm{cb}_{y'}).
\]
Take \(v_y=1\) at the other nodes. Its lift
\(\widetilde v_y\) at a fine site is the value of
\(v_y\) at the first node reached by ascending through
the block centres. The transformed
increment \(B_y\) is defined by
\[
 e^{iB_y(c)}=
 v_y(c_-)e^{iB(c)}V^\natural(c)
                  v_y(c_+)^{-1}V^\natural(c)^{-1}.
\]
Let \(\mathcal H_y,u_y\) be its Landau response and
node gauge. Removing the comb part leaves a response
controlled by the curvature of the coarse datum
and by the exponentially weak influence of rows
outside this fixed box.

To restore the desired leading potential, put
\(A'_y=\operatorname{Ad}(\mathsf f_y)A'\).
Average its bare translates over a small fixed
distance \(\delta_{\rm g}\), obtaining
\(a_y=\mathsf S_{\rm g}A'_y\), and let \(\phi_y\)
be its comb integral from \(y\). If \(P_\nu(x)\)
replaces the first \(\nu\) coordinates of \(x\)
by those of \(y\), and \(I_\nu\) is the signed
sum along that comb segment, then
\[
 \phi_y=\sum_\nu I_\nu[(a_y)_\nu],\qquad
 \nabla_\mu\phi_y=(a_y)_\mu-r_\mu,\qquad
 r_\mu=\sum_{\nu<\mu}I_\nu[(da_y)_{\nu\mu}].
\]
Thus the failure to be a primitive is controlled
by \(da_y\). Mollification gives
\[
 \|\mathsf S_{\rm g}X\|\le\|X\|,\qquad
 \sup|\mathsf S_{\rm g}X-X|\le C\delta_{\rm g}\|X\|,
 \qquad
 \sup|\nabla\mathsf S_{\rm g}F|
                           \le2\delta_{\rm g}^{-1}\sup|F|.
\]
These formulas retain control of a difference of
the primitive defect without demanding an extra
derivative of the original potential.

With \(n_y=e^{-i\phi_y}\), the local gauge is
\[
 \widehat{\mathsf k}_y
   =\mathsf f_y^{-1}n_yu_y(y)\mathsf f_y^\circ
                                  u_y^{-1}\widetilde v_y .
\]
From right to left, its factors remove the node
comb, pass to the Landau form, put its real
background near one, normalize at \(y\), insert
the primitive, and return to the gauge of \(U\).
All factors are holomorphic in the independent
complex variables. They are real gauges at
real variables, and
\(\widehat{\mathsf k}_y(y)=1\).

\emph{The quantitative local defects.}
The estimates needed from this construction can be
written with four sizes. Let \(\varrho_{\rm a}\)
bound \(A'\) and its first difference,
\(\varrho_{\rm d}\) bound \(\mathrm d_UA'\),
\(\varrho_{\rm p}\) bound the coarse plaquette
increment in the node box, and
\(\varrho_{\rm e}\) bound the weighted influence
of the lower rows. Set
\begin{align*}
 \varrho_{\rm s}={}&
 C\mathsf R_{\rm g}(\varrho_{\rm p}+\alpha_0)
 +C\mathsf R_{\rm g}e^{-\delta_{\mathrm G}\mathsf R_{\rm g}/2}
                                            \varrho_{\rm a}
 +\varrho_{\rm e},\\
 \varrho_{\rm E}={}&
 \frac C{\delta_{\rm g}}
       (\varrho_{\rm d}+\varrho_{\rm f}\varrho_{\rm a})
 +\frac C{\delta_{\rm g}^2}\varrho_{\rm a}^2
 +C(\varrho_{\rm s}+\varrho_{\rm f}),\\
 \Delta={}&C\delta_{\rm g}\varrho_{\rm a}+\varrho_{\rm E}.
\end{align*}
For
\[
 A^{(y)}=\tfrac1{i\eta_k}
           \log(\Phi^\natural(B)^{\widehat{\mathsf k}_y}U^{-1}),
\]
the local estimate has the specific leading term
\begin{equation}\label{eq:generated-local-defect}
 \operatorname{Ad}(\mathsf f_y)A^{(y)}
       =a_y+\mathsf E_y,\qquad
 \|a_y\|<(1+C\varrho_{\rm f})\varrho_{\rm a},
 \quad \|\mathsf E_y\|\le\varrho_{\rm E},
 \quad \sup|A^{(y)}-A'|\le\Delta .
\end{equation}
The first two terms of \(\varrho_{\rm E}\) are
the primitive defect and its nonlinear correction;
the remaining terms are the reduced response and
the small real gauges. This is the quantitative
local estimate to establish, together with its
node identity, rather than an assertion that
some gauge has a bounded potential.

\emph{Joining the gauges without losing the leading constant.}
Use the \(C^{1,1}\) partition of unity
\[
 \psi_y(x)=\prod_{\mu=1}^4\psi(x_\mu-y_\mu),\qquad
 \psi(s)=
 \begin{cases}
 1-3s^2+2|s|^3,&|s|\le1,\\
 0,&|s|>1.
 \end{cases}
\]
It has bounded first and second differences,
\(\sum_y\psi_y=1\), \(\sum_y\nabla\psi_y=0\), and
\(\psi_y(y')=\delta_{yy'}\) at the coarse sites.
Set
\begin{equation}\label{eq:xrg-quaternion-gauge}
 \mathsf k_C(x)=\mathcal N\left(
                         \sum_y\psi_y(x)\widehat{\mathsf k}_y(x)\right).
\end{equation}
On the real quaternion span,
\(m^*m=(\det m)1\) and \(\mathcal N(m)=m/|m|\).
It therefore takes real combinations back into
\(\SU(2)\). Near one local gauge the argument
is that gauge times \(1+\) a small matrix, so
the branch of \((\det m)^{-1/2}\) is fixed
by its value one. The same formula is
holomorphic at complex arguments and is
bi-covariant under determinant-one factors.

The local joining calculation gives
\begin{equation}\label{eq:generated-joined-defect}
 \|A^{\mathrm g}\|_U
       <\varrho_{\rm a}
                   +K_{\rm j}(C\delta_{\rm g}\varrho_{\rm a}
                                                   +\varrho_{\rm E}).
\end{equation}
Its leading part is the average of the \(A^{(y)}\).
In its gradient, the differentiated weights have
sum zero, so only differences of neighbouring
leading potentials remain. These are bounded by
\(\Delta\), rather than by
\(\varrho_{\rm a}\) once for every box.
Normalization produces quadratic errors; an
explicit factor \(\eta_k\) costs no new difference
because \(\|\eta_kX\|\le2\sup|X|\).
The node identity of the local gauges and
\(B=\mathcal Q_k(U;A')\) on the top rows give
\(\mathsf k_C=1\) at the coarse sites.
The construction thus preserves the attached
stored data there.

For the actual generated datum, write \(\mathcal B^{\rm o}\)
for the exterior datum increment on the tower rows of the
reduced determining sequence, and set
\(\alpha_i=gC_i^{\rm rad}(\log x)^q\), \(i=0,1\),
at the scale in question. Then
\[
 \varrho_{\rm a}=\alpha_1,\quad
 \varrho_{\rm d}\le C(\alpha_0+\alpha_1^2),\quad
 \varrho_{\rm p}\le C\alpha_0,\quad
 \varrho_{\rm e}\le
 C[e^{-\delta_{\mathrm G}M}\sup|\mathcal B^{\rm o}|
                      +e^{-c\delta_{\mathrm G}M\check R_k}].
\]
The boundary datum is restricted by
\(Ce^{-\delta_{\mathrm G}M}\sup|\mathcal B^{\rm o}|
\le c\beta_{\mathrm s}\alpha_0\).
Choose constants in their operative order:
first \(K_{\rm j}\); then
\(\delta_{\rm g}\) with
\(K_{\rm j}\delta_{\rm g}\ll\beta_{\mathrm s}\);
then \(\mathsf R_{\rm g}\) with
\(K_{\rm j}\mathsf R_{\rm g}
 e^{-\delta_{\mathrm G}\mathsf R_{\rm g}/2}
\ll\beta_{\mathrm s}\); then \(M\) larger than this
node box; then the ratio \(C_1^{\rm rad}/C_0^{\rm rad}\);
and finally the coupling ceiling.
The ratio of the radii pays every fixed multiple
of \(\alpha_0\). The ceiling pays
\(C\alpha_1^2\) and the distant inner response.
Inserting these choices in
\eqref{eq:generated-local-defect} and
\eqref{eq:generated-joined-defect} yields
\eqref{eq:xrg-generated-gauge-bound}.
The quantitative local and joining estimates,
including their covariance and analyticity,
are the analytic content needed for this last
deduction; the quaternion normalization alone
does not supply them.

The generated field also
supplies the complex
regularity needed to sum
old marginal terms. A
generic pair in the full
analytic domain has no
H\"older control of adjacent
curvatures. It gives only
the exponent zero in the
following local bound.
For a pointed family $G$
of level $j=k-\nu$ with
coefficient $e$, suppose on
a window of $l$ cube layers
the pair has a potential
regular at scale $L^\nu$
with amplitude $a_{\mathrm p}$
and H\"older exponent
$\beta'>0$. The required
derivative bounds are
\[
 |\mathbf A|\le a_{\mathrm p}L^{-\nu},\quad
 |\nabla\mathbf A|\le a_{\mathrm p}L^{-2\nu},\quad
 [\nabla\mathbf A]_{\beta'}
        \le a_{\mathrm p}L^{-(2+\beta')\nu},\quad
 |\mathbf J|\le a_{\mathrm p}L^{-3\nu}.
\]
Here the current is still
independent. For the sum of
domains through $z$ that
stay inside that window and
inside $Y$, the pointed
subtraction gives
\begin{equation}\label{eq:xrg-complex-pointed-window}
 C\|G\|(a_{\mathrm p}/\alpha_0)^2L^{-4\nu}
  \left[L^{-\beta'\nu}
       +e^{-\kappa D_Y/2}+e^{-\kappa(l-1)/2}\right],
\end{equation}
where $D_Y$ is the shortest
domain length from $z$ to
$Y^c$. The $L^{-4\nu}$
pays the number of base
points; only the positive
H\"older exponent leaves
age decay. Domains leaving
the window are paid
separately by
$e^{-c\kappa L^\nu}$.
Using the full analytic
domain without a generated
configuration would set
$\beta'=0$ and let the
scale sum grow with $k$.
This is why the extension
used in the component
integrals is chosen through
the generated field.

\subsection{Keeping stored gauges and passing between splittings}

A stored fluctuation is a real integration variable in the axial
gauge in which it was created. A later response matches new
averages only up to a node factor \(v\). Acting on the stored
variable by \(v^{-1}\) would move its real integration set when
the response parameters are complex. We instead compensate the
triple, prescribing the gauge at both the nodes and the sites
where stored data are read.

\emph{Stored sites and the prescribed values.}
Let \(y\) be an inner top of level \(m\ge1\), with pillar
\(\Pi(y)=B^m(y)\). The \(L^4\) sites of level \(m-1\) in that
pillar form its set of stored sites \(\mathsf A(y)\). Their
blocks \(\Pi'(a)=B^{m-1}(a)\) partition the pillar; for \(m=1\)
they are single sites. Write
\[
 \tau[\mathbf U](a)=(\mathcal M^{m-1}\mathbf U)(y,a),
 \qquad \tau[\mathbf U](y)=1.
\]
This is the group average of the shortest coordinate-path
transports from \(y\) to \(a\). Its transformation law is
\[
 \tau[\mathbf U^g](a)
       =g(y)\tau[\mathbf U](a)g(a)^{-1}.
\]
The condition \(\tau[U](a)=1\) says that the corresponding
averaged real field is in the block axial gauge in which the
level-\(m-1\) fluctuation was stored.

For a shift \((\mathbf A,\mathbf K)\), let
\(v[P;\mathbf A](y)\) be its accumulated node factor. Its smallness
is measured by
\(\varepsilon_y(\mathbf A)=\sup_{\Pi^*(y)}\xi_m|\mathbf A|\),
where \(\xi_m\) is the length of a level-\(m\) block in the
current units and \(\Pi^*(y)\) is the determining neighbourhood
of that pillar. The needed node-factor estimate is
\[
 |v^{\pm1}-1|\le c_{\rm r}\varepsilon_y(\mathbf A).
\]
At a stored site define the pin
\begin{equation}\label{eq:xrg-pinned-gauge}
 \begin{split}
 u^\delta(y)&=v[P;\mathbf A](y)^{-1},\\
 u^\delta(a)=\mathsf p[P;\mathbf A](a)
   &:=\tau[\mathbf U](a)^{-1}v[P;\mathbf A](y)^{-1}
                          \tau[e^{i\eta\mathbf A}\mathbf U](a).
 \end{split}
\end{equation}
In particular \(\mathsf p(y)=v(y)^{-1}\). The pins depend on
the triple only through its complex configuration on
\(\Pi^*(y)\). Their required bound is
\[
 |\mathsf p^{\pm1}-1|\le c_{\rm m}\varepsilon_y(\mathbf A),
                         \qquad c_{\rm m}=c_{\rm m}(L).
\]
They are one at zero shift and are conjugated by \(g(a)\) under
a real gauge transformation. On the domain of the complex
averages and their logarithms, they are holomorphic in the
complex configuration and shift. Continuity in the real factor
is required as well.

The depth-one geometry of stored bonds is important here. An
endpoint of a stored bond is either a top itself or one of the
sites of \(\mathsf A(y)\) one level below its top. The pillar,
transport and axial-gauge condition agree in all the stage and
main-step structures while that field remains stored. Restricting
to a finer determining sequence can lower the endpoint's layer
index; it does not turn it into a site deeper than one level
below a top. Thus these prescribed values cover the sites at
which the stored variables are read.

\emph{Extending the pins inside their cells.}
Choose the tops whose stored sites are to be pinned. Every inner
top of level at least two is eligible. A top of level one is
eligible only when the bare sites within \(2\eta\) of its pillar
lie in pillars of level at most one. For a pinned pillar the
turning cells are its blocks \(\Pi'(a)\), centered at the stored
sites. For every remaining inner pillar there is a single
turning cell, centered at its top.

On a turning cell \(B^q(a)\), take the covariant frame \(T_a[U]\)
associated with its local regular-gauge class. The frame is
normalised at \(a\) and transports a Lie-algebra value at \(a\)
to a Lie-algebra value at \(x\). Its required covariance is
\[
 T_a[U^g](x)
   =\operatorname{Ad}(g(x))T_a[U](x)
                              \operatorname{Ad}(g(a))^{-1}.
\]
It is this regular frame, with controlled differences, that enters
the construction. Transport along a chosen coordinate path does
not in general have the second-difference bound needed to control
the derivative of the turned potential.

Fix an even \(C^2\) bump equal to one at zero and supported in
\([-1/3,1/3]\). Its four-coordinate product, at the scale of the
cell, is denoted by \(\mathrm{bp}_a(x)\). For independent complex
parameters \(\zeta(a)\), put
\begin{equation}\label{eq:xrg-turning-cells}
 \mu_\zeta(x)=\mathrm{bp}_a(x)T_a[U](x)\zeta(a),
       \qquad u_\zeta=e^{-i\mu_\zeta}
\end{equation}
on that cell, and put \(\mu_\zeta=0\) elsewhere. On a cell
consisting of one site use \(\mu_\zeta(a)=\zeta(a)\).
The inner thirds are disjoint. The turn is therefore defined
without choosing an order for its factors.

The pinned parameters are
\[
 \zeta^{\rm pin}(a)=\frac1i\log\mathsf p(a)^{-1},
          \qquad u^\delta=u_{\zeta^{\rm pin}}.
\]
At an unpinned top this reads
\(\zeta^{\rm pin}(y)=\frac1i\log v(y)\).
We use these formulas on a common open small-shift domain where
all the tops involved are levelled: their local shift size is
strictly below the fixed logarithm threshold. At a pinned
level-one top the size is measured on the neighbourhood extending
\(2\eta\) beyond its pillar. Unused pillars have the identity
turn. Analyticity and the Lipschitz estimates below are used with
a margin inside this domain; they do not assert continuity of a
rule that switches abruptly to the identity at the levelling
threshold.

Define the turned shift by
\begin{equation}\label{eq:xrg-turned-shift}
 e^{i\eta\mathbf A^\zeta}\mathbf U
           =(e^{i\eta\mathbf A}\mathbf U)^{u_\zeta},
 \qquad
 \mathbf J+\mathbf K^\zeta
           =\operatorname{Ad}(u_\zeta)(\mathbf J+\mathbf K).
\end{equation}
The first identity defines \(\mathbf A^\zeta\) by the local
logarithm, keeping the original real factor fixed.
On a cell of scale \(\xi_q\), the triple must satisfy the local
regular-gauge condition together with
\(\xi_q|A'|,\xi_q^2|\nabla_UA'|<c\). The local estimates to verify are
\begin{align*}
 \xi_q|\mathbf A^\zeta|
   &\le2\xi_q|\mathbf A|+C|\zeta|,\\
 \xi_q^2|\nabla_{\mathbf U}\mathbf A^\zeta|
   &\le4\xi_q^2|\nabla_{\mathbf U}\mathbf A|
                       +C(\varepsilon+|\zeta|),\\
 \xi_q^3|\mathbf K^\zeta|
   &\le(1+C|\zeta|)\xi_q^3|\mathbf K|
                                      +C|\zeta|\mathsf j_0.
\end{align*}
Here \(\xi_q|\mathbf A|\le\varepsilon\), and
\(\xi_q^3|\mathbf J|\le\mathsf j_0\) is the local current
bound. For a one-site cell the same bounds are used at scale
\(\eta\), without a smoothness requirement on its site value.

The matrix calculation behind the potential estimate is the
local map
\[
 \mathsf F(w,X,Z)=\log(e^{-iw}e^Xe^{iw+Z}).
\]
For small arguments it is holomorphic, vanishes when \(X=Z=0\),
and satisfies
\[
 |\mathsf F|\le2(|X|+|Z|),\qquad
 |\Delta\mathsf F|
 \le C\max(|X|+|Z|)|\Delta w|
                         +4(|\Delta X|+|\Delta Z|).
\]
Applied bond by bond, the covariant differences of the bump and
frame give the first two estimates. The current estimate follows
from the second identity in \eqref{eq:xrg-turned-shift}. The
uniform regular-frame and local-logarithm estimates are the
analytic inputs to this calculation; prescribing the values at
the centers alone would not give the derivative bound.

Let \(\|\cdot\|_y\) denote the local scaled shift norm controlling
these potential, first-difference and current quantities. Let
\(\operatorname{nb}(y)\) consist of \(y\) and the pinned
level-one pillars within \(2\eta\) of its pillar, and set
\(\|\cdot\|_y^+=\max_{y'\in\operatorname{nb}(y)}
\|\cdot\|_{y'}\). For a top of level at least two this
neighbourhood consists just of that top. Passing from a pinned
pillar to its turning cells costs a fixed change of scale by
\(L\). The bound needed for the pinned shift is consequently
\begin{equation}\label{eq:xrg-pinned-shift-bound}
 \|(\mathbf A^\delta,\mathbf K^\delta)\|_y
   \le C_{\rm m}\|(\mathbf A,\mathbf K)\|_y^+,
                 \qquad C_{\rm m}=O(L+L^2c_{\rm m}).
\end{equation}
It is required uniformly in the label, level and number of
pillars. The same estimate holds in Lipschitz form on a half
ball in the common small-shift domain. It also bounds every free
turn with \( |\zeta(a)|\le2c_{\rm m}\varepsilon'_{y(a)}\),
using the larger of the shift size and the local bounds
\(\varepsilon'\) on the right. For an unpinned pillar the
corresponding constant is \(O(1+c_{\rm r})\).

\emph{Locality, domain room and preservation of the stored gauge.}
The turn at a bond over \(y\) reads only the pillars in
\(\operatorname{nb}(y)\). It is the identity outside the inner
thirds of the turning cells, and on every pillar where the shift
is zero. In particular, shifts whose nonzero local
neighbourhoods are disjoint are compensated separately. For
fixed \(U\), the construction is holomorphic in the complex
potential, current, shift and free turn parameters on their
common domain. It is real at real data, commutes with conjugation,
and is carried along by translations and the lattice isometries
preserving the pillar geometry. These statements use the same
properties of the node factors, local frames and logarithms.

The local domain test uses the bound for every free turn, before
the pinned point is substituted. If \(\mathsf s'_y\) denotes
the available room expressed in the local shift norm, require
\[
 C_{\rm m}\|(\mathbf A,\mathbf K)\|_y^+\le\mathsf s'_y
\]
at each pillar meeting the target. Then the compensated triple
and all the free turns in that polydisc remain in the old term's
domain. For a slowly varying room profile the maximum over
\(\operatorname{nb}(y)\) costs only \(e^{C\delta}\), where
\(\delta\) is its variation rate. The enlarged test that keeps
a level-one pinned top levelled costs an additional factor \(L\):
its size is bounded by \(Le^{C\delta}\) times the local room
profile. This test must stay strictly below its logarithm
threshold throughout the parameter polydisc. Checking the room inequality in
the pointwise, gauge, tower and core clauses is still necessary;
it is not a claim that arbitrary complex gauge copies belong
to the same domain.

Finally the prescribed values give the exact identity
\[
 \begin{split}
 \tau[(e^{i\eta\mathbf A}\mathbf U)^{u^\delta}](a)
 &=v(y)^{-1}\tau[e^{i\eta\mathbf A}\mathbf U](a)
                                              \mathsf p(a)^{-1}\\
 &=\tau[\mathbf U](a).
 \end{split}
\]
At an adapted real representative the last transport is one.
The stored field is therefore read in its original gauge and its
real coordinate has not moved. Although the frame extending the
turn can depend on the real splitting, the values prescribed at
tops and pinned stored sites depend only on the pair and the
shift. Equality of the interpolated terms for two splittings
uses the common holomorphic family below, together with this
agreement at the sites carrying data.


Some terms are first formed on an
axial slice. If $\mathfrak u[\mathbf U]$
is the holomorphic gauge carrying
the pair to that slice, their extension
has the form
\[
 F^\natural(\mathfrak e)
                   =F(\mathfrak u\star\mathfrak e).
\]
In $\mathfrak u\star\mathfrak e$,
the complex pair is transformed by
$\mathfrak u[\mathbf U]$, but the
real factor and stored data use the
real gauge $\mathfrak u[U]$.
The extension is defined only where
this transformed point belongs to
the original analytic domain.
It need not contain a fixed complex
neighbourhood of every real field.

For example, a constant imaginary
pure-gauge potential of size $\alpha$
can have its contour transport over
the entire tower concentrated on
one bond by the axial compensator.
Its resulting potential there has
size $\alpha L^k$. Along this
direction the allowed radius can
be as small as
$C\alpha_{1,k}L^{-k}$.
Thus a full complex orbit cannot
be used to argue that every gauge
representative is in the same
analytic domain. The component
construction uses its explicit
pinned points and domain margins.

For the noncollar terms, equality
at two splittings of one pair is
proved along a common holomorphic
family. Take one splitting and write
\[
 A'_z=A'_{\mathrm R}+zA'_{\mathrm I},\qquad
 \mathbf J_z=\mathbf J_{\mathrm R}
                                   +z\mathbf J_{\mathrm I}.
\]
For $z=it$, $0\le t\le1$, the
current is
$\frac{1+t}2\mathbf J+
 \frac{1-t}2\mathbf J^*$.
It stays in its convex current
ball. In the other splitting choose
the potential holomorphically so
that its real factor gives this
same pair. At real $z$ both
potentials are real.

The substituted-triple rule is
first formed for free gauge-turn
parameters $\zeta$ in a polydisc;
the actual rule evaluates them at
the pinned point
$\zeta^{\mathrm{pin}}$.
At real $z$ and real $\zeta$ the
two substituted pairs differ by
a real gauge function equal to
one on the datum. The old term,
already a function of its pair
and invariant under real gauges,
takes the same value on them.
The constructions remain defined
on radii increased by a fixed
factor and land inside the true
domains. This margin gives an
open complex neighbourhood of the
chord and of the parameter disc.
The identity theorem first in
$z$ and then in $\zeta$ proves
equality at $z=i$ and at the
pinned parameters. This argument
does not require the entire
analytic domain to be connected;
it uses this specified common
family. The technical point is to
verify the common-domain margins
along this family; the identity
theorem then gives the equality.

At an evaluation point, with
real stored variables,
$\lambda=1$ and decoupling
parameters equal to indicator
functions, the substituted pair
is the pair of the old regular
critical field. Real gauge
invariance and the just established
splitting identity therefore
identify the localised formulas
with the original integrals.
For collars, the construction
keeps the explicit triple and
its stored-data integral as a
single operation; a global
complex-gauge extension is not
substituted for it.

\subsection{Extracting the factors for failed tests}

The gain assigned to a failed test must come from the weight that
is still present at that place in the history. In particular, a
plaquette test is paid for by every real configuration with the
recorded averages, not only by the regular background used to
define the test. At a creation of step $t$, the paying Wilson
coefficient is $x_{t-1}$. During a later healing, the old action
already spent on the component must be excluded from the new
payment.

\subsubsection*{Energy of the averages and the three tests.}

For a set \(S\) of level-\(t\) plaquettes, let \(\mathcal N_0(S)\)
be the union, over \(p'\in S\), of the bare plaquettes whose
corners lie in the blocks beneath the corners of that single
plaquette \(p'\). The estimate is
\begin{equation}\label{eq:energy-of-averages}
 \sum_{p\in\mathcal N_0(S)}a_p(U)
 \ge\min\left\{
   \tfrac14\sum_{p'\in S}|\partial\mathcal M^tU(p')-1|^2,
   \tfrac12e_*\right\},\qquad e_*=e_*(L)>0.
\end{equation}
It applies to every real \(U\), without a regularity assumption.
In particular its use under a constrained integral does not restrict
the paying configuration to the regular minimizer.

We give the nonlinear iteration. In the comb gauge the derivative
of one averaging step is the straight average \(Q\), and its
plaquette coboundary is given by the discrete Stokes formula
\[
 (Qa)(\partial p')
   =L^{-4}\sum_{z\in B(y_1)}a(\partial S_z).
\]
Here \(y_1\) is one corner of \(p'\), and \(S_z\) is the
\(L\)-by-\(L\) plaquette square starting at \(z\).
Taking absolute values gives a nonnegative plaquette matrix \(P\).
Each square has \(L^2\) plaquettes, so a row sums to \(L^2\).
For a fixed fine plaquette, the starts in its two plane directions
have \(L^2\) possibilities in total across the adjacent coarse
blocks; the transverse coordinates are fixed. Its column therefore
sums to \(L^2/L^4=L^{-2}\). Consequently
\[
 \|P\|_{2\to2}\le1,\qquad \|P\|_{1\to1}\le L^{-2}.
\]
All these norms use counting measure at the respective levels.

Put \(f_l=|\partial\mathcal M^lU-1|\). Whenever the plaquettes
in the four corner blocks are sufficiently small, the comb-gauge
formula and analyticity of \(\mathcal M\) give
\[
 f_{l+1}\le P_lf_l+q_l,\qquad q_l\ge0,\qquad
 \|q_l\|_1\le C(L)\|f_l\|_2^2.
\]
The correction is quadratic because the displayed Stokes average
is the linear term. Its single-plaquette remainder is bounded by
the square of the largest curvature in those blocks; replacing
that square by the sum of squared curvatures and summing uses only
a bounded number of block overlaps.

These estimates can be restricted to the ancestors of \(S\).
Starting from \(S\), include at each preceding level all plaquettes
read by the comb-and-square formula for the next set. The bare
ancestor set is contained in \(\mathcal N_0(S)\): every path and
square stays in the blocks below the corners in question.
Restricting \(P_l\) discards nonnegative entries, so the same row,
column and operator bounds hold. Let \(E_l=\|f_l\|_2^2\) on these
ancestor sets. For a product of \(m\) restricted matrices,
\[
 \|P_{l+m-1}\cdots P_l v\|_2
 \le \|P_{l+m-1}\cdots P_l v\|_1
 \le L^{-2m}\|v\|_1.
\]
Iteration of the nonlinear inequality therefore yields
\begin{equation}\label{eq:average-energy-duhamel}
 E_t^{1/2}\le E_0^{1/2}
       +C(L)\sum_{l=0}^{t-1}L^{-2(t-1-l)}E_l.
\end{equation}
The linear part is marginal in dimension four. The nonlinear
corrections have an \(\ell^1\) norm, and propagation of that norm
is strictly contracting.

Choose \(e_*\) so small that
\[
 \frac{2C(L)}{1-L^{-2}}\sqrt{e_*}\le\sqrt2-1
\]
and that \(\sqrt{2e_*}\) lies inside every local small-curvature
chart used above. Suppose \(E_0\le e_*\).
If \(E_l\le2E_0\) at the preceding levels,
\eqref{eq:average-energy-duhamel} gives
\[
 E_t^{1/2}\le E_0^{1/2}
                  +\frac{2C(L)}{1-L^{-2}}E_0
             \le \sqrt{2E_0}.
\]
Starting at \(l=0\), this proves \(E_l\le2E_0\) for every
level, and simultaneously justifies all the small-curvature
formulas used in the iteration. The constants do not grow with \(t\).

Finally \(E_0=2\sum a_p(U)\) on the bare ancestor set. If the
action on \(\mathcal N_0(S)\) is below \(e_*/2\), the preceding
bound gives
\[
 \sum_{p'\in S}|\partial\mathcal M^tU(p')-1|^2
   \le 2E_0
   \le4\sum_{p\in\mathcal N_0(S)}a_p(U).
\]
If that action is at least \(e_*/2\), the other alternative in
\eqref{eq:energy-of-averages} applies. This proves the estimate
on the full real constraint set, including configurations which
were not assumed regular at the outset.

Each payment is most conveniently checked in contrapositive form.
If its paying quantity is below both the required threshold cost
and a fixed small constant, the data must be regular, the
localised orbit and contour variables must exist, and the test
must pass. This formulation also treats the complementary outcome
where a tested object is undefined.

For a plaquette test, failure at a test cube of step $t$ requires
a nearby coarse plaquette with
$|\partial V_t-1|\ge c\varepsilon_t$. Otherwise the local
background estimate places the orbit inside the plaquette test.
Together with \eqref{eq:energy-of-averages}, the required
contrapositive gives paying action at least
$c\varepsilon_t^2$ for every real field with these averages.

For the approximate fluctuation test, the Wilson action must be
combined with the retained gauge-fixing term. Write
$\mathrm{GF}_j(Y;V_j)$ for the part of the nonnegative soft
gauge-fixing action located at the blocks in $Y$, so its weight
is $\exp[-x_j\mathrm{GF}_j(Y;V_j)]$. A gauge transformation
at a site of level $t-1$ which is not a level-$t$ site can change
$V_{t-1}$ relative to the axial representative of the localised
background without changing its action or $V_t$. Action alone
therefore cannot control that difference. The local implication
to establish is
\[
 \sum_{p\text{ near the test cube}}a_p(U)
       +\mathrm{GF}_{t-1}(\text{near the test cube};V_{t-1})
                            \ge c\delta_{t-1}^2
\]
on the failed approximate test, again for all compatible real
fields. The gauge fixing supplies the directions invisible to
the Wilson action.

For the true fluctuation, put
$\mathcal A_j^{\mathrm s}=g_j^{-1}\mathcal A_j$, where
$\mathcal A_j$ is the unscaled logarithmic fluctuation. Coercivity
on the constrained fluctuation space gives the payment
\begin{equation}\label{eq:failed-fluctuation-payment}
 \frac12\langle\mathcal A_{t-1}^{\mathrm s},
                 \Delta^{(t-1)}\mathcal A_{t-1}^{\mathrm s}\rangle
                     \ge c\frac{\delta_{t-1}^2}{g_{t-1}^2}.
\end{equation}
The corresponding stored-field estimate uses
\eqref{eq:kept-coercivity}, which must be proved for the
locally retained conditional form. Its validity is not inherited
from Schur positivity after arbitrary truncation. Collar and rim
failures have the same payment, with constants determined by the
fixed fractions $1/(4\Theta_1)$ and $\theta_c$ of their thresholds.

The contrapositive must also be checked after the change of
variables on the strip. There its shift vanishes when the
underlying unscaled integration variable is small. Consequently
an undefined tested logarithm cannot be counted as a failed
small-field test while that variable remains below its threshold.
The lower bound is then applied to the integration variable,
with its actual kept Gaussian form. This is the local existence
and coercivity check needed to include the undefined outcomes
in the same estimate.

The one-step flow bounds compare $x_{t-1}$ and $x_t$ by a fixed
additive constant, without requiring an individual increment to
be positive. Thus
\[
 x_{t-1}\varepsilon_t^2
       =\frac{x_{t-1}}{x_t}p_0(g_t)^2
       \asymp p_0(g_t)^2,
 \qquad
 x_{t-1}\delta_{t-1}^2
       =(A_1/A_0)^2p_0(g_{t-1})^2
       \asymp p_0(g_t)^2.
\]
The threshold ratios are fixed before the coupling ceiling.
Once the three local implications above have been established,
each created cell therefore supplies
$\exp[-cp_0(g_t)^2]$. The coefficient used for this conclusion
is the actual $x_{t-1}$ in the paying weight; comparison with
$x_t$ is only a subsequent estimate.

\subsubsection*{Disjoint shares of the action and the integration reserve.}

At a later level $k$, a piece $(t,D)$ consists of a component
of the region produced by step $t$ on which no subsequent
large-field operation has acted. Give it the bare plaquettes within one cell of its
failed plaquette or approximate-fluctuation test cells.
The buffers between new tests and the hull of an older region
make these shares disjoint. On the share of $(t,D)$ the spatial
coefficient is exactly $x_{t-1}$: the older profiles have
completed their transition and the newer profiles have not begun
it. Extending $U_k$ beyond its first determining region by the
stored bare field, the allocation is
\begin{equation}\label{eq:creation-action-shares}
 \begin{split}
 A(x_k(\cdot),U_k)&\ge A_{\mathrm{cr}}(U_k)\ge0,\\
 A_{\mathrm{cr}}(U_k)
       &:=\sum_{(t,D)}x_{t-1}
                    \sum_{p\in\text{share of }(t,D)}a_p(U_k).
 \end{split}
\end{equation}
All other plaquettes retain their nonnegative action. The
allocation uses one Wilson weight, together with disjoint sets
of retained gauge-fixing factors and Gaussian variables.

The fields entering the tests and the paying action are related
by the averaging constraints. On the recorded set $\Gamma_{t-1}$,
\[
 \mathcal M^{t-1}U_k=V_{t-1},\qquad
                         V_t=\mathcal M(V_{t-1}).
\]
The first identity is part of the background datum and the second
is imposed by the averaging delta function. Hence the payment
estimate for every real field with those averages applies to
$U_k$, although the test itself was formed with a localised
background of step $t$. An estimate only at that background
would not justify \eqref{eq:creation-action-shares} in the
component integrals.

Split each paying exponent into fixed halves: one half gives
the small factor and the other remains for integration. In
particular the coupling-dependent compact-group normaliser
satisfies
\[
 \int e^{-\frac12x(1-\frac12\operatorname{Re}\tr u)}\,du
 \le C\int_{|u-1|<\epsilon}
              e^{-x(1-\frac12\operatorname{Re}\tr u)}\,du,
 \qquad 0<\epsilon\le1,\quad x\epsilon^2\ge1.
\]
Both integrals have order $x^{-3/2}$. Thus using half an
exponent for payment leaves an integrable weight against its
own normaliser. The stored Gaussian and gauge-fixing factors
are treated with the same reserve. This allocation, followed by
the pointwise product majorant in
\eqref{eq:mrg-product-majorant}, is what permits gains at different
places and levels to be multiplied.

\subsubsection*{New failures during a healing.}

Inside an operation of level $m$, the tests use the sharper
thresholds $\delta'_j$. A tree gauge there has length at most
$C\check R_j^C$. A large bond in that gauge implies only that
some nearby plaquette exceeds
$c\check R_j^{-C}\delta'_j$. The resulting gain has the weaker
form
\[
 \exp[-c\check R_j^{-C}p_1(g_j)^2]
                           \le\exp[-p_0(g_j)].
\]
The second condition in \eqref{eq:threshold-exponents} pays
this power loss. It does not give another full creation deposit
of order $p_0^2$.

Put $m_\flat=m-N_m$. The old \emph{floor} is the union of the
level-$m_\flat$ cells in $\Omega_{m_\flat}^c\cap X$ that do
not touch $\Omega_{m_\flat}$. It is only the interior part of
that old core, not all of $X\cap Z_{m_\flat}$. The plaquettes
there, and the gauge-fixing and stored Gaussian factors at levels
below $m_\flat$, have already paid for the history preceding
the wait. They are unavailable for a new failure.

Suppose now that the operation has a failed preparation test and
leaves the drop-out $X$. Let $X^+$ be $X$ enlarged by two layers
of level-$m$ unit cubes, and consider real stored variables with
the recorded averages and test values. The old-core datum on
$\Gamma_l$ must have plaquette size at most
$C\check R_l^2\varepsilon_l$ in its own level units. On these
compatible evaluations the required fresh payment is
\begin{equation}\label{eq:fresh-failure-payment}
 \begin{split}
 &c x_m\!\!\sum_{\substack{p\subset X^+\\
                         \text{no corner in the floor}}}\!\!a_p(U)
       +\sum_{j=m_\flat}^{m-1}x_j\mathrm{GF}_j(X;V_j)
       +c\sum_{j=m_\flat}^{m-1}
                           \|\mathcal A_j^{\mathrm s}\|_X^2\\
 &\hspace{2cm}\ge(\log x_m)^{2\mathfrak p_1-C\mathfrak r}
                             \ge p_0(g_m).
 \end{split}
\end{equation}
Here $\|\mathcal A_j^{\mathrm s}\|_X^2$ is the sum of squared
Euclidean norms over the level-$j$ stored bonds in $X$ carrying
the retained Gaussian form, with one orientation per bond and
counting measure. The fields are expressed in their respective
level units. The Wilson sum excludes every plaquette with a corner
in the floor, and both other sums start at $m_\flat$. The
estimate must hold for every real $U$ compatible with the
stored variables on $X^+$. This is the uniform version needed
when a later containing drop-out is evaluated at those data.

For the contrapositive, let the fresh quantity on the left of
\eqref{eq:fresh-failure-payment} be small, and put
\[
 \varsigma=Cg_m(\log x_m)^{\mathfrak p_1-4\mathfrak r}
                   \asymp C\delta'_m\check R_m^{-4}.
\]
The smallness threshold has order
\[
 g_m^{-2}\varsigma^2
                   \asymp(\log x_m)^{2\mathfrak p_1-8\mathfrak r}.
\]
Apply \eqref{eq:energy-of-averages} to the stored plaquettes whose
bare ancestors lie outside the floor. Their plaquette deviations
are then bounded by \(C\varsigma\). The kept gauge fixing and
Gaussian coercivity control the remaining contour and fluctuation
coordinates. At each stage the datum used by the test is thus
regular on the side of the waiting collar where that test is made.
The inverse couplings remain comparable throughout this one wait,
since its length is \(O(\check R_m)=o(x_m)\).

The background estimate required at this point retains the distance
from the old floor. If \(\hat U\) is any of the stage backgrounds
and \(j(p)\) is the greatest layer index whose region contains a
corner of \(p\), it has the form
\begin{equation}\label{eq:mass-stage-screening}
 |\partial\hat U(p)-1|
 \le C\left(\varsigma+\check R_m^2\varepsilon_m
                          e^{-\delta_{\mathrm G}M\check R_m}\right)
                              L^{-2j(p)}
\end{equation}
on the region where the operation acts, outside the old core.
Local regularity supplies the first term; locality across the
waiting collar supplies the second. The collar has at least one
cell of width in the local units and separates the tested region
from the unspecified old field in the floor.
The exponential makes the second term smaller than \(\varsigma\)
at small coupling. This is a local background estimate: it does
not assert small action or small curvature in the excluded floor.

We spell out its use in the most demanding test, the main-step
test \(|\tfrac1i\log V'(b)|<\delta'_m\).
Let \(W_l=\mathcal M^lU\) be the averages of the compatible real
configuration and \(W_l^0=\mathcal M^l\tilde U_0\) those of the
basic background. On the shells the two fields have plaquette
deviations at most \(C\varsigma\), by the preceding estimate.
They agree on the imposed exterior datum and, in the selected
tree gauge, on the tree bonds. Close a bond \(b\) by a walk
\(\ell_b\) of tree and exterior bonds. The variable \(V'(b)\) is
conjugate to the quotient
\[
 W(\ell_b)\,W^0(\ell_b)^{-1}.
\]
For unitary holonomies,
\[
 |V'(b)-1|
 \le |W(\ell_b)-1|+|W^0(\ell_b)-1|.
\]
The shell-loop estimate applies to both terms. A loop made of
\(\mathsf n\) straight pieces in a shell of side
\(s\le CM\check R_m\) satisfies
\[
 \left|\prod_{b\in\gamma}W(b)-1\right|
 \le C\mathsf n s^2
                       \sup_{p\ {\rm in\ the\ shell}}|\partial W(p)-1|.
\]
The shell is a box with a strictly interior box removed; in at
least three dimensions this is the quantitative loop-filling
estimate used here. It is applied within the shells of the
chosen tree, rather than across the unknown floor. Consequently
\[
 \left|\tfrac1i\log V'(b)\right|
       \le CM^2\check R_m^2\varsigma
       \le C'M^2\check R_m^{-2}\delta'_m
       <\delta'_m
\]
after reducing the coupling ceiling. Two powers of \(\check R_m\)
pay for the loop area and the remaining two powers absorb its
fixed constants.

The earlier approximate tests use the one-step averaging rigidity,
and the fluctuation, collar and rim tests use the kept Gaussian
form. Thus the same chain is action of a compatible field,
regularity of its stored averages, local regularity of the test
background, and finally its contour or bond variable. Verification
of \eqref{eq:mass-stage-screening} for the actual stage backgrounds,
the stated shell-loop bound for their chosen trees, and the
local kept-form lower bound are the model-specific inputs to this
contrapositive. No gain is extracted from the main quadratic
form \(q_X\). The weakest scale obtained from the loop argument is
\[
 \check R_m^{-8}p_1(g_m)^2
       \asymp(\log x_m)^{2\mathfrak p_1-8\mathfrak r};
\]
the strict inequality
\(2\mathfrak p_1-8\mathfrak r>\mathfrak p_0\)
makes it larger than \(p_0(g_m)\). The number of places at which
the operation can fail costs only its logarithm and is paid by a
fixed part of this surplus.

The allocation separates nested operations. Every Wilson plaquette
used by an earlier operation underneath \(X\) lies in the floor,
and the gauge-fixing and Gaussian factors used underneath it have
levels below \(m_\flat\). The new Wilson sum has no corner in that
floor, and the new sums of forms begin at \(m_\flat\).
This is why a new failure cannot be paid by the old creation deposit.

For the mass of the drop-out \(X\) itself, the indicator of the
test that made \(X\) a drop-out is left unused. Its own-mass
estimate below is obtained from its earlier creations and the
carrying Wilson action on real compatible evaluations. The fresh
failure estimate is used when \(X\) is a mark inside a later
component, where the compatible field is defined on \(X^+\)
and the required fresh plaquettes and forms are available.

\subsection{The genealogy of a retained component}

Each unhealed component at one scale has a unique successor containing
its coarse hull. The predecessors of a fixed terminal component form
a tree. Write $n_t^{\mathrm{cr}}$ for its number of created cells at
level $t$ and $n_t^{\mathrm{mk}}$ for its number of marked components.
If $C(u)$ is the region at a node $u$, let $|C(u)|_{n(u)}$ count its
cells at that node's level. With a nonincreasing cost sequence
$\widehat r_t$, the required combinatorial estimate is
\begin{equation}\label{eq:genealogy-count}
 \sum_u\widehat r_{n(u)}|C(u)|_{n(u)}
 \le C\sum_t
       (n_t^{\mathrm{cr}}+n_t^{\mathrm{mk}})
                  N_t\widehat r_t.
\end{equation}
Increasing $N_t$ by a fixed additive constant is understood here.
The bound is independent of the depth of the genealogy.

The count separates large components from stretches with only
boundedly many cells. Coarsening contracts the occupied-cell skeleton
of a large component over a fixed number of steps. A light component
can persist without an event for only $O(N_t)$ steps, because it then
enters the class on which the large-field operation acts. The events
are creations, successors of marks, and joinings. A joining is assigned
to a leaf below one of its nonpreferred predecessors. Fixing a first
predecessor at each branching makes this assignment injective: the
first departure from that predecessor on the upward path identifies
the branching. Every leaf contains a creation. Thus joinings do not
require a new, unpaid gain.

For the analytic estimates take
\[
 \widehat r_t=CM^4\check R_t^5,
 \qquad N_t\widehat r_t\le CM^4\check R_t^6.
\]
The four powers from cell volume, the logarithmic normalisation per
site, and the possible waiting time are all included. Choosing
$\mathfrak p_0>6\mathfrak r$ ensures that even a mark's weaker gain
$p_0(g_t)$ exceeds this cost. More explicitly, the excess of gains
over the summed local costs has the lower bound
\begin{equation}\label{eq:genealogy-gain}
 \sum_t cp_0(g_t)^2n_t^{\mathrm{cr}}
 +\sum_t p_0(g_t)n_t^{\mathrm{mk}}
 -\sum_u C\log x_{n(u)}(M\check R_{n(u)})^4|C(u)|_{n(u)}
 \ge 2\sum_t p_0(g_t)n_t^{\mathrm{cr}}.
\end{equation}
The constants in this inequality must be chosen once for the whole
run. It is the quantitative margin used both to sum histories and
to retain decay in their age.

The combinatorial estimate separates a geometric input
from the ensuing summation. Let $o_n$ be the occupied
cell count in the large part of the genealogy at level
$n$, and let $l_n$ count its light components. For a
fixed number $s_0$ of steps, the occupied-skeleton
estimate required by the construction is
\[
 o_n\le\tfrac12o_{n-s_0}+Cl_{n-s_0}
                     +C\sum_{t=n-s_0+1}^nn_t^{\mathrm{cr}}.
\]
The factor $1/2$ concerns occupied cells, not just a
linear diameter. This contraction is the geometric estimate
to verify. Multiply it by
the nonincreasing $\widehat r_n$ and sum in $n$.
The delayed occupied sum is bounded by the original
one, so its half is absorbed on the left. The creation
sum involves only $s_0$ levels. This bounds the total
weighted occupied volume by weighted light stretches
and creations.

A light stretch ends at an event or at an operation
after at most $CN_t$ steps. Monotonicity bounds its
cost by $CN_t\widehat r_t$ at its first event level.
Joinings can be assigned to distinct creation leaves
by choosing one preferred predecessor at each branch
and following a nonpreferred one down to its first
leaf. A leaf is charged a bounded number of times,
and its earlier event cost is no smaller. Summing
these stretch costs gives \eqref{eq:genealogy-count}.
This derivation explains why an estimate of diameter
alone or a cost sequence increasing with level would
not give that conclusion.

It also gives the credit for a component which has
not yet completed its wait. Continue its hulls
fictitiously, with no new creation, until it enters
the operation class. Apply the same count to this
extended genealogy, within the levels at which the
flow bounds are available. The unspent part of the
creation gains bounds every future volume cost of
that continuation. A credit cannot be spent twice:
on a later joining it is assigned through the new
genealogy, and on a marked outcome the new failed
test supplies the weaker $p_0$ gain for the next wait.

The deposit constant in $cp_0(g_t)^2$ is chosen after
the fixed fractions in the collar and rim tests.
Shrinking $g_t$ cannot fix a deposit constant larger
than the available gain of one such test: both have
the same logarithmic degree. After this constant
has been fixed, the strict inequality
$\mathfrak p_0>6\mathfrak r$ makes
\[
 (C+c_a)M^4\check R_t^6=o(p_0(g_t))
             \quad(g_t\longrightarrow0)
\]
for each fixed age reserve $c_a$. A fraction
$\mathsf w<1$ of creation deposits can also be
withheld, since $(1-\mathsf w)cp_0^2$ still exceeds
all these costs and $2p_0$. The ceiling is chosen
after $c_a$ and $\mathsf w$.

\subsection{Uniform masses and local exponentiation}
\label{sec:mass-exponentiation}

\subsubsection*{The integrated history estimate}
\phantomsection\label{hyp:history-bounds}

Consider an operation of level \(i\) on a component \(X\).
Let \(\omega\) be its local history and cubic frame, \(\xi\) its
stored variables and \(\mathbf p\) the exterior field and current
pair. Once the local test-payment, retained-form and background
estimates above hold, the mass estimate is
\begin{equation}\label{eq:patch-mass}
 \frac1{|W_4|}\sum_\omega e^{cv(\omega)}
 \int^*\sup_{\mathbf p\in\mathcal U_i(X)}
           |e^{\Phi_\omega(\xi;\mathbf p)}|\,d\mu_\omega(\xi)
 \le e^{-(2-c')p_0(g_i)},\qquad 0<c'<1.
\end{equation}
The loss \(c'\) is chosen sufficiently small to leave the reserve used
in the subsequent local exponentiation, and the number \(c\ge0\) is
fixed before the coupling ceiling.
The volume \(v\) counts a stored site of level \(j\) with weight
\(\log x_j\), counting each pair consisting of a level and a site
only a bounded number of times. It is additive when the local
history is cut at its inner operations.
The upper integral \(\int^*\) permits a supremum that is not
known to be Borel. The construction below gives Borel majorants
depending only on the variables of this component.
The positive measure \(\mu_\omega\) contains the fixed
stored-variable and auxiliary reference measures, together with
the history's test factors. The exponential weights are included
in \(\Phi_\omega\). This separation is retained before taking the
exterior-pair supremum.

\emph{A majorant at fixed stored variables.}
For a healed component or an island, the main-step formula separates
the exponent as
\begin{equation}\label{eq:mass-pair-dependence}
 \Phi_\omega(\xi;\mathbf p)
  =\Phi_\omega^\circ(\xi)
      -\tfrac12q_X(\xi;\mathbf p)+V_X(\xi;\mathbf p),
 \qquad q_X=\|\zeta^{1/2}\mathrm d\mathsf H B\|^2\ge0.
\end{equation}
Here \(\zeta=\phi_X(1-\phi_{\mathrm{core}})\) is the
nonnegative shell cutoff in \eqref{eq:main-ten-term-identity},
at the level of this component.
The norm square is formed at the real glued background belonging
to the \(G\)-valued factor of a splitting of \(\mathbf p\).
This is not a positivity assertion for a quadratic form at a
complex background. The last two summands can depend on the
splitting; their sum is the invariant function of the pair.
Every other factor belongs to \(\Phi_\omega^\circ\). These factors
include the old core Wilson action, the kept gauge-fixing and
Gaussian weights, chart densities and elimination factors, and
inner operations evaluated at real pairs determined by \(\xi\).

\emph{The quadratic factor of an island.}
For a second-class component \(X\) of level \(i\), retained in a
label of level \(k'\ge i\), the two pair-dependent functions in
\eqref{eq:mass-pair-dependence} have different carriers and different
bounds. The quadratic factor can be kept as one function with carrier
\(X\). Indeed, on the recorded support the sharper main-step test gives
\[
 |B(b)|<x_i^{1/2}\delta'_i
                  =A_1(\log x_i)^{\mathfrak p_1}.
\]
The quadratic background estimate needed here is
\(q_X\le C\mathsf n\sup|B|^2\), where the number of bonds entering
the form satisfies \(\mathsf n\le C\check R_i^5\). Consequently
\begin{equation}\label{eq:mass-island-quadratic}
 0\le q_X
 \le C(\log x_i)^{2\mathfrak p_1+5\mathfrak r}
 \le C(\log x_i)^{2\mathfrak p_0-3\mathfrak r-2}.
\end{equation}
The last inequality uses
\(\mathfrak p_1\le\mathfrak p_0-4\mathfrak r-1\).
Here, as above, the background used to define the norm square is
the real glued background of the chosen splitting. Its exterior
datum on the bonds inside \(X\) is \(\mathcal M^i(U)\), determined
by \(U|_X\). The quadratic function reads no other member of the
triple. Its uniform bound is required for the regular real factors
in the evaluation domains, including
\( |\partial U-1|<2\alpha_{0,i}\eta_i^2\).

This statement concerns the sum of the walk terms. In the expansion
\[
 q_X=\sum_{\omega_1,\omega_2}
 \langle \mathrm d\mathsf H_{\omega_1}B,
                 \zeta\,\mathrm d\mathsf H_{\omega_2}B\rangle,
\]
each summand is formed at a glued minimizer on all of \(X\).
Even if its displayed walks occupy a smaller set, that minimizer
can read the exterior datum and stored variables throughout \(X\).
Thus one cannot assign that smaller set as the carrier of the
summand. After the absolutely convergent walk sum is taken,
\eqref{eq:mass-island-quadratic} bounds the single carrier-\(X\)
factor.

The component remains outside the small-field region of level
\(i+1\). The waiting rule prevents a later operation from remaking
that scale while the component remains standing: an operation
meeting its enlargement would contain \(X\) and would have to wait
past level \(i\). Hence every cube of \(X^\sim\) stays outside the
action-sequence region \(\Omega^{\rm a}_{i+1}\). The hull of \(X\)
among its stored components has a tree with one corner, so its
reduced hull length is zero. This permits the bound for
\(-q_X/2\) to be filed as an oscillation of order
\((\log x_i)^{2\mathfrak p_0}\) on \(X\).

The function \(V_X\) instead has carrier \(X^\sim\). A hull for
this carrier must also traverse its collar. More importantly, its
available bound is proportional to the number of stored sites,
as in \eqref{eq:mass-collar-volume}; that number is not bounded
by a function of \(x_i\) alone. The main step's Wilson contribution
to \(V_X\) consists of differences between its backgrounds. The
new label's Wilson action at the new background, including the
plaquettes in \(X\), is already part of the new action. It is not
an additional term of \(V_X\). This distinction is needed both
when summing the Wilson differences and when using the mass
estimate: the quadratic factor has a levelwise bound, whereas
\(V_X\) is paid by the stored-history volume.


The collar estimate for the localized terms is
\begin{equation}\label{eq:mass-collar-volume}
 \sup_{\mathbf p\in\mathcal U_i(X)}|V_X(\xi;\mathbf p)|
                      \le C_V N_{\rm st}(\omega),
\end{equation}
where stored sites are counted at their own levels.
In this estimate older regular terms are summed as pointed totals
with both subtractions. At age \(\nu\), their current-volume count
\(L^{4\nu}\) is canceled by the quadratic field factor
\(L^{-4\nu}\). The remaining weights have a convergent sum of the
form
\[
 \sum_{\nu\ge0}(1+\nu)^C
 \left[L^{-\beta_{\mathrm H}\nu}
             +L^{-\nu}+L^{4\nu}e^{-c\kappa L^\nu}\right]<\infty.
\]
These terms come from the regular potential on a generated
configuration, the surface of its localization box, and domains
leaving that box. In the old core, the part at zero response is
put in \(\Phi_\omega^\circ\); the remaining anchored response
decreases by at least \((4L^8)^{-1}\) per level, exceeding the
\(L^4\) growth of the number of finer sites. The Wilson terms are
combined with their net profiles before this estimate. Separate
Wilson differences would have a cost increasing with the age.
These are the main-step estimates behind
\eqref{eq:mass-collar-volume}.

Discarding only the nonnegative \(q_X\) gives the explicit bound
\[
 \sup_{\mathbf p}|e^{\Phi_\omega(\xi;\mathbf p)}|
 \le b_\omega(\xi)
 :=\exp\!\left(\Re\Phi_\omega^\circ(\xi)
                              +C_VN_{\rm st}(\omega)\right).
\]
This is a Borel function of \(\xi\), independent of every other
component's variables. Thus the supremum has been bounded before
any integral is taken. Taking the supremum of the integral would
not supply this function for the later substitutions.

\emph{The real inner background and the available action.}
At this point \(\xi\) lies only in the support \(S_\omega\) of a
single local operation. It need not be part of the support of a
whole effective term. The inner background must therefore be
constructed from the local tests still present in that operation.
Up to the depth \(i_\flat\), these tests have not been altered by
later preparations. At the innermost box, the half-threshold
plaquette tests of the last shell replace the tests that were
removed. On the remaining datum the field is \(e^{ig_iB}\), with
\(|g_iB|<\delta'_i\ll\varepsilon_i\).
The local background theorem then gives the real
\(U''(\xi)\), compatible with the stored averages on the enlarged
old core and regular in each layer on one further enlargement.

Keep its old-core Wilson factor in \(\Phi_\omega^\circ\).
On every share used by the outer genealogy its weight is at least
\(cx_t\) at the level \(t\) of that event. Indeed, the core cutoff
equals one there, and the spatial inverse coupling is a convex
combination of the corresponding running coefficients. At a
point entering the large-field region the coupling stops running
within \(O(\check R_t)\) levels; the drift over that interval is
\(O(\check R_t)=o(x_t)\). One cannot replace this observation by a
lower bound on all later running couplings.

Let \(G_j(\xi)\) be the nonnegative gauge-fixing action left at
level \(j\), and let \(\gamma_{\rm p}\) be the retained-form lower
bound. The disjoint payment allocation gives
\begin{equation}\label{eq:mass-outer-payment}
 A(\upsilon,U'')+\sum_jx_jG_j(\xi)
           +\tfrac14\gamma_{\rm p}
                    \sum_j\|\mathcal A_j^{\rm s}\|^2
 \ge c_{\rm cr}\sum_t p_0(g_t)^2n_t^{\rm cr}
                         +\sum_t p_0(g_t)n_t^{\rm mk}.
\end{equation}
Here \(\upsilon\) is the retained core Wilson weight; the genealogy
is cut at the inner islands. The creation and mark gains in this
formula use disjoint Wilson plaquettes, gauge-fixing factors and
stored Gaussian bonds. A newly failed operation uses no plaquette
with a corner in its old floor, and no form below its waiting
depth. Thus its gain is distinct from the older creation gains.

Before splitting any paying exponent, bound all remaining local
terms, including terms quadratic in bounded stored fields, by
their costs per site. Then one fixed part of the paying exponent
gives \eqref{eq:mass-outer-payment}, and the remaining part is used
for integration. In particular a gauge-fixing normalizer is
integrated against a fixed fraction of its own exponent:
\[
 \int e^{-\frac12x(1-\frac12\Re\tr u)}\,du
 \le C\int_{|u-1|<\epsilon}
               e^{-x(1-\frac12\Re\tr u)}\,du,
 \qquad x\epsilon^2\ge1.
\]
Both sides have order \(x^{-3/2}\). Stored Gaussian coercivity
similarly leaves a product Gaussian envelope. No further analytic
factor is charged to this integration reserve after it has been
used.

For pieces \((t,D)\) not yet reached by an operation, this gives
the pointwise product
\begin{equation}\label{eq:mrg-product-majorant}
 e^{-A_{\mathrm{cr}}}\prod_{(t,D)}w_{t,D}
 \le\prod_{\mathrm{creations}}e^{-cp_0(g_t)^2}
                     \prod_{(t,D)}\bar w_{t,D}(\xi_{t,D}),
 \qquad
 \int\bar w_{t,D}\le e^{C\#_t(D)}.
\end{equation}
The variables \(\xi_{t,D}\) of different pieces are disjoint,
and \(\#_t(D)\) counts sites, rather than large-field cells.
The functions \(\bar w_{t,D}\) contain no background parameter.
Their integral cost has no additional \(\log x_t\) per site;
the separate vacuum and scaling constants are included in the
numerical history account.

\emph{Induction over the inner islands.}
For an inner island \(Y\) of level \(m<i\), the pair
\(\mathbf q_Y(\xi)\) at which its factor is evaluated is real and
belongs to its fixed analytic core. The margin follows from
\[
 \frac{C\check R_m^2\varepsilon_m}{\alpha_{0,m}}
       \le C(\log x_m)^{\mathfrak p_0+2\mathfrak r-q}
                          \longrightarrow0.
\]
Replace this factor by its Borel majorant at fixed inner stored
variables. Its sum and integral are at most the already proved
mass \(\mathsf M(m,Y)\le1\). Additivity of \(v\) makes these
inner mass sums separate. In this induction the estimates at real
pairs are the inputs; the complex-pair supremum is a consequence
of \eqref{eq:mass-pair-dependence}, rather than a new real
positivity assumption.

The final island test \(\chi_m^c\) is outside that inner mass.
It supplies, in the parent operation, a failed ordinary plaquette
test of the new field on \(Y\), away from its boundary margin.
Its gain \(cp_0(g_m)^2\) is drawn from the parent's core Wilson
weight. Thus the island is a creation leaf in the cut genealogy,
while its entire interior is paid by its own mass. There is no
second use of its interior deposits.

Use \eqref{eq:genealogy-count} and \eqref{eq:genealogy-gain} on
this cut genealogy. All constants from normalizers, chart
Jacobians, integration and the unrecorded outcomes are charged as
a fixed multiple of the stored-site volume. The weight \(cv\)
increases that multiple. The strict inequality
\(\mathfrak p_0>6\mathfrak r\) makes the event gains pay the
volume and its entropy at a coupling ceiling chosen after \(c\).
The remaining gain is at least
\[
 2\sum_t p_0(g_t)n_t^{\rm cr}.
\]
Every leaf of the cut genealogy contains a created cell, including
the creation just supplied by an inner island's final test.
The entropy is summed from the terminal region down: for example,
\[
 \sum_{Z_{n-1}\subset\mathsf{pa}^{-1}Z_n}
          L^{-4|Z_{n-1}|_{n-1}}
 \le(1+L^{-4})^{L^4|Z_n|_n}\le e^{|Z_n|_n}.
\]
Between events the geometry is deterministic. The possible place
where a preparation is interrupted costs its logarithm and is
paid by a fixed part of the weakest mark gain. The ledger already
includes these entropy costs. The lower flow bound gives,
uniformly for every creation level \(t\le i\),
\(2p_0(g_t)\ge(2-c')p_0(g_i)\) after the coupling ceiling is
reduced. This proves \eqref{eq:patch-mass}.

\emph{The own mass of a drop-out.}
A drop-out has no real inner background of its own and its
operation contains no outer Wilson factor. Define its envelope
instead over real evaluations \(U\) compatible with \(\xi\).
The evaluations have the recorded layer regularity, the stored
averages, all recorded tests except the one that made this
component a drop-out, and a real core pair at every inner island.
Include the carrying Wilson factor
\(e^{-A(\upsilon\mathbf1_X,U)}\) on bare plaquettes whose
corners all lie in \(X\). The free form and Jacobian parameters
range only over their one-sided coercivity and upper bounds.

The allocation \eqref{eq:mass-outer-payment} holds for every such
evaluation, by the all-real-field version of
\eqref{eq:energy-of-averages}. Its right-hand side and the
remaining Gaussian and compact-group envelopes are independent
of \(U\). They therefore give a Borel majorant depending only
on \(\xi\), before the supremum over evaluations is taken.
Integrate this envelope by the same island induction, entropy
and ledger. This proves the same own-mass bound for the
drop-out without using its defining failed test.

The supremum must remain fibred over \(\xi\): an arbitrary
independent real field would include a flat field and remove the
Wilson payment. Complex evaluations of this Wilson factor are
not used. Nor is the root's fresh failure estimate applied after
discarding its indicator. That failure is paid when the drop-out
stands inside a later containing component, whose compatible
field is defined on \(X^+\) and whose fresh paying window is
available.

\emph{Uniformity at later levels and multiplication of masses.}
The domains \(\mathcal U_i(X)\), and the real evaluation fibres
for a drop-out, are fixed from the mark and the data through
level \(i\). To see why they contain later arguments, retain the
tests even at an interruption between their insertion and the
following integration. On a belt the new plaquette and
approximate-fluctuation tests control the nonpivot bonds; the
averaging constraint supplies the pivot bond. The layer estimates
are
\[
 |\partial U_k-1|\le C\check R_j^2\varepsilon_j\eta_j^2,\qquad
 |J(U_k)|_{\text{level-}j}\le C\check R_j^2\varepsilon_j.
\]
There is no loss in \(k-j\): a more distant layer's influence
decreases exponentially while its threshold grows at most as a
square root, and
\(\sup_{a\ge0}(1+a)^{1/2}e^{-(a-1)^+}<\infty\).
The preceding radius margin places each pair in its fixed core.
For a drop-out the carrying configuration, its spatial Wilson
weight and the actual forms and Jacobians are consequently
members of the prescribed evaluation fibre.

For a full term, cut at its outermost marks. Replace each marked
operation by the own-variable Borel envelope above; use
\eqref{eq:mrg-product-majorant} on the pieces outside those
marks. All remaining variable groups are disjoint after the
averaging constraints have been solved. Tonelli now integrates
the product and sums each inner history into its own mass.
This establishes the product estimate under the actual
background substitutions. Bounds on the integral at each fixed
exterior argument alone would not justify this step.

After healing, $X$ belongs to the small-field
region at every earlier level. The new action
must contain all regular terms there, and
all remainders already available to this
step, even on domains that previously met
$Z_j\cap X$ and had never been integrated.
Complete the action by adding those terms
at the new background. Subtract precisely
the same terms inside the local functions
$V_{\mathfrak X}(Y)$. The terms actually
present before healing enter there at
their substituted triples. The part of
$R^{(k)}$ to be made by the present
exponentiation is not included in this
completion. This makes the operation
noncircular.

Terms made by the ordinary step and by
the stages are treated in the same way.
Far from a treated component they are
carried unchanged. Near it use the
corresponding term formed for the history
with that component's past erased,
adding it to the new action and
subtracting it in $V_{\mathfrak X}$.
A domain touching a healed collar is
enlarged to contain the entire collar.
This convention makes every such domain
contain each collar wholly or not at all.

It is essential to combine the Wilson
pieces before this localisation.
With the weights of
\eqref{eq:xrg-stage-weights}, their sum is
\begin{equation}\label{eq:xrg-profile-cancellation}
 -\sum_{i\le k}\varphi_i b_i
     +\sum_{m=1}^{N_k}(x^{(m)}-x^{(m-1)})
                         =-\sum_{i\le k}\varphi_i''b_i.
\end{equation}
The last profile $\varphi_k$ has
disappeared from the total on the
right. The coupling difference,
counterterm and final stage difference
must use one cover so their response
to a change of this sorting profile
cancels. If they were localised
individually first, erasing $X$ could
leave a trace a fraction of a large
cell away from its $M$-cube collar.
The subsequent family would no longer
be a polymer family in the healed
collars.

For noncollar domains the resulting
bound is
\begin{equation}\label{eq:xrg-completion-bound}
 |V_{\mathfrak X}(Y)|
       \le\alpha'e^{-(1+2\beta_{\mathrm s})\kappa
                                              d^{\mathrm{out}}(Y)},
 \qquad
 \alpha'=C\max\{B_{\mathbf T},M^4E_0,C_{\mathrm S}\}+C\ge1.
\end{equation}
The length $d^{\mathrm{out}}$ exempts
the standing island hulls and the
healed collars. At a collar only the
volume bound is used, and that term
stays inside the component operation.
A level-$k$ interaction crossing a
collar can enter in full, with no
counterpart at the other background.
Its size is of order $E_0$, which
does not go to zero with the coupling.
This is the reason for the bounded
amplitude $\alpha'\ge1$.

The old boundary bound $B$ does not
enter $\alpha'$. Outside a collar its
contribution is multiplied by
\begin{equation}\label{eq:xrg-old-boundary-smallness}
 C(\log x_k)^C\left[
 g_k^{-1}e^{-cM\check R_k}e^{CL\check R_k}
 +e^{-c\kappa L\check R_k}
 +e^{-c\bar\kappa\check R_k}e^{CL\check R_k}\right].
\end{equation}
The three summands come from a
background difference, a long domain
crossing the component, and decoupling
across its width. Lower-level position
entropy gives the two factors
$e^{CL\check R_k}$.
Choose $M$ and $\bar\kappa$ after
$L$, so that the negative exponents
dominate that entropy. Since
$\check R_k\ge(\log(x_k-C_\sharp))^{\mathfrak r}$
with $\mathfrak r\ge2$, they then
dominate the inverse coupling and
all logarithmic factors as well.
For fixed $B$, decreasing the
coupling ceiling makes its entire
contribution bounded independently
of $B$. Inside the collar the old
terms are charged as volume instead.

The faithful localisation rules give
the exact locality relation
\begin{equation}\label{eq:xrg-healed-locality}
 V_{\mathfrak X'}(Y)=V_{\mathfrak X}(Y)
\end{equation}
whenever $\mathfrak X\subset\mathfrak X'$
have the same outside standing data
and $Y$ meets no added healed collar.
The rules are linear in classes of
terms, depend only on the nearby
history, and do not separately read
the sorting data whose cancellation
was used in
\eqref{eq:xrg-profile-cancellation}.
This is the local identity used to
extract the empty family. A decay
estimate alone would not replace it.

The sum over components is next converted back to an exponential
of local terms. Some terms attached to islands have bounded, rather
than small, amplitude. They must be extracted before the expansion.
Keep the collar interactions inside their integral operations.
If $\mathfrak X$ is a family of healed components, enlarge each
remaining interaction domain $Y$ by every healed collar touching it and call
the result $Y^{\mathfrak X}$. Subtract the empty-family terms only
after grouping them on these enlarged domains:
\[
 W_{\mathfrak X}(Y')=
 V_{\mathfrak X}(Y')-
       \sum_{Y^{\mathfrak X}=Y'}V_\varnothing(Y),
 \qquad W_\varnothing=0.
\]
The local exponentiation identity then has the form
\begin{align}\label{eq:empty-family}
 &\sum_\mathfrak X\Big(\prod_{X\in\mathfrak X}T'_k(X)\Big)
                           \exp\sum_YV_{\mathfrak X}(Y)
 \notag\\
 &\quad=\exp\sum_YV_\varnothing(Y)
   \sum_\mathfrak X\Big(\prod_{X\in\mathfrak X}T'_k(X)\Big)
                   \exp\sum_{Y'}W_{\mathfrak X}(Y').
\end{align}
The product of integral operations acts on the entire exponential
following it. At retained exterior argument $p$, one operation is
$[T'_k(X)F](p)=|W_4|^{-1}\sum_\omega
\int e^{\Phi_\omega(\xi;p)}F(\xi;p)\,d\mu_\omega(\xi)$.
In \eqref{eq:empty-family}, $V_\varnothing$ is evaluated at the retained
argument, and the grouped differences remain inside the component
integrals. The resummation must retain the domain enlargement,
localisation and these evaluation conventions.
Every difference on the right touches a healed collar. Each remaining polymer
then contains a healed component and has its small mass available.
The cost of nearby bounded interactions and position counting is
absorbed by that mass, using a bound of order
$C\check R_k^4\ll p_0(g_k)$. Expanding before extracting
$V_\varnothing$ would create polymers made only of bounded island
terms and would not give a convergent small-activity expansion.

The resulting exponential is a formula for sums of histories with
the same standing regions and marks. It is not the action of each
individual history. The regular and remainder interactions in the
interior are constructed once on the infinite lattice and restricted
to each nonwrapping domain. Boundary terms may retain the label and
its stored data.

More quantitatively, the mass left
after the empty family is extracted
must pay the expansion cost
\begin{equation}\label{eq:xrg-exponentiation-budget}
 1+\big[C\kappa+e^{C\alpha'+C\kappa}\big]
                        (C\check R_k)^4\le cp_0(g_k).
\end{equation}
The fourth power counts the current
$M$-cubes in the collared component,
the tree that traverses the region
ignored by its hull length, and the
possible positions of its attachments.
The amplitude $\alpha'$ is fixed
before the coupling ceiling, which
therefore depends on it. The bound
is possible because
$\mathfrak p_0>4\mathfrak r$.

After the Mayer and polymer sums,
the new local exponent decomposes as
\begin{equation}\label{eq:local-exponentiation-bound}
 \mathbf R'^{(k)}=V_\varnothing+\mathbf R'_\sharp,
 \qquad
 |\mathbf R'_\sharp(Y)|
   \le Ce^{-(1+c)p_0(g_k)}
             e^{-(1+\beta_{\mathrm s}/2)\kappa d_k^\natural(Y)}.
\end{equation}
The empty-family function vanishes
except near a new island. One current
cell inside the new small-field region,
the result is the new remainder
$R^{(k)}$; the rest is boundary data.
The stated exponent belongs to the
sum over histories with the same
standing part. A history containing
a particular healed component and
the history with its whole past
removed are different summands in
this sum. Neither summand by itself
has the final exponent.

At a new island the relevant bound
is an oscillation, not an absolute
bound on its exponent. With its
factor supports $R_f$, it has the form
\begin{equation}\label{eq:xrg-island-oscillation}
 \sum_{R_f=R}\operatorname{osc}_{\mathbf p}
                   \Phi_{\omega,f}
       \le C(\log x_k)^{2\mathfrak p_0}
                            e^{-\kappa d_k^\natural(R)}.
\end{equation}
For the whole collar the length is
zero. Away from the old core there
are only $(\log x_k)^{C\mathfrak r}$
sites at their own scales. The
regular terms there are grouped as
pointed totals with both subtractions,
and each Wilson term is taken at its
net weight. The quadratic form of
the main integral is at most
$CA_1^2\check R_k^5
 (\log x_k)^{2\mathfrak p_1}$.
The exponent margins make both
contributions at most
$C_V(\log x_k)^{2\mathfrak p_0}$.
The old core can contain arbitrarily
many stored levels. There the
anchored amplitude decreases by
at least $(4L^8)^{-1}$ per level,
more than paying the $L^4$ site
count. Its zero-response value is
in $\Phi^\circ_\omega$, and only
its oscillation is summed.
Choosing the constant in
\eqref{eq:xrg-island-oscillation}
after $C_V$ closes this part of
the island presentation.

Boundary activities use a hull length rather than
ordinary domain length. Its exempt sets are the
complete level-$j$ hulls of islands still standing
in the label. A tree must meet every ordinary cube
and must touch each chosen exempt hull contained
in its domain, but it need not fill that hull.
The minimum length divided by $M$ is
$d_j^\natural$. For a collar of a new island it
can be zero, consistently with an interaction of
bounded amplitude on the whole collar.

This length is not freely subadditive under
joining across an exempt hull. Two terms meeting
opposite sides of a large exempt box can each have
zero hull cost, while a joined tree connecting
both sides has positive length of order its side.
The construction uses an anchored joining estimate:
the domains being joined must each contain a cube
meeting the integration region. Entropy is pinned
at such a cube outside the exempt hulls. These
anchor conditions belong to the local expansion
used in \eqref{eq:empty-family}.

An island of bounded activity is kept as an
integral operation. Normalising it a second time
and expanding its logarithm would consume its
small mass and leave an order-one activity.
Likewise, a boundary slot is either carried
unchanged or replaced by a newly made slot.
Its bound is the maximum of the carried and new
bounds, not their sum at each level. If a later
component absorbs an island, its donor history
removes that island's entire earlier record
before the new hull length is assigned. These
conventions are needed for a single boundary
constant to remain uniform in depth.

\subsection{Closing the common induction}

The object carried by the induction is a
presentation of a measure, rather than
the measure alone. Write it as the
array
\begin{equation}\label{eq:xrg-inductive-array}
 \mathfrak P_k=
 \big((x_j),
       (E^{(j)},\widehat E^{(j)},b_j),
       (R^{(j)},\widehat R^{(j)}),
       (B^{(j)}_{\mathcal S}),
       (\mu_\omega,\Phi^\circ_\omega,
                                 (\Phi_{\omega,f})_f)\big)_{j\le k}.
\end{equation}
The companion families are part of the
array on all admissible tori at once.
So are the origin, domain and old
analytic space of every boundary
term, and the measures and factors
of standing islands. Two different
arrays can represent the same density.
The gain at later backgrounds and
the local dependence on a history
cannot be recovered from that density
or from a bound on each family alone.

The inductive clauses have distinct
mathematical roles. First, the regular
families are identified with the
families produced from the earlier
array by the full-lattice small-field
step, on both radius systems. Their
coefficients are
$b_j=b[\widehat E^{(j)}]=x_{j-1}-x_j$.
Their pointed totals have the lattice
covariance needed for subtraction.
Second, every activity has its
specified domain, reality and gauge
invariance. It is a function of its
pair at the admissible splittings;
near stored tests it has the switch
structure already described. The
actual evaluation triple of each
carried term belongs to that domain
when its tests pass.

Third, the array is coherent under
local charts matching the history
trace. For an island this means
matching its Borel measure and set
of stored variables, its factors
and its stored-only exponent, and
the factorwise identity
\[
 \Phi_{\iota\omega,\iota f}
           (\iota\xi;\iota\mathbf p)
                =\Phi_{\omega,f}(\xi;\mathbf p)
\]
whenever $\iota$ matches the data
on $R_f$. Its mass is bounded at
the level where the island was made;
only its oscillation in the outer
pair is bounded. No bound on the
large negative quantity
$\Phi^\circ_\omega$ is imposed.
The drop-out mass is obtained from
its compatible real evaluation set,
not by assigning it these island
clauses. Finally, each term together
with its stored variables is a
positive measure at real data.

For the source construction one
separates these weights from the
real choices used to define tests
and reference measures. Fix a
real frame $\theta_j$ with the
same weak two-sided drift bounds
and nested cube sizes. Make all
thresholds, scaling, gauge fixing
and auxiliary measures with
$\theta_j$, and form the geometric
size with $\min_{i\le j}\theta_i$.
The weight $x_j$ is allowed in
a tube
\[
 |x_j-\theta_j|\le t_j,
             \qquad t_j\asymp c\theta_j(\log\theta_j)^{-C}.
\]
This is wide in absolute terms
but makes the weight discrepancy
a small potential at the chosen
radii. In the unsourced run of
this section $\theta=x$, so
the deviation is zero. Freezing
these real choices is the
interface used in \cref{sec:sources};
it does not change the exact
mass-preserving kernels.

The regular input load can now
be checked quantitatively. If
$w_{n,j}$ is the squared ratio
of output to input radius, then
\begin{equation}\label{eq:xrg-uniform-input}
 \begin{gathered}
 w_{n,j}\asymp
       \frac{\theta_j}{\theta_n}
       \left(\frac{\log\theta_n}{\log\theta_j}\right)^{2q}
                         \le C(n+2-j),\\
 \sum_{j\le n}\omega_{n+1-j}\|E^{(j)}\|_{j,\kappa}
       +\sum_{j\le n}w_{n,j}\|R^{(j)}\|_{j,\kappa}
                           \le N_\star.
 \end{gathered}
\end{equation}
The two ratios in the first line
are not simultaneously harmful.
If $\theta_j\ge\theta_n$, the
logarithmic ratio is at most
one, while the flow bounds
$\theta_j/\theta_n$ by a
constant times $1+n-j$.
If $\theta_j<\theta_n$, the
lower flow bound gives
$\theta_n\le\theta_j+C_\sharp$,
so the logarithmic ratio is
uniformly bounded. The first
sum in the second line uses
pointed subtraction; the second
uses the remainder power and
the lower drift. This is the
bounded input required by the
arbitrary-list step.

Every new cluster is classified
as marked if one of its vertices
descends from a remainder, an
earlier marked vertex or its
counterterm. Otherwise it is
unmarked. This gives
\begin{equation}\label{eq:xrg-marked-chain}
 E^{(j)}=E^{(j),\mathrm u}+E^{(j),\mathrm m},
 \qquad b_j=b_j^{\mathrm u}+b_j^{\mathrm m},
 \qquad b_j^{\mathrm u}=b_j^0+r_j,
 \qquad b_j^{\mathrm m}=q_j.
\end{equation}
Each family carries its own
counterterm. At level $k$, let
$k_1$ be the last index with
$L^{-(k-j)}M\check R_j\le L$,
or zero if no such index exists.
The regular action passed to
the next integral is
\begin{equation}\label{eq:xrg-next-regular-action}
 \begin{split}
 \widetilde{\mathbf E}_k
 &=\sum_{\sharp=\mathrm u,\mathrm m}\sum_{j\le k}
       \left[\sum_X(E^{(j),\sharp}(X,U)
                         -E^{(j),\sharp}(X,1))
                                 -b_j^\sharp A(U)\right]\\
 &\quad+\sum_{j\le k_1}\sum_X
                         [R^{(j)}(X,U)-R^{(j)}(X,1)].
 \end{split}
\end{equation}
The unmarked chain is autonomous
on fixed radii. It supplies
$r_j$ without the inverse
coupling cost of a shrinking
radius. The marked chain has
running radii and the small
remainder load; it supplies
$q_j$. The constant-total
weight-deviation family has
coefficient zero in either
chain. This implements one
coupling prescription on both
systems of radii.

The remainder made at a step
has two origins. The ordinary
step moves the young remainders
not yet included in
\eqref{eq:xrg-next-regular-action}.
Their number is
$O(\log\log x_k)$ and each
carries $g_j^{\kappa_0}$.
The displacement ratio
$\mathsf s_k$ pays their
expansion constants and the
square of this number, leaving
at most $\frac12g_k^{\kappa_0}$.
The healed-component expansion
contributes at most
$e^{-p_0(g_k)}\le\frac12g_k^{\kappa_0}$.
Their sum therefore satisfies
the same remainder bound at
level $k$.

Boundary data are assigned to
slots by their level, domain
and origin: the ordinary step,
regular points next to regions,
remainders near islands,
stages, and exponentiation.
A carried slot is literally
unchanged. A made slot is
bounded by the constants of
the present expansion. Thus
the common bound is a maximum,
\[
 B=\max\{1,B_{\mathbf T},CM^4E_0,C_{\mathrm S},\alpha'+C\}.
\]
The made-term constants do
not involve this $B$:
in the ordinary step it
enters through $B\mathsf s_k$;
in the stages through the
negative logarithmic switch
gain; outside a main-step
collar through
\eqref{eq:xrg-old-boundary-smallness};
inside it only through the
volume term paid under its
component mass. The coupling
ceiling is chosen after all
these constants.

An island exemption ends if
a later operation absorbs it.
The old hull length can then
increase. If a domain contains
that absorbed hull, one of
its cubes meets the treated
region, so its old term is
consumed in $V(Y)$. Its
completed slot uses the term
of the donor history with
the entire past of the treated
components erased. Superposition
of histories makes this a
valid history with the stronger
hull weight. Terms not consumed
are carried with exactly their
old length. This accounts for
the apparent increase of an
old weight without asserting
a stronger bound for an
unchanged old function.

The spatial rate is spent in
a fixed order,
\begin{equation}\label{eq:xrg-rate-ladder}
 (1+4\beta_{\mathrm s})\kappa
 >(1+3\beta_{\mathrm s})\kappa
 >(1+2\beta_{\mathrm s})\kappa
 >(1+\tfrac12\beta_{\mathrm s})\kappa
 >\kappa.
\end{equation}
These are the rates for the
ordinary step, stages,
completed differences,
exponentiation and the
returned array. Each gap
pays an entropy sum. It is
not spent afresh at every
old level of a carried term.

The actual induction order is
\[
 E^{(k+1)},\widehat E^{(k+1)}
 \ \longrightarrow\ b_{k+1}
 \ \longrightarrow\ x_{k+1}
 \ \longrightarrow\ R^{(k+1)}
       \text{ and the new boundary operations}.
\]
The first formation uses only
the old prefix and its aged
remainders. The coefficient
and flow check are therefore
made before the next
large-field geometry uses
$x_{k+1}$. Along a fixed
frame the additional tube
condition for $x_{k+1}$ must
be checked separately; the
local step does not supply it.
For $\theta=x$, the flow
and nested-size estimates
give that check with zero
deviation.
The local expansion and the
history estimates thereby form
one induction, with the same
couplings and analytic domains.

\subsection{Why the final bounds are uniform in the cutoff}

The ordinary remainder of level $j$ is marginal after subtracting
its constant term. Its scale sum uses the running coupling, not a
geometric scaling gain. In an old layer $n$ the analyticity radii
introduce the two-level factor
\[
 \varpi_{jn}=\max\{1,(x_j/x_n)^2\},\qquad
 \sum_{j\le n}\varpi_{jn}g_j^{\kappa_0}
                  \le Cg_n^{\kappa_0-2}.
\]
This is one reason the flow bound and the small-field induction must
be established together. The positive drift towards the ultraviolet
makes these inverse powers summable. The regular terms, after the
coupling subtraction, also have a geometric gain in $n-j$.

Here the remainder sum is an elementary consequence of
the flow once the stated norm bound is available.
Put $p=\kappa_0/2>3$ and $h=n-j$. For $x_n$
sufficiently large compared with $C_\sharp$,
\eqref{eq:flow} gives $x_j\ge c(x_n+h)$. Therefore
\[
 \begin{split}
 \sum_{j\le n}\varpi_{jn}g_j^{\kappa_0}
 &\le\sum_{j\le n}x_j^{-p}
                +x_n^{-2}\sum_{j\le n}x_j^{-(p-2)}\\
 &\le Cx_n^{1-p}
                +Cx_n^{-2}x_n^{3-p}
 \le Cg_n^{\kappa_0-2}.
 \end{split}
\]
The comparison with an integral is valid since
$p-2>1$. Thus the choice $\kappa_0\ge7$ suffices
for this sum. No geometric forgetting of remainders
has been used. At fixed coupling the same argument
would give a factor equal to the number of levels.

After paying for the local interactions, normalisations and entropy,
each terminal retained component still has a factor
$e^{-cp_0(g_k)-\kappa_1|C|_k}$. The constant $\kappa_1$ dominates
the exponential count of connected sets through a cell. Healed
components have disappeared from the label, so the final sum is over
cells of the current scale. Summing independently over unhealed
labels at all earlier scales would instead introduce a factor of the
form
\[
 \exp\!\left[
 |\mathbb T^{(k)}|\sum_{t\le k}
             L^{4(k-t)}e^{-cp_0(g_t)}\right],
\]
which need not remain bounded as $k$ grows. Removing small old
components is what avoids this accumulation.

The stronger history estimates retain some of the excess in
\eqref{eq:genealogy-gain}. For a live lineage define its weighted
volume by summing $2M^4\check R_j^5$ per occupied level-$j$ cell,
and let its creation sum be the total of the factors
$cp_0(g_j)^2$ allocated to its creations. A numerical history weight
$\Theta_{\mathrm{age}}$ is a fixed fraction of that creation sum plus $c_a$ times the
weighted volume. Inside each retained operation the corresponding
part $\Theta_{\mathrm{age},X}(\omega)$ must already be included in the mass bound:
\[
 \frac1{|W_4|}\sum_\omega e^{cv(\omega)+\Theta_{\mathrm{age},X}(\omega)}
       \int\sup_p|e^{\Phi_\omega(\xi;p)}|\,d\mu_\omega(\xi)
 \le e^{-(2-c')p_0(g_i)}.
\]
Increasing only the coefficient of $v$ in \eqref{eq:patch-mass}
does not give this assertion: $v$ charges $\log x_j$ per stored
site, whereas the volume mark has an additional factor of order
$\check R_j$. The creation-versus-volume estimate must leave this
extra margin before it is used later.

The precise numerical mark uses the unhealed lineage.
For a component $X$ live at level $k$, let its lineage
regions be $\mathsf Z_k(X)=X$ and, recursively,
the unions of live predecessors contained in the next
region. For a creation component $W$, let $d'_j(W)$ be its tree
length measured in level-$j$ large-field cells. With $|C|_j$ counting cells, set
\begin{equation}\label{eq:mrg-prepaid-mark}
 \begin{gathered}
 \mathrm{Vol}_k(X)=\sum_{(j,C)\text{ in the lineage}}
                           2M^4\check R_j^5|C|_j,\\
 \mathrm{Dep}_k(X)=\sum_{(j,W)\text{ created in the lineage}}
                      cp_0(g_j)^2(d'_j(W)+1),\\
 \Theta_{\mathrm{age}}^{(k)}(h)=
       \sum_{X\text{ live at }k}
       [\mathsf w\mathrm{Dep}_k(X)+c_a\mathrm{Vol}_k(X)].
 \end{gathered}
\end{equation}
The lineage continues under marked operations.
Deposits used in the mass of such an operation
cannot be withheld again at a later stage.
This is why its share of the numerical mark
stands inside its original mass sum.

\emph{The volume of a prescribed trace.}
The weighted genealogy count suffices here; a separate choice of
potential at every lineage member is unnecessary. Put
\(\mathsf r_i=M^4\check R_i^5\), with
\(\check R_i=R(\min_{j\le i}x_j)\) as above. Thus
\(\mathsf r_i\), unlike the individual \(x_i\), is
nonincreasing. For a member \((i,C)\), let
\(\mathrm{pred}(i,C)\) be its live predecessors and let
\(\mathrm{Act}_{i-1}\) record predecessors on which the preceding
operation acted but left a live component. The event count is
\[
 e(i,C)=(|\mathrm{pred}(i,C)|-1)^+
       +|\mathrm{pred}(i,C)\cap\mathrm{Act}_{i-1}|
       +\sum_{W\text{ created in }C}(d'_i(W)+1).
\]
These are, respectively, merges, failed operations and creations.
Cut the genealogy at its creations and at its failed operations.
Between two cuts its occupied size contracts at an ordinary
change of partition. A drop in \(\check R_i\) does not contract
the partition, but two consecutive drops are excluded by the
lower flow bound at the chosen coupling ceiling. Once the
descendant has bounded size, its remaining wait is at most
\(C\check R_i\). Applying \eqref{eq:genealogy-count} with the
nonincreasing weight \(\mathsf r_i\) therefore gives
\begin{equation}\label{eq:trace-volume-events}
 \sum_{(i,C)\in\mathfrak l}\mathsf r_i|C|_i
 \le c_G\sum_{(i,C)\in\mathfrak l}
                     \mathsf r_i\check R_i e(i,C).
\end{equation}
This is the geometric estimate used below. In particular the
price assigned to one event is of order
\(M^4\check R_i^6\), including the wait. Counting only the
volume of its first level would miss this factor.

\emph{The analytic majorant with the trace fixed.}
Record a trace \(\mathfrak l\) by its occupied regions, creation
sets and members of \(\mathrm{Act}_i\). Its volume and deposit
are those in \eqref{eq:mrg-prepaid-mark}. Retain the numerical
weight
\begin{equation}\label{eq:trace-numerical-weight}
 \mathrm{wt}(\mathfrak l)=
 \exp\!\left[
 C_{\mathrm{chg}}\mathrm{Vol}(\mathfrak l)
 -\mathrm{Dep}(\mathfrak l)
 -2\sum_{(i,W)\text{ created}}p_0(g_i)
 -\sum_i|\mathrm{Act}_i|p_0(g_i)\right].
\end{equation}
The constant \(C_{\mathrm{chg}}\) is fixed before the coupling
ceiling. Normalizations and the local vacuum constants cost
\(C\log x_i\) per stored site; compact-group and Gaussian
integrations cost \(C\) per stored site; and data not recorded
in the trace have exponential entropy per cell. These costs
are bounded by the volume term in
\eqref{eq:trace-numerical-weight}. The creation deposits and
the two displayed linear gains are retained, rather than
absorbed into that term.

There is an important order to this estimate. First sum the
histories of components that have already healed and absorb
them into the local exponent by the exact operation
\(\mathbf R\). Only then take a majorant of the terms with
the prescribed live traces. Majorizing every healed component
separately would leave an independent sum of small weights
at every old scale; that is precisely the nonsummable
accumulation excluded above.

For the remaining marked components, unfold their mass
estimates as far as their creations. At an inner island use
the Borel envelope in the proof of \eqref{eq:patch-mass};
outside the islands use \eqref{eq:mrg-product-majorant}.
Keep the trace fixed during this induction, so the gained
deposits, the failed-operation gains and the occupied-site
costs remain numerical factors. Each failed test uses its
assigned Wilson, gauge-fixing or Gaussian share. The shares
belong to disjoint sets of plaquettes, factors and bonds.
The product envelopes depend on disjoint groups of the
remaining integration variables; their Tonelli integrals
thus multiply. At no stage is a product of separate
conditional integral norms taken.

For an unmarked small-field piece, the supported exponent
is bounded above by minus its Wilson action plus \(C\) per
site. Its Haar integral is at most one after that extensive
factor has been extracted. Applying this to all pieces not
in a prescribed trace supplies \(e^{C \mathsf V_k}\), where
\(\mathsf V_k=|\mathbb T^{(k)}|\). All the other factors have just
been assigned either to their trace or to the common vacuum
constant. Thus, writing
\(\mathsf t(h)=\tau_h(\mathcal V_k)\ge0\) for the full
unsourced mass, the finite induction gives
\begin{equation}\label{eq:history-trace-mass}
 e^{-E_{(k)}}\sum_{h:\,\text{live traces }\mathcal L}
                  \mathsf t(h)
 \le e^{C \mathsf V_k}\prod_{\mathfrak l\in\mathcal L}
                                      \mathrm{wt}(\mathfrak l).
\end{equation}
The identity \(\sum_h\mathsf t(h)=Z\) uses the exact history
kernels. The estimate uses their exact resummation together
with the local envelope and payment estimates above. It
does not assert independence of the component measures.

\emph{Summing traces without spending the mark.}
Fix the top \(X\) at level \(k\). Given \(\mathsf Z_i\),
there are at most \(2^{|\mathsf Z_i|_i}\) choices of its
creation set, \(2^{|\mathsf Z_i|_i}\) choices of
\(\mathrm{Act}_i\), and
\(2^{L^4|\mathsf Z_i|_i}\) choices of
\(\mathsf Z_{i-1}\). Set
\[
 \mathrm{ent}(\mathfrak l)
       =(L^4+2)\log2\sum_i|\mathsf Z_i|_i.
\]
Starting at \(X\) and summing these choices successively
shows
\[
 \sum_{\mathfrak l:\,\mathrm{top}(\mathfrak l)=X}
                          e^{-\mathrm{ent}(\mathfrak l)}
 \le1.
\]
Consequently it suffices to bound
\[
 \mathrm{bal}(\mathfrak l)=
       \log\mathrm{wt}(\mathfrak l)
       +\mathsf w\mathrm{Dep}(\mathfrak l)
       +c_a\mathrm{Vol}(\mathfrak l)
       +\mathrm{ent}(\mathfrak l).
\]
Fix \(0\le\mathsf w<1\), \(c_a\ge0\), and a desired cell
penalty \(c_{\mathrm{cell}}\). Equation
\eqref{eq:trace-volume-events}, with its constants enlarged
to include the entropy per cell, gives event prices
\[
 c_i^{\mathrm{ev}}=C_*\mathsf r_i\check R_i,\qquad
 (C_{\mathrm{chg}}+c_a)\mathrm{Vol}
       +\mathrm{ent}+c_{\mathrm{cell}}|X|_k
 \le\sum_{(i,C)}c_i^{\mathrm{ev}}e(i,C).
\]
At sufficiently small coupling these satisfy
\begin{equation}\label{eq:trace-event-prices}
 c_i^{\mathrm{ev}}\le
       \tfrac12(1-\mathsf w)c\,p_0(g_i)^2,
 \qquad
 c_i^{\mathrm{ev}}\le\tfrac12p_0(g_i).
\end{equation}
The first inequality pays for new creations from the free
part of their quadratic deposit. The second pays for events
without a quadratic deposit. It uses the full condition
\(\mathfrak p_0>6\mathfrak r\), since
\(\mathsf r_i\check R_i\le C(\log x_i)^{6\mathfrak r}\).
The weaker inequality
\(2\mathfrak p_0>6\mathfrak r\) would pay for creations
but would not pay for these other events.

A predecessor in \(\mathrm{Act}_{i-1}\) pays its event price
from its own linear gain at level \(i-1\). This is allowed
because \(c_i^{\mathrm{ev}}\) is nonincreasing in \(i\).
For merges, view the lineage as a tree and assign to each
member its oldest creation root. At a member with \(q\)
predecessors, send its \(q-1\) merge units to the oldest
roots of the predecessors except for the root inherited
by the member. This map is injective: a root that has been
used is no longer inherited and cannot be used again.
The oldest root of the whole tree is never used.
If the merge is at level \(i\) and its assigned root at
level \(n<i\), the lower flow bound gives
\[
 c_i^{\mathrm{ev}}
       \le\tfrac12p_0(g_i)
       \le2\bigl(p_0(g_i)-1\bigr)
       \le2p_0(g_n).
\]
Thus its price is paid by that root's linear creation gain.
This argument uses the tree structure, rather than a bound
on the overlap of geometric hulls.

If \(j_0(\mathfrak l)\) is the oldest creation level,
the resulting balance is
\begin{equation}\label{eq:trace-age-balance}
 \begin{split}
 \mathrm{bal}(\mathfrak l)+c_{\mathrm{cell}}|X|_k
 \le{}&
 -\tfrac12(1-\mathsf w)
       \sum_{(i,W)\text{ created}}
            cp_0(g_i)^2(d'_i(W)+1)\\
 &-\tfrac12\sum_i|\mathrm{Act}_i|p_0(g_i)
       -2p_0(g_{j_0(\mathfrak l)}).
 \end{split}
\end{equation}
In particular, the mark has been retained inside the original
mass calculation, and
\begin{equation}\label{eq:trace-top-sum}
 \mathrm{Wt}(X):=
 \sum_{\mathrm{top}(\mathfrak l)=X}
       \mathrm{wt}(\mathfrak l)
       e^{\mathsf w\mathrm{Dep}(\mathfrak l)
                       +c_a\mathrm{Vol}(\mathfrak l)}
 \le e^{-c_{\mathrm{cell}}|X|_k}.
\end{equation}
The unused oldest-root gain in
\eqref{eq:trace-age-balance} is also available when the
trace is subsequently healed. No deposit already consumed
by an inner operation is withheld a second time.

The global bounds use one common vacuum constant,
so the label-dependent constants cannot be
discarded. If $E_{(k)}$ is the vacuum constant of
the first $k$ steps, their required split is
\[
 E+\mathcal E_{k,\mathcal S}
   =\mathfrak E_k+\sum_{j\le k}\mathfrak E'_j,
 \quad
 -\big(\tfrac92\log x_k+C\big)|\mathbb T^{(k)}|
       \le\mathfrak E_k\le C|\mathbb T^{(k)}|,
 \quad
 |\mathfrak E'_j|\le Cr_j|Z_j|_j.
\]
The remaining-scale constant is
$\mathfrak E_k=E-E_{(k)}$ and $\mathfrak E_K=0$.
Its logarithmic sign comes from the free-bond and
gauge-fixed-site count. Four bond directions and
one site gauge parameter leave three physical
polarisations, giving $-\frac92\log x$ per
eliminated site for $\SU(2)$. The local Gaussian
numbers, though independent of the coupling,
are included once in this split. An uncontrolled
constant per bare bond would not be a permitted
normalisation error.

To pass from \eqref{eq:history-trace-mass} to stability, group
histories by the set of their live traces and use additivity
of \(\Theta_{\mathrm{age}}\). Equation
\eqref{eq:trace-top-sum} gives
\[
 e^{-E_{(k)}}\sum_h\mathsf t(h)
                         e^{\Theta_{\mathrm{age}}^{(k)}(h)}
 \le e^{C \mathsf V_k}
       \sum_{\mathfrak C}\prod_{X\in\mathfrak C}
                                      e^{-c_{\mathrm{cell}}|X|_k}.
\]
Here \(\mathfrak C\) ranges over compatible sets of connected
unions of level-\(k\) cells. Dropping compatibility increases
the sum. The graph of cells has degree at most
\(\Delta=3^4-1\). A connected set of \(n\) cells containing
a fixed cell can be encoded by a traversal of a spanning
tree of length \(2(n-1)\), so there are at most
\(\Delta^{2(n-1)}\) such sets. Choose
\(c_{\mathrm{cell}}=\log(2\Delta^2)\). Then
\[
 \sum_{X\ni\mathfrak c}e^{-c_{\mathrm{cell}}|X|_k}\le1,
 \qquad
 \sum_{\mathfrak C}\prod_{X\in\mathfrak C}
                                      e^{-c_{\mathrm{cell}}|X|_k}
 \le\prod_X(1+e^{-c_{\mathrm{cell}}|X|_k})
 \le e^{n_k}\le e^{\mathsf V_k},
\]
where \(n_k\) is the number of level-\(k\) large-field cells.
Thus the extra cost of summing live traces is at most one
per current site. The dials \(\mathsf w,c_a\) change the
coupling ceiling needed in \eqref{eq:trace-event-prices};
they do not change this extensive constant. This calculation
is applied at the terminal level for the deposit mark, and
at the stated stopped levels for the volume mark.

After the volume and label sums, the required
upper bound and floor are
\begin{equation}\label{eq:mrg-marked-stability}
 \sum_he^{-E}\mathsf t(h)e^{\Theta_{\mathrm{age}}^{(K)}(h)}
       \le e^{E_+\ell^4},\qquad
 \sum_he^{-E}\mathsf t(h)\ge e^{-C_{\mathrm{low}}\ell^4},
 \qquad E_++C_{\mathrm{low}}\le C(1+\log x_K).
\end{equation}
For $\mathsf w=0$ the same bounds hold at each
stopped level, using $E_{(k)}$, $x_k$ and
$|\mathbb T^{(k)}|$. Write
$\widetilde\nu(h)=\mathsf t(h)/\sum_{h'}\mathsf t(h')$
for its normalised positive history masses. Dividing
the two bounds gives
\begin{equation}\label{eq:mrg-mark-moment}
 \sum_h\widetilde\nu(h)
                e^{\Theta_{\mathrm{age}}^{(k)}(h)}
       \le \exp[C(1+\log x_k)|\mathbb T^{(k)}|].
\end{equation}
This extensive moment bound is not yet a
volume-independent estimate for a specified cell.

The first logarithmic power in the floor is
essential. Retain the never-large label and
restrict each current bond to a Haar ball of
radius $s\asymp g_k$. Its Haar mass is at least
$cs^3$. The lift of this coarse field, equal
to one within blocks and equal to the coarse
transport across them, has all the required
averages. Under the regularity estimate its
background passes the tests and has action at
most $Cs^2|\mathbb T^{(k)}|$. Consequently
\[
 \int\chi_k e^{-x_kA(U_k)}\,dV_k
 \ge\exp\!\left[-\left(4\log\frac1{cs^3}
                        +Cx_ks^2\right)
                        |\mathbb T^{(k)}|\right].
\]
Since $x_ks^2=O(1)$, this gives
$C_{\mathrm{low}}(x_k)=C+6\log x_k$.
Using balls of radius $\varepsilon_k$ would
instead charge $p_0(g_k)^2$ per site and
would lose the floor used in the age argument.

\subsubsection*{From the extensive mark to rarity at one cell}

The extensive bound \eqref{eq:mrg-mark-moment} is used in two
separate probability arguments. First one obtains rarity of a live
cell, and then decay in the age of its lineage. The first target is
\begin{equation}\label{eq:xrg-live-rarity}
 \widetilde\nu(\mathfrak c\subset\mathcal Z_k)
                 \le Ce^{-cp_0(g_k)}.
\end{equation}
The event concerns the history law, rather than a requirement that
the bare field be regular on a fixed physical box.

Write $\mathcal E_{\mathfrak c}^{(k)}(K')$ for the event that a
level-$k$ cell remains live through a terminal level $K'\ge k$.
A lineage live at $k$ either is healed later or remains live at the
terminal level. In the first case its operation mass has a gain
$Ce^{-cp_0(g_k)-cp_0(g_j)}$ if healing occurs at level $j$.
Summing the masses with the lineage trace fixed, and then over $j$,
gives $Ce^{-cp_0(g_k)}$. This is where the prescribed-trace
product bound above is used; one cannot sum an unnormalised
component weight without the remaining history masses.

For the never-healed part, the terminal mark charges every cell
in the lineage at each level. In particular,
\[
 \sum_{\mathfrak c}
    \mathbf1_{\mathcal E_{\mathfrak c}^{(k)}(K')}
       \le \frac{\Theta_{\mathrm{age}}^{(K')}}
                    {c_a\mathsf r_k},
       \qquad \mathsf r_k=M^4\check R_k^5.
\]
Jensen's inequality applied to the exponential moment gives a
bound on the expectation of the numerator by the terminal volume
times $C(1+\log x_{K'})$. The live event at level $k$ is invariant
under translations by level-$k$ cells. The never-healed event need
not be, since it also reads the coarser partitions. Thus one averages
the latter event over the level-$k$ cells and keeps a maximum only
for the later-healed part. The resulting first-moment bound is
\begin{equation}\label{eq:xrg-live-first-moment}
 \widetilde\nu(\mathfrak c\subset\mathcal Z_k)
 \le Ce^{-cp_0(g_k)}+C_{g'}L^{-4(K'-k)},
\end{equation}
where $g'$ is the terminal reference coupling; the constant is
independent of the torus side. The volume ratio, rather than
translation invariance of the never-healed event, supplies the
last factor.

\emph{Continuation on the same bare lattice.}
Fix \(0<\mathsf w<1\) and \(c_a>0\), and then a small anchor
coupling \(\gamma_\star\). Write \(X_\star=\gamma_\star^{-2}\)
and \(K_\star=K(\beta,\gamma_\star)\). If the original terminal
side is \(2L^m\), a continuation to a first-passage cutoff
\(K'=K(\beta,g')\), \(g\le g'\le\gamma_\star\), has terminal
side \(2L^{m-(K'-K)}\). Call this target admissible when
that side satisfies \eqref{eq:flow-volume-floor} at \(g'\).
The original target is admissible. There are finitely many
possible cutoffs between \(K\) and \(K_\star\), so choose an
admissible target \(g^\dagger\) with largest cutoff and set
\[
 K^\dagger=K(\beta,g^\dagger),\quad
 p^\dagger=K^\dagger-K,\quad
 m^\dagger=m-p^\dagger,\quad
 X^\dagger=(g^\dagger)^{-2}.
\]
The new bare side is still \(2L^{m^\dagger+K^\dagger}
=2L^{m+K}\). Moreover, every intermediate torus satisfies
\eqref{eq:prefix-volume}, since the terminal volume condition
implies all its predecessors.

The probabilities of old prefix events are unchanged by
this continuation. Indeed a step uses only the preceding
lattices and partitions, its coupling prefix and the earlier
outcomes. None of these changes. Positivity and exact
composition of its kernels give, history by history,
\[
 \sum_{h':\,h'|_k=h}\mathsf t(h')=\mathsf t(h).
\]
The total mass is the same bare Wilson partition function.
After the change of physical units, the old live event
therefore has the same probability in the prolonged run.

Put
\[
 c_{\mathrm{gap}}=
       \left\lceil C_\sharp/\lambda_\sharp\right\rceil+1.
\]
There are two possibilities. If
\(K^\dagger+c_{\mathrm{gap}}\le K_\star\), the lower drift
in \eqref{eq:flow} guarantees a record low at an index
\(K''\in(K^\dagger,K^\dagger+c_{\mathrm{gap}}]\).
A record low is a first-passage cutoff: take the target
inverse coupling \(x_{K''}-B_\sharp\), or the anchor if this
number is below \(X_\star\). If
\(m^\dagger\ge m_\star(g^\dagger)+c_{\mathrm{gap}}\),
where \(2L^{m_\star(g')}=\ell_\star(g')\), that target
would still meet the volume floor. In fact
\[
 m^\dagger-(K''-K^\dagger)
       \ge m_\star(g^\dagger)\ge m_\star(g'').
\]
This contradicts maximality. Hence in this case
\begin{equation}\label{eq:continuation-volume-case}
 m^\dagger<m_\star(g^\dagger)+c_{\mathrm{gap}},\qquad
 N^\dagger:=16L^{4m^\dagger}
       \le C(L,M)\bigl(\log(X^\dagger+B_\sharp)\bigr)^{4\mathfrak r}.
\end{equation}
If \(K^\dagger+c_{\mathrm{gap}}>K_\star\), the upper
drift bound and the terminal anchor window instead give
\begin{equation}\label{eq:continuation-coupling-case}
 p^\dagger>
       \frac{g^{-2}-X_\star-B_\sharp}{B_\sharp}
                         -c_{\mathrm{gap}}.
\end{equation}
Thus either the new torus has polylogarithmic volume in its
terminal coupling, or the continuation supplies a number
of levels of order \(g^{-2}\).

\emph{The never-healed contribution.}
Use the terminal exponential moment of the prolonged datum,
with
\(\mathsf B(y)=C(1+\log y)\):
\[
 \mathbb E^\dagger e^{\Theta_{\mathrm{age}}^{(K^\dagger)}}
       \le e^{\mathsf B(x_{K^\dagger})N^\dagger}.
\]
A cell surviving from level \(k\) to \(K^\dagger\) lies
in one terminal component's lineage. Its first nonempty
level contains a creation: otherwise that level would have
an earlier live predecessor. The deposit of this creation
and the lower flow bound imply, on the survival event,
\[
 \Theta_{\mathrm{age}}^{(K^\dagger)}
       \ge\mathsf w c\,\bigl(p_0(g_k)-1\bigr)^2.
\]
In the volume case \eqref{eq:continuation-volume-case},
the extensive cost is at most
\[
 C\bigl(\log(X^\dagger+B_\sharp)\bigr)^{4\mathfrak r+1}.
\]
Choose the anchor after the dials so this is at most one
quarter of the displayed deposit, uniformly for
\(X^\dagger\ge X_\star\). This is possible since
\(2\mathfrak p_0>4\mathfrak r+1\). Exponential Markov
then bounds survival at each specified cell by
\(e^{-c_{\mathrm r}p_0(g_k)}\), with a fixed positive
\(c_{\mathrm r}\) below the rate allowed by healing.

In the coupling case use the volume part of the same mark.
Distinct terminal components have disjoint lineage regions,
so each surviving level-\(k\) cell is counted once and
\[
 \sum_{\mathfrak c}
       \mathbf1_{\mathcal E_{\mathfrak c}^{(k)}(K^\dagger)}
 \le
 \frac{\Theta_{\mathrm{age}}^{(K^\dagger)}}
                           {2c_aM^4\check R_k^5}.
\]
There are
\(n_k=N^\dagger L^{4(K^\dagger-k)}/(M\check R_k)^4\)
level-\(k\) cells. Jensen's inequality and the moment bound
therefore yield
\begin{equation}\label{eq:continuation-survival-density}
 \frac1{n_k}\sum_{\mathfrak c}
       \widetilde\nu^\dagger
           (\mathcal E_{\mathfrak c}^{(k)}(K^\dagger))
 \le
 \frac{\mathsf B(x_{K^\dagger})
                         L^{-4(K^\dagger-k)}}{2c_a\check R_k}.
\end{equation}
The torus volume has cancelled. This is an average over cells;
no translation invariance of the survival event is used.

To compare its depth with \(p_0(g_k)\), put \(u=K-k\) and
\(\bar x_u=g^{-2}+B_\sharp(u+1)\). Then
\[
 x_k\le\bar x_u,\qquad
 K^\dagger-k=p^\dagger+u,\qquad
 x_{K^\dagger}\le g^{-2}+B_\sharp.
\]
With \(\mathfrak P(y)=A_0(\log y)^{\mathfrak p_0}\), choose
\(u_R(g)\) large enough that, for every \(u\ge0\),
\[
 4(u_R(g)+u)\log L
 \ge c_{\mathrm r}\mathfrak P(\bar x_u)
            +\log^+\frac{\mathsf B(g^{-2}+B_\sharp)}{2c_a}.
\]
This requires only \(u_R(g)\le C(\log g^{-2})^{\mathfrak p_0}\).
Indeed \(\bar x_u\le(g^{-2}+B_\sharp)(1+u)\), and the
part of \(C(\log(1+u))^{\mathfrak p_0}\) exceeding
\(4u\log L\) has a finite supremum independent of \(g\).
By \eqref{eq:continuation-coupling-case}, decreasing the
original coupling ceiling makes \(p^\dagger\ge u_R(g)\).
Equation \eqref{eq:continuation-survival-density} is then
at most \(e^{-c_{\mathrm r}p_0(g_k)}\).

Finally the live event at level \(k\) is translation
covariant, although its never-healed part need not be.
Average its decomposition into healed and never-healed
fates over all level-\(k\) cells. The healed contribution
has the already obtained bound at every cell; the preceding
average bounds the other contribution. Prefix consistency
transports the result to the original run and proves
\eqref{eq:xrg-live-rarity}, uniformly over the volume floor.
The order of constants is the construction and the desired
rate, then the dials, then the anchor, then the original
small-coupling ceiling. The background estimates throughout
this argument use local determining sets; demanding bare
regularity throughout the torus would not give this
volume-uniform statement.

For decay in age, let
$\mathrm{Live}_{k,\ge\alpha}(\mathfrak c)$ mean that the specified cell
lies in a live component at level $k$ whose unhealed lineage has age
at least $\alpha$. The estimate needed for two points is the weighted
tail bound for the age of a live component,
\begin{equation}\label{eq:lineage-age}
 \sum_{\alpha\ge0}L^{b\alpha}\,
 \widetilde\nu\big(\mathrm{Live}_{k,\ge\alpha}(\mathfrak c)\big)
 \le C_b e^{-cp_0(g_k)}.
\end{equation}
Its constants and coupling ceiling are chosen after a fixed rate
$b\ge8$.

\subsubsection*{The lineage and its survival witnesses.}

The age needed in the two-point estimate is the age of an
unhealed lineage. If $j_0$ is its first nonempty level, that age
is $k-j_0$. For comparison, let $Z_t^{\mathrm{lab}}$ be the
large-field region retained at scale $t$ in the label of the
current history. The label age of a live component $C$ is
\[
 \mathrm{age}^{\mathrm{lab}}(C)
       =k-\min\{t:C\cap Z_t^{\mathrm{lab}}\ne\varnothing\},
\]
where regions at different scales are viewed as point sets in the
same torus. Its seed level is the earliest step at which a
created region is geometrically contained in $C$, whether or not
that region was subsequently healed; the seed age is $k$ minus
this level. Every old stored region in the label has an unhealed
predecessor chain to the current component. Thus
\[
 \mathrm{age}^{\mathrm{lab}}(C)
       \le k-j_0\le\mathrm{age}^{\mathrm{seed}}(C).
\]
The second inequality can be strict: a region can have been healed
and its former position covered by a new live component. Such a
region contributes to the seed age but supplies no surviving
predecessor and no deposit to the new lineage. The argument below
therefore uses survival, rather than a sum over all old creations
near the current cell.

For a level-$n$ cell $z$, write $\mathcal E_z^{(n)}(k)$ for the
event that it is live at every level $n,\ldots,k$. The local
probability estimate is made on tiles of side
$l_\square\asymp M\check R_k$, in level-$k$ units. The side is
a power of $L$ and spans at least sixteen level-$k$ cells. The
faces are bond mirrors, and a tiling is translated by whole
level-$k$ cells. A cell of any finer level is called deep when
its containing level-$k$ cell lies in the central three quarters
of the tile.

For an anchor $z$ deep in one tile and distinct other tiles
$\square_1,\ldots,\square_r$, the bounds to establish are
\begin{equation}\label{eq:mrg-history-chessboard}
 \begin{split}
 \widetilde\nu(\mathcal E_z^{(n)}(k))
       &\le l_\square^{-4}q_a^{c(k-n+1)},\\
 \sum_{c_i\text{ deep in }\square_i}
 \widetilde\nu\!\left(\mathcal E_z^{(n)}(k)
                 \cap\bigcap_{i=1}^r\{c_i\subset\mathcal Z_k\}\right)
       &\le l_\square^{-4}q_a^{c(k-n+1)}\epsilon^r,
             \qquad q_a=e^{-c_a},\quad c=1/16.
 \end{split}
\end{equation}
The constant $c$ is fixed before $c_a$; it comes from repeating
a pattern in one tile out of $2^4$ and taking the chessboard
root. The first inequality is also required on tori too small
for this tiling. That case needs a separate verification; the
large-torus reflection argument does not cover it by itself.

The extensive moment bound has the correct surplus for a repeated
pattern. For this bound set $\mathsf w=0$ and stop the run at
level $k$. For an event $F$ of that stopped run,
\[
 \widetilde\nu(F)
 \le\exp\!\left[C(1+\log x_k)|\mathbb T^{(k)}|
                   -\inf_F\Theta_{\mathrm{age}}^{(k)}\right].
\]
A cell surviving from $n$ to $k$ is charged in the lineage volume
at every one of those levels. Multiple surviving cells are charged
once for each distinct cell of each level, with weight
$2c_aM^4\check R_j^5$. The difficulty is to apply reflection
positivity to a local event determined by those histories.

The histories use the bare field and auxiliary draws from the
positive kernels. Realise each possible draw up to level $k$ by
an independent uniform variable, attached to the location of that
kernel. Tile data consist of the bare field restricted to the
tile and these attached variables. Reflection positivity applies
to functions of the bare field. Accordingly the witness construction
must eliminate the dependence on unspecified exterior draws and
then integrate the interior auxiliary variables in a way compatible
with products over distinct tiles.

\subsubsection*{Agreement inside a tile and threshold windows.}

Compare two runs with identical tile data. The local agreement
statement is to be proved by induction on the level and, within
a level, on the completed stages of the operation. It is needed
only under conditions imposed on the first run inside the tile;
the second run has no restriction outside it. Two exceptional
behaviours are removed: a tested scaled fluctuation can lie very
near its threshold, or a long cluster of trigger cells can reach
through the layers reserved for the comparison.

Let $\delta_{\mathrm{loc}}>0$ be the decay rate for the difference
of backgrounds determined by identical data in a region. At a
test of step $j$, with fluctuation scaled by $g_i$, $i\le j$,
the exceptional window has width
\begin{equation}\label{eq:history-decoupling-window}
 \begin{gathered}
 \frac{2\varepsilon_j^{\mathrm{dec}}}{g_i},\qquad
 \varepsilon_j^{\mathrm{dec}}
             =Ce^{-\delta_{\mathrm{loc}}D_j}\varepsilon_j,\\
 \delta_{\mathrm{loc}}D_j
   =C+4\log l_\square+(4\log L+\log2)(k-j)
                  +O(N_j\log L+\log p_0(g_j)).
 \end{gathered}
\end{equation}
Here $N_j=\check R_j$ is the waiting depth. The last terms allow
for the finite losses of the waiting window and tests. Their bounds
have to be established before $D_j$ is
fixed. The displayed distance gives the factor
$l_\square^{-4}(2L^4)^{-(k-j)}$ in the window width. There are
at most $Cl_\square^4L^{4(k-j)}$ possible level-$j$ sites in the
tile. After the remaining scale factors and the conditional-density
bound for the tested coordinate have been included, their product
is summable as $\sum_{j\le k}2^{-(k-j)}$. A width estimate alone
does not give the probability of a window: that requires the
uniform density estimate in the same scaled coordinate.

\emph{The conditional estimate at a sharp test.}
The density bound can be reduced to a one-coordinate estimate,
provided that the actual test order and the analytic margins are
retained. Identify $\mathfrak{su}(2)$ with $\R^3$, and let
$a$ be the scaled coordinate at one test. Write $\mathfrak t$
for its finite set of positive thresholds. Choose $W>0$ so that
the thresholds are at least $2W$ and are pairwise $2W$ apart.
The two shells at $t\in\mathfrak t$ are
$t\le|a|\le t+W$ and $t-W\le|a|\le t$.
Suppose that $\mathcal F$ is real-valued, $\psi\ge0$ is nondecreasing with the distance
from $t$ on each radial segment of either shell, and that
$|\partial_r\mathcal F|\le2/W$ on the shells where
$\psi\ne0$. Then, for $0\le u\le W$,
\begin{equation}\label{eq:revised-shell-bound}
 \int_{\operatorname{dist}(|a|,\mathfrak t)\le u}
                  e^{\mathcal F(a)}\psi(a)\,da
 \le C\frac uW\int_{\R^3}e^{\mathcal F(a)}\psi(a)\,da.
\end{equation}
Indeed, on a radial segment of length $W$, the exponential
varies by at most $e^2$, and the polar Jacobian $r^2$ varies
by at most a fixed factor since $t\ge2W$. The average of
$\psi$ on the part nearest the threshold is at most its
average on the whole segment. Integrate over directions and
sum the disjoint shells. No regularity of $\psi$ is used.

At the stage of a test $\sigma$ in the positive real run,
disintegrate its law with respect to all other coordinates $s$. Its contribution
has the form
\[
 Z^{-1}e^{\mathcal F_s(a)}\chi_s^{\rm st}(a)\,\Pi(ds)\,da,
\]
where $\chi_s^{\rm st}$ is the product of the sharp tests
already present. To exclude windows at later tests, multiply
this density by the conditional probability of their avoidance
and of the specified live-history event $E$. Thus the factor
in \eqref{eq:revised-shell-bound} is
\[
 \psi_s(a)=Z^{-1}\chi_s^{\rm st}(a)
                         \mathbb P(E\mid s,a).
\]
The event concerns the histories still live at this stage;
components already integrated out impose no further condition.

The conditional argument requires an interval-reading property
for the later transition kernels as well as for the labels and
rules: after its test, the coordinate must enter each of them
only through its recorded threshold interval. The rule that
labels record intervals does not by itself establish this
property for a kernel that also reads a continuous stored field.
It must be checked for the actual conditional law of the run.
Under this property, the remaining sharp-test checks are as follows. Other sharp
functions either test $|a|$ at a threshold separated by $2W$,
or test $|e^{ig_i a}v-1|$ at such a threshold. In the latter
case the centre displacement must be at most $cg_iW$,
and the error of the small exponential chart must have the
same bound. These conditions make the functions constant on
each shell. For a split test, they follow from the interior
placement of its test cubes and the local background comparison;
the test geometry is that of
Ba\l aban\cite[equation~(3.19)]{Balaban1988} and
\cite[equations~(1.56)--(1.58)]{Balaban1989I}.
A completion test on a bond previously controlled by a passed
cube test uses the stored averaged coordinate; both endpoints
are in the old small-field region, so the same interior
comparison applies. For a completion test in the remaining
strip, the stage data are tested first, and the broader
approximate fluctuation test equals one on these shells.
These are the actual support and centre conditions used by
the test construction; they are needed in addition to
separation of the thresholds.

There is an exception when the same bond also belongs to a
collar test. The narrow approximate collar test is performed
before the completion test. In the scaled coordinate its
radius is $2\delta'_i/g_i$, whereas a completion shell has
$|a|\ge c\delta_i/g_i$. The two centres differ by
$o(\delta_i)$ in the unscaled coordinate, and
$\delta'_i/\delta_i\to0$. Consequently the earlier passed
collar indicator is zero on the completion shells. The product
$\psi_s$, rather than its future-probability factor separately,
therefore vanishes there. At a collar coordinate itself, the
broader test has radius $2\delta'_i/g_i$, separated from
the threshold $\delta'_i/g_i$; the local background
displacement is smaller than the shell width.

If a later window tests the same level and bond at the same
threshold, its avoidance indicator is nondecreasing with
distance from that threshold. All other later decisions are
constant on the shell. Averaging over the later coins with
the interval-dependent transition laws preserves this
monotonicity. This proves the required property of
$\psi_s$, including the repeated-bond cases.

The slope estimate must use the same scaled coordinate. The
quadratic gradient is an exponentially decaying kernel applied
to fields bounded by $Cp_0(g_j)$, so it is bounded by
$Cp_0(g_j)$. For the interaction, Cauchy's estimate divides
the local supremum by its available analytic margin.
Ordinary localization domains through a bond have a bounded
exponentially weighted sum. Preparatory domains have instead
the bound $C(\log x_j)^{4\mathfrak r}$, because the small
stored boxes have zero cost in their length. At a collar
coordinate the analytic margin is at least a fixed multiple
of $\delta'_i/g_i\ge1$. Hence their contribution is at most
$C(\log x_j)^{4\mathfrak r}$, which is bounded by
$Cp_0(g_j)$ since $\mathfrak p_0>4\mathfrak r$.
Using only a margin of order $1/p_0$ for these domains would
give an extra logarithmic factor. The slope bound follows
under these actual localization and margin estimates.

Choose $W_\sigma=c/p_0(g_i)$. The flow comparison within
the waiting window bounds $p_0(g_j)$ by a fixed multiple of
$p_0(g_i)$, so the slope condition in
\eqref{eq:revised-shell-bound} holds after fixing $c$.
Integrating in $s$ gives a loss at most
$\eta_\sigma=Cu_\sigma/W_\sigma$ when the window of
width $u_\sigma$ at $\sigma$ is removed. Remove the windows
in decreasing order of their stages. At each removal the
remaining event has the form just considered, and therefore
\[
 \mathbb P(E\text{ with all windows avoided})
 \ge \prod_\sigma(1-\eta_\sigma)\mathbb P(E)
 \ge e^{-2\sum_\sigma\eta_\sigma}\mathbb P(E)
 \quad(\eta_\sigma\le\tfrac12).
\]
The site count and the choice of $D_j$ in
\eqref{eq:history-decoupling-window} make the sum over one
tile bounded. This supplies the stated conditional window
estimate once the interval-reading property of the conditional
law and the uniform local background, slope and analytic-margin
bounds have been verified for the run.

The comparison uses a separation of $CD_j$ level-$j$ lattice
units. In level-$j$ cells this is
$\lceil CD_j/(M\check R_j)\rceil$ layers. Their physical width
in current, level-$k$ units is multiplied by $L^{-(k-j)}$.
These losses, together with the widths needed to exclude long
trigger clusters, must have a summable total fitting inside a
fixed fraction of the tile. The fine-level comparisons therefore
use narrower and narrower layers in current units, despite the
increasing number of tested fine sites.

Away from the exceptional windows, a background error bounded by
$\varepsilon_j^{\mathrm{dec}}$ cannot change the outcome of the
corresponding test. The local agreement argument must then identify
the current and stored block fields, the label regions, and every
live component meeting the smaller tile. Such a component is to
be contained in the original tile; its label, inner history and
stored block fields are common to both runs.

The assertion does not identify all stored fluctuation fields.
On bonds examined by a passed fluctuation test, those fields are
formed using an averaged background and can still depend on
exterior data. The later operation uses them in one test; on the
relevant bonds that test must pass in both runs. This is part of
the agreement verification, rather than a consequence of equality
of the stored block fields. One also needs the determining set to
start from the local region $\Omega_0=\Omega_1^\sim$. Requiring
regularity at every bare plaquette of the torus would let a distant
rough plaquette change every background and would destroy the
claimed agreement.

\subsubsection*{Trigger patterns and bare-field witnesses.}

A trigger at level $J$ is a cell in $\Omega_J^c$, or a cell
carrying a tested bond where the scaled fluctuation exceeds a
lower threshold. If a long trigger cluster occurs, select a
pattern of
\begin{equation}\label{eq:history-trigger-pattern}
 m=1+2(k-J),\qquad
 c_{\mathrm{adj}}(J)=
       \left\lceil\frac{CD_J}{M\check R_J}\right\rceil+C,
 \qquad
 m_0(J)=6(c_{\mathrm{adj}}(J)+1)m
\end{equation}
level-$J$ cells. Any two pattern cells are separated by at least
$2c_{\mathrm{adj}}(J)+3$ cells, and each lies within
$m_0(J)/2$ cells of the first. The step $c_{\mathrm{adj}}(J)$
is the largest advance allowed by the decoupling comparison.
The factor six leaves room to select the $m$ separated cells
inside a cluster of diameter exceeding $m_0(J)$.
Each selected cell either lies within two cells of a large-field
cell or carries a lower-threshold fluctuation.

The same layers used in the comparison must accommodate these
patterns. In current units their required bound has the form
\[
 \sum_{J\le k}M\check R_Jm_0(J)L^{-(k-J)}
                              \le l_\square/32.
\]
The tile side is chosen with this fixed reserve. Establishing the
agreement statement with these actual widths is the geometric
and analytic verification needed here.

For a survival event deep inside the tile, the comparison has
two outcomes. Either the good tile data determine survival
independently of the exterior run, or the first failure of
agreement is accompanied by a window or one of the patterns
above. This is the proposed reduction to local witnesses.
After the auxiliary draws have been treated, it must produce
finite families of functions $0\le f_{i,\iota}\le1$ of the bare
field in tile $i$, with a uniformly bounded cost per tile.
For survival events $F_i$ read in distinct tiles, the needed
domination and its chessboard consequence are
\begin{equation}\label{eq:history-witness-domination}
 \widetilde\nu\!\left(\bigcap_{i=0}^r F_i\right)
 \le C^{r+1}\sum_{\iota_0,\ldots,\iota_r}
                  \left\langle\prod_{i=0}^rf_{i,\iota_i}\right\rangle
 \le\prod_{i=0}^r\left(C\sum_{\iota_i}
                                \|f_{i,\iota_i}\|_{\mathrm{ch}}\right).
\end{equation}
The second inequality applies reflection positivity to these
bare-field functions. Its use is legitimate only after the first
inequality and the uniform witness construction have been
established. The names of history events alone supply neither.
In particular, independence of the original auxiliary variables
does not by itself prove the first inequality after the tests
and conditional kernels have been applied.

There are three estimates for a repeated witness pattern.
A trigger fluctuation with no large-field cell within two cells
is bounded using its conditional Gaussian law in its box. A cell
whose component stays live until $k$ is paid by the volume mark
of the run stopped at $k$. A cell whose component is healed before
$k$ is paid by that component's lineage mass. The last class
must remain in the witness bound even though it no longer appears
in the terminal label. Conditional-density estimates handle the
exceptional windows. These estimates and local agreement are the
unfinished verifications behind
\eqref{eq:history-witness-domination}.

\subsubsection*{The chessboard root and the sum over patterns.}

The root of a repeated witness probability changes the extensive
factor into
\[
 \exp[C(1+\log x_k)l_\square^4]
\]
per tile. A pattern repeated in one of the $2^4$ classes of tiles
retains $q_a^c$ per charged level, with $c=1/16$. The surplus
from one occupied level is much larger than this cost, since
\[
 l_\square^4\log x_k\ll
          \mathsf r_k:=M^4\check R_k^5.
\]
Here $l_\square\asymp M\check R_k$ and
$\check R_k\gg\log x_k$. A fixed part of
$q_a^{c\mathsf r_k}$ pays the extensive factor, the number of
current cells in one tile and the factor $\epsilon^{-1}$, leaving
the factor $l_\square^{-4}$ in
\eqref{eq:mrg-history-chessboard}.

The remaining pattern count is performed level by level. At
level $J$, the first cell has at most
$Cl_\square^4L^{4(k-J)}$ positions. Each of the other $m-1$
cells is confined within $m_0(J)/2$ cells of the first. Dropping
the separation condition for this upper bound gives
\begin{equation}\label{eq:history-pattern-count}
 \#\{\text{patterns of level }J\}
       \le Cl_\square^4L^{4(k-J)}
                        (m_0(J)+1)^{4(m-1)},
       \qquad m_0(J)\le C(x_k+k-J)^3.
\end{equation}
Apart from the first-cell positions, its logarithm is bounded by
$Cm\log(x_k+k-J)$. The mark gives a factor
$q_a^{\mathsf r_J}$ for each of the $m=1+2(k-J)$ cells before
the fixed shares are spent. The lower flow bound, and the use of
the running minimum in $\check R_J$, give fixed constants
$c_1>0$ and $C$ for which
\[
 \mathsf r_J\ge
 M^4\big(\log(x_k+c_1(k-J)-C)\big)^{5\mathfrak r}.
\]
Thus a fixed fraction of the cell weights dominates that logarithmic
entropy uniformly in $k-J$, after the coupling ceiling has been
chosen. The comparison is uniform in the old level because the
inverse coupling grows at least linearly with $k-J$.

A pattern at $J$ together with survival from $n$ up to $J-1$
retains the exponent
\[
 (J-n)^++m=(J-n)^++1+2(k-J)
                    \ge(k-n+1)+(k-J).
\]
The first part gives the required survival factor
$q_a^{c(k-n+1)}$. The second part is compared with the positions
of the first pattern cell:
\[
 L^{4(k-J)}q_a^{c(k-J)}
      =(L^4q_a^c)^{k-J}\le2^{-(k-J)},
 \qquad L^4q_a^c\le\tfrac12.
\]
Summing over $J$ is therefore geometric. The same surplus can
leave a factor $\epsilon$ for every additional tile with a deep
live cell. Once the witness reduction is verified, these counting steps yield
\eqref{eq:mrg-history-chessboard}; none of their constants depend
on the number of tiles of the torus.

\subsubsection*{From survival to the weighted age series.}

Fix a component containing the specified current cell and having
lineage age $a$. At its first level $n=k-a$, choose an anchor
in its lineage region. It remains live through all levels
$n,\ldots,k$ and belongs to the current connected component.
A tile contains at most $Cl_\square^4L^{4a}$ anchors, each with
survival bound $l_\square^{-4}q_a^{c(a+1)}$. Thus nearby anchors
cost at most $CL^{4a}q_a^{c(a+1)}$.

For far anchors, a connected live component crosses a chain of
current occupied tiles. Use the $2^4$ tilings whose face planes
are displaced by $\pm\lfloor n_\square/4\rfloor$ cells from a
corner of the level-$k$ cell containing the anchor, where
$n_\square\ge16$ is the number of current cells along a tile
side. The anchor is deep in all these tilings, and every current
cell is deep in a tile of at least one of them. From a chain
of length $i$, one tiling contains at least $i/16$ deep cells.
A tile meets at most $16$ tiles of the original board, so these
include cells in at least $i/16^2$ distinct tiles. For
$i\ge i_0:=2\cdot16^2$, at least $i/i_0$ tiles differ from the
anchor tile.

The second inequality in
\eqref{eq:mrg-history-chessboard} bounds those tiles by
$\epsilon^{i/i_0}$. If $C_{\mathrm{ent}}$ is the entropy per
link of a chain, choose
$\epsilon=(2C_{\mathrm{ent}})^{-i_0}$. The extra cost, including
the choice of tiling and chains, is then at most
$16C_{\mathrm{ent}}^i\epsilon^{i/i_0}=16\cdot2^{-i}$.
It is summable against the polynomial number of tiles in distance
shells. Hence $CL^{4a}q_a^{c(a+1)}$ bounds all anchors, including
the distant ones.

Choose the fixed rate $c_a$ so that
$L^4q_a^c\le(2L^b)^{-2}$. This gives a pure age bound of order
$(2L^b)^{-2a}$. Its geometric mean with the live-cell rarity
bound \eqref{eq:xrg-live-rarity} gives
\[
 \widetilde\nu(\mathrm{Live}_{k,=a}(\mathfrak c))
       \le Ce^{-cp_0(g_k)}(2L^b)^{-a},
\]
with the constant in the first exponential decreased if necessary.
The conversion from exact age to tail age can be written as
\[
 \begin{split}
 \sum_{\alpha\ge0}L^{b\alpha}
       \widetilde\nu(\mathrm{Live}_{k,\ge\alpha}(\mathfrak c))
 &\le Ce^{-cp_0(g_k)}
       \sum_{a\ge0}(2L^b)^{-a}
                         \sum_{\alpha=0}^aL^{b\alpha}\\
 &\le C_be^{-cp_0(g_k)}\sum_{a\ge0}2^{-a}.
 \end{split}
\]
This proves the counting deduction \eqref{eq:lineage-age} from
the local witness and rarity estimates. The constants and the
ceiling depend on the requested fixed $b\ge8$. One finite
$c_a$ is not asserted to serve all rates simultaneously.

There is also a shorter estimate on the reference torus used in
\cref{sec:two-point}. Split the lineage according to whether it
is later healed. For the healed part sum its operation masses
only over traces of age at least $a$. For the never-healed part
use the terminal volume mark: each level with a lineage member
contributes at least $\bar c_a=2c_aM^4$, so the charges add
without a sum over anchor positions. The resulting two-term
bound is
\begin{equation}\label{eq:two-term-age}
 \begin{split}
 \widetilde\nu(\mathrm{Live}_{k,\ge a}(\mathfrak c))
 &\le Ce^{-\bar c_a a}e^{-cp_0(g_k)}\\
 &\quad+C\exp\!\left[C(1+\log x_{K'})|\mathbb T^{(K')}|
          -2p_0(g_k)-\bar c_a(a+K'-k+1)\right],
 \end{split}
\end{equation}
where $K'$ is the last level of an admissible prolongation,
including the original run as one choice. The first term has
no volume factor. The second retains a terminal volume cost
and consequently is not a volume-independent substitute for
\eqref{eq:mrg-history-chessboard}.

On the reference torus the terminal size is polylogarithmic.
Its extensive exponent has degree $4\mathfrak r+1$ in the
logarithm, whereas $p_0$ has the higher degree fixed earlier.
With the same terminal comparisons as in the rarity argument,
this cost is absorbed by $2p_0(g_k)$. Choosing
$\bar c_a\ge b\log L+\log2$ then gives
\eqref{eq:lineage-age} on that reference torus. The reflection
comparison in \cref{sec:two-point} transfers its two-point
consequence. The general history bound at all volumes still
uses the local witnesses above. Both methods split off the
later-healed lineages; a mark counting only components surviving
to one fixed terminal level would omit them.

The history estimates and the coupling bounds \eqref{eq:flow}
are to be proved in one induction on common analyticity domains.
Choose the exponents first, then $L$, $M_1$, the activity decay
rate $\kappa$, and $M$, in the order of \cref{sec:rg}.
Fix the interaction, threshold and history reserves before the
coupling ceiling. The coupling
calculation checks the next inverse coupling before the next
large-field step.


\section{Sources and uniform bounds on correlations}\label{sec:sources}

The observable is followed through the integration by changing the
inverse coupling. We keep the choices of the real construction fixed,
carry the source through the local expansion, and separate its scalar
energies from the normalized generating function. Locality removes the
ultraviolet divergence from separated mixed derivatives; chessboard
comparison removes the volume dependence. A graded version of this
comparison gives factorial growth in the number of insertions.

\subsection{The exact sourced identity}

For a plaquette weight $v$, write
\[
 A(v,U)=\sum_pv(x(p))a_p(U).
\]
Fix real functions $f_i$ and put
\[
 w=c+\sum_{i=1}^nz_if_i,
 \qquad \|w\|_\flat=|c|+\sum_i|z_i|\|f_i\|_\sharp,
 \qquad \|f\|_\sharp=\max\{\|f\|_\infty,CM\Lip f\}.
\]
The source radius $\rho>0$ is fixed before $g$ is made small.
Inserting $e^{A(w,U)}$ changes $e^{-A(\beta,U)}$ to
$e^{-A(\beta-w,U)}$. There is no power of the lattice spacing in this
identity. For $c=0$, its normalized integral is
\[
 \Gamma(z)=\frac{Z(w)}{Z(0)}
    =\left\langle e^{\sum_i z_i\cO_a(f_i)}\right\rangle.
\]
At fixed cutoff the variables are bounded and $\Gamma$ is entire. Its
logarithm is defined near zero, since $\Gamma(0)=1$, and
\[
 C_n(f_1,\ldots,f_n)=\partial_1\cdots\partial_n\log\Gamma(0).
\]
For $n\ge2$ this is unchanged by subtracting any constants from the
individual observables.

\begin{proposition}[Frozen source transport]\label{prop:frozen-source-transport}
Fix a real run to level $k$. Keep its tests, backgrounds, localizations,
gauge fixings, covariances, and reference probability laws fixed when
the source is introduced. There are source-independent kernels $K_h$
such that
\[
 \tau_h^w(dV)=\int e^{A(w,U)}e^{-A(\beta,U)}K_h(U,dV)\,dU,
 \qquad \sum_h\tau_h^w(\mathcal V_k)=Z(w).
\]
The unsourced measures are positive. The sourced measures are entire
complex measures, and their total-mass identity is valid for complex
sources.
\end{proposition}

\begin{proof}
Every operation of the fixed run is linear in the density. The averaging
constraint has total integral one in the new block variable. A field test
and its complementary alternative sum to one. Gauge fixing has total
integral one over the gauge group, for the gauge-invariant density to
which it is applied. A change of variables includes its Jacobian.
Finally, replacing a component measure by its mass times a fixed
probability measure preserves the mass of that component. Composing
these operations gives $K_h$, and summing their alternatives gives the
identity on total integrals. This proof is an identity of linear
functionals and does not use positivity of the input density.

At finite cutoff the integration variables range over compact group
spaces and $A(w,U)$ is a finite linear combination of bounded functions.
The power series of its exponential is locally uniformly integrable in
the source parameters. It can therefore be integrated term by term
against each fixed kernel, which proves the asserted holomorphy.
\end{proof}

Source derivatives use the total-mass identity of
\cref{prop:frozen-source-transport}. Observables inside a replaced box
are transported as in \eqref{eq:mrg-history-insertion}.

\subsection{The running source and its local variation}

Let $\varphi_j$ be the smooth cutoff of the level-$j$ action
region $\Lambda_j$, fixed by the action sequence, with
$0\le\varphi_j\le1$, equal to one in its interior and zero
outside. The unsourced inverse-coupling array is
\[
 x_k(p)=\beta-\sum_{j\le k}b_j\varphi_j(p).
\]
Except at the final two layers of a prepared sequence, the
inter-region margins exceed the smoothing length. Thus
$\varphi_{j'}(p)>0$ for $j'>j$ implies $\varphi_j(p)=1$.
If $m$ is the last positive profile, the array is
$(1-\varphi_m(p))x_{m-1}+\varphi_m(p)x_m$; if none is
positive, it is $\beta$.

At the last stage of an operation of level $t\ge2$, the
two final profiles can overlap in the thin shells. Put
$a=\varphi_{t-1}(p)$ and $b=\varphi_t(p)$. Their contribution
is the convex combination
\begin{align*}
 &(1-a)(1-b)x_{t-2}+a(1-b)x_{t-1}+abx_t\\
 &\hspace{30mm}+(1-a)b(x_t+x_{t-2}-x_{t-1}).
\end{align*}
The additional vertex is $x_t+b_{t-1}$. Its positivity
uses the one-step case of \eqref{eq:flow}, namely
$b_{t-1}\ge\lambda_\sharp-C_\sharp$.
For a first-passage prefix define
\begin{equation}\label{eq:revised-profile-minimum}
 x_k^{\min}=\min\bigl(\{x_j:0\le j\le k\}
                \cup\{x_t+b_{t-1}:2\le t\le k\}
                                                \bigr).
\end{equation}
The second set is empty when $k<2$. Since $x_j>g^{-2}$
along this prefix, $x_k^{\min}\ge g^{-2}-C_\sharp$.
Consequently the same positive lower bound controls every
label and stage profile once the coupling ceiling is small.
Convexity alone does not give the lower bound for the
additional vertex.

On the common regular region a constant source has running shift
\[
 \zeta_k(\zeta)=\zeta+\sum_{j\le k}[b_j(\zeta)-b_j(0)].
\]
For a varying source the same marginal coefficients are evaluated at
local values. Let $h_y^{(j)}$ be the nonnegative tents, with
$\sum_yh_y^{(j)}(p)=1$ and physical range at most $CML^{j-K}$.
Choose a bare plaquette $p_y$ for each tent centre, independently of the
orientation of the plaquette $p$ at which the array is evaluated. Then
\[
 \zeta_k(p;w)=w(x(p))+\sum_{j\le k}\varphi_j(p)
       \sum_yh_y^{(j)}(p)[b_j(w(x(p_y)))-b_j(0)].
\]
The sourced Wilson array is $x_k(p)-\zeta_k(p;w)$.

The differentiated marginal estimate required from the induction is
\begin{equation}\label{eq:source-flow}
 |b_{k+1}(\zeta)-b_{k+1}(0)|\le a_{[K-k-1]}|\zeta|,
 \qquad
 \sum_{s<K}a_{[s]}
 \le C\sum_{k<K}\left[g_k^3(\log x_k)^{\mathfrak p_0}
                  +\rho^{-1}g_k^{\kappa_0-4}\right]\le\frac13.
\end{equation}
Its two terms have different origins. The small-field part is analytic
in the earlier inverse-coupling shifts on discs of radii comparable to
$x_j$. Cauchy's inequality divides its nonconstant bound
$Cg_k(\log x_k)^{\mathfrak p_0}$ by a radius of order $x_k$.
The remainder left by a processed component is only known on a fixed
source disc; its size is $Cg_k^{\kappa_0-4}$, and differentiation costs
$\rho^{-1}$. Neither term acquires a factor equal to the number of
earlier data variables: the Schwarz estimate is applied in the maximum
norm of the joint polydisc.

Here is the elementary analytic estimate used in this statement. Suppose
the data coordinates have radii $r_l$, and a holomorphic map $G$ is
bounded by $A$ on the corresponding polydisc. For a displacement $v$
put $m=\max_l|v_l|/r_l$. The function
$t\mapsto G(tv)-G(0)$ is bounded by $2A$ on $|t|<m^{-1}$ and
vanishes at zero. Schwarz's lemma gives
\[
 |G(v)-G(0)|\le2Am\qquad(m<1).
\]
This estimate follows from one complex variable and has no factor
equal to the number of coordinates. The same argument applies at a
nonzero centre if the indicated polydisc about that centre remains in
the analytic class. In the source induction the earlier shifts have
sizes $O(|\zeta|)$; their radii are $cx_j$, and the prefix couplings
are bounded below by a constant multiple of the current $x_k$.
Thus $m\le C|\zeta|/x_k$. Multiplying by
$A=Cg_k(\log x_k)^{\mathfrak p_0}$ gives the first term in
\eqref{eq:source-flow}. The independent remainder is handled on its
fixed disc, not on the growing polydisc.

The lower flow bound $x_k\ge g^{-2}+c(K-k)-C$ makes these sensitivities
summable. For example, comparison with an integral bounds the first sum
by $Cg(1+|\log g|)^{\mathfrak p_0}$. The remainder sum tends to zero
when $\kappa_0>6$. Making $g$ small after $\rho$ gives
\[
 |\zeta_k(p;w)|\le\frac43\sup|w|,
 \qquad \frac23\le\zeta'_k(0)\le\frac43.
\]
The latter derivative is real, since the run is real at real sources.

In particular, at a real source of the class every Wilson
coefficient is at least
$x_k^{\min}-\frac43\sup|w|>0$, once the coupling ceiling
has been chosen after the source radius. This minimum includes
the prepared-stage vertices. Reality of the other factors is
preserved by the real changes of variables and integrations:
the elimination determinants have their positive real values,
the Haar density is positive in the small chart, and the
remaining factors are exponentials of real functions.

\subsubsection{The source data and their comparison with constants}
\label{sec:source-data-induction}

We describe the induction at a first-passage datum, with the choices
of the real run frozen at its own weights $x_j$. Put $r=2\rho$.
General sources are constructed on a class of radius $r$, while
constant sources are constructed on the doubled disc $D(0,2r)$.
The same two radii are used at every level. The inputs from
\cref{sec:rg,sec:large-fields} are the uniform one-step bounds on the
common complex data class, their locality at each carrier, and the
component mass with its stored-volume reserve. The argument below
specifies how these estimates propagate the sources; in particular,
it does not obtain a new source disc by shrinking the previous one.

At depth $s$ use the norms
\[
 \|f\|_{\sharp,s}=\max\{\|f\|_\infty,CM L^{-s}\Lip f\},\qquad
 \|w\|_{\flat,s}=|c|+\sum_i|z_i|\|f_i\|_{\sharp,s}.
\]
The source class used at level $m$ is
$\mathcal W_{r,m}=\{\|w\|_{\flat,K-m}<r\}$.
It decreases as $m$ increases. For a region $Y$ define
\[
 \mathsf l_j(Y;w)=CM L^{j-K}\Lip_Y w,\qquad
 \mathsf m_j(Y;w,\zeta)=
 \max\{\sup_Y|w-\zeta|,\mathsf l_j(Y;w)\}.
\]
The Lipschitz seminorm here is taken over the bare plaquettes in the
enlarged region, with their physical distances. The quantities are
monotone in $Y$ and $j$. On a fixed region,
$\mathsf l_i=L^{-(j-i)}\mathsf l_j$ for $i\le j$.
The tree defining $d_j(X)$ joins the plaquettes of $X^{\sim1}$
within distance $(d_j(X)+C)ML^{j-K}$. With the constant in
$\mathsf l_j$ fixed sufficiently large, this gives
\begin{equation}\label{eq:source-geometric-gain}
 \begin{split}
 \mathsf m_j(X^{\sim1};w,w(p_z))
       &\le e^{d_j(X)}\mathsf l_j(X^{\sim1};w),\\
 \mathsf l_j(Y;w)&\le L^{-(m-j)}\|w\|_{\flat,K-m}
                        \quad(j\le m).
 \end{split}
\end{equation}
Here $z\in X$ and its chosen plaquette $p_z$ lies in the indicated
enlargement. This is the factor that makes the spatial variation of
the source summable over the older levels.

Use the pointed form of the regular activities from \cref{sec:rg}.
At each site $z$ of level $j$, retain the entire constant-source
family $E^{(j),w(p_z)}(X,\mathbf p,z)$, its infinite-lattice
companion, and its own coefficient $b_j(w(p_z))$. Its counterterm
is attached to the same base point through $h_z^{(j)}$. The defect
is
\begin{equation}\label{eq:source-pointed-defect}
 D^{(j),w}(X,\mathbf p,z)
 =E^{(j),w}(X,\mathbf p,z)
                  -E^{(j),w(p_z)}(X,\mathbf p,z).
\end{equation}
The Euclidean conditions on a pointed family and its coefficient
are linear. Families at different sites may therefore carry
different constant sources while each retains the pointwise
coupling cancellation. Keeping only the full sum of the families
would lose this cancellation at the individual base points.

Suppose that, for every constant $\zeta\in D(0,r)$, the regular
activity has the slope
\begin{equation}\label{eq:source-regular-slope}
 |E^{(j),w}-E^{(j),\zeta}|(X,\mathbf p,z)
 \le\frac{\mathsf m_j(X^{\sim1};w,\zeta)}r
                                  e^{-(\kappa+1)d_j(X)}.
\end{equation}
Then \eqref{eq:source-geometric-gain}, with $\zeta=w(p_z)$,
shows that the defect of a source in $\mathcal W_{r,m}$ satisfies
\[
 |D^{(j),w}(X,\mathbf p,z)|
                   \le L^{-(m-j)}e^{-\kappa d_j(X)}.
\]
Defects have no further coupling subtraction. Their geometric
factor is consequently retained from the time they are made;
$\sum_{j\le m}L^{-(m-j)}$ is bounded independently of $m$.
The separate remainder sum is paid by $g_j^{\kappa_0}$.

The other entries of the source data are the bare weight $u=w$,
the remainder families, the boundary terms and the excess
exponents in the stored component operations. On a complex data
class of level $k$ the bounds have the form
\begin{equation}\label{eq:source-complex-data-bounds}
 \begin{gathered}
 \sup|u|+\sum_{j\le k}\sup_z|\widehat b_j(z)-b_j(0)|\le3r,\qquad
 \|G^{(j)}\|\le E_0+1,\qquad
 \|D^{(j)}\|\le L^{-(k+1-j)},\\
 \|H^{(j)}\|\le Cg_j^{\kappa_0},\qquad
 |B^{(j)}(X)|\le C\bar B_j(X),\qquad
 |\Psi_\omega|\le2c_\Phi\mathsf n(\omega).
 \end{gathered}
\end{equation}
The norms include domain decay and the supremum over base points.
The coefficient $\widehat b_j(z)$ is attached to site $z$; for
the actual source it equals $b_j(w(p_z))$. Here $G$ denotes
a constant-source pointed family, $H$ a remainder, and
$\bar B_j$ is the fixed boundary majorant of the real run.
The count $\mathsf n(\omega)$ adds the sites stored at their own
levels, including the inner component histories once each. Bounds
on the oscillations of excess factors are also retained, with the
carrier and hull-length majorants of \cref{sec:large-fields}.
The Gaussian structures, the reference measures and all the numbers
of the operations are those of the real run. The shift enters the
interaction and the Wilson ledger, not those structures.

The summable coefficient estimate \eqref{eq:source-flow} bounds
the first quantity in \eqref{eq:source-complex-data-bounds} for
the actual source data by $\frac43\sup|w|$. The equality of the
bare weights is plaquette by plaquette:
\[
 A(x_k(\cdot)-\zeta_k(\cdot;w),U)
  +\sum_{j\le k}A\!\left(
      \varphi_j\sum_z h_z^{(j)}b_j(w(p_z)),U\right)
                         =A(\beta-w,U).
\]
Thus the data with variable sources represent the original sourced
measure, with the same normalization conventions as the unsourced
run.

\emph{Localized lines of data.}
Fix a region $W$ for the next partition, a source $w$ in its
level class, and an arbitrary constant $\zeta\in D(0,r)$.
Put $\mathsf m=\mathsf m_{k+1}(W;w,\zeta)$.
For $|t|\mathsf m<r$, interpolate the ordinary entries inside
$W$ between the data of $\zeta$ and those of $w$, and leave the
entries outside $W$ at $\zeta$. Interpolate a whole pointed
family and its coefficient at each site. An indexed ordinary
term is changed only when its required enlargement lies in $W$;
a boundary term made by a sum over sites is treated summand by
summand. Denote these localized data by $\mathfrak L_t^W$.

The regular, remainder and boundary slopes towards the same
constant keep these entries inside their complex bounds. For
example, the frozen constant-source family has size at most
$E_0$, and its interpolation costs at most
$|t|\mathsf m/r\le1$. The defects vanish at the constant end
of the line, so at every older level their sizes are bounded by
\[
 |t|r^{-1}\mathsf l_j(W;w)
                       \le L^{-(k+1-j)}.
\]
The coefficient functions are controlled on $D(0,2r)$.
A Cauchy estimate on this doubled disc gives twice their radial
slope on $D(0,r)$; together with
$\sum_j a_{[K-j]}\le\frac13$, this places the interpolated
weight and coefficients within the bound $3r$. All these estimates
use $r/\mathsf m$, independently of the distance of $w$ from the
boundary of its source class.

The exponents of the stored operations require a different
construction. If $W$ cuts a component, direct interpolation of
its factors would select only some carriers. A bound on the total
exponent and bounds on the oscillations of individual factors do
not bound this partial sum: adding constants to those factors and
subtracting their total from the core changes neither of the two
bounds. We instead construct the operation exponents on each line
by the same recursion that constructs them in the RG step.

This recursion uses only older lines. A region for the partition
at level $k+2$ is also a region for the partition at level $k+1$;
the next source class is contained in the preceding class; and
$\mathsf m_{k+1}\le\mathsf m_{k+2}$. Hence the next line's
parameter disc is contained in the older disc on which its input
operation exponents have already been constructed. Apply the
step to those exponents and to the prescribed ordinary data.
For the next comparison, reset the ordinary entries to their
affine patchwork. Only the operation exponents are retained from
this output construction; equality of the entire interpolated
systems at intermediate $t$ is not required.

The exponents so produced are holomorphic in $t$. At $t=0$ they
are those of the constant source. When $\mathsf m<r$, at $t=1$
locality identifies
an excess factor with its value for $w$ whenever its carrier
enlargement lies in $W$, and identifies the core when the twice
enlarged component lies in $W$. These are exactly the endpoint
identities needed for a subsequent local comparison. The uniform
step bounds give
\[
 |\Psi_\omega(t)|\le2c_\Phi\mathsf n(\omega),\qquad
 \sum_{f:R_f=R}\operatorname{osc}\psi_f(t)
                              \le c_\times\bar M_i(R),
\]
where $\bar M_i$ is the fixed carrier majorant. No partial
exponent is estimated by the modulus of its full sum.

\emph{The regular output and its slope.}
The small-field part of the step produces a pointed family
$E^{(k+1)}[\mathfrak L]$ before its constant-source extension
is chosen. Its Gaussian normalizer $N^{\circ(k+1)}$ is common
to every system in the complex class. Here the circle retains
the local values at the trivial pair. The family normalized to
vanish there is
\[
 N^{(k+1)}(X,\mathbf p,z)
 =N^{\circ(k+1)}(X,\mathbf p,z)
                     -N^{\circ(k+1)}(X,\mathbf1,z).
\]
These local numbers are real and independent of the source and
of the running coupling; their field Hessian and marginal
coefficient are zero. In the unnormalized family they are
counted once among the vacuum values. The separate numerical
Gaussian factor counts, at each free bond, the three-dimensional
scaling and the unit Gaussian integral. It does not contain
the determinant of the fluctuation form: that determinant
remains in the localized normalizer. This is the same division
used in \cref{sec:rg}.
The uniform step estimate
has the quantitative form
\begin{equation}\label{eq:source-box-half-bound}
 \left|E^{(k+1)}[\mathfrak L](X,\mathbf p,z)
                     -N^{\circ(k+1)}(X,\mathbf p,z)\right|
 \le C\vartheta_k e^{-\kappa' d_{k+1}(X)}
 \le\tfrac12 e^{-\kappa'd_{k+1}(X)},
 \qquad \kappa'\ge\kappa+1,
\end{equation}
where one can take
\[
 \vartheta_k=
 \frac{\delta_k}{\min_\iota\alpha_{\iota,k}}
             +M^4g_kp_0(g_k)^3+x_k^{-1/4}.
\]
The first term is of order
$(\log x_k)^{\mathfrak p_0-q}$ and tends to zero.
The remaining terms also tend to zero. Since the real run uses
its own Gaussian weights, there is no discrepancy between the
frozen weight and the real weight in this calculation. All
constants are chosen before decreasing the coupling ceiling.

For the activity on $X$, choose $W=X^{\sim1}$.
If $\mathsf m=0$, locality identifies its output with that of the
constant source. Suppose next that $0<\mathsf m<r$ and set
$T=r/\mathsf m>1$. At a fixed point of the field domain the function
\[
 \phi(t)=
 \big(E^{(k+1)}[\mathfrak L_t^W]-N^{\circ(k+1)}\big)
                                      (X,\mathbf p,z)
\]
is holomorphic for $|t|<T$ and has modulus at most
$\frac12e^{-\kappa'd_{k+1}(X)}$ there. Its values at zero and
one are the constant-source and variable-source outputs. Schwarz's
lemma applied to $\phi(t)-\phi(0)$ yields
\[
 |\phi(1)-\phi(0)|
  \le\frac{\mathsf m}r e^{-\kappa'd_{k+1}(X)}.
\]
If $\mathsf m\ge r$, apply the two independent supremum bounds
to the actual source outputs: their difference is at most
$e^{-\kappa'd_{k+1}(X)}$, which is no larger than the same right
side multiplied by $\mathsf m/r$. This case needs no evaluation
of the localized line at $t=1$.
This recovers \eqref{eq:source-regular-slope} at level
$k+1$ with constant one. The subtraction of the common normalizer
is essential: the bound on the full activity is $E_0$, and its
direct use would give slope $2E_0$ rather than the constant needed
by the defect estimate.

Choose the extension at site $z$ to be the output for the constant
$w(p_z)$, with coefficient $b_{k+1}(w(p_z))$.
The difference from the actual output has the required defect
bound by \eqref{eq:source-geometric-gain}. On a localized line the
temporary extension has defect at most two in the same norm,
which is the bound admitted during the large-field operation.
The exact measure identity of the step then gives the sourced
format at the new level. The constant-source extension is formed
after the small-field output and before this last operation; its
definition therefore uses no unknown future remainder.

\emph{Remainders, boundary terms and stored operations.}
The uniform output bound for the new remainder is
\[
 |R^{(k+1),w}(X)|
 \le\left(\tfrac12g_{k+1}^{\kappa_0}
                   +e^{-p_0(g_{k+1})}\right)e^{-\kappa d_{k+1}(X)}
 \le g_{k+1}^{\kappa_0}e^{-\kappa d_{k+1}(X)}.
\]
The last inequality holds below one fixed coupling ceiling.
Apply the same line argument to this remainder and to each new
boundary term. Their slopes are twice their supremum majorants
times $\mathsf m/r$. Carried terms retain their previous slopes.
When boundary terms are sums of pointed activities, the spare
unit of domain decay pays the number of base points in $X$.

For a processed component the new part of its excess exponent has
the expression in \cref{sec:source-component-majorant} below.
Its bound is $c_{\rm new}$ times the number of newly displayed
sites. The older exponents use the additive count
$\mathsf n(\omega')$, so
\[
 |\Psi_\omega|
 \le c_{\rm new}\sum_j\#_j(Z_j(\omega)\cap X)
                 +c_\Phi\sum_{\omega'}\mathsf n(\omega')
 \le\max\{c_{\rm new},c_\Phi\}\mathsf n(\omega).
\]
Thus the nesting takes a maximum of constants. On a localized
line the same calculation returns
$\max\{c_{\rm new},2c_\Phi\}=2c_\Phi$ when
$c_\Phi=2c_{\rm new}$. Applying Schwarz on the twice enlarged
component gives slope $2c_{\rm new}=c_\Phi$ for its new part.
The inner parts retain their slopes towards the same arbitrary
constant $\zeta$, since $\mathsf m$ is monotone in the region.
No change of base constant is made at an inner component.

The oscillation norm of a family of factors is a sum of suprema,
so it is useful to spell out the analytic step. For
$0<\mathsf m<r$, choose finitely
many factors, two allowed field arguments for each, and phases
that align their differences between $t=1$ and $t=0$.
The phased sum of the corresponding differences at $t$ is one
holomorphic function with the carrier majorant as its bound.
Schwarz's lemma gives twice that majorant times $\mathsf m/r$.
Taking the chosen field arguments arbitrarily close to their
suprema, then increasing the finite set of factors, proves the
same slope in the sum-of-oscillations norm. This uses holomorphy
in $t$; stored near variables may remain merely Borel.
For $\mathsf m\ge r$ the two independent supremum bounds give
the same slope with the factor $\mathsf m/r$; for $\mathsf m=0$
the equality follows from locality.

At level zero the sole new datum is $u=w$ and the operation
excesses vanish. At level $k+1$, first construct the constant-source
increment from the coefficients and remainders through level $k$.
The coefficient estimate \eqref{eq:source-flow} keeps its total
sensitivity below $\frac13$. The constant-source whole-lattice
family and its coefficient then supply the extension just used.
The sourced step returns the new regular, remainder, boundary and
operation data with the same bounds, while locality identifies
the variable-source output with the frozen-source output wherever
the source is locally constant. This completes the source-data
induction under the uniform complex one-step estimates. It uses
the same radius at every level and introduces no factor equal to
the number of previous levels, sites, or nested components.

\subsection{Complex sources and the positive component majorant}
\label{sec:source-component-majorant}

On a regular field, the tent decomposition of the action gives
\[
 |A(h_y^{(k)}\zeta_k,\mathbf U)|
   \le C\alpha_{0,k}^2\sup|\zeta_k|
   \le C\rho g_k^2(\log x_k)^{2q}.
\]
This is a small analytic local perturbation. The domination on
real fields uses the complete profile minimum in
\eqref{eq:revised-profile-minimum}. Put
\[
 \varsigma=\frac{4\sup|w|}{3x_k^{\min}}<1.
\]
For a nonnegative plaquette weight $\chi$ and a real $U$,
\begin{equation}\label{eq:revised-source-domination}
 \begin{split}
 e^{-(1+\varsigma)A(x_k(\cdot)\chi,U)}
 &\le\left|e^{-A((x_k(\cdot)-\zeta_k(\cdot;w))\chi,U)}\right|\\
 &\le e^{-(1-\varsigma)A(x_k(\cdot)\chi,U)}\le1.
 \end{split}
\end{equation}
Indeed, $a_p(U)\ge0$ and
$|\operatorname{Re}\zeta_k(p;w)|\le\frac43\sup|w|
\le\varsigma x_k(p)$. Also $\varsigma\le C\rho g^2$
on the fixed source class.
The Wilson factor therefore retains its failed-test penalty. The
Gaussian and gauge-fixing penalties retain their original form because
their defining choices have not changed.

Stored variables require an integrated estimate. Let $\omega$ be the
local history of a processed component and $\Psi_\omega$ its source
excess exponent. A layer stored at level $j$ is regular at that level:
the fourth power of its fine spacing bounds the plaquette action and
cancels the reciprocal fourth power in the plaquette count per stored
site. The Wilson shift, the changes of interactions, and the excesses
of swallowed components consequently have the bound
\[
 |\Psi_\omega|\le C(1+\rho)\#(\omega).
\]
The count includes each stored site at its own level. Counts add when
components are swallowed. One fixed constant can thus be used throughout
the nesting of histories; it is not summed once more over the levels.

For a processed component the excess has three terms:
\[
 \Psi_\omega
  =A(\hat\zeta\phi_{\mathrm c}\mathbf1_X,U''(\xi))
       +[V^w-V^0](X^\sim)+\sum_{\omega'}\Psi_{\omega'}.
\]
Here $\phi_{\mathrm c}$ is the old-core cutoff, $U''(\xi)$ is the
real background reconstructed from stored variables, $V$ is the
interaction in the component operation, and $\omega'$ ranges over
swallowed islands and drop-outs. On the layer stored at level $j$,
regularity gives
\[
 |\partial U''-\mathbf1|
      \le C\check R_j^2\varepsilon_j\eta_j^2\le C\eta_j^2.
\]
Hence $a_p(U'')\le C\eta_j^4$, and multiplication by the
$C\eta_j^{-4}$ plaquettes per stored site proves the Wilson part of
the extensive bound. On the original bare layer compactness gives
$a_p\le2$. The interaction difference is bounded by the source-data
slope, and the last sum is treated by the additive site count. The
new inductive constant is the maximum of the old constant and the
constant for the new layer.

A drop-out is represented differently. Its carrying term retains the
Wilson weight over $X$; the remaining operation consists of the frozen
choices. Its excess exponent contains the swallowed-component excesses,
while the carrying Wilson factor records the source of the drop-out
itself. Allowed evaluations of the background are fibred over the
stored variables: they must be layerwise regular and give the recorded
test values, except for the test whose failure caused the drop-out.
Replacing this set by independent evaluation sets would allow a flat
field and could erase the Wilson penalty. The mass majorant uses
precisely the compatible evaluations.

The reserve in \hyperref[hyp:history-bounds]{the weighted history and live-age estimates} is required in the form
\[
 \mathsf M(i,X)=\frac1{|W_4|}\sum_\omega
       \int e^{c_r\mathsf v(\omega)}
          \sup_{\mathbf p}|e^{\Phi_\omega(\xi;\mathbf p)}|
                                 \,d\mu_\omega(\xi)
       \le e^{-(2-c')p_0(g_i)},
 \qquad
 \mathsf v(\omega)=\sum_{\text{stored sites at }j}\log x_j.
\]
The allowed evaluations $\mathbf p$ respect the stored background data.
The supremum is inside both the positive integral and the history sum.
After choosing $g$ so that $c_r\log x_j\ge2C(1+\rho)$, the modulus of
$e^{\Psi_\omega}$ is absorbed in the reserve. In particular,
\[
 |\mathbf T'_i(X;w)F|\le\|F\|_\infty\mathsf M(i,X).
\]
This estimate is valid although the sourced integrand is complex.

For a whole label, factorization uses pointwise component majorants
depending on disjoint groups of stored variables, as in
\eqref{eq:patch-mass}. Bounds on a conditional integral with an arbitrary
fixed exterior cannot simply be multiplied. The disjoint pointwise
majorants are what permit the positive integrations to factor. The
remaining normalization and missing-vacuum costs are
$O(1+\log x_j)$ per displayed site. The creation deposits and the raised
reserves pay these costs while retaining the stated component gains.

We spell out the resulting label estimate. After replacing the
integrands of processed-component operations by their absolute
majorants at the reserved weights, its form is
\begin{align*}
 &\int\mathbf T'_k(\mathcal S)
       [\chi_{k,\mathcal S}e^{-A(x_k(\cdot),U_k(\mathbf V))}]\,dV_k\\
 &\quad\le
   \exp\left[C\sum_{j\le k}(1+\log x_j)\#_j(Z_j)\right]
      \prod_{(t,\mathfrak c)}e^{-cp_0(g_t)}
      \prod_{(i,X)}\mathsf M(i,X).
\end{align*}
The first product runs over creation cells outside the displayed
processed components. The second runs only over outermost processed
components. Creations swallowed by one such component have already
paid for its mass and cannot be counted a second time. If a label
still displays part of a drop-out's interior, use a reserve constant
larger than the one spent on the source excess. The extra reserve pays
for that displayed interior. This constant is chosen after $\rho$ and
before the upper bound on $g$.

The site count and logarithm in the source charge cost at most
$CM^4\check R_j^5$ per level-$j$ cell. The original creation deposit
is $cp_0(g_j)^2$. The lineage credit estimate in
\cref{sec:large-fields} pays the cumulative charge until the next
component operation. At its parameter choices the weaker surviving
exponent $cp_0(g_j)$ still pays
$CM^4\check R_j^5\log x_j$ with room left. The weighted history
sum thus retains, for every final component $C$, a factor
\[
 \exp[-cp_0(g_k)-\kappa_1|C|_k].
\]
Its exponential cell decay pays the connected-set entropy. These are
uses of the weighted history estimate; an unweighted component mass
bound alone would not give them.

\subsection{Local source energies and the logarithm}
\label{sec:source-energies}

The scalar energy is fixed by the values of the constructed
activities at the trivial pair. To show both its locality and its
normalization, keep the pointed regular family and write
\[
 \Delta E^{(j)}(X,z;w)
    =E^{(j),w}(X,\mathbf1,z)-E^{(j)}(X,\mathbf1,z),\qquad
 \Delta R^{(j)}(X;w)
    =R^{(j),w}(X,\mathbf1)-R^{(j)}(X,\mathbf1).
\]
Set
\begin{equation}\label{eq:source-energy-local-definition}
 e_j(X;w)=\sum_{z\in X}\Delta E^{(j)}(X,z;w)
                         +\Delta R^{(j)}(X;w),\qquad e_0=0.
\end{equation}
The Gaussian normalizer is the same in the two regular families
and cancels in their difference. These values are therefore values
of the interaction output, with the same subtraction as in
\eqref{eq:source-box-half-bound}. In particular, the
coupling-independent local numbers retained by $N^{\circ(j)}$
cancel as well; they are not added a second time to the source
energy. Their unsourced sum remains part of the vacuum energy.

The slopes towards the constant zero give
\[
 |\Delta E^{(j)}(X,z;w)|
 \le r^{-1}\|w\|_{\flat,K-j}e^{-(\kappa+1)d_j(X)},\qquad
 |\Delta R^{(j)}(X;w)|
 \le2r^{-1}\|w\|_{\flat,K-j}g_j^{\kappa_0}e^{-\kappa d_j(X)}.
\]
The number of base points is at most $CM^4(d_j(X)+1)$.
Using $(d+1)e^{-d}\le1$, and decreasing the decay rate by a fixed
fraction when enlarging a domain, gives
\begin{equation}\label{eq:source-local-energy-bound}
 |e_j(X;w)|\le CM^4r^{-1}\|w\|_{\flat,K-j}
                       e^{-\kappa_e d_j(X)},\qquad
 \kappa_e>\kappa_{\mathcal D}.
\end{equation}
Here $\kappa_{\mathcal D}$ is the entropy rate of localization
domains through a fixed cube. Each energy is holomorphic in its
level source class, real at real sources, zero at $w=0$, and
depends only on $w$ in $X^{\sim1}$.

\emph{The scalar share of a label.}
Let $\mathcal Z_j(\mathcal S)$ be the pairs $(z,X)$ of regular
terms retained by a label, and let $\mathcal R_j(\mathcal S)$
be its retained remainder domains. Define
\[
 T_j(\mathcal S;w)=
 \sum_{(z,X)\in\mathcal Z_j(\mathcal S)}
                       \Delta E^{(j)}(X,z;w)
 +\sum_{X\in\mathcal R_j(\mathcal S)}\Delta R^{(j)}(X;w).
\]
Denote by $\mathcal E_k(\mathcal S;w)$ the scalar that appears
with a minus sign in the exponent of its term. The numerical
scaling factors, gauge normalizers and Gaussian constants are
source independent. Subtracting the two scalar constants thus
gives the exact identities
\begin{equation}\label{eq:source-label-share}
 \begin{split}
 \mathcal E_k(\mathcal S;w)
    &=\mathcal E_k(\mathcal S;0)-\sum_{j\le k}T_j(\mathcal S;w),\\
 T_j(\mathcal S_\emptyset;w)&=\sum_Xe_j(X;w),
 \end{split}
\end{equation}
where $\mathcal S_\emptyset$ is the all-small-field label.
Boundary terms carry no additional scalar subtraction.

A term missing from the all-small-field share contains a cube of
the enlarged large-field region $Z_j^+$. Pin it at such a cube
and sum its domain decay there. The tree bound and
\eqref{eq:source-local-energy-bound} give
\begin{equation}\label{eq:source-missing-share}
 |T_j(\mathcal S_\emptyset;w)-T_j(\mathcal S;w)|
                                      \le C\#_j(Z_j^+).
\end{equation}
The count is at the level of the missing term, including its
base-point count. No sum over bare sites of a fixed physical
region is used in this estimate.

For a run stopped at $k$, put
\[
 E_{(k)}=-\mathcal E_k(\mathcal S_\emptyset;0),\qquad
 \mathcal E_{(k)}(w)=\sum_{j\le k}\sum_Xe_j(X;w),\qquad
 \widetilde Z_k(w)=e^{-E_{(k)}-\mathcal E_{(k)}(w)}Z(w).
\]
At $k=K$ these are $E$, $\mathcal E$ and $\widetilde Z$.
The uniform estimates to obtain from the weighted history sum are
\begin{equation}\label{eq:source-stability-floor}
 |\widetilde Z_k(w)|\le e^{C_r(1+\log x_k)\mathsf V_k},\qquad
 \widetilde Z_k(0)\ge e^{-C(1+\log x_k)\mathsf V_k},\qquad
 \mathsf V_k=\#\mathbb T^{(k)}.
\end{equation}
The permitted sources belong to $\mathcal W_{r,k}$. At first
passage $x_K\in(g^{-2},g^{-2}+B_\sharp]$, so the constants in
the final-level estimate are independent of $K$, $\beta$ and
the torus side. Write $C_g=C(1+\log g^{-2})$ after fixing the
source radius.

\emph{Upper stability.}
Take moduli inside every stored-variable integral, using the
positive component majorants of
\cref{sec:source-component-majorant}. The source changes the
non-Wilson part of the label exponent by at most a constant
depending on $r$ per determining bond. Its stored operation
excesses are absorbed by raising the fixed volume reserves.
The weighted label estimate of \cref{sec:large-fields} consequently
has the sourced form
\begin{equation}\label{eq:source-weighted-label-sum}
 \sum_{\mathcal S}
 \big|e^{\mathcal E_k(\mathcal S;w)}\tau^w_{\mathcal S}\big|
                                            (\mathcal V_k)
 \exp\!\left[\sum_{j\le k}
       C(1+\log x_{j-1})\#_j(Z_j^+)\right]
                            \le e^{c(r)\mathsf V_k}.
\end{equation}
The only uncharged extensive count is the determining set of the
last lattice, whose size is at most $C\mathsf V_k$. All older displayed
bonds lie in large-field regions, where the creation payments and
component masses pay their charge. The constant $c(r)$ is
independent of the coupling. The uniform product and history
estimates needed for this summation are the same ones used in
the unsourced construction, with larger fixed reserves.

The remaining normalization is explicit. By
\eqref{eq:source-label-share},
\begin{align*}
 &-E_{(k)}-\mathcal E_{(k)}(w)-\mathcal E_k(\mathcal S;w)\\
 &\qquad=-[E_{(k)}+\mathcal E_k(\mathcal S;0)]
       -\sum_{j\le k}[T_j(\mathcal S_\emptyset;w)
                                      -T_j(\mathcal S;w)].
\end{align*}
The first bracket compares the label's numerical vacuum constant
with that of the all-small-field label. Scaling Jacobians and
gauge normalizers are omitted in its large-field region; their
logarithms cost $C(1+\log x_{j-1})$ per omitted site. The source
shares have the smaller cost \eqref{eq:source-missing-share}.
Their sum is bounded above by the weight already present in
\eqref{eq:source-weighted-label-sum}, together with
$C(1+\log x_k)\mathsf V_k$ at the last lattice. Sum the exact sourced
identity over labels. This proves the upper bound in
\eqref{eq:source-stability-floor}.

\emph{The lower bound at zero source.}
Every unsourced term is positive, so retain the all-small-field
label alone. Restrict its last block field to
$|V(b)-\mathbf1|<cg_k$ at each bond, with $c$ fixed small.
Background regularity places this set in the kept small-field
domain and bounds its Wilson action by $C\mathsf V_k$ after multiplication
by $x_k=g_k^{-2}$. The remaining normalized exponent is also
bounded by $C\mathsf V_k$. A Haar ball in $\SU(2)$ of radius $cg_k$
has volume at least $c'g_k^3$, and there are at most $4\mathsf V_k$
independent bond variables. Its volume therefore costs at most
$C(1+\log x_k)\mathsf V_k$. This proves the lower bound in
\eqref{eq:source-stability-floor}. Balls of the larger threshold
radius $g_kp_0(g_k)$ would allow action $Cp_0(g_k)^2\mathsf V_k$ and
would give the weaker floor with that logarithmic power.

At a fixed depth $s$, stop at $k=K-s$. All source and energy
identities use the same local data through this level. The bound
$x_{K-s}\le g^{-2}+B_\sharp(s+1)$ replaces $C_g$ by $C(g,s)$;
the number of sites is $\mathsf V_{K-s}$. This gives the depth version
without changing a constant according to the total cutoff.

\begin{proposition}[Separated derivatives of the scalar energy]
\label{prop:source-energy-derivatives}
Under these local source-energy estimates, for $n\ge2$ and supports at
mutual distance at least $d>0$,
\[
 |\partial_1\cdots\partial_n\mathcal E(0)|
 \le C\ell^4(1+d^{-4})(1+\log^+d^{-1})
                (n/\rho)^n\prod_i\|f_i\|_\sharp.
\]
If one support lies in a ball of radius $R$, the factor depending on
$\ell$ can be replaced by a constant depending on $R$.
\end{proposition}

\begin{proof}
Locality makes a mixed derivative vanish unless the enlarged domain
meets every support. In particular it must join two of them, so
$d_j(X)\ge dL^{K-j}/M-C$. On the polydisc with radii
$\rho/(n\|f_i\|_\sharp)$, Cauchy's inequality bounds that derivative by
$(n/\rho)^n\prod_i\|f_i\|_\sharp$ times the activity bound.
Summing domains through a fixed cube first uses up part of their
exponential decay. The remaining estimate is
\[
 \begin{aligned}
 &(n/\rho)^n\prod_i\|f_i\|_\sharp\\
 &\qquad\cdot\sum_{j\le K}C\ell^4(L^{K-j}/M)^4
                               e^{-c[dL^{K-j}/M-C]_+}.
 \end{aligned}
\]
For $L^{K-j}/M\le C\max(1,d^{-1})$ there are
$O(1+\log^+d^{-1})$ scales, and the summand is bounded by
$C\ell^4\max(1,d^{-4})$. At the finer scales the exponential beats
the fourth power, and their geometric-scale sum has the same bound.
This proves the first assertion. A nonzero derivative must also meet
the selected support. Anchoring the domain there replaces the
full-volume cube count by
$(CRL^{K-j}/M+C)^4$ and proves the local version. For a single
derivative there is no second support to force a long domain; the
argument therefore gives no cutoff-independent uncentered mean.
\end{proof}

\begin{lemma}[A local logarithm estimate]\label{lem:source-logarithm}
Suppose $h$ is holomorphic on $|z|_c=\sum_ic_i|z_i|<1$, with
$c_i>0$, $h(0)=1$, and $|h|\le Q$. Then $h$ has a logarithm, zero at
zero, on $|z|_c<(Q+1)^{-1}$, and
\[
 |\partial_1\cdots\partial_n\log h(0)|
       \le\log2\,[2(Q+1)n]^n\prod_ic_i.
\]
\end{lemma}

\begin{proof}
The function $(h-1)/(Q+1)$ is bounded by one and vanishes at zero.
Restricting it to each complex line through zero and applying Schwarz's
lemma gives $|h(z)-1|\le(Q+1)|z|_c$. Thus $|h-1|<1$ on the smaller
domain, where the power series for $\log(1+(h-1))$ defines the desired
branch. On its half-radius domain this series has absolute sum at most
$\sum_{m\ge1}2^{-m}/m=\log2$. The polydisc of radii
$[2(Q+1)nc_i]^{-1}$ lies in that domain. Cauchy's formula on a slightly
smaller closed polydisc, followed by a limit of its radii, proves the
estimate.
\end{proof}

Apply the lemma to $h=e^{-\mathcal E}\Gamma$, with
$c_i=\rho^{-1}\|f_i\|_\sharp$ after a harmless enlargement of $Q$.
Near zero,
\[
 \log\Gamma=\mathcal E+\log h.
\]
The energy estimate and the lemma give a cumulant bound independent of
the cutoff. With global stability alone $Q$ is exponential in $\ell^4$,
and the resulting bound is exponential in the volume. The next
comparison removes this volume dependence.

\subsection{Chessboard comparison and one common normalizer}

Partition a torus into cells of side $\mathsf l$, with faces at bond
mirrors and $2\mathsf l\le\ell$. For $F$ localized in one cell let
$F^{[\square]}$ be its mirror copies, conjugated in odd cells, and put
\[
 \|F\|_{\mathrm{ch},\ell}
  =\left\langle\prod_\square F^{[\square]}\right\rangle^{(\mathsf l/\ell)^4}.
\]
The repeated expectation is nonnegative. The reflected
Cauchy--Schwarz inequalities give the chessboard estimate
\cite[Theorems~4.1--4.3]{FILS1978}:
\[
 \left|\left\langle\prod_\square F_\square\right\rangle\right|
          \le\prod_\square\|F_\square\|_{\mathrm{ch},\ell}.
\]
The cyclic and reflected Cauchy--Schwarz properties are supplied by
\hyperref[hyp:gibbs-reflections]{the lattice reflection argument}. For an exponential source the repeated
expression is a quotient of partition functions, and the root replaces
an extensive volume factor by its cost per cell.

We give the Gibbs comparison at fixed spacing. For a bounded real $H$
in one cell, reflection invariance of the Gibbs state and its
conditional equation imply
\begin{equation}\label{eq:source-gibbs-surface}
 \nu(H)\le\log\|e^H\|_{\mathrm{ch},\ell}
                   +C\beta\mathsf l^4a_K^{-3}\ell^{-1}.
\end{equation}
Indeed, sum the reflected copies of $H$ over the cells of a large cube.
Their expectations are all $\nu(H)$. Jensen's inequality bounds the
expectation of this sum by the logarithm of the expectation of its
exponential. In the conditional Gibbs integral, replacing the exterior
configuration by periodic boundary conditions changes only plaquettes
meeting the boundary. There are $O((\ell/a_K)^3)$ of these, and the
Wilson action changes by at most $C\beta$ for each one. Divide by the
number $(\ell/\mathsf l)^4$ of cells to obtain
\eqref{eq:source-gibbs-surface}. The limit $\ell\to\infty$ is taken
at this fixed $K,\beta$.

Here is the linearization needed for bounded cell functions. For real
$G$ and $t>0$ small, set $H=\log(1+tG)$. If
$b_\ell=\|1+tG\|_{\mathrm{ch},\ell}$, then
$b_\ell\ge\exp[\nu(H)-o(1)]$. The seminorm triangle inequality gives
$b_\ell\le1+t\|G\|_{\mathrm{ch},\ell}$. Consequently, first along
any chosen volume sequence and then as $t\downarrow0$,
\[
 \nu(G)\le\liminf_\ell\|G\|_{\mathrm{ch},\ell}.
\]
For complex $F$, choose a phase so that
$\nu(\operatorname{Re}(e^{-i\alpha}F))=|\nu(F)|$ and use the
seminorm bound on the real part. For a product of finitely many cell
functions, put them in one coarse cell and apply this comparison there.
The fine-board chessboard inequality then gives
\begin{equation}\label{eq:source-gibbs-product}
 \left|\nu\left(\prod_\square F_\square\right)\right|
  \le\liminf_{r\to\infty}\prod_\square
                          \|F_\square\|_{\mathrm{ch},\ell_r}.
\end{equation}
The assertion is along every admissible sequence $\ell_r\to\infty$.
This is needed when a holomorphic normalizer is selected along one such
sequence. The right side can also be bounded by the product of the
individual limsups. The surface error has been removed before any
cutoff limit is taken.

If the product vanishes identically its comparison is immediate.
Otherwise choose a common board shift and a power $L^{k_0}$
so that all occupied fine cells lie in one cell $B$ of side
$l'=L^{k_0}\mathsf l$. As $L$ is odd, the coarse face planes are
also fine face planes. Put
\[
 F=\prod_\square F_\square,\qquad
 G=\operatorname{Re}(e^{-i\alpha}F),\qquad
 \Phi=\log(1+tG),\qquad
 0<t\le\tfrac12\|F\|_\infty^{-1}.
\]
Then $\nu(\Phi)\ge t|\nu(F)|-t^2\|F\|_\infty^2$.
There are $N'=(\ell/l')^4$ coarse cells. Expanding the repeated
exponential before estimating it gives
\begin{align*}
 \left\langle\prod_{\widehat\square}
                   (1+tG)^{[\widehat\square]}\right\rangle
  &=\sum_{Y}t^{|Y|}
      \left\langle\prod_{\widehat\square\in Y}
                                      G^{[\widehat\square]}\right\rangle\\
  &\le\sum_Yt^{|Y|}
          \left(\prod_\square
                  \|F_\square\|_{\mathrm{ch},\ell}\right)^{|Y|}
   =\left(1+t\prod_\square
                  \|F_\square\|_{\mathrm{ch},\ell}\right)^{N'}.
\end{align*}
Indeed, each reflected $G$ is one half the sum of a product of copies
of the $F_\square$ and its complex conjugate, with phases of
modulus one. After expanding these two choices in each occupied coarse
cell, every term has one factor in each of distinct fine cells.
The fine-board chessboard estimate bounds its absolute expectation
by the indicated product. The factors $2^{-|Y|}$ cancel the number
of conjugation choices. No convexity inequality has been applied to
the original complex product.

Apply the real Gibbs comparison in the coarse cells and divide its
logarithm by $N'$. It follows that
\[
 t|\nu(F)|-t^2\|F\|_\infty^2
  \le\log\left(1+t\prod_\square
                          \|F_\square\|_{\mathrm{ch},\ell}\right)
       +C\beta(l')^4a_K^{-3}\ell^{-1}.
\]
At this fixed cutoff take the lower limit along any chosen
$\ell_r\to\infty$, divide by $t$, and let $t\downarrow0$.
This proves the product domination with its lower limit. Both
$\|F\|_\infty$ and $l'$ can depend on this finite collection of
functions; their contributions disappear in the stated order of
limits. In particular the proof applies to a normalizer chosen on a
subsequence of tori and does not impose a cutoff-uniform bound on the
unnormalized products.

The same fine-board inequality gives the rule used for columns of
insertions: if $F$ is a product of fine-cell functions inside $B$,
then
\[
 \|F\|_{\mathrm{ch},\ell;B}
          \le\prod_\square\|F_\square\|_{\mathrm{ch},\ell}.
\]
Indeed its reflected coarse array is such a product over fine cells,
and taking the coarse root of the fine-board estimate gives this
inequality. The rule keeps the actual complex fine-cell functions,
rather than replacing them by pointwise moduli.

Now take a source supported in $m$ unit balls a distance at least three
from the faces of one cell. Fix one reference torus $T_{\ell'}$ with
$\ell'\ge2\mathsf l$. For its mirror-periodic continuation $W$, the
finite-torus transfer estimate used in the source comparison is
\begin{equation}\label{eq:source-torus-transfer}
 0\le Z_\ell(W)\le Z_{\ell'}(W)^{(\ell/\ell')^4},
 \qquad \ell\ge\ell'\ge2\mathsf l.
\end{equation}
It applies when $W\circ\theta=\overline W$ at each face and $W$
vanishes on plaquettes having a bond crossing a face. Taking the
chessboard root transfers the numerator to the reference torus.
Upper stability there and the floor on $T_\ell$ leave a vacuum-density
difference. Local vacuum energies agree on nonwrapping domains, and
the decaying wrapping tails bound this difference independently of the
cutoff.

In detail, with $N_\ell=(\ell/\mathsf l)^4$, the exponential
seminorm and transfer inequality give
\begin{align*}
 \|e^{\cO_a(w)}\|_{\mathrm{ch},\ell}
  &=\left[\frac{Z_\ell(W)}{Z_\ell(0)}\right]^{1/N_\ell}\\
  &\le\exp\left\{(\mathsf l/\ell')^4
                         \operatorname{Re}\mathcal E_{\ell'}(W)
        +\mathsf l^4\left[
             C_g+\frac{E_{K,\ell',\beta}}{\ell'^4}
                    -\frac{E_{K,\ell,\beta}}{\ell^4}\right]\right\}.
\end{align*}
Upper source stability is used only in the numerator on $T_{\ell'}$;
the unsourced floor is used in the denominator on $T_\ell$. The
vacuum densities may each diverge with the cutoff. Their difference
is what the wrapping-domain estimate bounds. Estimating them
separately would lose the comparison.

To define the normalizer, first divide the localized scalar energies
according to their enlarged domains. A domain belongs to the first
class if $X^{\sim1}$ is contained in one board cell; all other
domains belong to the second class. Terms on which the source is zero
vanish, since $e_j(X;0)=0$. Thus a term in the second class can
contribute only if its enlarged domain meets both a source ball and
the boundary of a cell. This division is made for the scalar energies
themselves, before their absolute values are taken.

Write $N'=(\ell'/\mathsf l)^4$. On the reference torus, the
within-cell part of $\mathcal E_{\ell'}(W)$ is a sum over the
$N'$ cells and the levels $j$. For an even cell its summand is
holomorphic in the original variables $z_1,\ldots,z_n$. For an
odd cell it is holomorphic in $\bar z_1,\ldots,\bar z_n$.
Conjugate the value of every odd-cell summand and then divide the
sum by $N'$. The result is $\mathcal E^\circ(z)$. For example,
if the odd-cell summand is
\[
 \phi(\bar z)=\sum_\alpha c_\alpha\bar z^\alpha,
 \qquad
 \overline{\phi(\bar z)}
                  =\sum_\alpha\bar c_\alpha z^\alpha.
\]
The outer conjugation changes no real part and restores holomorphy
jointly in every source variable. Also $\mathcal E^\circ(0)=0$.
Locality is used here in its full form: when $X^{\sim1}$ is in
one cell, the energy depends on the source in that cell alone. No
values of the mirror source on neighbouring cells enter its summand.

This definition gives the exact real-part cancellation
\[
 \frac1{N'}\operatorname{Re}\mathcal E_{\ell'}(W)
              -\operatorname{Re}\mathcal E^\circ(z)
  =\frac1{N'}\operatorname{Re}
       \sum_j\sum_{X:\ X^{\sim1}\text{ meets a cell face}}
                                                   e_j(X;W).
\]
The individual within-cell sums can grow with $K$. In particular,
their first source derivatives contain the cutoff-dependent mean.
They disappear from this difference identically, so they never have
to be estimated by a cutoff-independent sum of moduli.

We estimate the right side by charging a nonzero face domain to a
source ball that it meets. If there are several choices, one can
charge it to all of them; this only enlarges the upper bound. Put
$t_j=L^{K-j}/M$. A ball of fixed radius contains at most
$(Ct_j+C)^4$ cubes of the level-$j$ partition that can serve as
anchors for such a domain. Its distance from the faces is at least
three, so a face domain through one of these anchors has
\[
 d_j(X)\ge ct_j-C.
\]
Use part of $e^{-\kappa_ed_j(X)}$ to sum domains through their
anchor, and retain the rest for this joining distance. The tree
bound then gives the contribution
\[
 C(Ct_j+C)^4e^{-c[ct_j-C]_+}
\]
for one ball at this level. The $t_j$ form a geometric sequence.
The polynomial factor is dominated by the exponential at its large
values; only a bounded number of the remaining values lie below any
fixed threshold. Thus the full cost per ball is
\[
 C_{\mathrm{face}}
 \le C\sum_{i\ge0}(CL^i/M+C)^4
                           e^{-c[cL^i/M-C]_+}<\infty.
\]
The count over reflected cells is $N'$ times this bound for each
of the $m$ source balls. Division by $N'$ leaves at most
$mC_{\mathrm{face}}$. The estimate is independent of the cutoff
and of the number of cells on the reference torus.

Subtract $\operatorname{Re}\mathcal E^\circ$ in the preceding
seminorm comparison. The face estimate and the bounded difference of
vacuum densities give
\begin{equation}\label{eq:source-common-normalizer}
 |e^{-\mathcal E^\circ(z)}\Gamma_{\mathfrak s}(z)|
       \le\mathsf Q,
 \qquad \mathsf Q=e^{mC_{\mathrm{face}}+C_g\mathsf l^4},
 \qquad \|w\|_\flat<\rho.
\end{equation}
The constant $C_g$ absorbs the bounded vacuum-density difference.
Here $\mathfrak s$ is any torus of side at least $\ell'$ or any
periodic Gibbs state at the same first-passage datum. For a Gibbs
state, apply \eqref{eq:source-gibbs-product} to the exponential
source and then the same finite-torus estimate. The normalizer was
constructed on $T_{\ell'}$ and is unchanged when the volume or the
state changes.

There are two distinct uses of complex variables in this argument.
The mirror continuation $W$ contains both $z$ and $\bar z$;
$\mathcal E_{\ell'}(W)$ itself need not be holomorphic in $z$.
Its real part enters the bound and is compared with the real part
of the holomorphic function $\mathcal E^\circ$. In contrast,
the original generating function $\Gamma_{\mathfrak s}(z)$ is
holomorphic in $z$ at fixed cutoff. Its insertion is a bounded
function of finitely many group variables, even in an infinite-volume
Gibbs state. Hence the normalized generating function in
\eqref{eq:source-common-normalizer} is holomorphic on the original
joint source domain. The modulus comparison has not changed the
variables in which its derivatives are taken.

The whole finite source list is used when
$\mathcal E^\circ(z_1,\ldots,z_n)$ is formed. A within-cell
energy may depend on several of these variables and retains that
joint dependence. In particular the construction does not replace
it by a sum of one-variable normalizers. Restricting some variables
to zero gives the normalizer for that restricted source in the same
analytic function. Applying the local logarithm estimate to two
states therefore produces
\[
 \partial_1\cdots\partial_n\log\Gamma_{\mathfrak s}(0)
  =\partial_1\cdots\partial_n\mathcal E^\circ(0)
       +\partial_1\cdots\partial_n
          \log(e^{-\mathcal E^\circ}\Gamma_{\mathfrak s})(0).
\]
The first term is exactly the same for the two states. Its mixed
derivative, including any cutoff-dependent part, cancels on
subtraction. Agreement only on the coordinate axes would not imply
this cancellation of mixed derivatives.

\begin{proposition}[Differences of means and cumulants]
\label{prop:source-state-difference}
For tests supported in a fixed unit ball and any two allowed states or
tori at one first-passage datum,
\[
 |C_n[\mathfrak s](F_1,\ldots,F_n)
       -C_n[\mathfrak s'](F_1,\ldots,F_n)|
       \le C_n(g)\prod_i\|F_i\|_\sharp,\qquad n\ge1.
\]
No separation or sign condition is required. In particular, for $f$ in
a ball of radius $R$ and two periodic Gibbs states,
\[
 |\nu(\cO_a(f))-\nu'(\cO_a(f))|\le C_1(R)\|f\|_\sharp.
\]
\end{proposition}

\begin{proof}
Apply \cref{lem:source-logarithm} to both sides of
\eqref{eq:source-common-normalizer}, on the same source domain. Each
cumulant is the corresponding derivative of $\mathcal E^\circ$ plus
a uniformly bounded derivative of a local logarithm. Subtraction
cancels the derivative of the normalizer exactly. This proof includes
$n=1$, although either individual mean can diverge with $K$. A fixed
tent partition into unit balls proves the radius-$R$ version.
\end{proof}

\subsection{A uniform bound at each fixed order}

The per-ball bound used for common extraction and the time-translation
semigroup requires independence from the positions of the support
balls. Their diameter as one configuration can be arbitrarily large.
There is consequently no single bounded torus containing all possible
configurations to which the global logarithm estimate can be applied.
The chessboard estimate treats the occupied cells separately.

\begin{proposition}[The per-ball correlation bound]\label{prop:per-ball-bound}
For $n\ge2$, supports at mutual
distance at least $d>0$, and each support contained in a radius-$R$
ball, one has
\[
 |C_n(f_1,\ldots,f_n)|
       \le C_n^\infty(d;R)\prod_i\|f_i\|_\sharp.
\]
The constants are independent of the cutoff and the positions. The
torus version requires $\ell\ge Cn(R+1)$; in periodic Gibbs states
there is no volume restriction. They may be chosen nonincreasing in
$d$.
\end{proposition}

\begin{proof}
Take $\mathsf l$ to be the least power of $L$ larger than
$C n(R+1)$, with $C$ fixed sufficiently large, and require
$2\mathsf l\le\ell$. Vary the offset of the board through one
cell. For each coordinate, a given ball excludes an interval of
offsets of length at most $C(R+1)$ if it is to stay a distance at
least $R+2$ from every cell face. The union for $n$ balls excludes
at most $Cn(R+1)$ in that coordinate. The choice of $\mathsf l$
therefore leaves a permissible offset in all four coordinates.
On the lattice one counts offsets in steps of $a_K$; the extra endpoint
cost is at most $Cna_K$ and is absorbed by the same choice. This is
one shift of the whole board.

Group the sources in each occupied cell and denote their sum by
$w_\square$. Several of the original balls may lie in one cell;
we keep their source variables together there. There are $m\le n$
occupied cells. The torus chessboard bound and source stability give
\begin{equation}\label{eq:per-ball-cell-product}
 |\Gamma(z)|
   \le\prod_{\square\ \mathrm{occupied}}
       \exp\left[(\mathsf l/\ell)^4
                   \operatorname{Re}\mathcal E(w_\square^{\mathrm{ch}})
                      +C_g\mathsf l^4\right].
\end{equation}
The norm restriction is still the original joint source restriction;
each cell source uses a subset of its variables. Its reflected array
has norm at most twice that restriction.

We next separate the divergent one-copy energies from the bounded
interaction between copies. For pieces $v_\iota$ in balls of radius
$R$ at distance at least $D>0$, locality gives
\[
 \left|\mathcal E\left(\sum_\iota v_\iota\right)
                         -\sum_\iota\mathcal E(v_\iota)\right|
            \le C(D,R)\#\{\text{balls}\}.
\]
For a domain meeting no more than one piece the summand is exactly
zero. A remaining domain must join two pieces and has length at
least $cD L^{K-j}/M-C$ at level $j$. If it meets $r$ pieces,
the difference of its joint and individual energies is bounded by
$(r+1)Ce^{-\kappa_ed_j(X)}$ and can be charged to those $r$
pieces. For one charged ball the sum is bounded by
\[
 C\sum_{j\le K}(CRL^{K-j}/M+C)^4
                 e^{-c[DL^{K-j}/M-C]_+},
\]
which is finite uniformly in the cutoff. This explains both the
bounded defect and its linear dependence on the number of pieces.
In the reflected array, different occupied copies lie on opposite
sides of a face with the fixed buffer $R+2$, so this defect has
a bound depending on $R$ only.

Use the cell averages of the one-copy energies as a holomorphic
normalizer $\mathcal E_{\underline f}$. As before, odd-cell outer
conjugation restores holomorphy without changing the real parts.
Subtracting it in \eqref{eq:per-ball-cell-product} cancels every
one-copy energy exactly; the copy-interaction defect remains bounded.
We obtain
\[
 |e^{-\mathcal E_{\underline f}(z)}\Gamma(z)|
        \le Q(n,R)^m\le Q(n,R)^n,
 \qquad Q(n,R)=e^{C(R)+C_g\mathsf l^4}.
\]
The scalar normalizer has the ball form of the mixed-derivative
estimate in \cref{prop:source-energy-derivatives}. In the reflected
copies, a domain producing all the source derivatives must still
meet every original support after folding back into the cell.
Reflection folding does not increase distances, so the original
separation $d$ is a lower bound on the joining distance used in
that estimate.

Apply \cref{lem:source-logarithm} to this normalized function and
add the scalar mixed derivative. This proves the claimed fixed-order
bound. Replacing $d$ by a smaller separation in all the preceding
estimates gives nonincreasing constants.

For a periodic Gibbs state use \eqref{eq:source-gibbs-product} and
construct the same normalizer on a single auxiliary reference torus.
Choose it large enough to contain the finite configuration with room
to spare. The finite-torus transfer inequality brings all repeated
numerators to that reference torus, and the bounded vacuum-density
difference adds a fixed cell cost. The source-energy and copy-defect
estimates are local. Thus the resulting constant depends on
$n,d,R,g$ but not on the auxiliary side or the diameter of the
configuration. This proves the state version.
\end{proof}

This bound is sufficient for extraction and for translated
self-pairings. Its growth in $n$ is much larger than factorial:
$\mathsf l\asymp n(R+1)$, while the logarithm bound involves
$(2Q(n,R)^n+2)^n$. The graded estimate below uses unit cells for
all orders and records one factor for each active source, which avoids
this loss.

\subsection{Graded source bounds and factorial growth}

Use unit cells. Let each $v_i$ be supported in a ball of radius $\eta$,
with $\eta\ge4a_K$, and at distance at least $\eta$ from their faces.
Assume that the different supports have mutual distance at least $d>0$.
The source induction requires a holomorphic function $\mathcal E''$,
zero at zero, for which
\[
 h(z)=e^{-\mathcal E''(z)}
                 \left\langle e^{\sum_i z_i\cO_a(v_i)}\right\rangle
\]
has the graded bounds
\begin{equation}\label{eq:graded}
 |h(z)|\le Q(\eta)^{|B|}\quad\text{when }z_i=0\ (i\notin B),
 \qquad Q(\eta)\le C(g)\eta^{-p_1},
 \qquad \sum_i|z_i|\|v_i\|_\sharp<\rho.
\end{equation}
The exponent $p_1$ is independent of $n$. The same energy locality
estimate controls the mixed derivatives of $\mathcal E''$.

\subsubsection*{The source induction.}
\phantomsection\label{hyp:sourced-estimates}
The construction in \cref{sec:source-data-induction} propagates the
source estimates together on $\|w\|_\flat<2\rho$, with the real
choices fixed, once the uniform complex one-step estimates hold. The
inductive data include the local analytic activities, their scalar
energies, the stability bound and floor, and the differentiated
coefficient estimate \eqref{eq:source-flow}. For the two-point comparison
we additionally require \eqref{eq:witness-source-slope} for the full
renormalized first derivative, including its residual quadratic
component, on these same complex source domains. Spatial source
differences are measured relative to a local constant, and the stored
excesses use the raised history reserves. These are parts of one local induction, using the
same domains and coefficient sequence as the unsourced construction.

For the volume comparison use \eqref{eq:source-torus-transfer} on its
compatible cells and tori: the reflected source satisfies
$W\circ\theta=\overline W$ at each face, and vanishes on plaquettes
with a bond crossing a face. The graded estimate \eqref{eq:graded}
then records the cost of each active source with one holomorphic
normalizer. Its local mixed-derivative bounds, the exponent $p_1$,
and the constants in the induction are uniform in the number of
insertions, cutoff, torus side and positions, within the displayed
source-norm and support restrictions.

The cancellation underlying \eqref{eq:graded} can be organized
on the fixed source domain, rather than separately for its coordinate
faces. Put $c_i=\|v_i\|_\sharp$ and write
$w_B=\sum_{i\in B}z_iv_i$. The common restriction
$\sum_i|z_i|c_i<\rho$ gives
\[
 \|w_B\|_\flat\le\sum_{i\in B}|z_i|c_i<\rho,
 \qquad
 \|(w_B|_\square)^{\mathrm{ch}}\|_\flat<2\rho
             \quad\text{for every occupied cell }\square.
\]
For the last inequality the copies of that cell source are grouped
into their two parity arrays; their number does not enter a supremum
norm. Every
restriction to an active set $B$ therefore stays in the same
source class used to construct the scalar energies.

Group the active pieces in their occupied unit cells. In each cell
keep its entire source, including all its active variables, when
forming the within-cell part of the normalizer. A localized energy
can meet several pieces in the same cell. Its joint dependence is
cancelled with that same energy in $\mathcal E''$; decomposing
the normalizer further into one term per source would leave these
mixed contributions uncancelled. As with
$\mathcal E^\circ$, outer conjugation of the odd-cell terms
makes this a single holomorphic function of all variables. On the
face $z_i=0$ for $i\notin B$, every term whose enlarged domain
misses all active supports is zero. Thus the restriction
$\mathcal E''|_B$ makes the cancellation for $w_B$ without
constructing a new normalizer.

Only the domains meeting a face remain in the normalized modulus
bound. Charge one to an active piece whose support it meets.
For a piece of radius $\eta$ at distance at least $\eta$ from
the faces, a domain contributing at level $j$ has
\[
 d_j(X)\ge c\eta t_j-C,
 \qquad t_j=L^{K-j}/M.
\]
There are at most $(C\eta t_j+C)^4$ possible anchor cubes near
that piece. After the tree sum through the anchor, the cost per
piece is bounded by
\[
 C\sum_{j\le K}(C\eta t_j+C)^4
                            e^{-c[c\eta t_j-C]_+}.
\]
At the fine scales $\eta t_j\ge1$ this is a convergent
geometric-scale sum of exponentially decaying quantities. At the
coarser scales the anchor count is bounded and each scale costs at
most a constant. There are
$O(1+\log(M/\eta))$ such scales, independently of $K$.
Consequently
\[
 C_{\mathrm{face}}(\eta)
                   \le C(1+\log(M/\eta)).
\]
If a domain can be charged to several active pieces, the sum of its
charges overcounts it. This yields an upper bound proportional to
$|B|$, without a factor for inactive variables. The same counting
applies to the difference between a joint local energy and its
cell contributions: a domain meeting $r$ contributing cells costs
at most $(r+1)Ce^{-\kappa_ed_j(X)}$, which is bounded by twice
the sum of its $r$ charges.

The stability cost of a unit cell is $C_g$. There are at most
$|B|$ occupied cells. Within-cell energies have already cancelled,
and the face costs just found give
\[
 \log|e^{-\mathcal E''(z_B)}\Gamma(z_B)|
      \le |B|\bigl[C_g+2C_{\mathrm{face}}(\eta)\bigr].
\]
This is \eqref{eq:graded}, with
$Q(\eta)=\exp[C_g+2C_{\mathrm{face}}(\eta)]$.
Since $M$ is fixed before the small-coupling ceiling, the logarithmic
face cost gives $Q(\eta)\le C(g)\eta^{-p_1}$ with an exponent
independent of $g,n$ and the positions. In the case $B=\varnothing$
the source, scalar energy and normalized logarithm all vanish, and
the normalized generating function equals one.

This calculation requires the uniform local source-energy estimates
for the whole reflected array and their compatibility under restriction
to every active list. Their propagation uses the same complex data construction as the
regular, boundary and stored-excess terms, including its recursively
produced operation exponents. The uniform one-step domain and
majorant estimates for that construction remain the analytic inputs. The scalar terms cancel as in the construction of
\eqref{eq:source-common-normalizer}, before the face costs are summed.

\begin{proposition}[From graded estimates to cumulant growth]
\label{prop:graded-factorial}
Once the graded source estimates have been established, there are fixed
exponents $p_2,q_0$ such that, for $n\ge2$, tests $f_i\in C_c^1(\R^4)$
supported in radius-$R$ balls, $R>0$, with mutual support distance
at least $d>0$, satisfy
\begin{equation}\label{eq:local-growth}
 |C_n(f_1,\ldots,f_n)|
   \le A(d;R)B(R)^n(n!)^{p_2}\prod_i\|f_i\|_\sharp,
 \qquad A(d;1)\le C\max\{1,d^{-q_0}\}.
\end{equation}
The estimate holds at every cutoff, on all tori above a fixed volume
floor and in every periodic Gibbs state. Its constants are independent
of the positions and the diameter of the configuration.
\end{proposition}

\begin{proof}
Write $c_i=\|v_i\|_\sharp$. For a nonempty block $B$, set all other
sources equal to zero and apply Cauchy's formula with radii
$\rho/(|B|c_i)$. The graded estimate gives
\[
 |\partial_Bh(0)|\le(|B|Q/\rho)^{|B|}\prod_{i\in B}c_i.
\]
Since $h(0)=1$, its logarithmic derivative is the finite algebraic
identity
\[
 \partial_1\cdots\partial_n\log h(0)
  =\sum_\pi(-1)^{|\pi|-1}(|\pi|-1)!
                        \prod_{B\in\pi}\partial_Bh(0).
\]
There is no use here of a zero-free domain whose radius is uniform in
$n$. For completeness, its combinatorial sum can be evaluated by an
ordinary generating function. If $a_m=m^m$, then $a_m/m!\le e^m$,
and the exponential formula for partitions gives
\begin{align*}
 \sum_{\pi}(|\pi|-1)!\prod_{B\in\pi}|B|^{|B|}
 &=n![t^n]\left[-\log\left(1-\sum_{m\ge1}
                                      \frac{m^m}{m!}t^m\right)\right]\\
 &\le n![t^n]\left[-\log\left(1-\frac{et}{1-et}\right)\right]\\
 &=(n-1)!e^n(2^n-1).
\end{align*}
All series in this coefficient comparison have nonnegative
coefficients. It follows that
\begin{equation}\label{eq:cumulant-combinatorics}
 |\partial_1\cdots\partial_n\log h(0)|
       \le(n-1)!2^n(eQ/\rho)^n\prod_ic_i.
\end{equation}
The power of $Q$ is exactly $n$, because each source index appears in
one block of the partition.

Suppose first that $a_K\le c/n$, with the absolute constant chosen
small. Use a tent partition at a mesh $\eta$ comparable to $n^{-1}$
and not smaller than $4a_K$. A function in a radius-$R$ ball has at
most $(CnR+C)^4$ pieces. Each piece has sharp norm at most
$CMn\|f_i\|_\sharp$. For each resulting $n$-tuple, choose one
common lattice shift of the board which keeps all its pieces away from
unit-cell faces. To see the choice, vary the board offset over one
unit cell. A fixed piece rules out a face strip of width $C\eta$ in
each coordinate. The union for $n$ pieces has relative measure at
most $Cn\eta$. The proportionality constant in $\eta$ can be made
small enough that a permitted offset remains. Counting offsets on
the lattice instead of in the continuous cell gives the same estimate
with an additional $Cna_K$, still small. This shifts the board, or
equivalently all entries by the same lattice vector; it does not
translate the entries independently.

There are at most $(CnR+C)^{4n}$ tuples of pieces. In
\eqref{eq:cumulant-combinatorics}, the norm cost is $(CMn)^n$ and
the graded cost is $Q(\eta)^n\le C(g)^nn^{p_1n}$. The elementary
inequality $n^n\le e^nn!$ absorbs all these costs in a geometric
factor times a fixed power of $n!$.

The remaining term is $\partial_1\cdots\partial_n\mathcal E''$.
Its locality estimate, as in
\cref{prop:source-energy-derivatives}, has a fixed separation power,
and its factor $n^n$ is another fixed factorial cost. The logarithm
and scalar-energy contributions together give
\eqref{eq:local-growth}. In particular the separation exponent does
not acquire a factor $n$ at the level of connected functions.

If $a_K>c/n$, no cutting is needed. The elementary bound
$0\le a_p\le2$ gives
\[
 |\cO_a(f_i)|\le C(R/a_K+1)^4\|f_i\|_\infty
                  \le C(nR+1)^4\|f_i\|_\infty.
\]
For bounded variables $|X_i|\le b_i$, the same partition formula,
with moments bounded by $\prod_{i\in B}b_i$, gives
$|\kappa(X_1,\ldots,X_n)|\le(n-1)!2^n\prod_ib_i$.
This has the required factorial form, uniformly in the bare coupling.

Finally, for a periodic Gibbs state construct the normalizer on one
auxiliary torus containing the whole configuration with room to spare.
The within-cell cancellation and face costs are local, so their
constants do not depend on that side or on the diameter of the
configuration. Apply \eqref{eq:source-gibbs-product} on the selected
volume sequence. The same graded bound and preceding algebra then
prove the state estimate.
\end{proof}


\section{The positive two-point function}\label{sec:two-point}

The calculation has three parts. First an exact source identity is
read at two observation levels. The local calculation is then compared
with a free curvature covariance on one reference torus. Finally,
the common source normalizer transfers the estimate to periodic Gibbs
states. The running coupling fixes the logarithmic scale, while the
live-age estimate controls contributions from stored finer fields.

\subsection{The test pair and the observation levels}

Fix $\vartheta>0$. Let $f,h$ be nonzero, nonnegative Lipschitz
functions supported in balls of radius $d$, with the distance between
their supports in $[d,6d]$, and suppose
\[
 \int f\ge\vartheta\|f\|_\infty d^4,
 \qquad \Lip f\le\frac{\|f\|_\infty}{\vartheta d},
\]
and the same inequalities for $h$. We use
\[
 I(f,h)=\iint f(x)h(y)|x-y|^{-8}\,dx\,dy.
\]
Every pair of points in the supports is at distance at most $10d$,
and all are at distance at least $d$. The mass lower bounds and the
volume upper bounds therefore give
\[
 c\vartheta^2\|f\|_\infty\|h\|_\infty
       \le I(f,h)\le C\|f\|_\infty\|h\|_\infty.
\]
The conclusion required later is
\begin{equation}\label{eq:two-point-comparison}
 c(g,\vartheta)\frac{I(f,h)}{(\log(1/d))^2}
 \le C_2(f,h)\le
 C(g,\vartheta)\frac{I(f,h)}{(\log(1/d))^2},
 \qquad 0<d\le d_0(g,\vartheta),
\end{equation}
with $c(g,\vartheta)>0$ and $g$ small for the fixed shape.

Choose the smallest $s=s(d)$ for which
\[
 C_W(s)L^{-s}\le d,
 \quad C_W(s)=C_{W,0}p_0(\bar x_s^{-1/2})^2,
 \quad \bar x_s=g^{-2}+B_\sharp(s+1).
\]
Require $K-s\ge n_0\ge1$ and set
\[
 k_0=K-s,\qquad \hat d=L^sd,\qquad
 x=x_{k_0}=g(d)^{-2}.
\]
Minimality and monotonicity of $C_W$ give
$C_W(s)\le\hat d<LC_W(s)$. Thus the separation in block units
tends to infinity, while $\log C_W(s)=o(s)$. In particular,
\[
 \frac{s\log L}{\log(1/d)}\longrightarrow1.
\]
The flow estimate and the terminal coupling window imply, uniformly in
first-passage data,
\[
 x_{K-s}=x_K+c_0s+o(s),\qquad
 \frac{g(d)^{-2}}{\log(1/d)}\longrightarrow\frac{11}{3\pi^2}.
\]
The source response obeys
\[
 \zeta'_{k_0}=1+\sum_{j\le k_0}b_j'(0),\qquad
 |\zeta'_{k_0}-1|
   \le C\left[g(d)p_0(g(d))+
                    \rho^{-1}g(d)^{\kappa_0-6}\right]=o(1).
\]
These bounds have no volume dependence.

The second observation level is
\[
 k_1=k_0+\mathsf n,\qquad
 \mathsf n=\lceil2\log_L\bar x_s\rceil.
\]
First fix $g$, and then take $d$ small so that $\mathsf n\le s$.
This ensures $k_1\le K$. The lower and upper flow envelopes give
$\bar x_s\le Cx$ and $x_j\asymp x$ for $k_0\le j\le k_1$,
since $\mathsf n=O(\log x)=o(x)$.

At physical depth $s$ the relevant source norm is
\[
 \|f\|_{\sharp,s}
       =\max\{\|f\|_\infty,CML^{-s}\Lip f\}
       =\|f\|_\infty
\]
when $\vartheta C_W(s)\ge CM$. At level $k\ge k_0$ this norm is
at most $L^{k-k_0}\|f\|_\infty$. Each source derivative at $k_1$
can consequently cost $L^{\mathsf n}$. This loss determines how
many extra levels are needed.

These choices separate three requirements. The scale $k_0$ leaves
$\hat d\to\infty$ lattice units between the supports, enough for
exponential localization errors to beat powers of the running coupling.
The difference $k_1-k_0$ is only logarithmic in $x$, so the coupling
does not change by a fixed fraction on this segment. Nevertheless its
volume rescaling is large: $L^{\mathsf n}\asymp\bar x_s^2$.
Four powers of that length will be gained by translated pairs, whereas
only two are lost in source differentiation. Making $d$ small at fixed
$g$ ensures all three requirements and $k_1\le K$ simultaneously.
They impose a depth floor, not an additional limit $g\to0$ during
the calculation.

\subsection{One exact identity at both levels}

Sum histories with the same standing part before writing their effective
exponent. The positive unsourced weight gives a probability law
$\mathsf P^{(k)}$ on the standing history, the block field, and stored
variables. Its standing-history marginal is the image of the unsourced
history law. Let $\Phi^{(k)}(z)$ be the sourced exponent minus its
unsourced value. Then $\Phi^{(k)}(0)=0$ and frozen source transport gives
\[
 \Gamma(z)=\mathbb E_{\mathsf P^{(k)}}e^{\Phi^{(k)}(z)}.
\]
Write $\Phi_i^{(k)}=\partial_i\Phi^{(k)}(0)$ and
$\Phi_{12}^{(k)}=\partial_1\partial_2\Phi^{(k)}(0)$.
The exact two-point formula is
\begin{equation}\label{eq:witness-exact-identity}
 C_2(f_1,f_2)
  =\operatorname{Cov}_{\mathsf P^{(k)}}
                   (\Phi_1^{(k)},\Phi_2^{(k)})
       +\mathbb E_{\mathsf P^{(k)}}\Phi_{12}^{(k)}.
\end{equation}
Indeed, $\partial_i\Gamma(0)=\mathbb E\Phi_i$ and
$\partial_{12}\Gamma(0)=\mathbb E(\Phi_1\Phi_2+\Phi_{12})$;
subtracting $\partial_1\Gamma(0)\partial_2\Gamma(0)$ proves the
identity. It applies at every level under consideration. The mixed
derivative is present even when the original source is linear.

There are two joint measures in this argument, with different variables.
We describe their common marginal. Fix a class $[h]$ of histories with
one standing part, and let $\xi$ denote the fields still stored at
level $k$. Its unsourced effective measure has the form
\[
 \tau_{[h]}(dV)=\int e^{\mathcal A_{[h]}(V,\xi)}
                                  \,d\lambda_{[h]}(V,\xi),
\]
where $\lambda_{[h]}$ is the positive measure formed from the kept
gauge-fixing factors, characteristic functions, averaging constraints
and Haar measures. Summing the displayed integrand over classes and
dividing by $Z(0)$ defines $\mathsf P^{(k)}$. In particular,
\[
 \mathsf P^{(k)}([h],dV,d\xi)
       =Z(0)^{-1}e^{\mathcal A_{[h]}(V,\xi)}
                                      d\lambda_{[h]}(V,\xi).
\]
The exponent belongs to the sum over histories in the class. Healed
histories have already been summed into it; assigning this same
exponent to each separate history would count them again.

Frozen kernels also give the measure on the bare configuration, an
individual history and its final block field. Project this measure
to $([h],V)$, and test it on gauge-invariant functions of $V$.
Its value is $Z(0)^{-1}\tau_{[h]}(dV)$, the same as the marginal of
$\mathsf P^{(k)}$. Projecting further to $[h]$ gives the image of
the original history probability $\widetilde\nu$. Thus an event of
the standing regions or their lineage ages has the same probability
in all these descriptions. A history-independent gauge-invariant
function of $V$, such as $X_i$ below, has the same marginal as well.
The stored-field source excesses are evaluated in $\mathsf P^{(k)}$,
where their variables are present.

As in \eqref{eq:mrg-history-insertion}, the bare-field coupling gives
$V=\mathcal M^kU$ only away from processed components. The two-point
calculation uses the sourced exponents in \eqref{eq:witness-exact-identity}
at $k_0$ and $k_1$, with their respective laws $\mathsf P^{(k)}$.

At $k_0$ define the frozen-tent interpolation by
\[
 \zeta'_{k_0}\tilde f_i(p)
   =f_i(x(p))+\sum_{j\le k_0}b_j'(0)
                           \sum_yh_y^{(j)}(p)f_i(x(p_y)).
\]
It is defined globally, independently of the histories. Where the action
profiles all equal one, this is exactly the derivative of the running
source. Since every tent samples points at physical distance at most
$CML^{-s}$ from $p$, subtraction of
$\zeta'_{k_0}f_i(x(p))$ gives
\[
 |\tilde f_i(p)-f_i(x(p))|
     \le\frac{CM}{\vartheta\hat d}
                   \sum_{j\le k_0}|b_j'(0)|\|f_i\|_\infty.
\]
The supports enlarge by at most $CML^{-s}$.

Localize the background around the support and put a regularity window
$\chi_{k_0}$ on each localization cube. The resulting main insertion is
\[
 X_i(V)=\sum_p\tilde f_i(p)\chi_{k_0}(\square(p))
                         a_p(U_{k_0,\square(p)}(V)).
\]
This is one function of $V$, the same in every history. Differentiating
the effective exponent decomposes $\Phi_i^{(k_0)}$ into a constant,
$\zeta'_{k_0}X_i$, irrelevant activities, contributions of live
components, and the error in localization. The window gives
\[
 |X_i|\le Cx^{-1}p_0(g(d))^2\hat d^4\|f_i\|_\infty.
\]
For the full irrelevant part, the two-point comparison requires the
differentiated regular-family bound specified in
\hyperref[hyp:sourced-estimates]{the local source induction}:
\begin{equation}\label{eq:witness-source-slope}
 Cx^{-1}(\log x)^{3\mathfrak p_0-q}
                           \hat d^4\|f_i\|_\infty.
\end{equation}
The factor $x^{-1}$ is essential here. Its source is the dependence of
the regular family on the coupling array, weighted by
$CM^4\alpha_{0,j}^2$, and the source-slope induction. A Cauchy
estimate on a fixed disc of radius $\rho$ would give no such factor.
The slope induction in \hyperref[hyp:sourced-estimates]{the local source induction} has size
$x_j^{-1}(\log x_j)^{\mathfrak p_0+q}$; the spatially varying part
has an additional $L^{-(k_0-j)}$. Its short-scale recurrence closes
because $\mathsf n^2(\log x)^{\mathfrak p_0-q}=o(1)$.

The covariance of an irrelevant term with a main term has, relative to
$x^{-2}\prod_i\|f_i\|_\infty$, a coefficient bounded by
\[
 Cp_0(g(d))^2\hat d^8(\log x)^{3\mathfrak p_0-q}
        =O((\log x)^{21\mathfrak p_0-q})=o(1).
\]
This explains the choice $q>21\mathfrak p_0$. The covariance of two
irrelevant terms is smaller by the same estimates. The constants in
the first derivatives do not contribute to a covariance.

The cubic part of the required estimate has a simple power count.
Let $\alpha$ be a field-analyticity radius comparable to
$\alpha_{0,k_0}$ and $\alpha_{1,k_0}$, with comparison constants
fixed by the analytic class. The source perturbation of the Wilson
action has size $O(\alpha^2)$ per localization cube on this domain.
For a differentiated activity that retains this weight, its cubic
Taylor remainder on the window has size
\[
 \alpha^2(\varepsilon_{k_0}/\alpha)^3
       =O\bigl(x^{-1}(\log x)^{3\mathfrak p_0-q}\bigr).
\]
This calculation does not cover the whole insertion. After its
constant and its multiple of the plaquette action have been
subtracted, a quadratic component remains. The expansion
\eqref{eq:deep-window-taylor} and the following pointed sum retain
both the quadratic H\"older error and the cubic remainder, as well
as the tail of the omitted domains.
Applying that cancellation to the differentiated families requires
their source-response bound on the same complex domains.
Until this estimate and its applicability are established,
\eqref{eq:witness-source-slope} is an additional input to the
two-point comparison. A Cauchy estimate on a fixed bare-source
disc does not supply its factor $x^{-1}$.

\subsection{Live components and mixed derivatives}

For disjoint supports the action terms which depend on the source at
one sampled plaquette have zero mixed derivative: their derivative
contains $f_1(x(p_y))f_2(x(p_y))=0$. A regular local activity has a
mixed derivative only if its enlarged domain meets both supports.
It then has tree length at least $c\hat d/M$. Domain summation and
Cauchy's estimate bound these terms by
\[
 C\hat d^4e^{-c\kappa\hat d/M}
                         \prod_i\|f_i\|_\infty.
\]
Since $\hat d\asymp p_0(g(d))^2$, this is smaller than every fixed
power of $g(d)$.

For live components only their means are small. A component of lineage
age $\alpha$ may contain fields from $\alpha+1$ finer levels, with
$O(L^{4(\alpha+1)})$ plaquettes per level-$k$ cell. One first
derivative can therefore grow like $L^{4\alpha}$; two derivatives
or a product of two first derivatives can grow like $L^{8\alpha}$.
Use \hyperref[hyp:history-bounds]{the weighted history and live-age estimates} with a rate at least eight, increasing it
to absorb any polynomial in $\alpha$:
\[
 \sum_{\alpha\ge0}L^{8\alpha}
    \widetilde\nu(\mathrm{Live}_{k,\ge\alpha}(\mathfrak c))
                      \le Ce^{-cp_0(g_k)}.
\]
The amplitude and probability use the same unhealed lineage age. A
region healed and covered again has not survived through the
intervening levels.

At $k_0$ at most $C\hat d^4$ cells are in a fixed multiple of the
support neighbourhood. Let $\alpha_{\mathrm c}$ be the largest live
lineage age touching this neighbourhood. The union bound, followed by
the nonnegative age sum, gives
\[
 \mathbb E_{\mathsf P^{(k_0)}}
       [\mathbf1_{\mathrm{contact}}L^{8\alpha_{\mathrm c}}]
       \le C\hat d^4e^{-cp_0(g(d))}.
\]
It pays the crude amplitude factors
$(\log x)^CL^{C\check R_{k_0}}$ and any fixed power of $x$,
because the exponent defining $p_0$ exceeds that defining $\check R$.
This controls both the live part of $\mathbb E\Phi_{12}^{(k_0)}$
and the covariances containing live first derivatives. At $k_1$ the
same reasoning is applied with its own law $\mathsf P^{(k_1)}$.
There are only $O(1)$ touching cells there. The extra factor
$L^{\mathsf n}\le L\bar x_s^2$ per source derivative is a power of
$x$ and is again paid by the age estimate. Each argument concerns one
selected observation level, so there is no union over the full run.

We give the separate bounds entering this conclusion. Let
$\mathcal L_i$ be the live contribution to $\Phi_i^{(k_0)}$ and
$\mathcal L_{12}$ its mixed derivative. On a contact class of
largest lineage age $a$, the local history and source estimates to
establish are
\begin{align*}
 |\mathcal L_{12}|
 &\le(\log x)^C
       [e^{-c\hat d}+L^{C\check R_{k_0}}L^{8a}]
                                      \|f_1\|_\infty\|f_2\|_\infty,\\
 |\mathcal L_i|
 &\le(\log x)^C
       [e^{-c\hat d}+L^{C\check R_{k_0}}(a+2)^CL^{4a}]
                                                   \|f_i\|_\infty.
\end{align*}
The exponent $4a$ counts the stored plaquette sum for one derivative;
$8a$ counts two such sums. If a component holds an island, the stored
core begins at the first level of the unhealed lineage. Its label age
can restart, but that only decreases the exponent relative to $a$.
The local terms outside contact have the displayed exponentially small
localization contribution.

Choose the age rate in \eqref{eq:lineage-age} strictly larger than
eight. Then $(a+2)^CL^{8a}\le C_bL^{ba}$, and the same contact
count gives the weighted estimates with that polynomial inserted.
Equivalently the extra geometric factor in the pointwise age tail
absorbs the polynomial. Thus
\begin{align*}
 |\mathbb E\mathcal L_{12}|
 &\le(\log x)^C
       [e^{-c\hat d}+C\hat d^4L^{C\check R_{k_0}}
                                       e^{-cp_0(g(d))}]
                                      \prod_i\|f_i\|_\infty,\\
 \mathbb E|\mathcal L_i|
 &\le(\log x)^C
       [e^{-c\hat d}+C\hat d^4L^{C\check R_{k_0}}
                                       e^{-cp_0(g(d))}]
                                                   \|f_i\|_\infty,\\
 \mathbb E|\mathcal L_i|^2
 &\le(\log x)^C
       [e^{-c\hat d}+C\hat d^4L^{C\check R_{k_0}}
                                       e^{-cp_0(g(d))}]
                                                 \|f_i\|_\infty^2.
\end{align*}
The generic constants and the coefficient of $\check R_{k_0}$ have
been enlarged in the last line. Multiplication by any fixed power of
$x$ still leaves a quantity tending to zero. Indeed
$\hat d$ is a fixed power of $\log x$ of exponent
$2\mathfrak p_0$, while $p_0(g(d))$ has exponent
$\mathfrak p_0$ and $\check R_{k_0}$ has the smaller exponent $\mathfrak r$.

There are two different covariance uses. For the two live first
derivatives, Cauchy--Schwarz bounds the product moment by the square
root of their second moments, and bounds their means by their first
moments. Hence
\[
 |\operatorname{Cov}(\mathcal L_1,\mathcal L_2)|
                         =o(x^{-2})\prod_i\|f_i\|_\infty.
\]
For a live first derivative and a bounded regular factor $B_i(V)$,
use instead
\[
 |\operatorname{Cov}(\mathcal L_1,B_2)|
                          \le2\|B_2\|_\infty\mathbb E|\mathcal L_1|.
\]
The main and irrelevant regular factors have supremum bounded by
$Cx^{-1}(\log x)^C\hat d^4\|f_2\|_\infty$.
The first-moment bound therefore gives the same $o(x^{-2})$ result.
Using a pointwise supremum for the live factor here would lose the
age gain. Together with the mixed-derivative estimate, these are the
two uses of the age theorem at $k_0$: one for the mean mixed
derivative, one for the first-derivative covariances. The third use is
at $k_1$ for its rough first-derivative covariance, under its own
standing-history law and its own contact age. No comparison of
unhealed ages between the two observation levels is required.

At the finer level \eqref{eq:witness-exact-identity} thus locates the
two-point function in the covariance of the main insertions. At the
coarser level the sourced slope bounds and the preceding live estimates
give instead
\[
 |\operatorname{Cov}_{\mathsf P^{(k_1)}}
             (\Phi_1^{(k_1)},\Phi_2^{(k_1)})|
             \le o(1)x^{-2}\prod_i\|f_i\|_\infty.
\]
Its mixed derivative is a scalar source-energy derivative plus a
field-dependent term. The latter has an initial bound
$CL^{2\mathsf n}$ per unit volume at level $k_1$.

Here the common translated-array device gains a volume factor. Translate
the pair over a grid of spacing comparable to $d$ and give the two
members of each copy the same independent sign. Averaging in the signs
removes all mixed-copy derivatives. The array has the same supremum
and Lipschitz bounds as one copy. There are
$\asymp L^{4\mathsf n}\hat d^{-4}$ copies per unit volume at
$k_1$. Division by this number leaves the field-dependent error
\[
 CL^{-2\mathsf n}\hat d^4
       \le C\bar x_s^{-4}\hat d^4=o(x^{-2}).
\]
The surviving scalar number is the translated average of the mixed
derivatives of the local source energies at the trivial block field.
This reduces the nonlinear calculation to a scalar recursion; it does
not by itself evaluate that recursion.

More explicitly, for a translated-pair grid $\mathfrak P$ at level
$k_0$, the scalar is
\[
 c_{12}=\frac1{|\mathfrak P|}\sum_{v\in\mathfrak P}
   \left.\partial_1\partial_2
     \sum_{j\le k_1}\sum_X
       e_j(X;z_1\tau_vf_1+z_2\tau_vf_2)\right|_{z=0}.
\]
To verify the division used above, let
$F_i^\epsilon=\sum_{v\in\mathfrak P}\epsilon_v\tau_vf_i$.
The mixed derivative of any holomorphic sourced expression is bilinear
in the two source directions at zero. Its average at
$(F_1^\epsilon,F_2^\epsilon)$ is therefore the sum of the mixed
derivatives at the coincident translated pairs, since
$\mathbb E\epsilon_v\epsilon_{v'}=\mathbf1_{v=v'}$.
Common translation invariance makes the bare two-point function of
each pair equal to the original value. The field-dependent part of
the effective expression has a bound per coarse unit volume, so
division by the number of pairs per such volume gives
$L^{-2\mathsf n}\hat d^4$. This averaging is done in the original
source variables. It neither replaces the reference law nor conditions
the bare observable on the block field.

\subsection{Free curvature and the reference-torus comparison}

The free comparison has three parts. First identify the curvature
covariance independently of the gauge used for the potential. Then
prove decay for one constrained Gaussian step, with constants
independent of the number of fine sites in a block. Finally sum the
steps below the observation scale. The second part is the regularity
argument: an operator bound in a Hilbert norm alone does not control
the curvature at one bare plaquette.

\emph{Averages and conditional Gaussian projections.}
Work first with one colour and inverse coupling one. Let
$\mathcal E^2_{\rm ex}$ be the exact two-cochains of the bare torus,
with the counting inner product, and let $P_{\rm ex}$ be its
orthogonal projection. A standard Gaussian $F$ in this space has
covariance $P_{\rm ex}$. Indeed, in an orthonormal basis of
$\mathcal E^2_{\rm ex}$ its coefficients are independent standard
normal variables. The discrete Hodge identity gives
\[
 P_{\rm ex}=d\mathcal L_1^+d^*,\qquad
 \mathcal L_1=d^*d+dd^*,
\]
where the superscript $+$ denotes the inverse on the orthogonal
complement of the kernel. The harmonic one-cochains on the torus
are constant; they do not contribute to $F$.

Denote the linearized contour average by $Q$, the straight average
by $Q^{\rm st}$, and the block mean of site functions by $Q'$.
The difference of the first two is a coarse gradient. More
explicitly, if $\bar A(y,x)$ is the integral along the contour from
the centre $y$ to $x$, then
\[
 \varphi_A(y)=L^{-4}\sum_{x\in B(y)}\bar A(y,x),\qquad
 QA=Q^{\rm st}A-d\varphi_A.
\]
For $A=d\lambda$ this reads
\[
 Q^{\rm st}d\lambda=dQ'\lambda,
 \qquad Qd\lambda=d(\lambda|_{\rm coarse}).
\]
Stokes' formula on each translated coarse square therefore gives
\[
 dQ=dQ^{\rm st}=Q^{\rm pl}d.
\]
The centred blocks ensure that both potential averages send a
constant one-cochain to $L$ times that cochain. These identities
persist under composition through $j$ levels.

The straight average is onto. Assign $B(c)$ to every fine bond
crossing the face of a coarse bond $c$, and zero to the other
bonds. There are $L^3$ crossing bonds, each appearing in $L$
of the straight paths, so the average is
$L^{-4}L^3L\,B(c)=B(c)$. It follows by the coboundary identity
that the plaquette average maps exact bare curvatures onto the
exact coarse curvatures.

We also need the norm of the plaquette average to be uniform in
$j$. For a coarse plaquette $P$, write
\[
 (Q_j^{\rm pl}F)(P)=L^{-4j}\sum_p m_P(p)F(p).
\]
Here $m_P(p)\le L^{2j}$ and $\sum_pm_P(p)=L^{6j}$.
Cauchy--Schwarz bounds this square by the sum of $|F(p)|^2$ over
its carrier. Each bare plaquette belongs to at most four such
carriers, whence $\|Q_j^{\rm pl}\|\le2$.

Put
\[
 \mathcal Y_j=\mathcal E^2_{\rm ex}
             \ominus\ker(Q_j^{\rm pl}|_{\mathcal E^2_{\rm ex}}),
 \qquad
 \mathsf S_j=Q_j^{\rm pl}P_{\rm ex}(Q_j^{\rm pl})^*,
 \qquad
 \mathfrak E_j=P_{\rm ex}(Q_j^{\rm pl})^*\mathsf S_j^+.
\]
Thus $\mathfrak E_jh$ is the exact bare curvature of least norm
whose plaquette average is $h$, and
\[
 \mathcal H_j=\mathfrak E_jd_{(j)},\qquad
 \mathsf M_j=\mathcal H_j^*\mathcal H_j
            =d_{(j)}^*\mathsf S_j^+d_{(j)}
            \ge\tfrac14d_{(j)}^*d_{(j)}.
\]
The lower bound follows from $\|\mathsf S_j\|\le4$ on its range.
In particular $\mathcal H_j$ vanishes on closed one-cochains.

These definitions agree with minimizing the potential action
under either potential average. To see this, suppose
$Q_j^{\rm pl}dA=d_{(j)}B$. The difference
$Q_j^{\rm st}A-B$ is closed. Write it as $d\lambda+\kappa$,
where $\kappa$ is constant. Subtract from $A$ a gradient of a
blockwise constant extension of $\lambda$, and subtract
$L^{-j}\kappa$. This leaves $dA$ unchanged and makes
$Q_j^{\rm st}A=B$. For the contour average, use an extension
with prescribed values at the coarse sites instead. Consequently
\[
 \{dA:Q_jA=B\}=\{dA:Q_j^{\rm st}A=B\}
 =\{F'\in\mathcal E^2_{\rm ex}:
                              Q_j^{\rm pl}F'=d_{(j)}B\}.
\]
This identifies $\mathcal H_jB$ with the curvature of the linear
minimizer in either convention.

Conditioning the potential modulo closed forms is therefore
equivalent, for curvature, to conditioning $Q_j^{\rm pl}F$.
The latter determines $P_{\mathcal Y_j}F$. In an orthonormal
basis adapted to $\mathcal Y_j\oplus\mathcal Y_j^\perp$ the two
Gaussian components are independent, so the conditional mean is
$P_{\mathcal Y_j}F$. The spaces decrease with $j$. At depth $s$,
with $k=K-s$, this proves
\[
 C_{K,s}^{\rm bg}=P_{\mathcal Y_k},\qquad
 C_K^{\rm free}-C_{K,s}^{\rm bg}
   =\sum_{j=1}^{k}\mathcal C_{(j)},\qquad
 \mathcal C_{(j)}=P_{\mathcal Y_{j-1}}-P_{\mathcal Y_j}.
\]
The slices are independent Gaussian components. Coarse translations
permute the averaging blocks and commute with these projections;
general bare translations need not do so.

\emph{A uniform bound for the linear minimizer.}
We give the regularity argument for the centred averages. It is
the flat-field minimizer method of Ba\l aban\cite{Balaban1984},
with the averaging operators retained explicitly. Work in level-$j$
units. The fine mesh is $\eta=\eta_j$, coarse blocks have side one,
and fine-lattice norms use quadrature measure $\eta^4$. Differences
contain $\eta^{-1}$, and
$\mathcal L=\nabla^*\nabla=d^*d+dd^*$ on one-forms. Potential
angles are rescaled by $\eta^{-1}$. If $\mathsf H^{\rm st}$ is
the resulting minimum-action response, then
\[
 \mathcal H_j(p,b)=\eta_j^2
                    (d\mathsf H^{\rm st}\delta_b)(p).
\]
The factor $\eta_j^2$ is an area factor, rather than a decay
estimate. The remaining bound must be uniform in $\eta$.

Start with the scalar inverse
\[
 G'=(\mathcal L+Q'^*Q')^{-1},\qquad
 \Sigma=Q'G'^2Q'^*,\qquad
 R=I-G'Q'^*\Sigma^{-1}Q'G'.
\]
The block Poincar\'e inequality is uniform in the mesh:
\[
 \|\lambda-Q'^*Q'\lambda\|^2
       \le\tfrac14\|d\lambda\|^2.
\]
Indeed, the first nonzero Neumann eigenvalue of a unit block is
$4\eta^{-2}\sin^2(\pi\eta/2)\ge4$. It follows that
$\mathcal L+Q'^*Q'\ge c(\mathcal L+1)$.

The coarse matrix $\Sigma$ has a uniform lower bound as well.
Choose on each block a smooth bump $e_y$ of block mean one, zero
outside that block, with $\|\mathcal Le_y\|\le C$. Products of
the centred functions $\cos^2\pi(x_\mu-y_\mu)$, normalized to
mean one, have these properties. Their vanishing to second order
at block faces controls the differences across the faces too.
For $\varphi_\omega=\sum_y\omega(y)e_y$,
\[
 \|\omega\|^2
 =\langle Q'\varphi_\omega,\omega\rangle
 =\langle(\mathcal L+Q'^*Q')\varphi_\omega,
                                      G'Q'^*\omega\rangle
 \le C\|\omega\|\,\|G'Q'^*\omega\|.
\]
Thus $\Sigma\ge cI$. The case $j=0$ needs no bump construction,
since $\mathcal H_0=d$.

The operator $I-R$ is the orthogonal projection onto
$\operatorname{Ran}(G'Q'^*)$, and $R$ projects onto
$\mathcal L\ker Q'$. For $\lambda_0\in\ker Q'$ the two
spaces are orthogonal, since
\[
 \langle G'Q'^*\omega,\mathcal L\lambda_0\rangle
 =\langle\omega,Q'G'(\mathcal L+Q'^*Q')\lambda_0\rangle=0;
\]
their dimensions are complementary. This also identifies the
condition $Rd^*a=0$ as fixing the gauge modes of zero block mean.

Define
\[
 \mathsf K_v=d^*d+dRd^*+Q^{{\rm st}*}Q^{\rm st},\qquad
 \mathsf G^{\rm st}=\mathsf K_v^{-1}.
\]
We prove
\[
 \langle a,\mathsf K_va\rangle
 =\|da\|^2+\|Rd^*a\|^2+\|Q^{\rm st}a\|^2
 \ge c(\|\nabla a\|^2+\|a\|^2).
\]
The gauge part not seen by the curl needs particular care. Let
$\Upsilon a$ be the average of $a$ over the fine bonds crossing
the face between two neighbouring blocks. The divergence theorem,
integration along lines, and the block Poincar\'e inequality give
respectively
\[
 Q'd^*=d_c^*\Upsilon,\qquad
 \|(Q^{\rm st}-\Upsilon)a\|\le C\|\nabla a\|,\qquad
 \|a\|^2\le C(\|\nabla a\|^2+\|Q^{\rm st}a\|^2).
\]
For the middle estimate a value in the straight average is compared
with the crossing bond on the same line. Their distance is at most
one; Cauchy--Schwarz bounds the difference by the differences along
that line. The last estimate first subtracts the mean on each block
and then recovers that mean from its straight average.

Choose $\lambda_0\in\ker Q'$ with
$\mathcal L\lambda_0=Rd^*a$, and put $a'=a-d\lambda_0$.
The block inequality gives
\[
 \|d\lambda_0\|+\|\nabla d\lambda_0\|
                                      \le C\|Rd^*a\|.
\]
Now $u=d^*a'$ belongs to $\operatorname{Ran}(G'Q'^*)$ and
has mean zero. Set $\psi=\mathcal L^+u$ and
$a''=a'-d\psi$. Then $d^*a''=0$ and
$\|\nabla a''\|=\|da\|$. Write $\psi_c=Q'\psi$. Since
$Q^{\rm st}d\lambda_0=0$,
\[
 Q^{\rm st}a=Q^{\rm st}a''+d_c\psi_c.
\]
The flux of $a''$ is divergence free. Hence
\[
 \langle Q^{\rm st}a'',d_c\psi_c\rangle
 =\langle(Q^{\rm st}-\Upsilon)a'',d_c\psi_c\rangle,
 \qquad
 \|d_c\psi_c\|\le C(\|da\|+\|Q^{\rm st}a\|).
\]

It remains to recover $d\psi$ from these coarse differences.
Write $u=G'Q'^*\omega$. The lower bound for $\Sigma$ implies
$\|\omega\|\le C\|u\|$, and
\[
 \mathcal Lu=Q'^*\omega_0,\qquad
 \omega_0=\omega-Q'u=Q'\mathcal Lu=d_c^*\Upsilon du.
\]
Consequently
\[
 \|u\|^2=\langle\mathcal Lu,\psi\rangle
        =\langle\Upsilon du,d_c\psi_c\rangle.
\]
The scalar equation bounds $\|\mathcal Lu\|$ and $\|du\|$
by $C\|u\|$. The flux estimate then bounds
$\|\Upsilon du\|$ by the same quantity, giving
$\|u\|\le C\|d_c\psi_c\|$. Also
\[
 \begin{split}
 \|d\psi\|^2
 &=\langle\psi-Q'^*Q'\psi,u\rangle
                    +\langle d_c\psi_c,\Upsilon d\psi\rangle\\
 &\le\tfrac12\|d\psi\|\,\|u\|
   +\|d_c\psi_c\|(\|d\psi\|+C\|u\|).
 \end{split}
\]
Since $\|\nabla d\psi\|=\|u\|$, this proves
$\|d\psi\|+\|\nabla d\psi\|\le C\|d_c\psi_c\|$.
The vector block inequality now bounds $a''$; adding the bounds
for all three pieces proves the asserted gap for $\mathsf K_v$.

The constrained minimizer has the formula
\[
 \mathsf H^{\rm st}
 =\mathsf G^{\rm st}Q^{{\rm st}*}
       (Q^{\rm st}\mathsf G^{\rm st}Q^{{\rm st}*})^{-1}.
\]
To check the gauge condition, test
$\mathsf K_vf=Q^{{\rm st}*}v$ against $d\lambda_0$ with
$Q'\lambda_0=0$. The curl and average terms vanish and the
remaining identity gives $Rd^*f=0$. On the fibre
$Q^{\rm st}a=B$, every potential can be moved to this gauge
without changing its curl. Minimizing the form of $\mathsf K_v$
therefore minimizes $\|da\|^2+\|B\|^2$ on that fibre. In
particular
\[
 (Q^{\rm st}\mathsf G^{\rm st}Q^{{\rm st}*})^{-1}
                                      =I+\mathsf M_j.
\]
For a uniform upper bound on this inverse, use smooth vector bumps
concentrated near the common face in each coarse-bond tube. Their
matrix of straight averages has three nonnegative entries in the
bond direction, of sum one, with its central entry at least
$3/4$. It is thus a uniformly invertible perturbation of its
diagonal. Combining the bumps gives $a_B$ with
$Q^{\rm st}a_B=B$ and $\|\nabla a_B\|\le C\|B\|$.
The variational formula yields
$Q^{\rm st}\mathsf G^{\rm st}Q^{{\rm st}*}\ge cI$.

We next retain exponential localization through these inverses.
Conjugate by $e^\varphi$, where $\varphi$ has a small fixed
Lipschitz constant. For the local gradients and averaging maps,
the conjugation error is $O(\|\nabla\varphi\|_\infty)$ between
the scales defined by $(\mathcal L+1)^{\pm1/2}$. The lower
bounds just proved permit a Neumann series successively for
\[
 \mathcal L+1,\quad \mathcal L+Q'^*Q',\quad
 \Sigma,\quad \mathsf K_v,\quad
 Q^{\rm st}\mathsf G^{\rm st}Q^{{\rm st}*}.
\]
Their constants and an allowed decay rate $\theta_\sharp>0$
are independent of $j$, the fine mesh and the torus. For
$\varphi(z)=\delta|z-x_b|$ at a smaller rate this gives
\[
 \|e^\varphi v\|_{\ell^2({\rm coarse})}
       +\|e^\varphi f\|_{L^2(\eta^4)}\le C,
 \quad
 v=(Q^{\rm st}\mathsf G^{\rm st}Q^{{\rm st}*})^{-1}\delta_b,
 \quad f=\mathsf H^{\rm st}\delta_b.
\]

The improvement to a pointwise bound uses the equation, including
the gauge projection. Set $\Xi=Q'G'd^*$. Then
\[
 d^*d+dRd^*=\mathcal L-\Xi^*\Sigma^{-1}\Xi,
\]
so that
\[
 (\mathcal L+1)f=h_1+f,\qquad
 h_1=Q^{{\rm st}*}v-Q^{{\rm st}*}Q^{\rm st}f
                                  +\Xi^*\Sigma^{-1}\Xi f.
\]
Every term in $h_1$ is the image of coarse-lattice data. For the
last term the rows of $\Xi$ are $du_y$, with
$u_y=G'Q'^*\delta_y$. These rows are controlled first. The
scalar equation is
\[
 (\mathcal L+1)u_y
        =\mathbf1_{B(y)}+u_y-Q'^*Q'u_y.
\]
The non-$u_y$ terms on the right are obtained from uniformly
bounded, exponentially localized coarse data. The scalar
resolvent and its first differences then give
\[
 |u_y(z)|+|\nabla u_y(z)|\le Ce^{-c|z-y|}.
\]
Here are the two norm improvements used for this assertion and
then for $f$. If $\mathsf g$ is the kernel of
$(\mathcal L+1)^{-1}$, its weighted kernel norms satisfy
\[
 \|e^{\theta|\cdot|}\mathsf g\|_p\le C_p\quad(p<2),
 \qquad
 \|e^{\theta|\cdot|}\nabla\mathsf g\|_p\le C_p\quad(p<4/3).
\]
All norms include the quadrature factor $\eta^4$. The heat-kernel
representation proves their uniformity. The one-dimensional heat
kernel is a limit of convolutions of three-point log-concave
kernels; Cauchy--Binet preserves log-concavity under convolution.
An exponential tilt preserves it too. A summable log-concave
sequence is unimodal, so the sum of its absolute first differences
is twice its maximum. Together with the heat-kernel supremum
bound, this controls the first-difference kernel. Shifting the
Fourier contour introduces only $e^{Ct\theta^2}$, absorbed by
the resolvent factor $e^{-t}$ for small $\theta$. Interpolation
and integration in $t$ give exactly $p<2$ for the resolvent and
$p<4/3$ for its gradient. Periodization retains the bounds after
part of the exponential rate is spent.

Weighted Young inequalities now use
\[
 (p,q,r)=(8/5,2,8),\qquad (p,q,r)=(8/7,8,\infty).
\]
The first application puts the solution in weighted $L^8$.
In the scalar equation the remaining forcing is the sum of this
solution and bounded coarse images, so it too lies in weighted
$L^8$. The second application bounds the solution and its first
differences. Thus it proves the estimate for $u_y$. Its gradient
bound controls the rows of $\Xi$; their exponential decay and
the coarse inverse bound for $\Sigma^{-1}$ give the same
bounded coarse-image forcing in the equation for $f$. Applying
the two improvements there yields
$|f|+|\nabla f|\le Ce^{-\delta|z-x_b|}$.
This verifies the forcing estimate needed by the second Young
inequality, rather than applying it to an arbitrary $L^2$ source.
Returning to bare plaquette units proves
\begin{equation}\label{eq:free-minimizer-decay}
 |\mathcal H_j(p,b)|\le C\eta_j^2
               e^{-\delta_0|x(p)-x_b|/a_{K-j}}.
\end{equation}
The rate $\delta_0$ can be fixed independently of $j$ and the
torus. We will use a smaller rate depending on the fixed block
factor $L$ for the constrained slices.

\emph{The constrained covariance and the choice of gauge.}
At step $j$, write $d'=d_{(j-1)}$. In each block fix the
lexicographic tree from its centre, and let
\[
 \Omega^1_{\rm tr}
 =\{B:B=0\text{ on tree bonds},\ Q^{\rm st}B=0\}.
\]
On this space
\[
 \|d'B\|^2\ge c_L\|B\|^2,
\]
where $c_L>0$ is a negative power of $L$. Inside a block a bond
value is recovered by summing curvature along the strips between
the bond and the tree. Starting with the fourth direction, which
vanishes on the tree, one recovers the third component along
$e_4$, the second along $e_4,e_3$, and the first along
$e_4,e_3,e_2$. Each formula uses $O(L)$ plaquettes, and a
plaquette is used in at most $CL^3$ such formulas. On a common
face, differences between parallel crossing bonds are given by
curvature and the recovered in-block bonds. The equation
$Q^{\rm st}B=0$ determines their mean. Cauchy--Schwarz in these
formulas gives the stated inequality.

The map $d'$ is a bijection from $\Omega^1_{\rm tr}$ onto
the exact curvatures with zero plaquette average. For surjectivity,
remove the closed defect of the straight average as above and then
set the tree values to zero by a gradient whose block mean is zero.
This preserves the average constraint. The Poincar\'e inequality
gives injectivity. For each coarse bond the average constraint
can be solved for its single pivot bond. Thus there is a local
coordinate map $\mathsf C_{\rm tr}$ from the remaining bonds
onto $\Omega^1_{\rm tr}$, with range $O(L)$ and
$\|\mathsf C_{\rm tr}B'\|\ge\|B'\|$.

Let $\iota$ be the inclusion of this tree-gauge space. Conditioning
a Gaussian on a linear constraint restricts its quadratic form to
the kernel of that constraint. The Schur complement therefore gives
\[
 C_j^{\rm tr}
 =\iota(\iota^*\mathsf M_{j-1}\iota)^{-1}\iota^*
 =\mathsf C_{\rm tr}
  (\mathsf C_{\rm tr}^*\mathsf M_{j-1}\mathsf C_{\rm tr})^{-1}
                                      \mathsf C_{\rm tr}^*,
 \qquad
 \mathcal C_{(j)}=\mathcal H_{j-1}C_j^{\rm tr}\mathcal H_{j-1}^*.
\]
The inverse is uniformly positive at fixed $L$, since its form is
at least $c_L/4$.

The soft-gauge potential covariance need not equal $C_j^{\rm tr}$.
The identity needed here is
\[
 d'C_j^{\rm soft}d'^*=d'C_j^{\rm tr}d'^*.
\]
Indeed, both constraints give the same set of curvatures, and on
this set the quadratic action is
$\langle d'B,\mathsf S_{j-1}^+d'B\rangle$. The fibre of the
contour-average constraint over one curvature is an orbit of the
block gauge group. The soft gauge-fixing factor has the same
normalized integral on every such orbit: its quadratic part is
the square of the contour potential, which changes by translation
along the orbit. Integrating the gauge variables leaves the
curvature Gaussian unchanged. The hard tree gauge gives that same
Gaussian through the preceding bijection. Since $\mathcal H_{j-1}$
vanishes on closed forms, this is precisely the equality required
in the slice formula.

\emph{Decay of a slice and the sum over fine scales.}
Put $a'=a_{K-j+1}$. From
\eqref{eq:free-minimizer-decay},
\[
 |\mathsf M_{j-1}(b,b')|
 \le\sum_p|\mathcal H_{j-1}(p,b)\mathcal H_{j-1}(p,b')|
 \le Ce^{-c|x_b-x_{b'}|/a'}.
\]
The factors $\eta_{j-1}^2$ in the two kernels give
$\eta_{j-1}^4$, exactly canceled by the number
$O(\eta_{j-1}^{-4})$ of bare plaquettes in one unit volume at
this level. The finite range of $\mathsf C_{\rm tr}$ preserves
exponential decay of its compressed form.

For completeness, a positive matrix $A\ge\gamma I$ with
$|A(b,b')|\le C_Ae^{-2\alpha d(b,b')}$ has an exponentially
decaying inverse. Conjugating by
$e^{\tau d(\cdot,b_0)}$ changes it by a matrix whose row and
column sums are at most
\[
 C_A\tau\sup_b\sum_{b'}d(b,b')e^{-\alpha d(b,b')}.
\]
Choose $\tau$ so this is at most $\gamma/2$. A Neumann series
then bounds the conjugated inverse by $2/\gamma$, and its
$b_0$-column decays as $e^{-\tau d(b,b_0)}$.
Apply this to the compressed tree-gauge form. Composing its
inverse with the two minimizer kernels, and spending half the
rate on the triangle inequality and half on the two bond sums,
gives
\[
 |\mathcal C_{(j)}(p,q)|
 \le C\eta_{j-1}^4
          e^{-\delta|x(p)-x(q)|/a_{K-j+1}}.
\]
Here $C,\delta>0$ may depend on $L$, but not on $j,K$ or the
torus. The constrained inverse can have a slower rate than the
unconstrained minimizer.

Let $u=x(p)-x(q)$ be the shortest torus representative and
$t=|u|/a_s\ge1$. With $k=K-s$ and $i=k-j+1$,
\[
 a'=a_sL^{-i},\qquad
 \eta_{j-1}^4=L^{-4k}L^{4i}.
\]
Consequently
\[
 |C_K^{\rm free}-C_{K,s}^{\rm bg}|(p,q)
 \le CL^{-4k}\sum_{i\ge1}L^{4i}e^{-\delta tL^i}
 \le C(a_K/a_s)^4e^{-ct}.
\]
The sum converges because the exponential in $L^i$ dominates
$L^{4i}$. Beyond a fixed lower bound for $t$, its consecutive
terms decrease geometrically. Factoring out
$\sigma_{pq}=a_K^4\pi^{-2}|u|^{-4}$ gives
\[
 |C_K^{\rm free}-C_{K,s}^{\rm bg}|(p,q)
       \le Ct^4e^{-ct}\sigma_{pq},
       \qquad t\ge cC_W(s).
\]
The reference room is chosen beyond the point at which this last
factor decreases.

\emph{The continuum tensor and the lattice error.}
The continuum Green function is $\Phi(u)=(4\pi^2|u|^2)^{-1}$.
Applying the two curls to $\delta_{\mu\nu}\Phi$ gives
\[
 \mathbb K(u)=\pi^{-2}|u|^{-4}
                    \Lambda^2(I-2\hat u\hat u^{\mathsf T}),
 \qquad \hat u=u/|u|.
\]
The reflection $I-2\hat u\hat u^{\mathsf T}$ is an orthogonal
involution. Its exterior square is therefore an orthogonal
involution on the six-dimensional two-form space, and
\[
 \tr\mathbb K(u)^2=6\pi^{-4}|u|^{-8}.
\]

Here is the lattice comparison in the form required below. On the
unit lattice let $p_t$ be the product heat kernel and let
$p_t^{(n)}$ be its periodization on a torus of side $n$. The
mean-zero Green function is
\[
 \Phi_n=\int_0^\infty(p_t^{(n)}-n^{-4})\,dt.
\]
The curvature covariance consists of second differences of
$\Phi_n$. In one dimension the Fourier symbol is
$\omega_1(\theta)=2-2\cos\theta$, and
$\omega_1(\theta)-\theta^2=O(\theta^4)$. Fourier integration
by parts bounds the second differences of the four-dimensional
heat kernel, for $t\ge1$, by
\[
 Ct^{-3}\langle |z|/\sqrt t\rangle^{-N},
\]
whereas replacing them by the corresponding continuum derivatives
costs at most
$Ct^{-7/2}\langle |z|/\sqrt t\rangle^{-N}$ for $t\ge1$.
The one-sided difference symbol has a first-order relative error;
this is the origin of the half-power loss in $t$. For $t\le1$
the random-walk series gives arbitrary polynomial decay away from
the origin.

Split the Green-function integral at $n^2$. The infinite-lattice
heat-kernel error integrates to $C|z|^{-5}$. For the nonzero
periodic images for $1\le t\le n^2$, the bound
\[
 \sum_{w\in n\mathbb Z^4\setminus\{0\}}
       |\nabla^2p_t(z+w)|\le Ct^{-3}(\sqrt t/n)^6
\]
integrates to $Cn^{-4}$. On $0<t\le1$ the random-walk estimate
gives an image contribution $O(n^{-N})$ for any fixed $N$.
Above time $n^2$, the nonzero torus
modes have a spectral gap of order $n^{-2}$, and their second
differences are bounded by $Cn^{-6}e^{-ct/n^2}$. Their integral,
and the continuum heat tail, are again $O(n^{-4})$. Thus, away
from the diagonal,
\[
 |\nabla^2\Phi_n(z)-\partial^2\Phi(z)|
                                    \le C(|z|^{-5}+n^{-4}).
\]
Moving from the base points to the plaquette barycentres changes
the tensor by another $C|z|^{-5}$. With $n=\ell/a_K$ and
$2a_K\le|u|\le\ell/4$, rescaling therefore yields
\[
 |C_K^{\rm free}(p,q)-a_K^4\mathbb K(u)_{pq}|
 \le C\left[\frac{a_K}{|u|}
                 +\left(\frac{|u|}{\ell}\right)^4\right]\sigma_{pq}.
\]
For its application to the two supports, the actual distances
satisfy $d/2\le|u|\le C_\vartheta d$, with $\vartheta$
the fixed shape parameter. The witness choice gives
$d\ge C_W(s)a_s$. Fixing the floor on $C_W$ so that
$C_W\ge4$, and taking $d$ small at fixed $g,\vartheta$,
gives
\[
 2a_K\le\tfrac12C_W(s)a_s\le|u|
                         \le C_\vartheta d\le\ell_\star/4.
\]
Thus the distances lie in the interval on which the lattice
comparison was proved, and $t=|u|/a_s\ge C_W(s)/2$.
Enlarge the fixed floor once more so that $t^4e^{-ct}$ is
decreasing on this range. Combining the lattice and fine-slice
estimates yields the relative error
\[
 C\left[\frac{a_K}{d}
        +\left(\frac d{\ell_\star}\right)^4
        +C_W(s)^4e^{-cC_W(s)/2}\right].
\]
The constants may depend on the fixed shape. They are
independent of the cutoff and of the positions of the
plaquettes within the two supports.

The scalar source recursion must still be compared with this free
calculation. At one step its two quantities are the covariance of sums
of background plaquette energies and the mean of a bilinear form in
background curvatures. They must be estimated after removal of the
analytic interaction and comparison with the flat background, with
errors summable through the steps. Both quantities are required: they
generate the cross terms when the total covariance is squared. We
isolate their quantitative consequence.

The free Gaussian algebra explains exactly which two terms the
nonlinear one-step estimates must reproduce. Let $M$ be a fixed
background curvature and $\xi$ a centred Gaussian fluctuation of
covariance $A$. For $Q_p=|M_p+\xi_p|^2/2$, Wick's formula gives
\begin{align*}
 \mathbb E_\xi Q_p&=\tfrac12|M_p|^2+\tfrac12\tr A_{pp},\\
 \operatorname{Cov}_\xi(Q_p,Q_q)
      &=M_p^{\mathsf T}A_{pq}M_q
                           +\tfrac12\tr(A_{pq}A_{qp}).
\end{align*}
The first covariance term is the bilinear-curvature quantity; the
second is the fluctuation half-square covariance. If the background
is itself an independent centred Gaussian of covariance $B$, the
total covariance formula adds the covariance of the conditional means.
The result is
\[
 \tfrac12\tr(B_{pq}B_{qp})
       +\tr(A_{pq}B_{qp})
       +\tfrac12\tr(A_{pq}A_{qp})
       =\tfrac12\tr((A+B)_{pq}(A+B)_{qp}).
\]
The middle term cannot be discarded when slices are summed. This
calculation evaluates the free recursion; it supplies no bound for
replacing a nonlinear fluctuation integral by it. Such a replacement
requires both one-step estimates and their accumulated errors. The
resulting quantitative target is the following comparison.

The small parameters available in this last comparison can be made
explicit. In the fluctuation integral at a level $j$, the localized
interaction splits into the new interaction and the variation of the
earlier activities. Their bounds are
\[
 \begin{split}
 |\mathsf W_j^{\mathrm{int}}(Y;B)|
 &\le C\delta_j\Bigl(1+M^{-4}\sum_{b\subset Y}|B(b)|^2\Bigr)
                                      e^{-c\kappa d_j(Y)},\\
 |\mathsf W_j^{\mathrm{old}}(Y;B)|
 &\le C\mathsf s_jN_{\omega,j}e^{-c\kappa d_j(Y)}.
 \end{split}
\]
Here the Gaussian controls the quadratic growth in $B$, while the
ratio $\mathsf s_j=\delta_j/\alpha_{0,j}$ controls the old terms.
The quantity $N_{\omega,j}$ is their weighted regular and remainder
norm from \cref{sec:rg}.
The tests add a tail of size $e^{-cp_0(g_j)^2}$. Since $x_j\asymp x$
on the segment and $N_{\omega,j}$ stays bounded, the total of these
local perturbation amplitudes satisfies
\[
 \sum_{j=k_0}^{k_1-1}
 \bigl(\delta_j+\mathsf s_jN_{\omega,j}+e^{-cp_0(g_j)^2}\bigr)
 \le C(\log x)\Bigl[
 x^{-1/2}(\log x)^{\mathfrak p_0}
 +(\log x)^{\mathfrak p_0-q}
 +e^{-cp_0(x^{-1/2})^2}\Bigr]=o(1).
\]
Thus the short segment leaves room for summable perturbations. The
Cauchy--Taylor calculation must give the corresponding estimates with
the two energy insertions and with the curvature bilinear insertion,
with the coefficients dictated by the source recursion and errors
measured against the required two-point scale $x^{-2}$.
At real data the subsequent normalized integrations preserve a
supremum bound. To use this in the comparison, the inserted recursion
must carry each defect through those integrations with its correct
coefficient. The local amplitude bound alone
does not give this relative inserted estimate; this is the calculation
needed for the comparison below.

\subsubsection*{The relative comparison on the reference torus.}
\phantomsection\label{hyp:relative-two-point}
The comparison can be organized by carrying bounded functions through
the actual unsourced run. This treats the two insertions together and
retains the terms between different Gaussian slices. Let
$\lambda_{k_0}$ be the block-field marginal of
$\mathsf P^{(k_0)}$, tested on gauge-invariant functions, and write
$\phi_1=\tilde f$, $\phi_2=\tilde h$ and $\mathsf O_a=X_a$.
In the rest of this calculation put
\[
 \mathsf g=x^{-1/2},\qquad p_0=p_0(\mathsf g),\qquad
 \mathsf r=\hat d,\qquad \eta_j=L^{-j}.
\]
Thus $\mathsf g$ is the observation-scale coupling; the reference
coupling $g$ remains fixed. In particular
$\mathsf r\asymp p_0^2$, $\mathsf r\ge C\check R_{k_0}$ and
$\sup|\mathsf O_a|\le Cx^{-1}\mathsf r^4p_0^2
\|\phi_a\|_\infty$. The comparison will be uniform in the
first-passage data, the position of the pair and the fixed shape class.

\emph{The positive coupling of the run.}
Freeze every test, profile, reference law and gauge choice at its
unsourced value. The operations in \eqref{eq:history-kernels} then
carry $G\rho_{k_0}$ as well as $\rho_{k_0}$ for every bounded
gauge-invariant Borel function $G$. A test selects one share of a
decomposition of unity; integration, a coordinate change and a
reference draw are positive kernels; merging labels adds their
measures. A gauge-fixing move can be realized by drawing a gauge
transformation from its normalized orbit law. This realization also
preserves positivity and mass and leaves gauge-invariant insertions
unchanged.

Retain the intermediate stored variables. For a term $h$ of level $j$,
let $\hat\tau_h$ be its positive unsourced measure and
$\lambda_h[G]$ the measure carrying $G$. Positivity gives
\[
 |\lambda_h[G]|\le\|G\|_\infty\hat\tau_h,\qquad
 G^{[j]}:=\frac{d\lambda_h[G]}{d\hat\tau_h},\qquad
 |G^{[j]}|\le\|G\|_\infty.
\]
If $\mathfrak Q_j$ is the positive map of a step, define
\[
 \mathbb E_j F=
 \frac{d\mathfrak Q_j(F\hat\tau^{(j-1)})}
      {d\hat\tau^{(j)}}.
\]
The input here may be a disjoint union of terms that are subsequently
merged. In either case $\mathbb E_j$ is positive,
$\mathbb E_j1=1$, and it contracts the supremum norm. The chain rule
for Radon--Nikodym derivatives gives
$G^{[j]}=\mathbb E_jG^{[j-1]}$.

At level $k_1$ condition these measures on the terminal class and its
block field $W$. Classes retain the integral representation and
the exponent used below. If the partition into classes is
refined to record cleanness, that representation must continue
to hold on every resulting class; measurability of cleanness
alone does not imply this. The restriction of the positive
bracket below supplies the required representation.
This gives a probability kernel $\Pi$ and a law
$\mathbb Q$ with
\begin{equation}\label{eq:inserted-run-coupling}
 \int G\,d\lambda_{k_0}
       =\int (\Pi G)(\mathsf h,W)\,d\mathbb Q(\mathsf h,W).
\end{equation}
This is a coupling defined by the run. Where a large-field operation
has drawn new variables, $W$ need not equal
$\mathcal M_{k_0\to k_1}V$. No identification with conditioning on
that deterministic average is used below.

\emph{Clean classes and the local law to verify.}
Around the pair choose, at each level $j-1$, a read set
$\mathcal S_j$ and a larger box $\mathcal B_j$ whose distance from
the read set is at least $\mathsf D$ spacings. Include the localization
cubes of the bare plaquette functions in the read set at the first
step. The straightening of the constraint and the elimination of pivot
bonds have finite range, so a local insertion reads only
$y|_{\mathcal S_j}$ after that change of variables. The number of free
bonds in the box satisfies
\begin{equation}\label{eq:inserted-box-count}
 n_j:=\#\mathcal B_j
       \le C\big(\mathsf rL^{-(j-1-k_0)}+\mathsf D\big)^4.
\end{equation}
Take $\mathsf D$ of order $\mathsf r$. For $d$ sufficiently small at
fixed $g$, all the boxes and their required collars embed in the
reference torus: their largest physical radius is at most
$C(\mathsf r+\mathsf D+\check R_{k_0})a_{s-\mathsf n}$, which tends
to zero with $d$. All subsequent local Gaussian comparisons
use this nonwrapping range. Declaring every history bad when
an exceptional-mass upper bound exceeds one gives only a
vacuous bound; it supplies no Gaussian identity on a wrapping
box. The exact measure decomposition is kept throughout.

A clean history has no live region meeting these boxes, no creation or
large-field operation there along the segment, and no such contact in
the preceding window of $C\check R_{k_0}$ levels. The cell bounds in
\cref{sec:large-fields} give the estimate
\begin{equation}\label{eq:inserted-unclean-weight}
 \epsilon_{\rm uc}\le
 C L^{4\mathsf n}e^{C\check R_{k_0}}(1+\mathsf n)
    (\mathsf r+\check R_{k_0})^4e^{-cp_0}.
\end{equation}
Indeed, count cells at their own levels before applying their
large-field probabilities. The $C\check R_{k_0}$ earlier levels cost
$e^{C\check R_{k_0}}$. Extending this exclusion through all earlier
levels would produce a divergent cell count. Regions healed farther
below the segment must remain summed in the local remainders.

\emph{Restricting the positive bracket.}
The clean restriction can be made before the healed components are
exponentiated. Use a fixed buffered neighbourhood of the pair
containing every read box. At the levels of the clean window, except
its lowest level, exclude each first-class component whose collar
meets that neighbourhood. Also require the retained live regions to
avoid it. This is a slightly smaller clean set than the one obtained
by inspecting the live regions alone. Its excluded components have a
live predecessor within the buffer enlarged by \(C\check R_t\);
the preceding cell count therefore still gives
\eqref{eq:inserted-unclean-weight}. No condition is imposed on
components healed at the lowest level of the window or earlier.

Fix a level \(t\), an output standing part \(o\), and all its retained
variables. Write \(\mathfrak e\) for the common outer datum. The
standing-part identity before exponentiation has the form
\begin{equation}\label{eq:clean-bracket-retained-identity}
 \sum_{(h')^\flat=o}m_{h'}
 =(\varsigma_o)_*
     \bigl(F_o(\mathfrak e)K_t(\mathfrak e)\mu_o(d\mathfrak e)\bigr),
\end{equation}
where \(F_o>0\), \(\mu_o\) is positive, and
\[
 K_t(\mathfrak e)=
 \sum_{\mathfrak X\in\mathfrak F(o)}
 \int\prod_{X\in\mathfrak X}
       \nu_X(\mathfrak e|_{X^\sim};d\xi_X)\,
        \exp\!\sum_Y V_{\mathfrak X}(Y;\mathfrak e,\xi).
\]
The measures \(\nu_X\), at the fixed retained datum, include the sum
over the histories inside \(X\)
and the normalized frame sum. This is the first, unnormalized
representation of a healed component; there is no division by its
mass. The map \(\varsigma_o\) retains every variable on which a later
move depends.

Two facts about this identity are used. Every string producing \(o\)
occurs exactly once in the displayed family sum, and the formula
holds under the retained-variable integral, before taking its image.
Consequently the restriction excluding the forbidden collars is
obtained by deleting those families from this same bracket. Every
remaining summand has its original measure and its original
\(V_{\mathfrak X}\). Denote the restricted bracket by
\(K_t^{\rm cl}\). Equivalently, give every forbidden component a
factor \(z\) in the finite family sum. The constant coefficient of
this polynomial is \(K_t^{\rm cl}\). Coefficient selection is made
in \eqref{eq:clean-bracket-retained-identity}, so the selected identity
still holds against every bounded Borel function of the retained
variables. No RG step is applied to a newly restricted exponent.

Take out the contribution of the empty family, using the same
regrouping in both brackets:
\[
 K_t=e^{\sum_Y V_\emptyset(Y)}B_t,\qquad
 K_t^{\rm cl}=e^{\sum_Y V_\emptyset(Y)}B_t^{\rm cl}.
\]
Here
\[
 B_t=\sum_{\mathfrak X\in\mathfrak F(o)}
     \int\prod_{X\in\mathfrak X}
       \nu_X(\mathfrak e|_{X^\sim};d\xi_X)\,
                 e^{\sum_Y W_{\mathfrak X}(Y;\mathfrak e,\xi)},
 \qquad W_\emptyset=0.
\]
The domains and the functions \(W_{\mathfrak X}\) are unchanged on
every retained family. Both brackets contain the positive empty term
one, and all their summands are nonnegative at real data. Thus
\(1\le B_t^{\rm cl}\le B_t\). In particular their logarithms are real
and well defined. The possibly extensive empty-family exponential
cancels in their quotient.

The abstract bracket expansion applies to both \(B_t\) and
\(B_t^{\rm cl}\) with the same majorant. Indeed, deletion is implemented
by replacing the component measure by zero on the forbidden
components. It changes neither compatibility nor the consistency
identities of the remaining interactions, and it decreases every
absolute component-mass bound. Write their localized logarithms as
\[
 \log B_t=\sum_Y L_t(Y),\qquad
 \log B_t^{\rm cl}=\sum_Y L_t^{\rm cl}(Y).
\]
The local exponentiation estimate
\eqref{eq:local-exponentiation-bound} in
\cref{sec:mass-exponentiation} gives, for some fixed \(c_*>0\),
\begin{equation}\label{eq:clean-bracket-local-log-bound}
 |L_t(Y)|+|L_t^{\rm cl}(Y)|
 \le C e^{-(1+c_*)p_0(g_t)}
                  e^{-\kappa_*d^\natural_t(o;Y)},
 \qquad \kappa_*>c_{\rm E}.
\end{equation}
It also gives locality of each term. One can see the latter directly
from the finite-subset formula for the exponential of \(L_t(Y)\):
it uses only brackets restricted to subregions of \(Y\).
Therefore
\[
 L_t(Y)=L_t^{\rm cl}(Y)
 \quad\hbox{unless \(Y\) contains a deleted collared component.}
\]
The logarithm branches agree since they are obtained from the same
absolutely convergent connected expansion.

Let \(A_t\) be the set of level-\(t\) cubes at which a collar can
witness intersection with the forbidden neighbourhood, including the
fixed buffer. Put \(c_t=\#A_t\). A deleted collar lies in the output
small-field region and meets none of its fixed exempt hulls. Hence
any cube witnessing its intersection with \(A_t\) is free. It is not
necessary that the distinguished corner or the lowest cube of the
component be near the pair. For every changed \(L_t(Y)\) choose such
a witnessing cube \(q\in Y\cap A_t\). The rooted entropy bound for
the hull length is
\[
 \sum_{Y\ni q}e^{-\kappa_*d_t^\natural(o;Y)}\le C
 \quad\hbox{for every free cube $q$ and $\kappa_*>c_{\rm E}$.}
\]
It is uniform in the sizes of the fixed exempt hulls. Together with
\eqref{eq:clean-bracket-local-log-bound} it gives
\begin{align}
 0\le\Delta_t(\mathfrak e)
 &:=\log\frac{K_t(\mathfrak e)}{K_t^{\rm cl}(\mathfrak e)}
   =\sum_Y\bigl(L_t(Y)-L_t^{\rm cl}(Y)\bigr)\notag\\
 &\le C e^{-(1+c_*)p_0(g_t)}
       \sum_{\substack{q\in A_t\\q\ \mathrm{free}}}
       \sum_{Y\ni q}e^{-\kappa_*d^\natural_t(o;Y)}
 \le Cc_t e^{-p_0(g_t)}.
 \label{eq:clean-bracket-log-quotient}
\end{align}
This is a pointwise estimate at the common retained datum.
The last exponent is weakened to avoid choosing a particular value
of the spare constant \(c_*\). The difference is rooted at forbidden
components before summation; subtracting two extensive bounds on
the full logarithms would not give this conclusion.

\emph{The clean density at every prefix.}
Let \(j_0\) be the lowest level of the clean window. At \(j_0\) the
retained standing part decides whether its live region is clean;
all its earlier healed histories remain summed, so its clean density
is one whenever that standing part is admitted. For \(t>j_0\), write
\(a_t=Cc_t e^{-p_0(g_t)}\). At a clean output standing part,
deleting the forbidden components in the one-step bracket also
excludes the unclean predecessor standing parts. An offending
retained region can disappear from the outer data only inside a
component healed at this step. If it meets the excluded
neighbourhood, that component is among the deleted ones.

Suppose inductively that, on each admitted predecessor retained space,
the prefix-clean measure lies between \(e^{-e_{t-1}}\) and one times
the original resummed measure. Apply the original positive step
kernels to these inequalities, restricting to the current clean
strings. Their image from the original resummed inputs is represented
by \(K_t^{\rm cl}\); by
\eqref{eq:clean-bracket-log-quotient} it lies between \(e^{-a_t}\)
and one times the full output term. The prefix-clean output therefore
lies between \(e^{-(e_{t-1}+a_t)}\) and one times that term.
The argument holds before every subsequent marginalization and with
merges interpreted on the disjoint union of predecessor spaces.
Taking images preserves both inequalities.

More explicitly, let \(\delta_t\in[0,a_t]\) be the negative logarithm
of the density of the one-step-restricted output relative to the full
output, after taking the common retained-variable image. If
\(\mathbb E_t^{\rm ref}\) is the backward kernel of that restricted
step with its original resummed inputs, then its prefix-clean density
has the exact recursion
\[
 D_t=\delta_t-\log\mathbb E_t^{\rm ref}(e^{-D_{t-1}}),
 \qquad 0\le D_t\le a_t+e_{t-1}.
\]
At a merge, \(D_{t-1}\) is read on the disjoint union of predecessor
spaces. Thus the recursion does not identify two different
conditional laws or discard any stored variable.

Radon--Nikodym derivatives consequently give, for every clean standing
part \(h\) of level \(j\), on the same retained-variable space,
\begin{equation}\label{eq:inserted-clean-density}
 d\hat\tau_h^{\rm cl}=e^{-D_h}\,d\hat\tau_{[h]},
 \qquad 0\le D_h\le e_j,\qquad
 e_j=\sum_{j_0<t\le j}a_t
       \le C\sum_{j_0<t\le j}c_t e^{-p_0(g_t)}.
\end{equation}
One may pull this bounded Borel derivative back through a reference
representation \(e^{A_{[h]}}\mu_{[h]}\).
Only its bounded Borel values will be used in the conditional kernels.
The constants are independent of the retained variables, the torus,
the cutoff and the standing part.

The cube count is
\[
 c_t\le C(\mathsf r+\check R_{k_0})^4L^{4(k_1-t)}
       \quad(j_0<t\le k_1),\qquad c_t=0\quad(t\le j_0).
\]
The flow gives \(p_0(g_t)\ge cp_0\) in this finite window.
Since \(k_0-j_0\le C\check R_{k_0}\), geometric summation yields
\[
 \mathsf n e_{k_1}
 \le C L^{4\mathsf n}e^{C\check R_{k_0}}(1+\mathsf n)
       (\mathsf r+\check R_{k_0})^4e^{-cp_0}.
\]
Thus its cost is already included in
\eqref{eq:inserted-unclean-weight}. The normalized clean kernels have
the form
\begin{equation}\label{eq:inserted-clean-kernel}
 \mathbb E_jF=\mathbb E_j^{\rm ref}(Fw_j),\qquad
 e^{-e_{j-1}}\le w_j\le e^{e_{j-1}},\qquad
 \mathbb E_j^{\rm ref}w_j=1,
\end{equation}
as follows by normalizing the two common retained-variable measures.

\emph{Normalising the clean kernels.}
The pointwise clean-class representation is used before taking any
conditional law. Keep the same variables on both sides of
\eqref{eq:inserted-clean-density}, and at level $j-1$ write
\[
 d\mu_{j-1}^{\rm cl}=r_{j-1}\,d\bar\mu_{j-1},\qquad
 r_{j-1}=e^{-D_{j-1}},\qquad e^{-e_{j-1}}\le r_{j-1}\le1.
\]
The predecessor space is a disjoint union when several standing
parts merge. Let $\mathfrak Q_j^{\rm cl}$ be the original positive kernel of
step $j$, restricted to the strings that stay clean. Form the
reference joint measure
\[
 dJ_j^{\rm ref}(y,z)=d\bar\mu_{j-1}(y)\mathfrak Q_j^{\rm cl}(y,dz).
\]
It uses the original resummed input, its original counterterms and
the unchanged kernels of the run. The corresponding clean joint
measure is $r_{j-1}(y)dJ_j^{\rm ref}(y,z)$. Disintegrate the first
joint measure over its output $z$, and call its backward probability
kernel $Q_j^{\rm ref}(z,dy)$. The clean backward kernel is then
\[
 Q_j(z,dy)=
 \frac{r_{j-1}(y)}{Q_j^{\rm ref}r_{j-1}(z)}
                                  Q_j^{\rm ref}(z,dy).
\]
Its denominator lies in $[e^{-e_{j-1}},1]$. This proves
\eqref{eq:inserted-clean-kernel} with
$w_j=r_{j-1}/Q_j^{\rm ref}r_{j-1}$. Positivity of the lower
bound in \eqref{eq:inserted-clean-density} also makes the clean
and resummed measures equivalent on each retained space. The
backward kernels can therefore be used on common sets of full
measure when they are composed. Merging predecessor terms changes
the disintegration, but none of these bounds.

In particular $Q_j\ge e^{-e_{j-1}}Q_j^{\rm ref}$ as kernels of
positive measures. Composing the finitely many backward kernels
gives the lower bound with
$P_0=\exp(-\sum_{k_0<j\le k_1}e_{j-1})$.
Keep the actual final disintegration $\mathbb E^W$ of the clean
class over $W$. This positive normalized mean preserves the
inequality. It is not replaced by a separately normalized reference
terminal law.
Consequently
\[
 \Pi=P_0\Pi^{\rm ref}+(1-P_0)\Pi^{\rm rest},
\]
where $\Pi^{\rm rest}$ is a probability kernel when $P_0<1$;
when $P_0=1$ the two kernels agree. The reference kernel is the
composition of the $Q_j^{\rm ref}$, followed by that same final
conditional mean. A coin with probability $P_0$, independent of
the output class and $W$, realizes this decomposition. The total
exceptional mass is at most the unclean mass plus
$1-P_0\le\sum_{k_0<j\le k_1}e_{j-1}$.
Only bounded Borel values of $D_h$
are used. No derivative or analyticity of this density is required.


The resulting clean measure $\mathbb Q_{\rm cl}$, kernel
$\Pi_{\rm cl}$ and exceptional measure $\lambda_{\rm bad}$ satisfy
\begin{equation}\label{eq:inserted-clean-sum}
 \int G\,d\lambda_{k_0}
 =\int\Pi_{\rm cl}G(W)\,d\mathbb Q_{\rm cl}(W)
                  +\int G\,d\lambda_{\rm bad},\qquad
 \lambda_{\rm bad}(1)\le C\epsilon_{\rm uc}.
\end{equation}
The cost $\sum_{k_0<j\le k_1}e_{j-1}$ is absorbed in
\eqref{eq:inserted-unclean-weight}, with its unspecified positive
constant $c$ decreased if necessary. The kernel on the clean part
uses $\mathbb E_j^{\rm ref}$ only. In the following calculation
$\mathbb E_j$ denotes these reference kernels. Write
$F_{\rm ref}^{[j]}=\mathbb E_j\cdots\mathbb E_{k_0+1}F$ for their
successive transport; this differs from the original transport
$F^{[j]}$ before the exceptional part has been removed.

\emph{The small-field integral on a read set.}
The reference kernel in \eqref{eq:inserted-clean-kernel} has the
following small-field form on a clean box. Retain a bounded gauge-invariant
factor $F$ while performing the changes of variables. Criticality
removes the linear term of the action; the nonlinear straightening
changes only the pivots. For a small-field integral its normalized
form is
\begin{equation}\label{eq:inserted-small-field-law}
 \frac{\displaystyle\int_{\mathcal W_j}
 F\big(\mathfrak v(y;V_j)\big)e^{P_j(y;V_j)}
                         \,d\gamma_{C^{(j-1)}(V_j)}(y)}
 {\displaystyle\int_{\mathcal W_j}
 e^{P_j(y;V_j)}\,d\gamma_{C^{(j-1)}(V_j)}(y)},
 \quad
 \mathcal W_j=\prod_{b\ {\mathrm{free}}}\{|y_b|<\mathsf b_j\},
\end{equation}
where $\mathsf b_j\asymp p_0$ and
\[
 \mathfrak v(y;V_j)=
 \exp\!\left(i\left[g_{j-1}\mathsf Cy
                   -\widetilde{\mathsf D}(g_{j-1}\mathsf Cy)\right]\right)
                       V_{j-1}^{\rm bg}(V_j).
\]
The function $\widetilde{\mathsf D}$ begins at second order.
The product runs over the free bonds of that small-field integral.
The covariance is that of the actual step, before comparison with its
flat value. The formula retains $F$ in the changes of variables of
Ba\l aban\cite[Section~2, equations~(2.1)--(2.16)]{Balaban1987}.

\emph{The window and the gauge fixing.}
The reference kernel is formed from a resummed standing part,
rather than from a single inner history or from a new step applied
to the exponent $A_{[h]}-D_h$. The standing-part identity keeps all
variables on which a later move depends. Its strings differ in the
components healed at this step. On a clean class these components
lie outside the excluded neighbourhood. The box around the read
set is therefore inside the ordinary small-field integration
region.

The strip substitution is performed before the inner variables
are conditioned. On the collar and rim it leaves the variables
unchanged. In its cutoff zone the substitution switches strictly
between the thresholds $\delta_{j-1}$ and $2\delta_{j-1}$.
Where it has not yet reached the localized background, the
approximate test at $2\delta_{j-1}$ equals one for either background.
The kept test factors adjacent to the integration region
equal one on the carried support and are removed there; after pivot
elimination the other kept tests depend on the new field only
outside that region. Each pivot value is determined from its own
averaging constraint, rather than from the global averaged background.
The nonlinear straightening changes only the
pivot variables. The remaining restriction on each free inner bond
is precisely $|y_b|<\mathsf b_j$. Freezing all exterior variables
thus leaves the product window in
\eqref{eq:inserted-small-field-law}; it does not introduce a further
history-dependent inner indicator.

Condition first on the recorded frame as well. Every operation
outside the box is positive and reads only its retained variables.
Disintegrating these variables therefore gives a positive measure
$m(d\mathsf e)$ of exterior data. If $Z_{\mathsf e}>0$ is the
inner normalising integral, a normalized expectation has the form
\[
 \frac{\int m(d\mathsf e)Z_{\mathsf e}
                    E_{\mathsf e}[F\circ\mathfrak v_{\mathsf e}]}
      {\int m(d\mathsf e)Z_{\mathsf e}}.
\]
It is a mixture of the normalized local laws, with their
partition functions included in the mixing weights. The same
argument applies to merges and to the final conditioning on $W$.
One may then sum over frames. Replacing their normalization factors
by those of a single frame would change the measure.

A gauge-fixing move is realized as a probability kernel on the
gauge orbit. It preserves every gauge-invariant inserted function.
For such a function the invariant measure identity permits the
gauge factor to be written in the small-field integral. Under the
joint gauge action on the background and fluctuation coordinates,
each free-bond coordinate is rotated by an adjoint action on
$\mathfrak{su}(2)$, which preserves its Euclidean norm. The balls
in the product window and the Haar changes of variables are
therefore unchanged. Pivot coordinates remain determined by their
covariant averaging constraints.

The main terms in the carried expansion, including the contracted
one-line and two-line terms, are gauge-invariant functions of the
current field. Apply the displayed integral formula to these terms.
An arbitrary remainder is carried only by its positive normalized
kernel, which contracts its supremum norm. There is no need to
identify a gauge-fixing kernel with multiplication by a gauge
factor on a non-invariant remainder.

Finally, use the convention which redraws every healed coupling
profile and compensates its change in the Wilson counterterm at
each plaquette. The total bare Wilson coefficient is unchanged.
On the clean box the redrawn profiles all equal one, so the local
coefficient is the running $x_{j-1}$. Healed components below the
finite exclusion window remain summed in the original deterministic
$R^{(i)}$ families. No condition excluding all earlier creations,
and no conditioning on an unreconstructed deficit history, is used.

The exterior of $\mathcal B_j$ has to be removed from
\eqref{eq:inserted-small-field-law}. Freezing it gives a mixture of
laws on the product window with energies
\[
 H_{\mathsf e}(y)=\tfrac12\langle y,Q^{\mathsf e}_{II}y\rangle
       +\langle Q^{\mathsf e}_{IO}y_O(\mathsf e),y\rangle
       -P_{\mathsf e}(y),\qquad I=\mathcal B_j.
\]
The reference Gaussian in the box is the marginal of the full
reference Gaussian. Its precision is the Schur complement
$Q_{II}-Q_{IO}Q_{OO}^{-1}Q_{OI}$, rather than $Q_{II}$.
The RG gap, exponential decay and the small Hessian of $P_j$ make
this reference energy uniformly convex after conjugation by a small
exponential weight. The difference from $H_{\mathsf e}$ has gradient
bounded at a bond $b$ by
$C\mathsf b_j e^{-c\,\mathrm{dist}(b,O)}$.

Here is the weighted comparison used for this step. Let $H_0$ be
the reference energy on a product of convex bond windows, let $H_1$
be the energy with exterior datum fixed, and put
$T_\varphi z=(e^{\varphi(b)}z_b)_b$. The lower gap and the small
interaction Hessian must give, for every sufficiently slowly varying
$\varphi$, the uniform bound
\[
 \langle z,T_\varphi H_0''T_\varphi^{-1}z\rangle
                                      \ge\lambda|z|^2,
 \qquad\lambda>0.
\]
Write $\delta_b=\sup|\partial_b(H_1-H_0)|$ on the window.
First extend the energies off the product window. Retract each bond
coordinate smoothly into a compact interior set, identically on a
slightly smaller window, and compose the interaction with this
retraction. Choose common smooth convex walls whose positive
block-diagonal Hessians outside that smaller window dominate the
extra Hessian terms of the retraction. The quadratic Gaussian part is
retained. The extended energies are smooth and coercive; their force
difference has at most linear growth, and the common walls cancel
from it. Their block-diagonal Hessians commute with $T_\varphi$.
The wall limit is taken first at a fixed smaller window and only
then as that window increases to the original one.

Couple the Hamiltonian flows of $H_0,H_1$ with common momenta and
set $v=T_\varphi(q^1-q^0)$. After smooth convex walls have been
introduced at each bond, their Hessians are bounded by a finite
number $\Lambda$. For a sufficiently short time $\tau$, the
weighted coercivity and the integral form of the flow equation give
\[
 \mathbb E|v(\tau)|
 \le(1-\tfrac12\lambda\tau^2+C\Lambda^2\tau^4)
                      \mathbb E|v(0)|
 +\tfrac12\tau^2(1+C\Lambda\tau^2)
            \nu_{H_1}|T_\varphi\nabla(H_1-H_0)|.
\]
Refresh the common Gaussian momentum after each interval. The two
position marginals retain their Gibbs laws, since each Hamiltonian
flow preserves its Gibbs-position and Gaussian-momentum product.
Iteration and then $\tau\downarrow0$ bound a Lipschitz expectation
difference by $\ell_G/\lambda$ times the expected weighted force.
The at-most-linear growth and the common coercive walls permit
dominated convergence in this expectation as the walls increase.
On the resulting hard window it is bounded by the supremum force
components $\delta_b$. Thus
\[
 |\nu_{H_1}G-\nu_{H_0}G|
 \le\frac{\ell_G}{\lambda}
                       \Big(\sum_be^{2\varphi(b)}\delta_b^2\Big)^{1/2},
 \quad
 |G(y)-G(y')|\le\ell_G|T_\varphi(y-y')|.
\]
The weighted lower bound is retained in these limits, since the
walls act separately on the bond coordinates.

The product shape is also needed to differentiate a marginal
density without differentiating its integration domain. If $m_a$
is the read-set marginal and $\nu_a^u$ the conditional law of the
other bonds at read-set value $u$, then
\[
 \partial_b\log(m_1/m_0)(u)
 =-\nu_1^u\partial_b(H_1-H_0)
               -(\nu_1^u-\nu_0^u)\partial_bH_0.
\]
Apply the preceding expectation estimate to the last difference,
using $\varphi=-\mu\,\mathrm{dist}(\cdot,b)$. The weighted
Lipschitz norm of $\partial_bH_0$ is bounded by the exponentially
decaying mixed Hessian. The weighted force norm is at most
$C\mathsf b_je^{-c\mathsf D}$ when $b$ belongs to the read set.
Integrating this derivative along segments of the read-set window
and using that both densities have integral one gives the density
ratio in $[e^{-\Xi_j},e^{\Xi_j}]$, with
\begin{equation}\label{eq:inserted-box-decoupling}
 \Xi_j\le C n_jp_0^2e^{-c\mathsf D}.
\end{equation}
For the polymer part of the force, a domain reaching the exterior
spends part of its exponential decay on the distance to that
exterior. A domain lying inside the box instead uses the locality
of the background entering its activity. In both cases the decay
left over sums the differentiated domains. The same localized
analytic expansion is needed for the Hessian bound in the convex
comparison. A small supremum for the whole exponent would not
give either spatial estimate by itself.

\emph{The exponent on the box.}
On the scaled fluctuation window write the exponent as the sum of
the new interaction, the increments of the regular activities and
the increments of the regenerated activities. Boundary activities
are anchored outside the excluded neighbourhood and contribute
through the exterior-force estimate. We use the redrawn-profile
convention already fixed above, so that no coupling deficit remains
inside a clean box.

After their critical value and Hessian have been extracted, the new
action and gauge-fixing remainders start at third order
in the unscaled displacement $g_{j-1}\mathsf Cy$. Their action
coefficient is $x_{j-1}=g_{j-1}^{-2}$. Thus on
$|y_b|<\mathsf b_j\asymp p_0$ they contribute at most
$C\mathsf g p_0^3$ per current unit volume. The Haar and
change-of-variable factors have an explicit factor of
$g_{j-1}$ and obey the same coarser bound. More precisely, their
odd and even parts have sizes
$C\mathsf g p_0^3$ and $C\mathsf g^2p_0^4$, respectively.
The odd part is harmless for the bounded-function comparison
below; the finer parity estimates are not needed there.

For an earlier regular activity at level $i$, the displacement
lies in a complex ball of radius
$\alpha_{\cdot,i}\asymp g_i(\log x_i)^q$. Its $m$th Taylor
term is therefore bounded by its activity norm times
\[
 \left(\frac{Cp_0}{(\log x)^q}\right)^m.
\]
There is no additional positive power of $\mathsf g$ in this
quotient. At a flat field the first derivative vanishes: constant
gauge transformations fix the field and act by the adjoint
representation on each direction, and $\mathfrak{su}(2)$ has no
fixed vector. Moving the derivative to a regular nearby background
costs one additional background-curvature-to-radius factor.
Consequently both the first-order contribution and the quadratic
contribution are bounded by
$C(\log x)^{2(\mathfrak p_0-q)}$ per current unit volume.

This estimate must be summed after the two subtractions in
\eqref{eq:rg-old-interaction}. At age $j-1-i$ there are
$L^{4(j-1-i)}$ old domains per current unit volume. The two
curvature factors each supply $L^{-2(j-1-i)}$, which compensates
this count exactly and leaves no decay in the age. The additional
summable weight comes from subtracting the action coefficient at
the pointed tents before forming the localized complex increment.
The tents add to one in the interior of the box. The resulting
irrelevant remainder has the weighted bound already obtained for
the old interaction; its geometric age factor absorbs the
polynomial change of the radii with age. This is a bound for the
localized increment after cancellation of the pointed totals,
rather than an assertion that each unsubtracted old term has the
gain separately.

For $R^{(i)}$, the flat value is subtracted but there is no
coupling counterterm. Two Cauchy derivatives cost at most $Cx_i$,
and the two field factors contribute at most $Cp_0^2$. After
the same cancellation of domain counts and curvature scales, the
remaining sum is
\[
 Cp_0^2\sum_{i<j}x_i^{1-\kappa_0/2}
                      \le Cp_0^2x^{2-\kappa_0/2}.
\]
For the last inequality use the lower flow bound
$x_i\ge x_{j-1}+c(j-1-i)-C$ and compare the sum with
$\int_0^\infty(x+ct)^{1-\kappa_0/2}\,dt$.
It converges precisely in the range used here, $\kappa_0>4$.
No summable geometric age factor for these terms has been assumed.

Adding these bounds over the $n_j$ free coordinates in the box
gives
\begin{equation}\label{eq:inserted-tilt-size}
 \theta_j:=\sup|P_j|
 \le Cn_j\left[
  \mathsf g p_0^3+(\log x)^{2(\mathfrak p_0-q)}
                     +p_0^2x^{2-\kappa_0/2}\right].
\end{equation}
The analyticity margin must be retained while obtaining the
Hessian. The new interaction is localized in domains $Y$ and is
analytic on a fluctuation ball of radius $2\mathsf b_j$; two
Cauchy derivatives give $C\mathsf g p_0$. The localized old
increments are analytic on radius
$\alpha_{\cdot,i}/(Cg_{j-1})$, larger than the window. Two
derivatives give $C(\log x)^{-2q}$ times their summable norm;
the $R$ sum is smaller at $\kappa_0>4$. Summing only domains
containing both differentiated bonds retains exponential decay
in their separation. Thus the mixed-Hessian row bound needed in
the convex comparison is small, uniformly in the box size. This
uses the complex increment estimate before restriction to real
fields; differentiating a bound stated only at real backgrounds
would not justify it.

\emph{Five comparisons for the read-set law.}
Let $I=\mathcal B_j$ be the box and let
$S=\mathcal S_j\subset I$ be its read set. The Gaussian in
the box comparison is the marginal of the full Gaussian on $I$.
Its precision is the Schur complement already computed above.
Let $\gamma_j^0$ be the flat Gaussian marginal on the free
coordinates of $S$, with the same pivot convention, and put
\[
 \mathcal W_j^{\rm read}
       =\prod_{b\in S}\{|y_b|<\mathsf b_j\}.
\]
The law $\gamma_j^0$ has one identical spatial covariance in
each colour. Its upper and lower covariance bounds are uniform
in $j$ and the volume, by the compressed quadratic-form bounds.

Let $\Psi$ be bounded and Borel on this read-set window.
Extend it to the other read-set values by zero when comparing
unrestricted Gaussian laws. First the density estimate
\eqref{eq:inserted-box-decoupling} removes the exterior data
at cost at most
$(e^{\Xi_j}-1)\|\Psi\|_\infty$. The law left on $I$ is
the box Gaussian restricted to its product window and tilted
by $P_j$.

Second remove that tilt. For real data,
$|P_j|\le\theta_j$ gives
\[
 \left|\frac{e^{P_j}}{\langle e^{P_j}\rangle_0}-1\right|
                        \le e^{2\theta_j}-1\le C\theta_j,
\]
where $\langle\cdot\rangle_0$ is the windowed Gaussian
expectation and $\theta_j$ is small. The change of a bounded
expectation is therefore at most
$C\theta_j\|\Psi\|_\infty$.

Third remove the window on the entire box. Each free coordinate
has bounded Gaussian variance, so a union bound gives
\[
 t_j\le Cn_je^{-cp_0^2}
\]
for its complement. Replacing conditioning on an event of
probability at least $1-t_j$ by the unconditioned law changes
a bounded expectation by at most $2t_j\|\Psi\|_\infty$.

Fourth rotate the fluctuation into the axial frame and replace
its Gaussian covariance on $S$ by the flat one. Locality and
analytic dependence of the covariance, applied to the small
axial potential, give the entrywise error
\[
 \varpi_j\le C[\mathsf r\mathsf g p_0+e^{-c\mathsf D}].
\]
Here the first term is the size of the potential in that frame;
the second is the error of localization. The resolvent identity
is used on the box with its retained regularity margin. In
dimension $m=3\#S$, the Hilbert--Schmidt error is at most
$Cm\varpi_j$. To see why it controls a bounded expectation,
normalize by the flat covariance and diagonalize the resulting
relative covariance. For eigenvalues close to one, comparing
the finite Gaussian densities gives total variation bounded by
the Euclidean norm of their deviations. If that norm is bounded
below by a fixed constant, the trivial total variation bound
gives the same inequality after its constant is enlarged.
The uniform lower covariance bound therefore gives the cost
$Cn_j\varpi_j\|\Psi\|_\infty$.

Finally restore the read-set window in the flat Gaussian.
Its complement has the same Gaussian tail bound, and the cost
is again at most $Ct_j\|\Psi\|_\infty$. The five comparisons
prove
\begin{equation}\label{eq:inserted-read-comparison}
 \left|\mathbb E_j\Psi(\widetilde y_j)
              -\gamma_j^0(\Psi\mid\mathcal W_j^{\rm read})\right|
 \le T_j\|\Psi\|_\infty,
 \quad
 T_j\le C[(e^{\Xi_j}-1)+\theta_j+t_j+n_j\varpi_j].
\end{equation}
The constants depend on the fixed geometric data and the fixed
choice of radii, not on the cutoff, the class or the exterior
datum. Positivity of the exterior and frame mixtures preserves
all five bounds. This comparison applies to an arbitrary bounded
read-set function; it uses the localized potential and the exact
window as well as Gaussian covariance decay.

\emph{Lowering a plaquette and telescoping the generating function.}
For a bare plaquette $p$, let $P_p^{(j)}(A)$ be its holonomy
minus one in the Landau representative of the localized
background $U_{j,\square}(e^{iA})$. The input potential is given
on the level-$j$ bonds of a cube about the pair. The cube has the
fixed localization margin, and the next read set includes all
the bonds on which this background depends. At $j=k_0$ this
includes the finitely many test cubes used in its definition.

The continuation from the trivial datum is analytic on a fixed
ball $\|A\|<c$, independent of $j$. Its curvature bound is
$C\eta_j^2\|A\|$, with the same margin in this complex ball.
At zero its derivative is the linear minimizer computed above.
Taylor's formula therefore gives
\[
 P_p^{(j)}(0)=0,\qquad
 DP_p^{(j)}(0)A=i(\mathcal H_jA)(p),\qquad
 |P_p^{(j)}(A)-i(\mathcal H_jA)(p)|
                                       \le C\eta_j^2\|A\|^2.
\]
In particular the constant is not obtained by composing a
one-step regularity constant $j$ times. It comes from the
background problem at level $j$ on this one cube. The same
continuation and Cauchy estimate apply to its localized version.

We explain the one-level comparison, including the frame error.
Set $a_*=C\mathsf r\mathsf g p_0$. The segment is taken so
that $a_*\ll1$. At real data put the level-$j$ field in the
axial gauge of the box and extend this gauge transformation
constantly on its blocks. In this frame write the fluctuation as
\[
 V_{j-1}=e^{ia}V_{j-1}^{\rm bg}(V_j),\qquad
 a=g_{j-1}\mathsf C\widetilde y_j
                -\widetilde{\mathsf D}(g_{j-1}\mathsf C\widetilde y_j).
\]
The current fields, the background input and their Landau
representatives must all be taken in their regularity windows
on the common interior cube of the two localizations; the
smaller representative is restricted there, never extended.
The background response bound gives potential
size at most $Ca_*$ there; its fine bond angles are at most
$C\eta_{j-1}a_*$. These are the uniform small-representative
bounds used in the following comparison.

The straightening begins at second order. Analytic dependence
of the pivot elimination on the background and the
Baker--Campbell--Hausdorff expansion consequently give the
linear fluctuation $g_{j-1}\mathsf C^0\widetilde y_j$, with
an error $O(a_*^2)$ in potential norm. Here $\mathsf C^0$
is the flat pivot elimination. Taylor expansion of
$P_p^{(j-1)}$ about the background, and replacement of its
derivative there by its derivative at zero, therefore leave
\[
 ig_{j-1}(\mathcal H_{j-1}\mathsf C^0\widetilde y_j)(p)
\]
as the first-order term, with error bounded by
$C\eta_{j-1}^2a_*^2$.

For the remaining background term use nesting before
localization:
$U_{j-1}(V_{j-1}^{\rm bg}(V_j))$ and $U_j(V_j)$ belong to
the same critical orbit. Localization changes this comparison
by a field whose curvature is at most
$C\eta_{j-1}^2\mathsf g p_0e^{-c\mathsf D}$ on the cube
of the pair. The orbit identity is modulo a spatial gauge
transformation; it is not an exact identity between the chosen
potentials, nor does the gauge have to be constant.

Here the small-representative bounds give the required
approximation by one constant rotation. If two real bond
fields $W,Z$ on the common cube obey $Z=W^u$ and their bond
angles are bounded by $C\eta a_*$, their gauge relation along
a bond yields
\[
 |u(x+\eta e_\mu)-u(x)|\le C\eta a_*.
\]
Summing along paths from a fixed root $x_0$ gives
$|u(x)-u(x_0)|\le C\mathsf r a_*$. A plaquette angle is
at most $C\eta_{j-1}^2\mathsf g p_0$, so replacing its
conjugating element $u(x)$ by the common element $u(x_0)$
costs at most
\[
 C\eta_{j-1}^2\mathsf g p_0\,\mathsf r a_*
                     \le C\eta_{j-1}^2a_*^2.
\]
For the localized comparison one first accounts for the small
localization field and then uses this gauge calculation for
the nested orbit. This last comparison concerns only the two
backgrounds at the fixed output $V_j$; its root rotation is
chosen before the current fluctuation is integrated. It may
depend on $V_j$, but not on $y_j$. Absorb it into the common
frame defining $\widetilde y_j$, once for all the plaquettes.
The result is the deterministic identity
\begin{equation}\label{eq:inserted-plaquette-lowering}
 P_p^{(j-1)}(A_{j-1})
 =ig_{j-1}(\mathcal H_{j-1}\mathsf C^0\widetilde y_j)(p)
                         +P_p^{(j)}(A_j)+R_j(p),
 \qquad |R_j(p)|\le\eta_{j-1}^2d_{\rm st},
\end{equation}
where
\[
 d_{\rm st}\le C[(\mathsf r\mathsf g p_0)^2
                          +\mathsf g p_0e^{-c\mathsf D}].
\]
The remaining freedom is this one constant frame. Constant
gauge covariance of the Landau continuation and colour
isotropy of the flat Gaussian allow it to be removed by the
group average in the generating function below. The error
has already paid for the spatial variation of the gauge;
no exact constant-gauge identification of the two localized
representatives is required.

For $a=1,2$ choose complex variables with radii
\begin{equation}\label{eq:inserted-generating-radii}
 |\lambda_a|\le\Lambda_a
       :=\frac{cx}{\mathsf r^4p_0^2\|\phi_a\|_\infty}.
\end{equation}
The constant $c$ is small and fixed. For independent flat windowed
fluctuations $\xi_{j'}$, set
\[
 m_{\le j}(p;\xi)=
 i\sum_{k_0<j'\le j}g_{j'-1}
                    (\mathcal H_{j'-1}\mathsf C^0\xi_{j'})(p),
\]
\[
 Z_j(P)=\int\prod_{k_0<j'\le j}d\gamma_{j'}^0(\xi_{j'}\mid\mathcal W_{j'}^{\rm read})
 \exp\!\left\{\frac12\sum_{a=1}^2\lambda_a\sum_p
                  \phi_a(p)|m_{\le j}(p;\xi)+P(p)|^2\right\}.
\]
Define $Y_j(V_j)$ by putting $P=(P_p^{(j)}(A_j))_p$ and averaging
over the common constant gauge rotation. These are bounded,
gauge-invariant Borel functions. They remain functions of the same
kind even after the original profiles occupy less than one spacing:
only the plaquette functions of the smooth minimizer are carried.
The normalization \eqref{eq:inserted-generating-radii} gives
$|Y_j|\le C$ and
\[
 |Z_j(P')-Z_j(P)|
 \le \frac{C}{\eta_{k_0}^2\mathsf g p_0}
                                 \max_p|P'(p)-P(p)|.
\]
On the clean part $Y_{k_0}=
\exp(\lambda_1\mathsf O_1+\lambda_2\mathsf O_2)$.
The definition of $Z_j$ also gives the exact one-step integration
identity obtained by adding its next independent flat fluctuation.
Use \eqref{eq:inserted-plaquette-lowering} first and then
\eqref{eq:inserted-read-comparison}. This yields
\[
 \|\mathbb E_jY_{j-1}-Y_j\|_\infty
 \le C\left[
 L^{-2(j-1-k_0)}\frac{d_{\rm st}}{\mathsf g p_0}+T_j\right].
\]
The outer positive normalized maps contract every accumulated error.
Consequently, for $F=e^{\lambda_1\mathsf O_1+\lambda_2\mathsf O_2}$,
\begin{equation}\label{eq:inserted-hybrid-telescope}
 \|F_{\rm ref}^{[k_1]}-Y_{k_1}\|_\infty
 \le C\sum_{j=k_0+1}^{k_1}
 \left[L^{-2(j-1-k_0)}
                         \frac{d_{\rm st}}{\mathsf g p_0}+T_j\right].
\end{equation}
There is no volume factor $L^{4\mathsf n}$ in this telescoping sum.
The terminal background is smaller by $L^{-2\mathsf n}$ in
plaquette units. Symmetry of the centered Gaussian removes its
linear contribution to $Z_{k_1}$, so replacing it by zero costs
$CL^{-4\mathsf n}$. Removing the Gaussian windows costs $Ct^{1/2}$,
where $t\le C\sum_jn_je^{-cp_0^2}$; Cauchy--Schwarz applies because
the generating function with twice the small parameter still has a
bounded Gaussian integral.

The resulting full Gaussian covariance, for one colour, is
\[
 \mathcal C^x=\sum_{j=k_0+1}^{k_1}x_{j-1}^{-1}\mathcal C_{(j)}.
\]
For $\Phi_\lambda=\lambda_1\phi_1+\lambda_2\phi_2$, its generating
function is
\begin{equation}\label{eq:inserted-gaussian-determinant}
 Z^0(\lambda)=\det(1-\Phi_\lambda\mathcal C^x)^{-3/2}.
\end{equation}
The power $3/2$ counts the three colours of $\SU(2)$. Positivity and
$\mathcal C^x(p,p)\le Cx^{-1}\eta_{k_0}^4$ give
\[
 \| (\mathcal C^x)^{1/2}\Phi_\lambda
                         (\mathcal C^x)^{1/2}\|_1
 \le\sum_a|\lambda_a|\sum_p|\phi_a(p)|\mathcal C^x(p,p)
 \le Cc p_0^{-2}.
\]
Thus the determinant logarithm is its convergent trace series, and
$c_1\le|Z^0|\le C_1$ on the polydisc. Conditional on the clean
part, \eqref{eq:inserted-hybrid-telescope} proves a bound independent
of the terminal class and $W$. Integrating it in
\eqref{eq:inserted-clean-sum} therefore gives
\begin{equation}\label{eq:inserted-generating-comparison}
 \left|\frac{\int e^{\lambda_1\mathsf O_1+\lambda_2\mathsf O_2}
                                  d\lambda_{k_0}}{Z^0(\lambda)}-1\right|
 \le C\delta,
 \quad
 \delta:=\epsilon_{\rm uc}+L^{-4\mathsf n}+t^{1/2}
       +\sum_j\left[
       L^{-2(j-1-k_0)}\frac{d_{\rm st}}{\mathsf g p_0}+T_j\right].
\end{equation}
In particular this argument makes no use of the conditional law on a
fibre of $\mathcal M_{k_0\to k_1}$. The exceptional set of terminal
data can be chosen independently of $\lambda$: first take a
countable dense set in the polydisc, and then use continuity of the
bounded generating functions.

\emph{The relative error and the mixed Gaussian slices.}
For small $\delta$, take the logarithm of the ratio in
\eqref{eq:inserted-generating-comparison} and apply the mixed Cauchy
estimate. It gives
\begin{equation}\label{eq:inserted-cumulant-comparison}
 \left|\operatorname{Cov}_{\lambda_{k_0}}(\mathsf O_1,\mathsf O_2)
       -\frac32\sum_{p,q}\phi_1(p)\phi_2(q)
                                      \mathcal C^x(p,q)^2\right|
 \le C\delta\Lambda_1^{-1}\Lambda_2^{-1}
 \le C\delta\mathsf r^8p_0^4x^{-2}
                       \|\phi_1\|_\infty\|\phi_2\|_\infty.
\end{equation}
The Gaussian coefficient follows by differentiating
$\frac32\sum_{m\ge1}m^{-1}\tr(\Phi_\lambda\mathcal C^x)^m$.
In particular
\[
 \mathcal C^x(p,q)^2
   =\sum_{j,j'}x_{j-1}^{-1}x_{j'-1}^{-1}
                            \mathcal C_{(j)}(p,q)\mathcal C_{(j')}(p,q).
\]
The mixed slices $j\ne j'$ are retained. They are exactly the terms
which, in the free conditional-covariance calculation above, come
from the bilinear curvature contribution and the covariance of the
conditional means.

The bound on $\delta$ must beat the mixed Cauchy loss
$\mathsf r^8p_0^4=O((\log x)^{20\mathfrak p_0})$.
The local estimates above give its contributions, before that loss,
as
\begin{align*}
 &C[\mathsf r^2\mathsf g p_0+e^{-c\mathsf D}],\\
 &C\mathsf n\mathsf r^4
       [\mathsf g p_0^3+p_0^2x^{2-\kappa_0/2}],\\
 &C\mathsf n\mathsf r^4(\log x)^{2(\mathfrak p_0-q)},\\
 &C\mathsf n\mathsf r^4
       [\mathsf r\mathsf g p_0+e^{-c\mathsf D}],\\
 &C\mathsf n\mathsf r^4p_0^2e^{-c\mathsf D},\\
 &C\big(\mathsf n\mathsf r^4e^{-cp_0^2}\big)^{1/2},
       \qquad CL^{-4\mathsf n},\\
 &C L^{4\mathsf n}e^{C\check R_{k_0}}(1+\mathsf n)
       (\mathsf r+\check R_{k_0})^4e^{-cp_0}.
\end{align*}
They represent, respectively, deterministic lowering, the ordinary
and remainder tilts, the earlier regular activities, covariance
replacement, exterior decoupling, Gaussian windows, the terminal
background and exceptional histories. With the already chosen
$\mathsf n=\lceil2\log_L\bar x_s\rceil$, every contribution tends
to zero after multiplication by the Cauchy loss. Powers of
$\mathsf g$, and the remainder power $x^{2-\kappa_0/2}$, beat all
logarithmic factors. The terminal term is
$O(x^{-8})$ up to fixed flow-comparison constants. The exceptional
term is small since $\mathfrak p_0>\max\{1,\mathfrak r\}$.
The only contribution with neither a power of $\mathsf g$ nor an
exponential is bounded after the loss by
\[
 C(\log x)^{1+30\mathfrak p_0-2q}=o(1).
\]
Thus the earlier choice $q>21\mathfrak p_0$ is sufficient here too.

Finally replace $\mathcal C^x$ by
$x^{-1}C_{K,s}^{\rm bg}$. The drift costs
$C(B_\sharp\mathsf n+C_\sharp)/x$, the omitted terminal covariance
costs $C(\mathsf rL^{-\mathsf n})^4$, and localization costs
$C\mathsf r^4e^{-c\mathsf D}$. These are comparisons of the explicit
Gaussian covariance sums. The slice envelopes sum over the two
separated supports with a bounded total, so these three errors
carry no mixed Cauchy loss. Combining
\eqref{eq:inserted-cumulant-comparison} with the earlier source,
localization and live-component estimates yields the target
\[
 \left|C_2(f,h)-\zeta'_{k_0}{}^2\frac32x^{-2}
       \sum_{p,q}\tilde f(p)\tilde h(q)
                          C_{K,s}^{\mathrm{bg}}(p,q)^2\right|
 \le\epsilon_{g,\vartheta}(d)x^{-2}
                         \|f\|_\infty\|h\|_\infty,
 \qquad\epsilon_{g,\vartheta}(d)\longrightarrow0.
\]
Here $g$ and $\vartheta$ are fixed, $K-s(d)\ge n_0$, and $d$ is
small enough that the prescribed segment and boxes fit in the run
and in the reference torus. The clean-class density and kernel
comparison above uses the original retained-variable representation
of the run and its small-field product window. The quantitative
comparison also uses the uniform regularity, locality, complex
increment and history bounds required by the preceding
renormalisation and source arguments.

\begin{proposition}[The reference-torus two-point function]
\label{prop:reference-two-point}
The reference-torus comparison gives
\[
 C_2(f,h)=\zeta'_{k_0}{}^2\frac9{\pi^4}
                      g(d)^4I(f,h)(1+\varepsilon),
 \qquad \sup|\varepsilon|\longrightarrow0\quad(d\downarrow0)
\]
over the specified data and shape class.
\end{proposition}

\begin{proof}
At separation of order $d$, the slice error tends to zero because
$\hat d\to\infty$. The lattice error tends to zero since
$a_K/d=L^{-(K-s)}/\hat d\le L^{-n_0}/\hat d\to0$ and
$d/\ell_\star(g)\to0$. The frozen-tent interpolation changes a
test by $o(\|f\|_\infty)$ and its support by $o(d)$. The shape
bounds control the resulting sum error relative to $I(f,h)$. The
Riemann-sum error is bounded by $Ca_K/(\vartheta^3d)$, which tends
to zero under the same scale condition. Thus
\[
 \sum_{p,q}\tilde f(p)\tilde h(q)
                         C_{K,s}^{\mathrm{bg}}(p,q)^2
                         =\frac6{\pi^4}I(f,h)(1+o(1)).
\]
For jointly Gaussian vectors $F,G$, Wick's identity gives
$\operatorname{Cov}(|F|^2/2,|G|^2/2)
=\tfrac12\sum_{\mu,\nu}(\mathbb EF_\mu G_\nu)^2$.
There are three colour components for $\SU(2)$. Consequently the
coefficient is $(3/2)(6/\pi^4)=9/\pi^4$, as displayed.
The absolute error in \hyperref[hyp:relative-two-point]{the reference-torus comparison} becomes a relative
error because $I(f,h)$ has a positive shape-dependent lower bound.
\end{proof}

\subsection{Transfer from one torus to every phase}

\begin{proposition}[Local transfer between states]
\label{prop:witness-volume-transfer}
At one first-passage datum, let $\mathfrak s,\mathfrak s'$ be two permitted
tori or periodic Gibbs states. Suppose $a_K\le d\le c$ and the real
functions $f_i$ are supported in one ball of radius $5d$, with
$\Lip f_i\le\|f_i\|_\infty/(\vartheta d)$. There are no sign or
separation restrictions. Then
\[
 |C_n[\mathfrak s](f_1,\ldots,f_n)
       -C_n[\mathfrak s'](f_1,\ldots,f_n)|
       \le C_n(g,\vartheta)d^{4-n}\prod_i\|f_i\|_\infty.
\]
\end{proposition}

\begin{proof}
Let $\mathsf D$ be the difference of cumulants. The common-normalizer
estimate in \cref{prop:source-state-difference} gives
$|\mathsf D(F_1,\ldots,F_n)|\le C_n\prod_i\|F_i\|_\sharp$
for arrays in a unit ball. By \hyperref[hyp:gibbs-reflections]{the lattice symmetry and reflection argument}, both states
are invariant under the full lattice translations, so $\mathsf D$ is
invariant under a common lattice translation.

Choose a lattice grid of $N\ge cd^{-4}$ common translates of the
radius-$5d$ support ball, with mutually disjoint, positively separated
copies in a unit ball. This is possible with grid spacing a lattice
multiple in $[Cd,(C+1)d]$, since $a_K\le d$. For $n\ge2$ take
independent signs $\epsilon_v^{(i)}$, $1\le i<n$, and set
$\epsilon_v^{(n)}=\prod_{i<n}\epsilon_v^{(i)}$. Define
\[
 F_i=\sum_v\epsilon_v^{(i)}\tau_vf_i.
\]
Disjoint supports give $\|F_i\|_\sharp\le\|f_i\|_\sharp$.
Expanding $\mathsf D$ by multilinearity, the mean of its sign factor
is zero unless all translation indices coincide. Indeed, each
independent sign family pairs its index with the index of the final
entry. Hence
\[
 \mathbb E_\epsilon\mathsf D(F_1,\ldots,F_n)
          =\sum_v\mathsf D(\tau_vf_1,\ldots,\tau_vf_n)
          =N\mathsf D(f_1,\ldots,f_n).
\]
For $n=1$ the same identity follows by taking the unsigned sum and
using linearity. The array estimate now gives
\[
 |\mathsf D(f_1,\ldots,f_n)|
   \le Cd^4\prod_i\|f_i\|_\sharp
   \le C_n(g,\vartheta)d^{4-n}\prod_i\|f_i\|_\infty.
\]
\end{proof}

For $n=2$ the transfer cost is
$C_2(g,\vartheta)d^2\|f\|_\infty\|h\|_\infty$.
By \cref{prop:reference-two-point} the positive leading term is
comparable to $(\log(1/d))^{-2}\|f\|_\infty\|h\|_\infty$.
The transfer cost divided by this term tends to zero. Choosing $d$
small after $g$ gives a uniformly small relative error, for instance
$|\varepsilon|\le1/2$, in every periodic Gibbs state.

Together with the uniform flow normalization this proves
\eqref{eq:two-point-comparison} in every such state, subject to the
cutoff-depth condition. At each fixed $d$ that condition holds
eventually in a sequence $K\to\infty$. The two-sided bound is a fixed
closed condition on the two-point value and therefore passes to any
common continuum subsequence of pure phases. Recentring an observable
changes no second cumulant.

For the profiles fixed in \cref{thm:main}, the rescaling
$x=dX$, $y=dY$ gives
\[
 I(f_d,h_d)=d^{-8}I(f_1,h_1),\qquad I(f_1,h_1)>0.
\]
They obey one shape condition independent of $d$. Thus, with $g$ small
for this fixed shape, \eqref{eq:two-point-comparison} is exactly
\eqref{eq:nontrivial}. No uniform smallness threshold for all possible
bump shapes is required.


\section{The continuum limit and diagonal reflections}\label{sec:limits}

The volume limit is taken at fixed lattice spacing. The source estimates
then permit one subsequence in the lattice spacing on which all
separated correlations converge. We also need the same extraction for
a Gibbs state centred at the mean of a different state. This will
provide the coordinates for the phase decomposition in
\cref{sec:phases}.

\subsection{Gibbs states at fixed spacing}

At fixed $K,\beta$, enumerate the unoriented bonds of the infinite
lattice. The configuration space is a countable product of compact
copies of $\SU(2)$ and is compact and metrizable. Periodically extend
the measures on tori with sides tending to infinity. Their laws have
weakly convergent subsequences. Denote their limit points by
$\mathcal G^{\mathrm{tor}}_{K,\beta}$.

For a finite set $W$ of bonds, the conditional Gibbs kernel integrates
the bonds of $W$ with density proportional to
\[
 \exp\left[-\beta\sum_{p:p\cap W\ne\varnothing}a_p(U)\right]
\]
against their Haar measure, keeping the exterior bonds fixed. Only
finitely many exterior bonds occur. Its normalizing integral is strictly
positive and continuous in those bonds. The kernel therefore maps local
continuous functions to local continuous functions. For a sufficiently
large torus the conditional identity against a fixed local exterior
function is the infinite-lattice identity on that finite collection of
bonds. It passes to a weak limit. Approximation and the monotone class
argument extend it to the conditional Gibbs identity. This proves that
the limit points are Gibbs states.

The periodic Wilson measures are invariant under every lattice
translation and under the cubic group $W_4$. These identities involve
local continuous functions and persist in a weak limit. In particular
$\mathcal G^{\mathrm{tor}}_{K,\beta}$ is nonempty and consists of
translation-invariant Gibbs states. The larger class
$\mathcal G^{\mathrm{per}}_{K,\beta}$ consists of Gibbs states
invariant under some finite-index subgroup of lattice translations.
We do not identify these two classes. The pure states needed later
belong to the latter class by
\hyperref[hyp:gibbs-decomposition]{the Gibbs ergodic decomposition}.

At this fixed spacing $\cO_a(f)$ is a bounded local continuous
function for compactly supported $f$. Products and cumulants are finite
combinations of integrals of such functions. Thus a weak torus limit
preserves every fixed moment and cumulant. Every estimate uniform on
large tori passes to $\mathcal G^{\mathrm{tor}}_{K,\beta}$.
The comparison of \cref{sec:sources}, under
\hyperref[hyp:gibbs-reflections]{the lattice symmetry and reflection argument}, gives the estimates on the larger class.

\subsection{One subsequence for every separated correlation}

The estimates for tests in small separated balls imply bounds on
any fixed separated support sets. It is useful to make this
extension before choosing countable test families.

\begin{lemma}[Bounds on fixed separated supports]\label{lem:fixed-support-bounds}
Fix a reference coupling in the range of the state-uniform source
and one-point estimates. Let \(A_1,\ldots,A_n\) be compact,
pairwise disjoint subsets of \(\R^4\). There is a constant
\(C(A_1,\ldots,A_n)\), independent of the first-passage datum
and of \(\rho,\nu\in\mathcal G^{\mathrm{per}}_{K,\beta}\), such that,
for real \(f_i\in C_c^1(\R^4)\) with \(\operatorname{supp}f_i\subset A_i\),
\begin{equation}\label{eq:fixed-support-moments}
 \left|\rho\left(\prod_{i=1}^n
       [\cO_a(f_i)-\nu(\cO_a(f_i))]\right)\right|
       \le C(A_1,\ldots,A_n)\prod_i\|f_i\|_\sharp,
\end{equation}
for all sufficiently small \(a=L^{-K}\).
The same assertion holds for cumulants of order at least two.
The constant can be chosen using only the number of entries,
a positive lower bound on their support separation and
radii of balls containing the individual support sets.
\end{lemma}

\begin{proof}
For \(n=1\), the assertion is the one-point transfer.
For \(n\ge2\), choose \(R>0\) so that distinct \(A_i\)
have distance at least \(6R\), and put \(d=2R\).
Fix a smooth partition of unity by translates whose
supports lie in balls of radius \(R\), and retain for
each \(A_i\) only the finitely many pieces that meet it.
Balls belonging to distinct entries have distance at
least \(6R-4R=d\): each ball can extend at most
\(2R\) from its point of intersection with \(A_i\).
Write
\[
 f_i=\sum_m f_{i,m},\qquad f_{i,m}=\chi_mf_i.
\]
Every piece is real and belongs to \(C_c^1\); its support
is contained both in \(A_i\) and in its assigned ball.
The product rule in the sharp norm gives
\begin{equation}\label{eq:support-partition-norm}
 \sum_m\|f_{i,m}\|_\sharp\le K_i\|f_i\|_\sharp.
\end{equation}
The constants are fixed by the partitions. A fixed translated
partition at this scale gives a bound for their number and
derivatives in terms of \(R\) and the radius of a ball
containing \(A_i\). Thus the \(K_i\) have the dependence
claimed in the lemma.

Choose an \(L\)-adic \(\varpi>0\) large enough that all the
containing balls lie in
\((-\varpi+d/2,\varpi-d/2)^4\). For all sufficiently large
\(K\), one has \(a\le d/2\) and \(\varpi/a\in\mathbb N\).
Every tuple obtained by selecting one piece of each entry
therefore belongs to a common small-ball support class.
The state-uniform cumulant estimates apply to every subset
of such a tuple, and the one-point transfer bounds its
first cumulants. The finite moment--cumulant formula gives
\[
 \left|\rho\left(\prod_i
 [\cO_a(f_{i,m_i})-\nu(\cO_a(f_{i,m_i}))]\right)\right|
       \le c_n(d;R)\prod_i\|f_{i,m_i}\|_\sharp .
\]
The constant is independent of the enclosing cube and
the two states. Multilinearity and
\eqref{eq:support-partition-norm} now give
\[
 \sum_{m_1,\ldots,m_n}
           c_n(d;R)\prod_i\|f_{i,m_i}\|_\sharp
 \le c_n(d;R)\prod_iK_i\|f_i\|_\sharp .
\]
This proves the moment assertion. The same finite expansion
directly proves the cumulant assertion. The lattice
conditions were required only eventually for these fixed
support sets.
\end{proof}

In particular, for tuples \(\underline f,\underline h\)
whose corresponding entries have supports in the same
sets \(A_i\), multilinearity gives
\begin{equation}\label{eq:fixed-support-continuity}
 \begin{split}
 |M_a^{\rho,\nu}(\underline f)-M_a^{\rho,\nu}(\underline h)|
 \le{}&C(A_1,\ldots,A_n)\sum_i\|f_i-h_i\|_\sharp\\
 &\qquad\qquad\cdot
        \prod_{l\ne i}(\|f_l\|_\sharp+\|h_l\|_\sharp),
 \end{split}
\end{equation}
where \(M_a^{\rho,\nu}\) denotes the reference-centred
moment in \eqref{eq:fixed-support-moments}. This estimate
holds for all states under consideration with one constant.
It is the continuity estimate used below, including when
the support of an entry is too large to lie in one
small ball separated from the others.

\begin{proposition}[Common extraction]\label{prop:common-extraction}
Let $(K_j,\beta_j)$ be first-passage data with $K_j\to\infty$,
and choose $\nu_j\in\mathcal G^{\mathrm{tor}}_{K_j,\beta_j}$.
There is one subsequence on which the cumulants converge for every
order $n\ge2$ and every separated $C_c^1$ tuple on $\R^4$. The
limits satisfy the same per-ball and factorial bounds.

It is enough that the bounds hold eventually for each fixed order,
radius and positive separation. The $j$th family may be defined only
for tests in $B(0,r_j)$ with $r_j\to\infty$. The subsequence can
also be chosen to make any prescribed countable list of bounded scalar
sequences converge.

The same assertion holds for moments in arbitrary
$\rho_j\in\mathcal G^{\mathrm{per}}_{K_j,\beta_j}$, with the
observables centred at $\nu_j$, provided the state-uniform bounds and
the one-point transfer apply to both states.
\end{proposition}

\begin{proof}
The bounds in \cref{lem:fixed-support-bounds} hold eventually
for every fixed separated support list, uniformly in the states.
Equation \eqref{eq:fixed-support-continuity} gives the corresponding
continuity estimate on common support sets. Their separation is
fixed throughout each approximation.

Choose all finite lists of bounded open finite unions of rational
boxes whose closures are pairwise disjoint. In each support list
choose a countable family of real smooth tests dense in the
\(C^1\) norm for compactly supported \(C^1\) tests in that list.
There are countably many tuples over all orders and lists.
The resulting family has the following support-controlled
approximation property: for a separated tuple \((f_i)\)
and any open neighbourhoods \(O_i\) of its supports,
one can approximate it in the sharp norm by enumerated
tuples \((f_i^{(l)})\) with
\(\operatorname{supp}f_i^{(l)}\subset O_i\).
Indeed first choose a rational support list with
\(\operatorname{supp}f_i\) in its interior and its closure
in \(O_i\), retaining positive separation, and then use
that list's dense family.

Equation \eqref{eq:fixed-support-moments} bounds every
enumerated coordinate. Successively extract a convergent
subsequence for the first tuple, then the second, and so
on. The diagonal subsequence converges on all of them.


Every enumerated tuple is in the domain of the $j$th family and has
its uniform bound once $j$ is sufficiently large. Removing its finite
exceptional prefix changes no extraction. Thus the same diagonal
procedure applies to the stated growing domains and eventual bounds.
Interleave any additional countable bounded scalar sequences with the
test coordinates to select their limits at the same time.

Given any separated $C_c^1$ tuple, enlarge its supports slightly to
one of the support lists just chosen, retaining a positive separation.
Approximate each entry in the sharp norm by the corresponding dense
tests. The telescoping estimate controls the difference uniformly in
$j$. For any prescribed accuracy, choose the approximation first and
then $j$ large enough that its tuple values are Cauchy. The original
tuple values are consequently Cauchy. The same estimate makes their
limit independent of the approximants and proves continuity and the
claimed bounds. The use of $C^1$ tests is essential to this particular
countable extraction: the space of all Lipschitz functions is not
separable in the sharp norm.

For the final assertion, the one-point transfer gives
\[
 |\rho_j(\cO_{a_j}(f))-\nu_j(\cO_{a_j}(f))|
                       \le C_1(R)\|f\|_\sharp.
\]
Extract these linear forms along the same countable family and include
their diagonal selections. A moment in $\rho_j$, centred at $\nu_j$,
is a polynomial in these mean differences and in the higher cumulants
of $\rho_j$. Every factor has the required uniform bound and
continuity. The identical diagonal and approximation argument therefore
gives convergence of all these moments on one subsequence.
\end{proof}

The same argument identifies the topology of the possible limit
families. Fix the per-ball bounds at every order and at rational
positive separations and radii. Choose the dense test collections
closed under rational linear combinations within their support sets,
and enumerate the resulting separated tuples as $X^1,X^2,\ldots$.
If $M_q$ is the corresponding uniform bound for the $q$th value, the
evaluation map takes a family into the compact countable product
\[
 \prod_{q\ge1}[-M_q,M_q].
\]
On bounded families its coordinate topology is metrized by
\[
 d(C,C')=\sum_{q\ge1}2^{-q}
                      \min\{1,|C(X^q)-C'(X^q)|\}.
\]
The support-preserving telescoping estimate implies two facts. A
family is determined by these coordinates, and coordinate convergence
extends to convergence at every separated $C_c^1$ tuple. To check
the second fact, first approximate the tuple at a fixed positive
separation, uniformly for the whole bounded class, and only then use
coordinate convergence. Thus no uniform continuity estimate as
$d\downarrow0$ is needed.

The multilinearity and compatibility conditions on these coordinates
are closed finite relations. A coordinate limit satisfies them on the
dense collections; the same approximation argument extends them to
all tuples. For a new tuple, that approximation proves convergence of
the original families themselves at the tuple. Pass to the limit in
its original bound, with its actual radius and separation; no
continuity of the constants in those parameters is required.
Consequently the set of families satisfying
the prescribed bounds and fixed closed identities is a closed subset
of the displayed compact product. It is compact and metrizable, not
just sequentially bounded. The choice of rational support regions
and the retention of separation in the approximation are what make
this countable-coordinate description valid. This topology will also
be used for the limiting coordinates of phases.

For the reference state, the centred limit moments are
\begin{equation}\label{eq:centred-moments}
 S^{\mathrm c}(f_1,\ldots,f_n)
   =\sum_{\substack{\pi\text{ a partition of }\{1,\ldots,n\}\\
                         |B|\ge2\text{ for every }B\in\pi}}
                       \prod_{B\in\pi}C_{|B|}((f_i)_{i\in B}),
 \qquad S^{\mathrm c}(\varnothing)=1,
 \qquad S^{\mathrm c}(f)=0.
\end{equation}
This is the finite-cutoff moment identity for
$\cO_{a_j}(f)-\nu_j(\cO_{a_j}(f))$, followed by a limit. It requires
no convergence of either uncentred mean.

Several further consequences of the extraction will be used. First,
a fixed closed relation among finitely many correlation values passes
to the limit. This applies to equalities, uniform intervals and finite
positive quadratic forms. If a relation has varying coefficients or
bounds, one must first establish its limiting closed relation. Second,
moving tests can be included whenever their sharp-norm difference tends
to zero and their supports retain a common positive separation margin.
The telescoping estimate proves this without any further extraction.
Third, every volume-first family can be approximated by a sequence of
tori whose sides exceed any prescribed floors. At stage $j$ choose a
torus approximating the first $j$ countable test values of $\nu_j$ to
within $1/j$, with side above that floor; the same continuity argument
then gives the family on every separated tuple.

In making this last diagonal choice, both the cutoff $K_j$ and bare
coupling $\beta_j$ are held fixed while its torus side is chosen.
At fixed spacing the finitely many required observables and products
are local continuous functions, so weak approximation of $\nu_j$
approximates their moments. Cumulants are finite polynomials in those
moments. If moments centred at another state are also being retained,
include the required mean coordinates in this finite list. Choose
the side beyond the admissibility floor of that datum and every
prescribed additional floor. After this choice at each $j$, the
approximation error tends to zero on every enumerated coordinate.
The common extraction estimate extends it to the original continuum
family on all separated tuples. This construction asserts no
interchange of limits for an entire unselected sequence; it constructs
one diagonal sequence representing the already chosen volume-first
limit.

For clarity, continuity under translations is quantitative. For
$f\in C_c^1$ put
\[
 \varpi_f(\delta)=\max\left\{
        \|\nabla f\|_\infty\delta,
        CM\sup_{|x-y|\le\delta}|\nabla f(x)-\nabla f(y)|\right\}.
\]
Then $\|\tau_vf-\tau_{v'}f\|_\sharp\le\varpi_f(|v-v'|)$
and $\varpi_f(\delta)\to0$. For finitely many moving entries the
same telescoping formula gives uniform continuity on support sets with
a fixed separation margin. This will allow lattice translations and
moving lattice mirrors to pass to continuum translations and mirrors.

In the case of two different centring states, full lattice symmetry
also determines the form of the limiting first moment. Let
$\delta_j^{\mathrm{mean}}=\rho_j(a_p)-\nu_j(a_p)$, independent of the plaquette $p$ by
\hyperref[hyp:gibbs-reflections]{the lattice symmetry and reflection argument}, and put $b_j^{\mathrm{mean}}=6a_j^{-4}\delta_j^{\mathrm{mean}}$. Then
\[
 \rho_j(\cO_{a_j}(f))-\nu_j(\cO_{a_j}(f))
     =b_j^{\mathrm{mean}}\left(\frac{a_j^4}{6}\sum_pf(x(p))\right).
\]
The parenthesis tends to $\int f$: there are six unoriented
plaquette directions in four dimensions. Fixing one smooth $f_0$ of
integral one, the mean-transfer bound makes $b_j^{\mathrm{mean}}$ bounded, and the
common extraction makes it converge. Thus the limiting first moment is
$b\int f$, with $|b|$ bounded uniformly in the chosen states. The
higher cumulants remain those of $\rho_j$; they are not averages of
component cumulants.


\subsection{Inherited symmetries and coordinate reflections}

Let $\sigma$ be one of the limiting moment functionals just constructed.
It may be the reference-centred functional $S^{\mathrm c}$, or the
limit in a state $\rho_j$ with subtraction taken from $\nu_j$. For
the latter we use the mean-transfer bound as well as the source bounds.
Throughout this subsection the state symmetries and positivity required
on the lattice are those of \hyperref[hyp:gibbs-reflections]{the lattice symmetry and reflection argument}.

Consider the lattice isometry $\mathsf r(x)=wx+v$, where $w\in W_4$
and $v\in a_K\mathbb Z^4$. Pulling back bond configurations by
$\mathsf r$ preserves Haar measure and the Wilson action. The plaquette
observable is sampled at the plaquette centre. Thus relabelling its
sum gives the exact identity
\[
 \cO_a(f)\circ\mathsf r^*=\cO_a(\mathsf rf),
 \qquad (\mathsf rf)(x)=f(\mathsf r^{-1}x).
\]
The same identity holds for the centred representatives, since the
subtracted state is invariant. Hence every finite-cutoff moment and
cumulant has the lattice covariance identity.

The lattices are nested. A vector in
$\mathfrak D_L^4$, where $\mathfrak D_L=\bigcup_Ka_K\mathbb Z$,
belongs to every sufficiently fine lattice. Covariance under this dense
translation group and under $W_4$ therefore passes to the common
subsequence as a closed relation. To pass from the dense group to all
translations, use the translation modulus established above. If
$v_m\in\mathfrak D_L^4$ tends to $v$, each entry satisfies
\[
 \|\tau_{v_m}f-\tau_vf\|_\sharp\le\varpi_f(|v_m-v|)\longrightarrow0.
\]
The translated tuple retains its separation and lies in support balls
of uniformly bounded radius. Telescoping its entries in the per-ball
estimate proves $\sigma(\tau_vX)=\lim_m\sigma(\tau_{v_m}X)
=\sigma(X)$. The same argument applies to a finite collection of
different moving translations as long as a positive separation margin
is retained. Reality, multilinearity and permutation symmetry pass
through the finite moment--cumulant identities.

We recall the coordinate reflection calculation, because it identifies
the positive form whose limits we use. At fixed $K,\beta>0$, take a
bond mirror $\Sigma=\{x_4=c\}$ with
$c\in a_K(\mathbb Z+\tfrac12)$ and its partner on the torus. Write
each crossing bond as
\[
 U_b=(W_b^-)^{-1}W_b^+,
\]
with $W_b^-,W_b^+$ independent Haar variables. Their quotient is Haar,
so this replacement leaves integrals unchanged. It splits the remaining
variables into two independent, identically distributed collections
$\varsigma_-,\varsigma_+$ exchanged by reflection. For a crossing
plaquette with crossing bonds $b_1,b_2$, its positive-side path gives
a matrix $\Xi_p(\varsigma_+)$, and
\[
 \tr U_p=\tr\big(\Xi_p(\varsigma_+)
                                  \Xi_p(\varsigma_-)^*\big).
\]
In particular its trace is a sum of a matrix entry on one side times
the conjugate of the corresponding entry on the other. The real part
has the same property after entries and conjugate entries are both
included. Expanding the exponential of the crossing Wilson weights
therefore gives, for a bounded function $F$ on the positive side,
\begin{equation}\label{eq:coordinate-rp-expansion}
 Z\left\langle\overline{F\circ\theta_\Sigma^*}F\right\rangle
 =e^{-\beta|\mathsf p_\times|}
   \sum_{n\ge0}\frac{(\beta/4)^n}{n!}
       \sum_{\vec\gamma}
       \left|\int F e^{-\beta A_+}
                         \xi_{\vec\gamma}\,d\varsigma_+\right|^2
 \ge0.
\end{equation}
Here $\xi_{\vec\gamma}$ runs over products of entries of the
$\Xi_p$ and their conjugates, and $A_+$ is the action from plaquettes
lying entirely on the positive side. Every coefficient is nonnegative.
The exponential series converges absolutely on the compact integration
space, so the expansion and integrations can be interchanged. A
plaquette whose centre is in the open positive side has all its corners
there. Thus a test supported on that side produces an observable of
precisely the variables allowed in this calculation.

A site mirror has a different decomposition. No plaquette has points
on both open sides: the plaquettes either lie in the mirror or belong
to one closed side. Condition on all bonds lying in the mirror.
Given these bonds, the remaining bond integrations on the two sides
are independent and are exchanged by reflection. The product of
the two reflected integrals is the squared modulus of one side's
integral, multiplied by the positive mirror weight. Integration over
the mirror bonds proves positivity for functions of one closed side.
This proves the finite-volume coordinate site and bond assertions
without identifying the two factorizations.

For the continuum passage use the bond mirror
$\Sigma_j=\{x_4=a_j/2\}$, whose partner escapes compact sets in
the volume limit. Its action on a test is
$\theta_{\Sigma_j}f=\tau_{a_je_4}\theta f$, with
$\theta(x',x_4)=(x',-x_4)$. If all positive-time tests stay a fixed
distance from zero, their reflected tuple remains separated. The
translation modulus and per-ball estimates replace
$\theta_{\Sigma_j}$ by $\theta$ at an error tending to zero.
Consequently each finite reflection-positive quadratic form has a
nonnegative limit. In the periodic Gibbs-state route the site mirror
at zero is available directly from \hyperref[hyp:gibbs-reflections]{the lattice symmetry and reflection argument}.
The passage from finite tori to the full periodic Gibbs class is
explained in \hyperref[hyp:gibbs-reflections]{the lattice symmetry and
reflection argument}.

\subsection{The Hilbert space and the time-translation semigroup}

Let $\mathcal D_+$ be the complex vector space of finite formal linear
combinations of separated tuples of real smooth tests in $\{x_4>0\}$, including the
empty tuple. If $F=\sum_ic_iX_i$, write
$\overline F=\sum_i\overline c_iX_i$ and
$\Theta F=\theta\overline F$. Concatenation is used only when the
combined supports are separated. In particular $\Theta F\cdot G$
is defined for all $F,G\in\mathcal D_+$: the reflected and unreflected
supports lie in opposite open half-spaces. The form
\[
 \langle F,G\rangle_\sigma=\sigma(\Theta F\cdot G)
\]
is Hermitian and positive semidefinite. Denote its null space by
$\mathcal N$ and the completion of $\mathcal D_+/\mathcal N$ by
$\mathcal H_\sigma$. The empty tuple gives a unit vector $\Omega$.

Before taking the quotient, define
\[
 T_t^\circ F=\tau_{te_4}F\quad(t\ge0),
 \qquad U^\circ(v)F=\tau_{(v,0)}F\quad(v\in\R^3).
\]
They preserve $\mathcal D_+$ and fix the empty tuple. Translation
invariance gives
\begin{align*}
 \langle F,T_t^\circ G\rangle_\sigma
 &=\sigma(\Theta F\cdot\tau_{te_4}G)
   =\sigma(\tau_{-te_4}\Theta F\cdot G)
   =\langle T_t^\circ F,G\rangle_\sigma,\\
 \langle U^\circ(v)F,U^\circ(v)G\rangle_\sigma
 &=\langle F,G\rangle_\sigma.
\end{align*}
The second identity already implies that spatial translations preserve
the null space. For time translations we need a uniform bound for
translated self-pairings.

Let $X,Y$ be two fixed positive-time tuples and let $d_0$ be the least
separation in $\theta X\cdot Y$. For $t\ge0$, shifting $Y$ forward
does not decrease any cross separation; separations within its entries
remain fixed. Each entry remains in a ball of its original radius,
though the centres move. The per-ball constants are independent of
these centres and may be chosen nonincreasing in the minimum
separation. The moment--cumulant formula therefore gives
\[
 \sup_{t\ge0}|\sigma(\theta X\cdot\tau_{te_4}Y)|<\infty.
\]
For a functional centred at a different state, the singleton blocks
are bounded by the one-point transfer; the same conclusion follows.
There are only finitely many tuple pairs in any $F\in\mathcal D_+$,
so
\begin{equation}\label{eq:bounded-time-self-pairing}
 c_F:=\sup_{t\ge0}|\langle F,T_t^\circ F\rangle_\sigma|<\infty.
\end{equation}
The global torus estimate, with its volume factor, would not imply
this bound. This is a use of the uniform per-ball source estimate.

\begin{proposition}[Reconstruction of the translation operators]
\label{prop:reflection-semigroup}
The maps $T_t^\circ$ and $U^\circ(v)$ induce a strongly continuous
semigroup $\mathcal T_t$ of positive self-adjoint contractions and a
strongly continuous unitary representation $\mathcal U(v)$ on
$\mathcal H_\sigma$. There is a unique self-adjoint $H\ge0$ with
$\mathcal T_t=e^{-tH}$, and $H\Omega=0$. The spatial generators
$P_1,P_2,P_3$ commute strongly with $H$ and with each other.
\end{proposition}

\begin{proof}
Set $a(t)=\|T_t^\circ F\|_\sigma$. Symmetry and the semigroup law
give
\[
 a(t)^2=\langle F,T_{2t}^\circ F\rangle_\sigma
       \le\|F\|_\sigma a(2t).
\]
Iterating Cauchy--Schwarz $r$ times and using
$a(2^rt)^2=\langle F,T_{2^{r+1}t}^\circ F\rangle_\sigma\le c_F$
yields
\[
 a(t)\le\|F\|_\sigma^{1-2^{-r}}c_F^{2^{-r-1}}.
\]
If $\|F\|_\sigma>0$, take $r\to\infty$. If
$\|F\|_\sigma=0$, the first Cauchy--Schwarz inequality already gives
$a(t)=0$. Thus $a(t)\le\|F\|_\sigma$ for every $F$.
Time translation preserves $\mathcal N$, descends to a contraction,
and extends uniquely to the completion. Its symmetry makes the bounded
operator self-adjoint. The semigroup law also gives
\[
 \mathcal T_t=\mathcal T_{t/2}^*\mathcal T_{t/2}\ge0.
\]

The moving-test estimate gives continuity at zero of
$\langle F,T_t^\circ F\rangle_\sigma$ for every $F\in\mathcal D_+$.
Consequently
\[
 \|\mathcal T_t[F]-[F]\|^2
 =\langle F,T_{2t}^\circ F\rangle_\sigma
     -2\langle F,T_t^\circ F\rangle_\sigma+\langle F,F\rangle_\sigma
 \longrightarrow0.
\]
Density and the contraction bound extend this convergence to all
vectors. The semigroup law then proves strong continuity at every
$t\ge0$. The spatial isometries have inverses $U^\circ(-v)$ and
thus extend to unitaries. The identical correlation-continuity
calculation gives strong continuity in $v$.

Let $E$ be the spectral measure of the positive contraction
$\mathcal T_1$, supported in $[0,1]$. Positivity and uniqueness of
the positive square root give
\[
 \mathcal T_{2^{-r}}=\int_{[0,1]}x^{2^{-r}}\,dE(x).
\]
Every such operator annihilates the range of $E(\{0\})$.
Strong continuity as $r\to\infty$ therefore implies
$E(\{0\})=0$. The self-adjoint spectral integral
\[
 H=\int_{(0,1]}-\log x\,dE(x)
\]
is nonnegative. Functional calculus identifies $e^{-tH}$ with
$\mathcal T_t$ first at positive dyadic $t$ and then at every $t$
by strong continuity. Conversely $e^{-H'}=\mathcal T_1$ for a
nonnegative self-adjoint $H'$ fixes its spectral measure under the
one-to-one map $E\mapsto e^{-E}$, so $H'=H$.
In particular every $\mathcal T_t$, $t>0$, has trivial kernel.
Since $\mathcal T_t\Omega=\Omega$, its spectral measure at
$\Omega$ is concentrated at zero energy, and $H\Omega=0$.

Spatial translations commute with $\mathcal T_1$ on the dense
quotient and hence everywhere. A bounded operator commuting with this
self-adjoint operator and its adjoint commutes with its spectral
projections. Thus $\mathcal U(v)$ commutes with every spectral
projection of $H$. The commuting spatial group has its joint spectral
measure, and the joint spectral theorem gives a single projection-valued
measure $\mathcal P$ for $(H,P_1,P_2,P_3)$, with
\[
 e^{i\tau H}\mathcal U(v)
    =\int_{[0,\infty)\times\R^3}
              e^{i(\tau E+v\cdot q)}\,d\mathcal P(E,q).
\]
This is the asserted strong commutation and spectral half-space.
\end{proof}

The countable support-controlled test families used in the extraction
are dense for the reflection norm as well. Approximate each tuple
entry with its support in a fixed separated positive-time set. The
per-ball estimate bounds the difference of every inner product, hence
the reflection norm of the difference tends to zero. Finite linear
combinations with rational real and imaginary coefficients form a
countable dense subset of $\mathcal H_\sigma$. The Hilbert space is
therefore separable.

Every $w\in W_3$ preserves the positive half-space, commutes with
$\theta$, and preserves $\sigma$. It induces a unitary fixing
$\Omega$ and commuting with $\mathcal T_t$ and $H$. The covariance
$w\tau_v=\tau_{wv}w$ gives
$\mathcal U(w)\mathcal U(v)\mathcal U(w)^*=\mathcal U(wv)$.
By uniqueness of the joint spectral measure the joint spectrum is
invariant under $(E,q)\mapsto(E,wq)$.

\subsection{Lattice representatives of the reflection form}

For $F=\sum_ic_iX_i\in\mathcal D_+$ let
\[
 \mathbb F_j=\sum_ic_i\prod_{f\in X_i}
                         [\cO_{a_j}(f)-\nu_j(\cO_{a_j}(f))].
\]
Define $\mathbb G_j$ similarly. In the reference-state volume-first
route the scalar product is
\begin{equation}\label{eq:lattice-reflection-limit}
 \langle F,G\rangle_\sigma
      =\lim_j\rho_j\left(
          \overline{\mathbb F_j\circ\theta^*}\,\mathbb G_j\right).
\end{equation}
The same formula holds along approximating tori with the bond mirror
$\theta_{\Sigma_j}$ in place of $\theta$. The difference between
these two mirrors tends to zero by the moving-test estimate above.
The reference subtraction is compatible with the reflected
representative because $\nu_j$ has the lattice reflection symmetry.

For a time $t\in\mathfrak D_L$, $t\ge0$, the translated
representative equals the lattice time translation of $\mathbb G_j$
at every sufficiently fine cutoff. Thus
$\langle F,\mathcal T_tG\rangle$ is a limit of lattice matrix
elements. The same holds for spatial translations in
$\mathfrak D_L^3$. An inequality involving finitely many such
matrix elements is a closed relation on a finite list of moments.
It passes to the continuum on this dense family, and then to all
vectors by boundedness. This is the passage used below for the
lattice energy--momentum inequality.

At this stage the spectrum is only in the half-space
$[0,\infty)\times\R^3$ and the zero-energy space can have dimension
larger than one. The spectral theorem gives the strong limit
\[
 \mathcal T_t\longrightarrow P_0
        =\mathcal P(\{0\}\times\R^3)\qquad(t\to\infty).
\]
The diagonal mirrors, rotation argument, and phase selection supply the
additional properties needed later. The construction of the semigroup
itself applies to $\R^4$; on a fixed torus an arbitrary positive time
translation need not preserve the chosen reflection slab.

\subsection{Diagonal mirrors and the Gibbs-state transfer}

At spacing $a$, the coordinate mirror families are
$\{x_\mu=c\}$, $c\in\tfrac a2\mathbb Z$. They include site and
bond mirrors. The diagonal families are
$\{x_\mu\pm x_\lambda=c\}$, $c\in a\mathbb Z$, for $\mu<\lambda$.
There are four coordinate and twelve diagonal directions. The
comparison concerns the full periodic Gibbs class, including the
states selected in an ergodic decomposition.
\subsubsection*{Symmetries of periodic Gibbs states.}
\phantomsection\label{hyp:gibbs-reflections}
The lattice argument seeks invariance under every lattice translation
and cubic isometry, and reflection positivity at the coordinate
site/bond and diagonal mirrors, for the full class
$\mathcal G^{\mathrm{per}}_{K,\beta}$. The finite-volume kernels,
temporal-layer calculation and pressure comparison below are the
steps of this argument. Using the full Gibbs class is essential when
an ergodic component is selected later.

Use tori adapted to each mirror and temporal layers to compare their
reflection forms with a periodic Gibbs state. The period lattice
matters: a cubic torus need not split along a diagonal mirror. The
pressure comparison below applies to the full Gibbs class, whereas
the original reference state was obtained from cubic tori.

\subsubsection*{Positive crossing kernels on adapted tori.}
For a chosen mirror $\Sigma$, use a period lattice rectangular in an
orthogonal frame containing its normal. A diagonal normal requires a
frame rotated through $\pi/4$ in the corresponding coordinate plane.
The periods must preserve the bond lattice and permit the two sides
to be reflected into each other. These are geometric conditions on
the finite-volume approximation, independent of the ultraviolet
estimates.

Choose the slab height to be an even integer in the
integer layer coordinate, at least two, and large enough
that the fixed test support lies in the slab.
Each side of a cubic approximating torus has an even
number of sites. The adapted periods also satisfy
the support-dependent no-short-vector condition used
in \eqref{eq:pressure-tangent}. These parity and lower
size conditions admit arbitrarily large compatible
tori at each fixed spacing and coupling.

A diagonal mirror cuts a crossing plaquette into two paths of length
two. Write $v$ for the ordered product on one path and $v'$ for the
reflected product from the other side. Apart from its constant
factor, the crossing Wilson weight is
\[
 \exp\!\left(\frac\beta2\operatorname{Re}\tr(vv'^{-1})\right).
\]
Since $v'$ is unitary,
\[
 \frac12\operatorname{Re}\tr(vv'^{-1})
 =\frac14\sum_{i,j}
       \big(v_{ij}\overline{v'_{ij}}
                   +\overline{v_{ij}}v'_{ij}\big).
\]
Each summand is a function of the variables on one side times the
complex conjugate of its reflected counterpart. Every power of the
sum has the same property, with nonnegative coefficients. The
exponential series is uniformly convergent on the compact bond space,
so the crossing weight is a uniform limit of positive reflected sums.

The plaquette weights lying wholly on one side are absorbed into the
functions on that side. The decoupled Haar functional is reflection
positive; for bond variables fixed by the mirror the two factors are
integrated against the same Haar variable. Thus the positive-kernel
criterion of \cite[Theorems~2.1--2.2, pp.~5--6]{FILS1978} applies.
It gives finite-volume reflection positivity in the adapted frame.
The same calculation applies to coordinate mirrors, with the usual
site or bond placement. It does not by itself identify which
infinite-volume Gibbs states have the resulting inequalities.

\subsubsection*{Temporal layers and trace contractions.}
Fix a coordinate direction as time and put the finite-volume measure
in temporal gauge. A layer variable consists of the spatial bonds in
one time slice. For a finite block of those bonds the coupling of two
neighbouring layers has kernel
\begin{equation}\label{eq:expanded-layer-kernel}
 k_\beta(u,u')=
 \prod_{b}\exp\!\left[-\beta\left(1-\tfrac12
                         \tr(u_bu_b'^{-1})\right)\right].
\end{equation}
It is real, symmetric, continuous and strictly positive on the
compact product of copies of $\SU(2)$. The layer space for a bond
block $\mathcal B$ is $L^2(\SU(2)^{\mathcal B},\prod_bdu_b)$.
The positivity needed for its
convolution operator is spectral, not merely pointwise. It follows
from the character expansion used in
\cite[DESY76/54, \S II.B, Proposition~1, pp.~8--9]{Luscher1977}.
For one bond, with $f$ the defining representation, write
\[
 e^{-\beta}\exp\!\left[\frac\beta4
                      (\chi_f+\chi_{\bar f})\right]
 =\sum_\rho a_\rho\chi_\rho,
 \qquad
 a_\rho=e^{-\beta}\sum_{m,n\ge0}
       \frac{(\beta/4)^{m+n}}{m!n!}N_\rho(m,n).
\]
Here $N_\rho(m,n)$ is the multiplicity of $\rho$ in
$f^{\otimes m}\otimes\bar f^{\otimes n}$. Every term is
nonnegative. Each irreducible representation of $\SU(2)$ occurs in
one of these tensor powers, so $a_\rho>0$ for $\beta>0$.
Peter--Weyl orthogonality makes convolution by this kernel a positive
scalar $a_\rho/\dim\rho$ on the matrix entries of $\rho$. Tensoring
over the bonds gives a positive operator $\mathcal C$ with no zero
eigenvalue on any finite collection of Peter--Weyl blocks. Its inverse
need not be bounded on the full layer space; only the inverse on such
finite-dimensional subspaces is used below.

After the spatial plaquette weights and the gauge projection have
been included, write $\mathsf L\ge0$ for the torus transfer operator.
For a torus with $s$ layers, the partition function is
$\tr\mathsf L^s$. If $\Xi_1,\ldots,\Xi_s$ are contractions on the
layer Hilbert space, the trace H\"older inequality gives
\begin{equation}\label{eq:expanded-layer-holder}
 \big|\tr(\Xi_1\mathsf L\cdots\Xi_s\mathsf L)\big|
 \le\prod_{i=1}^s\|\Xi_i\mathsf L\|_s
 \le\|\mathsf L\|_s^s
 =\tr\mathsf L^s.
\end{equation}
Here $\|\cdot\|_s$ is the Schatten norm, and the finite-volume
transfer representation makes the displayed traces finite. A block operator is extended by the identity on the other layer
bonds before it is used in this trace. Identity operators are allowed
among the $\Xi_i$.

The position of a local insertion matters. It is placed at the centre
of the bond kernel. If
\[
 k_{\Xi}(u,u')=
  \big(\mathcal C^{1/2}\Xi\mathcal C^{1/2}\big)(u,u')
\]
is continuous and positive, its insertion is multiplication by
$e^{A(u,u')}=k_\Xi(u,u')/k_\beta(u,u')$. The two half-kernels are
absorbed into the layer weights on their respective sides. This is
the factorisation that produces one $\mathsf L$ between consecutive
operators in \eqref{eq:expanded-layer-holder}. An arbitrary bounded
function of a pair of layers does not automatically arise from a
contraction in this way.

\subsubsection*{The lower pressure tangent.}
Let $A$ be real, continuous and local, with support in a cube of side
$ra$. Let $\nu$ be a periodic Gibbs state, and choose a subgroup of
its periods, denoted by $\mathsf P$, over which to repeat $A$. Keep
$\mathsf P$ fixed. On a torus with period lattice
$\Gamma\subset\mathsf P$ and fundamental domain $D$, let $\mathbb A$ be the resulting
finite array. Require
\[
 |v|_\infty>(r+1)a\qquad(0\ne v\in\Gamma),
\]
so the copies and their supports have the intended nonwrapping
interpretation. Write $|D|$ for the number of sites, $|\mathbb A|$
for the number of copies of $A$, and $|\partial D|$ for the number of
plaquettes meeting both sides of the boundary.

The DLR equation represents the restriction of $\nu$ to $D$ as a
mixture of finite-box measures with fixed exterior bonds. Replacing
these exterior bonds by periodic boundary data changes only the
crossing plaquettes. Since $0\le a_p\le2$, their Wilson energy costs
at most a constant times $\beta|\partial D|$. Copies of $A$ meeting
that boundary contribute at most their number times
$|t|\|A\|_\infty$. A support cube and its translates meet at most the
fixed power of $r$ appearing in the boundary count. Jensen's
inequality for the conditional measure therefore gives the lower
tangent
\begin{equation}\label{eq:pressure-tangent}
 \frac1{|D|}\log\frac{Z_D(t\mathbb A)}{Z_D(0)}
 \ge t\frac{|\mathbb A|}{|D|}\nu(A)
 -Cr^5(\beta+|t|\|A\|_\infty)
\frac{|\partial D|}{|D|}
                         \qquad(t\in\R).
\end{equation}
For a fixed local observable this is a surface error. The spacing,
coupling and support radius are held fixed while $D$ grows. Thus its
last term tends to zero along boxes of bounded aspect ratio, in either
the coordinate or the adapted frame. No estimate uniform as
$a\to0$ is needed at this step: the volume is removed first.

For an admissible centred-kernel insertion coming from a contraction,
repeat it with a fixed period cell in compatible layers and apply \eqref{eq:expanded-layer-holder}. The corresponding
normalised finite-volume partition function is at most one. Combining
this upper bound with \eqref{eq:pressure-tangent} at $t=1$ for the
actual logarithmic insertion $A$ gives, after the volume limit,
$0\ge\delta\nu(A)$, where
$\delta=\lim|\mathbb A|/|D|>0$ is the fixed copy density. Hence
\begin{equation}\label{eq:expanded-layer-insertion}
 \nu(A)\le0
\end{equation}
for each such real local logarithmic insertion. This conclusion uses
the contraction representation of the entire inserted array after
the residual gauge projection. It does not assert that
$e^{tA}$ comes from a contraction for every real $t$, or that every
unprojected contraction is an admissible insertion.

\subsubsection*{Antisymmetric functions of adjacent layers.}
The preceding inequality has a useful two-sided consequence. Let
$Y(u,u')$ be continuous, real and invariant under simultaneous gauge
transformations in the two layers, with
\[
 Y(u',u)=-Y(u,u').
\]
First suppose that $Yk_\beta$ is a polynomial in matrix entries. Its
integral operator $\mathsf K_Y$ has finite Peter--Weyl support. Kernel
antisymmetry gives $\mathsf K_Y^*=-\mathsf K_Y$ on the real layer
space. Consequently
\[
 \mathsf o=\mathcal C^{-1/2}\mathsf K_Y\mathcal C^{-1/2}
\]
is well defined on that finite-dimensional support and is a real
antisymmetric operator, extended by zero on its orthogonal
complement. The operators $\Xi_\varepsilon=e^{\varepsilon\mathsf o}$
are orthogonal contractions for every real $\varepsilon$. The
simultaneous gauge invariance of $Y$ makes both $\mathcal C$ and
$\mathsf K_Y$ intertwine the layer gauge action. The same is true
of $\mathsf o$ and $\Xi_\varepsilon$. Thus these particular
contractions are compatible with the residual gauge projection in
the transfer representation.

Their inserted kernels have the uniform expansion
\[
 \mathcal C^{1/2}\Xi_\varepsilon\mathcal C^{1/2}
 =\mathcal C+\varepsilon\mathsf K_Y+O(\varepsilon^2).
\]
The remainder is a continuous kernel on the same finite-dimensional
blocks. Since $k_\beta$ has a positive minimum on its compact domain,
the new kernel is positive for both signs of sufficiently small
$\varepsilon$. Its ratio to $k_\beta$ has a real logarithm
\[
 A_\varepsilon=\varepsilon Y+O(\varepsilon^2),
 \qquad
 \|A_\varepsilon-\varepsilon Y\|_\infty
                                      \le C_Y\varepsilon^2.
\]
Using \eqref{eq:expanded-layer-insertion} for positive and negative
$\varepsilon$ gives
\[
 \varepsilon\nu(Y)\le C_Y\varepsilon^2.
\]
Divide by $\varepsilon$ on each side of zero and let it tend to zero.
It follows that $\nu(Y)=0$.

To remove the polynomial restriction, approximate $Yk_\beta$
uniformly by matrix-entry polynomials, average under the simultaneous
gauge action, and antisymmetrise under $u\leftrightarrow u'$. These
operations preserve polynomial finite-dimensionality and do not
increase the approximation error. Division by the fixed positive
kernel gives functions tending uniformly to $Y$. Their expectations
are zero by the argument just given, and hence so is $\nu(Y)$.
Thus every admissible antisymmetric observable of two adjacent layers
has mean zero.

\subsubsection*{Determining a stack from two layers.}
The adjacent-layer identity determines the full layer law through the
local Gibbs equations. We adapt the two-layer determination argument of
\cite[Lemma~103.2]{YMCompanion},
reproducing it here in four dimensions. Fix $\beta>0$,
choose $x_4$ as time, and let $\mathcal B_3$ be the positively oriented
bonds of the spatial lattice $a\mathbb Z^3$. Write
\[
 \mathcal U=\SU(2)^{\mathcal B_3},\qquad
 \Omega_{\mathrm{lay}}=\mathcal U^{\mathbb Z}.
\]
Thus an entire spatial layer is an element of $\mathcal U$; it has
infinitely many bonds. All resampling sets below are finite.
Let $u_t$ be the spatial bonds at time $ta$, and let $w_t(x)$ be the
vertical bond from $(x,ta)$ to $(x,(t+1)a)$. For a fixed $t_0\in\mathbb Z$
define the temporal gauge by
\[
 \widehat w_{t_0}(x)=1,\qquad
 \widehat w_{t+1}(x)=\widehat w_t(x)w_t(x),\qquad
 \widetilde u_t=\widehat w_t\cdot u_t,
\]
where $(h\cdot u)(x,y)=h(x)u(x,y)h(y)^{-1}$. The recursion determines
$\widehat w_t$ in both directions. Every coordinate of the map
$\mathsf{tg}:U\mapsto\widetilde u$ is continuous and depends on
finitely many original bonds. All vertical bonds become $1$.
In particular
\[
 \widetilde u_{t_0}=u_{t_0},\qquad
 \widetilde u_{t_0+1}=w_{t_0}\cdot u_{t_0+1}.
\]
Write $U[\widetilde u]$ for the bond field with these spatial layers
and all vertical bonds equal to $1$. A gauge-invariant continuous
local $F$ satisfies
\[
 F(U)=\widetilde F(\mathsf{tg}(U)),\qquad
 \widetilde F(\widetilde u)=F(U[\widetilde u]).
\]

For finite $S\subset\mathcal B_3$, denote the kernel
\eqref{eq:expanded-layer-kernel} on this bond set by $k_S$. Let
$V_S(u)$ be the sum of the spatial plaquette actions meeting $S$.
Each spatial bond belongs to four spatial plaquettes, and every
plaquette action belongs to $[0,2]$. Hence
\[
 0\le V_S\le8|S|,\qquad e^{-2\beta|S|}\le k_S\le1.
\]
If $v\in\SU(2)^S$, let $\widetilde u^{(S,t;v)}$ replace the bonds
$S$ in layer $t$ by $v$. The local layer resampling kernel is
\begin{equation}\label{eq:temporal-gauge-chain-kernel}
 \begin{split}
 q_{S,t}(v;\widetilde u)
 &=k_S(\widetilde u_{t-1}|_S,v)
   e^{-\beta V_S(v\vee\widetilde u_t|_{S^c})}
   k_S(v,\widetilde u_{t+1}|_S),\\
 (\widetilde\gamma_{S,t}\psi)(\widetilde u)
 &=\frac{\int\psi(\widetilde u^{(S,t;v)})
                   q_{S,t}(v;\widetilde u)\,dv}
          {\int q_{S,t}(v;\widetilde u)\,dv}.
 \end{split}
\end{equation}
Here $dv$ is probability Haar measure. Both $q_{S,t}$ and its
normalising integral lie between $e^{-12\beta|S|}$ and $1$.
Consequently $\widetilde\gamma_{S,t}$ maps continuous local functions
to continuous local functions.

The image $\rho=\mathsf{tg}_*\nu$ of any Wilson Gibbs state satisfies
\begin{equation}\label{eq:temporal-gauge-chain-dlr}
 \rho(\widetilde\gamma_{S,t}\psi)=\rho(\psi)
 \quad(S\subset\mathcal B_3\text{ finite},\ t\in\mathbb Z,
                      \ \psi\text{ continuous and local}).
\end{equation}
Indeed, resample only the original spatial bonds corresponding to
$(S,t)$. This leaves every $\widehat w_s$ unchanged. The substitution
$v=\widehat w_t\cdot U_{S,t}$ preserves Haar measure. The plaquettes
meeting these bonds are precisely the spatial plaquettes meeting
$S$ and the two vertical plaquettes at each bond of $S$. Their
weight is \eqref{eq:temporal-gauge-chain-kernel}, since a vertical
plaquette becomes
$1-\frac12\tr(\widetilde u_t(b)\widetilde u_{t+1}(b)^{-1})$.
The original DLR equation therefore gives
\eqref{eq:temporal-gauge-chain-dlr}. This is a local resampling
identity; no conditional independence of entire infinite layers is
being asserted.

The equations are covariant under every index reflection
\[
 (R_\ell\widetilde u)_t=\widetilde u_{\ell-t}.
\]
The spatial potential is unchanged and the symmetry of $k_S$
exchanges the two neighbouring factors in
\eqref{eq:temporal-gauge-chain-kernel}. Thus $R_{\ell*}\rho$ obeys
the same equations whenever $\rho$ does.

\begin{lemma}\label{lem:layer-law-determination}
Let $\beta>0$. Two probability laws on $\Omega_{\mathrm{lay}}$
satisfying \eqref{eq:temporal-gauge-chain-dlr} and agreeing on two
or more consecutive entire layers agree on all layers.
\end{lemma}

\begin{proof}
Suppose that the two laws agree on layers $m,\ldots,n$, with $m<n$.
We first extend agreement to layer $n+1$. Fix a nonempty finite
$S\subset\mathcal B_3$ and use the coordinates
\[
 u'=\widetilde u_{n+1}|_S,\qquad v=\widetilde u_n|_S,\qquad
 c=(\widetilde u_n|_{S^c},\widetilde u_{n-1},\ldots,
                                      \widetilde u_m)\in Y.
\]
The space $Y$ is a compact countable product of copies of $\SU(2)$.
Put
\[
 \Phi(v;c)=k_S(\widetilde u_{n-1}|_S,v)
       e^{-\beta V_S(v\vee\widetilde u_n|_{S^c})},\qquad
 N(u',c)=\int k_S(v,u')\Phi(v;c)\,dv.
\]
The condition $m<n$ ensures that $\widetilde u_{n-1}$ is among the
coordinates of $c$. These functions are continuous and local, and
\[
 \Phi\ge e^{-10\beta|S|},\qquad N\ge e^{-12\beta|S|}.
\]
Let $\rho_i'$ be the image of law $i$ on $(u',c)$ and define the
finite positive measure $\sigma_i=N^{-1}\rho_i'$. Its mass is at
most $e^{12\beta|S|}$. Equation
\eqref{eq:temporal-gauge-chain-dlr} at $(S,n)$ gives
\begin{equation}\label{eq:layer-law-resampling}
 \rho_i(\psi)=\int\sigma_i(du',dc)\int
       k_S(v,u')\Phi(v;c)\psi(u',v,c)\,dv
\end{equation}
for continuous local functions of the displayed coordinates.
Take
\[
 \psi(u',v,c)=\frac{\varphi(v)\chi(c)}{\Phi(v;c)},
\]
where $\varphi$ is a finite Peter--Weyl sum on $\SU(2)^S$ and
$\chi$ is continuous and local on $Y$. This is a bounded continuous
local function of the already equal layers $m,\ldots,n$.
Subtracting the two resampling identities gives
\[
 \int (\mathcal C_S\varphi)(u')\chi(c)
                              \,d(\sigma_1-\sigma_2)(u',c)=0,
\]
where $\mathcal C_S$ is convolution by $k_S$.

Every Peter--Weyl block of $\mathcal C_S$ has a strictly positive
scalar eigenvalue, as proved above. Thus $\mathcal C_S$ maps the
finite Peter--Weyl sums onto themselves: divide each of the finitely
many coefficients by its eigenvalue. These sums are uniformly dense
in $C(\SU(2)^S)$, and the continuous local functions are uniformly
dense in $C(Y)$. Finite sums of their products are uniformly dense
in $C(\SU(2)^S\times Y)$. Since $\sigma_1-\sigma_2$ has finite
total variation, the last displayed identity implies
$\sigma_1=\sigma_2$.
There is no limit of inverse operators in this argument. The set
$S$ and each Peter--Weyl test are fixed when their finitely many
eigenvalues are divided out; density is applied to the resulting
finite measures.

Equation \eqref{eq:layer-law-resampling} now gives equality on all
continuous local functions of $(u',v,c)$. Given a local test on
layers $m,\ldots,n+1$, choose $S$ to contain its finitely many
coordinates on layer $n+1$. This proves equality on the entire
block $m,\ldots,n+1$. Reflecting the indices about $m+n$ and using
covariance of the chain equations proves extension to layer $m-1$.
Iteration in both directions gives agreement on every finite block.
Continuous local functions determine a measure on the compact
product $\Omega_{\mathrm{lay}}$, proving the result.
\end{proof}

\begin{proposition}\label{prop:periodic-bond-reflections}
Every periodic Wilson Gibbs state on $a\mathbb Z^4$ is invariant
under every coordinate bond reflection. It is consequently invariant
under translations by $2a\mathbb Z^4$.
\end{proposition}

\begin{proof}
We first explain the residual gauge action, so that the adjacent-layer
identity applies to a full two-layer marginal. The DLR equations
give local gauge invariance of every Gibbs state. At one site,
resampling its eight incident bonds and making the corresponding
left and right Haar substitutions leaves every plaquette trace
unchanged. Thus the DLR kernel gives the same expectation before and
after that gauge transformation. Composition proves invariance for
finitely supported gauge transformations, which suffices for any
local observable.

Under an original gauge transformation $g_t(x)$, the temporal gauge
satisfies
\[
 \widehat w_t^{\,g}=g_{t_0}\widehat w_tg_t^{-1},\qquad
 \widetilde u_t^{\,g}=g_{t_0}\cdot\widetilde u_t.
\]
Hence $\rho=\mathsf{tg}_*\nu$ is invariant under the simultaneous
spatial gauge action on all layers. For a continuous local
$H(\widetilde u_{t_0},\widetilde u_{t_0+1})$, average over the gauge
variables at the endpoints of its spatial bonds. Call the result
$H^{\mathrm{inv}}$. The average preserves its expectation under
$\rho$, remains local, and commutes with exchanging the two layers.
The adjacent-layer argument, applied to the real and imaginary
parts of
\[
 H^{\mathrm{inv}}(u,u')-H^{\mathrm{inv}}(u',u),
\]
therefore proves that the entire two-layer marginal of $\rho$ is
invariant under exchange. This use of the earlier insertion
calculation requires only its gauge-compatible contractions.

Let $R=R_{2t_0+1}$. Both $\rho$ and $R_*\rho$ obey the local chain
equations. On layers $t_0,t_0+1$ their marginals are exchanged, so
they agree by the preceding paragraph. Lemma
\ref{lem:layer-law-determination} gives $R_*\rho=\rho$. In
particular this equality holds for every observable on any fixed
finite stack, with no restriction on its height.

For the original bond reflection
$\theta=\theta_{e_4,(t_0+1/2)a}$ one has
\[
 \theta U[\widetilde u]=U[R\widetilde u].
\]
The factorisation through temporal gauge therefore gives
$\nu(\theta F)=\nu(F)$ for every gauge-invariant continuous local
$F$. The local gauge average below
extends this equality to all continuous local functions. These
functions determine the bond-field measure, so $\theta_*\nu=\nu$.

The same argument applies with any coordinate chosen as time.
Equivalently, apply it to the coordinate permutation of $\nu$;
covariance of the Wilson specification makes this another periodic
Gibbs state, without assuming that $\nu$ is permutation invariant.
The two bond reflections at $x_\mu=a/2$ and $x_\mu=3a/2$ compose to
translation by $2ae_\mu$. This proves the asserted periods.

If $\beta=0$, the DLR equation on the finite support of a local
function sets its expectation equal to its product Haar integral.
It determines the unique product Haar Gibbs state, for which all
these invariances hold as well.
\end{proof}

All steps were at fixed spacing and coupling. Periodicity entered
the adjacent-layer pressure calculation; the local chain argument
requires neither a torus limit representation of the state nor
ergodicity. It therefore applies to every member of
$\mathcal G^{\mathrm{per}}_{K,\beta}$, including its periodic ergodic
components. The diagonal mirror inequalities are supplied by the
adapted-torus pressure comparison below, after the common periods
$2a\mathbb Z^4$ have been established.

For the reflection inequalities below, continuous local tests also
suffice for bounded Borel local tests. Fix a mirror reflection
$\theta$ and a finite bond set supporting $F$ on one closed side.
Approximate $F$ by continuous functions $F_j$ of those same bonds in
$L^2$ for the sum of the corresponding marginal laws of $\nu$ and
$\theta_*\nu$. Then both $F_j\to F$ and
$\theta F_j\to\theta F$ in $L^2(\nu)$. Cauchy--Schwarz gives
\[
 \begin{split}
 \big|\nu(\overline{\theta F_j}F_j)
             -\nu(\overline{\theta F}F)\big|
 &\le\|\theta(F_j-F)\|_2\|F_j\|_2
       +\|\theta F\|_2\|F_j-F\|_2
 \longrightarrow0.
 \end{split}
\]
Thus positivity passes to $F$. In particular the weakly closed
continuous-test inequalities used in the pressure comparison define
the full local reflection-positive class.

\subsubsection*{From bond invariance to all mirror inequalities.}
To pass from bond reflection invariance on gauge-invariant local
polynomials to all local functions, use local gauge invariance from
the DLR equations.
For a local $F$, let
\[
 F^{\mathrm{inv}}(U)=\int F(U^u)\,du,
\]
where the Haar average is over gauge transformations at the endpoints
of bonds in the support of $F$. This is still a local function,
$\nu(F^{\mathrm{inv}})=\nu(F)$, and the average commutes with
reflection after the endpoint set is reflected. Equality on
gauge-invariant functions therefore gives equality on all local
functions. Compositions of neighbouring bond reflections then give
translations by $2a\mathbb Z^4$. At this point those translations are
periods of the state, whatever its originally specified periods were.

Now fix an equidistant mirror family. Repeat a real local $A$ so that
each copy lies in a slab bounded by adjacent mirrors and the period
group takes each slab to the next. On a compatible adapted torus the
chessboard estimate folds $\mathbb A$ into a reflection-symmetric
array $\mathbb A^{\mathrm f}$ and gives
\[
 Z_D(t\mathbb A)\le Z_D(t\mathbb A^{\mathrm f}).
\]
Write $\delta_D=|\mathbb A|/|D|$, and use the same number of
copies in the folded array. For $t>0$, convexity gives the explicit
upper tangent
\[
 \frac1{|D|}\log\frac{Z_D(t\mathbb A^{\mathrm f})}{Z_D(0)}
 \le \frac{t}{|D|}\frac{d}{dt}\log Z_D(t\mathbb A^{\mathrm f})
 =t\delta_D\,\nu_{D,t}^{\mathrm f}(A).
\]
Here $\nu_{D,t}^{\mathrm f}$ is the finite-volume measure with the
folded perturbation. The last equality reads $A$ in its chosen
representative slab: the other copies have the same expectation by
the period and reflection symmetries of the folded array. Combining
this bound, the folding inequality and \eqref{eq:pressure-tangent}
yields
\[
 \nu(A)\le\nu_{D,t}^{\mathrm f}(A)
 +\frac{Cr^5(\beta+t\|A\|_\infty)}{t\delta_D}
                          \frac{|\partial D|}{|D|}.
\]
The repetition group is fixed, so $\delta_D$ tends to a strictly
positive density. Let the volume tend to infinity at fixed $t$.
A subsequential limit $\nu_t^{\mathrm f}$ retains the folded-array
reflection inequalities and satisfies the Gibbs specification with
that fixed periodic perturbation. Next let $t\downarrow0$. On every
fixed finite region the added bounded local interaction tends to
zero, so the conditional kernels tend to the unperturbed Wilson
kernels. Weak limit points are therefore Wilson Gibbs states.
Reflection positivity and mirror invariance are closed, giving a
comparison state in the class
\[
 \mathcal C_\Sigma=
 \{\rho\in\mathcal G^{\mathrm{per}}_{K,\beta}:
       \rho\circ\theta_\Sigma=\rho,
       \ \rho\hbox{ is reflection positive at }\Sigma\}.
\]
The compatible slabs and periods are fixed during this argument.
The surface error is removed before its factor $1/t$ can cause a
problem. Weak limits preserve the finite-region DLR identities and
the fixed periods of the folded array. The temporal reflection
argument then gives their $2a\mathbb Z^4$ periods as well. The
conclusion is
\begin{equation}\label{eq:expanded-mirror-comparison}
 \nu(A)\le\sup_{\rho\in\mathcal C_\Sigma}\rho(A).
\end{equation}
The comparison state may depend on $A$. The compatible slabs and
periods are chosen before the upper-tangent calculation, and remain
fixed while the volume is removed.

The slab restriction must also be removed. Translate a given local
$A$, possibly straddling $\Sigma$, by a vector in $2a\mathbb Z^4$
into one side and take the slabs wide enough to contain that
translate. Bond invariance makes this translation a period of both
the original and the comparison states. Thus the comparison also
holds for the untranslated $A$. Reflection positivity is a family of
closed linear inequalities in the state; mirror invariance and the
DLR equations are closed as well. By the bond-reflection argument,
every periodic Gibbs state is already invariant under
$2a\mathbb Z^4$. Conversely such a Gibbs state is periodic. The
class can therefore be written using this single fixed period group,
whose invariance is closed under weak limits on the compact bond
configuration space. The finite-range DLR identities are also closed,
so this Gibbs class is compact and convex. Its closed mirror-invariant,
reflection-positive subclass $\mathcal C_\Sigma$ is consequently
compact and convex.
If $\nu$ lay outside it, separation would give a real local
continuous $A$ with
$\nu(A)>\sup_{\rho\in\mathcal C_\Sigma}\rho(A)$, contrary to
\eqref{eq:expanded-mirror-comparison}.

This pressure and separation argument treats local tests crossing
the mirror as well as tests on one side. Together with
\cref{prop:periodic-bond-reflections}, it gives the lattice symmetries
and mirror inequalities used in the source comparison and phase
selection.

Finally, site and bond mirrors whose positions differ by $a/2$
generate translation by $a$ in a coordinate direction. Coordinate
and diagonal mirrors generate the cubic isometries. Applying the
comparison to these mirror families gives the lattice symmetries
claimed at the start of this argument.


\subsection{The lattice spectral inequality}

The coordinate bond mirror
$\{x_4=-a/2\}$ gives a positive contraction $T_a$, and spatial
translations give commuting unitaries $U_j=e^{iaP_j}$.
The site-mirror space at $\{x_4=0\}$ has a
positive contraction $T_a^{\mathrm s}$. The relation between the two
reflection forms is
\[
 \langle[F],[G]\rangle
 =\langle[F]^{\mathrm s},T_a^{\mathrm s}[G]^{\mathrm s}\rangle.
\]
Thus $[F]\mapsto(T_a^{\mathrm s})^{1/2}[F]^{\mathrm s}$ identifies the
bond space with the closure of the range of $(T_a^{\mathrm s})^{1/2}$.
This closed subspace is $(\ker T_a^{\mathrm s})^\perp$, on which the restriction
of $T_a^{\mathrm s}$ is injective. Consequently $T_a$ is injective and
$T_a=e^{-aH}$ defines a nonnegative self-adjoint operator $H$.
The site space may have a zero-transfer subspace; none of the arguments
below identifies that subspace with finite energy.

\begin{proposition}[The lattice cone]\label{prop:lattice-cone}
The coordinate and diagonal reflection operators satisfy
\begin{equation}\label{eq:lattice-cone}
 (1-T_a)^2\ge T_a(2-U_j-U_j^*).
\end{equation}
Equivalently, their joint spectral variables satisfy
\[
 \cosh(aE)+\cos(aq_j)\ge2.
\]
\end{proposition}

\begin{proof}
The argument compares the coordinate and diagonal reflection pictures
of the same correlation. Fix $j$ and write $U=U_j$. The displacement
$ka(e_4+e_j)+la(e_4-e_j)$ is represented in the coordinate picture by
$T_a^{k+l}U^{k-l}$. In the picture of a diagonal mirror, the first
diagonal step is a positive contraction $T_+$ and the other is a
unitary $V_-$. The two commute because they represent translations.

Choose an integer $R\ge1$ and let $A$ be local in
$\{0\le x_4\le Ra,\ |x_j|\le Ra\}$, and first translate its
reflection vector by $Ra$ in time.
Its joint spectral measure for $(T_a,U)$ is a finite positive measure
on $[0,1]\times S^1$. Let $\varsigma$ be its image under
$(t,\zeta)\mapsto t\zeta$. Expanding a polynomial $Q$ in monomials
expresses $\int|Q|^2\,d\varsigma$ as a finite sum of correlations of
the displacements just described. For the diagonal mirror
$\Sigma_d=\{x_4+x_j=Ra\}$, write $\theta_d$ for its reflection and
$\theta_l$ for reflection in $\{x_4=-a/2\}$. The outer vectors are
\[
 \psi''=[\tau_{2Rae_4}A]_d,\qquad
 \psi'=[\theta_d\theta_l A]_d.
\]
Both representatives lie on the same allowed side of $\Sigma_d$.
Expanding the polynomial gives the identity
\[
 \int|Q|^2\,d\varsigma
 =\langle\psi',Q(T_+)\overline Q(V_-)\psi''\rangle,
\]
where $\overline Q$ denotes coefficientwise conjugation of $Q$.
The resulting bound is
\begin{equation}\label{eq:slit-disc-bound}
 \int_{\overline{\mathbb D}}|Q(z)|^2\,d\varsigma(z)
 \le C_A\sup_{[0,1]}|Q|\sup_{S^1}|Q|.
\end{equation}
Indeed the diagonal expression is a pairing with
$Q(T_+)\overline Q(V_-)$ between those two vectors; the spectral
theorem bounds its operator norm by the product of the two suprema.
The number $C_A$ is the product of the outer-vector norms and is
independent of $Q$.

It remains to extract the support information in this inequality.
For $s>0$ put
\[
 f_s(z)=\frac{\cos\!\left(s\operatorname{Log}
             \frac{1+\sqrt z}{1-\sqrt z}\right)}{\cosh(\pi s/2)},
 \qquad
 \phi(z)=\left|\operatorname{Im}\operatorname{Log}
                    \frac{1+\sqrt z}{1-\sqrt z}\right|.
\]
Changing the branch of $\sqrt z$ changes the logarithm's sign and
does not change its cosine. Its power series at zero contains only
integer powers of $z$, so $f_s$ is holomorphic in the disc. The
fraction inside the logarithm maps the right half of the square-root
disc to the right half-plane. In particular $0\le\phi<\pi/2$.
The elementary formulas for the cosine of a complex number give
\[
 |f_s(z)|\le1,\qquad
 \sup_{0\le z<1}|f_s(z)|\le\frac1{\cosh(\pi s/2)},\qquad
 |f_s(z)|\ge\frac{\sinh(s\phi(z))}{\cosh(\pi s/2)}.
\]
For $0<\epsilon<1$ the Taylor polynomials of
$f_s((1-\epsilon)z)$ converge uniformly on the closed unit disc.
Since $\varsigma$ is finite, they may be inserted in
\eqref{eq:slit-disc-bound} and then replaced by their uniform limit.
It follows that
\[
 \int|f_s((1-\epsilon)z)|^2\,d\varsigma(z)
 \le\frac{C_A}{\cosh(\pi s/2)}.
\]
On a compact subset where $\phi((1-\epsilon)z)>\pi/4+\delta$,
the lower bound for the integrand has order
$e^{(-\pi/2+2\delta)s}$, whereas the right side has order
$e^{-\pi s/2}$. Letting $s\to\infty$ shows that the measure of
that compact subset is zero. Then let $\epsilon\downarrow0$ and
exhaust the part of the closed disc where $\phi>\pi/4$ by compact subsets.
Hence
\[
 \operatorname{supp}\varsigma\subset
     \{z\in\overline{\mathbb D}:|1+z|^2\ge4|z|\}.
\]
Here the boundary curve follows directly from the logarithm formula:
$\phi=\pi/4$ is the inverse image of the rays of arguments
$\pm\pi/4$; eliminating $\sqrt z$ gives $|1+z|^2=4|z|$.

For $z=t\zeta$ one has
\[
 |1+t\zeta|^2-4t=(1-t)^2-t|1-\zeta|^2.
\]
Thus every time-translated local vector used above has spectral
measure supported where $(1-t)^2\ge t|1-\zeta|^2$.
For each local $A$, choose $R$ large enough for its support. The
spectral projection of the complementary open set annihilates
$T_a^R[A]$. It commutes with $T_a^R$, which is injective, so it
annihilates $[A]$ as well. Local vectors are dense; hence this
projection vanishes. The spectral theorem gives \eqref{eq:lattice-cone}.
\end{proof}

At a fixed spacing this implies a linear cone with a constant smaller
than one. Taking $aq_j$ in $[-\pi,\pi]$, the curved lower boundary is
\[
 aE\ge\operatorname{arccosh}(2-\cos(aq_j)),\qquad
 E\ge k_*|q_j|,\qquad
 k_*=\frac2\pi\log(1+\sqrt2).
\]
The positive constant $k_*$ is sufficient in the spectral estimate for
clustering. Complex rotation requires the constant-one inequality in
the continuum, for which one retains the curved bound until after
the cutoff is removed.

\subsection{The constant-one cone in the continuum}

\begin{proposition}[The continuum coordinate cones]\label{prop:continuum-cone}
Take a common limit of the separated moments, also allowing evaluation in
states $\rho_j$ with subtraction at reference states $\nu_j$.
The reflection-space semigroup and spatial translations satisfy
\begin{equation}\label{eq:cone}
 H\ge|P_j|\qquad(j=1,2,3).
\end{equation}
\end{proposition}

\begin{proof}
The reflection-space construction is the positive-time translation
construction of \cite[\S4.1, pp.~90--92]{OS1973}, applied to the
moments on separated tests. We use its bounded semigroup operators,
rather than asserting convergence of the unbounded lattice generators.

By \cref{prop:lattice-cone},
$\sinh(aE/2)\ge|\sin(aq_j/2)|$ on the lattice spectrum.
For an integer $k\ge1$, the inequalities
$\sinh(ku)\ge k\sinh u$ for $u\ge0$ and
$|\sin(kv)|\le k|\sin v|$ give the same assertion with $a$ replaced
by $ka$. Put $t=ka$. The polynomial operator inequality also holds
in the site picture: the bond/site identification proves it on
$(\ker T_a^{\mathrm s})^\perp$, and on $\ker T_a^{\mathrm s}$ it
reduces to $1\ge0$. In that site picture, evaluated at a local
reflection vector represented by $F$, it is the finite inequality
\begin{equation}\label{eq:five-point}
 u(0)-4u(te_4)+u(2te_4)
       +u(te_4+te_j)+u(te_4-te_j)\ge0,
\end{equation}
where $u(v)=\rho(\overline{\theta F}\,\tau_vF)$.
All the correlations here have separated reflected supports.

Fix a positive $L$-adic physical displacement $t$. At sufficiently fine
cutoffs it is an integer multiple of the lattice spacing. The lattice
representatives of $F$ are exactly covariant under these translations,
because the density is evaluated at plaquette barycentres. Moment
convergence therefore passes every term of \eqref{eq:five-point} to
the continuum. This also holds for the common reference-state
subtraction: the same finite reflection inequality holds in each
evaluating state $\rho_j$. Strong continuity extends the resulting
bounded-operator inequality from $L$-adic $t$ to every $t>0$.

Let $\mathcal P$ be the continuum joint spectral measure of $(H,P_j)$.
The inequality says that
\[
 (1-e^{-tE})^2
       -e^{-tE}(2-e^{itq}-e^{-itq})\ge0
\]
for $\mathcal P$-almost every $(E,q)$, for every positive $L$-adic
$t$. Use the countability of these $t$ to choose one common full-measure
spectral set. At each of its points divide by $4e^{-tE}(t/2)^2$ and
let $t\downarrow0$. This gives $E^2\ge q^2$, and $E\ge0$ gives
$E\ge|q|$. This proves \eqref{eq:cone} without any exchange of an
unbounded generator with the cutoff limit.
\end{proof}

\section{From cubic symmetry to rotations}\label{sec:rotations}

The lattice group supplies quarter-turns. To obtain the target theorem
we must recover every angle. Work first in the plane with coordinates $(x_4,x_1)$.
The argument combines three estimates: the
constant-one continuum cone, an insertion estimate independent of
the number of other fields, and a short-distance two-point bound
that improves a pure power by a factor tending to zero. Diagonal
positivity supplies one further sign at the end of the argument.

For a component with mean $b\int f$, apply this argument to the
centred functional. Restoring the invariant mean by
\eqref{eq:restore-mean} preserves all its rotation identities.
In this section write $\sigma$ for a centred limiting moment
functional and $\|F\|_{(c)}=\|[\tau_{-ce_4}F]\|$ whenever the
translated polynomial remains supported at positive times.

\subsection{An insertion estimate uniform in the other fields}

\phantomsection\label{hyp:insertion-bound}
The first step is to bound a local insertion between arbitrary
reflection vectors. Let $0<\eta\le1$. Place $h$ in a ball of radius $c\eta$ near the
time-zero plane, with $|x_4|\le\eta/8$, and let $F,G$ be finite
polynomials in separated tests supported in $x_4>\eta/4$. The
required estimate is
\begin{equation}\label{eq:rotation-insertion}
 |\sigma(\Theta F\cdot(h)\cdot G)|
 \le C(g)\eta^{-p_*}\|h\|_\sharp
                    \|F\|_{(\eta/4)}\|G\|_{(\eta/4)}.
\end{equation}
The constants must be independent of the degrees of $F$ and $G$.

The proof combines source normalization with repeated reflections. First
subtract the scalar source energies separately in each cell of side
$l\asymp\eta$. Locality then makes the interaction between cells cost
only a power of $l^{-1}$. Next place the normalized insertion in a long
reflected column. Taking a root removes the two endpoint observables,
including their potentially large supremum norms. The second operation
is what makes the answer independent of the number of outer fields.
The quantitative work is to carry the one-cell cancellation through
all scales $a_K\ll l\le1$ with the same source induction; a bound for
unit cells alone would not suffice.

We carry out the argument for real $h$. Splitting a complex test into
its real and imaginary parts gives the same conclusion by linearity,
with a harmless fixed enlargement of the constant.

\subsubsection{Source energies on a subunit cell}

Take a chessboard with side $l=L^{-i}\asymp\eta$ and
$a_K\le l/L^2$, whose faces are bond mirrors. Suppose first that
$h$ is supported in the middle of one cell, with a fixed positive
fraction of $l$ separating its support from every face. Let $h^+$
and $h^-$ be its reflected arrays on even- and odd-parity cells,
respectively. The repeated complex source is
\[
 W_z=zh^++\bar z h^-.
\]
Copies belonging to either array have disjoint supports. Their
supremum and Lipschitz bounds are those of $h$, including at a
cell face because the source vanishes in a neighbourhood of that
face. Hence
\[
 \|W_z\|_\flat\le2|z|\|h\|_\sharp.
\]
This bound is independent of $N=(\ell/l)^4$, the number of
cells. We use two array functions on the source domain, rather than
$N$ source variables with their norms added together.

Apply upper stability and the unsourced floor on $T_\ell$
itself. The vacuum factor is the same in numerator and denominator,
so it cancels before any estimate of vacuum densities. The quotient
and its chessboard root give
\[
 \|e^{z\cO_a(h)}\|_{\mathrm{ch},\ell}
  =\left[\frac{Z_\ell(W_z)}{Z_\ell(0)}\right]^{1/N}
  \le\exp\left[
       \frac1N\operatorname{Re}\mathcal E_\ell(W_z)+C_gl^4
           \right],
 \qquad |z|\|h\|_\sharp<c\rho.
\]
The repeated numerator is nonnegative by reflection positivity.
This application of stability requires neither a reference torus
nor the finite-torus transfer estimate.

For an enlarged domain $X^{\sim1}$ let its copy set consist of
the reflected supports of $h$ that it meets. A term with an empty
copy set is zero. Split the scalar energies into those whose copy
set has one element and those whose copy set has at least two
elements. For a term of the first kind, locality leaves a
holomorphic function of $z$ if the copy is even, and of $\bar z$
if it is odd. Conjugate the latter term, sum the resulting
holomorphic functions and divide by $N$. Denote the result by
$E_{h;\ell}(z)$. It is zero at zero and satisfies the exact
cancellation
\[
 \frac1N\operatorname{Re}\mathcal E_\ell(W_z)
                     -\operatorname{Re}E_{h;\ell}(z)
   =\frac1N\operatorname{Re}
       \sum_j\sum_{X:\ X^{\sim1}\text{ meets at least two copies}}
                                                    e_j(X;W_z).
\]
One-copy domains can cross a face, provided they reach no second
support; all remain in $E_{h;\ell}$. Their cutoff-dependent scalar
contributions cancel as in the construction of
\eqref{eq:source-common-normalizer}.

We now sum the domains in the remaining class. Put
$t_j=L^{K-j}/M$, the inverse physical side of a localization cube.
At a fine scale, $lt_j\ge1$, a domain meeting two copies has
$d_j(X)\ge clt_j-C$. Charge it to one of the copies it meets
and anchor it in one of the at most $(Clt_j+C)^4$ cubes near
that support. The tree sum gives, after division by $N$,
\[
 C\sum_{lt_j\ge1}(Clt_j+C)^4e^{-c[clt_j-C]_+}\le C.
\]
This is a geometric-scale sum. Once $lt_j$ exceeds a fixed
constant its exponential decay dominates the fourth-power count;
the remaining terms form a bounded finite segment of that sequence.

At a coarse scale, $lt_j<1$, count all localization cubes on the
torus before dividing by $N$. The local energy bound and the tree
bound give a cost $C\ell^4t_j^4$ for all domains of that level.
Thus their cost per cell is $C(lt_j)^4$, and
\[
 \frac1N\sum_{lt_j<1}C\ell^4t_j^4
                 =C\sum_{lt_j<1}(lt_j)^4\le C.
\]
This global count is stronger than anchoring separately at every
copy at a scale whose cubes contain many cells. In the latter
count the additive constant in $(Clt_j+C)^4$ would cost one
constant at each remaining level. Counting the cubes first is the
reason that the subunit-cell cost is summable even as $l$ tends
to zero.

Combining the two scale ranges with the stability term yields
\begin{equation}\label{eq:rotation-cell-stability}
 \left\|e^{z\cO_a(h)-E_{h;\ell}(z)}\right\|_{\mathrm{ch},\ell}
       \le Q(l),\qquad Q(l)\le C(g)l^{-C},
 \qquad |z|\|h\|_\sharp<c\rho.
\end{equation}
The displayed polynomial bound is sufficient for the insertion
argument; the two scalar-energy sums above are bounded even before
allowing this polynomial loss. Their use requires the local
source-energy and stability estimates uniformly for the two reflected
array functions. Establishing that shared array induction is the
source construction needed here. The cancellation and counts then
give the cell estimate at the actual side $l$.

The difference from the graded estimate for many active sources is
useful. There the coarse face domains are charged separately to
active pieces of radius $\eta$, which gives a logarithmic number
of bounded contributions per piece. Here the same one-copy pattern
fills every cell, and the global cube count divided by $N$ gains
the factor $(lt_j)^4$ at the coarse scales. The normalizers for
the two arguments serve different purposes and need not agree.

The geometric part of the chessboard argument is unchanged at
$l<1$. Dilation by $L^i$ sends this board to a unit board and
the lattice to the one of cutoff $K-i\ge2$. It preserves the
seminorm, reflection inequalities and fixed-cutoff Gibbs domination.
It is not used to transfer an RG stability estimate from another
datum: the stability and energy calculation just made was performed
directly on the original torus.

At a fixed $K$, the functions $E_{h;\ell}$ are bounded on compact
subdiscs, uniformly in $\ell$. To see the volume independence,
sum the local energy bounds through each localization cube and
divide their total $C\ell^4t_j^4$ by $N$. This bounds the
contribution of level $j$ by $Cl^4t_j^4$. The sum over the
finitely many levels is finite at this $K$, though it can increase
with $K$. The same bound holds on a slightly larger source disc
than a given compact subdisc. Montel's theorem therefore gives a
sequence $\ell_r\to\infty$ and a holomorphic limit $E_h$, with
locally uniform convergence and $E_h(0)=0$. The sequence and limit
may depend on $K,\beta,h$ and $l$. They are chosen from finite
tori before a Gibbs state or an outer reflection vector is selected.

For $D_z=e^{z\cO_a(h)-E_h(z)}$, homogeneity gives the exact
relation
\[
 \|D_z\|_{\mathrm{ch},\ell_r}
  =|e^{E_{h;\ell_r}(z)-E_h(z)}|
       \left\|e^{z\cO_a(h)-E_{h;\ell_r}(z)}
                                      \right\|_{\mathrm{ch},\ell_r}.
\]
The first factor tends to one locally uniformly. Hence
\begin{equation}\label{eq:insertion-selected-normalizer}
 q_{D_z}:=\limsup_r\|D_z\|_{\mathrm{ch},\ell_r}\le Q(l).
\end{equation}
This is the selected sequence throughout the column argument. No
comparison of $E_{h;\ell}$ with $E_h$ is needed outside it.
Gibbs product domination in \cref{sec:sources} holds along every
chosen volume sequence, and so applies along this one for every
periodic Gibbs state.

\subsubsection{The column estimate at a fixed cutoff}

We now bound insertion of $D_z$ between arbitrary outer reflection
vectors by its chessboard seminorm. The outer factors enter through
endpoint supremum norms at this fixed cutoff.

Fix a periodic Gibbs state $\nu'$ at this cutoff. For functions above
a mirror at time $t$, use the positive form
\[
 \langle B,A\rangle_{\nu',t}
                =\nu'(\overline{B\circ\theta_t^*}A).
\]
Let $D$ be supported in a cell immediately above this mirror, between
times $t$ and $t+l$. On the space of bounded local functions above
the mirror put
\[
 \widehat D A=D\,\tau_{le_4}A,
 \qquad D^\dagger=\overline{D\circ\theta_{t+l/2}^*}.
\]
Reflection followed by translation shows that the formal adjoint of
$\widehat D$ is $\widehat{D^\dagger}$. In fact reflection of
$D^\dagger$ in the lower face is the conjugate of $D$ translated
by $-l$, and the same translation carries the reflected upper
observable to the required side. Translation invariance of $\nu'$
then gives
\[
 \langle B,\widehat DA\rangle_{\nu',t}
              =\langle\widehat{D^\dagger}B,A\rangle_{\nu',t}.
\]
All expressions here are finite products of bounded local functions;
boundedness on the Hilbert quotient has not been assumed.

Set $K_D=\widehat{D^\dagger}\widehat D$ and
$m_n=\langle A,K_D^nA\rangle_{\nu',t}$. These numbers are
nonnegative. The reflected Cauchy--Schwarz inequality makes them
log-convex:
$m_n^2\le m_{n-1}m_{n+1}$. For odd $n$ use the vectors
$K_D^{(n-1)/2}A$ and $K_D^{(n+1)/2}A$ in the original form. For
even $n$, use the corresponding two vectors after applying
$\widehat D$. Their scalar product is $m_n$, and their squared
norms are the two neighbouring moments. It follows that
\[
 m_1^n\le m_0^{n-1}m_n.
\]
The formula remains valid if a moment is zero, by Cauchy--Schwarz.

The expression for $m_n$ is a column of $2n$ successive factors,
alternating $D$ and its reflected conjugate, with $A$ at one end
and its reflected conjugate at the other. Choose a coarse chessboard
whose side $L^{k_0}l$ contains the support of $A$ in its first cell
above the mirror, and whose faces are also fine-board faces. Take $n$
through multiples of $L^{k_0}$. The ends of the column then lie on
coarse faces, with the two endpoint functions in the exterior coarse
cells and every $D$ factor inside the column.

For each fixed $n$, apply Gibbs product domination along the selected
sequence $\ell_r\to\infty$. Endpoint seminorms are bounded by
$\|A\|_\infty$. On the interior coarse cells, the fine-to-coarse
comparison bounds the seminorm of the product of fine-cell factors by
the product of their fine seminorms. Reflection and conjugation leave
that seminorm unchanged. Thus
\[
 m_n\le\|A\|_\infty^2q_D^{2n},
 \qquad q_D=\limsup_r\|D\|_{\mathrm{ch},\ell_r}.
\]
Combining this inequality with log-convexity gives
\begin{equation}\label{eq:reflection-column}
 \|\widehat DA\|_{\nu',t}^{2n}
       \le\|A\|_{\nu',t}^{2n-2}\|A\|_\infty^2q_D^{2n}.
\end{equation}
Take the $2n$th root and then let $n\to\infty$ through the allowed
sequence. At this fixed cutoff $A$ is bounded, so its supremum to
power $1/n$ tends to one. Consequently
\[
 \|\widehat DA\|_{\nu',t}\le q_D\|A\|_{\nu',t},\qquad
 |\nu'(\overline{B\circ\theta_t^*}D\tau_{le_4}A)|
                   \le q_D\|B\|_{\nu',t}\|A\|_{\nu',t}.
\]
The limit order matters. The volume limit is taken for each fixed
$K,n$, then the column length tends to infinity at that fixed $K$,
and only later is $K$ removed. For a polynomial of degree $d$ in
the lattice density, $\|A\|_\infty$ can grow like $a_K^{-4d}$.
Its disappearance before the cutoff limit is what yields a bound
independent of the outer degrees. The coarse side may depend on $A$,
but no such dependence survives in this inequality.

\subsubsection{Differentiation, centring, and the continuum limit}

Apply the column estimate to the holomorphic family $D_z$ from
\eqref{eq:insertion-selected-normalizer}. Its scalar matrix elements
are holomorphic and bounded by
$Q(l)\|B\|_{\nu',t}\|A\|_{\nu',t}$ on
$|z|\|h\|_\sharp<c\rho$. Cauchy's formula at zero therefore gives
\[
 \left|\nu'\left(\overline{B\circ\theta_t^*}
       [\cO_a(h)-E_h'(0)]\tau_{le_4}A\right)\right|
 \le C\rho^{-1}Q(l)\|h\|_\sharp
                              \|B\|_{\nu',t}\|A\|_{\nu',t}.
\]
The same normalizing function works in any second periodic Gibbs state
$\nu$. Gibbs domination and the local logarithm estimate, with one
source, give
\[
 |\nu(\cO_a(h))-E_h'(0)|
                       \le C\rho^{-1}Q(l)\|h\|_\sharp.
\]
Subtracting this scalar instead changes the insertion form by at most
the same bound: the intervening lattice time translation is a
contraction in the reflection norm. We have obtained the estimate
with centring at $\nu$, even when the expectation is in $\nu'$.

For this lattice contraction use the norm argument of
\cref{prop:reflection-semigroup} at integer time steps. Its iterated
Cauchy--Schwarz bound requires only bounded translated self-pairings,
which here are at most the squared supremum of the fixed bounded
local function. No continuum strong-continuity assertion is needed
before taking this cutoff limit.

To use this in a physical slab of width $\eta$, choose $l$ comparable
to $\eta$ with room on either side of the inserted support. A bounded
number of tents of mesh comparable to $\eta$ puts the pieces of $h$
in the middle of such cells. The sum of their sharp norms costs at
most a fixed power of $\eta^{-1}\|h\|_\sharp$. The translations
needed to fit the cells consume only part of the slab. Remaining
positive time translations are contractions, so the norms of the outer
vectors are bounded by their norms after the shift $-(\eta/4)e_4$.
The power $Q(l)\le C(g)l^{-C}$ absorbs these fixed geometric losses.

Finally take the common continuum subsequence. At fixed outer
polynomials every matrix element and every squared outer norm is a
finite collection of separated moments. The state-uniform per-ball
bounds and mean transfer make their small geometric shifts
equicontinuous and pass the inequality to the limiting form. The
constant in the result comes from the cell estimate and the source
radius; it is independent of the number of fields in the outer
polynomials. This is the column and limit argument for
\eqref{eq:rotation-insertion}; the quantitative source input is the local energy and stability
induction for the reflected arrays, used in the direct derivation
of \eqref{eq:rotation-cell-stability}.


\subsection{From smooth insertions to plates}

The insertion estimate and the constant-one cone from
\cref{sec:limits} turn smooth local fields into bounded planar
insertion forms. We spell out the time margins, since these are
needed again when several insertions are multiplied.

Write planar displacements as $v=(v^0,v^1)$, where $v^0$ is time,
with planar time unit vector $e_0=(1,0)$,
and put $v_\pm=v^0\pm iv^1$. The joint spectral representation is
\[
 e^{-v^0H+iv^1P_1}
 =\exp\!\left(-v_+\frac{H-P_1}{2}
              -v_-\frac{H+P_1}{2}\right).
\]
By \eqref{eq:cone} this is a holomorphic contraction on
\begin{equation}\label{eq:rotation-tube}
 \mathbb V=\{v\in\mathbb C^2:
                    \operatorname{Re}v_+>0,
                    \operatorname{Re}v_->0\}.
\end{equation}
The insertion estimate defines a bounded operator between the
time-translated reflection vectors. Keeping its scale explicit
prevents using up the time margin twice.

In detail, let $h$ be supported in a ball about zero of radius
$c\eta$. At scale $\eta/2$ define $B_{h;\eta/2}$ by
\[
 \left\langle[\tau_{-(\eta/8)e_4}F'],
    B_{h;\eta/2}[\tau_{-(\eta/8)e_4}G']\right\rangle
       =\sigma(\Theta F'\cdot(h)\cdot G'),
\]
for $F',G'$ supported in $x_4>\eta/8$. Then
$\|B_{h;\eta/2}\|\le C\eta^{-p_*}\|h\|_\sharp$.
For $F,G$ supported in $x_4>\eta/4$, put
$\Phi=[\tau_{-(\eta/4)e_4}F]$ and
$\Psi=[\tau_{-(\eta/4)e_4}G]$. The exact translation identity is
\begin{equation}\label{eq:translated-insertion}
 \sigma(\Theta F\cdot(\tau_xh)\cdot G)
 =\left\langle\Phi,
 T_4(\tfrac\eta8e_4+x)B_{h;\eta/2}
 T_4(\tfrac\eta8e_4-x)\Psi\right\rangle,
\end{equation}
initially for real $x$ with $|x_4|<\eta/8$. Here
$T_4(z)=e^{-z_4H+i\sum_{j=1}^3z_jP_j}$. Since each $|P_j|\le H$,
this operator is a contraction when
$\operatorname{Re}z_4\ge\sum_j|\operatorname{Im}z_j|$.
Consequently \eqref{eq:translated-insertion} continues to
\[
 |\operatorname{Re}x_4|
            +\sum_{j=1}^3|\operatorname{Im}x_j|<\eta/8,
\]
which contains polydiscs of radius $c\eta$ in all four variables.
The operator at scale $\eta$ would not give the identity after the
translation, because the outer fields then have only half that
separation available.

\begin{lemma}[The integral norm of a plate]\label{lem:plate-integral}
In the reflection space just constructed, fix
$\varphi\in C_c^\infty(\R^2)$ in the transverse coordinates $(x_2,x_3)$.
For $\alpha\in C_c^\infty(\R^2)$ supported in $|x_4|<c\eta$,
with $c$ fixed sufficiently small, one has
\begin{equation}\label{eq:plate-bound}
 |\sigma(\Theta F\cdot(\alpha\otimes\varphi)\cdot G)|
 \le C_\varphi\eta^{-w_*}\|\alpha\|_{L^1(\R^2)}
                  \|F\|_{(\eta/4)}\|G\|_{(\eta/4)},
 \qquad w_*=p_*+5.
\end{equation}
The insertion form has a real-analytic representative in the planar
position. Smooth approximate identities concentrating at a point within
the allowed time strip, with uniformly bounded $L^1$ norms, converge
to its value there. At a fixed smaller time margin, the point limit
can be taken in operator norm.
\end{lemma}

\begin{proof}
The insertion estimate initially controls one derivative of a test
function. We remove that derivative by a local parametrix, using the
holomorphic translations in \eqref{eq:translated-insertion} to pay for
the derivatives that the parametrix introduces.

In four dimensions a fundamental solution of $\Delta^3$ is a constant
multiple of $|x|^2\log|x|$. It is $C^1$ at zero. Let $r=c_1\eta$,
where $c_1$ is small enough that the ball of radius $2r$ fits inside
the available translation margin. After a smooth cutoff at scale $r$,
and subtraction of the multiple of $|x|^2\log r$, choose
$n_r\in C_c^1(\R^4)$ such that
\[
 \Delta^3 n_r=\delta_0-s_r,
 \qquad
 h=(\Delta^3h)*n_r+h*s_r.
\]
The function $s_r$ is smooth and supported away from zero in a ball
of radius $2r$. Rescaling the cutoff gives
\[
 \|n_r\|_\sharp\le CMr,\qquad
 \|s_r\|_\sharp\le CMr^{-5}.
\]
Indeed $n_r(rx)$ is $r^2$ times a fixed $C^1$ function; its first
derivative costs $r^{-1}$. Applying six derivatives produces the
$r^{-4}$ scaling of $s_r$, and its first derivative costs one further
power. The subtraction of $|x|^2\log r$ removes a harmless polynomial
term and makes these scaling bounds uniform as $r\downarrow0$.

Fix the outer vectors and write $L(h)=\sigma(\Theta F\cdot(h)\cdot G)$.
The insertion estimate extends this form from smooth tests to compactly
supported $C^1$ tests by $C^1$ approximation. Smooth approximations to
$n_r$ may be chosen with their supports in the same small ball. Their
translation forms have uniform bounds on every smaller polydisc and
converge there to the holomorphic form for $n_r$.
For every real centre $y$ whose time coordinate lies in the smaller
strip, the map $z\mapsto L(\tau_{y+z}n_r)$ is holomorphic on a
four-dimensional polydisc of radius comparable to $\eta$.
By the insertion bound its modulus is at most
\[
 C\eta^{-p_*}\|n_r\|_\sharp
                 \|F\|_{(\eta/4)}\|G\|_{(\eta/4)}.
\]
The constants are independent of the spatial centre $y$: spatial
translations in the reflection space are unitary. Cauchy's estimates
therefore bound its derivatives of total order six by the same
expression times $C\eta^{-6}$.

Apply $L$ to the parametrix identity and integrate by parts six times.
The derivatives of $h$ move onto the holomorphic function of its
translation centre. The first term is bounded by
\[
 C\eta^{-p_*}\eta^{-6}\|n_r\|_\sharp\|h\|_{L^1}
                 \|F\|_{(\eta/4)}\|G\|_{(\eta/4)};
\]
the second is bounded directly by the insertion estimate with $s_r$.
Since $r=c_1\eta$, both bounds are at most
\[
 C\eta^{-p_*-5}\|h\|_{L^1(\R^4)}
                 \|F\|_{(\eta/4)}\|G\|_{(\eta/4)}.
\]
Taking $h=\alpha\otimes\varphi$ proves \eqref{eq:plate-bound},
with the fixed $L^1$ norm of the transverse label absorbed into
$C_\varphi$.

The same representation expresses $L$ as integration against a
real-analytic function of the centre on every smaller strip: the
sixth derivatives of the translated $n_r$ form and the translated
$s_r$ form both have the holomorphic extensions just used. Their
integral representation identifies the density of $L$ there. After
integrating in the transverse variables against $\varphi$, it gives
a real-analytic planar density. A smooth approximate identity with
bounded integral norm consequently converges to that density at
every point in the strip. This justifies pointlike planar insertions,
which we call plates. They remain smeared in $x_2,x_3$.

To obtain operator-norm convergence, keep a smaller strip in the same
holomorphic representation. One further Cauchy estimate bounds the
first planar derivative of its density $k_{F,G}$ by
\[
 |\nabla_y k_{F,G}(y)|
 \le C_\varphi\eta^{-w_*-1}
                 \|F\|_{(\eta/4)}\|G\|_{(\eta/4)}.
\]
The constant is independent of the outer vectors. Choose nonnegative
profiles $\alpha_\epsilon$ of integral one, supported in
$B(0,\epsilon)$, with $\epsilon<c\eta$. Integrating
$k_{F,G}(y)-k_{F,G}(0)$ against $\alpha_\epsilon$ bounds the
difference of the smeared and point forms by the right-hand side
above times $\epsilon$. The translated reflection vectors are dense.
Thus the corresponding bounded operators differ in norm by at most
$C_\varphi\eta^{-w_*-1}\epsilon$, with the same outer time
translations in both forms.
\end{proof}

All four complex translation variables are used in this proof,
although the final position of a plate has only two coordinates.

\subsection{A finite Fourier expansion in the rotation angle}

The next step is analytic continuation in the angle. The operator
product gives a holomorphic kernel in each ordered configuration
chart. Cubic symmetry identifies its quarter-turned copies. The
continuation problem is to put those charts together without losing
the bounds as an ordering changes.

The passage from a single bounded insertion to an ordered kernel uses
the available time intervals explicitly. Fix transverse labels
$\varphi_1,\ldots,\varphi_n$ and real planar positions
$y_1,\ldots,y_n$ in strict time order. Put
$v_i=y_{i+1}-y_i$, and initially suppose that every $v_i$ lies in a
fixed narrow sector about $e_0$. Choose $0<\eta\le1$ so that
\[
 \operatorname{Re}(v_i)_\pm>\eta
 \qquad(1\le i<n).
\]
There is one common scale for this finite configuration. It is chosen
below the smallest intervening tube margin, not separately according
to the degrees of the products on the two sides of an insertion.

For each $i$, approximate a planar point mass at zero by smooth
profiles of integral one and use \cref{lem:plate-integral}. The
corresponding forms, with the outer vectors translated by
$-\eta e_4/4$, converge in operator norm to a bounded operator
$B_{i,\eta}$ with
\begin{equation}\label{eq:expanded-plate-operator}
 \|B_{i,\eta}\|\le C_{\varphi_i}\eta^{-w_*}.
\end{equation}
The two time intervals of length $\eta/4$ are included in this
operator. This is a definition through its bounded form; it does not
assert that an unsmeared field at time zero is a bounded operator.
The range of the translated reflection vectors is dense, so the form
determines the operator uniquely. The integral-norm estimate and its
degree independence are what permit arbitrary separated products on
either side during this construction.

Between two such operators the unused displacement is
$v_i-\eta e_0/2$. It remains in $\mathbb V$. Translation of the
smooth insertion forms gives the corresponding ordered product.
At fixed $\eta$, the plate approximations converge in operator norm,
and the intervening translation operators are contractions. Telescoping
the finite product therefore permits the point limit and gives
\begin{equation}\label{eq:expanded-ordered-product}
 K_n(v_1,\ldots,v_{n-1})
 =\left\langle\Omega,
 B_{1,\eta}\mathcal T(v_1-\tfrac\eta2e_0)
 B_{2,\eta}\cdots
 \mathcal T(v_{n-1}-\tfrac\eta2e_0)B_{n,\eta}\Omega
 \right\rangle,
 \qquad \mathcal T(v)=e^{-v^0H+iv^1P_1}.
\end{equation}
The end intervals disappear against $\Omega$, since $H\Omega=0$.
Each intermediate factor is a holomorphic contraction. Thus the
right-hand side is holomorphic while all residual displacements have
positive light-cone margins.

Different permitted choices of $\eta$ give the same expression at
real ordered configurations, since both expressions equal the same
separated correlation. Their holomorphic continuations agree on
overlaps. On each compact subset of $\mathbb V^{n-1}$ one can choose
$\eta$ below its smallest margin, which defines the ordered kernel
throughout that tube. This is compatibility within one ordering of
the insertions. The compatibility of differently ordered angular
charts is the continuation step discussed below.

There are $n$ bounded insertion operators and $n-1$ contractions in
\eqref{eq:expanded-ordered-product}. Taking their norms gives
\[
 |K_n(v)|\le\left(\prod_{i=1}^n C_{\varphi_i}\right)\eta^{-w_*n}.
\]
Choose $\eta$ a fixed small multiple of the minimum tube margin,
capped by one. Consequently
\begin{equation}\label{eq:kernel-margin}
 |K_n(v)|\le C_{\varphi_1,\ldots,\varphi_n}
 \min\{1,\min_{i,\pm}\operatorname{Re}(v_i)_\pm\}^{-w_*n}.
\end{equation}
The exponent is linear in the number of insertions because each
insertion costs one power in \eqref{eq:expanded-plate-operator}.
A bound whose exponent or constant depended on the degrees of the
outer polynomials would not give this conclusion.

The precise constant in the continuum cone enters here. For a real
time-axis displacement $v=(r,0)$, an imaginary rotation has
\[
 R_{is}v=(r\cosh s,\;ir\sinh s),\qquad
 (R_{is}v)_\pm=re^{\mp s}.
\]
Both light-cone coordinates stay positive for every real $s$, though
their minimum decreases like $e^{-|s|}$. If one knew only
$H\ge k_*|P_1|$ with $k_*<1$, the contraction domain would instead
require
$\operatorname{Re}v^0\ge k_*^{-1}|\operatorname{Im}v^1|$.
The same rotated point leaves that domain when
$\tanh|s|>k_*$. Thus a fixed weaker lattice cone is insufficient for
the arbitrarily long imaginary rotations used to remove high angular
frequencies. The bounded five-point limit in \cref{sec:limits}
supplies the constant-one cone before this argument is applied.


\subsubsection*{Continuing across ordered charts.}
\phantomsection\label{hyp:angular-continuation}
We first continue the kernel for a configuration in a narrow sector.
Its finite angular equation then extends to each ordered tube. This
gives analytic coefficients on every configuration of distinct planar
points, including the regular polygons used below. The final
identification with correlations at projection ties is made at the
level of tests in the last subsection.

\begin{lemma}[Finite angular type]\label{lem:finite-angular-type}
Let $n\ge2$, fix real transverse labels and put
$N_n=\lfloor w_*n/4\rfloor$.
For a strictly time-ordered configuration whose successive differences
lie in a common sector of aperture $2\alpha<\pi/2$, the angular
kernel has an entire continuation with period $\pi/2$, and
\begin{equation}\label{eq:finite-harmonics}
 K_n(R_zv)=\sum_{|m|\le N_n}c_m(v)e^{4imz},\qquad
 |K_n(R_{t+is}v)|\le C_v e^{w_*n|s|}.
\end{equation}
The bound is uniform in $t$ over one real period. The coefficients
extend to real-analytic functions on the configurations of distinct
planar points. They are multilinear in the labels, invariant under
common translations and simultaneous permutations of points and
labels, and satisfy
\begin{equation}\label{eq:harmonic-covariance}
 c_m(R_t\mathbf y)=e^{4imt}c_m(\mathbf y).
\end{equation}
\end{lemma}

\begin{proof}
For a narrow configuration, let $r$ be the smallest length of its
successive differences. In a coordinate frame whose axis is closest
to the rotated sector, every real difference has time component at
least
\[
 \delta=r\cos(\alpha+\pi/4)>0.
\]
The same bound, decreased once, holds in a small product neighbourhood
of the configuration. In each of the four frames the ordered kernel
defines a holomorphic angular function on a vertical strip of width
$\pi-2\alpha$. Consecutive strips have centres separated by $\pi/2$,
so they cover the complex plane and overlap. Their overlaps contain
real intervals. On those intervals the kernels are densities of the
same correlation: a change of frame is a cubic isometry and any
reordering permutes the labels as well. Their holomorphic extensions
therefore agree on the overlaps. They give a single entire function
with period $\pi/2$.

Imaginary rotation multiplies the two light-cone coordinates by
$e^s$ and $e^{-s}$. The smallest margin in the chosen frame is thus
at least $\delta e^{-|s|}$. Equation \eqref{eq:kernel-margin} gives
the bound in \eqref{eq:finite-harmonics}, with a constant uniform
over a real period. For
\[
 c_m=\frac2\pi\int_0^{\pi/2}K_n(R_tv)e^{-4imt}\,dt,
\]
Cauchy's theorem on a rectangle, with its vertical edges cancelled
by periodicity, gives
\[
 |c_m|\le C_v e^{w_*n|s|+4ms}.
\]
Let $s\to-\infty$ for $m>w_*n/4$, and let $s\to+\infty$ for
$m<-w_*n/4$. These coefficients vanish. Fourier uniqueness and the
identity theorem give the finite expansion.

We now extend its coefficients beyond the narrow configurations.
For $\upsilon\in\frac\pi2\mathbb Z$ and a permutation $\pi$, define
the ordered tube
\[
 \mathbb D^\pi_\upsilon=
 \{\mathbf z:
    R_{-\upsilon}(z_{\pi(i+1)}-z_{\pi(i)})\in\mathbb V,
                                             \ 1\le i<n\}.
\]
Read the product kernel in this frame and ordering, with labels
$\varphi_{\pi(1)},\ldots,\varphi_{\pi(n)}$. Each tube is open and
convex, stable under common translations, entrywise real part, and
$R_{is}$ for real $s$. A nonempty intersection of two tubes therefore
contains a nonempty open set of real configurations. In small product
neighbourhoods of those real configurations, both kernels represent
the same functional on plate profiles. Continuous densities of the
same functional agree there. The identity theorem on the convex
intersection makes the kernels agree throughout it.

They consequently define one holomorphic kernel $\mathbf K_n$ on
$\mathbb D_n=\bigcup_{\upsilon,\pi}\mathbb D^\pi_\upsilon$. Its
real domain is
\[
 \mathrm{Ord}_n=
 \{\mathbf y:\text{all }y_i^0\text{ are distinct}\}
 \ \cup\
 \{\mathbf y:\text{all }y_i^1\text{ are distinct}\}.
\]
The kernel on this union is already compatible with simultaneous
permutations and translations.

For fixed $s$, the angular equation on a narrow configuration is a
holomorphic identity in the positions. More explicitly, on a tube
put $f(s;\mathbf z)=\mathbf K_n(R_{is}\mathbf z)$. The expansion
just proved gives
\[
 \prod_{m=-N_n}^{N_n}(\partial_s+4m)f(s;\mathbf y)=0
\]
on the nonempty open set of real narrow configurations in that tube.
Every derivative on the left is holomorphic in $\mathbf z$.
The identity theorem extends the equation to the entire tube, for
each real $s$. In angular notation its equation is
\begin{equation}\label{eq:outline-angular-ode}
 \prod_{m=-N_n}^{N_n}(\partial_z-4im)F(z)=0.
\end{equation}
The roots are distinct. Solving the equation in $s$ gives
\[
 \mathbf K_n(R_{is}\mathbf z)
       =\sum_{|m|\le N_n}c_m(\mathbf z)e^{-4ms}.
\]
The coefficients are fixed linear combinations of the first
$2N_n+1$ derivatives at $s=0$, and hence are holomorphic. On tube
intersections their expansions agree for every $s$; independence of
the exponentials makes their coefficients agree as well.

To pass from imaginary to real angles, take $\mathbf y\in
\mathrm{Ord}_n$ and choose one of its tubes. For $t$ in an open
interval containing zero, $R_t\mathbf y$ stays in this tube. Since
$R_{is}$ preserves the tube, the map
$t+is\mapsto\mathbf K_n(R_{t+is}\mathbf y)$ is holomorphic on
that whole vertical strip. On the imaginary axis it equals
$\sum_m c_m(\mathbf y)e^{4im(t+is)}$. The identity theorem and
uniqueness of the expansion at $R_t\mathbf y$ yield
\eqref{eq:harmonic-covariance} for those small real $t$.

Let $\mathrm{Gen}_n$ be the open dense set of configurations of
distinct planar points on which no two pair differences are
perpendicular. Density follows by excluding
finitely many zero sets of nonzero polynomials. This set is rotation
invariant and lies in $\mathrm{Ord}_n$: a time tie and a spatial
tie would give perpendicular differences, or the same pair would
coincide. Its entire real rotation orbit therefore lies in
$\mathrm{Ord}_n$. The locally constant function
\[
 t\longmapsto e^{-4imt}c_m(R_t\mathbf y)
\]
is constant on that orbit. Thus \eqref{eq:harmonic-covariance}
holds for all real $t$ when $\mathbf y\in\mathrm{Gen}_n$.
Continuity and density extend it whenever both endpoints belong to
$\mathrm{Ord}_n$. This also makes the angular continuations agree
on every finite chain of overlapping ordered charts. Compactness of
one real angular period supplies a finite chart cover of a generic
orbit and positive margins on its closed subcharts.

Finally, for any configuration of distinct planar points, choose
$t$ such that $R_t\mathbf y\in\mathrm{Ord}_n$. Only finitely many
angles modulo $\pi$ are excluded by the pair differences. Define
\[
 \widetilde c_m(\mathbf y)
       =e^{-4imt}c_m(R_t\mathbf y).
\]
The covariance just proved makes this independent of the chosen
angle. An angle admissible at one point stays admissible in a
neighbourhood, so the definition is real analytic there. It extends
$c_m$ and preserves its permutation, translation and rotation
identities. We henceforth write $c_m$ for this extension. The space
of distinct planar points is connected, being the complement of
finitely many codimension-two coincidence subspaces. Thus a
coefficient vanishing on all narrow configurations vanishes
identically. This proves both the claimed extension and the
compatibility of its maximum frequency with the definition using
narrow configurations.
\end{proof}

Let $M(n)$ be the largest absolute index of a nonzero coefficient,
over all positions and labels, with $M(n)=0$ if only the constant
coefficient remains. The lemma gives $M(n)\le w_*n/4$. For real
labels, reality gives $c_{-m}=\overline{c_m}$. Orders zero and one
are constant by normalization, translation invariance and the form
of the first moment, so $M(0)=M(1)=0$. The analytic extension
is essential at a regular polygon: its chords include perpendicular
differences, so its rotation orbit need not stay in
$\mathrm{Ord}_n$. The next argument uses the extended coefficients
at precisely these configurations.

\subsection{Convexity and the polygon argument}

Imaginary rotation preserves the reflected pairing. Write
$\theta(v^0,v^1)=(-v^0,v^1)$ for time reflection. Reversing the
order of the reflected insertions gives the induced map
$\iota(v)=-\theta(\overline v)$ on consecutive displacements:
\[
 \iota(v^0,v^1)=(\overline{v^0},-\overline{v^1}),
\]
and $R_{is}\iota=\iota R_{is}$ for real $s$. The ordered product and
its reflected product can therefore be continued through the same
positive pairing. In particular, for any finite list of separated
half-configurations $Z_1,\ldots,Z_N$ with their transverse labels,
the reflected Gram matrix $G(s)$ at imaginary angle $is$ satisfies
\begin{equation}\label{eq:expanded-boost-gram}
 \sum_{i,j=1}^N\overline{z_i}z_jG_{ij}(s)\ge0
 \qquad(z\in\mathbb C^N,\ s\in\mathbb R).
\end{equation}
Initially this is read for half-configurations with distinct internal
times, where the ordered products are defined. Approximate an arbitrary
finite list of separated half-configurations by lists of this kind.
The continued coefficients are continuous, so every matrix entry at a
fixed imaginary angle converges. The cone of positive semidefinite
finite matrices is closed; hence \eqref{eq:expanded-boost-gram}
holds also at equal internal times. This includes the polygon halves
used below.

Split one configuration into groups of $k$ and $l$ insertions across
the time mirror. Let $G_{12}(s)$ be their pairing, and let
$G_{11}(s),G_{22}(s)$ be the pairings obtained by reflecting and
doubling each group. Positivity of the two-by-two Gram matrix gives
\[
 |G_{12}(s)|^2\le G_{11}(s)G_{22}(s).
\]
Choose the positions and labels in an open set where the highest
coefficient of the $k+l$ kernel does not vanish. Its finite angular
expansion then grows at rate $4M(k+l)$ in one of the two imaginary
directions. The doubled kernels grow at rates at most $4M(2k)$ and
$4M(2l)$. Divide the inequality by the corresponding exponentials
and let $|s|$ tend to infinity in that direction. A larger rate on the
left would contradict the inequality. Hence
\begin{equation}\label{eq:degree-convexity}
 2M(k+l)\le M(2k)+M(2l).
\end{equation}
No lower bound on a doubled kernel is used: if its leading
coefficient vanishes at the chosen configuration, its smaller rate
only strengthens this restriction. In particular the nonnegative
integers $m_k=M(2k)$ form a convex sequence, since the choice of group
sizes $k-1,k+1$ gives $2m_k\le m_{k-1}+m_{k+1}$.
Here $m_0=0$, and the operator bound already gives $m_k\le Ck$.

There is a second use of \eqref{eq:expanded-boost-gram}. At even
order $2k$, let $M=M(2k)$ and take its leading coefficient
$c_*=c_{-M}$. For a fixed finite list of $k$-point
half-configurations, multiply the Gram matrix by $e^{-4Ms}$ and let
$s\to+\infty$. Every entry has a finite expansion, so it has a
limit. Positivity is closed under this finite matrix limit:
\begin{equation}\label{eq:expanded-leading-positive}
 \sum_{i,j}\overline{z_i}z_j
      c_*(\theta Z_i,Z_j)\ge0.
\end{equation}
Thus the highest angular coefficient itself is a positive kernel
across the time mirror, even though the other coefficients need not
have that property.

The covariance of the coefficients describes exactly how this
positivity changes when the mirror is turned. Let
$\theta_t=R_t\theta R_{-t}$. For the rotated half-configurations,
\[
 c_*(\theta_tR_tZ_i,R_tZ_j)
      =e^{-4iMt}c_*(\theta Z_i,Z_j).
\]
Multiplication by $e^{4iMt}$ therefore makes the reflected kernel
positive across the line turned by $t$. This conclusion concerns the
leading coefficient with its specified scalar phase. It does not
claim reflection positivity of the full functional at a new planar
mirror before rotation invariance has been proved.

Take now a regular $2k$-gon, with the same real transverse label at
every vertex. Rotation through $\pi/k$ permutes this labelled set.
Permutation compatibility leaves its coefficient unchanged, whereas
coefficient covariance multiplies it by $e^{-4iM\pi/k}$. Therefore
\[
 (1-e^{-4iM\pi/k})c_*(\hbox{regular labelled }2k\hbox{-gon})=0.
\]
If the polygon value is nonzero, it follows that
\begin{equation}\label{eq:polygon}
 k\mid2M(2k).
\end{equation}
If divisibility fails, the polygon value is zero. In a suitable
edge-bisector reflection it is a diagonal value of the phased
positive kernel above. For any positive kernel $Q$,
\[
 |Q(Z,W)|^2\le Q(Z,Z)Q(W,W).
\]
Thus $Q(Z,Z)=0$ annihilates its whole row, not merely its value at
one polygon. This gives a specific zero row with the equal-label
restriction inherited from the polygon. It is the starting point of
the geometric propagation, rather than its conclusion.


\subsubsection*{Propagating the zero row.}
\phantomsection\label{hyp:polygon-propagation}
Suppose that $k\nmid2M(2k)$, and put $M=M(2k)$ and $c_*=c_{-M}$.
The preceding polygon identity says that $c_*$ vanishes at every
regular $2k$-gon whose vertices carry one common real transverse
label. We show that it vanishes for all labels and all distinct
positions, contradicting its choice. When $k=1$ divisibility is
automatic, so assume $k\ge2$.

Call a finite labelled set $Q$ a zero set if $c_*$ vanishes at every
$2k$-point labelled configuration containing $Q$, with arbitrary
remaining distinct positions and real labels. A regular $2k$-gon
with common label $\varphi$ is a zero set. Real analyticity of the
coefficients on the full distinct-point domain is used throughout;
regular polygons themselves can have simultaneous ordering ties.

We first record the consequence of the positive-kernel zero row.
Let $\mathfrak l$ be a line and $Q'$ a set of at most $k$ labelled
points in one of its open half-planes. Suppose that
$Q'\cup\theta_{\mathfrak l}Q'$ contains a zero set. Complete $Q'$
by arbitrary additional points $W$ in that half-plane to a
$k$-point labelled configuration $Z=Q'\cup W$. Give the new points
arbitrary real labels. The reflected diagonal
$\theta_{\mathfrak l}Z\cup Z$ contains the zero set, so its
$c_*$-value is zero. After multiplication by the fixed phase
specified above for this mirror, it is a zero diagonal of a positive
kernel. Its whole row vanishes. Therefore
\[
 c_*(\theta_{\mathfrak l}Q'\cup
          \theta_{\mathfrak l}W\cup T)=0
\]
for every $k$-point labelled configuration $T$ in the same open
half-plane, with its positions and real labels independent of $W$.

For each fixed choice of these remaining labels, their
$2k-|Q'|$ positions vary in a nonempty open set. The space of
distinct planar positions avoiding the finitely many fixed points
of $\theta_{\mathfrak l}Q'$ is connected: one may move points
successively along planar paths avoiding the other finitely many
points. Real analyticity extends the displayed zero to that entire
space. Thus $\theta_{\mathfrak l}Q'$ is a zero set. Interchanging
the half-planes proves the same for $Q'$. The arbitrary labels of
$W$ and $T$ have been released by the row identity; they have not
been inferred from an equal-label diagonal by polarization.

Now fix a regular $2k$-gon with cyclically numbered vertices
$v_i$, $i\in\mathbb Z/(2k)$, and common label $\varphi$.
We prove the following reduction: if every consecutive arc of
$j$ vertices is a zero set, then every consecutive arc of
$q=\lceil j/2\rceil$ vertices is a zero set, for $2\le j\le2k$.
Choose such a shorter arc
\[
 Q'=\{v_{i_0+1},\ldots,v_{i_0+q}\}.
\]
Let $\mathfrak l$ be the perpendicular bisector of the polygon edge
$[v_{i_0},v_{i_0+1}]$. It passes through the centre, meets no
vertex, and acts by
\[
 \theta_{\mathfrak l}v_{i_0+i}=v_{i_0+1-i}.
\]
The $k$ vertices $v_{i_0+1},\ldots,v_{i_0+k}$ lie in one open
half-plane, and $q\le k$. The union
$\theta_{\mathfrak l}Q'\cup Q'$ is a consecutive arc of $2q$
vertices. Since $2q\ge j$, it contains a zero arc of length $j$.
The zero-row consequence just proved makes $Q'$ a zero set.

Starting with the whole polygon, repeat this reduction:
\[
 2k\longrightarrow k\longrightarrow\lceil k/2\rceil
                       \longrightarrow\cdots\longrightarrow1.
\]
It works for every integer $k$, with no power-of-two restriction.
Each single vertex with label $\varphi$ is therefore a zero set.
Any planar point can be a vertex of a suitably translated regular
$2k$-gon, and $\varphi$ was arbitrary. Every other real label was
already arbitrary in the zero-set property. Hence $c_*$ vanishes
for all real labels and all distinct planar positions.
Multilinearity extends the conclusion to complex labels. This
contradicts its nonvanishing and proves \eqref{eq:polygon}.

The remaining frequency argument is arithmetic.

For completeness, suppose $m_0=0$, the $m_k$ are integers, their
successive increments are nondecreasing, $m_k\le Ck$, and
$k\mid2m_k$. Every increment is at most $C$: a larger increment
would persist by monotonicity and violate the linear bound. The
increments, being integers, are therefore eventually constant, say
$D$. Write $m_k=Dk-b$ for all large $k$. Divisibility says that each
such $k$ divides the fixed integer $2b$, hence $b=0$. Since earlier
increments are at most $D$, their sum can equal $Dk$ only if every
one is $D$. Therefore $m_k=km_1$. Combining this with
\eqref{eq:degree-convexity} gives
\begin{equation}\label{eq:degree-linearity}
 M(2k)=kM(2),\qquad M(n)\le\tfrac n2M(2).
\end{equation}
The divisibility condition restricts the highest index at each
order. It is stronger than the assertion that angular frequencies
are multiples of four. Cauchy--Schwarz supplies convexity, while
the polygon argument supplies this extra restriction.

\subsection{The ultraviolet estimate for two plates}

The angular base case requires only $t^8\Gamma(t)\to0$. We first
derive it from the two-point estimate of \cref{sec:two-point}, and
then obtain it from a weaker bound supplied by source stability. The
second derivation separates rotations from the precise nontriviality
calculation.

The equal-label kernel is read from the same separated
two-point distribution as the ordered kernels.
It retains positivity at the diagonal mirror
from \cref{sec:limits}; this is the premise used
when identifying its diagonal reflection picture below.

For equal real transverse labels the two-point kernel has the
positive spectral representation
\[
 K_2(v)=\int_{E\ge|q|}e^{-v^0E+iv^1q}\,d\rho(E,q),
 \qquad \rho\ge0.
\]
Set \(\Gamma(t)=\int e^{-tE}\,d\rho\).
The ordered plate kernel is finite at positive unit
separation, so \(\Gamma(1)=K_2(e_0)<\infty\).
Hence \(e^{-E}\rho\) is finite, and
\(\rho(\{E\le N\})\le e^N\Gamma(1)\);
these sets exhaust its nonnegative-energy carrier.
Thus the measure is sigma-finite before the
Tonelli computations below.
The upper two-point estimate
from \cref{sec:two-point}, for nonnegative bumps of maximum one
whose size and separation are comparable to $d$, has the form
\[
 0\le C_2(f,h)\le u(d)\|f\|_\infty\|h\|_\infty,
 \qquad u(d)\longrightarrow0.
\]
The vanishing of $u$ comes from the running coupling. It must be
available at the fixed shape used in the decomposition of plates.

Choose a nonzero smooth bump $0\le\beta\le1$ supported in $[1,2]$,
and the cubic B-spline
$\gamma=c_\gamma\mathbf1_{[-1/2,1/2]}^{*4}$, normalized by
$\|\gamma\|_\infty=1$. For $\kappa\in[1,2]$, put
\[
 \alpha_{t,\kappa}(s,y)=\beta(s/t)\gamma(8y/(\kappa t)),
 \qquad f_{t,\kappa}=\alpha_{t,\kappa}\otimes\varphi.
\]
These planar profiles have width and height comparable to $t$,
are placed at positive time of order $t$, and have maximum at most
one. With sufficiently small transverse cells, their reflected
pairs fit in balls of radius $t$ and have separation between $t$
and $6t$. The comparison needed below is
\begin{equation}\label{eq:profile-comparison}
 \Gamma(Ct)\le Ct^{-4}\sup_\kappa\|[f_{t,\kappa}]\|^2.
\end{equation}
To verify it, write $\operatorname{sinc}u=\sin u/u$, with value one
at zero. The spatial Fourier transform of $\gamma$ is
$c_\gamma\operatorname{sinc}(q/2)^4$, while
\[
 \int\beta(s/t)e^{-sE}\,ds\ge cte^{-2tE}
 \qquad(E\ge0).
\]
The Fourier--Laplace representation of the reflected norm gives
\[
 \|[f_{t,\kappa}]\|^2
 \ge ct^4\int e^{-4tE}
             \operatorname{sinc}(\kappa tq/16)^8\,d\rho(E,q).
\]
The spline is $C^2$ and can be used by smooth approximation in
$C^1$, retaining the positive-time support. The sharp-norm bounds
give convergence of the reflected norms; the transforms are
dominated by $Ct^2e^{-tE}$, so the spectral formula passes to the
limit as well. No finite undamped mass of $\rho$ is required.

The dilation average satisfies
\[
 J(u):=\int_1^2\operatorname{sinc}(\kappa u)^8\,d\kappa
             \ge c(1+|u|)^{-8}.
\]
Indeed, $J$ is continuous and strictly positive on each compact
interval, since the dilated transforms have no common zero. For
$|u|\ge2\pi$,
\[
 J(u)\ge2^{-8}|u|^{-9}
                  \int_{|u|}^{2|u|}\sin^8v\,dv
       \ge c|u|^{-8},
\]
using the positive mean of $\sin^8v$. Tonelli's theorem and the cone
$|q|\le E$ now give
\[
 \int_1^2\|[f_{t,\kappa}]\|^2\,d\kappa
 \ge ct^4\int\frac{e^{-4tE}}{(1+t|q|)^8}\,d\rho(E,q)
 \ge ct^4\Gamma(5t).
\]
Here $(1+u)^8\le Ce^u$ for $u\ge0$. This proves
\eqref{eq:profile-comparison} with fixed constants.

For a fixed transverse label, a profile $f_{t,\kappa}$ decomposes
into $O(t^{-2})$ transverse cells, each a difference of controlled
nonnegative four-dimensional bumps. The reflection norm of each
piece is bounded by $Cu(ct)^{1/2}$. The triangle inequality yields
\[
 \|[f_{t,\kappa}]\|\le Ct^{-2}u(ct)^{1/2}.
\]
Together with \eqref{eq:profile-comparison} this gives
\begin{equation}\label{eq:little-o}
 t^8\Gamma(t)\le Cu(ct)\longrightarrow0.
\end{equation}
The exponent eight contains both the planar transform loss and
the squared count of transverse cells. No estimate on a point
insertion is inferred directly from a bound on an arbitrary test.

\subsubsection*{A bound sufficient for rotations without the relative comparison.}

The limit in \eqref{eq:little-o} needs less than the precise asymptotic
law of \cref{sec:two-point}. We give a second derivation from the
coupling dependence of the source bounds. In particular, the relative
Gaussian comparison is needed for \cref{it:nontrivial}, but need not
be used to prove rotations.

We first record that dependence. With $L,M,\rho$, the analytic reserves
and the activity rates fixed, the constants in the local source-energy
bound and the comparison of vacuum densities are independent of the
reference coupling $0<g'\le\gamma$. At terminal inverse coupling $y$,
the upper stability and lower-floor exponents can be taken to satisfy
\begin{equation}\label{eq:terminal-stability-dependence}
 E_+(2\rho;y)+C_{\rm low}(y)+\delta_{\rm vac}
       \le C(1+\rho+\log y).
\end{equation}
This is the normalization used in the source argument: each omitted
Gaussian factor costs $C(1+\log x_j)$ per site, the weighted history
estimate pays those costs, and the final small ball of radius
$x_k^{-1/2}$ has bounded action and Haar-volume cost $C\log x_k$.
The constants are chosen before the coupling ceiling. The auxiliary
torus may have to grow as $g'$ decreases. Its side does not enter the
per-cell estimate \eqref{eq:per-ball-cell-product}; division by the
number of reflected cells leaves the same cell volume. Thus, for two
tests in unit balls at a fixed positive separation, that proof gives
\begin{equation}\label{eq:per-ball-coupling-dependence}
 |C_2(f,h)|\le C e^{C(1+\log(g'^{-2}+B_\sharp))}
                         \|f\|_\sharp\|h\|_\sharp
\end{equation}
in every periodic Gibbs state at first-passage data for $g'$. This
retains the uniform source and history estimates required earlier;
zero-source stability alone would not give it.

\begin{lemma}[A subpower bound at the observation scale]
\label{lem:subpower-two-point}
Fix a shape of two nonnegative bumps whose size and separation are
comparable to $d$, where $0<d\le d_0<1$. The source bounds imply, at all sufficiently fine
first-passage cutoffs,
\begin{equation}\label{eq:subpower-two-point}
 |C_2(f,h)|\le u_*(d)\|f\|_\infty\|h\|_\infty,
 \qquad
 u_*(d)=C\bigl(g^{-2}+B_\sharp(C+\log_L(1/d))\bigr)^C.
\end{equation}
The cutoff threshold may depend on $d$, but the bound does not.
In particular $d^2u_*(d)\to0$. The same bound holds in continuum
limits along these states.
\end{lemma}

\begin{proof}
Call $k$ a record level if $x_k<x_i$ for all $i<k$. The first-passage
level $K$ is a record level. If $i$ is one record level, the lower
drift bound shows that another occurs within
$j_0=1+\lceil C_\sharp/\lambda_\sharp\rceil$ steps, unless the run
has already ended. Choose $b$ on the lower side of $\log_L(1/d)$,
within $j_0+1$ of it, so that $K-b$ is a record level and
$L^{-j_0-1}\le L^bd\le1$. Take $d_0\le L^{-j_0-2}$, so $b\ge1$.
It is possible once
$K\ge C+j_0+\lceil\log_L(1/d)\rceil$.

Set $g_b^{-2}=x_{K-b}-B_\sharp$. For $b>0$ this is positive,
$g_b<g$, and $(K-b,\beta)$ is a first-passage datum for $g_b$:
all earlier $x_i$ exceed $x_{K-b}$. The upper drift bound gives
$x_{K-b}\le g^{-2}+B_\sharp(b+1)$. The coupling sequence is the
same sequence in either choice of units, since its defining local
operations and extraction do not depend on the terminal reference
coupling.

The dilation $x\mapsto L^bx$ takes a spacing-$L^{-K}$ Gibbs state
to one at spacing $L^{-(K-b)}$, with the same Wilson weight. It also
gives the exact equality
\[
 C_2(f,h)=C_2(f^{[b]},h^{[b]}),\qquad
 f^{[b]}(x)=f(L^{-b}x).
\]
The sup norms do not change, and
$\Lip(f^{[b]})=L^{-b}\Lip f\le C\|f\|_\infty$,
with a constant depending on the fixed shape and $j_0$ only.
Apply \eqref{eq:per-ball-coupling-dependence} at $g_b$ to this pair
at fixed separation. This gives \eqref{eq:subpower-two-point}.
For fixed $d$ the cutoff requirement is eventually satisfied, so
the estimate passes to the continuum before $d$ decreases to zero.
\end{proof}

To keep the decay between transverse cells, we use a point kernel in
four dimensions. The construction of that kernel is included here
because a bound on a smeared plate alone does not give it.

\begin{lemma}[The point kernel and its axis bound]
\label{lem:four-dimensional-point-kernel}
Before proving rotations, the two-point function has a real-analytic
kernel $D$ off the origin, invariant under cubic isometries, with
\begin{equation}\label{eq:four-dimensional-spectral-kernel}
 D(te_4+z)=\int e^{-tE+ip\cdot z}\,d\rho_4(E,p),\qquad t>0,
 \qquad \rho_4\ge0,
 \qquad |p_j|\le E.
\end{equation}
Here $z\in\R^3$ is embedded as a spatial vector. Its axis value
$s_a(t)=D(te_4)$ is finite, nonnegative and
nonincreasing, and
\begin{equation}\label{eq:axis-domination}
 |D(x)|\le s_a(|x|_\infty)\le s_a(|x|/2),\qquad x\ne0.
\end{equation}
If $C_2(\theta h,h)\le u(d)\|h\|_\infty^2$ for the fixed-shape
mirror bumps of size $d$, then
\begin{equation}\label{eq:axis-from-mirror-bumps}
 s_a(t)\le Ct^{-8}u(ct).
\end{equation}
\end{lemma}

\begin{proof}
Use the joint spectral measure $\Pi$ of $H,P_1,P_2,P_3$ on the
one-field subspace. For positive-time tests write
$\widetilde f(E,p)=\int f(x)e^{-x_4E+ip\cdot\vec x}\,dx$ and
$m_f=\langle[f],\Pi(\cdot)[f]\rangle$. Convolution of two such
tests gives
$[h*f]=\widetilde h(H,P)[f]=\widetilde f(H,P)[h]$, hence
\[
 |\widetilde h|^2m_f=|\widetilde f|^2m_h.
\]
Choose a nonnegative time bump translated strictly into positive
time, and in each spatial variable use
$b=(v*v^\vee)^{*2}$, where $v(x)= e^{-x}\mathbf1_{[0,1]}(x)$
and $v^\vee(x)=v(-x)$.
This is a compactly supported $C^1$ shape and
\[
 \widehat b(q)=\frac{|1-e^{-1+iq}|^4}{(1+q^2)^2}
                     \ge c(1+q^2)^{-2}>0.
\]
The product test $f_*$ therefore has a nowhere-zero transform on
$E\ge0$. Set $\rho_4=|\widetilde f_*|^{-2}m_{f_*}$.
Here \(m_{f_*}(\mathbb R^4)=\|[f_*]\|^2<\infty\).
The sets where \(|\widetilde f_*|\ge1/N\)
exhaust the spectral carrier and have
\(\rho_4\)-mass at most \(N^2\|[f_*]\|^2\).
This makes \(\rho_4\) sigma-finite despite
division by a transform that need not be
bounded away from zero globally.

The exchange identity gives $m_f=|\widetilde f|^2\rho_4$ for
every positive-time test. Its carrier satisfies the three coordinate
cone inequalities because $\Pi$ does.

Apply this identity to
$h_t(x)=a(x_4/t)\prod_{j=1}^3b(x_j/t)$, translating the fixed
time shape $a$ high enough that $(\theta h_t,h_t)$ belongs to a
fixed mirror-pair shape class. The time transform is bounded below
by $ct e^{-CtE}$; the spatial transforms are bounded below by
$ct(1+t^2p_j^2)^{-2}$. The cone absorbs all polynomial losses into
an exponential, so
\[
 \|[h_t]\|^2\ge ct^8\int e^{-C'tE}\,d\rho_4.
\]
The left side is finite by the per-ball bounds. Thus every positive
time damps $\rho_4$ to a finite measure. The spectral integral in
\eqref{eq:four-dimensional-spectral-kernel} is well-defined and
holomorphic when the imaginary spatial displacement is smaller than
the positive time margin. Polarization and the spectral norm formula
identify it with the original two-point distribution on $x_4>0$.
Cubic changes of frame extend it over $x\ne0$; the extensions agree
on overlaps because they represent the same distribution. This also
proves real analyticity there. The mirror-pair bound in the same
inequality gives \eqref{eq:axis-from-mirror-bumps}, after changing
fixed constants.

Taking absolute values in the spectral integral gives
$|D(te_4+z)|\le s_a(t)$. A cubic isometry can put the coordinate
of largest modulus in the positive time direction. Monotonicity of
the Laplace transform proves \eqref{eq:axis-domination}.
\end{proof}

For the fixed transverse label $\varphi$, the plate kernel is
\[
 \Gamma(t)=\int_{\R^2}D(te_4+\zeta)
                         (\varphi\star\varphi)(\zeta)\,d\zeta,
 \qquad
 (\varphi\star\varphi)(\zeta)
       =\int\varphi(y)\varphi(y+\zeta)\,dy.
\]
Its spectral measure is the image of
$|\widehat\varphi(p_2,p_3)|^2\rho_4$ under $(E,p)\mapsto(E,p_1)$;
in particular it is the same positive measure defining $\Gamma$
above. Cauchy--Schwarz gives
$|\varphi\star\varphi|\le\|\varphi\|_2^2$. Split its compact
support into $|\zeta|_\infty\le t$, the annuli
$2^it<|\zeta|_\infty\le2^{i+1}t$ with $2^it\le\delta$, and the remainder.
The annular area is $O(2^{2i}t^2)$, while
\eqref{eq:axis-from-mirror-bumps} bounds the kernel by
$C(2^it)^{-8}u(c2^it)$. Hence
\begin{equation}\label{eq:transverse-decay-growth}
 \Gamma(t)\le C_\varphi\left[
 t^{-6}\sum_{2^it\le\delta}2^{-6i}u(c2^it)
                      +s_a(\delta)\right].
\end{equation}
After multiplication by $t^8$, the sum is at most
$C\sup_{0<d\le c\delta}d^2u(d)$, because its remaining weights
are $2^{-8i}$. First choose $\delta$ small, then let $t\to0$.
Thus $d^2u(d)\to0$ implies \eqref{eq:little-o}. Taking
$u=u_*$ from \cref{lem:subpower-two-point} supplies this condition
from the source estimates alone. The coordinate cone and diagonal
positivity are still required in the subsequent harmonic argument.

In the common angular continuation, at a fixed separation $v=re_0$, move the angle contour in
\eqref{eq:finite-harmonics}. Positivity and the cone bound the
coefficient by
\[
 |c_m(re_0)|\le e^{4ms}\Gamma(cr e^{-|s|}).
\]
The power bound eliminates $|m|>2$, and the strict little-$o$ in
\eqref{eq:little-o} eliminates $|m|=2$. A mere $O(t^{-8})$ bound
would leave precisely that endpoint unresolved.

\subsection{The last two harmonics and diagonal positivity}

The remaining possible terms are indexed by $m=0,\pm1$. At
$v=re_0$ put $a_m=c_m(re_0)$. Imaginary rotations give the transform
$a_0+a_1e^{-4s}+a_{-1}e^{4s}$. At the diagonal angle $\pi/4$ the
nonconstant terms have the opposite sign. Diagonal reflection
positivity represents that rotated two-point kernel by another
positive measure $\rho_{\mathrm d}$.

Let $b_{\pm1}$ be the Lorentz boosts of the planar energy--momentum
variables with rapidities $\pm1$. The identity of the two transforms
at the three remaining harmonics yields the required measure identity
\begin{equation}\label{eq:measure-comparison}
 (\cosh4-1)\rho_{\mathrm d}
 +(b_1)_*\rho+(b_{-1})_*\rho
 =(\cosh4+1)\rho.
\end{equation}
For the uniqueness step, each measure has nonnegative energy and
finite positive-time damping. For $\rho$ and $\rho_{\mathrm d}$ this
follows from the plate bounds; the boosted measures satisfy it
because their energies are at least $e^{-1}E$. Let $M$ be the
difference of the two sides of \eqref{eq:measure-comparison}. For
each fixed real spatial displacement $y$, its transform
\[
 F_M(z,y)=\int e^{-zE+iyq}\,dM(E,q)
\]
is holomorphic on $\operatorname{Re}z>0$. Under the stated angular
continuation, the three-harmonic identity makes it zero on the
common real ordering sector $t>0$, $|y|<ct$, for a fixed $c>0$.
For this fixed $y$, the identity theorem extends its vanishing
from $t>|y|/c$ to the whole right half-plane. Thus, for every
$t_0>0$, the two finite positive measures obtained by damping the
sides of \eqref{eq:measure-comparison} by $e^{-t_0E}$ have the same
joint Fourier transform, obtained by setting $z=t_0+iu$. Fourier
uniqueness \cite[Theorem, p.~1]{Grimmett2004} makes these damped
measures equal, and hence proves \eqref{eq:measure-comparison}.
This uses only nonnegative energy for $\rho_{\mathrm d}$.
Positivity now gives
$\rho_{\mathrm d}\le C\rho$, so $\rho_{\mathrm d}$ is supported
in the same cone. Its transform can therefore also be evaluated at
all the imaginary rotations under consideration. It follows that
\begin{equation}\label{eq:two-signs}
 a_0-a_1e^{-4s}-a_{-1}e^{4s}\ge0\quad(s\in\R),
 \qquad a_{-1}\ge0.
\end{equation}
The second inequality is positivity of the leading Gram coefficient
on a reflected time-axis pair with equal labels. Multiply the first
inequality by $e^{-4s}$ and let $s\to+\infty$ to get $a_{-1}\le0$.
Thus $a_{-1}=0$. Translation and the covariance of the coefficient
extend its vanishing from that axis to every pair, polarisation
extends it to unequal labels, and reality gives $a_1=0$.
Consequently the two-point argument gives $M(2)=0$. The polygon
propagation and \eqref{eq:degree-linearity} then pass this conclusion
to every $n$. Thus the ultraviolet estimate settles the base case,
while the geometric argument carries it through the hierarchy.

\subsection{From ordered plates to separated tests}

The analytic coefficients above were extended through configurations
with projection ties. We must identify their sum with the original
correlation on test functions there. Equality of densities only on
the generic set would not exclude a distribution supported on its
complement. The following local translation identity supplies the
needed passage at the level of tests.

Take two clusters $A,B$ on opposite sides of a coordinate slab of
width $2\delta>0$, smaller than their normal separation. Remove the
two intervals of length $\delta$ and let $a,b$ be the resulting
reflection vectors. For a vector $u$ tangent to the slab,
\begin{equation}\label{eq:expanded-slab-sliding}
 \sigma(A\cdot\tau_{ru}B)
 =\langle a,e^{-2\delta H+iru\cdot P}b\rangle
 =\int e^{-2\delta E+iru\cdot q}
                                \,d\varsigma_{A,B}(E,q).
\end{equation}
The complex measure $\varsigma_{A,B}$ has finite total variation.
The coordinate cones give $|u\cdot q|\le |u|_1E$, so the expression
is holomorphic on $|\operatorname{Im}r|\,|u|_1<2\delta$. On every
smaller strip the exponential also dominates each fixed derivative.
For a different coordinate normal, use its reflection space.
This formula applies while the normal gap remains open, regardless
of ties in a tangential projection. Plate point limits are taken at
a fixed smaller gap using the insertion estimates already proved.

Put $\widetilde K_n(\mathbf y)=\sum_m c_m(\mathbf y)$, with the
extended analytic coefficients. For fixed transverse labels and
pairwise disjoint compact planar profiles, the identity to prove is
\[
 \sigma((\alpha_i\otimes\varphi_i)_i)
   =\int\prod_i\alpha_i(y_i)\,
                         \widetilde K_n(\mathbf y)\,d\mathbf y.
\]
It holds whenever the profiles are ordered in time or in the planar
spatial coordinate, by the construction of $\mathbf K_n$. We prove
it on a neighbourhood of every configuration of distinct planar
points, by induction on the number of pairs tied in their spatial
coordinate $y^1$.

If $y_i^1=y_j^1$ and $y_i\ne y_j$, then $y_i^0\ne y_j^0$.
Choose a time level strictly between these times, avoiding all the
other times. It divides the configuration into an earlier cluster
and a later one. Small product neighbourhoods retain a positive
time gap. Slide the later cluster by $re_1$. For a sufficiently
small positive $r_0$, the slide removes every spatial tie between
the two clusters and creates no new one. Ties within each cluster
are unchanged, so the number of tied pairs decreases.

The original correlation is real analytic in $r$ by
\eqref{eq:expanded-slab-sliding}. The integral of
$\widetilde K_n$ against the translated profiles is real analytic
as well: the time gap prevents collisions during the slide, and a
real-analytic function on the distinct-point domain has a bounded
holomorphic extension on a neighbourhood of each compact set of
positions used in the integral. Near $r_0$ the two expressions agree
by the induction hypothesis, for every choice of profiles in the
chosen product neighbourhoods. The identity theorem in $r$ gives
agreement at zero. To globalise, cover the compact product of the
profile supports by finitely many products of discs where the local
identity holds. Split each profile by a smooth partition with supports
small enough that every nonzero product term lies in one member of
this cover. Multilinearity proves the formula for the original
profiles. It is an identity of distributions on each separated
product neighbourhood; thus no additional term
supported only at a simultaneous time and spatial ordering tie
remains.

Once the polygon and two-point arguments have shown $M(n)=0$, this
formula and a change of planar variables give rotation invariance
for plates with distinct planar supports:
\[
 \begin{split}
 \sigma(((\alpha_i\circ R_{-t})\otimes\varphi_i)_i)
 &=\int\prod_i\alpha_i(y_i)\,
                             \widetilde K_n(R_t\mathbf y)\,d\mathbf y\\
 &=\int\prod_i\alpha_i(y_i)\,
                             \widetilde K_n(\mathbf y)\,d\mathbf y.
 \end{split}
\]

Consider next a configuration of distinct points in $\R^4$ whose
planar projections are distinct. In small product neighbourhoods,
the planar projections lie in disjoint discs. Fourier approximation
on product boxes, multiplied by fixed product cutoffs, approximates
each smooth test in $C_c^\infty$ by finite sums
$\alpha(x_4,x_1)\varphi(x_2,x_3)$ with supports in those
neighbourhoods. Multilinearity, the plate identity and separated-test
continuity give rotation invariance on such neighbourhoods.

It remains to allow coincident planar projections while the full
points remain distinct. Suppose one pair has the same projection.
Its points differ in $x_2$ or $x_3$. A hyperplane of constant value
of that coordinate separates them and, after choosing the level
away from all other points, separates the configuration into two
clusters. Retain a positive normal gap in small neighbourhoods and
slide one cluster by $re_4$. A small nonzero displacement removes
all planar coincidences between the clusters without creating new
ones. The number of coincident pairs therefore decreases.

Fix the rotation angle $t$ and compare
$\sigma(R_tX)$ with $\sigma(X)$ during this slide. The transverse
separating hyperplane is preserved by $R_t$, and both $e_4$ and
$R_te_4$ are tangent to it. Both expressions are thus real analytic
in $r$ by \eqref{eq:expanded-slab-sliding}, in the corresponding
coordinate picture. They agree near the displaced configuration by
induction, and hence at the original one. This again is an identity
for arbitrary tests in product neighbourhoods, rather than only an
identity of point values away from the tie set. Induction and a
partition of unity cover every configuration outside the full
coincidence diagonals.

The separated-test continuity bound finally extends the identity to
the $C^1$ tests used earlier. Rotations in this plane, their cubic
conjugates in all coordinate planes, and a coordinate reflection
generate $\mathrm O(4)$. The moment--cumulant formula transfers the
identity between moments and connected functions. Thus, once the
sharp cone and degree-independent insertion bounds are available,
the ordered angular kernels, their analytic coefficient extension,
the polygon argument and the two-point base case give rotation
invariance on every separated tuple. The uniform construction
estimates used to obtain those inputs remain the analytic tasks
identified earlier.

\section{Selecting a phase and proving clustering}\label{sec:phases}

The centred limit can contain several translation-invariant phases.
Even if it is Euclidean invariant and reflection positive, correlations
between distant groups need not factor. Phase selection therefore
requires information beyond the compactness argument of
\cref{sec:limits}. The useful estimate is a rate of approach to the
translation-invariant plateau, uniform both in the lattice spacing
and in the Gibbs state.

\subsection{The finite-lattice decomposition and its coordinates}

\subsubsection*{The Gibbs ergodic decomposition.}
\phantomsection\label{hyp:gibbs-decomposition}
The ergodic decomposition at a fixed spacing is a property of the
Wilson Gibbs specification. To apply the Gibbs decomposition theorem,
regard one lattice site as carrying the four positively oriented
outgoing bonds. Its state space is $E=\SU(2)^4$, a compact standard
Borel space, and the configuration space is $E^{\mathbb Z^4}$. The
plaquette interaction is bounded and has finite range. Its normalized
finite-region kernels are translation covariant, and the reference
state is translation invariant. The Gibbs decomposition theorem
\cite[Theorem~14.17, p.~298]{Georgii2011} therefore decomposes it
into translation-ergodic states satisfying
these same Wilson DLR kernels. Selecting a component changes the
state, not the interaction or the conditional law in a finite region.
Every finite bond set is contained in finitely many of these site
variables. Conditioning the finite-site Wilson kernel further on the
unused outgoing bonds recovers the original bond kernel, so this
encoding preserves the specification used throughout the manuscript.

This standard Gibbs theorem supplies the decomposition used here.
Reflection positivity and the uniform estimates for its components
are obtained by applying the lattice reflection and source arguments
to the full Gibbs class. The uniform clustering estimate below is
what permits pure states to vary with the cutoff.

In this section a pure state means a translation-ergodic Gibbs state.
Let $\mathcal I$ be the sigma-algebra of translation-invariant events.
The Gibbs ergodic decomposition gives
\[
 \nu=\int\rho\,D_\nu(d\rho),\qquad
 E_\nu[A\mid\mathcal I](U)=\rho_U(A)
\]
for bounded local $A$, with $D_\nu$ concentrated on ergodic Gibbs states.
The conditional states are shift invariant and belong to the full
periodic Gibbs class. The source estimates therefore apply to almost
every ergodic component with the same constants.

For the reference sequence $\nu_j$ define
\begin{equation}\label{eq:phase-centring}
 X_j^{\nu_j}=\prod_{f\in X}
       [\cO_{a_j}(f)-\nu_j(\cO_{a_j}(f))].
\end{equation}
The subtraction is common to all states at datum $j$. The source
normalizer gives
\begin{equation}\label{eq:phase-mean-bound}
 |\rho(\cO_{a_j}(f))-\nu_j(\cO_{a_j}(f))|
       \le C_1(R)\|f\|_\sharp
\end{equation}
for support in a ball of radius $R$, uniformly in $j$ and $\rho$.
This bounds the first coordinates even when the two uncentred means
diverge separately.

\begin{proposition}[Compact component coordinates]\label{prop:component-coordinates}
Let $S^{\mathrm{ref}}$
be a common volume-first limit of the reference-centred moments.
There are a compact metrisable coordinate space $\mathsf X$, a
probability measure $\lambda$ of full support, and continuous
coordinates $\xi\mapsto\sigma_\xi(X)$ for every separated tuple,
such that
\begin{equation}\label{eq:component-barycentre}
 S^{\mathrm{ref}}(X)=\int_{\mathsf X}\sigma_\xi(X)\,d\lambda(\xi).
\end{equation}
Each $\sigma_\xi$ is a limit of the moments in pure states $\rho_j$,
with the subtraction at $\nu_j$ as in \eqref{eq:phase-centring}.
For any finite list of tuples,
\begin{align}\label{eq:component-plateaux}
 \int_{\mathsf X}\prod_{i=1}^q\sigma_\xi(X^i)\,d\lambda(\xi)
 =\lim_{|v_2|\to\infty}\cdots\lim_{|v_q|\to\infty}
 S^{\mathrm{ref}}(X^1\cdot\tau_{v_2}X^2\cdots\tau_{v_q}X^q).
\end{align}
The separated correlations of each $\sigma_\xi$ satisfy the uniform
clustering estimate in \eqref{eq:cluster}.
\end{proposition}

\begin{proof}
Use a countable family \((X_m)_{m\ge1}\) of real separated
tuples with the support-controlled approximation property
constructed above. Include the empty tuple. At datum \(j\)
evaluate
\[
 \rho\longmapsto
                   (\rho((X_m)_j^{\nu_j}))_{m\ge1}.
\]
For each fixed coordinate, \eqref{eq:fixed-support-moments}
gives a bound \(B_m\) for all sufficiently large \(j\),
uniformly in \(\rho\). The finitely many earlier values
are bounded too: at fixed spacing each field product
is a bounded local continuous function. Enlarge \(B_m\)
to include them. Thus one compact product of intervals
contains all the evaluation measures.


The product $\mathsf K=\prod_m[-B_m,B_m]$ is compact and metrisable.
For example the metric
\[
 d(u,v)=\sum_{m\ge1}2^{-m}
       \min\{1,|u_m-v_m|\}
\]
induces its product topology. Push $D_{\nu_j}$ forward to a probability
measure $\widehat D_j$ on this product and pass to a weakly convergent
subsequence with limit $\lambda$. Put $\mathsf X=\operatorname{supp}\lambda$.
For each coordinate the defining decomposition gives
\[
 \int u_m\,d\widehat D_j(u)
       =\nu_j((X_m)_j^{\nu_j}).
\]
Both sides have limits, proving \eqref{eq:component-barycentre} first
on the countable family.

Fix $\xi\in\mathsf X$. Every open neighbourhood of $\xi$ has positive
$\lambda$-mass and hence positive $\widehat D_j$-mass for all
sufficiently large $j$. Its preimage contains a pure state because
$D_{\nu_j}$ is concentrated on such states. Choose successively
neighbourhoods restricting the first $N$ coordinates to accuracy
$1/N$, and let $N$ increase slowly with $j$. This gives pure $\rho_j$
whose coordinates converge to $\xi$.

For an arbitrary separated tuple \(X=(f_i)\), choose
pairwise disjoint compact neighbourhoods \(A_i\) of its
supports and approximate it using the countable family
with all approximants supported in these same \(A_i\).
The estimate \eqref{eq:fixed-support-continuity} is
uniform in \(j\) and the pure evaluating state.
Consequently the limit of \(\rho_j(X_j^{\nu_j})\)
exists, is independent of the chosen pure-state
sequence converging to \(\xi\), and will be denoted
\(\sigma_\xi(X)\).

Each approximating tuple is a coordinate, hence is
continuous in \(\xi\). The same estimate controls its
error uniformly over all \(\xi\in\mathsf X\), so the
function \(\xi\mapsto\sigma_\xi(X)\) is continuous.
It is real, multilinear and symmetric, and takes
the value one at the empty tuple. Coordinate
functions separate \(\mathsf X\), since they are
its product coordinates. Integrating the uniform
approximation also extends
\eqref{eq:component-barycentre} to every tuple.
These statements hold at every point of the support;
no tuple-dependent exceptional set is introduced.

Here is the additional step needed for products of coordinates.
The lattice estimate \eqref{eq:lattice-cluster}, proved below using
the shifted Fourier filter, is uniform in the spacing and the
state when the products have fixed physical supports. In a pure
state its plateau is the product of the two expectations. Thus,
for fixed tuples $A,B$, the difference between
$\rho_j(A_j^{\nu_j}\tau_vB_j^{\nu_j})$ and
$\rho_j(A_j^{\nu_j})\rho_j(B_j^{\nu_j})$ tends to zero as
$|v|\to\infty$, uniformly in $j$ and in the pure evaluating state.
The reflection sizes, rather than the divergent supremum norms,
give the uniform bound needed here.

Integrate this inequality against $D_{\nu_j}$. At each fixed large
displacement the reference correlation converges to
$S^{\mathrm{ref}}(A\cdot\tau_vB)$, while the integral of the product
of the two coordinates converges by weak convergence on $\mathsf K$.
The uniform error permits the two limits to be interchanged. This
proves \eqref{eq:component-plateaux} for $q=2$. Repeat with a product
of the first $q-1$ translated groups as $A$ and the last group as $B$.
Those earlier translations are kept fixed during the last distant-group
limit. Their containing cube, and hence the constant, may depend on
them, but remains independent of $j$. Induction proves the stated
iterated limits without requiring a uniform bound in all displacements
at once.

The same pure-state inequality passes directly to each chosen sequence
$\rho_j$ and gives \eqref{eq:cluster} for $\sigma_\xi$. The constants
are independent of $\xi$. Moment inequalities, reflection positivity
and the source bounds pass to these limits as closed finite conditions.
All coordinates therefore concern the same continuum family.
\end{proof}

The image measure is intrinsic to the chosen reference hierarchy.
To see this, do not fix the subsequence used for weak compactness.
Every convergent subsequence of $\widehat D_j$ has the mixed moments
in \eqref{eq:component-plateaux}: its proof uses the already fixed
limit $S^{\mathrm{ref}}$ and the state-uniform plateau error. Thus,
for every coordinate monomial $u_{m_1}\cdots u_{m_q}$, all weak limit
points have the same integral. Finite linear combinations of such
monomials form an algebra containing the constants and separating
points of the compact product $\mathsf K$. They are uniformly dense
in $C(\mathsf K)$ by Stone--Weierstrass. Two probability measures
with these same monomial moments therefore agree on every continuous
function and are equal.

It follows that $\widehat D_j$ converges along the entire given
sequence of reference states whose correlations converge to
$S^{\mathrm{ref}}$. This is uniqueness of the coordinate law for
that hierarchy. It is not uniqueness of the continuum hierarchy as
the bare data vary. Compactness supplies a limit point; the plateau
identities identify every such point and remove only the additional
subsequence in the coordinate-measure construction.

The description also does not depend on the initially chosen dense
family of tuples. Given two such families, use their union as a third
countable family. Equicontinuity determines the missing coordinates
continuously from either dense family. The two projection maps from
the common support are therefore continuous and injective, as well
as onto the respective supports: a compact image carrying the full
image measure must contain its support, and conversely every image
of a support point belongs to that support. A continuous bijection
from a compact space to a Hausdorff space is a homeomorphism. Under
these identifications, the measures and all tuple evaluations agree.
This is the compact moment-problem meaning of the component space;
it is stronger than an arbitrary barycentric representation.


For completeness the plateau identities identify these functionals with
translation components of the reference hierarchy, rather than just
with an arbitrary convex decomposition. In its reflection space let
$P_0$ be the projection onto $\ker H$. The spectral theorem gives
$e^{-tH}\to P_0$ strongly as $t\to\infty$. Taking a time plateau in
\eqref{eq:component-plateaux} gives
\[
 \langle[F],P_0[G]\rangle
       =\int_{\mathsf X}\overline{\sigma_\xi(F)}
                            \sigma_\xi(G)\,d\lambda(\xi).
\]
Hence $P_0[F]\mapsto(\xi\mapsto\sigma_\xi(F))$ extends isometrically
from the invariant subspace into $L^2(\lambda)$ and sends the vacuum
to the constant one. Every compact separated tuple can be translated
far enough in the positive time direction that all its entries have
positive-time support. Translation invariance of each $\sigma_\xi$
preserves its value. Thus all coordinate functions occur in the image
of the positive-time polynomial vectors. They separate points of
$\mathsf X$, contain the constants, and are closed under conjugation.
Their linear span has a uniform closure closed under multiplication:
for positive-time tuples $X,Y$, uniform clustering approximates
$\sigma_\xi(X)\sigma_\xi(Y)$ by
$\sigma_\xi(X\cdot\tau_{te_4}Y)$, uniformly in $\xi$, as $t\to\infty$.
For sufficiently large $t$ this concatenated tuple is separated and
still has positive-time support. Stone--Weierstrass
therefore makes that closure $C(\mathsf X)$, which is dense in
$L^2(\lambda)$. The isometry is consequently onto.

Each limiting pure-state functional clusters, so its own invariant
subspace is generated by its vacuum. This is the translation-ergodic
component property used in \cref{thm:main}. Selecting any point of
$\mathsf X$ supplies the required component. The representing pure
lattice states belong to the enlarged periodic Gibbs class; they need
not themselves be limits of cubic-torus measures at fixed spacing.

\subsection{The plateau estimate and reflection sizes}

The lattice estimate measures the approach to the invariant plateau
using reflected two-copy correlations. These reflection sizes stay
bounded for products of centred densities at fixed physical
separation, even when their pointwise supremum norms diverge as the
spacing tends to zero. This makes the rate useful for the continuum
limit.

Let $a$ be fixed, let $r\ge1$ be an integer, and let $A,B$ be bounded
bond observables supported in $[-ra,ra]^4$. For $x\in a\mathbb Z^4$
the required estimate is
\begin{align}\label{eq:lattice-cluster}
 &\left|\nu(\overline A\,\tau_xB)
 -\nu\!\left(\overline{E_\nu[A\mid\mathcal I]}
                         E_\nu[B\mid\mathcal I]\right)\right|
 \notag\\
 &\hspace{12mm}\le C\mathcal N_r(A)^{1/2}\mathcal N_r(B)^{1/2}
                         (ra/|x|)^2,
 \qquad |x|\ge Cra.
\end{align}
The constant $C$ is absolute. To define the reflection size
$\mathcal N_r(A)$, let $\theta_\mu$
reflect $x_\mu$ at zero, and let $\gamma$ range over the sixteen
coordinate sign changes. Put
\begin{equation}\label{eq:reflection-size}
 \mathcal N_r(A)=\max_{\mu,\gamma}
 \nu\!\left(\overline{\theta_\mu\gamma A}\,
                         \tau_{2rae_\mu}\gamma A\right).
\end{equation}
These are nonnegative reflected two-copy correlations. Their maximum
is at most $\|A\|_\infty^2$, but that weaker estimate will not be
used for density products in the continuum limit.

Choose a coordinate with $|x_\mu|$ maximal and use it as time.
After exchanging $A,B$ if necessary, assume $x_\mu>0$.
In the site-mirror Hilbert space put
\[
 \Phi_A=[\tau_{rae_\mu}\theta_\mu A],\qquad
 \Phi_B=[\tau_{rae_\mu}B],\qquad
 n=x_\mu/a-2r.
\]
Then the correlation is a pairing of these vectors with
$T_a^nU(x_\perp)$, and their squared norms are among the reflection
sizes. Write $P_0$ for the projection onto $\ker(1-T_a)$.
The lattice cone with constant $k_*>0$ puts zero-energy vectors at
zero momentum, so $P_0U(v)=P_0$. Reflection positivity and the
ergodic averages identify the zero-energy pairing with
\[
 \langle\Phi_A,P_0\Phi_B\rangle
   =\nu\!\left(\overline{E_\nu[A\mid\mathcal I]}
                              E_\nu[B\mid\mathcal I]\right).
\]
This is the plateau subtracted in \eqref{eq:lattice-cluster}.
It must be removed before estimating the decaying spectral part.

\subsection{The spectral estimate behind the inverse-square rate}

One may first split $A$ into its even and odd parts under
$\theta_\mu$. The estimate controls each part by the
reflection size of $A$. If $\theta_\mu A=\varepsilon_A A$, with
$\varepsilon_A\in\{1,-1\}$, the definition of $\Phi_A$ gives, for
an integer $m\ge1$ and $w\in\mathbb Z^4$ with $w_\mu=0$,
\[
 \nu(\overline A\,\tau_{a(2rme_\mu+w)}A)
 =\varepsilon_A\langle\Phi_A,
       T_a^{2r(m-1)}U(aw)\Phi_A\rangle.
\]
Thus time $2rma$ corresponds to the power $t^{m-1}$ in the
spectral measure of $\Phi_A$, where $t=e^{-2raE}$.
The filter must combine both time signs into the transform of a
nonnegative function of $T_a$, with this shift and the reflection
parity taken into account.
\subsubsection*{A filter for the shifted time powers.}
\phantomsection\label{hyp:spectral-filter}

The required coefficients can be chosen independently of the state,
spacing and support size. Use normalized Haar measure on the circle.
Fix two nonnegative even smooth functions $h_+,h_-$ of integral one,
supported respectively where
\[
 \tfrac12\le\cos\alpha\le\tfrac{\sqrt3}{2},
 \qquad
 -\tfrac{\sqrt3}{2}\le\cos\alpha\le-\tfrac12.
\]
Put $W_0=h_+-h_-$ and
$w_m=\int W_0(\alpha)e^{-im\alpha}\,d\alpha/(2\pi)$.
These coefficients are real and even, $w_0=0$, and
$\sum_m|w_m|<\infty$, by twice integrating by parts.
The symbol $W_0$ vanishes on a fixed neighbourhood of zero.

For $0<t<1$, the Poisson kernel
\[
 P_t(\alpha)=\frac{1-t^2}{1-2t\cos\alpha+t^2}
           =\sum_{m\in\mathbb Z}t^{|m|}e^{im\alpha}
\]
gives the shifted transfer weight
\[
 \Psi(t)=2\sum_{m\ge1}w_m t^{m-1}
        =t^{-1}\int W_0(\alpha)P_t(\alpha)
                                      \,\frac{d\alpha}{2\pi}.
\]
The series defines a continuous function on $[0,1]$, with the
$m=1$ term interpreted as a constant also at $t=0$.
To check its sign and size, let $D_t(\alpha)=1-2t\cos\alpha+t^2$.
Since both $h_\pm$ have integral one,
\[
 \Psi(t)=2(1-t^2)\iint
 \frac{h_+(\alpha)h_-(\beta)
                       (\cos\alpha-\cos\beta)}
      {D_t(\alpha)D_t(\beta)}
                    \,\frac{d\alpha\,d\beta}{(2\pi)^2}.
\]
On these supports the cosine difference is at least one, and
$D_t\le4$. Consequently
\[
 \Psi(t)\ge\tfrac18(1-t^2),\qquad
 \Psi(e^{-s})\ge\tfrac18\tanh s\quad(s\ge0).
\]
Continuity includes both endpoints in these inequalities. In
particular $\Psi\ge0$, and $\Psi(1)=W_0(0)=0$.

For a field of reflection parity $\varepsilon_A$, take
$b_m^{\mathrm{filt}}=\varepsilon_Aw_m$ for $m\ne0$.
Reflection invariance and $\theta_\mu A=\varepsilon_A A$ identify
the positive and negative old-time correlations at the same
transverse displacement. Combining them with the preceding
shifted-power identity gives exactly $\Psi(t)$, since
$\varepsilon_A^2=1$. The transverse symbol of the same sequence
is $W_{\varepsilon_A}:=\varepsilon_AW_0$, so it vanishes near zero, and its absolute
coefficient sum has one fixed bound for both parities. Thus one
sequence has both required spectral readings, including $t=0$.

Use the full joint spectral measure of $T_a^{2r}$ and transverse
translations on $\Phi_A$, with transfer variable $t\in[0,1]$ and
$q\in\mathbb T^3=(\R/2\pi\mathbb Z)^3$ measured per lattice step.
The endpoint $t=0$ is retained in the coefficient identity. The next
calculation shows that this filter gives a projected weighted density
bounded by $Cr^3\mathcal N_r(A)$. Only after that calculation can
zero transfer and the invariant endpoint $t=1$ be separated for the
energy integral.

The role of the filter is best seen before the final energy integral.
Keep $e_\mu$ as the original time direction, and let
$\bar\mu_A$ be the projection to $q\in\mathbb T^3$ of the weighted
spectral measure of the full bounded transfer spectrum, including
$t=0$. The filter represents its Fourier coefficients in the form
\begin{equation}\label{eq:expanded-filter-coefficients}
 \widehat{\bar\mu}_A(w)
 =\sum_{m\ne0}b_m^{\mathrm{filt}}\,
        \nu\!\left(\overline A\,
                    \tau_{a(2rm e_\mu+w)}A\right),
 \qquad \sum_{m\ne0}|b_m^{\mathrm{filt}}|\le C,
\end{equation}
with the appropriate parity convention. In this formula $w_\mu=0$
and both $m$ and $w$ are in lattice units. In the original
site-reflection space, the positive-time terms contain powers
$T_a^{2r(m-1)}$, not $T_a^{2rm}$. Thus the coefficients must produce
the specified $\Psi$ with that shift. Reflection combines the two
signs of $m$; neither an arbitrary summable sequence nor an arbitrary
nonnegative function of $T_a$ gives the required identity.

Absolute summability permits the series to be evaluated in each
reflection picture. For $|w|_\infty<4r$, the reflected norm bound on
each correlation gives
\[
 |\widehat{\bar\mu}_A(w)|\le C\mathcal N_r(A).
\]
There are at most $Cr^3$ such displacements. For the remaining
displacements choose $\lambda\ne\mu$ with
$|w_\lambda|=|w|_\infty$ and take $e_\lambda$ as time, reversing
its sign if necessary. The support of $A$ still lies in a cube of
radius $r$ in this picture, and its reflection vectors have squared
norms at most $\mathcal N_r(A)$. Their new time separation is
\[
 n_\lambda=|w_\lambda|-2r\ge\tfrac12|w|_\infty.
\]
All translations in the old time direction are now spatial
unitaries. Therefore the sum over $m$ in
\eqref{eq:expanded-filter-coefficients} becomes the bounded spectral
multiplier
\[
 W_{\varepsilon_A}(2rQ_\mu)=\sum_{m\ne0}b_m^{\mathrm{filt}} e^{2irmQ_\mu},
 \qquad \|W_{\varepsilon_A}(2rQ_\mu)\|\le\sum_m|b_m^{\mathrm{filt}}|,
\]
where $Q_\mu$ is the old-direction momentum measured per lattice
step. This is a different spectral reading of the same scalar
correlation series, not a comparison of the two Hilbert spaces by
an assumed rotation.

The required symbol vanishes on a fixed neighbourhood of zero on
the circle. Consequently its support after the substitution $2rq$
has $|q|\ge c/r$ in the principal momentum interval. There can be
other zero intervals because of periodicity; these only remove more
spectral support. On the remaining support, the cone in the new
picture gives $aE_\lambda\ge c/r$. All other transverse
translations are unitary and commute with this joint spectral
multiplier. Hence
\[
 \|T_{\lambda,a}^{n_\lambda}
                   W_{\varepsilon_A}(2rQ_\mu)\|
 \le C e^{-cn_\lambda/r}.
\]
Cauchy--Schwarz for the two reflection vectors now gives
\begin{equation}\label{eq:expanded-filter-decay}
 |\widehat{\bar\mu}_A(w)|
       \le C\mathcal N_r(A)e^{-c|w|_\infty/r}
       \qquad(|w|_\infty\ge4r).
\end{equation}
The choice of $\lambda$ can be any of the three directions
transverse to the original time. This is where all four reflection
pictures, rather than one energy--momentum inequality, enter the
density estimate.

The number of lattice displacements in the shell
$j\le|w|_\infty<j+1$ is at most $C(j+1)^2$. Combining the near
count with \eqref{eq:expanded-filter-decay} yields
\[
 \sum_{w\in\mathbb Z^3}|\widehat{\bar\mu}_A(w)|
 \le C\mathcal N_r(A)
       \left[r^3+\sum_{j\ge4r}(j+1)^2e^{-cj/r}\right]
 \le Cr^3\mathcal N_r(A).
\]
Thus these coefficients have an absolutely convergent Fourier
series. The continuous function represented by that series has the
same Fourier coefficients as the finite measure $\bar\mu_A$.
Equality on trigonometric polynomials, followed by uniform density,
identifies the measure with that density. In particular
\[
 \left\|\frac{d\bar\mu_A}{dq}\right\|_\infty
                         \le Cr^3\mathcal N_r(A).
\]

\begin{proposition}[Decay to the invariant plateau]\label{prop:plateau-decay}
The coefficient construction above and the reflection calculation give \eqref{eq:lattice-cluster}, with a constant
independent of the state and lattice spacing.
\end{proposition}

\begin{proof}
Use the time direction and vectors $\Phi_A,\Phi_B$ chosen above.
First suppose $A$ has definite reflection parity. Let $\varsigma_A$
be the full joint spectral measure of $T_a^{2r}$ and transverse
translations on $\Phi_A$. Its transfer variable is $t\in[0,1]$;
momentum $q\in[-\pi,\pi]^3$ is measured per lattice step. Put
\[
 d\mu_A(t,q)=\Psi(t)\,d\varsigma_A(t,q),
 \qquad \bar\mu_A=\hbox{its projection to }q.
\]
The Fourier summation above gives
$ d\bar\mu_A/dq\le Cr^3\mathcal N_r(A)$ on this full spectrum.

The site transfer operator can have a kernel. This contributes at
$t=0$, and the shifted $m=1$ term in the filter identity can see it.
For the decay estimate, however, $n\ge1$, so $T_a^n$ annihilates
this subspace. The plateau is the separate endpoint $t=1$, removed
by $1-P_0$. We may therefore restrict the remaining integral to
$0<t<1$ and set
\[
 E=-\frac{1}{2ra}\log t>0.
\]
The projected weighted measure on this restriction is a positive
submeasure of $\bar\mu_A$ and inherits its density bound. On this
finite-energy subspace the lattice cone gives $aE\ge c|q|$ and
\[
 \Psi(t)=\Psi(e^{-2raE})\ge c\tanh(2raE).
\]
On positive energies the weight is strictly positive. For $q\ne0$
the cone consequently bounds the integrand after removal of the weight:
\[
 \frac{e^{-naE}}{\Psi(e^{-2raE})}
       \le C\frac{e^{-cn|q|}}{\tanh(cr|q|)}.
\]
The projected weighted measure is absolutely continuous, so the point
$q=0$ contributes no weighted mass. Strict positivity of the weight at
positive energy then excludes positive-energy mass over that point as well. The unweighted
zero-energy atom has already been removed by $1-P_0$.
Integration gives
\begin{align}\label{eq:spectral-cluster}
 \langle\Phi_A,T_a^n(1-P_0)\Phi_A\rangle
 &\le Cr^3\mathcal N_r(A)
       \int_{\mathbb T^3}\frac{e^{-cn|q|}}{\tanh(cr|q|)}\,dq
 \notag\\
 &\le C\mathcal N_r(A)
                 \big((r/n)^2+(r/n)^3\big).
\end{align}
Here is the last integration explicitly. For $u>0$,
$1/\tanh u\le C(1+u^{-1})$. Enlarge the momentum cube to
$\R^3$ for the resulting nonnegative majorant and use radial
coordinates. Up to absolute constants the two integrals are
\[
 \int_0^\infty s^2e^{-cns}\,ds=Cn^{-3},\qquad
 \frac1r\int_0^\infty se^{-cns}\,ds=Cr^{-1}n^{-2}.
\]
Multiplication by $r^3$ gives exactly the two powers in
\eqref{eq:spectral-cluster}. In the range $n\ge Cr$ the cubic term
is bounded by a constant times the quadratic term.

For general $A$, split into its even and odd reflection parts. The
reflection seminorms of each part are at most those entering
$\mathcal N_r(A)$: use the triangle inequality for the corresponding
reflection norm and the coordinate sign changes in its definition.
Applying the estimate to the two parts and then the triangle
inequality changes only an absolute constant. The same applies to $B$.
Since $T_a$ and the transverse translations commute, Cauchy--Schwarz
for the positive operator $T_a^n(1-P_0)$ gives
\[
 |\langle\Phi_A,T_a^nU(x_\perp)(1-P_0)\Phi_B\rangle|
 \le C\mathcal N_r(A)^{1/2}\mathcal N_r(B)^{1/2}(r/n)^2.
\]
The removed projection has precisely the invariant plateau identified
above. When $|x|\ge Cra$, choosing a largest coordinate gives
$n=x_\mu/a-2r\asymp|x|/a$. This proves
\eqref{eq:lattice-cluster}.
\end{proof}

The exponent two comes from the three transverse momentum dimensions
and the first-order zero in the lower filter bound. The proof uses a
positive lower cone constant at fixed spacing, and uses all three
transverse directions when the filter is read with a different time
axis. It makes no positive-gap assertion. The constant-one continuum
cone, which was essential for rotations, is not needed in this radial
integration.


\subsection{Closedness of pure states and the continuum passage}


At a fixed spacing, the uniform rate also makes the set of ergodic
periodic Gibbs states weakly closed. Suppose $\rho_i\to\rho$ with
each $\rho_i$ ergodic. For a continuous local $A$, apply
\eqref{eq:lattice-cluster} to a distant translate of $A$ in each
$\rho_i$. Its plateau is $|\rho_i(A)|^2$. Pass to the weak limit
at each fixed displacement, and then send that displacement to
infinity. The plateau formula for $\rho$ gives
\[
 \rho(|E_\rho[A\mid\mathcal I]|^2)-|\rho(A)|^2=0.
\]
For a countable dense family of local $A$, this makes the invariant
conditional state constant almost surely; thus $\rho$ is ergodic.
This is why support points of $D_\nu$ can be selected as pure
states, rather than merely almost-everywhere representatives.

For the continuum, fix separated tuples
\(X=(f_1,\ldots,f_n)\) and \(Y=(h_1,\ldots,h_m)\).
The required reflection-size bound must be uniform
in the selected pure state, not just finite for
each of its limiting functionals.

\emph{The reflected support geometry.}
Choose one \(R>0\) so that distinct supports within either
\(X\) or \(Y\) have distance at least \(6R\), ignoring
an internal separation condition for a one-entry tuple.
Put \(d=2R\) and use this radius in both partitions of
\cref{lem:fixed-support-bounds}. Thus the containing balls
of pieces from distinct entries of either tuple are
separated by \(d\). For \(X\), write
\[
 f_i=\sum_\alpha f_{i,\alpha},\qquad
 X^\alpha=(f_{1,\alpha_1},\ldots,f_{n,\alpha_n}),\qquad
 L_X=\sum_{\alpha_1,\ldots,\alpha_n}
                         \prod_i\|f_{i,\alpha_i}\|_\sharp.
\]
Then \(L_X=\prod_i\sum_\alpha\|f_{i,\alpha}\|_\sharp
\le(\prod_iK_i)\prod_i\|f_i\|_\sharp\).
Make the corresponding choices for \(Y\).
Use one \(L\)-adic \(\varpi>0\) large enough that
all these balls lie in
\((-\varpi+d/2,\varpi-d/2)^4\).
These choices depend on the original supports alone.

For every coordinate direction \(\mu\), sign change
\(\gamma\), and two piece tuples \(X^\alpha,X^\beta\),
consider
\[
 \theta_\mu\gamma X^\alpha\cdot
                        \tau_{2\varpi e_\mu}\gamma X^\beta.
\]
Within each half the balls have separation \(d\).
Between the two halves, the \(\mu\)-coordinates
of the first half are at most \(\varpi-d/2\),
and those of the second are at least
\(\varpi+d/2\). Thus the two halves also have
separation \(d\). All balls are contained in
\((-3\varpi+d/2,3\varpi-d/2)^4\).
The whole \(2n\)-entry tuple belongs to one
small-ball support class. This checks the
per-ball estimate for the reflected copies,
even when a whole original support cannot
fit into a small separated ball.

For all sufficiently large \(j\),
\(r_j=\varpi/a_j\) is an integer and
\(a_j\le d/2\). Put
\(A_j=X_j^{\nu_j}\). Multilinearity of the
centred field gives
\[
 A_j=\sum_\alpha (X^\alpha)_j^{\nu_j}.
\]
The reference state is invariant under the
lattice symmetries, so the same centring
is used in the reflected copies.
The state-uniform moment estimate on the
preceding small-ball class gives
\begin{equation}\label{eq:uniform-reflected-size}
 \begin{split}
 \mathcal N_{r_j}(A_j)
 &\le c_{2n}(d;R)
       \sum_{\alpha,\beta}
       \prod_i\|f_{i,\alpha_i}\|_\sharp
       \prod_i\|f_{i,\beta_i}\|_\sharp\\
 &=c_{2n}(d;R)L_X^2=:C_X^2.
 \end{split}
\end{equation}
The bound holds for every evaluating
\(\rho_j\in\mathcal G^{\mathrm{per}}_{K_j,\beta_j}\),
with the reflection size taken in that state.
The same argument gives
\(\mathcal N_{r_j}(Y_j^{\nu_j})\le C_Y^2\).
The constants can depend on the fixed reference
coupling and support geometry, but not on \(j\),
the evaluating pure phase, \(\xi\), or the
chosen volume-first limit. Bounds by
\(\|A_j\|_\infty^2\) would instead diverge.

\emph{Continuity and clustering on arbitrary supports.}
For any fixed disjoint support neighbourhoods,
\eqref{eq:fixed-support-continuity} passes to
every component sequence. It proves the same
local continuity estimate for each \(\sigma_\xi\),
uniformly in \(\xi\), including for the tuples
with reflected copies just used.
It also supplies the fixed-support local bounds
needed by the Hilbert-space and translation
construction; small-ball support restrictions
are not imposed on its polynomial vectors.

Apply \eqref{eq:lattice-cluster} in a pure
evaluating state. Its invariant conditional
expectations are constants, so the plateau
is the product of its two moments.
Equations \eqref{eq:lattice-cluster} and
\eqref{eq:uniform-reflected-size} give, at
lattice displacements \(v_j\),
\[
 \left|\rho_j(X_j^{\nu_j}\tau_{v_j}Y_j^{\nu_j})
       -\rho_j(X_j^{\nu_j})\rho_j(Y_j^{\nu_j})\right|
       \le C C_XC_Y\varpi^2|v_j|^{-2},
 \qquad |v_j|\ge C\varpi .
\]
For fixed sufficiently large \(v\), take
\(v_j\in a_j\mathbb Z^4\) tending to \(v\).
Their tuples retain common disjoint support
neighbourhoods. Translation continuity of
\(C_c^1\) functions in the sharp norm and
\eqref{eq:fixed-support-continuity} permit
these moving lattice representatives to pass
to the same limit as the fixed continuum tests.
This yields \eqref{eq:cluster} with
\[
 r_{X,Y}=C\varpi,\qquad
 C_{X,Y}=C C_XC_Y\varpi^2,
\]
after enlarging the fixed constants if necessary.
Both are uniform over the component space
and the volume-first limits at the chosen
reference coupling.

Equivalently, one may first apply the bound
to the finitely many small-ball pairs and
sum their cluster errors. Indeed, multilinearity
gives exactly
\[
 \begin{split}
 &\sigma_\xi(X\cdot\tau_vY)
                         -\sigma_\xi(X)\sigma_\xi(Y)\\
 &\qquad=\sum_{\alpha,\beta}
 \left[\sigma_\xi(X^\alpha\cdot\tau_vY^\beta)
             -\sigma_\xi(X^\alpha)\sigma_\xi(Y^\beta)\right].
 \end{split}
\]
The partitions and the common containing
cube are fixed once for \(X,Y\), before
varying the state or taking a large displacement.
Thus the sum of the class constants is finite
and has precisely the uniformity needed for
the component plateau identities.

The connected two-point bound passes
independently through the cumulants of
\(\rho_j\). Its positivity cannot be
explained solely by variation of the
component mean.


\subsection{The component mean and the invariant plateau}

The subtraction in \eqref{eq:phase-centring} is at the reference
state, not at the selected pure state. Choose a fixed smooth
integral-one test $f_0$. Translation invariance and the first-mean
transfer give a bounded real number
\[
 b_\xi=\sigma_\xi((f_0)),\qquad
 \sigma_\xi((f))=b_\xi\int f.
\]
It is continuous on $\mathsf X$, since tuple evaluations are
continuous. The common reference mean is zero, so the barycentre
identity gives $\int b_\xi\,d\lambda=0$. On the lattice its value
is represented by
\[
 b_\xi=\lim_j6a_j^{-4}
                  \big(\rho_j(a_p)-\nu_j(a_p)\big).
\]
There are six unoriented plaquette planes in four dimensions. Their
expectation differences are equal by the lattice symmetries; the
Riemann sum of the fixed test then gives the displayed factor.
Neither uncentred expectation has to converge separately.

It is useful to distinguish two different second moments. For
separated $f,h$, define the component's connected covariance by
\[
 C_{2,\xi}(f,h)=\sigma_\xi((f,h))
                      -b_\xi^2\left(\int f\right)\left(\int h\right).
\]
Since the reference first moment is zero, its ordinary and connected
second moments agree. Averaging the component identity gives
\begin{equation}\label{eq:expanded-phase-variance}
 S^{\mathrm{ref}}_2(f,h)
 =\int_{\mathsf X}C_{2,\xi}(f,h)\,d\lambda(\xi)
  +\kappa\left(\int f\right)\left(\int h\right),
 \qquad \kappa=\int_{\mathsf X}b_\xi^2\,d\lambda(\xi).
\end{equation}
The first term contains connected component covariances. Using
ordinary component second moments there would count their means
twice. In particular the barycentre itself remains
$S^{\mathrm{ref}}_2(f,h)=\int\sigma_\xi((f,h))\,d\lambda$.

The uniform component clustering bound permits the large-translation
limit to pass through this integral. It gives
\[
 \lim_{|v|\to\infty}S^{\mathrm{ref}}_2(f,\tau_vh)
       =\kappa\left(\int f\right)\left(\int h\right).
\]
Thus variation of the component mean can create a nonzero reference
plateau even though every component clusters. Conversely the
positive short-distance covariance in \cref{thm:main} is the
connected covariance of each selected phase, obtained from the
state-uniform two-point comparison. It is not produced solely by
this variation of means. Equality of the means alone would not show
that the components agree on higher tuples.

\subsubsection*{Time clustering and the vacuum.}
In the reflection space of one component the spectral theorem gives
$e^{-tH}\to P_0$ strongly, where $P_0$ projects onto $\ker H$.
For positive-time polynomial vectors,
\[
 \lim_{t\to\infty}\langle[F],e^{-tH}[G]\rangle
       =\langle[F],P_0[G]\rangle.
\]
Time clustering identifies the left side with
$\overline{\sigma_\xi(F)}\sigma_\xi(G)$.
Since $\langle\Omega,[G]\rangle=\sigma_\xi(G)$, these are exactly
the matrix elements of the rank-one projection
$|\Omega\rangle\langle\Omega|$. The polynomial vectors are dense
and both operators are bounded; consequently
$P_0=|\Omega\rangle\langle\Omega|$. Conversely, if this projection
identity holds, the same strong limit gives time clustering on all
such vectors. This equivalence uses no rate. The rate in
\eqref{eq:lattice-cluster} is needed earlier to preserve clustering
when the cutoff and the evaluating pure state both vary.

There is an explicit reason that the cone and all the mirror
positivities cannot replace that uniform estimate. Start with a
clustering component functional $\sigma$ and shift its field by a
constant $c\ne0$. On a tuple $X=(f_1,\ldots,f_k)$ set
\begin{equation}\label{eq:expanded-constant-shift}
 \sigma_c(X)=\sum_{T\subset\{1,\ldots,k\}}
       c^{|T|}\prod_{i\in T}\int f_i\;
       \sigma((f_i)_{i\notin T}),
 \qquad \widetilde\sigma=\tfrac12(\sigma+\sigma_c).
\end{equation}
The shift commutes with translations, rotations and reflection: a
test is replaced by the corresponding field insertion plus the real
constant $c\int f$. Reflection positivity follows by applying the
original positive form to this shifted polynomial. It holds at each
of the coordinate and diagonal mirrors. Averaging the two positive
forms preserves their positivity. The bounded five-point inequality
used for the cone is linear in the functional for every fixed test
polynomial, so it is preserved by this averaging as well.

The reconstruction used for this example applies to the
mixture itself. To check this without assuming it is
one chosen lattice approximation, first realize
\(\sigma_c\) using the same approximating states
with their centring constants decreased by
\(c\,a_j^4/6\). The added lattice constant is
\[
 c_j(f)=\frac{c\,a_j^4}{6}\sum_p f(x(p))
                              \longrightarrow c\int f.
\]
The sum ranges over the six plaquette families.
Multiplying out their field products gives exactly
\eqref{eq:expanded-constant-shift}. Thus \(\sigma_c\)
has the same kind of local continuity bounds,
symmetry and reflection form as \(\sigma\).

For each fixed polynomial the contraction inequality
\[
 \sigma(\theta\overline Z\cdot\tau_{2te_4}Z)
                         \le\sigma(\theta\overline Z\cdot Z)
\]
and the five-point inequality are linear
inequalities in the functional. They pass
from the two shifted functionals to their
equal mixture, as do their invariances
and reflection positivity. The local
continuity bounds pass with the larger
constant. Symmetry, multilinearity and
normalization at the empty tuple pass too.

These properties construct the reflection
Hilbert space of \(\widetilde\sigma\), its
strongly continuous contraction semigroup
and spatial unitary translations by the
same argument used above. The five-point
inequality can then be applied to this
reconstruction to obtain the coordinate
cones. This uses no ergodicity of the
mixture. The equivalence between time
clustering and a one-dimensional vacuum
space uses only the reconstruction and
the strong limit \(e^{-tH}\to P_0\), so
it applies to \(\widetilde\sigma\) as well.
The nonzero plateau below therefore does
prove a second vacuum direction for this
mixture, while making no such claim for
the selected clustering components.

Write $\sigma((f))=b\int f$. At large time separation the two
components give the mean products $b^2\int f\int h$ and
$(b+c)^2\int f\int h$. The mean of their equal mixture is
$(b+c/2)\int f$. Hence
\begin{align*}
 &\lim_{t\to\infty}
 \big[\widetilde\sigma((f,\tau_{te_4}h))
           -\widetilde\sigma((f))\widetilde\sigma((h))\big]\\
 &\qquad=
 \left[\frac{b^2+(b+c)^2}{2}
                       -\left(b+\frac c2\right)^2\right]
                         \left(\int f\right)\left(\int h\right)
 =\frac{c^2}{4}\left(\int f\right)\left(\int h\right).
\end{align*}
For a positive-time test of nonzero integral this is a nonzero
invariant covariance. Its reflection vector has an invariant
projection not proportional to the vacuum. Thus the mixture has
at least two invariant vectors despite satisfying the mirror
positivities and the cone. Selecting pure approximants and carrying
their clustering through a state-uniform rate is the step that
excludes this phenomenon for the component in \cref{thm:main}.

\subsubsection*{The round spectral cone after rotations.}
Once rotation invariance of the component has been proved, its
spatial rotations act on the reflection space by
\[
 U(R)[F]=[(R\oplus1)F],\qquad R\in\mathrm O(3).
\]
The reflected form is unchanged because these rotations commute
with time reflection. The action is therefore isometric on the
dense polynomial vectors and extends to a unitary representation.
Continuity in the tests gives strong continuity. It fixes the vacuum,
commutes with $H$, and rotates spatial translations:
\[
 U(R)e^{iv\cdot P}U(R)^{-1}=e^{i(Rv)\cdot P}.
\]
The joint energy--momentum spectrum is consequently invariant under
spatial rotations. Applying this to the already proved coordinate
inequality $H\ge|P_1|$ gives $H\ge|u\cdot P|$ for every unit vector
$u\in\R^3$. On the joint spectrum,
$\sup_{|u|=1}|u\cdot p|=|p|$, and hence its support lies in
$\{E\ge|p|\}$. This round cone is a consequence of the established
rotation invariance. The earlier rotation argument used the
constant-one inequalities in the individual coordinate planes and
does not use this conclusion as an input.


\section{Extension to flat Schwartz tests}\label{sec:extension}

We pass from separated tuples to the flat Schwartz space in
\eqref{eq:flat-space}. The state-uniform cumulant estimates control
one derivative per entry and a fixed power of the support separation.
Flatness supplies the derivatives needed to sum near coincidences.
We keep the factorial and derivative counts, since the growth clause
of \cref{thm:main} requires more than qualitative distributional
extension. The construction leaves contact terms unspecified.

\subsection{Distributions away from coincidences}

\begin{proposition}[The local distribution]\label{prop:configuration-distribution}
Let $T$ be a multilinear functional on smooth separated tuples. Suppose
it satisfies a product sharp-norm bound whenever the entries have
supports in prescribed mutually separated compact sets. Then it defines
a unique distribution on the open configuration space
\[
 (\R^4)^n_{\ne}=\{(x_1,\ldots,x_n):x_i\ne x_j\text{ if }i\ne j\}.
\]
On each separated product region its order is at most $Cn$. It agrees
with $T$ on tensor products supported in that region.
\end{proposition}

\begin{proof}
Let $F$ have compact support in $\prod_iO_i$, where the $O_i$ are
mutually separated open sets. Choose $\psi_i\in C_c^\infty(O_i)$
equal to one on the corresponding support projections. Put the supports
in the interior of a cube of side $P$ and extend $F$ periodically,
using that it is zero near the cube boundary. Its Fourier series gives
\[
 F=\sum_{\xi\in(2\pi/P)\mathbb Z^{4n}}\widehat F(\xi)
             \bigotimes_i[\psi_i(x_i)e^{i\xi_i\cdot x_i}].
\]
The sharp norm of the $i$th factor is at most $C_i(1+|\xi_i|)$.
Consequently the summand after applying $T$ is bounded by a constant
times $|\widehat F(\xi)|\prod_i(1+|\xi_i|)$.
Integration by parts bounds the Fourier coefficient by
$C_k(1+|\xi|)^{-k}\|F\|_{C^k}$. Taking $k>5n$, for example an
even integer at most $6n+2$, makes the sum absolutely convergent:
the $n$ sharp derivatives consume $n$ powers, and the frequency
lattice has dimension $4n$.

Apply the convergent series as the definition of $T(F)$ on the product
region. It is a distribution of the stated finite order. On a tensor,
Fourier partial sums converge in the sharp norm of each factor; the
multilinear bound therefore proves agreement with the original $T$.
Finite sums of tensors are dense in the smooth topology on such a
product region. Two local definitions which agree on tensors consequently
agree on their overlap. A partition of unity on the configuration space
patches them to a distribution, and the same density proves uniqueness.
\end{proof}

We now keep track of the constants needed near the diagonals. For tests
supported in unit balls, \eqref{eq:local-growth} has the form
\[
 |C_n(f_1,\ldots,f_n)|
       \le C(g)^n(n!)^{p_2}\max\{1,d^{-q_0}\}
                                    \prod_i\|f_i\|_\sharp,
 \qquad n\ge2,
\]
where $q_0$ does not depend on $n$. For component-centred moments,
the moment--cumulant formula sums over partitions with block sizes at
least two. There are at most $n!$ such partitions and at most $n/2$
blocks. Also $\prod_{B\in\pi}|B|!\le n!$. Hence
\[
 |S_n^0(f_1,\ldots,f_n)|
   \le C(g)^n(n!)^{p_2+1}
           \max\{1,d^{-q_0n/2}\}\prod_i\|f_i\|_\sharp.
\]
Rounding the separation exponent upward makes it a nonnegative integer
$a_n\le Cn$. The separation loss is therefore linear in the order
and the growth remains a fixed power of $n!$.

\subsection{A quantitative extension on flat tests}

\begin{proposition}[Flat-test extension]\label{prop:flat-test-extension}
Suppose $n\ge2$ and the separated functional $T$ has the bound
\[
 |T(f_1,\ldots,f_n)|
   \le c_T\max\{1,d^{-a}\}\prod_i\|f_i\|_\sharp
\]
for supports in unit balls at mutual distance at least $d$, where
$a\ge0$ is an integer. Its configuration-space distribution has a
unique continuous extension to $\cS_{\ne}(\R^{4n})$, with
\begin{equation}\label{eq:extension}
 |\widetilde T(\varphi)|
       \le C^{n+a}(n!)^Cc_T\,|\varphi|_{Cn+a}.
\end{equation}
The absolute constants can depend on the fixed sharp-norm convention.
\end{proposition}

\begin{proof}
Write $r(x)=\min_{i<j}|x_i-x_j|$. Use dyadic scales
$\delta=1,1/2,1/4,\ldots$, with the first region covering $r\ge1$
and the subsequent regions covering $r\asymp\delta$. One can make
the smooth shell cutoffs directly from a fixed monotone radial function
which vanishes for radius at most one and is one for radius at least
two: take products over all pairs and subtract the products at
successive scales. Their sum is one off the diagonals. At the inner
shell of scale $\delta$, all pair distances are at least $c\delta$
and one pair distance is at most $C\delta$.

Refine each shell by products of cubes of side $c\delta$, with a fixed
small $c$. The slightly enlarged cubes in any product meeting the shell
still have mutual distance at least $c'\delta$. Their diameters are
less than one, so the assumed tuple bound applies there. Fixed smooth
cutoffs in the scaled cubes permit the Fourier construction in
\cref{prop:configuration-distribution}. Choose the one-variable
cutoffs of Gevrey order two, with derivative bounds $C^k(k!)^2$.
Only $Cn$ derivatives of them occur in the Fourier estimate. Products
over the $O(n^2)$ pairs and the $n$ cubes consequently cost at most
$C^n(n!)^C$: differentiating $k\le Cn$ times introduces at most
$(Cn^2)^k$ choices, again a fixed factorial power. No derivative of
order $a$ is applied to these cutoffs.

Here is the estimate on one scaled product cube. Rescale it to a fixed
product cube and apply the Fourier series just described. The
separation factor is at most $C^a\delta^{-a}$. The first derivative
in each sharp norm costs at most $C\delta^{-1}$, giving
$C^n\delta^{-n}$. Thus for its smooth cutoff $\upsilon$ the local
distribution obeys
\[
 |T(\upsilon\varphi)|
  \le C^{n+a}(n!)^Cc_T\delta^{-a-n}
       \max_{|\gamma|\le Cn}
                     \delta^{|\gamma|}\sup_{\mathrm{cube}}
                                        |\partial^\gamma\varphi|.
\]
The frequency sum is over a fixed scaled cube. For example, using
$(1-\Delta)^{3n+1}$ in the Fourier coefficient estimate costs
$(Cn)^{Cn}$ in its coefficients, another fixed power of $n!$.
This accounts for the dimension dependence in the displayed bound.

Flatness supplies as many powers of $\delta$ as needed. On an inner
shell choose a pair at distance at most $C\delta$ and move one of
its points to the other. Every derivative of $\varphi$ vanishes at
the resulting point of that diagonal. Taylor's formula, with $s$
further derivatives applied to $\varphi$ itself, gives
\[
 |\partial^\gamma\varphi(x)|
   \le C^{s+Cn}\delta^s
          |\varphi|_{Cn+s}(1+|x|^2)^{-(Cn+s)/2},
 \qquad |\gamma|\le Cn,
\]
after increasing the absolute derivative constant. Along the Taylor
segment only one four-dimensional coordinate moves a distance
$O(\delta)\le O(1)$, so the Schwartz weights are comparable with
a factor $C^{Cn+s}$. The Taylor factorial in the denominator may be
discarded. In particular this step costs $C^s$, rather than $s!$.
For the outer shell take $\delta=1$ and use the direct Schwartz
bound, which has the same form without a small-scale gain.

We sum over cube positions using the Schwartz weight. A product grid
in $4n$ dimensions with mesh $c\delta$ satisfies, for an exponent
larger than $4n$ by a fixed multiple of $n$,
\[
 \sum_z(1+|x_z|^2)^{-(Cn+s)/2}
                 \le C^{n+s}\delta^{-4n}.
\]
This follows by comparison with the integral over the product cubes;
the integral of the radial weight is bounded by $C^n$ when its
exponent is at least $6n$. We have used the coarser count
$\delta^{-4n}$; the smaller count obtained after separating a common
translation is not needed. The complete scale bound is therefore
\[
 \sum_z|T(\upsilon_{\delta,z}\varphi)|
   \le C^{n+a+s}(n!)^Cc_T
                 \delta^{-a-5n+s}|\varphi|_{Cn+s}.
\]
Choose $s=a+6n+1$. A positive power of $\delta$ remains, and the
dyadic sum is bounded by an absolute constant. The derivative index
$Cn+s$ is $a+C'n$ after renaming the absolute constant. The factor
$C^s$ is absorbed by $C^{n+a}$, proving \eqref{eq:extension}.
The argument has differentiated the geometric cutoffs only $O(n)$
times; the further $a$ derivatives belong to the flat test. Thus no
additional factorial in $a$ has been concealed in the estimate.

The absolutely convergent sum defines the extension. The density
statement proved below shows that it is uniquely determined by its
values on the separated tensors, and hence by the original functional.
\end{proof}

\begin{lemma}[Density on the flat test spaces]\label{lem:flat-test-density}
Finite sums of smooth compactly supported separated tensors are dense
in $\cS_{\ne}(\R^{4n})$ for its Schwartz topology. For tests flat
outside the positive-time configuration space, the tensor entries can
all be chosen with compact support in $\{x_4>0\}$.
\end{lemma}

\begin{proof}
We describe the approximations in an arbitrary fixed seminorm
$|\cdot|_k$; repeating the choices for increasing $k$ gives a
sequence converging in the Schwartz topology. Let $\chi_R$ be a
smooth cutoff equal to one on the ball of radius $R$ and zero beyond
radius $2R$. Its derivatives satisfy
$|\partial^\alpha\chi_R|\le C_\alpha R^{-|\alpha|}$.
Leibniz's rule and the extra Schwartz decay give, for every fixed
$m>0$,
\[
 |(1-\chi_R)\varphi|_k\le C_{k,m}R^{-m}|\varphi|_{k+m}.
\]
This removes infinity first. The cutoff does not affect flatness.

Next let $A_\varepsilon$ be a smooth cutoff which is zero if some
pair distance is at most $\varepsilon$ and one if all pair distances
are at least $2\varepsilon$. It can be the product of the pair
cutoffs used in the dyadic decomposition. For fixed $n,k$ its
derivatives of order at most $k$ are bounded by
$C_{n,k}\varepsilon^{-k}$. On the support of
$1-A_\varepsilon$, at least one point lies within
$C\varepsilon$ of a pair diagonal. Flatness and Taylor's formula,
applied to any derivative of $\chi_R\varphi$, give a factor
$\varepsilon^s$ for any chosen $s$. Thus, with $s>k$,
\[
 |(1-A_\varepsilon)\chi_R\varphi|_k
             \le C_{n,k,s,R}\varepsilon^{s-k}
                                      |\chi_R\varphi|_{k+s}
             \longrightarrow0.
\]
The large derivatives of the cutoff have therefore been paid for by
flatness. There is no claim of a uniform bound for tests with nonzero
jets on a diagonal.

The remaining test has compact support a positive distance from the
diagonals. Cover its support by finitely many products of mutually
separated open cubes and choose a subordinate smooth partition of
unity. On each product region the Fourier construction of
\cref{prop:configuration-distribution} approximates the localized test
by finite sums of smooth tensors, with supports in the same separated
open sets. Choose the Fourier truncation large enough for the fixed
$C^k$ accuracy on this compact support. Summing the finitely many
pieces proves the first assertion.

For the positive-time assertion, after the cutoff at infinity also
cut off the region where any time coordinate is at most
$\varepsilon$. Smoothness and vanishing outside the open half-space
imply that all derivatives vanish on its boundary. The same Taylor
and Leibniz estimates make this cutoff error tend to zero in each
fixed seminorm. The separated product cover can now be chosen inside
the positive half-space, as can all Fourier tensor factors. These
operations prove density in the required positive-time flat space.
\end{proof}

Applying the proposition with $a=a_n\le Cn$ and the moment bound
above gives constants independent of $n$ such that
\[
 |S_n^0(\varphi)|\le A_{\mathrm c}B_{\mathrm c}^n(n!)^p
                                  |\varphi|_{kn}.
\]
The zero and first centred moments are $1$ and $0$, respectively.
This proves the required quantitative extension for a component
centred at its own mean.

\subsection{Restoring the component mean}

The selected phase retains the subtraction of
the reference state. Its first moment can be $b\int f$, with $|b|$
uniformly bounded by one-point transfer. Its higher cumulants are the
selected phase's cumulants. Consequently its own centred moments
$S_n^0$ need not be the reference mixture's centred moments
$S_n^{\mathrm c}$ in \eqref{eq:centred-moments}.

For $A\subset\{1,\ldots,n\}$ define
$\psi_A=\int\varphi\,dx_A$. Integration preserves Schwartz decay
and flatness on every diagonal of the remaining variables: every
derivative in those variables still vanishes on that diagonal before
the integration. The separated-tuple binomial identity extends as
\begin{equation}\label{eq:restore-mean}
 S_n(\varphi)
       =\sum_{A\subset\{1,\ldots,n\}}
                    b^{|A|}S_{n-|A|}^0(\psi_A),\qquad S_0^0=1.
\end{equation}

Increase $k$ so that $k>4$. If $y$ denotes the remaining variables,
then
\[
 (1+|x|^2)^{kn/2}
  \ge (1+|y|^2)^{k(n-|A|)/2}
                         \prod_{i\in A}(1+|x_i|^2)^{k/2}.
\]
Each factor on the right is bounded by the corresponding power of
$1+|x|^2$, and the exponents sum to $kn/2$. Differentiating under
the integral, and using this inequality, yields
\[
 |\psi_A|_{k(n-|A|)}
       \le C_{\int}^{|A|}|\varphi|_{kn},
 \qquad C_{\int}=\int_{\R^4}(1+|y|^2)^{-k/2}\,dy<\infty.
\]
For $A$ the whole index set this includes the integral of $\varphi$
against the constant zero-variable functional. Since
$(n-|A|)!\le n!$, \eqref{eq:restore-mean} and the binomial theorem
give the growth bound with
\[
 B=B_{\mathrm c}+|b|C_{\int}
\]
and an enlarged prefactor to include $S_0^0$. The constant mean changes
only the geometric growth factor. Its uniform bound therefore permits
the same four growth constants for all components under consideration.

\subsection{Passage of the Euclidean properties}

The seminorms have the useful tensor inequality
\[
 |\varphi\otimes\psi|_k\le|\varphi|_k|\psi|_k.
\]
Indeed $1+|x|^2+|y|^2\le(1+|x|^2)(1+|y|^2)$, and each derivative
of the tensor is the product of one derivative of each factor.
This inequality and the growth bound control every term of a finite
reflection-positive quadratic form under test-function approximation.

Take a smooth Schwartz test flat outside the positive-time configuration
space. It can first be cut off at infinity, then near the coincidence
sets, and then near the time boundary, all in its Schwartz topology.
The boundary cutoff is allowed because every derivative vanishes where
one time coordinate is zero, by smoothness of its zero continuation.
After these cutoffs the tensor approximation can be made within a
finite union of separated positive-time product regions.

For two such tests, reflection of the first followed by tensor product
with the second is flat on the cross diagonals as well as the original
diagonals. A negative-time point and a positive-time point cannot
coincide; at their common boundary all derivatives vanish. The
reflection-positive quadratic form on the tensor approximants is thus
defined entirely on separated tuples. It is nonnegative by the
previous sections. Every one of its finitely many terms converges by
the growth bound, so its limit is nonnegative.

More explicitly, for a sequence with degrees at most $N$ the
quadratic form contains only distributions $S_{m+n}$ with
$m+n\le2N$. Their growth bounds require derivatives of order at
most $2kN$. Apply \cref{lem:flat-test-density} to each test in a
Schwartz seminorm at least this large. The tensor inequality bounds
each difference of reflected products by the sum of one factor
difference times bounded factor seminorms. Every term of the finite
quadratic form therefore tends to its original value. Positivity on
the separated tensor approximants proves positivity on the test
sequence. No estimate uniform in an infinite number of degrees is
needed for this passage.

The identical density argument passes symmetry and the Euclidean
invariance established in \cref{sec:rotations} to the flat test space.
For each fixed Euclidean transformation the pullback is continuous in
the Schwartz topology and preserves flatness. The separated identities
therefore determine the identities of the extended distributions.
This gives the test-space assertions of \cref{thm:main} while leaving
their values on general tests at coincidences unspecified.


\subsection{Assembly of the construction}\label{sec:assembly}

We now put the proof outline in the order needed for
\cref{thm:main}. This also specifies which limits must be shared
by the different parts of the argument.

First carry out the renormalisation, history and source constructions
on the same finite prefixes. The scalar coefficients extracted from
the effective density determine the next inverse coupling before the
next large-field operation. Their bounded one-loop partial-sum error,
the small ordinary correction and the summable remainder correction
give the two-sided drift. For bare inverse couplings tending to
infinity, stop when the run first enters the fixed reference window.
The stopping indices $K_j$ tend to infinity, so $a_j=L^{-K_j}\to0$.
All endpoint couplings stay in one small-coupling range. This avoids
an exact endpoint tuning or an assertion that the chosen sequence is
unique.

At each such datum take the volume limit first. Compactness of the
bond configuration space gives a reference Wilson Gibbs state $\nu_j$.
The pressure and reflection arguments are designed to give the same
coordinate and diagonal reflection forms in this state and in every
Gibbs state used later. The common scalar source normalizer then
controls the centred correlations uniformly in these choices. Select
one subsequence on a countable dense family of separated tuples;
the per-ball equicontinuity bounds extend convergence to all separated
tuples. Call the resulting hierarchy $S^{\mathrm{ref}}$.

Next decompose each $\nu_j$ into translation-ergodic Gibbs states.
Georgii's decomposition theorem applies to the same Wilson
specification, as explained in \cref{sec:phases}. Map these states
into the compact space of correlation coordinates. The uniform
plateau estimate identifies mixed coordinate moments with iterated
limits in which groups of insertions are translated far apart.
Compact moment determinacy then identifies a support point of the
limiting coordinate law with a component of $S^{\mathrm{ref}}$.
It can be approached by pure lattice states $\rho_j$, evaluated with
the original subtraction at $\nu_j$. Passing to a further subsequence
preserves the already chosen reference hierarchy. Denote the
component by $\sigma$.

The purpose of the state-uniform estimates is visible here. A
reflection inequality or a two-point comparison known only for the
reference tori would not automatically hold in $\rho_j$. Bounds for
the whole admitted Gibbs class allow the same limits to carry
reflection positivity, reality, permutation symmetry, translations
and cubic symmetries to $\sigma$. The bounded five-point transfer
inequalities give the constant-one coordinate cones without passing
unbounded generators through the limit.

Apply the rotation argument to the component centred at its own
mean. The cell normalization and long-column limit give the
degree-independent insertion bound. The plate construction and
ordered products give the analytic kernels and their linear angular
growth. The ordered-tube differential equation gives global analytic
coefficients, and slab translations identify them with the original
functional at projection ties. The polygon zero-row induction gives
divisibility at every order. The weaker ultraviolet estimate from
source stability eliminates the high two-point harmonics;
diagonal reflection positivity removes the remaining two-point
harmonics. Convexity together with the polygon divisibility then
removes them at every order. Support-controlled approximation by
plates returns the rotation identity to ordinary separated tests.

The common first-derivative comparison bounds the component's first
moment. Translation invariance identifies it as $b\int f$, with
$b\in\R$ finite. Restoring this constant mean preserves the rotation
identities and the reflection form. The graded source bound, the
moment--cumulant formula and the diagonal-flat extension give
\eqref{eq:growth}, with one derivative order linear in $n$.
The positive-time approximation in this section passes reflection
positivity to the finite Schwartz sequences in \eqref{eq:rp}.
This is the route to clauses~\ref{it:invariance}--\ref{it:growth}
on the stated test spaces.

For clustering, apply the lattice plateau estimate to representatives
of two fixed tuples in $\rho_j$. Their physical support cube is fixed,
and their reflection sizes are uniformly bounded by the separated
moment estimates. Their plateaux are products because the states are
ergodic. The inverse-square error survives the cutoff limit and gives
\eqref{eq:cluster} for the selected component. No factorization of
the reference mixture is needed.

For the last clause, apply \cref{prop:reference-two-point} and
\cref{prop:witness-volume-transfer} to the fixed
bump pair. The positive leading term has size $g(d)^4d^{-8}$, with
$g(d)^4\asymp(\log(1/d))^{-2}$. The state-transfer error is at most
$Cd^2\|f_d\|_\infty\|h_d\|_\infty$; its ratio to the leading term
is $O(d^2(\log(1/d))^2)$, which tends to zero. Decreasing $d_0$
absorbs this error uniformly. Reference subtraction does not change
a second cumulant, and the finite mean product cancels in the
connected difference. This gives the lower and upper bounds in
\eqref{eq:nontrivial}.

\subsection{Parameters and compatibility}\label{sec:remaining-tasks}

The main construction cannot be assembled from unrelated estimates
that happen to have the same form. Its small-field, large-field and
source representations must be the representations of one measure.
The exact kernels fix the common measure identity. The analytic
representations must additionally be constructed on one compatible
array of domains; equality of the represented measures alone does not
identify those arrays.
At a finite prefix, all tests and stored-field operations use the same
geometry and real choices. A source changes only the weight on which
the frozen operations act. The marginal coefficient is extracted from
that common density, and the history contribution is included before
the next coupling is fixed. These are finite induction steps, not
conditions imposed on an eventual continuum limit.

After fixing the test, remainder and analytic exponents, choose $L$
and $M_1$, then the activity decay rate $\kappa$, and then the
larger interaction scale $M$ and interaction bound $E_0$, as in
\cref{sec:rg}. The further geometric enlargements and analytic
reserves must accommodate the layered inverses, localization domains,
shell preparation and reflected source arrays. The logarithmic test
exponents must both pay the
failed sharper tests and leave room for the generated fields; the
live-age rate is chosen strictly above eight so that it absorbs two
stored plaquette sums and their polynomial factors. The ceiling for
the running coupling is chosen last, after the fixed two-point shape
and source radii. Sending the cutoff to infinity does not change any
of these fixed choices.

Several different limit orders occur, and none can replace another.
The normalizer is selected along a volume sequence at fixed cutoff.
The column length then tends to infinity at that fixed cutoff, so its
endpoint supremum disappears before it can grow with the number of
outer insertions. The common continuum subsequence is selected only
after the state-uniform estimates are available. In the four spectral
pictures the shifted filter identity uses the full bounded transfer
spectrum; the zero-transfer part is removed only by a later positive
time step, and the invariant plateau is removed separately. Finally,
the cluster limits used to identify the component law are iterated
with the earlier translations fixed.

These compatibility checks explain the role of the detailed estimates
throughout the outline. The analytic work still to complete includes
the uniform RG and history bounds on the actual evaluation domains
and the common complex one-step, generated-background and
complex-increment estimates used by the source and relative
two-point comparisons. The differentiated estimate
\eqref{eq:witness-source-slope} also remains to be established for
the full insertion, including the residual quadratic component,
on the required complex source domains.
Their local analytic verifications include
the seam right-inverse energy estimate, the scale-weighted auxiliary
system bounds, the retained-form coercivity comparison, and the
generated-gauge and stage estimates on their common domains. The local source propagation, energy
normalization, component majorants and free curvature comparison
have been developed above under those inputs. The history-window
argument also requires the interval-reading property of the actual
later conditional laws. Recording interval labels and separating
sharp thresholds do not by themselves establish that property.
The fixed-spacing reflection transfer and the angular and polygon
arguments have been given explicitly. Their continuum application
still depends on the uniform construction estimates supplying the
sharp cone, insertion bounds and source stability. The weaker
short-distance bound separates the rotation argument from the
relative two-point comparison. Completing these analytic estimates
and conditional-law verifications would turn this construction into a proof of the target
theorem.

\begin{thebibliography}{FILS78}
\bibitem[B84]{Balaban1984}
T.~Ba\l aban,
\emph{Propagators and renormalization transformations for lattice gauge
theories. I},
Comm. Math. Phys. \textbf{95} (1984), 17--40.
\href{https://doi.org/10.1007/BF01215753}{doi:10.1007/BF01215753}.

\bibitem[B85]{Balaban1985}
T.~Ba\l aban,
\emph{Propagators for lattice gauge theories in a background field},
Comm. Math. Phys. \textbf{99} (1985), 389--434.
\href{https://doi.org/10.1007/BF01240355}{doi:10.1007/BF01240355}.

\bibitem[B87]{Balaban1987}
T.~Ba\l aban,
\emph{Renormalization group approach to lattice gauge field theories.
I. Generation of effective actions in a small field approximation and
a coupling constant renormalization in four dimensions},
Comm. Math. Phys. \textbf{109} (1987), 249--301.
\href{https://doi.org/10.1007/BF01215223}{doi:10.1007/BF01215223}.

\bibitem[B88]{Balaban1988}
T.~Ba\l aban,
\emph{Convergent renormalization expansions for lattice gauge theories},
Comm. Math. Phys. \textbf{119} (1988), 243--285.
\href{https://doi.org/10.1007/BF01217741}{doi:10.1007/BF01217741}.

\bibitem[B89I]{Balaban1989I}
T.~Ba\l aban,
\emph{Large field renormalization. I. The basic step of the $R$ operation},
Comm. Math. Phys. \textbf{122} (1989), 175--202.
\href{https://doi.org/10.1007/BF01257412}{doi:10.1007/BF01257412}.

\bibitem[B89]{Balaban1989}
T.~Ba\l aban,
\emph{Large field renormalization. II. Localization, exponentiation,
and bounds for the $R$ operation},
Comm. Math. Phys. \textbf{122} (1989), 355--392.
\href{https://doi.org/10.1007/BF01238433}{doi:10.1007/BF01238433}.

\bibitem[FILS78]{FILS1978}
J.~Fr\"ohlich, R.~Israel, E.~H.~Lieb and B.~Simon,
\emph{Phase transitions and reflection positivity.
I. General theory and long range lattice models},
Comm. Math. Phys. \textbf{62} (1978), 1--34.
\href{https://doi.org/10.1007/BF01940327}{doi:10.1007/BF01940327}.

\bibitem[G04]{Grimmett2004}
G.~Grimmett,
\emph{Probability and measure, Hand-out 4:
Fourier transforms: uniqueness and inversion},
University of Cambridge lecture handout (2004), pp.~1--2.
\url{https://www.statslab.cam.ac.uk/~grg/teaching/uniq.pdf}.

\bibitem[KP86]{KoteckyPreiss1986}
R.~Koteck\'y and D.~Preiss,
\emph{Cluster expansion for abstract polymer models},
Comm. Math. Phys. \textbf{103} (1986), 491--498.
\href{https://doi.org/10.1007/BF01211762}{doi:10.1007/BF01211762}.

\bibitem[OS73]{OS1973}
K.~Osterwalder and R.~Schrader,
\emph{Axioms for Euclidean Green's functions},
Comm. Math. Phys. \textbf{31} (1973), 83--112.
\href{https://doi.org/10.1007/BF01645738}{doi:10.1007/BF01645738}.

\bibitem[OS75]{OS1975}
K.~Osterwalder and R.~Schrader,
\emph{Axioms for Euclidean Green's functions II},
with an appendix by S.~Summers,
Comm. Math. Phys. \textbf{42} (1975), 281--305.
\href{https://doi.org/10.1007/BF01608978}{doi:10.1007/BF01608978}.
\bibitem[L77]{Luscher1977}
M.~L\"uscher, \emph{Construction of a selfadjoint, strictly positive
transfer matrix for Euclidean lattice gauge theories}, Commun. Math.
Phys. \textbf{54} (1977), 283--292.
\href{https://doi.org/10.1007/BF01614090}{doi:10.1007/BF01614090}.
Author preprint DESY76/54.

\bibitem[G11]{Georgii2011}
H.-O.~Georgii,
\emph{Gibbs Measures and Phase Transitions}, 2nd extended ed.,
De Gruyter Studies in Mathematics~9, Walter de Gruyter,
Berlin--New York, 2011.
\href{https://doi.org/10.1515/9783110250329}{doi:10.1515/9783110250329}.

\bibitem[YM]{YMCompanion}
\emph{Scalar fields from three-dimensional lattice Yang--Mills theory},
companion manuscript, \S103, ``Symmetries of periodic Gibbs states'',
Lemma~103.2.
\href{https://github.com/ethanwsussman/Yang--Mills-LLM-experiment/blob/1b92bd0347e3b8763cd8ecc1f96dc3c2a92dd310/YM3/YMM1.pdf}{Repository edition}.

\end{thebibliography}
\end{document}
